%% Space Visual Dictionary
\documentclass[12pt,a4paper,twoside]{book}

\usepackage[utf8]{inputenc}
\usepackage[T1]{fontenc}
\usepackage{lmodern}
\usepackage{microtype}

\usepackage[top=2.5cm, bottom=2.5cm, left=2.5cm, right=2.5cm]{geometry}
\usepackage{multicol}
\setlength{\columnsep}{1.2em}
\setlength{\multicolsep}{0pt}

\usepackage{fancyhdr}
\pagestyle{fancy}
\fancyhf{}
\fancyhead[LE,RO]{\leftmark}
\fancyhead[LO,RE]{\thepage}
\setlength{\headheight}{14.5pt}
\renewcommand{\headrulewidth}{0.4pt}

\usepackage{xcolor}
\definecolor{catAD}{RGB}{0,90,150}
\definecolor{catAE}{RGB}{50,120,200}
\definecolor{catRM}{RGB}{180,30,30}
\definecolor{catD}{RGB}{20,140,130}
\definecolor{catSS}{RGB}{120,40,160}

\usepackage[style=altlist,nonumberlist,nopostdot]{glossaries}
\makeglossaries

\usepackage[hidelinks]{hyperref}

\newcommand{\dictentry}[3]{%
  \par\addvspace{3pt}%
  \hangindent=0pt
  \hangafter=0
  \noindent{\color{#3}\textbf{#1}}\hspace{0.5em}{#2}\par
}

\newcommand{\booktitle}{Space Visual Dictionary}
\title{\Huge \booktitle}
\author{}
\date{}

% Glossary Entries for Index


\newglossaryentry{Abort}{ sort={Abort},
  name={Abort},
  description={\textcolor{catAE}{\textbf{[Aerospace]}}},
}
\newglossaryentry{Aerobraking}{ sort={Aerobraking},
  name={Aerobraking},
  description={\textcolor{catAE}{\textbf{[Aerospace]}}},
}
\newglossaryentry{Ejection-Seat}{ sort={Ejection Seat},
  name={Ejection Seat},
  description={\textcolor{catAE}{\textbf{[Aerospace]}}},
}
\newglossaryentry{Instrument-Landing-System}{ sort={Instrument Landing System},
  name={Instrument Landing System},
  description={\textcolor{catAE}{\textbf{[Aerospace]}}},
}
\newglossaryentry{JATO}{ sort={JATO},
  name={JATO},
  description={\textcolor{catAE}{\textbf{[Aerospace]}}},
}
\newglossaryentry{Parachute}{ sort={Parachute},
  name={Parachute},
  description={\textcolor{catAE}{\textbf{[Aerospace]}}},
}
\newglossaryentry{Traffic-Collision-Avoidance-System}{ sort={Traffic Collision Avoidance System},
  name={Traffic Collision Avoidance System},
  description={\textcolor{catAE}{\textbf{[Aerospace]}}},
}
\newglossaryentry{Terrain-Awareness-and-Warning-System}{ sort={Terrain Awareness and Warning System},
  name={Terrain Awareness and Warning System},
  description={\textcolor{catAE}{\textbf{[Aerospace]}}},
}
\newglossaryentry{Barometric-Altimeter}{ sort={Barometric Altimeter},
  name={Barometric Altimeter},
  description={\textcolor{catAE}{\textbf{[Aerospace]}}},
}
\newglossaryentry{Capstan-Winch}{ sort={Capstan Winch},
  name={Capstan Winch},
  description={\textcolor{catAE}{\textbf{[Aerospace]}}},
}
\newglossaryentry{Cockpit-Voice-Recorder}{ sort={Cockpit Voice Recorder},
  name={Cockpit Voice Recorder},
  description={\textcolor{catAE}{\textbf{[Aerospace]}}},
}
\newglossaryentry{Directional-Gyroscope}{ sort={Directional Gyroscope},
  name={Directional Gyroscope},
  description={\textcolor{catAE}{\textbf{[Aerospace]}}},
}
\newglossaryentry{Electronic-Flight-Instrument-System}{ sort={Electronic Flight Instrument System},
  name={Electronic Flight Instrument System},
  description={\textcolor{catAE}{\textbf{[Aerospace]}}},
}
\newglossaryentry{Flight-Data-Recorder}{ sort={Flight Data Recorder},
  name={Flight Data Recorder},
  description={\textcolor{catAE}{\textbf{[Aerospace]}}},
}
\newglossaryentry{Flight-Envelope}{ sort={Flight Envelope},
  name={Flight Envelope},
  description={\textcolor{catAE}{\textbf{[Aerospace]}}},
}
\newglossaryentry{Flight-Level}{ sort={Flight Level},
  name={Flight Level},
  description={\textcolor{catAE}{\textbf{[Aerospace]}}},
}
\newglossaryentry{Gyroscopic-Precession}{ sort={Gyroscopic Precession},
  name={Gyroscopic Precession},
  description={\textcolor{catAE}{\textbf{[Aerospace]}}},
}
\newglossaryentry{Heading-Indicator}{ sort={Heading Indicator},
  name={Heading Indicator},
  description={\textcolor{catAE}{\textbf{[Aerospace]}}},
}
\newglossaryentry{Horizontal-Situation-Indicator}{ sort={Horizontal Situation Indicator},
  name={Horizontal Situation Indicator},
  description={\textcolor{catAE}{\textbf{[Aerospace]}}},
}
\newglossaryentry{Machmeter}{ sort={Machmeter},
  name={Machmeter},
  description={\textcolor{catAE}{\textbf{[Aerospace]}}},
}
\newglossaryentry{Maneuvering-Speed}{ sort={Maneuvering Speed},
  name={Maneuvering Speed},
  description={\textcolor{catAE}{\textbf{[Aerospace]}}},
}
\newglossaryentry{Maximum-Certified-Altitude}{ sort={Maximum Certified Altitude},
  name={Maximum Certified Altitude},
  description={\textcolor{catAE}{\textbf{[Aerospace]}}},
}
\newglossaryentry{Maximum-Landing-Weight}{ sort={Maximum Landing Weight},
  name={Maximum Landing Weight},
  description={\textcolor{catAE}{\textbf{[Aerospace]}}},
}
\newglossaryentry{Maximum-Takeoff-Weight}{ sort={Maximum Takeoff Weight},
  name={Maximum Takeoff Weight},
  description={\textcolor{catAE}{\textbf{[Aerospace]}}},
}
\newglossaryentry{Minimum-Control-Speed}{ sort={Minimum Control Speed},
  name={Minimum Control Speed},
  description={\textcolor{catAE}{\textbf{[Aerospace]}}},
}
\newglossaryentry{Operating-Limit-Speed}{ sort={Operating Limit Speed},
  name={Operating Limit Speed},
  description={\textcolor{catAE}{\textbf{[Aerospace]}}},
}
\newglossaryentry{Radius-of-Action}{ sort={Radius of Action},
  name={Radius of Action},
  description={\textcolor{catAE}{\textbf{[Aerospace]}}},
}
\newglossaryentry{Rate-of-Climb}{ sort={Rate of Climb},
  name={Rate of Climb},
  description={\textcolor{catAE}{\textbf{[Aerospace]}}},
}
\newglossaryentry{Vne}{ sort={Vne},
  name={Vne},
  description={\textcolor{catAE}{\textbf{[Aerospace]}}},
}
\newglossaryentry{Vno}{ sort={Vno},
  name={Vno},
  description={\textcolor{catAE}{\textbf{[Aerospace]}}},
}
\newglossaryentry{Vfe}{ sort={Vfe},
  name={Vfe},
  description={\textcolor{catAE}{\textbf{[Aerospace]}}},
}
\newglossaryentry{Vle}{ sort={Vle},
  name={Vle},
  description={\textcolor{catAE}{\textbf{[Aerospace]}}},
}
\newglossaryentry{Vmca}{ sort={Vmca},
  name={Vmca},
  description={\textcolor{catAE}{\textbf{[Aerospace]}}},
}
\newglossaryentry{Weather-Radar}{ sort={Weather Radar},
  name={Weather Radar},
  description={\textcolor{catAE}{\textbf{[Aerospace]}}},
}
\newglossaryentry{Attitude-Heading-Reference-System}{ sort={Attitude Heading Reference System},
  name={Attitude Heading Reference System},
  description={\textcolor{catAE}{\textbf{[Aerospace]}}},
}
\newglossaryentry{Automatic-Direction-Finder}{ sort={Automatic Direction Finder},
  name={Automatic Direction Finder},
  description={\textcolor{catAE}{\textbf{[Aerospace]}}},
}
\newglossaryentry{Control-Loading}{ sort={Control Loading},
  name={Control Loading},
  description={\textcolor{catAE}{\textbf{[Aerospace]}}},
}
\newglossaryentry{Ditching}{ sort={Ditching},
  name={Ditching},
  description={\textcolor{catAE}{\textbf{[Aerospace]}}},
}
\newglossaryentry{Emergency-Locator-Transmitter}{ sort={Emergency Locator Transmitter},
  name={Emergency Locator Transmitter},
  description={\textcolor{catAE}{\textbf{[Aerospace]}}},
}
\newglossaryentry{Engine-Indicating-and-Crew-Alerting-System}{ sort={Engine Indicating and Crew Alerting System},
  name={Engine Indicating and Crew Alerting System},
  description={\textcolor{catAE}{\textbf{[Aerospace]}}},
}
\newglossaryentry{Fuel-Management}{ sort={Fuel Management},
  name={Fuel Management},
  description={\textcolor{catAE}{\textbf{[Aerospace]}}},
}
\newglossaryentry{Glide-Slope}{ sort={Glide Slope},
  name={Glide Slope},
  description={\textcolor{catAE}{\textbf{[Aerospace]}}},
}
\newglossaryentry{Ground-Proximity-Warning-System}{ sort={Ground Proximity Warning System},
  name={Ground Proximity Warning System},
  description={\textcolor{catAE}{\textbf{[Aerospace]}}},
}
\newglossaryentry{Landing-Distance-Available}{ sort={Landing Distance Available},
  name={Landing Distance Available},
  description={\textcolor{catAE}{\textbf{[Aerospace]}}},
}
\newglossaryentry{Load-Manifest}{ sort={Load Manifest},
  name={Load Manifest},
  description={\textcolor{catAE}{\textbf{[Aerospace]}}},
}
\newglossaryentry{Navigation-Fix}{ sort={Navigation Fix},
  name={Navigation Fix},
  description={\textcolor{catAE}{\textbf{[Aerospace]}}},
}
\newglossaryentry{Obstacle-Clearance}{ sort={Obstacle Clearance},
  name={Obstacle Clearance},
  description={\textcolor{catAE}{\textbf{[Aerospace]}}},
}
\newglossaryentry{Runway-Visual-Range}{ sort={Runway Visual Range},
  name={Runway Visual Range},
  description={\textcolor{catAE}{\textbf{[Aerospace]}}},
}
\newglossaryentry{Stall-Warning}{ sort={Stall Warning},
  name={Stall Warning},
  description={\textcolor{catAE}{\textbf{[Aerospace]}}},
}
\newglossaryentry{Taxiway}{ sort={Taxiway},
  name={Taxiway},
  description={\textcolor{catAE}{\textbf{[Aerospace]}}},
}
\newglossaryentry{Vertical-Speed-Indicator}{ sort={Vertical Speed Indicator},
  name={Vertical Speed Indicator},
  description={\textcolor{catAE}{\textbf{[Aerospace]}}},
}
\newglossaryentry{Wind-Component}{ sort={Wind Component},
  name={Wind Component},
  description={\textcolor{catAE}{\textbf{[Aerospace]}}},
}
\newglossaryentry{Crosswind-Component}{ sort={Crosswind Component},
  name={Crosswind Component},
  description={\textcolor{catAE}{\textbf{[Aerospace]}}},
}
\newglossaryentry{Go-Around}{ sort={Go-Around},
  name={Go-Around},
  description={\textcolor{catAE}{\textbf{[Aerospace]}}},
}
\newglossaryentry{Missed-Approach}{ sort={Missed Approach},
  name={Missed Approach},
  description={\textcolor{catAE}{\textbf{[Aerospace]}}},
}
\newglossaryentry{Holding-Pattern}{ sort={Holding Pattern},
  name={Holding Pattern},
  description={\textcolor{catAE}{\textbf{[Aerospace]}}},
}
\newglossaryentry{Traffic-Pattern}{ sort={Traffic Pattern},
  name={Traffic Pattern},
  description={\textcolor{catAE}{\textbf{[Aerospace]}}},
}
\newglossaryentry{Final-Approach}{ sort={Final Approach},
  name={Final Approach},
  description={\textcolor{catAE}{\textbf{[Aerospace]}}},
}
\newglossaryentry{Base-Leg}{ sort={Base Leg},
  name={Base Leg},
  description={\textcolor{catAE}{\textbf{[Aerospace]}}},
}
\newglossaryentry{Downwind-Leg}{ sort={Downwind Leg},
  name={Downwind Leg},
  description={\textcolor{catAE}{\textbf{[Aerospace]}}},
}
\newglossaryentry{Straight-In-Approach}{ sort={Straight-In Approach},
  name={Straight-In Approach},
  description={\textcolor{catAE}{\textbf{[Aerospace]}}},
}
\newglossaryentry{Circling-Approach}{ sort={Circling Approach},
  name={Circling Approach},
  description={\textcolor{catAE}{\textbf{[Aerospace]}}},
}
\newglossaryentry{Visual-Approach}{ sort={Visual Approach},
  name={Visual Approach},
  description={\textcolor{catAE}{\textbf{[Aerospace]}}},
}
\newglossaryentry{Instrument-Approach}{ sort={Instrument Approach},
  name={Instrument Approach},
  description={\textcolor{catAE}{\textbf{[Aerospace]}}},
}
\newglossaryentry{Precision-Approach}{ sort={Precision Approach},
  name={Precision Approach},
  description={\textcolor{catAE}{\textbf{[Aerospace]}}},
}
\newglossaryentry{Non-Precision-Approach}{ sort={Non-Precision Approach},
  name={Non-Precision Approach},
  description={\textcolor{catAE}{\textbf{[Aerospace]}}},
}
\newglossaryentry{Decision-Altitude}{ sort={Decision Altitude},
  name={Decision Altitude},
  description={\textcolor{catAE}{\textbf{[Aerospace]}}},
}
\newglossaryentry{Minimum-Descent-Altitude}{ sort={Minimum Descent Altitude},
  name={Minimum Descent Altitude},
  description={\textcolor{catAE}{\textbf{[Aerospace]}}},
}
\newglossaryentry{Localizer}{ sort={Localizer},
  name={Localizer},
  description={\textcolor{catAE}{\textbf{[Aerospace]}}},
}
\newglossaryentry{Glideslope-Antenna}{ sort={Glideslope Antenna},
  name={Glideslope Antenna},
  description={\textcolor{catAE}{\textbf{[Aerospace]}}},
}
\newglossaryentry{Marker-Beacon}{ sort={Marker Beacon},
  name={Marker Beacon},
  description={\textcolor{catAE}{\textbf{[Aerospace]}}},
}
\newglossaryentry{Distance-Measuring-Equipment}{ sort={Distance Measuring Equipment},
  name={Distance Measuring Equipment},
  description={\textcolor{catAE}{\textbf{[Aerospace]}}},
}
\newglossaryentry{VOR}{ sort={VOR},
  name={VOR},
  description={\textcolor{catAE}{\textbf{[Aerospace]}}},
}
\newglossaryentry{VORTAC}{ sort={VORTAC},
  name={VORTAC},
  description={\textcolor{catAE}{\textbf{[Aerospace]}}},
}
\newglossaryentry{Non-Directional-Beacon}{ sort={Non-Directional Beacon},
  name={Non-Directional Beacon},
  description={\textcolor{catAE}{\textbf{[Aerospace]}}},
}
\newglossaryentry{ADF}{ sort={ADF},
  name={ADF},
  description={\textcolor{catAE}{\textbf{[Aerospace]}}},
}
\newglossaryentry{Traffic-Information-Service}{ sort={Traffic Information Service},
  name={Traffic Information Service},
  description={\textcolor{catAE}{\textbf{[Aerospace]}}},
}
\newglossaryentry{Safety-Management-System}{ sort={Safety Management System},
  name={Safety Management System},
  description={\textcolor{catAE}{\textbf{[Aerospace]}}},
}
\newglossaryentry{Airworthiness-Directive}{ sort={Airworthiness Directive},
  name={Airworthiness Directive},
  description={\textcolor{catAE}{\textbf{[Aerospace]}}},
}
\newglossaryentry{Production-Certificate}{ sort={Production Certificate},
  name={Production Certificate},
  description={\textcolor{catAE}{\textbf{[Aerospace]}}},
}
\newglossaryentry{Deicing-System}{ sort={Deicing System},
  name={Deicing System},
  description={\textcolor{catAE}{\textbf{[Aerospace]}}},
}
\newglossaryentry{Anti-Icing-System}{ sort={Anti-Icing System},
  name={Anti-Icing System},
  description={\textcolor{catAE}{\textbf{[Aerospace]}}},
}
\newglossaryentry{Fire-Suppression}{ sort={Fire Suppression},
  name={Fire Suppression},
  description={\textcolor{catAE}{\textbf{[Aerospace]}}},
}
\newglossaryentry{Hydraulic-Pump}{ sort={Hydraulic Pump},
  name={Hydraulic Pump},
  description={\textcolor{catAE}{\textbf{[Aerospace]}}},
}
\newglossaryentry{Thermal-Anti-Ice}{ sort={Thermal Anti-Ice},
  name={Thermal Anti-Ice},
  description={\textcolor{catAE}{\textbf{[Aerospace]}}},
}
\newglossaryentry{Windshield-Heating}{ sort={Windshield Heating},
  name={Windshield Heating},
  description={\textcolor{catAE}{\textbf{[Aerospace]}}},
}
\newglossaryentry{Ground-Handling}{ sort={Ground Handling},
  name={Ground Handling},
  description={\textcolor{catAE}{\textbf{[Aerospace]}}},
}
\newglossaryentry{Pushback}{ sort={Pushback},
  name={Pushback},
  description={\textcolor{catAE}{\textbf{[Aerospace]}}},
}
\newglossaryentry{Turnaround}{ sort={Turnaround},
  name={Turnaround},
  description={\textcolor{catAE}{\textbf{[Aerospace]}}},
}
\newglossaryentry{Line-Maintenance}{ sort={Line Maintenance},
  name={Line Maintenance},
  description={\textcolor{catAE}{\textbf{[Aerospace]}}},
}
\newglossaryentry{Heavy-Maintenance}{ sort={Heavy Maintenance},
  name={Heavy Maintenance},
  description={\textcolor{catAE}{\textbf{[Aerospace]}}},
}
\newglossaryentry{Letter-Check}{ sort={Letter Check},
  name={Letter Check},
  description={\textcolor{catAE}{\textbf{[Aerospace]}}},
}
\newglossaryentry{A-Check}{ sort={A-Check},
  name={A-Check},
  description={\textcolor{catAE}{\textbf{[Aerospace]}}},
}
\newglossaryentry{C-Check}{ sort={C-Check},
  name={C-Check},
  description={\textcolor{catAE}{\textbf{[Aerospace]}}},
}
\newglossaryentry{D-Check}{ sort={D-Check},
  name={D-Check},
  description={\textcolor{catAE}{\textbf{[Aerospace]}}},
}
\newglossaryentry{Airframe}{ sort={Airframe},
  name={Airframe},
  description={\textcolor{catAE}{\textbf{[Aerospace]}}},
}
\newglossaryentry{Fuselage-Station}{ sort={Fuselage Station},
  name={Fuselage Station},
  description={\textcolor{catAE}{\textbf{[Aerospace]}}},
}
\newglossaryentry{Waterline}{ sort={Waterline},
  name={Waterline},
  description={\textcolor{catAE}{\textbf{[Aerospace]}}},
}
\newglossaryentry{Butt-Line}{ sort={Butt Line},
  name={Butt Line},
  description={\textcolor{catAE}{\textbf{[Aerospace]}}},
}
\newglossaryentry{Aircraft-Classification-Number}{ sort={Aircraft Classification Number},
  name={Aircraft Classification Number},
  description={\textcolor{catAE}{\textbf{[Aerospace]}}},
}
\newglossaryentry{Pavement-Classification-Number}{ sort={Pavement Classification Number},
  name={Pavement Classification Number},
  description={\textcolor{catAE}{\textbf{[Aerospace]}}},
}
\newglossaryentry{Runway-Strength}{ sort={Runway Strength},
  name={Runway Strength},
  description={\textcolor{catAE}{\textbf{[Aerospace]}}},
}
\newglossaryentry{FOD}{ sort={FOD},
  name={FOD},
  description={\textcolor{catAE}{\textbf{[Aerospace]}}},
}
\newglossaryentry{RAM}{ sort={RAM},
  name={RAM},
  description={\textcolor{catAE}{\textbf{[Aerospace]}}},
}
\newglossaryentry{Windshear}{ sort={Windshear},
  name={Windshear},
  description={\textcolor{catAE}{\textbf{[Aerospace]}}},
}
\newglossaryentry{Microburst}{ sort={Microburst},
  name={Microburst},
  description={\textcolor{catAE}{\textbf{[Aerospace]}}},
}
\newglossaryentry{Clear-Air-Turbulence}{ sort={Clear Air Turbulence},
  name={Clear Air Turbulence},
  description={\textcolor{catAE}{\textbf{[Aerospace]}}},
}
\newglossaryentry{Lee-Wave}{ sort={Lee Wave},
  name={Lee Wave},
  description={\textcolor{catAE}{\textbf{[Aerospace]}}},
}
\newglossaryentry{Mountain-Wave}{ sort={Mountain Wave},
  name={Mountain Wave},
  description={\textcolor{catAE}{\textbf{[Aerospace]}}},
}
\newglossaryentry{Wake-Turbulence}{ sort={Wake Turbulence},
  name={Wake Turbulence},
  description={\textcolor{catAE}{\textbf{[Aerospace]}}},
}
\newglossaryentry{Jet-Blast}{ sort={Jet Blast},
  name={Jet Blast},
  description={\textcolor{catAE}{\textbf{[Aerospace]}}},
}
\newglossaryentry{Propwash}{ sort={Propwash},
  name={Propwash},
  description={\textcolor{catAE}{\textbf{[Aerospace]}}},
}
\newglossaryentry{Slipstream-Effect}{ sort={Slipstream Effect},
  name={Slipstream Effect},
  description={\textcolor{catAE}{\textbf{[Aerospace]}}},
}
\newglossaryentry{Asymmetric-Thrust}{ sort={Asymmetric Thrust},
  name={Asymmetric Thrust},
  description={\textcolor{catAE}{\textbf{[Aerospace]}}},
}
\newglossaryentry{Critical-Engine}{ sort={Critical Engine},
  name={Critical Engine},
  description={\textcolor{catAE}{\textbf{[Aerospace]}}},
}
\newglossaryentry{Vmcg}{ sort={Vmcg},
  name={Vmcg},
  description={\textcolor{catAE}{\textbf{[Aerospace]}}},
}
\newglossaryentry{Balanced-Field-Length}{ sort={Balanced Field Length},
  name={Balanced Field Length},
  description={\textcolor{catAE}{\textbf{[Aerospace]}}},
}
\newglossaryentry{Takeoff-Distance-Required}{ sort={Takeoff Distance Required},
  name={Takeoff Distance Required},
  description={\textcolor{catAE}{\textbf{[Aerospace]}}},
}
\newglossaryentry{Accelerate-Stop-Distance}{ sort={Accelerate-Stop Distance},
  name={Accelerate-Stop Distance},
  description={\textcolor{catAE}{\textbf{[Aerospace]}}},
}
\newglossaryentry{Climb-Gradient}{ sort={Climb Gradient},
  name={Climb Gradient},
  description={\textcolor{catAE}{\textbf{[Aerospace]}}},
}
\newglossaryentry{Biplane}{ sort={Biplane},
  name={Biplane},
  description={\textcolor{catAD}{\textbf{[Aerodynamics]}}},
}
\newglossaryentry{Coanda-Effect}{ sort={Coanda Effect},
  name={Coanda Effect},
  description={\textcolor{catAD}{\textbf{[Aerodynamics]}}},
}
\newglossaryentry{Fineness-Ratio}{ sort={Fineness Ratio},
  name={Fineness Ratio},
  description={\textcolor{catAD}{\textbf{[Aerodynamics]}}},
}
\newglossaryentry{Lift-to-Drag-Ratio}{ sort={Lift-to-Drag Ratio},
  name={Lift-to-Drag Ratio},
  description={\textcolor{catAD}{\textbf{[Aerodynamics]}}},
}
\newglossaryentry{Moment-Coefficient}{ sort={Moment Coefficient},
  name={Moment Coefficient},
  description={\textcolor{catAD}{\textbf{[Aerodynamics]}}},
}
\newglossaryentry{Shock-Boundary-Layer-Interaction}{ sort={Shock-Boundary Layer Interaction},
  name={Shock-Boundary Layer Interaction},
  description={\textcolor{catAD}{\textbf{[Aerodynamics]}}},
}
\newglossaryentry{Stability-Derivative}{ sort={Stability Derivative},
  name={Stability Derivative},
  description={\textcolor{catAD}{\textbf{[Aerodynamics]}}},
}
\newglossaryentry{Aerodynamic-Balancing}{ sort={Aerodynamic Balancing},
  name={Aerodynamic Balancing},
  description={\textcolor{catAD}{\textbf{[Aerodynamics]}}},
}
\newglossaryentry{Aileron-Trim}{ sort={Aileron Trim},
  name={Aileron Trim},
  description={\textcolor{catAD}{\textbf{[Aerodynamics]}}},
}
\newglossaryentry{Body-Force}{ sort={Body Force},
  name={Body Force},
  description={\textcolor{catAD}{\textbf{[Aerodynamics]}}},
}
\newglossaryentry{Boundary-Layer-Suction}{ sort={Boundary Layer Suction},
  name={Boundary Layer Suction},
  description={\textcolor{catAD}{\textbf{[Aerodynamics]}}},
}
\newglossaryentry{Chordwise-Distribution}{ sort={Chordwise Distribution},
  name={Chordwise Distribution},
  description={\textcolor{catAD}{\textbf{[Aerodynamics]}}},
}
\newglossaryentry{Control-Volume}{ sort={Control Volume},
  name={Control Volume},
  description={\textcolor{catAD}{\textbf{[Aerodynamics]}}},
}
\newglossaryentry{Convergent-Divergent-Duct}{ sort={Convergent-Divergent Duct},
  name={Convergent-Divergent Duct},
  description={\textcolor{catAD}{\textbf{[Aerodynamics]}}},
}
\newglossaryentry{Cruise-Efficiency}{ sort={Cruise Efficiency},
  name={Cruise Efficiency},
  description={\textcolor{catAD}{\textbf{[Aerodynamics]}}},
}
\newglossaryentry{Direct-Simulation-Monte-Carlo}{ sort={Direct Simulation Monte Carlo},
  name={Direct Simulation Monte Carlo},
  description={\textcolor{catAD}{\textbf{[Aerodynamics]}}},
}
\newglossaryentry{Drag-Bucket}{ sort={Drag Bucket},
  name={Drag Bucket},
  description={\textcolor{catAD}{\textbf{[Aerodynamics]}}},
}
\newglossaryentry{Drag-Divergence}{ sort={Drag Divergence},
  name={Drag Divergence},
  description={\textcolor{catAD}{\textbf{[Aerodynamics]}}},
}
\newglossaryentry{Effective-Aspect-Ratio}{ sort={Effective Aspect Ratio},
  name={Effective Aspect Ratio},
  description={\textcolor{catAD}{\textbf{[Aerodynamics]}}},
}
\newglossaryentry{Elastic-Axis}{ sort={Elastic Axis},
  name={Elastic Axis},
  description={\textcolor{catAD}{\textbf{[Aerodynamics]}}},
}
\newglossaryentry{Equivalent-Body-of-Revolution}{ sort={Equivalent Body of Revolution},
  name={Equivalent Body of Revolution},
  description={\textcolor{catAD}{\textbf{[Aerodynamics]}}},
}
\newglossaryentry{Fin-Flutter}{ sort={Fin Flutter},
  name={Fin Flutter},
  description={\textcolor{catAD}{\textbf{[Aerodynamics]}}},
}
\newglossaryentry{Flow-visualization}{ sort={Flow visualization},
  name={Flow visualization},
  description={\textcolor{catAD}{\textbf{[Aerodynamics]}}},
}
\newglossaryentry{Freestream}{ sort={Freestream},
  name={Freestream},
  description={\textcolor{catAD}{\textbf{[Aerodynamics]}}},
}
\newglossaryentry{Hypersonic-Boundary-Layer}{ sort={Hypersonic Boundary Layer},
  name={Hypersonic Boundary Layer},
  description={\textcolor{catAD}{\textbf{[Aerodynamics]}}},
}
\newglossaryentry{Induced-Drag-Factor}{ sort={Induced Drag Factor},
  name={Induced Drag Factor},
  description={\textcolor{catAD}{\textbf{[Aerodynamics]}}},
}
\newglossaryentry{Interference-Coefficient}{ sort={Interference Coefficient},
  name={Interference Coefficient},
  description={\textcolor{catAD}{\textbf{[Aerodynamics]}}},
}
\newglossaryentry{Isentropic-Flow}{ sort={Isentropic Flow},
  name={Isentropic Flow},
  description={\textcolor{catAD}{\textbf{[Aerodynamics]}}},
}
\newglossaryentry{Laminar-Bucket}{ sort={Laminar Bucket},
  name={Laminar Bucket},
  description={\textcolor{catAD}{\textbf{[Aerodynamics]}}},
}
\newglossaryentry{Laminar-Separation-Bubble}{ sort={Laminar Separation Bubble},
  name={Laminar Separation Bubble},
  description={\textcolor{catAD}{\textbf{[Aerodynamics]}}},
}
\newglossaryentry{Lifting-Body}{ sort={Lifting Body},
  name={Lifting Body},
  description={\textcolor{catAD}{\textbf{[Aerodynamics]}}},
}
\newglossaryentry{Mach-Disk}{ sort={Mach Disk},
  name={Mach Disk},
  description={\textcolor{catAD}{\textbf{[Aerodynamics]}}},
}
\newglossaryentry{Mangler-Transform}{ sort={Mangler Transform},
  name={Mangler Transform},
  description={\textcolor{catAD}{\textbf{[Aerodynamics]}}},
}
\newglossaryentry{Mass-Flow-Ratio}{ sort={Mass Flow Ratio},
  name={Mass Flow Ratio},
  description={\textcolor{catAD}{\textbf{[Aerodynamics]}}},
}
\newglossaryentry{Mean-Camber-Line}{ sort={Mean Camber Line},
  name={Mean Camber Line},
  description={\textcolor{catAD}{\textbf{[Aerodynamics]}}},
}
\newglossaryentry{Minimum-Drag}{ sort={Minimum Drag},
  name={Minimum Drag},
  description={\textcolor{catAD}{\textbf{[Aerodynamics]}}},
}
\newglossaryentry{Mixed-Boundary-Layer}{ sort={Mixed Boundary Layer},
  name={Mixed Boundary Layer},
  description={\textcolor{catAD}{\textbf{[Aerodynamics]}}},
}
\newglossaryentry{Net-Thrust}{ sort={Net Thrust},
  name={Net Thrust},
  description={\textcolor{catAD}{\textbf{[Aerodynamics]}}},
}
\newglossaryentry{Non-Inertial-Reference-Frame}{ sort={Non-Inertial Reference Frame},
  name={Non-Inertial Reference Frame},
  description={\textcolor{catAD}{\textbf{[Aerodynamics]}}},
}
\newglossaryentry{Normal-Force}{ sort={Normal Force},
  name={Normal Force},
  description={\textcolor{catAD}{\textbf{[Aerodynamics]}}},
}
\newglossaryentry{Oblique-Flow}{ sort={Oblique Flow},
  name={Oblique Flow},
  description={\textcolor{catAD}{\textbf{[Aerodynamics]}}},
}
\newglossaryentry{Open-Jet-Tunnel}{ sort={Open Jet Tunnel},
  name={Open Jet Tunnel},
  description={\textcolor{catAD}{\textbf{[Aerodynamics]}}},
}
\newglossaryentry{Pitching-Moment-Coefficient}{ sort={Pitching Moment Coefficient},
  name={Pitching Moment Coefficient},
  description={\textcolor{catAD}{\textbf{[Aerodynamics]}}},
}
\newglossaryentry{Planform}{ sort={Planform},
  name={Planform},
  description={\textcolor{catAD}{\textbf{[Aerodynamics]}}},
}
\newglossaryentry{Pressure-Recovery}{ sort={Pressure Recovery},
  name={Pressure Recovery},
  description={\textcolor{catAD}{\textbf{[Aerodynamics]}}},
}
\newglossaryentry{Pressure-Surface}{ sort={Pressure Surface},
  name={Pressure Surface},
  description={\textcolor{catAD}{\textbf{[Aerodynamics]}}},
}
\newglossaryentry{Ramp-Angle}{ sort={Ramp Angle},
  name={Ramp Angle},
  description={\textcolor{catAD}{\textbf{[Aerodynamics]}}},
}
\newglossaryentry{Rankine-Hugoniot-Relations}{ sort={Rankine-Hugoniot Relations},
  name={Rankine-Hugoniot Relations},
  description={\textcolor{catAD}{\textbf{[Aerodynamics]}}},
}
\newglossaryentry{Reduced-Frequency}{ sort={Reduced Frequency},
  name={Reduced Frequency},
  description={\textcolor{catAD}{\textbf{[Aerodynamics]}}},
}
\newglossaryentry{Reference-Area}{ sort={Reference Area},
  name={Reference Area},
  description={\textcolor{catAD}{\textbf{[Aerodynamics]}}},
}
\newglossaryentry{Reynolds-Number-Scaling}{ sort={Reynolds Number Scaling},
  name={Reynolds Number Scaling},
  description={\textcolor{catAD}{\textbf{[Aerodynamics]}}},
}
\newglossaryentry{Roll-Rate}{ sort={Roll Rate},
  name={Roll Rate},
  description={\textcolor{catAD}{\textbf{[Aerodynamics]}}},
}
\newglossaryentry{Roughness-Element}{ sort={Roughness Element},
  name={Roughness Element},
  description={\textcolor{catAD}{\textbf{[Aerodynamics]}}},
}
\newglossaryentry{Schlieren-System}{ sort={Schlieren System},
  name={Schlieren System},
  description={\textcolor{catAD}{\textbf{[Aerodynamics]}}},
}
\newglossaryentry{Scoop}{ sort={Scoop},
  name={Scoop},
  description={\textcolor{catAD}{\textbf{[Aerodynamics]}}},
}
\newglossaryentry{Section-Drag-Coefficient}{ sort={Section Drag Coefficient},
  name={Section Drag Coefficient},
  description={\textcolor{catAD}{\textbf{[Aerodynamics]}}},
}
\newglossaryentry{Section-Lift-Coefficient}{ sort={Section Lift Coefficient},
  name={Section Lift Coefficient},
  description={\textcolor{catAD}{\textbf{[Aerodynamics]}}},
}
\newglossaryentry{Separation-Bubble}{ sort={Separation Bubble},
  name={Separation Bubble},
  description={\textcolor{catAD}{\textbf{[Aerodynamics]}}},
}
\newglossaryentry{Shock-Expansin-Theory}{ sort={Shock Expansin Theory},
  name={Shock Expansin Theory},
  description={\textcolor{catAD}{\textbf{[Aerodynamics]}}},
}
\newglossaryentry{Slat-Track}{ sort={Slat Track},
  name={Slat Track},
  description={\textcolor{catAD}{\textbf{[Aerodynamics]}}},
}
\newglossaryentry{Slipstream}{ sort={Slipstream},
  name={Slipstream},
  description={\textcolor{catAD}{\textbf{[Aerodynamics]}}},
}
\newglossaryentry{Span-Loading}{ sort={Span Loading},
  name={Span Loading},
  description={\textcolor{catAD}{\textbf{[Aerodynamics]}}},
}
\newglossaryentry{Spanwise-Flow}{ sort={Spanwise Flow},
  name={Spanwise Flow},
  description={\textcolor{catAD}{\textbf{[Aerodynamics]}}},
}
\newglossaryentry{Speed-Brake}{ sort={Speed Brake},
  name={Speed Brake},
  description={\textcolor{catAD}{\textbf{[Aerodynamics]}}},
}
\newglossaryentry{Subsonic-Leading-Edge}{ sort={Subsonic Leading Edge},
  name={Subsonic Leading Edge},
  description={\textcolor{catAD}{\textbf{[Aerodynamics]}}},
}
\newglossaryentry{Supersonic-Leading-Edge}{ sort={Supersonic Leading Edge},
  name={Supersonic Leading Edge},
  description={\textcolor{catAD}{\textbf{[Aerodynamics]}}},
}
\newglossaryentry{Supersonic-Area-Rule}{ sort={Supersonic Area Rule},
  name={Supersonic Area Rule},
  description={\textcolor{catAD}{\textbf{[Aerodynamics]}}},
}
\newglossaryentry{Surface-Pressure-Distribution}{ sort={Surface Pressure Distribution},
  name={Surface Pressure Distribution},
  description={\textcolor{catAD}{\textbf{[Aerodynamics]}}},
}
\newglossaryentry{Sweep-Effect}{ sort={Sweep Effect},
  name={Sweep Effect},
  description={\textcolor{catAD}{\textbf{[Aerodynamics]}}},
}
\newglossaryentry{Terminal-Velocity}{ sort={Terminal Velocity},
  name={Terminal Velocity},
  description={\textcolor{catAD}{\textbf{[Aerodynamics]}}},
}
\newglossaryentry{Thermal-Boundary-Layer}{ sort={Thermal Boundary Layer},
  name={Thermal Boundary Layer},
  description={\textcolor{catAD}{\textbf{[Aerodynamics]}}},
}
\newglossaryentry{Transonic-Dip}{ sort={Transonic Dip},
  name={Transonic Dip},
  description={\textcolor{catAD}{\textbf{[Aerodynamics]}}},
}
\newglossaryentry{Two-Dimensional-Flow}{ sort={Two-Dimensional Flow},
  name={Two-Dimensional Flow},
  description={\textcolor{catAD}{\textbf{[Aerodynamics]}}},
}
\newglossaryentry{Velocity-Potential}{ sort={Velocity Potential},
  name={Velocity Potential},
  description={\textcolor{catAD}{\textbf{[Aerodynamics]}}},
}
\newglossaryentry{Viscous-Drag}{ sort={Viscous Drag},
  name={Viscous Drag},
  description={\textcolor{catAD}{\textbf{[Aerodynamics]}}},
}
\newglossaryentry{Volume-Coefficient}{ sort={Volume Coefficient},
  name={Volume Coefficient},
  description={\textcolor{catAD}{\textbf{[Aerodynamics]}}},
}
\newglossaryentry{Vortex-Core}{ sort={Vortex Core},
  name={Vortex Core},
  description={\textcolor{catAD}{\textbf{[Aerodynamics]}}},
}
\newglossaryentry{Wind-Shear}{ sort={Wind Shear},
  name={Wind Shear},
  description={\textcolor{catAD}{\textbf{[Aerodynamics]}}},
}
\newglossaryentry{Zero-Lift-Angle}{ sort={Zero-Lift Angle},
  name={Zero-Lift Angle},
  description={\textcolor{catAD}{\textbf{[Aerodynamics]}}},
}
\newglossaryentry{Zero-Lift-Drag}{ sort={Zero-Lift Drag},
  name={Zero-Lift Drag},
  description={\textcolor{catAD}{\textbf{[Aerodynamics]}}},
}
\newglossaryentry{Aileron-Lock}{ sort={Aileron Lock},
  name={Aileron Lock},
  description={\textcolor{catAD}{\textbf{[Aerodynamics]}}},
}
\newglossaryentry{Adverse-Pressure-Gradient}{ sort={Adverse Pressure Gradient},
  name={Adverse Pressure Gradient},
  description={\textcolor{catAD}{\textbf{[Aerodynamics]}}},
}
\newglossaryentry{Boundary-Layer-Thickness}{ sort={Boundary Layer Thickness},
  name={Boundary Layer Thickness},
  description={\textcolor{catAD}{\textbf{[Aerodynamics]}}},
}
\newglossaryentry{Centerline-Pressure}{ sort={Centerline Pressure},
  name={Centerline Pressure},
  description={\textcolor{catAD}{\textbf{[Aerodynamics]}}},
}
\newglossaryentry{Characteristic-Mach-Angle}{ sort={Characteristic Mach Angle},
  name={Characteristic Mach Angle},
  description={\textcolor{catAD}{\textbf{[Aerodynamics]}}},
}
\newglossaryentry{Compressibility-Correction}{ sort={Compressibility Correction},
  name={Compressibility Correction},
  description={\textcolor{catAD}{\textbf{[Aerodynamics]}}},
}
\newglossaryentry{Couette-Flow}{ sort={Couette Flow},
  name={Couette Flow},
  description={\textcolor{catAD}{\textbf{[Aerodynamics]}}},
}
\newglossaryentry{Entrance-Length}{ sort={Entrance Length},
  name={Entrance Length},
  description={\textcolor{catAD}{\textbf{[Aerodynamics]}}},
}
\newglossaryentry{Favorable-Pressure-Gradient}{ sort={Favorable Pressure Gradient},
  name={Favorable Pressure Gradient},
  description={\textcolor{catAD}{\textbf{[Aerodynamics]}}},
}
\newglossaryentry{Hagen-Poiseuille-Flow}{ sort={Hagen-Poiseuille Flow},
  name={Hagen-Poiseuille Flow},
  description={\textcolor{catAD}{\textbf{[Aerodynamics]}}},
}
\newglossaryentry{Hydraulic-Diameter}{ sort={Hydraulic Diameter},
  name={Hydraulic Diameter},
  description={\textcolor{catAD}{\textbf{[Aerodynamics]}}},
}
\newglossaryentry{Kinematic-Viscosity}{ sort={Kinematic Viscosity},
  name={Kinematic Viscosity},
  description={\textcolor{catAD}{\textbf{[Aerodynamics]}}},
}
\newglossaryentry{Laminar-Turbulent-Transition}{ sort={Laminar-Turbulent Transition},
  name={Laminar-Turbulent Transition},
  description={\textcolor{catAD}{\textbf{[Aerodynamics]}}},
}
\newglossaryentry{Linearized-Theory}{ sort={Linearized Theory},
  name={Linearized Theory},
  description={\textcolor{catAD}{\textbf{[Aerodynamics]}}},
}
\newglossaryentry{Local-Mach-Number}{ sort={Local Mach Number},
  name={Local Mach Number},
  description={\textcolor{catAD}{\textbf{[Aerodynamics]}}},
}
\newglossaryentry{No-Slip-Condition}{ sort={No-Slip Condition},
  name={No-Slip Condition},
  description={\textcolor{catAD}{\textbf{[Aerodynamics]}}},
}
\newglossaryentry{Potential-Flow}{ sort={Potential Flow},
  name={Potential Flow},
  description={\textcolor{catAD}{\textbf{[Aerodynamics]}}},
}
\newglossaryentry{Stokes-Flow}{ sort={Stokes Flow},
  name={Stokes Flow},
  description={\textcolor{catAD}{\textbf{[Aerodynamics]}}},
}
\newglossaryentry{Stream-Function}{ sort={Stream Function},
  name={Stream Function},
  description={\textcolor{catAD}{\textbf{[Aerodynamics]}}},
}
\newglossaryentry{Taylor-Couette-Flow}{ sort={Taylor-Couette Flow},
  name={Taylor-Couette Flow},
  description={\textcolor{catAD}{\textbf{[Aerodynamics]}}},
}
\newglossaryentry{Thermal-Conductivity}{ sort={Thermal Conductivity},
  name={Thermal Conductivity},
  description={\textcolor{catAD}{\textbf{[Aerodynamics]}}},
}
\newglossaryentry{Turbulent-Reynolds-Stress}{ sort={Turbulent Reynolds Stress},
  name={Turbulent Reynolds Stress},
  description={\textcolor{catAD}{\textbf{[Aerodynamics]}}},
}
\newglossaryentry{Velocity-Profile}{ sort={Velocity Profile},
  name={Velocity Profile},
  description={\textcolor{catAD}{\textbf{[Aerodynamics]}}},
}
\newglossaryentry{Vorticity}{ sort={Vorticity},
  name={Vorticity},
  description={\textcolor{catAD}{\textbf{[Aerodynamics]}}},
}
\newglossaryentry{Guidance-Computer}{ sort={Guidance Computer},
  name={Guidance Computer},
  description={\textcolor{catRM}{\textbf{[Rockets and Missiles]}}},
}
\newglossaryentry{Launch-Escape-System}{ sort={Launch Escape System},
  name={Launch Escape System},
  description={\textcolor{catRM}{\textbf{[Rockets and Missiles]}}},
}
\newglossaryentry{Abort-System}{ sort={Abort System},
  name={Abort System},
  description={\textcolor{catRM}{\textbf{[Rockets and Missiles]}}},
}
\newglossaryentry{Ablative-Shield}{ sort={Ablative Shield},
  name={Ablative Shield},
  description={\textcolor{catRM}{\textbf{[Rockets and Missiles]}}},
}
\newglossaryentry{Ballistic-Coefficient}{ sort={Ballistic Coefficient},
  name={Ballistic Coefficient},
  description={\textcolor{catRM}{\textbf{[Rockets and Missiles]}}},
}
\newglossaryentry{Castor}{ sort={Castor},
  name={Castor},
  description={\textcolor{catRM}{\textbf{[Rockets and Missiles]}}},
}
\newglossaryentry{Centaur}{ sort={Centaur},
  name={Centaur},
  description={\textcolor{catRM}{\textbf{[Rockets and Missiles]}}},
}
\newglossaryentry{Launch-Clamp}{ sort={Launch Clamp},
  name={Launch Clamp},
  description={\textcolor{catRM}{\textbf{[Rockets and Missiles]}}},
}
\newglossaryentry{Movable-Nozzle}{ sort={Movable Nozzle},
  name={Movable Nozzle},
  description={\textcolor{catRM}{\textbf{[Rockets and Missiles]}}},
}
\newglossaryentry{Retrofire}{ sort={Retrofire},
  name={Retrofire},
  description={\textcolor{catRM}{\textbf{[Rockets and Missiles]}}},
}
\newglossaryentry{Stage}{ sort={Stage},
  name={Stage},
  description={\textcolor{catRM}{\textbf{[Rockets and Missiles]}}},
}
\newglossaryentry{Abort-Motor}{ sort={Abort Motor},
  name={Abort Motor},
  description={\textcolor{catRM}{\textbf{[Rockets and Missiles]}}},
}
\newglossaryentry{Attitude-Control-Motor}{ sort={Attitude Control Motor},
  name={Attitude Control Motor},
  description={\textcolor{catRM}{\textbf{[Rockets and Missiles]}}},
}
\newglossaryentry{Escape-Tower}{ sort={Escape Tower},
  name={Escape Tower},
  description={\textcolor{catRM}{\textbf{[Rockets and Missiles]}}},
}
\newglossaryentry{Pusher-Rocket}{ sort={Pusher Rocket},
  name={Pusher Rocket},
  description={\textcolor{catRM}{\textbf{[Rockets and Missiles]}}},
}
\newglossaryentry{ullage-Motor}{ sort={ullage Motor},
  name={ullage Motor},
  description={\textcolor{catRM}{\textbf{[Rockets and Missiles]}}},
}
\newglossaryentry{ullage}{ sort={ullage},
  name={ullage},
  description={\textcolor{catRM}{\textbf{[Rockets and Missiles]}}},
}
\newglossaryentry{Propellant-Gauging}{ sort={Propellant Gauging},
  name={Propellant Gauging},
  description={\textcolor{catRM}{\textbf{[Rockets and Missiles]}}},
}
\newglossaryentry{Mass-Fraction}{ sort={Mass Fraction},
  name={Mass Fraction},
  description={\textcolor{catRM}{\textbf{[Rockets and Missiles]}}},
}
\newglossaryentry{Nozzle-Efficiency}{ sort={Nozzle Efficiency},
  name={Nozzle Efficiency},
  description={\textcolor{catRM}{\textbf{[Rockets and Missiles]}}},
}
\newglossaryentry{Ignition-Overpressure}{ sort={Ignition Overpressure},
  name={Ignition Overpressure},
  description={\textcolor{catRM}{\textbf{[Rockets and Missiles]}}},
}
\newglossaryentry{Hard-Start}{ sort={Hard Start},
  name={Hard Start},
  description={\textcolor{catRM}{\textbf{[Rockets and Missiles]}}},
}
\newglossaryentry{Two-Combustion-Instability}{ sort={Two-Combustion Instability},
  name={Two-Combustion Instability},
  description={\textcolor{catRM}{\textbf{[Rockets and Missiles]}}},
}
\newglossaryentry{Acoustic-Mode}{ sort={Acoustic Mode},
  name={Acoustic Mode},
  description={\textcolor{catRM}{\textbf{[Rockets and Missiles]}}},
}
\newglossaryentry{Grain-Burn-Rate}{ sort={Grain Burn Rate},
  name={Grain Burn Rate},
  description={\textcolor{catRM}{\textbf{[Rockets and Missiles]}}},
}
\newglossaryentry{Web-Thickness}{ sort={Web Thickness},
  name={Web Thickness},
  description={\textcolor{catRM}{\textbf{[Rockets and Missiles]}}},
}
\newglossaryentry{Burn-Rate-Exponent}{ sort={Burn Rate Exponent},
  name={Burn Rate Exponent},
  description={\textcolor{catRM}{\textbf{[Rockets and Missiles]}}},
}
\newglossaryentry{Neutral-Burn}{ sort={Neutral Burn},
  name={Neutral Burn},
  description={\textcolor{catRM}{\textbf{[Rockets and Missiles]}}},
}
\newglossaryentry{Progressive-Burn}{ sort={Progressive Burn},
  name={Progressive Burn},
  description={\textcolor{catRM}{\textbf{[Rockets and Missiles]}}},
}
\newglossaryentry{Regressive-Burn}{ sort={Regressive Burn},
  name={Regressive Burn},
  description={\textcolor{catRM}{\textbf{[Rockets and Missiles]}}},
}
\newglossaryentry{Sliver}{ sort={Sliver},
  name={Sliver},
  description={\textcolor{catRM}{\textbf{[Rockets and Missiles]}}},
}
\newglossaryentry{Void-Fraction}{ sort={Void Fraction},
  name={Void Fraction},
  description={\textcolor{catRM}{\textbf{[Rockets and Missiles]}}},
}
\newglossaryentry{Free-Volume}{ sort={Free Volume},
  name={Free Volume},
  description={\textcolor{catRM}{\textbf{[Rockets and Missiles]}}},
}
\newglossaryentry{Impinging-Jet}{ sort={Impinging Jet},
  name={Impinging Jet},
  description={\textcolor{catRM}{\textbf{[Rockets and Missiles]}}},
}
\newglossaryentry{Coaxial-Swirl}{ sort={Coaxial Swirl},
  name={Coaxial Swirl},
  description={\textcolor{catRM}{\textbf{[Rockets and Missiles]}}},
}
\newglossaryentry{Turbopump-Assembly}{ sort={Turbopump Assembly},
  name={Turbopump Assembly},
  description={\textcolor{catRM}{\textbf{[Rockets and Missiles]}}},
}
\newglossaryentry{Oxidizer-Turbine}{ sort={Oxidizer Turbine},
  name={Oxidizer Turbine},
  description={\textcolor{catRM}{\textbf{[Rockets and Missiles]}}},
}
\newglossaryentry{Fuel-Turbine}{ sort={Fuel Turbine},
  name={Fuel Turbine},
  description={\textcolor{catRM}{\textbf{[Rockets and Missiles]}}},
}
\newglossaryentry{Staged-Combustion}{ sort={Staged Combustion},
  name={Staged Combustion},
  description={\textcolor{catRM}{\textbf{[Rockets and Missiles]}}},
}
\newglossaryentry{Gas-generator-Cycle}{ sort={Gas-generator Cycle},
  name={Gas-generator Cycle},
  description={\textcolor{catRM}{\textbf{[Rockets and Missiles]}}},
}
\newglossaryentry{Full-Flow-Staged-Combustion}{ sort={Full-Flow Staged Combustion},
  name={Full-Flow Staged Combustion},
  description={\textcolor{catRM}{\textbf{[Rockets and Missiles]}}},
}
\newglossaryentry{Oxidizer-Rich-Staged-Combustion}{ sort={Oxidizer-Rich Staged Combustion},
  name={Oxidizer-Rich Staged Combustion},
  description={\textcolor{catRM}{\textbf{[Rockets and Missiles]}}},
}
\newglossaryentry{Pressure-Fed-Engine}{ sort={Pressure-Fed Engine},
  name={Pressure-Fed Engine},
  description={\textcolor{catRM}{\textbf{[Rockets and Missiles]}}},
}
\newglossaryentry{Monopropellant-Thruster}{ sort={Monopropellant Thruster},
  name={Monopropellant Thruster},
  description={\textcolor{catRM}{\textbf{[Rockets and Missiles]}}},
}
\newglossaryentry{Bipropellant-Thruster}{ sort={Bipropellant Thruster},
  name={Bipropellant Thruster},
  description={\textcolor{catRM}{\textbf{[Rockets and Missiles]}}},
}
\newglossaryentry{Green-Propellant}{ sort={Green Propellant},
  name={Green Propellant},
  description={\textcolor{catRM}{\textbf{[Rockets and Missiles]}}},
}
\newglossaryentry{Ion-Thruster}{ sort={Ion Thruster},
  name={Ion Thruster},
  description={\textcolor{catRM}{\textbf{[Rockets and Missiles]}}},
}
\newglossaryentry{Hall-Effect-Thruster}{ sort={Hall Effect Thruster},
  name={Hall Effect Thruster},
  description={\textcolor{catRM}{\textbf{[Rockets and Missiles]}}},
}
\newglossaryentry{Pulsed-Plasma-Thruster}{ sort={Pulsed Plasma Thruster},
  name={Pulsed Plasma Thruster},
  description={\textcolor{catRM}{\textbf{[Rockets and Missiles]}}},
}
\newglossaryentry{Arcjet}{ sort={Arcjet},
  name={Arcjet},
  description={\textcolor{catRM}{\textbf{[Rockets and Missiles]}}},
}
\newglossaryentry{Resistojet}{ sort={Resistojet},
  name={Resistojet},
  description={\textcolor{catRM}{\textbf{[Rockets and Missiles]}}},
}
\newglossaryentry{Variable-Specific-Impulse-Magnetoplasma-Rocket}{ sort={Variable Specific Impulse Magnetoplasma Rocket},
  name={Variable Specific Impulse Magnetoplasma Rocket},
  description={\textcolor{catRM}{\textbf{[Rockets and Missiles]}}},
}
\newglossaryentry{Laser-Propulsion}{ sort={Laser Propulsion},
  name={Laser Propulsion},
  description={\textcolor{catRM}{\textbf{[Rockets and Missiles]}}},
}
\newglossaryentry{Nuclear-Thermal-Propulsion}{ sort={Nuclear Thermal Propulsion},
  name={Nuclear Thermal Propulsion},
  description={\textcolor{catRM}{\textbf{[Rockets and Missiles]}}},
}
\newglossaryentry{Nuclear-Electric-Propulsion}{ sort={Nuclear Electric Propulsion},
  name={Nuclear Electric Propulsion},
  description={\textcolor{catRM}{\textbf{[Rockets and Missiles]}}},
}
\newglossaryentry{Antimatter-Propulsion}{ sort={Antimatter Propulsion},
  name={Antimatter Propulsion},
  description={\textcolor{catRM}{\textbf{[Rockets and Missiles]}}},
}
\newglossaryentry{Tether-Propulsion}{ sort={Tether Propulsion},
  name={Tether Propulsion},
  description={\textcolor{catRM}{\textbf{[Rockets and Missiles]}}},
}
\newglossaryentry{Mass-Driver}{ sort={Mass Driver},
  name={Mass Driver},
  description={\textcolor{catRM}{\textbf{[Rockets and Missiles]}}},
}
\newglossaryentry{Launch-Loop}{ sort={Launch Loop},
  name={Launch Loop},
  description={\textcolor{catRM}{\textbf{[Rockets and Missiles]}}},
}
\newglossaryentry{Space-Elevator}{ sort={Space Elevator},
  name={Space Elevator},
  description={\textcolor{catRM}{\textbf{[Rockets and Missiles]}}},
}
\newglossaryentry{Orbital-Tug}{ sort={Orbital Tug},
  name={Orbital Tug},
  description={\textcolor{catRM}{\textbf{[Rockets and Missiles]}}},
}
\newglossaryentry{Cislunar-Transport}{ sort={Cislunar Transport},
  name={Cislunar Transport},
  description={\textcolor{catRM}{\textbf{[Rockets and Missiles]}}},
}
\newglossaryentry{Eclipse}{ sort={Eclipse},
  name={Eclipse},
  description={\textcolor{catSS}{\textbf{[Space Science]}}},
}
\newglossaryentry{Ephemeris}{ sort={Ephemeris},
  name={Ephemeris},
  description={\textcolor{catSS}{\textbf{[Space Science]}}},
}
\newglossaryentry{Oberth-Effect}{ sort={Oberth Effect},
  name={Oberth Effect},
  description={\textcolor{catSS}{\textbf{[Space Science]}}},
}
\newglossaryentry{Orbital-Debris}{ sort={Orbital Debris},
  name={Orbital Debris},
  description={\textcolor{catSS}{\textbf{[Space Science]}}},
}
\newglossaryentry{Rendezvous}{ sort={Rendezvous},
  name={Rendezvous},
  description={\textcolor{catSS}{\textbf{[Space Science]}}},
}
\newglossaryentry{Tsiolkovsky-Rocket-Equation}{ sort={Tsiolkovsky Rocket Equation},
  name={Tsiolkovsky Rocket Equation},
  description={\textcolor{catSS}{\textbf{[Space Science]}}},
}
\newglossaryentry{Aberration}{ sort={Aberration},
  name={Aberration},
  description={\textcolor{catSS}{\textbf{[Space Science]}}},
}
\newglossaryentry{Astrobiology}{ sort={Astrobiology},
  name={Astrobiology},
  description={\textcolor{catSS}{\textbf{[Space Science]}}},
}
\newglossaryentry{Asteroid-Mining}{ sort={Asteroid Mining},
  name={Asteroid Mining},
  description={\textcolor{catSS}{\textbf{[Space Science]}}},
}
\newglossaryentry{Atmospheric-Entry}{ sort={Atmospheric Entry},
  name={Atmospheric Entry},
  description={\textcolor{catSS}{\textbf{[Space Science]}}},
}
\newglossaryentry{Colony}{ sort={Colony},
  name={Colony},
  description={\textcolor{catSS}{\textbf{[Space Science]}}},
}
\newglossaryentry{Conjunction}{ sort={Conjunction},
  name={Conjunction},
  description={\textcolor{catSS}{\textbf{[Space Science]}}},
}
\newglossaryentry{Delta-V-Budget}{ sort={Delta-V Budget},
  name={Delta-V Budget},
  description={\textcolor{catSS}{\textbf{[Space Science]}}},
}
\newglossaryentry{Deorbit}{ sort={Deorbit},
  name={Deorbit},
  description={\textcolor{catSS}{\textbf{[Space Science]}}},
}
\newglossaryentry{Equatorial-Orbit}{ sort={Equatorial Orbit},
  name={Equatorial Orbit},
  description={\textcolor{catSS}{\textbf{[Space Science]}}},
}
\newglossaryentry{Habitable-Zone}{ sort={Habitable Zone},
  name={Habitable Zone},
  description={\textcolor{catSS}{\textbf{[Space Science]}}},
}
\newglossaryentry{Keplerian-Elements}{ sort={Keplerian Elements},
  name={Keplerian Elements},
  description={\textcolor{catSS}{\textbf{[Space Science]}}},
}
\newglossaryentry{Orbital-Plane}{ sort={Orbital Plane},
  name={Orbital Plane},
  description={\textcolor{catSS}{\textbf{[Space Science]}}},
}
\newglossaryentry{Probe}{ sort={Probe},
  name={Probe},
  description={\textcolor{catSS}{\textbf{[Space Science]}}},
}
\newglossaryentry{Radiation-Belt}{ sort={Radiation Belt},
  name={Radiation Belt},
  description={\textcolor{catSS}{\textbf{[Space Science]}}},
}
\newglossaryentry{Station-Keeping}{ sort={Station-Keeping},
  name={Station-Keeping},
  description={\textcolor{catSS}{\textbf{[Space Science]}}},
}
\newglossaryentry{Venus-Flyby}{ sort={Venus Flyby},
  name={Venus Flyby},
  description={\textcolor{catSS}{\textbf{[Space Science]}}},
}
\newglossaryentry{Aerocapture}{ sort={Aerocapture},
  name={Aerocapture},
  description={\textcolor{catSS}{\textbf{[Space Science]}}},
}
\newglossaryentry{Apsis}{ sort={Apsis},
  name={Apsis},
  description={\textcolor{catSS}{\textbf{[Space Science]}}},
}
\newglossaryentry{Argument-of-Periapsis}{ sort={Argument of Periapsis},
  name={Argument of Periapsis},
  description={\textcolor{catSS}{\textbf{[Space Science]}}},
}
\newglossaryentry{Longitude-of-Ascending-Node}{ sort={Longitude of Ascending Node},
  name={Longitude of Ascending Node},
  description={\textcolor{catSS}{\textbf{[Space Science]}}},
}
\newglossaryentry{Gravitational-Parameter}{ sort={Gravitational Parameter},
  name={Gravitational Parameter},
  description={\textcolor{catSS}{\textbf{[Space Science]}}},
}
\newglossaryentry{Two-Body-Problem}{ sort={Two-Body Problem},
  name={Two-Body Problem},
  description={\textcolor{catSS}{\textbf{[Space Science]}}},
}
\newglossaryentry{Restricted-Three-Body-Problem}{ sort={Restricted Three-Body Problem},
  name={Restricted Three-Body Problem},
  description={\textcolor{catSS}{\textbf{[Space Science]}}},
}
\newglossaryentry{Resonance}{ sort={Resonance},
  name={Resonance},
  description={\textcolor{catSS}{\textbf{[Space Science]}}},
}
\newglossaryentry{Secular-Perturbation}{ sort={Secular Perturbation},
  name={Secular Perturbation},
  description={\textcolor{catSS}{\textbf{[Space Science]}}},
}
\newglossaryentry{Periodic-Perturbation}{ sort={Periodic Perturbation},
  name={Periodic Perturbation},
  description={\textcolor{catSS}{\textbf{[Space Science]}}},
}
\newglossaryentry{J2-Effect}{ sort={J2 Effect},
  name={J2 Effect},
  description={\textcolor{catSS}{\textbf{[Space Science]}}},
}
\newglossaryentry{Precession}{ sort={Precession},
  name={Precession},
  description={\textcolor{catSS}{\textbf{[Space Science]}}},
}
\newglossaryentry{Nutation}{ sort={Nutation},
  name={Nutation},
  description={\textcolor{catSS}{\textbf{[Space Science]}}},
}
\newglossaryentry{Libration}{ sort={Libration},
  name={Libration},
  description={\textcolor{catSS}{\textbf{[Space Science]}}},
}
\newglossaryentry{Albedo}{ sort={Albedo},
  name={Albedo},
  description={\textcolor{catSS}{\textbf{[Space Science]}}},
}
\newglossaryentry{Bond-Albedo}{ sort={Bond Albedo},
  name={Bond Albedo},
  description={\textcolor{catSS}{\textbf{[Space Science]}}},
}
\newglossaryentry{Geometric-Albedo}{ sort={Geometric Albedo},
  name={Geometric Albedo},
  description={\textcolor{catSS}{\textbf{[Space Science]}}},
}
\newglossaryentry{Apparent-Magnitude}{ sort={Apparent Magnitude},
  name={Apparent Magnitude},
  description={\textcolor{catSS}{\textbf{[Space Science]}}},
}
\newglossaryentry{Absolute-Magnitude}{ sort={Absolute Magnitude},
  name={Absolute Magnitude},
  description={\textcolor{catSS}{\textbf{[Space Science]}}},
}
\newglossaryentry{Stellar-Parallax}{ sort={Stellar Parallax},
  name={Stellar Parallax},
  description={\textcolor{catSS}{\textbf{[Space Science]}}},
}
\newglossaryentry{Light-Curve}{ sort={Light Curve},
  name={Light Curve},
  description={\textcolor{catSS}{\textbf{[Space Science]}}},
}
\newglossaryentry{Spectroscopy}{ sort={Spectroscopy},
  name={Spectroscopy},
  description={\textcolor{catSS}{\textbf{[Space Science]}}},
}
\newglossaryentry{Blueshift}{ sort={Blueshift},
  name={Blueshift},
  description={\textcolor{catSS}{\textbf{[Space Science]}}},
}
\newglossaryentry{Doppler-Effect}{ sort={Doppler Effect},
  name={Doppler Effect},
  description={\textcolor{catSS}{\textbf{[Space Science]}}},
}
\newglossaryentry{Cosmic-Rays}{ sort={Cosmic Rays},
  name={Cosmic Rays},
  description={\textcolor{catSS}{\textbf{[Space Science]}}},
}
\newglossaryentry{Magnetic-Reconnection}{ sort={Magnetic Reconnection},
  name={Magnetic Reconnection},
  description={\textcolor{catSS}{\textbf{[Space Science]}}},
}
\newglossaryentry{Krmn-Line}{ sort={Kármán Line},
  name={Kármán Line},
  description={\textcolor{catSS}{\textbf{[Space Science]}}},
}
\newglossaryentry{Coronal-Hole}{ sort={Coronal Hole},
  name={Coronal Hole},
  description={\textcolor{catSS}{\textbf{[Space Science]}}},
}
\newglossaryentry{Solar-Maximum}{ sort={Solar Maximum},
  name={Solar Maximum},
  description={\textcolor{catSS}{\textbf{[Space Science]}}},
}
\newglossaryentry{Solar-Minimum}{ sort={Solar Minimum},
  name={Solar Minimum},
  description={\textcolor{catSS}{\textbf{[Space Science]}}},
}
\newglossaryentry{Sunspot}{ sort={Sunspot},
  name={Sunspot},
  description={\textcolor{catSS}{\textbf{[Space Science]}}},
}
\newglossaryentry{Solar-Cycle}{ sort={Solar Cycle},
  name={Solar Cycle},
  description={\textcolor{catSS}{\textbf{[Space Science]}}},
}
\newglossaryentry{Stellar-Remnant}{ sort={Stellar Remnant},
  name={Stellar Remnant},
  description={\textcolor{catSS}{\textbf{[Space Science]}}},
}
\newglossaryentry{Pulsar}{ sort={Pulsar},
  name={Pulsar},
  description={\textcolor{catSS}{\textbf{[Space Science]}}},
}
\newglossaryentry{Quasar}{ sort={Quasar},
  name={Quasar},
  description={\textcolor{catSS}{\textbf{[Space Science]}}},
}
\newglossaryentry{Event-Horizon}{ sort={Event Horizon},
  name={Event Horizon},
  description={\textcolor{catSS}{\textbf{[Space Science]}}},
}
\newglossaryentry{Singularity}{ sort={Singularity},
  name={Singularity},
  description={\textcolor{catSS}{\textbf{[Space Science]}}},
}
\newglossaryentry{Gamma-Ray-Burst}{ sort={Gamma-Ray Burst},
  name={Gamma-Ray Burst},
  description={\textcolor{catSS}{\textbf{[Space Science]}}},
}
\newglossaryentry{Gravitational-Wave}{ sort={Gravitational Wave},
  name={Gravitational Wave},
  description={\textcolor{catSS}{\textbf{[Space Science]}}},
}
\newglossaryentry{Gravitational-Lens}{ sort={Gravitational Lens},
  name={Gravitational Lens},
  description={\textcolor{catSS}{\textbf{[Space Science]}}},
}
\newglossaryentry{Exoplanet-Detection}{ sort={Exoplanet Detection},
  name={Exoplanet Detection},
  description={\textcolor{catSS}{\textbf{[Space Science]}}},
}
\newglossaryentry{Radial-Velocity-Method}{ sort={Radial Velocity Method},
  name={Radial Velocity Method},
  description={\textcolor{catSS}{\textbf{[Space Science]}}},
}
\newglossaryentry{Transit-Method}{ sort={Transit Method},
  name={Transit Method},
  description={\textcolor{catSS}{\textbf{[Space Science]}}},
}
\newglossaryentry{Direct-Imaging}{ sort={Direct Imaging},
  name={Direct Imaging},
  description={\textcolor{catSS}{\textbf{[Space Science]}}},
}
\newglossaryentry{Habitable-Exoplanet}{ sort={Habitable Exoplanet},
  name={Habitable Exoplanet},
  description={\textcolor{catSS}{\textbf{[Space Science]}}},
}
\newglossaryentry{Goldilocks-Zone}{ sort={Goldilocks Zone},
  name={Goldilocks Zone},
  description={\textcolor{catSS}{\textbf{[Space Science]}}},
}
\newglossaryentry{Biosignature}{ sort={Biosignature},
  name={Biosignature},
  description={\textcolor{catSS}{\textbf{[Space Science]}}},
}
\newglossaryentry{Panspermia}{ sort={Panspermia},
  name={Panspermia},
  description={\textcolor{catSS}{\textbf{[Space Science]}}},
}
\newglossaryentry{Fermi-Paradox}{ sort={Fermi Paradox},
  name={Fermi Paradox},
  description={\textcolor{catSS}{\textbf{[Space Science]}}},
}
\newglossaryentry{Drake-Equation}{ sort={Drake Equation},
  name={Drake Equation},
  description={\textcolor{catSS}{\textbf{[Space Science]}}},
}
\newglossaryentry{SETI}{ sort={SETI},
  name={SETI},
  description={\textcolor{catSS}{\textbf{[Space Science]}}},
}
\newglossaryentry{Von-Neumann-Probe}{ sort={Von Neumann Probe},
  name={Von Neumann Probe},
  description={\textcolor{catSS}{\textbf{[Space Science]}}},
}
\newglossaryentry{Dyson-Sphere}{ sort={Dyson Sphere},
  name={Dyson Sphere},
  description={\textcolor{catSS}{\textbf{[Space Science]}}},
}
\newglossaryentry{Kardashev-Scale}{ sort={Kardashev Scale},
  name={Kardashev Scale},
  description={\textcolor{catSS}{\textbf{[Space Science]}}},
}
\newglossaryentry{Light-Sail}{ sort={Light Sail},
  name={Light Sail},
  description={\textcolor{catSS}{\textbf{[Space Science]}}},
}
\newglossaryentry{Breakthrough-Starshot}{ sort={Breakthrough Starshot},
  name={Breakthrough Starshot},
  description={\textcolor{catSS}{\textbf{[Space Science]}}},
}
\newglossaryentry{ADS-B}{ sort={ADS-B},
  name={ADS-B},
  description={\textcolor{catD}{\textbf{[Drones]}}},
}
\newglossaryentry{Beyond-Visual-Line-of-Sight}{ sort={Beyond Visual Line of Sight},
  name={Beyond Visual Line of Sight},
  description={\textcolor{catD}{\textbf{[Drones]}}},
}
\newglossaryentry{Detect-and-Avoid}{ sort={Detect and Avoid},
  name={Detect and Avoid},
  description={\textcolor{catD}{\textbf{[Drones]}}},
}
\newglossaryentry{Electronic-Speed-Controller}{ sort={Electronic Speed Controller},
  name={Electronic Speed Controller},
  description={\textcolor{catD}{\textbf{[Drones]}}},
}
\newglossaryentry{Ground-Control-Station}{ sort={Ground Control Station},
  name={Ground Control Station},
  description={\textcolor{catD}{\textbf{[Drones]}}},
}
\newglossaryentry{On-Screen-Display}{ sort={On-Screen Display},
  name={On-Screen Display},
  description={\textcolor{catD}{\textbf{[Drones]}}},
}
\newglossaryentry{Received-Signal-Strength-Indicator}{ sort={Received Signal Strength Indicator},
  name={Received Signal Strength Indicator},
  description={\textcolor{catD}{\textbf{[Drones]}}},
}
\newglossaryentry{Collision-Avoidance-System}{ sort={Collision Avoidance System},
  name={Collision Avoidance System},
  description={\textcolor{catD}{\textbf{[Drones]}}},
}
\newglossaryentry{Datalink}{ sort={Datalink},
  name={Datalink},
  description={\textcolor{catD}{\textbf{[Drones]}}},
}
\newglossaryentry{Flight-Envelope-Protection}{ sort={Flight Envelope Protection},
  name={Flight Envelope Protection},
  description={\textcolor{catD}{\textbf{[Drones]}}},
}
\newglossaryentry{Geofencing}{ sort={Geofencing},
  name={Geofencing},
  description={\textcolor{catD}{\textbf{[Drones]}}},
}
\newglossaryentry{GPS-Denial}{ sort={GPS Denial},
  name={GPS Denial},
  description={\textcolor{catD}{\textbf{[Drones]}}},
}
\newglossaryentry{Gyro-Stabilization}{ sort={Gyro Stabilization},
  name={Gyro Stabilization},
  description={\textcolor{catD}{\textbf{[Drones]}}},
}
\newglossaryentry{LiDAR}{ sort={LiDAR},
  name={LiDAR},
  description={\textcolor{catD}{\textbf{[Drones]}}},
}
\newglossaryentry{Obstacle-Avoidance}{ sort={Obstacle Avoidance},
  name={Obstacle Avoidance},
  description={\textcolor{catD}{\textbf{[Drones]}}},
}
\newglossaryentry{Payload-Release}{ sort={Payload Release},
  name={Payload Release},
  description={\textcolor{catD}{\textbf{[Drones]}}},
}
\newglossaryentry{Sense-and-Avoid}{ sort={Sense and Avoid},
  name={Sense and Avoid},
  description={\textcolor{catD}{\textbf{[Drones]}}},
}
\newglossaryentry{Unmanned-Traffic-Management}{ sort={Unmanned Traffic Management},
  name={Unmanned Traffic Management},
  description={\textcolor{catD}{\textbf{[Drones]}}},
}
\newglossaryentry{Visual-Line-of-Sight}{ sort={Visual Line of Sight},
  name={Visual Line of Sight},
  description={\textcolor{catD}{\textbf{[Drones]}}},
}
\newglossaryentry{Circle-Mode}{ sort={Circle Mode},
  name={Circle Mode},
  description={\textcolor{catD}{\textbf{[Drones]}}},
}
\newglossaryentry{Follow-Me-Mode}{ sort={Follow Me Mode},
  name={Follow Me Mode},
  description={\textcolor{catD}{\textbf{[Drones]}}},
}
\newglossaryentry{Land-Mode}{ sort={Land Mode},
  name={Land Mode},
  description={\textcolor{catD}{\textbf{[Drones]}}},
}
\newglossaryentry{Takeoff-Mode}{ sort={Takeoff Mode},
  name={Takeoff Mode},
  description={\textcolor{catD}{\textbf{[Drones]}}},
}
\newglossaryentry{Data-Logger}{ sort={Data Logger},
  name={Data Logger},
  description={\textcolor{catD}{\textbf{[Drones]}}},
}
\newglossaryentry{Black-Box}{ sort={Black Box},
  name={Black Box},
  description={\textcolor{catD}{\textbf{[Drones]}}},
}
\newglossaryentry{Flight-Log}{ sort={Flight Log},
  name={Flight Log},
  description={\textcolor{catD}{\textbf{[Drones]}}},
}
\newglossaryentry{Telemetry-Link}{ sort={Telemetry Link},
  name={Telemetry Link},
  description={\textcolor{catD}{\textbf{[Drones]}}},
}
\newglossaryentry{Channel-Mapping}{ sort={Channel Mapping},
  name={Channel Mapping},
  description={\textcolor{catD}{\textbf{[Drones]}}},
}
\newglossaryentry{Receiver-Binding}{ sort={Receiver Binding},
  name={Receiver Binding},
  description={\textcolor{catD}{\textbf{[Drones]}}},
}
\newglossaryentry{Lost-Link}{ sort={Lost Link},
  name={Lost Link},
  description={\textcolor{catD}{\textbf{[Drones]}}},
}
\newglossaryentry{Low-Voltage-Cutoff}{ sort={Low Voltage Cutoff},
  name={Low Voltage Cutoff},
  description={\textcolor{catD}{\textbf{[Drones]}}},
}
\newglossaryentry{Battery-Monitoring}{ sort={Battery Monitoring},
  name={Battery Monitoring},
  description={\textcolor{catD}{\textbf{[Drones]}}},
}
\newglossaryentry{Arm-Switch}{ sort={Arm Switch},
  name={Arm Switch},
  description={\textcolor{catD}{\textbf{[Drones]}}},
}
\newglossaryentry{Motor-Protocol}{ sort={Motor Protocol},
  name={Motor Protocol},
  description={\textcolor{catD}{\textbf{[Drones]}}},
}
\newglossaryentry{IBUS}{ sort={IBUS},
  name={IBUS},
  description={\textcolor{catD}{\textbf{[Drones]}}},
}
\newglossaryentry{ELRS}{ sort={ELRS},
  name={ELRS},
  description={\textcolor{catD}{\textbf{[Drones]}}},
}
\newglossaryentry{FrSky}{ sort={FrSky},
  name={FrSky},
  description={\textcolor{catD}{\textbf{[Drones]}}},
}
\newglossaryentry{Mission-Planner-GCS}{ sort={Mission Planner GCS},
  name={Mission Planner GCS},
  description={\textcolor{catD}{\textbf{[Drones]}}},
}
\newglossaryentry{Video-Transmitter}{ sort={Video Transmitter},
  name={Video Transmitter},
  description={\textcolor{catD}{\textbf{[Drones]}}},
}
\newglossaryentry{Aerodynamics}{ sort={Aerodynamics},
  name={Aerodynamics},
  description={\textcolor{catAD}{\textbf{[Aerodynamics]}}},
}

\newglossaryentry{Aerodynamic-Center}{ sort={Aerodynamic Center},
  name={Aerodynamic Center},
  description={\textcolor{catAD}{\textbf{[Aerodynamics]}}},
}

\newglossaryentry{Aerodynamic-Force}{ sort={Aerodynamic Force},
  name={Aerodynamic Force},
  description={\textcolor{catAD}{\textbf{[Aerodynamics]}}},
}

\newglossaryentry{Aerodynamic-Heating}{ sort={Aerodynamic Heating},
  name={Aerodynamic Heating},
  description={\textcolor{catAD}{\textbf{[Aerodynamics]}}},
}

\newglossaryentry{Aeroelasticity}{ sort={Aeroelasticity},
  name={Aeroelasticity},
  description={\textcolor{catAD}{\textbf{[Aerodynamics]}}},
}

\newglossaryentry{Aileron}{ sort={Aileron},
  name={Aileron},
  description={\textcolor{catAD}{\textbf{[Aerodynamics]}}},
}

\newglossaryentry{Airfoil}{ sort={Airfoil},
  name={Airfoil},
  description={\textcolor{catAD}{\textbf{[Aerodynamics]}}},
}

\newglossaryentry{Angle-of-Attack}{ sort={Angle of Attack},
  name={Angle of Attack},
  description={\textcolor{catAD}{\textbf{[Aerodynamics]}}},
}

\newglossaryentry{Angle-of-Incidence}{ sort={Angle of Incidence},
  name={Angle of Incidence},
  description={\textcolor{catAD}{\textbf{[Aerodynamics]}}},
}

\newglossaryentry{Anhedral}{ sort={Anhedral},
  name={Anhedral},
  description={\textcolor{catAD}{\textbf{[Aerodynamics]}}},
}

\newglossaryentry{Aspect-Ratio}{ sort={Aspect Ratio},
  name={Aspect Ratio},
  description={\textcolor{catAD}{\textbf{[Aerodynamics]}}},
}

\newglossaryentry{Attached-Flow}{ sort={Attached Flow},
  name={Attached Flow},
  description={\textcolor{catAD}{\textbf{[Aerodynamics]}}},
}

\newglossaryentry{Autorotation}{ sort={Autorotation},
  name={Autorotation},
  description={\textcolor{catAD}{\textbf{[Aerodynamics]}}},
}

\newglossaryentry{Boundary-Layer}{ sort={Boundary Layer},
  name={Boundary Layer},
  description={\textcolor{catAD}{\textbf{[Aerodynamics]}}},
}

\newglossaryentry{Boundary-Layer-Separation}{ sort={Boundary Layer Separation},
  name={Boundary Layer Separation},
  description={\textcolor{catAD}{\textbf{[Aerodynamics]}}},
}

\newglossaryentry{Boundary-Layer-Transition}{ sort={Boundary Layer Transition},
  name={Boundary Layer Transition},
  description={\textcolor{catAD}{\textbf{[Aerodynamics]}}},
}

\newglossaryentry{Camber}{ sort={Camber},
  name={Camber},
  description={\textcolor{catAD}{\textbf{[Aerodynamics]}}},
}

\newglossaryentry{Center-of-Pressure}{ sort={Center of Pressure},
  name={Center of Pressure},
  description={\textcolor{catAD}{\textbf{[Aerodynamics]}}},
}

\newglossaryentry{Chord}{ sort={Chord},
  name={Chord},
  description={\textcolor{catAD}{\textbf{[Aerodynamics]}}},
}

\newglossaryentry{Chord-Line}{ sort={Chord Line},
  name={Chord Line},
  description={\textcolor{catAD}{\textbf{[Aerodynamics]}}},
}

\newglossaryentry{Circulation}{ sort={Circulation},
  name={Circulation},
  description={\textcolor{catAD}{\textbf{[Aerodynamics]}}},
}

\newglossaryentry{Compressible-Flow}{ sort={Compressible Flow},
  name={Compressible Flow},
  description={\textcolor{catAD}{\textbf{[Aerodynamics]}}},
}

\newglossaryentry{Critical-Angle-of-Attack}{ sort={Critical Angle of Attack},
  name={Critical Angle of Attack},
  description={\textcolor{catAD}{\textbf{[Aerodynamics]}}},
}

\newglossaryentry{Critical-Mach-Number}{ sort={Critical Mach Number},
  name={Critical Mach Number},
  description={\textcolor{catAD}{\textbf{[Aerodynamics]}}},
}

\newglossaryentry{Dihedral}{ sort={Dihedral},
  name={Dihedral},
  description={\textcolor{catAD}{\textbf{[Aerodynamics]}}},
}


\newglossaryentry{Drag}{ sort={Drag},
  name={Drag},
  description={\textcolor{catAD}{\textbf{[Aerodynamics]}}},
}

\newglossaryentry{Drag-Coefficient}{ sort={Drag Coefficient},
  name={Drag Coefficient},
  description={\textcolor{catAD}{\textbf{[Aerodynamics]}}},
}

\newglossaryentry{Dynamic-Stall}{ sort={Dynamic Stall},
  name={Dynamic Stall},
  description={\textcolor{catAD}{\textbf{[Aerodynamics]}}},
}

\newglossaryentry{Dynamic-Stability}{ sort={Dynamic Stability},
  name={Dynamic Stability},
  description={\textcolor{catAD}{\textbf{[Aerodynamics]}}},
}

\newglossaryentry{Dutch-Roll}{ sort={Dutch Roll},
  name={Dutch Roll},
  description={\textcolor{catAD}{\textbf{[Aerodynamics]}}},
}

\newglossaryentry{Elevator}{ sort={Elevator},
  name={Elevator},
  description={\textcolor{catAD}{\textbf{[Aerodynamics]}}},
}

\newglossaryentry{Elevon}{ sort={Elevon},
  name={Elevon},
  description={\textcolor{catAD}{\textbf{[Aerodynamics]}}},
}

\newglossaryentry{Expansion-Wave}{ sort={Expansion Wave},
  name={Expansion Wave},
  description={\textcolor{catAD}{\textbf{[Aerodynamics]}}},
}

\newglossaryentry{Flap}{ sort={Flap},
  name={Flap},
  description={\textcolor{catAD}{\textbf{[Aerodynamics]}}},
}


\newglossaryentry{Form-Drag}{ sort={Form Drag},
  name={Form Drag},
  description={\textcolor{catAD}{\textbf{[Aerodynamics]}}},
}

\newglossaryentry{Friction-Drag}{ sort={Friction Drag},
  name={Friction Drag},
  description={\textcolor{catAD}{\textbf{[Aerodynamics]}}},
}

\newglossaryentry{Hypersonic-Flow}{ sort={Hypersonic Flow},
  name={Hypersonic Flow},
  description={\textcolor{catAD}{\textbf{[Aerodynamics]}}},
}

\newglossaryentry{Incompressible-Flow}{ sort={Incompressible Flow},
  name={Incompressible Flow},
  description={\textcolor{catAD}{\textbf{[Aerodynamics]}}},
}

\newglossaryentry{Induced-Drag}{ sort={Induced Drag},
  name={Induced Drag},
  description={\textcolor{catAD}{\textbf{[Aerodynamics]}}},
}

\newglossaryentry{Induced-Velocity}{ sort={Induced Velocity},
  name={Induced Velocity},
  description={\textcolor{catAD}{\textbf{[Aerodynamics]}}},
}

\newglossaryentry{Interference-Drag}{ sort={Interference Drag},
  name={Interference Drag},
  description={\textcolor{catAD}{\textbf{[Aerodynamics]}}},
}

\newglossaryentry{Inviscid-Flow}{ sort={Inviscid Flow},
  name={Inviscid Flow},
  description={\textcolor{catAD}{\textbf{[Aerodynamics]}}},
}

\newglossaryentry{Irrotational-Flow}{ sort={Irrotational Flow},
  name={Irrotational Flow},
  description={\textcolor{catAD}{\textbf{[Aerodynamics]}}},
}

\newglossaryentry{Kutta-Condition}{ sort={Kutta Condition},
  name={Kutta Condition},
  description={\textcolor{catAD}{\textbf{[Aerodynamics]}}},
}

\newglossaryentry{Kutta-Joukowski-Theorem}{ sort={Kutta-Joukowski Theorem},
  name={Kutta-Joukowski Theorem},
  description={\textcolor{catAD}{\textbf{[Aerodynamics]}}},
}

\newglossaryentry{Laminar-Flow}{ sort={Laminar Flow},
  name={Laminar Flow},
  description={\textcolor{catAD}{\textbf{[Aerodynamics]}}},
}

\newglossaryentry{Lateral-Stability}{ sort={Lateral Stability},
  name={Lateral Stability},
  description={\textcolor{catAD}{\textbf{[Aerodynamics]}}},
}

\newglossaryentry{Leading-Edge}{ sort={Leading Edge},
  name={Leading Edge},
  description={\textcolor{catAD}{\textbf{[Aerodynamics]}}},
}

\newglossaryentry{Lift}{ sort={Lift},
  name={Lift},
  description={\textcolor{catAD}{\textbf{[Aerodynamics]}}},
}

\newglossaryentry{Lift-Coefficient}{ sort={Lift Coefficient},
  name={Lift Coefficient},
  description={\textcolor{catAD}{\textbf{[Aerodynamics]}}},
}

\newglossaryentry{Longitudinal-Stability}{ sort={Longitudinal Stability},
  name={Longitudinal Stability},
  description={\textcolor{catAD}{\textbf{[Aerodynamics]}}},
}

\newglossaryentry{Mach-Angle}{ sort={Mach Angle},
  name={Mach Angle},
  description={\textcolor{catAD}{\textbf{[Aerodynamics]}}},
}

\newglossaryentry{Mach-Cone}{ sort={Mach Cone},
  name={Mach Cone},
  description={\textcolor{catAD}{\textbf{[Aerodynamics]}}},
}

\newglossaryentry{Mean-Aerodynamic-Chord}{ sort={Mean Aerodynamic Chord},
  name={Mean Aerodynamic Chord},
  description={\textcolor{catAD}{\textbf{[Aerodynamics]}}},
}

\newglossaryentry{Momentum-Thickness}{ sort={Momentum Thickness},
  name={Momentum Thickness},
  description={\textcolor{catAD}{\textbf{[Aerodynamics]}}},
}

\newglossaryentry{Navier-Stokes-Equations}{ sort={Navier-Stokes Equations},
  name={Navier-Stokes Equations},
  description={\textcolor{catAD}{\textbf{[Aerodynamics]}}},
}

\newglossaryentry{Neutral-Point}{ sort={Neutral Point},
  name={Neutral Point},
  description={\textcolor{catAD}{\textbf{[Aerodynamics]}}},
}

\newglossaryentry{Normal-Shock}{ sort={Normal Shock},
  name={Normal Shock},
  description={\textcolor{catAD}{\textbf{[Aerodynamics]}}},
}


\newglossaryentry{Oblique-Shock}{ sort={Oblique Shock},
  name={Oblique Shock},
  description={\textcolor{catAD}{\textbf{[Aerodynamics]}}},
}

\newglossaryentry{Parasite-Drag}{ sort={Parasite Drag},
  name={Parasite Drag},
  description={\textcolor{catAD}{\textbf{[Aerodynamics]}}},
}

\newglossaryentry{Phugoid}{ sort={Phugoid},
  name={Phugoid},
  description={\textcolor{catAD}{\textbf{[Aerodynamics]}}},
}

\newglossaryentry{Pitching-Moment}{ sort={Pitching Moment},
  name={Pitching Moment},
  description={\textcolor{catAD}{\textbf{[Aerodynamics]}}},
}

\newglossaryentry{Pitot-Tube}{ sort={Pitot Tube},
  name={Pitot Tube},
  description={\textcolor{catAD}{\textbf{[Aerodynamics]}}},
}

\newglossaryentry{Prandtl-Meyer-Expansion}{ sort={Prandtl-Meyer Expansion},
  name={Prandtl-Meyer Expansion},
  description={\textcolor{catAD}{\textbf{[Aerodynamics]}}},
}

\newglossaryentry{Pressure-Drag}{ sort={Pressure Drag},
  name={Pressure Drag},
  description={\textcolor{catAD}{\textbf{[Aerodynamics]}}},
}


\newglossaryentry{Reattachment}{ sort={Reattachment},
  name={Reattachment},
  description={\textcolor{catAD}{\textbf{[Aerodynamics]}}},
}

\newglossaryentry{Relative-Wind}{ sort={Relative Wind},
  name={Relative Wind},
  description={\textcolor{catAD}{\textbf{[Aerodynamics]}}},
}

\newglossaryentry{Reynolds-Number}{ sort={Reynolds Number},
  name={Reynolds Number},
  description={\textcolor{catAD}{\textbf{[Aerodynamics]}}},
}

\newglossaryentry{Rolling-Moment}{ sort={Rolling Moment},
  name={Rolling Moment},
  description={\textcolor{catAD}{\textbf{[Aerodynamics]}}},
}

\newglossaryentry{Rotor}{ sort={Rotor},
  name={Rotor},
  description={\textcolor{catAD}{\textbf{[Aerodynamics]}}},
}

\newglossaryentry{Rudder}{ sort={Rudder},
  name={Rudder},
  description={\textcolor{catAD}{\textbf{[Aerodynamics]}}},
}

\newglossaryentry{Schlieren-Photography}{ sort={Schlieren Photography},
  name={Schlieren Photography},
  description={\textcolor{catAD}{\textbf{[Aerodynamics]}}},
}

\newglossaryentry{Separated-Flow}{ sort={Separated Flow},
  name={Separated Flow},
  description={\textcolor{catAD}{\textbf{[Aerodynamics]}}},
}

\newglossaryentry{Shock-Wave}{ sort={Shock Wave},
  name={Shock Wave},
  description={\textcolor{catAD}{\textbf{[Aerodynamics]}}},
}

\newglossaryentry{Short-Period-Mode}{ sort={Short Period Mode},
  name={Short Period Mode},
  description={\textcolor{catAD}{\textbf{[Aerodynamics]}}},
}

\newglossaryentry{Skin-Friction}{ sort={Skin Friction},
  name={Skin Friction},
  description={\textcolor{catAD}{\textbf{[Aerodynamics]}}},
}

\newglossaryentry{Skin-Friction-Drag}{ sort={Skin Friction Drag},
  name={Skin Friction Drag},
  description={\textcolor{catAD}{\textbf{[Aerodynamics]}}},
}

\newglossaryentry{Slat}{ sort={Slat},
  name={Slat},
  description={\textcolor{catAD}{\textbf{[Aerodynamics]}}},
}

\newglossaryentry{Sonic-Boom}{ sort={Sonic Boom},
  name={Sonic Boom},
  description={\textcolor{catAD}{\textbf{[Aerodynamics]}}},
}

\newglossaryentry{Spoiler}{ sort={Spoiler},
  name={Spoiler},
  description={\textcolor{catAD}{\textbf{[Aerodynamics]}}},
}

\newglossaryentry{Stall}{ sort={Stall},
  name={Stall},
  description={\textcolor{catAD}{\textbf{[Aerodynamics]}}},
}

\newglossaryentry{Static-Margin}{ sort={Static Margin},
  name={Static Margin},
  description={\textcolor{catAD}{\textbf{[Aerodynamics]}}},
}

\newglossaryentry{Static-Port}{ sort={Static Port},
  name={Static Port},
  description={\textcolor{catAD}{\textbf{[Aerodynamics]}}},
}

\newglossaryentry{Static-Stability}{ sort={Static Stability},
  name={Static Stability},
  description={\textcolor{catAD}{\textbf{[Aerodynamics]}}},
}

\newglossaryentry{Steady-Flow}{ sort={Steady Flow},
  name={Steady Flow},
  description={\textcolor{catAD}{\textbf{[Aerodynamics]}}},
}

\newglossaryentry{Subsonic-Flow}{ sort={Subsonic Flow},
  name={Subsonic Flow},
  description={\textcolor{catAD}{\textbf{[Aerodynamics]}}},
}

\newglossaryentry{Supersonic-Flow}{ sort={Supersonic Flow},
  name={Supersonic Flow},
  description={\textcolor{catAD}{\textbf{[Aerodynamics]}}},
}

\newglossaryentry{Sweep-Angle}{ sort={Sweep Angle},
  name={Sweep Angle},
  description={\textcolor{catAD}{\textbf{[Aerodynamics]}}},
}

\newglossaryentry{Thickness-Ratio}{ sort={Thickness Ratio},
  name={Thickness Ratio},
  description={\textcolor{catAD}{\textbf{[Aerodynamics]}}},
}


\newglossaryentry{Thrust-Vectoring}{ sort={Thrust Vectoring},
  name={Thrust Vectoring},
  description={\textcolor{catAD}{\textbf{[Aerodynamics]}}},
}

\newglossaryentry{Trailing-Edge}{ sort={Trailing Edge},
  name={Trailing Edge},
  description={\textcolor{catAD}{\textbf{[Aerodynamics]}}},
}

\newglossaryentry{Transonic-Flow}{ sort={Transonic Flow},
  name={Transonic Flow},
  description={\textcolor{catAD}{\textbf{[Aerodynamics]}}},
}

\newglossaryentry{Trim-Drag}{ sort={Trim Drag},
  name={Trim Drag},
  description={\textcolor{catAD}{\textbf{[Aerodynamics]}}},
}



\newglossaryentry{Turbulent-Flow}{ sort={Turbulent Flow},
  name={Turbulent Flow},
  description={\textcolor{catAD}{\textbf{[Aerodynamics]}}},
}

\newglossaryentry{Unsteady-Flow}{ sort={Unsteady Flow},
  name={Unsteady Flow},
  description={\textcolor{catAD}{\textbf{[Aerodynamics]}}},
}

\newglossaryentry{Upwash}{ sort={Upwash},
  name={Upwash},
  description={\textcolor{catAD}{\textbf{[Aerodynamics]}}},
}

\newglossaryentry{Vortex}{ sort={Vortex},
  name={Vortex},
  description={\textcolor{catAD}{\textbf{[Aerodynamics]}}},
}

\newglossaryentry{Vortex-Generator}{ sort={Vortex Generator},
  name={Vortex Generator},
  description={\textcolor{catAD}{\textbf{[Aerodynamics]}}},
}

\newglossaryentry{Wake}{ sort={Wake},
  name={Wake},
  description={\textcolor{catAD}{\textbf{[Aerodynamics]}}},
}

\newglossaryentry{Wall-Shear-Stress}{ sort={Wall Shear Stress},
  name={Wall Shear Stress},
  description={\textcolor{catAD}{\textbf{[Aerodynamics]}}},
}

\newglossaryentry{Washout}{ sort={Washout},
  name={Washout},
  description={\textcolor{catAD}{\textbf{[Aerodynamics]}}},
}

\newglossaryentry{Wave-Drag}{ sort={Wave Drag},
  name={Wave Drag},
  description={\textcolor{catAD}{\textbf{[Aerodynamics]}}},
}

\newglossaryentry{Wing}{ sort={Wing},
  name={Wing},
  description={\textcolor{catAD}{\textbf{[Aerodynamics]}}},
}

\newglossaryentry{Wing-Area}{ sort={Wing Area},
  name={Wing Area},
  description={\textcolor{catAD}{\textbf{[Aerodynamics]}}},
}

\newglossaryentry{Wing-Span}{ sort={Wing Span},
  name={Wing Span},
  description={\textcolor{catAD}{\textbf{[Aerodynamics]}}},
}

\newglossaryentry{Winglet}{ sort={Winglet},
  name={Winglet},
  description={\textcolor{catAD}{\textbf{[Aerodynamics]}}},
}

\newglossaryentry{Wing-Root}{ sort={Wing Root},
  name={Wing Root},
  description={\textcolor{catAD}{\textbf{[Aerodynamics]}}},
}

\newglossaryentry{Wing-Tip}{ sort={Wing Tip},
  name={Wing Tip},
  description={\textcolor{catAD}{\textbf{[Aerodynamics]}}},
}

\newglossaryentry{Yawing-Moment}{ sort={Yawing Moment},
  name={Yawing Moment},
  description={\textcolor{catAD}{\textbf{[Aerodynamics]}}},
}

\newglossaryentry{Adverse-Yaw}{ sort={Adverse Yaw},
  name={Adverse Yaw},
  description={\textcolor{catAD}{\textbf{[Aerodynamics]}}},
}

\newglossaryentry{Airbrake}{ sort={Airbrake},
  name={Airbrake},
  description={\textcolor{catAD}{\textbf{[Aerodynamics]}}},
}

\newglossaryentry{Boundary-Layer-Control}{ sort={Boundary Layer Control},
  name={Boundary Layer Control},
  description={\textcolor{catAD}{\textbf{[Aerodynamics]}}},
}

\newglossaryentry{Choked-Flow}{ sort={Choked Flow},
  name={Choked Flow},
  description={\textcolor{catAD}{\textbf{[Aerodynamics]}}},
}

\newglossaryentry{Computational-Fluid-Dynamics}{ sort={Computational Fluid Dynamics},
  name={Computational Fluid Dynamics},
  description={\textcolor{catAD}{\textbf{[Aerodynamics]}}},
}

\newglossaryentry{Contra-Rotating-Propeller}{ sort={Contra-Rotating Propeller},
  name={Contra-Rotating Propeller},
  description={\textcolor{catAD}{\textbf{[Aerodynamics]}}},
}

\newglossaryentry{Diffuser}{ sort={Diffuser},
  name={Diffuser},
  description={\textcolor{catAD}{\textbf{[Aerodynamics]}}},
}

\newglossaryentry{Displacement-Thickness}{ sort={Displacement Thickness},
  name={Displacement Thickness},
  description={\textcolor{catAD}{\textbf{[Aerodynamics]}}},
}

\newglossaryentry{Divergence}{ sort={Divergence},
  name={Divergence},
  description={\textcolor{catAD}{\textbf{[Aerodynamics]}}},
}

\newglossaryentry{Euler-Equations}{ sort={Euler Equations},
  name={Euler Equations},
  description={\textcolor{catAD}{\textbf{[Aerodynamics]}}},
}

\newglossaryentry{Finite-Element-Method}{ sort={Finite Element Method},
  name={Finite Element Method},
  description={\textcolor{catAD}{\textbf{[Aerodynamics]}}},
}

\newglossaryentry{Finite-Volume-Method}{ sort={Finite Volume Method},
  name={Finite Volume Method},
  description={\textcolor{catAD}{\textbf{[Aerodynamics]}}},
}

\newglossaryentry{Force-Balance}{ sort={Force Balance},
  name={Force Balance},
  description={\textcolor{catAD}{\textbf{[Aerodynamics]}}},
}

\newglossaryentry{Grid}{ sort={Grid},
  name={Grid},
  description={\textcolor{catAD}{\textbf{[Aerodynamics]}}},
}

\newglossaryentry{LES}{ sort={LES},
  name={LES},
  description={\textcolor{catAD}{\textbf{[Aerodynamics]}}},
}

\newglossaryentry{Mach-Number}{ sort={Mach Number},
  name={Mach Number},
  description={\textcolor{catAD}{\textbf{[Aerodynamics]}}},
}

\newglossaryentry{Mesh}{ sort={Mesh},
  name={Mesh},
  description={\textcolor{catAD}{\textbf{[Aerodynamics]}}},
}

\newglossaryentry{Particle-Image-Velocimetry}{ sort={Particle Image Velocimetry},
  name={Particle Image Velocimetry},
  description={\textcolor{catAD}{\textbf{[Aerodynamics]}}},
}

\newglossaryentry{Ramjet}{ sort={Ramjet},
  name={Ramjet},
  description={\textcolor{catAD}{\textbf{[Aerodynamics]}}},
}

\newglossaryentry{RANS}{ sort={RANS},
  name={RANS},
  description={\textcolor{catAD}{\textbf{[Aerodynamics]}}},
}

\newglossaryentry{DNS}{ sort={DNS},
  name={DNS},
  description={\textcolor{catAD}{\textbf{[Aerodynamics]}}},
}

\newglossaryentry{Scramjet}{ sort={Scramjet},
  name={Scramjet},
  description={\textcolor{catAD}{\textbf{[Aerodynamics]}}},
}

\newglossaryentry{Smoke-Visualization}{ sort={Smoke Visualization},
  name={Smoke Visualization},
  description={\textcolor{catAD}{\textbf{[Aerodynamics]}}},
}

\newglossaryentry{Static-Test}{ sort={Static Test},
  name={Static Test},
  description={\textcolor{catAD}{\textbf{[Aerodynamics]}}},
}

\newglossaryentry{Test-Section}{ sort={Test Section},
  name={Test Section},
  description={\textcolor{catAD}{\textbf{[Aerodynamics]}}},
}

\newglossaryentry{Transitional-Flow}{ sort={Transitional Flow},
  name={Transitional Flow},
  description={\textcolor{catAD}{\textbf{[Aerodynamics]}}},
}

\newglossaryentry{Turbulence-Model}{ sort={Turbulence Model},
  name={Turbulence Model},
  description={\textcolor{catAD}{\textbf{[Aerodynamics]}}},
}

\newglossaryentry{Vortex-Shedding}{ sort={Vortex Shedding},
  name={Vortex Shedding},
  description={\textcolor{catAD}{\textbf{[Aerodynamics]}}},
}

\newglossaryentry{Wind-Tunnel}{ sort={Wind Tunnel},
  name={Wind Tunnel},
  description={\textcolor{catAD}{\textbf{[Aerodynamics]}}},
}

\newglossaryentry{Aerothermal-Heating}{ sort={Aerothermal Heating},
  name={Aerothermal Heating},
  description={\textcolor{catAD}{\textbf{[Aerodynamics]}}},
}

\newglossaryentry{Ablation}{ sort={Ablation},
  name={Ablation},
  description={\textcolor{catAD}{\textbf{[Aerodynamics]}}},
}

\newglossaryentry{Re-entry}{ sort={Re-entry},
  name={Re-entry},
  description={\textcolor{catAD}{\textbf{[Aerodynamics]}}},
}

\newglossaryentry{Shock-Layer}{ sort={Shock Layer},
  name={Shock Layer},
  description={\textcolor{catAD}{\textbf{[Aerodynamics]}}},
}

\newglossaryentry{Thermal-Protection-System}{ sort={Thermal Protection System},
  name={Thermal Protection System},
  description={\textcolor{catAD}{\textbf{[Aerodynamics]}}},
}

\newglossaryentry{Pulsed-Detonation-Engine}{ sort={Pulsed Detonation Engine},
  name={Pulsed Detonation Engine},
  description={\textcolor{catAD}{\textbf{[Aerodynamics]}}},
}

\newglossaryentry{Turbine}{ sort={Turbine},
  name={Turbine},
  description={\textcolor{catAD}{\textbf{[Aerodynamics]}}},
}

\newglossaryentry{Compressor}{ sort={Compressor},
  name={Compressor},
  description={\textcolor{catAD}{\textbf{[Aerodynamics]}}},
}

\newglossaryentry{Spoileron}{ sort={Spoileron},
  name={Spoileron},
  description={\textcolor{catAD}{\textbf{[Aerodynamics]}}},
}

\newglossaryentry{Strake}{ sort={Strake},
  name={Strake},
  description={\textcolor{catAD}{\textbf{[Aerodynamics]}}},
}

\newglossaryentry{Canard}{ sort={Canard},
  name={Canard},
  description={\textcolor{catAD}{\textbf{[Aerodynamics]}}},
}

\newglossaryentry{Blended-Wing-Body}{ sort={Blended Wing Body},
  name={Blended Wing Body},
  description={\textcolor{catAD}{\textbf{[Aerodynamics]}}},
}

\newglossaryentry{Flying-Wing}{ sort={Flying Wing},
  name={Flying Wing},
  description={\textcolor{catAD}{\textbf{[Aerodynamics]}}},
}

\newglossaryentry{Delta-Wing}{ sort={Delta Wing},
  name={Delta Wing},
  description={\textcolor{catAD}{\textbf{[Aerodynamics]}}},
}

\newglossaryentry{Swept-Wing}{ sort={Swept Wing},
  name={Swept Wing},
  description={\textcolor{catAD}{\textbf{[Aerodynamics]}}},
}

\newglossaryentry{Forward-Swept-Wing}{ sort={Forward Swept Wing},
  name={Forward Swept Wing},
  description={\textcolor{catAD}{\textbf{[Aerodynamics]}}},
}

\newglossaryentry{Ogive}{ sort={Ogive},
  name={Ogive},
  description={\textcolor{catAD}{\textbf{[Aerodynamics]}}},
}

\newglossaryentry{Streamline}{ sort={Streamline},
  name={Streamline},
  description={\textcolor{catAD}{\textbf{[Aerodynamics]}}},
}

\newglossaryentry{Pathline}{ sort={Pathline},
  name={Pathline},
  description={\textcolor{catAD}{\textbf{[Aerodynamics]}}},
}

\newglossaryentry{Streakline}{ sort={Streakline},
  name={Streakline},
  description={\textcolor{catAD}{\textbf{[Aerodynamics]}}},
}

\newglossaryentry{Continuity-Equation}{ sort={Continuity Equation},
  name={Continuity Equation},
  description={\textcolor{catAD}{\textbf{[Aerodynamics]}}},
}

\newglossaryentry{Momentum-Equation}{ sort={Momentum Equation},
  name={Momentum Equation},
  description={\textcolor{catAD}{\textbf{[Aerodynamics]}}},
}

\newglossaryentry{Energy-Equation}{ sort={Energy Equation},
  name={Energy Equation},
  description={\textcolor{catAD}{\textbf{[Aerodynamics]}}},
}

\newglossaryentry{Bernoulli-Equation}{ sort={Bernoulli Equation},
  name={Bernoulli Equation},
  description={\textcolor{catAD}{\textbf{[Aerodynamics]}}},
}

\newglossaryentry{Camber-Line}{ sort={Camber Line},
  name={Camber Line},
  description={\textcolor{catAD}{\textbf{[Aerodynamics]}}},
}

\newglossaryentry{Mean-Line}{ sort={Mean Line},
  name={Mean Line},
  description={\textcolor{catAD}{\textbf{[Aerodynamics]}}},
}

\newglossaryentry{Pressure-Coefficient}{ sort={Pressure Coefficient},
  name={Pressure Coefficient},
  description={\textcolor{catAD}{\textbf{[Aerodynamics]}}},
}

\newglossaryentry{Stagnation-Point}{ sort={Stagnation Point},
  name={Stagnation Point},
  description={\textcolor{catAD}{\textbf{[Aerodynamics]}}},
}

\newglossaryentry{Stagnation-Pressure}{ sort={Stagnation Pressure},
  name={Stagnation Pressure},
  description={\textcolor{catAD}{\textbf{[Aerodynamics]}}},
}

\newglossaryentry{Static-Pressure}{ sort={Static Pressure},
  name={Static Pressure},
  description={\textcolor{catAD}{\textbf{[Aerodynamics]}}},
}

\newglossaryentry{Total-Pressure}{ sort={Total Pressure},
  name={Total Pressure},
  description={\textcolor{catAD}{\textbf{[Aerodynamics]}}},
}

\newglossaryentry{Aerodynamic-Twist}{ sort={Aerodynamic Twist},
  name={Aerodynamic Twist},
  description={\textcolor{catAD}{\textbf{[Aerodynamics]}}},
}

\newglossaryentry{Geometric-Twist}{ sort={Geometric Twist},
  name={Geometric Twist},
  description={\textcolor{catAD}{\textbf{[Aerodynamics]}}},
}

\newglossaryentry{Wash-In}{ sort={Wash-In},
  name={Wash-In},
  description={\textcolor{catAD}{\textbf{[Aerodynamics]}}},
}

\newglossaryentry{Oswald-Span-Efficiency}{ sort={Oswald Span Efficiency},
  name={Oswald Span Efficiency},
  description={\textcolor{catAD}{\textbf{[Aerodynamics]}}},
}

\newglossaryentry{Elliptical-Wing}{ sort={Elliptical Wing},
  name={Elliptical Wing},
  description={\textcolor{catAD}{\textbf{[Aerodynamics]}}},
}

\newglossaryentry{Taper-Ratio}{ sort={Taper Ratio},
  name={Taper Ratio},
  description={\textcolor{catAD}{\textbf{[Aerodynamics]}}},
}

\newglossaryentry{Wing-Loading}{ sort={Wing Loading},
  name={Wing Loading},
  description={\textcolor{catAD}{\textbf{[Aerodynamics]}}},
}

\newglossaryentry{Load-Factor}{ sort={Load Factor},
  name={Load Factor},
  description={\textcolor{catAD}{\textbf{[Aerodynamics]}}},
}

\newglossaryentry{V-n-Diagram}{ sort={V-n Diagram},
  name={V-n Diagram},
  description={\textcolor{catAD}{\textbf{[Aerodynamics]}}},
}

\newglossaryentry{Maneuverability}{ sort={Maneuverability},
  name={Maneuverability},
  description={\textcolor{catAD}{\textbf{[Aerodynamics]}}},
}

\newglossaryentry{Agility}{ sort={Agility},
  name={Agility},
  description={\textcolor{catAD}{\textbf{[Aerodynamics]}}},
}

\newglossaryentry{Turning-Radius}{ sort={Turning Radius},
  name={Turning Radius},
  description={\textcolor{catAD}{\textbf{[Aerodynamics]}}},
}

\newglossaryentry{Rate-of-Turn}{ sort={Rate of Turn},
  name={Rate of Turn},
  description={\textcolor{catAD}{\textbf{[Aerodynamics]}}},
}

\newglossaryentry{Coordinated-Turn}{ sort={Coordinated Turn},
  name={Coordinated Turn},
  description={\textcolor{catAD}{\textbf{[Aerodynamics]}}},
}

\newglossaryentry{Sideslip}{ sort={Sideslip},
  name={Sideslip},
  description={\textcolor{catAD}{\textbf{[Aerodynamics]}}},
}

\newglossaryentry{Ground-Effect}{ sort={Ground Effect},
  name={Ground Effect},
  description={\textcolor{catAD}{\textbf{[Aerodynamics]}}},
}

\newglossaryentry{Reynolds-Analogy}{ sort={Reynolds Analogy},
  name={Reynolds Analogy},
  description={\textcolor{catAD}{\textbf{[Aerodynamics]}}},
}

\newglossaryentry{Prandtl-Number}{ sort={Prandtl Number},
  name={Prandtl Number},
  description={\textcolor{catAD}{\textbf{[Aerodynamics]}}},
}

\newglossaryentry{Mach-Tuck}{ sort={Mach Tuck},
  name={Mach Tuck},
  description={\textcolor{catAD}{\textbf{[Aerodynamics]}}},
}

\newglossaryentry{Shock-Stall}{ sort={Shock Stall},
  name={Shock Stall},
  description={\textcolor{catAD}{\textbf{[Aerodynamics]}}},
}

\newglossaryentry{Buffeting}{ sort={Buffeting},
  name={Buffeting},
  description={\textcolor{catAD}{\textbf{[Aerodynamics]}}},
}

\newglossaryentry{Flutter}{ sort={Flutter},
  name={Flutter},
  description={\textcolor{catAD}{\textbf{[Aerodynamics]}}},
}

\newglossaryentry{Control-Reversal}{ sort={Control Reversal},
  name={Control Reversal},
  description={\textcolor{catAD}{\textbf{[Aerodynamics]}}},
}

\newglossaryentry{Gust-Load}{ sort={Gust Load},
  name={Gust Load},
  description={\textcolor{catAD}{\textbf{[Aerodynamics]}}},
}

\newglossaryentry{Maneuvering-Load}{ sort={Maneuvering Load},
  name={Maneuvering Load},
  description={\textcolor{catAD}{\textbf{[Aerodynamics]}}},
}

\newglossaryentry{Lift-Distribution}{ sort={Lift Distribution},
  name={Lift Distribution},
  description={\textcolor{catAD}{\textbf{[Aerodynamics]}}},
}

\newglossaryentry{Trailing-Vortex}{ sort={Trailing Vortex},
  name={Trailing Vortex},
  description={\textcolor{catAD}{\textbf{[Aerodynamics]}}},
}

\newglossaryentry{Starting-Vortex}{ sort={Starting Vortex},
  name={Starting Vortex},
  description={\textcolor{catAD}{\textbf{[Aerodynamics]}}},
}

\newglossaryentry{Horseshoe-Vortex}{ sort={Horseshoe Vortex},
  name={Horseshoe Vortex},
  description={\textcolor{catAD}{\textbf{[Aerodynamics]}}},
}

\newglossaryentry{Panel-Method}{ sort={Panel Method},
  name={Panel Method},
  description={\textcolor{catAD}{\textbf{[Aerodynamics]}}},
}

\newglossaryentry{Vortex-Lattice-Method}{ sort={Vortex Lattice Method},
  name={Vortex Lattice Method},
  description={\textcolor{catAD}{\textbf{[Aerodynamics]}}},
}

\newglossaryentry{Source-Panel-Method}{ sort={Source Panel Method},
  name={Source Panel Method},
  description={\textcolor{catAD}{\textbf{[Aerodynamics]}}},
}

\newglossaryentry{Doublet}{ sort={Doublet},
  name={Doublet},
  description={\textcolor{catAD}{\textbf{[Aerodynamics]}}},
}

\newglossaryentry{Rankine-Body}{ sort={Rankine Body},
  name={Rankine Body},
  description={\textcolor{catAD}{\textbf{[Aerodynamics]}}},
}

\newglossaryentry{Joukowski-Airfoil}{ sort={Joukowski Airfoil},
  name={Joukowski Airfoil},
  description={\textcolor{catAD}{\textbf{[Aerodynamics]}}},
}

\newglossaryentry{Thin-Airfoil-Theory}{ sort={Thin Airfoil Theory},
  name={Thin Airfoil Theory},
  description={\textcolor{catAD}{\textbf{[Aerodynamics]}}},
}

\newglossaryentry{Lifting-Line-Theory}{ sort={Lifting Line Theory},
  name={Lifting Line Theory},
  description={\textcolor{catAD}{\textbf{[Aerodynamics]}}},
}

\newglossaryentry{Lifting-Surface-Theory}{ sort={Lifting Surface Theory},
  name={Lifting Surface Theory},
  description={\textcolor{catAD}{\textbf{[Aerodynamics]}}},
}

\newglossaryentry{Sweep-Theory}{ sort={Sweep Theory},
  name={Sweep Theory},
  description={\textcolor{catAD}{\textbf{[Aerodynamics]}}},
}

\newglossaryentry{Slender-Body-Theory}{ sort={Slender Body Theory},
  name={Slender Body Theory},
  description={\textcolor{catAD}{\textbf{[Aerodynamics]}}},
}

\newglossaryentry{Area-Rule}{ sort={Area Rule},
  name={Area Rule},
  description={\textcolor{catAD}{\textbf{[Aerodynamics]}}},
}

\newglossaryentry{Supercritical-Airfoil}{ sort={Supercritical Airfoil},
  name={Supercritical Airfoil},
  description={\textcolor{catAD}{\textbf{[Aerodynamics]}}},
}

\newglossaryentry{Shock-Free-Airfoil}{ sort={Shock-Free Airfoil},
  name={Shock-Free Airfoil},
  description={\textcolor{catAD}{\textbf{[Aerodynamics]}}},
}

\newglossaryentry{Natural-Laminar-Flow}{ sort={Natural Laminar Flow},
  name={Natural Laminar Flow},
  description={\textcolor{catAD}{\textbf{[Aerodynamics]}}},
}

\newglossaryentry{Laminar-Flow-Control}{ sort={Laminar Flow Control},
  name={Laminar Flow Control},
  description={\textcolor{catAD}{\textbf{[Aerodynamics]}}},
}

\newglossaryentry{Riblet}{ sort={Riblet},
  name={Riblet},
  description={\textcolor{catAD}{\textbf{[Aerodynamics]}}},
}

\newglossaryentry{Dimple}{ sort={Dimple},
  name={Dimple},
  description={\textcolor{catAD}{\textbf{[Aerodynamics]}}},
}

\newglossaryentry{Base-Drag}{ sort={Base Drag},
  name={Base Drag},
  description={\textcolor{catAD}{\textbf{[Aerodynamics]}}},
}

\newglossaryentry{Boattail}{ sort={Boattail},
  name={Boattail},
  description={\textcolor{catAD}{\textbf{[Aerodynamics]}}},
}

\newglossaryentry{Bump}{ sort={Bump},
  name={Bump},
  description={\textcolor{catAD}{\textbf{[Aerodynamics]}}},
}

\newglossaryentry{Fence}{ sort={Fence},
  name={Fence},
  description={\textcolor{catAD}{\textbf{[Aerodynamics]}}},
}

\newglossaryentry{Wing-Fence}{ sort={Wing Fence},
  name={Wing Fence},
  description={\textcolor{catAD}{\textbf{[Aerodynamics]}}},
}

\newglossaryentry{Dogtooth}{ sort={Dogtooth},
  name={Dogtooth},
  description={\textcolor{catAD}{\textbf{[Aerodynamics]}}},
}

\newglossaryentry{Leading-Edge-Extension}{ sort={Leading Edge Extension},
  name={Leading Edge Extension},
  description={\textcolor{catAD}{\textbf{[Aerodynamics]}}},
}

\newglossaryentry{LEX-Fence}{ sort={LEX Fence},
  name={LEX Fence},
  description={\textcolor{catAD}{\textbf{[Aerodynamics]}}},
}

\newglossaryentry{Strakes}{ sort={Strakes},
  name={Strakes},
  description={\textcolor{catAD}{\textbf{[Aerodynamics]}}},
}

\newglossaryentry{Chine}{ sort={Chine},
  name={Chine},
  description={\textcolor{catAD}{\textbf{[Aerodynamics]}}},
}

\newglossaryentry{Belly-Flap}{ sort={Belly Flap},
  name={Belly Flap},
  description={\textcolor{catAD}{\textbf{[Aerodynamics]}}},
}

\newglossaryentry{Deceleron}{ sort={Deceleron},
  name={Deceleron},
  description={\textcolor{catAD}{\textbf{[Aerodynamics]}}},
}

\newglossaryentry{Flapperon}{ sort={Flapperon},
  name={Flapperon},
  description={\textcolor{catAD}{\textbf{[Aerodynamics]}}},
}

\newglossaryentry{Taileron}{ sort={Taileron},
  name={Taileron},
  description={\textcolor{catAD}{\textbf{[Aerodynamics]}}},
}

\newglossaryentry{Stabilator}{ sort={Stabilator},
  name={Stabilator},
  description={\textcolor{catAD}{\textbf{[Aerodynamics]}}},
}

\newglossaryentry{V-Tail}{ sort={V-Tail},
  name={V-Tail},
  description={\textcolor{catAD}{\textbf{[Aerodynamics]}}},
}

\newglossaryentry{Cruciform-Tail}{ sort={Cruciform Tail},
  name={Cruciform Tail},
  description={\textcolor{catAD}{\textbf{[Aerodynamics]}}},
}

\newglossaryentry{T-Tail}{ sort={T-Tail},
  name={T-Tail},
  description={\textcolor{catAD}{\textbf{[Aerodynamics]}}},
}

\newglossaryentry{H-Tail}{ sort={H-Tail},
  name={H-Tail},
  description={\textcolor{catAD}{\textbf{[Aerodynamics]}}},
}

\newglossaryentry{Twin-Tail}{ sort={Twin Tail},
  name={Twin Tail},
  description={\textcolor{catAD}{\textbf{[Aerodynamics]}}},
}

\newglossaryentry{Endplate}{ sort={Endplate},
  name={Endplate},
  description={\textcolor{catAD}{\textbf{[Aerodynamics]}}},
}

\newglossaryentry{Hoerner-Tip}{ sort={Hoerner Tip},
  name={Hoerner Tip},
  description={\textcolor{catAD}{\textbf{[Aerodynamics]}}},
}

\newglossaryentry{Raked-Wingtip}{ sort={Raked Wingtip},
  name={Raked Wingtip},
  description={\textcolor{catAD}{\textbf{[Aerodynamics]}}},
}

\newglossaryentry{Tip-Fence}{ sort={Tip Fence},
  name={Tip Fence},
  description={\textcolor{catAD}{\textbf{[Aerodynamics]}}},
}

\newglossaryentry{Wingtip-Vortex}{ sort={Wingtip Vortex},
  name={Wingtip Vortex},
  description={\textcolor{catAD}{\textbf{[Aerodynamics]}}},
}

\newglossaryentry{Contrail}{ sort={Contrail},
  name={Contrail},
  description={\textcolor{catAD}{\textbf{[Aerodynamics]}}},
}

\newglossaryentry{Vapor-Cone}{ sort={Vapor Cone},
  name={Vapor Cone},
  description={\textcolor{catAD}{\textbf{[Aerodynamics]}}},
}

\newglossaryentry{Shock-Diamond}{ sort={Shock Diamond},
  name={Shock Diamond},
  description={\textcolor{catAD}{\textbf{[Aerodynamics]}}},
}

\newglossaryentry{Mach-Wave}{ sort={Mach Wave},
  name={Mach Wave},
  description={\textcolor{catAD}{\textbf{[Aerodynamics]}}},
}

\newglossaryentry{Bow-Shock}{ sort={Bow Shock},
  name={Bow Shock},
  description={\textcolor{catAD}{\textbf{[Aerodynamics]}}},
}

\newglossaryentry{Attached-Shock}{ sort={Attached Shock},
  name={Attached Shock},
  description={\textcolor{catAD}{\textbf{[Aerodynamics]}}},
}

\newglossaryentry{Lambda-Shock}{ sort={Lambda Shock},
  name={Lambda Shock},
  description={\textcolor{catAD}{\textbf{[Aerodynamics]}}},
}

\newglossaryentry{Fish-Tail-Shock}{ sort={Fish-Tail Shock},
  name={Fish-Tail Shock},
  description={\textcolor{catAD}{\textbf{[Aerodynamics]}}},
}

\newglossaryentry{Detached-Shock}{ sort={Detached Shock},
  name={Detached Shock},
  description={\textcolor{catAD}{\textbf{[Aerodynamics]}}},
}

\newglossaryentry{Expansion-Fan}{ sort={Expansion Fan},
  name={Expansion Fan},
  description={\textcolor{catAD}{\textbf{[Aerodynamics]}}},
}

\newglossaryentry{Prandtl-Meyer-Function}{ sort={Prandtl-Meyer Function},
  name={Prandtl-Meyer Function},
  description={\textcolor{catAD}{\textbf{[Aerodynamics]}}},
}

\newglossaryentry{Method-of-Characteristics}{ sort={Method of Characteristics},
  name={Method of Characteristics},
  description={\textcolor{catAD}{\textbf{[Aerodynamics]}}},
}

\newglossaryentry{Supersonic-Wind-Tunnel}{ sort={Supersonic Wind Tunnel},
  name={Supersonic Wind Tunnel},
  description={\textcolor{catAD}{\textbf{[Aerodynamics]}}},
}

\newglossaryentry{Hypersonic-Wind-Tunnel}{ sort={Hypersonic Wind Tunnel},
  name={Hypersonic Wind Tunnel},
  description={\textcolor{catAD}{\textbf{[Aerodynamics]}}},
}

\newglossaryentry{Shock-Tunnel}{ sort={Shock Tunnel},
  name={Shock Tunnel},
  description={\textcolor{catAD}{\textbf{[Aerodynamics]}}},
}

\newglossaryentry{Arc-Heated-Wind-Tunnel}{ sort={Arc-Heated Wind Tunnel},
  name={Arc-Heated Wind Tunnel},
  description={\textcolor{catAD}{\textbf{[Aerodynamics]}}},
}

\newglossaryentry{Cryogenic-Wind-Tunnel}{ sort={Cryogenic Wind Tunnel},
  name={Cryogenic Wind Tunnel},
  description={\textcolor{catAD}{\textbf{[Aerodynamics]}}},
}

\newglossaryentry{Indraft-Wind-Tunnel}{ sort={Indraft Wind Tunnel},
  name={Indraft Wind Tunnel},
  description={\textcolor{catAD}{\textbf{[Aerodynamics]}}},
}

\newglossaryentry{Open-Circuit-Tunnel}{ sort={Open Circuit Tunnel},
  name={Open Circuit Tunnel},
  description={\textcolor{catAD}{\textbf{[Aerodynamics]}}},
}

\newglossaryentry{Closed-Circuit-Tunnel}{ sort={Closed Circuit Tunnel},
  name={Closed Circuit Tunnel},
  description={\textcolor{catAD}{\textbf{[Aerodynamics]}}},
}

\newglossaryentry{Low-Speed-Tunnel}{ sort={Low Speed Tunnel},
  name={Low Speed Tunnel},
  description={\textcolor{catAD}{\textbf{[Aerodynamics]}}},
}

\newglossaryentry{High-Speed-Tunnel}{ sort={High Speed Tunnel},
  name={High Speed Tunnel},
  description={\textcolor{catAD}{\textbf{[Aerodynamics]}}},
}

\newglossaryentry{Wall-Interference}{ sort={Wall Interference},
  name={Wall Interference},
  description={\textcolor{catAD}{\textbf{[Aerodynamics]}}},
}

\newglossaryentry{Blockage}{ sort={Blockage},
  name={Blockage},
  description={\textcolor{catAD}{\textbf{[Aerodynamics]}}},
}

\newglossaryentry{Turbulence-Intensity}{ sort={Turbulence Intensity},
  name={Turbulence Intensity},
  description={\textcolor{catAD}{\textbf{[Aerodynamics]}}},
}

\newglossaryentry{Flow-Angularity}{ sort={Flow Angularity},
  name={Flow Angularity},
  description={\textcolor{catAD}{\textbf{[Aerodynamics]}}},
}

\newglossaryentry{Sting-Mount}{ sort={Sting Mount},
  name={Sting Mount},
  description={\textcolor{catAD}{\textbf{[Aerodynamics]}}},
}

\newglossaryentry{Strut-Mount}{ sort={Strut Mount},
  name={Strut Mount},
  description={\textcolor{catAD}{\textbf{[Aerodynamics]}}},
}

\newglossaryentry{Magnetic-Suspension}{ sort={Magnetic Suspension},
  name={Magnetic Suspension},
  description={\textcolor{catAD}{\textbf{[Aerodynamics]}}},
}

\newglossaryentry{Pressure-Sensitive-Paint}{ sort={Pressure Sensitive Paint},
  name={Pressure Sensitive Paint},
  description={\textcolor{catAD}{\textbf{[Aerodynamics]}}},
}

\newglossaryentry{Temperature-Sensitive-Paint}{ sort={Temperature Sensitive Paint},
  name={Temperature Sensitive Paint},
  description={\textcolor{catAD}{\textbf{[Aerodynamics]}}},
}

\newglossaryentry{Infrared-Thermography}{ sort={Infrared Thermography},
  name={Infrared Thermography},
  description={\textcolor{catAD}{\textbf{[Aerodynamics]}}},
}

\newglossaryentry{Optical-Strain-Measurement}{ sort={Optical Strain Measurement},
  name={Optical Strain Measurement},
  description={\textcolor{catAD}{\textbf{[Aerodynamics]}}},
}


\newglossaryentry{Strain-Gauge}{ sort={Strain Gauge},
  name={Strain Gauge},
  description={\textcolor{catAD}{\textbf{[Aerodynamics]}}},
}

\newglossaryentry{Thermocouple}{ sort={Thermocouple},
  name={Thermocouple},
  description={\textcolor{catAD}{\textbf{[Aerodynamics]}}},
}


\newglossaryentry{External-Balance}{ sort={External Balance},
  name={External Balance},
  description={\textcolor{catAD}{\textbf{[Aerodynamics]}}},
}


\newglossaryentry{Half-Model-Testing}{ sort={Half Model Testing},
  name={Half Model Testing},
  description={\textcolor{catAD}{\textbf{[Aerodynamics]}}},
}

\newglossaryentry{Full-Model-Testing}{ sort={Full Model Testing},
  name={Full Model Testing},
  description={\textcolor{catAD}{\textbf{[Aerodynamics]}}},
}

\newglossaryentry{Dynamic-Testing}{ sort={Dynamic Testing},
  name={Dynamic Testing},
  description={\textcolor{catAD}{\textbf{[Aerodynamics]}}},
}

\newglossaryentry{Static-Testing}{ sort={Static Testing},
  name={Static Testing},
  description={\textcolor{catAD}{\textbf{[Aerodynamics]}}},
}

\newglossaryentry{Oscillatory-Testing}{ sort={Oscillatory Testing},
  name={Oscillatory Testing},
  description={\textcolor{catAD}{\textbf{[Aerodynamics]}}},
}

\newglossaryentry{Grid-Turbulence}{ sort={Grid Turbulence},
  name={Grid Turbulence},
  description={\textcolor{catAD}{\textbf{[Aerodynamics]}}},
}

\newglossaryentry{Hot-Wire-Anemometry}{ sort={Hot Wire Anemometry},
  name={Hot Wire Anemometry},
  description={\textcolor{catAD}{\textbf{[Aerodynamics]}}},
}

\newglossaryentry{Laser-Doppler-Velocimetry}{ sort={Laser Doppler Velocimetry},
  name={Laser Doppler Velocimetry},
  description={\textcolor{catAD}{\textbf{[Aerodynamics]}}},
}

\newglossaryentry{Phase-Doppler-Anemometry}{ sort={Phase Doppler Anemometry},
  name={Phase Doppler Anemometry},
  description={\textcolor{catAD}{\textbf{[Aerodynamics]}}},
}

\newglossaryentry{Pressure-Scanner}{ sort={Pressure Scanner},
  name={Pressure Scanner},
  description={\textcolor{catAD}{\textbf{[Aerodynamics]}}},
}

\newglossaryentry{Wind-Tunnel-Fan}{ sort={Wind Tunnel Fan},
  name={Wind Tunnel Fan},
  description={\textcolor{catAD}{\textbf{[Aerodynamics]}}},
}

\newglossaryentry{Honeycomb-Straightener}{ sort={Honeycomb Straightener},
  name={Honeycomb Straightener},
  description={\textcolor{catAD}{\textbf{[Aerodynamics]}}},
}

\newglossaryentry{Settling-Chamber}{ sort={Settling Chamber},
  name={Settling Chamber},
  description={\textcolor{catAD}{\textbf{[Aerodynamics]}}},
}

\newglossaryentry{Contraction-Cone}{ sort={Contraction Cone},
  name={Contraction Cone},
  description={\textcolor{catAD}{\textbf{[Aerodynamics]}}},
}

\newglossaryentry{Diffuser-Section}{ sort={Diffuser Section},
  name={Diffuser Section},
  description={\textcolor{catAD}{\textbf{[Aerodynamics]}}},
}

\newglossaryentry{Corner-Vanes}{ sort={Corner Vanes},
  name={Corner Vanes},
  description={\textcolor{catAD}{\textbf{[Aerodynamics]}}},
}

\newglossaryentry{Turning-Vanes}{ sort={Turning Vanes},
  name={Turning Vanes},
  description={\textcolor{catAD}{\textbf{[Aerodynamics]}}},
}

\newglossaryentry{Screens}{ sort={Screens},
  name={Screens},
  description={\textcolor{catAD}{\textbf{[Aerodynamics]}}},
}

\newglossaryentry{Acoustic-Testing}{ sort={Acoustic Testing},
  name={Acoustic Testing},
  description={\textcolor{catAD}{\textbf{[Aerodynamics]}}},
}

\newglossaryentry{Aeroacoustics}{ sort={Aeroacoustics},
  name={Aeroacoustics},
  description={\textcolor{catAD}{\textbf{[Aerodynamics]}}},
}

\newglossaryentry{Jet-Noise}{ sort={Jet Noise},
  name={Jet Noise},
  description={\textcolor{catAD}{\textbf{[Aerodynamics]}}},
}

\newglossaryentry{Airframe-Noise}{ sort={Airframe Noise},
  name={Airframe Noise},
  description={\textcolor{catAD}{\textbf{[Aerodynamics]}}},
}

\newglossaryentry{Propeller-Noise}{ sort={Propeller Noise},
  name={Propeller Noise},
  description={\textcolor{catAD}{\textbf{[Aerodynamics]}}},
}

\newglossaryentry{Rotor-Noise}{ sort={Rotor Noise},
  name={Rotor Noise},
  description={\textcolor{catAD}{\textbf{[Aerodynamics]}}},
}

\newglossaryentry{Fan-Noise}{ sort={Fan Noise},
  name={Fan Noise},
  description={\textcolor{catAD}{\textbf{[Aerodynamics]}}},
}

\newglossaryentry{Combustion-Noise}{ sort={Combustion Noise},
  name={Combustion Noise},
  description={\textcolor{catAD}{\textbf{[Aerodynamics]}}},
}

\newglossaryentry{Noise-Certification}{ sort={Noise Certification},
  name={Noise Certification},
  description={\textcolor{catAD}{\textbf{[Aerodynamics]}}},
}

\newglossaryentry{Noise-Reduction}{ sort={Noise Reduction},
  name={Noise Reduction},
  description={\textcolor{catAD}{\textbf{[Aerodynamics]}}},
}

\newglossaryentry{Chevron-Nozzle}{ sort={Chevron Nozzle},
  name={Chevron Nozzle},
  description={\textcolor{catAD}{\textbf{[Aerodynamics]}}},
}

\newglossaryentry{Hush-Kit}{ sort={Hush Kit},
  name={Hush Kit},
  description={\textcolor{catAD}{\textbf{[Aerodynamics]}}},
}

\newglossaryentry{Emission-Index}{ sort={Emission Index},
  name={Emission Index},
  description={\textcolor{catAD}{\textbf{[Aerodynamics]}}},
}

\newglossaryentry{Carbon-Footprint}{ sort={Carbon Footprint},
  name={Carbon Footprint},
  description={\textcolor{catAD}{\textbf{[Aerodynamics]}}},
}


\newglossaryentry{Hydrogen-Propulsion}{ sort={Hydrogen Propulsion},
  name={Hydrogen Propulsion},
  description={\textcolor{catAD}{\textbf{[Aerodynamics]}}},
}

\newglossaryentry{Electric-Propulsion}{ sort={Electric Propulsion},
  name={Electric Propulsion},
  description={\textcolor{catAD}{\textbf{[Aerodynamics]}}},
}

\newglossaryentry{Hybrid-Propulsion}{ sort={Hybrid Propulsion},
  name={Hybrid Propulsion},
  description={\textcolor{catAD}{\textbf{[Aerodynamics]}}},
}

\newglossaryentry{Solar-Powered-Flight}{ sort={Solar Powered Flight},
  name={Solar Powered Flight},
  description={\textcolor{catAD}{\textbf{[Aerodynamics]}}},
}

\newglossaryentry{High-Altitude-Flight}{ sort={High Altitude Flight},
  name={High Altitude Flight},
  description={\textcolor{catAD}{\textbf{[Aerodynamics]}}},
}

\newglossaryentry{Stratospheric-Flight}{ sort={Stratospheric Flight},
  name={Stratospheric Flight},
  description={\textcolor{catAD}{\textbf{[Aerodynamics]}}},
}

\newglossaryentry{Suborbital-Flight}{ sort={Suborbital Flight},
  name={Suborbital Flight},
  description={\textcolor{catAD}{\textbf{[Aerodynamics]}}},
}

\newglossaryentry{Atmospheric-Re-entry}{ sort={Atmospheric Re-entry},
  name={Atmospheric Re-entry},
  description={\textcolor{catAD}{\textbf{[Aerodynamics]}}},
}

\newglossaryentry{Re-entry-Corridor}{ sort={Re-entry Corridor},
  name={Re-entry Corridor},
  description={\textcolor{catAD}{\textbf{[Aerodynamics]}}},
}

\newglossaryentry{Re-entry-Angle}{ sort={Re-entry Angle},
  name={Re-entry Angle},
  description={\textcolor{catAD}{\textbf{[Aerodynamics]}}},
}

\newglossaryentry{Skip-Re-entry}{ sort={Skip Re-entry},
  name={Skip Re-entry},
  description={\textcolor{catAD}{\textbf{[Aerodynamics]}}},
}

\newglossaryentry{Ballistic-Re-entry}{ sort={Ballistic Re-entry},
  name={Ballistic Re-entry},
  description={\textcolor{catAD}{\textbf{[Aerodynamics]}}},
}

\newglossaryentry{Lifting-Re-entry}{ sort={Lifting Re-entry},
  name={Lifting Re-entry},
  description={\textcolor{catAD}{\textbf{[Aerodynamics]}}},
}

\newglossaryentry{Re-entry-Velocity}{ sort={Re-entry Velocity},
  name={Re-entry Velocity},
  description={\textcolor{catAD}{\textbf{[Aerodynamics]}}},
}

\newglossaryentry{Re-entry-Heating-Rate}{ sort={Re-entry Heating Rate},
  name={Re-entry Heating Rate},
  description={\textcolor{catAD}{\textbf{[Aerodynamics]}}},
}

\newglossaryentry{Peak-Heating}{ sort={Peak Heating},
  name={Peak Heating},
  description={\textcolor{catAD}{\textbf{[Aerodynamics]}}},
}

\newglossaryentry{Heat-Flux}{ sort={Heat Flux},
  name={Heat Flux},
  description={\textcolor{catAD}{\textbf{[Aerodynamics]}}},
}

\newglossaryentry{Convective-Heating}{ sort={Convective Heating},
  name={Convective Heating},
  description={\textcolor{catAD}{\textbf{[Aerodynamics]}}},
}

\newglossaryentry{Radiative-Heating}{ sort={Radiative Heating},
  name={Radiative Heating},
  description={\textcolor{catAD}{\textbf{[Aerodynamics]}}},
}

\newglossaryentry{Catalytic-Heating}{ sort={Catalytic Heating},
  name={Catalytic Heating},
  description={\textcolor{catAD}{\textbf{[Aerodynamics]}}},
}

\newglossaryentry{Thermal-Protection-Material}{ sort={Thermal Protection Material},
  name={Thermal Protection Material},
  description={\textcolor{catAD}{\textbf{[Aerodynamics]}}},
}

\newglossaryentry{Ablative-TPS}{ sort={Ablative TPS},
  name={Ablative TPS},
  description={\textcolor{catAD}{\textbf{[Aerodynamics]}}},
}

\newglossaryentry{Reusable-TPS}{ sort={Reusable TPS},
  name={Reusable TPS},
  description={\textcolor{catAD}{\textbf{[Aerodynamics]}}},
}


\newglossaryentry{Thermal-Barrier-Coating}{ sort={Thermal Barrier Coating},
  name={Thermal Barrier Coating},
  description={\textcolor{catAD}{\textbf{[Aerodynamics]}}},
}

\newglossaryentry{Hot-Structure}{ sort={Hot Structure},
  name={Hot Structure},
  description={\textcolor{catAD}{\textbf{[Aerodynamics]}}},
}

\newglossaryentry{Cold-Structure}{ sort={Cold Structure},
  name={Cold Structure},
  description={\textcolor{catAD}{\textbf{[Aerodynamics]}}},
}

\newglossaryentry{Nose-Cap}{ sort={Nose Cap},
  name={Nose Cap},
  description={\textcolor{catAD}{\textbf{[Aerodynamics]}}},
}

\newglossaryentry{Leading-Edge-TPS}{ sort={Leading Edge TPS},
  name={Leading Edge TPS},
  description={\textcolor{catAD}{\textbf{[Aerodynamics]}}},
}

\newglossaryentry{Tiles}{ sort={Tiles},
  name={Tiles},
  description={\textcolor{catAD}{\textbf{[Aerodynamics]}}},
}

\newglossaryentry{Reinforced-Carbon-Carbon}{ sort={Reinforced Carbon Carbon},
  name={Reinforced Carbon Carbon},
  description={\textcolor{catAD}{\textbf{[Aerodynamics]}}},
}

\newglossaryentry{Flexible-TPS}{ sort={Flexible TPS},
  name={Flexible TPS},
  description={\textcolor{catAD}{\textbf{[Aerodynamics]}}},
}

\newglossaryentry{Inflatable-TPS}{ sort={Inflatable TPS},
  name={Inflatable TPS},
  description={\textcolor{catAD}{\textbf{[Aerodynamics]}}},
}

\newglossaryentry{Aeroshell}{ sort={Aeroshell},
  name={Aeroshell},
  description={\textcolor{catAD}{\textbf{[Aerodynamics]}}},
}


\newglossaryentry{Buoyancy}{ sort={Buoyancy},
  name={Buoyancy},
  description={\textcolor{catAD}{\textbf{[Aerodynamics]}}},
}

\newglossaryentry{Clean-Configuration}{ sort={Clean Configuration},
  name={Clean Configuration},
  description={\textcolor{catAD}{\textbf{[Aerodynamics]}}},
}

\newglossaryentry{Decalage}{ sort={Decalage},
  name={Decalage},
  description={\textcolor{catAD}{\textbf{[Aerodynamics]}}},
}

\newglossaryentry{Disk-Loading}{ sort={Disk Loading},
  name={Disk Loading},
  description={\textcolor{catAD}{\textbf{[Aerodynamics]}}},
}

\newglossaryentry{Enstrophy}{ sort={Enstrophy},
  name={Enstrophy},
  description={\textcolor{catAD}{\textbf{[Aerodynamics]}}},
}

\newglossaryentry{Keel-Effect}{ sort={Keel Effect},
  name={Keel Effect},
  description={\textcolor{catAD}{\textbf{[Aerodynamics]}}},
}

\newglossaryentry{Rogallo-Wing}{ sort={Rogallo Wing},
  name={Rogallo Wing},
  description={\textcolor{catAD}{\textbf{[Aerodynamics]}}},
}

\newglossaryentry{Viscosity}{ sort={Viscosity},
  name={Viscosity},
  description={\textcolor{catAD}{\textbf{[Aerodynamics]}}},
}

\newglossaryentry{Aerospace-Engineering}{ sort={Aerospace Engineering},
  name={Aerospace Engineering},
  description={\textcolor{catAE}{\textbf{[Aerospace]}}},
}

\newglossaryentry{Aircraft}{ sort={Aircraft},
  name={Aircraft},
  description={\textcolor{catAE}{\textbf{[Aerospace]}}},
}

\newglossaryentry{Airplane}{ sort={Airplane},
  name={Airplane},
  description={\textcolor{catAE}{\textbf{[Aerospace]}}},
}

\newglossaryentry{Helicopter}{ sort={Helicopter},
  name={Helicopter},
  description={\textcolor{catAE}{\textbf{[Aerospace]}}},
}

\newglossaryentry{Rotorcraft}{ sort={Rotorcraft},
  name={Rotorcraft},
  description={\textcolor{catAE}{\textbf{[Aerospace]}}},
}

\newglossaryentry{Tiltrotor}{ sort={Tiltrotor},
  name={Tiltrotor},
  description={\textcolor{catAE}{\textbf{[Aerospace]}}},
}

\newglossaryentry{Gyrocopter}{ sort={Gyrocopter},
  name={Gyrocopter},
  description={\textcolor{catAE}{\textbf{[Aerospace]}}},
}

\newglossaryentry{Glider}{ sort={Glider},
  name={Glider},
  description={\textcolor{catAE}{\textbf{[Aerospace]}}},
}



\newglossaryentry{eVTOL}{ sort={eVTOL},
  name={eVTOL},
  description={\textcolor{catAE}{\textbf{[Aerospace]}}},
}

\newglossaryentry{Airship}{ sort={Airship},
  name={Airship},
  description={\textcolor{catAE}{\textbf{[Aerospace]}}},
}

\newglossaryentry{Balloon}{ sort={Balloon},
  name={Balloon},
  description={\textcolor{catAE}{\textbf{[Aerospace]}}},
}





\newglossaryentry{Spaceplane}{ sort={Spaceplane},
  name={Spaceplane},
  description={\textcolor{catAE}{\textbf{[Aerospace]}}},
}

\newglossaryentry{Hypersonic-Vehicle}{ sort={Hypersonic Vehicle},
  name={Hypersonic Vehicle},
  description={\textcolor{catAE}{\textbf{[Aerospace]}}},
}

\newglossaryentry{Flight-Mechanics}{ sort={Flight Mechanics},
  name={Flight Mechanics},
  description={\textcolor{catAE}{\textbf{[Aerospace]}}},
}

\newglossaryentry{Sideslip-Angle}{ sort={Sideslip Angle},
  name={Sideslip Angle},
  description={\textcolor{catAE}{\textbf{[Aerospace]}}},
}

\newglossaryentry{Flight-Path}{ sort={Flight Path},
  name={Flight Path},
  description={\textcolor{catAE}{\textbf{[Aerospace]}}},
}

\newglossaryentry{Glide-Ratio}{ sort={Glide Ratio},
  name={Glide Ratio},
  description={\textcolor{catAE}{\textbf{[Aerospace]}}},
}

\newglossaryentry{Ceiling}{ sort={Ceiling},
  name={Ceiling},
  description={\textcolor{catAE}{\textbf{[Aerospace]}}},
}

\newglossaryentry{Stall-Recovery}{ sort={Stall Recovery},
  name={Stall Recovery},
  description={\textcolor{catAE}{\textbf{[Aerospace]}}},
}

\newglossaryentry{Spin-Recovery}{ sort={Spin Recovery},
  name={Spin Recovery},
  description={\textcolor{catAE}{\textbf{[Aerospace]}}},
}

\newglossaryentry{Downwash}{ sort={Downwash},
  name={Downwash},
  description={\textcolor{catAE}{\textbf{[Aerospace]}}},
}

\newglossaryentry{Flow-Separation}{ sort={Flow Separation},
  name={Flow Separation},
  description={\textcolor{catAE}{\textbf{[Aerospace]}}},
}

\newglossaryentry{Fuselage}{ sort={Fuselage},
  name={Fuselage},
  description={\textcolor{catAE}{\textbf{[Aerospace]}}},
}

\newglossaryentry{Wing-Spar}{ sort={Wing Spar},
  name={Wing Spar},
  description={\textcolor{catAE}{\textbf{[Aerospace]}}},
}

\newglossaryentry{Rib}{ sort={Rib},
  name={Rib},
  description={\textcolor{catAE}{\textbf{[Aerospace]}}},
}

\newglossaryentry{Stringer}{ sort={Stringer},
  name={Stringer},
  description={\textcolor{catAE}{\textbf{[Aerospace]}}},
}

\newglossaryentry{Bulkhead}{ sort={Bulkhead},
  name={Bulkhead},
  description={\textcolor{catAE}{\textbf{[Aerospace]}}},
}

\newglossaryentry{Longeron}{ sort={Longeron},
  name={Longeron},
  description={\textcolor{catAE}{\textbf{[Aerospace]}}},
}

\newglossaryentry{Empennage}{ sort={Empennage},
  name={Empennage},
  description={\textcolor{catAE}{\textbf{[Aerospace]}}},
}

\newglossaryentry{Vertical-Stabilizer}{ sort={Vertical Stabilizer},
  name={Vertical Stabilizer},
  description={\textcolor{catAE}{\textbf{[Aerospace]}}},
}

\newglossaryentry{Horizontal-Stabilizer}{ sort={Horizontal Stabilizer},
  name={Horizontal Stabilizer},
  description={\textcolor{catAE}{\textbf{[Aerospace]}}},
}

\newglossaryentry{Fin}{ sort={Fin},
  name={Fin},
  description={\textcolor{catAE}{\textbf{[Aerospace]}}},
}

\newglossaryentry{Skin}{ sort={Skin},
  name={Skin},
  description={\textcolor{catAE}{\textbf{[Aerospace]}}},
}


\newglossaryentry{Monocoque}{ sort={Monocoque},
  name={Monocoque},
  description={\textcolor{catAE}{\textbf{[Aerospace]}}},
}

\newglossaryentry{Semi-Monocoque}{ sort={Semi-Monocoque},
  name={Semi-Monocoque},
  description={\textcolor{catAE}{\textbf{[Aerospace]}}},
}

\newglossaryentry{Aluminum-Alloy}{ sort={Aluminum Alloy},
  name={Aluminum Alloy},
  description={\textcolor{catAE}{\textbf{[Aerospace]}}},
}

\newglossaryentry{Titanium-Alloy}{ sort={Titanium Alloy},
  name={Titanium Alloy},
  description={\textcolor{catAE}{\textbf{[Aerospace]}}},
}

\newglossaryentry{Carbon-Fiber}{ sort={Carbon Fiber},
  name={Carbon Fiber},
  description={\textcolor{catAE}{\textbf{[Aerospace]}}},
}

\newglossaryentry{Fiberglass}{ sort={Fiberglass},
  name={Fiberglass},
  description={\textcolor{catAE}{\textbf{[Aerospace]}}},
}

\newglossaryentry{Kevlar}{ sort={Kevlar},
  name={Kevlar},
  description={\textcolor{catAE}{\textbf{[Aerospace]}}},
}

\newglossaryentry{Composite-Material}{ sort={Composite Material},
  name={Composite Material},
  description={\textcolor{catAE}{\textbf{[Aerospace]}}},
}

\newglossaryentry{Ceramic-Matrix-Composite}{ sort={Ceramic Matrix Composite},
  name={Ceramic Matrix Composite},
  description={\textcolor{catAE}{\textbf{[Aerospace]}}},
}

\newglossaryentry{Superalloy}{ sort={Superalloy},
  name={Superalloy},
  description={\textcolor{catAE}{\textbf{[Aerospace]}}},
}

\newglossaryentry{Shape-Memory-Alloy}{ sort={Shape Memory Alloy},
  name={Shape Memory Alloy},
  description={\textcolor{catAE}{\textbf{[Aerospace]}}},
}

\newglossaryentry{Ablative-Material}{ sort={Ablative Material},
  name={Ablative Material},
  description={\textcolor{catAE}{\textbf{[Aerospace]}}},
}

\newglossaryentry{Hydraulic-System}{ sort={Hydraulic System},
  name={Hydraulic System},
  description={\textcolor{catAE}{\textbf{[Aerospace]}}},
}

\newglossaryentry{Pneumatic-System}{ sort={Pneumatic System},
  name={Pneumatic System},
  description={\textcolor{catAE}{\textbf{[Aerospace]}}},
}

\newglossaryentry{Electrical-System}{ sort={Electrical System},
  name={Electrical System},
  description={\textcolor{catAE}{\textbf{[Aerospace]}}},
}

\newglossaryentry{Fuel-System}{ sort={Fuel System},
  name={Fuel System},
  description={\textcolor{catAE}{\textbf{[Aerospace]}}},
}

\newglossaryentry{Environmental-Control-System}{ sort={Environmental Control System},
  name={Environmental Control System},
  description={\textcolor{catAE}{\textbf{[Aerospace]}}},
}

\newglossaryentry{Cabin-Pressurization}{ sort={Cabin Pressurization},
  name={Cabin Pressurization},
  description={\textcolor{catAE}{\textbf{[Aerospace]}}},
}

\newglossaryentry{Landing-Gear}{ sort={Landing Gear},
  name={Landing Gear},
  description={\textcolor{catAE}{\textbf{[Aerospace]}}},
}

\newglossaryentry{Braking-System}{ sort={Braking System},
  name={Braking System},
  description={\textcolor{catAE}{\textbf{[Aerospace]}}},
}

\newglossaryentry{Anti-Skid-System}{ sort={Anti-Skid System},
  name={Anti-Skid System},
  description={\textcolor{catAE}{\textbf{[Aerospace]}}},
}

\newglossaryentry{Ice-Protection-System}{ sort={Ice Protection System},
  name={Ice Protection System},
  description={\textcolor{catAE}{\textbf{[Aerospace]}}},
}

\newglossaryentry{Flight-Computer}{ sort={Flight Computer},
  name={Flight Computer},
  description={\textcolor{catAE}{\textbf{[Aerospace]}}},
}

\newglossaryentry{Flight-Management-System}{ sort={Flight Management System},
  name={Flight Management System},
  description={\textcolor{catAE}{\textbf{[Aerospace]}}},
}

\newglossaryentry{Autopilot}{ sort={Autopilot},
  name={Autopilot},
  description={\textcolor{catAE}{\textbf{[Aerospace]}}},
}

\newglossaryentry{Fly-by-Wire}{ sort={Fly-by-Wire},
  name={Fly-by-Wire},
  description={\textcolor{catAE}{\textbf{[Aerospace]}}},
}

\newglossaryentry{Glass-Cockpit}{ sort={Glass Cockpit},
  name={Glass Cockpit},
  description={\textcolor{catAE}{\textbf{[Aerospace]}}},
}

\newglossaryentry{Head-Up-Display}{ sort={Head-Up Display},
  name={Head-Up Display},
  description={\textcolor{catAE}{\textbf{[Aerospace]}}},
}

\newglossaryentry{Multi-Function-Display}{ sort={Multi-Function Display},
  name={Multi-Function Display},
  description={\textcolor{catAE}{\textbf{[Aerospace]}}},
}

\newglossaryentry{Air-Data-Computer}{ sort={Air Data Computer},
  name={Air Data Computer},
  description={\textcolor{catAE}{\textbf{[Aerospace]}}},
}

\newglossaryentry{GPS}{ sort={GPS},
  name={GPS},
  description={\textcolor{catAE}{\textbf{[Aerospace]}}},
}

\newglossaryentry{GNSS}{ sort={GNSS},
  name={GNSS},
  description={\textcolor{catAE}{\textbf{[Aerospace]}}},
}

\newglossaryentry{INS}{ sort={INS},
  name={INS},
  description={\textcolor{catAE}{\textbf{[Aerospace]}}},
}

\newglossaryentry{Dead-Reckoning}{ sort={Dead Reckoning},
  name={Dead Reckoning},
  description={\textcolor{catAE}{\textbf{[Aerospace]}}},
}

\newglossaryentry{Waypoint}{ sort={Waypoint},
  name={Waypoint},
  description={\textcolor{catAE}{\textbf{[Aerospace]}}},
}

\newglossaryentry{Great-Circle-Route}{ sort={Great Circle Route},
  name={Great Circle Route},
  description={\textcolor{catAE}{\textbf{[Aerospace]}}},
}

\newglossaryentry{Magnetic-Heading}{ sort={Magnetic Heading},
  name={Magnetic Heading},
  description={\textcolor{catAE}{\textbf{[Aerospace]}}},
}

\newglossaryentry{True-Heading}{ sort={True Heading},
  name={True Heading},
  description={\textcolor{catAE}{\textbf{[Aerospace]}}},
}

\newglossaryentry{Indicated-Airspeed}{ sort={Indicated Airspeed},
  name={Indicated Airspeed},
  description={\textcolor{catAE}{\textbf{[Aerospace]}}},
}

\newglossaryentry{Calibrated-Airspeed}{ sort={Calibrated Airspeed},
  name={Calibrated Airspeed},
  description={\textcolor{catAE}{\textbf{[Aerospace]}}},
}

\newglossaryentry{True-Airspeed}{ sort={True Airspeed},
  name={True Airspeed},
  description={\textcolor{catAE}{\textbf{[Aerospace]}}},
}

\newglossaryentry{Ground-Speed}{ sort={Ground Speed},
  name={Ground Speed},
  description={\textcolor{catAE}{\textbf{[Aerospace]}}},
}

\newglossaryentry{Turbojet}{ sort={Turbojet},
  name={Turbojet},
  description={\textcolor{catAE}{\textbf{[Aerospace]}}},
}

\newglossaryentry{Turbofan}{ sort={Turbofan},
  name={Turbofan},
  description={\textcolor{catAE}{\textbf{[Aerospace]}}},
}

\newglossaryentry{Turboprop}{ sort={Turboprop},
  name={Turboprop},
  description={\textcolor{catAE}{\textbf{[Aerospace]}}},
}

\newglossaryentry{Turboshaft}{ sort={Turboshaft},
  name={Turboshaft},
  description={\textcolor{catAE}{\textbf{[Aerospace]}}},
}

\newglossaryentry{Pulsejet}{ sort={Pulsejet},
  name={Pulsejet},
  description={\textcolor{catAE}{\textbf{[Aerospace]}}},
}


\newglossaryentry{Afterburner}{ sort={Afterburner},
  name={Afterburner},
  description={\textcolor{catAE}{\textbf{[Aerospace]}}},
}



\newglossaryentry{Payload}{ sort={Payload},
  name={Payload},
  description={\textcolor{catAE}{\textbf{[Aerospace]}}},
}

\newglossaryentry{Fairing}{ sort={Fairing},
  name={Fairing},
  description={\textcolor{catAE}{\textbf{[Aerospace]}}},
}

\newglossaryentry{Booster}{ sort={Booster},
  name={Booster},
  description={\textcolor{catAE}{\textbf{[Aerospace]}}},
}


\newglossaryentry{Delta-v}{ sort={Delta-v},
  name={Delta-v},
  description={\textcolor{catAE}{\textbf{[Aerospace]}}},
}



\newglossaryentry{Staging}{ sort={Staging},
  name={Staging},
  description={\textcolor{catAE}{\textbf{[Aerospace]}}},
}

\newglossaryentry{Gravity-Turn}{ sort={Gravity Turn},
  name={Gravity Turn},
  description={\textcolor{catAE}{\textbf{[Aerospace]}}},
}

\newglossaryentry{Hot-Staging}{ sort={Hot Staging},
  name={Hot Staging},
  description={\textcolor{catAE}{\textbf{[Aerospace]}}},
}







\newglossaryentry{Docking-Port}{ sort={Docking Port},
  name={Docking Port},
  description={\textcolor{catAE}{\textbf{[Aerospace]}}},
}


\newglossaryentry{Heat-Shield}{ sort={Heat Shield},
  name={Heat Shield},
  description={\textcolor{catAE}{\textbf{[Aerospace]}}},
}

\newglossaryentry{Orbit}{ sort={Orbit},
  name={Orbit},
  description={\textcolor{catAE}{\textbf{[Aerospace]}}},
}

\newglossaryentry{Elliptical-Orbit}{ sort={Elliptical Orbit},
  name={Elliptical Orbit},
  description={\textcolor{catAE}{\textbf{[Aerospace]}}},
}

\newglossaryentry{Circular-Orbit}{ sort={Circular Orbit},
  name={Circular Orbit},
  description={\textcolor{catAE}{\textbf{[Aerospace]}}},
}






\newglossaryentry{Transfer-Orbit}{ sort={Transfer Orbit},
  name={Transfer Orbit},
  description={\textcolor{catAE}{\textbf{[Aerospace]}}},
}

\newglossaryentry{Hohmann-Transfer}{ sort={Hohmann Transfer},
  name={Hohmann Transfer},
  description={\textcolor{catAE}{\textbf{[Aerospace]}}},
}

\newglossaryentry{Orbital-Inclination}{ sort={Orbital Inclination},
  name={Orbital Inclination},
  description={\textcolor{catAE}{\textbf{[Aerospace]}}},
}



\newglossaryentry{Guidance}{ sort={Guidance},
  name={Guidance},
  description={\textcolor{catAE}{\textbf{[Aerospace]}}},
}

\newglossaryentry{Navigation}{ sort={Navigation},
  name={Navigation},
  description={\textcolor{catAE}{\textbf{[Aerospace]}}},
}

\newglossaryentry{Attitude-Control}{ sort={Attitude Control},
  name={Attitude Control},
  description={\textcolor{catAE}{\textbf{[Aerospace]}}},
}


\newglossaryentry{PID-Controller}{ sort={PID Controller},
  name={PID Controller},
  description={\textcolor{catAE}{\textbf{[Aerospace]}}},
}

\newglossaryentry{Kalman-Filter}{ sort={Kalman Filter},
  name={Kalman Filter},
  description={\textcolor{catAE}{\textbf{[Aerospace]}}},
}

\newglossaryentry{State-Estimation}{ sort={State Estimation},
  name={State Estimation},
  description={\textcolor{catAE}{\textbf{[Aerospace]}}},
}

\newglossaryentry{Sensor-Fusion}{ sort={Sensor Fusion},
  name={Sensor Fusion},
  description={\textcolor{catAE}{\textbf{[Aerospace]}}},
}

\newglossaryentry{Control-Law}{ sort={Control Law},
  name={Control Law},
  description={\textcolor{catAE}{\textbf{[Aerospace]}}},
}


\newglossaryentry{Space-Vacuum}{ sort={Space Vacuum},
  name={Space Vacuum},
  description={\textcolor{catAE}{\textbf{[Aerospace]}}},
}

\newglossaryentry{Space-Weather}{ sort={Space Weather},
  name={Space Weather},
  description={\textcolor{catAE}{\textbf{[Aerospace]}}},
}


\newglossaryentry{Van-Allen-Belts}{ sort={Van Allen Belts},
  name={Van Allen Belts},
  description={\textcolor{catAE}{\textbf{[Aerospace]}}},
}

\newglossaryentry{Cosmic-Ray}{ sort={Cosmic Ray},
  name={Cosmic Ray},
  description={\textcolor{catAE}{\textbf{[Aerospace]}}},
}

\newglossaryentry{Space-Plasma}{ sort={Space Plasma},
  name={Space Plasma},
  description={\textcolor{catAE}{\textbf{[Aerospace]}}},
}

\newglossaryentry{Space-Debris}{ sort={Space Debris},
  name={Space Debris},
  description={\textcolor{catAE}{\textbf{[Aerospace]}}},
}


\newglossaryentry{Micrometeoroid}{ sort={Micrometeoroid},
  name={Micrometeoroid},
  description={\textcolor{catAE}{\textbf{[Aerospace]}}},
}

\newglossaryentry{Additive-Manufacturing}{ sort={Additive Manufacturing},
  name={Additive Manufacturing},
  description={\textcolor{catAE}{\textbf{[Aerospace]}}},
}

\newglossaryentry{CNC-Machining}{ sort={CNC Machining},
  name={CNC Machining},
  description={\textcolor{catAE}{\textbf{[Aerospace]}}},
}

\newglossaryentry{Friction-Stir-Welding}{ sort={Friction Stir Welding},
  name={Friction Stir Welding},
  description={\textcolor{catAE}{\textbf{[Aerospace]}}},
}

\newglossaryentry{Electron-Beam-Welding}{ sort={Electron Beam Welding},
  name={Electron Beam Welding},
  description={\textcolor{catAE}{\textbf{[Aerospace]}}},
}

\newglossaryentry{Autoclave-Curing}{ sort={Autoclave Curing},
  name={Autoclave Curing},
  description={\textcolor{catAE}{\textbf{[Aerospace]}}},
}

\newglossaryentry{Resin-Transfer-Molding}{ sort={Resin Transfer Molding},
  name={Resin Transfer Molding},
  description={\textcolor{catAE}{\textbf{[Aerospace]}}},
}

\newglossaryentry{Filament-Winding}{ sort={Filament Winding},
  name={Filament Winding},
  description={\textcolor{catAE}{\textbf{[Aerospace]}}},
}

\newglossaryentry{Fatigue-Test}{ sort={Fatigue Test},
  name={Fatigue Test},
  description={\textcolor{catAE}{\textbf{[Aerospace]}}},
}

\newglossaryentry{Flutter-Test}{ sort={Flutter Test},
  name={Flutter Test},
  description={\textcolor{catAE}{\textbf{[Aerospace]}}},
}



\newglossaryentry{Structural-Test}{ sort={Structural Test},
  name={Structural Test},
  description={\textcolor{catAE}{\textbf{[Aerospace]}}},
}

\newglossaryentry{Flight-Test}{ sort={Flight Test},
  name={Flight Test},
  description={\textcolor{catAE}{\textbf{[Aerospace]}}},
}

\newglossaryentry{Qualification-Test}{ sort={Qualification Test},
  name={Qualification Test},
  description={\textcolor{catAE}{\textbf{[Aerospace]}}},
}

\newglossaryentry{Acceptance-Test}{ sort={Acceptance Test},
  name={Acceptance Test},
  description={\textcolor{catAE}{\textbf{[Aerospace]}}},
}

\newglossaryentry{Redundancy}{ sort={Redundancy},
  name={Redundancy},
  description={\textcolor{catAE}{\textbf{[Aerospace]}}},
}

\newglossaryentry{Fail-Safe-Design}{ sort={Fail-Safe Design},
  name={Fail-Safe Design},
  description={\textcolor{catAE}{\textbf{[Aerospace]}}},
}

\newglossaryentry{Fault-Tolerance}{ sort={Fault Tolerance},
  name={Fault Tolerance},
  description={\textcolor{catAE}{\textbf{[Aerospace]}}},
}

\newglossaryentry{MTBF}{ sort={MTBF},
  name={MTBF},
  description={\textcolor{catAE}{\textbf{[Aerospace]}}},
}

\newglossaryentry{Reliability}{ sort={Reliability},
  name={Reliability},
  description={\textcolor{catAE}{\textbf{[Aerospace]}}},
}

\newglossaryentry{Risk-Assessment}{ sort={Risk Assessment},
  name={Risk Assessment},
  description={\textcolor{catAE}{\textbf{[Aerospace]}}},
}

\newglossaryentry{Hazard-Analysis}{ sort={Hazard Analysis},
  name={Hazard Analysis},
  description={\textcolor{catAE}{\textbf{[Aerospace]}}},
}

\newglossaryentry{Flight-Safety}{ sort={Flight Safety},
  name={Flight Safety},
  description={\textcolor{catAE}{\textbf{[Aerospace]}}},
}


\newglossaryentry{Airworthiness}{ sort={Airworthiness},
  name={Airworthiness},
  description={\textcolor{catAE}{\textbf{[Aerospace]}}},
}

\newglossaryentry{Type-Certificate}{ sort={Type Certificate},
  name={Type Certificate},
  description={\textcolor{catAE}{\textbf{[Aerospace]}}},
}

\newglossaryentry{Supplemental-Type-Certificate}{ sort={Supplemental Type Certificate},
  name={Supplemental Type Certificate},
  description={\textcolor{catAE}{\textbf{[Aerospace]}}},
}




\newglossaryentry{DO-178C}{ sort={DO-178C},
  name={DO-178C},
  description={\textcolor{catAE}{\textbf{[Aerospace]}}},
}

\newglossaryentry{DO-254}{ sort={DO-254},
  name={DO-254},
  description={\textcolor{catAE}{\textbf{[Aerospace]}}},
}

\newglossaryentry{Urban-Air-Mobility}{ sort={Urban Air Mobility},
  name={Urban Air Mobility},
  description={\textcolor{catAE}{\textbf{[Aerospace]}}},
}

\newglossaryentry{Electric-Aircraft}{ sort={Electric Aircraft},
  name={Electric Aircraft},
  description={\textcolor{catAE}{\textbf{[Aerospace]}}},
}

\newglossaryentry{Hybrid-Electric-Propulsion}{ sort={Hybrid-Electric Propulsion},
  name={Hybrid-Electric Propulsion},
  description={\textcolor{catAE}{\textbf{[Aerospace]}}},
}

\newglossaryentry{Hydrogen-Aircraft}{ sort={Hydrogen Aircraft},
  name={Hydrogen Aircraft},
  description={\textcolor{catAE}{\textbf{[Aerospace]}}},
}

\newglossaryentry{Sustainable-Aviation-Fuel}{ sort={Sustainable Aviation Fuel},
  name={Sustainable Aviation Fuel},
  description={\textcolor{catAE}{\textbf{[Aerospace]}}},
}

\newglossaryentry{Autonomous-Flight}{ sort={Autonomous Flight},
  name={Autonomous Flight},
  description={\textcolor{catAE}{\textbf{[Aerospace]}}},
}


\newglossaryentry{Morphing-Wing}{ sort={Morphing Wing},
  name={Morphing Wing},
  description={\textcolor{catAE}{\textbf{[Aerospace]}}},
}


\newglossaryentry{Cockpit}{ sort={Cockpit},
  name={Cockpit},
  description={\textcolor{catAE}{\textbf{[Aerospace]}}},
}

\newglossaryentry{Canopy}{ sort={Canopy},
  name={Canopy},
  description={\textcolor{catAE}{\textbf{[Aerospace]}}},
}

\newglossaryentry{Nose-Gear}{ sort={Nose Gear},
  name={Nose Gear},
  description={\textcolor{catAE}{\textbf{[Aerospace]}}},
}

\newglossaryentry{Main-Gear}{ sort={Main Gear},
  name={Main Gear},
  description={\textcolor{catAE}{\textbf{[Aerospace]}}},
}

\newglossaryentry{Tailwheel}{ sort={Tailwheel},
  name={Tailwheel},
  description={\textcolor{catAE}{\textbf{[Aerospace]}}},
}

\newglossaryentry{Tricycle-Gear}{ sort={Tricycle Gear},
  name={Tricycle Gear},
  description={\textcolor{catAE}{\textbf{[Aerospace]}}},
}

\newglossaryentry{Brake}{ sort={Brake},
  name={Brake},
  description={\textcolor{catAE}{\textbf{[Aerospace]}}},
}

\newglossaryentry{Parking-Brake}{ sort={Parking Brake},
  name={Parking Brake},
  description={\textcolor{catAE}{\textbf{[Aerospace]}}},
}

\newglossaryentry{Control-Column}{ sort={Control Column},
  name={Control Column},
  description={\textcolor{catAE}{\textbf{[Aerospace]}}},
}

\newglossaryentry{Yoke}{ sort={Yoke},
  name={Yoke},
  description={\textcolor{catAE}{\textbf{[Aerospace]}}},
}

\newglossaryentry{Cyclic}{ sort={Cyclic},
  name={Cyclic},
  description={\textcolor{catAE}{\textbf{[Aerospace]}}},
}

\newglossaryentry{Collective}{ sort={Collective},
  name={Collective},
  description={\textcolor{catAE}{\textbf{[Aerospace]}}},
}

\newglossaryentry{Throttle}{ sort={Throttle},
  name={Throttle},
  description={\textcolor{catAE}{\textbf{[Aerospace]}}},
}

\newglossaryentry{Mixture-Control}{ sort={Mixture Control},
  name={Mixture Control},
  description={\textcolor{catAE}{\textbf{[Aerospace]}}},
}

\newglossaryentry{Propeller-Control}{ sort={Propeller Control},
  name={Propeller Control},
  description={\textcolor{catAE}{\textbf{[Aerospace]}}},
}

\newglossaryentry{Trim-Tab}{ sort={Trim Tab},
  name={Trim Tab},
  description={\textcolor{catAE}{\textbf{[Aerospace]}}},
}

\newglossaryentry{Balance-Tab}{ sort={Balance Tab},
  name={Balance Tab},
  description={\textcolor{catAE}{\textbf{[Aerospace]}}},
}

\newglossaryentry{Anti-Servo-Tab}{ sort={Anti-Servo Tab},
  name={Anti-Servo Tab},
  description={\textcolor{catAE}{\textbf{[Aerospace]}}},
}

\newglossaryentry{Mass-Balance}{ sort={Mass Balance},
  name={Mass Balance},
  description={\textcolor{catAE}{\textbf{[Aerospace]}}},
}

\newglossaryentry{Aerodynamic-Balance}{ sort={Aerodynamic Balance},
  name={Aerodynamic Balance},
  description={\textcolor{catAE}{\textbf{[Aerospace]}}},
}

\newglossaryentry{Horn-Balance}{ sort={Horn Balance},
  name={Horn Balance},
  description={\textcolor{catAE}{\textbf{[Aerospace]}}},
}

\newglossaryentry{Internal-Balance}{ sort={Internal Balance},
  name={Internal Balance},
  description={\textcolor{catAE}{\textbf{[Aerospace]}}},
}

\newglossaryentry{Boosted-Controls}{ sort={Boosted Controls},
  name={Boosted Controls},
  description={\textcolor{catAE}{\textbf{[Aerospace]}}},
}

\newglossaryentry{Power-Controls}{ sort={Power Controls},
  name={Power Controls},
  description={\textcolor{catAE}{\textbf{[Aerospace]}}},
}

\newglossaryentry{Control-Horn}{ sort={Control Horn},
  name={Control Horn},
  description={\textcolor{catAE}{\textbf{[Aerospace]}}},
}

\newglossaryentry{Pushrod}{ sort={Pushrod},
  name={Pushrod},
  description={\textcolor{catAE}{\textbf{[Aerospace]}}},
}

\newglossaryentry{Cable}{ sort={Cable},
  name={Cable},
  description={\textcolor{catAE}{\textbf{[Aerospace]}}},
}

\newglossaryentry{Bellerank}{ sort={Bellerank},
  name={Bellerank},
  description={\textcolor{catAE}{\textbf{[Aerospace]}}},
}

\newglossaryentry{Quadrant}{ sort={Quadrant},
  name={Quadrant},
  description={\textcolor{catAE}{\textbf{[Aerospace]}}},
}

\newglossaryentry{Linkage}{ sort={Linkage},
  name={Linkage},
  description={\textcolor{catAE}{\textbf{[Aerospace]}}},
}

\newglossaryentry{Control-Surface}{ sort={Control Surface},
  name={Control Surface},
  description={\textcolor{catAE}{\textbf{[Aerospace]}}},
}

\newglossaryentry{Primary-Control}{ sort={Primary Control},
  name={Primary Control},
  description={\textcolor{catAE}{\textbf{[Aerospace]}}},
}

\newglossaryentry{Secondary-Control}{ sort={Secondary Control},
  name={Secondary Control},
  description={\textcolor{catAE}{\textbf{[Aerospace]}}},
}

\newglossaryentry{High-Lift-Device}{ sort={High-Lift Device},
  name={High-Lift Device},
  description={\textcolor{catAE}{\textbf{[Aerospace]}}},
}

\newglossaryentry{Leading-Edge-Device}{ sort={Leading Edge Device},
  name={Leading Edge Device},
  description={\textcolor{catAE}{\textbf{[Aerospace]}}},
}

\newglossaryentry{Trailing-Edge-Device}{ sort={Trailing Edge Device},
  name={Trailing Edge Device},
  description={\textcolor{catAE}{\textbf{[Aerospace]}}},
}

\newglossaryentry{Fowler-Flap}{ sort={Fowler Flap},
  name={Fowler Flap},
  description={\textcolor{catAE}{\textbf{[Aerospace]}}},
}

\newglossaryentry{Split-Flap}{ sort={Split Flap},
  name={Split Flap},
  description={\textcolor{catAE}{\textbf{[Aerospace]}}},
}

\newglossaryentry{Slotted-Flap}{ sort={Slotted Flap},
  name={Slotted Flap},
  description={\textcolor{catAE}{\textbf{[Aerospace]}}},
}

\newglossaryentry{Krueger-Flap}{ sort={Krueger Flap},
  name={Krueger Flap},
  description={\textcolor{catAE}{\textbf{[Aerospace]}}},
}

\newglossaryentry{Handley-Page-Slat}{ sort={Handley Page Slat},
  name={Handley Page Slat},
  description={\textcolor{catAE}{\textbf{[Aerospace]}}},
}

\newglossaryentry{Blown-Flap}{ sort={Blown Flap},
  name={Blown Flap},
  description={\textcolor{catAE}{\textbf{[Aerospace]}}},
}

\newglossaryentry{Jet-Flap}{ sort={Jet Flap},
  name={Jet Flap},
  description={\textcolor{catAE}{\textbf{[Aerospace]}}},
}

\newglossaryentry{Circulation-Control-Wing}{ sort={Circulation Control Wing},
  name={Circulation Control Wing},
  description={\textcolor{catAE}{\textbf{[Aerospace]}}},
}

\newglossaryentry{Active-Flow-Control}{ sort={Active Flow Control},
  name={Active Flow Control},
  description={\textcolor{catAE}{\textbf{[Aerospace]}}},
}

\newglossaryentry{Load-Alleviation}{ sort={Load Alleviation},
  name={Load Alleviation},
  description={\textcolor{catAE}{\textbf{[Aerospace]}}},
}

\newglossaryentry{Gust-Alleviation}{ sort={Gust Alleviation},
  name={Gust Alleviation},
  description={\textcolor{catAE}{\textbf{[Aerospace]}}},
}

\newglossaryentry{Mode-Control-Panel}{ sort={Mode Control Panel},
  name={Mode Control Panel},
  description={\textcolor{catAE}{\textbf{[Aerospace]}}},
}

\newglossaryentry{Navigation-Display}{ sort={Navigation Display},
  name={Navigation Display},
  description={\textcolor{catAE}{\textbf{[Aerospace]}}},
}

\newglossaryentry{Primary-Flight-Display}{ sort={Primary Flight Display},
  name={Primary Flight Display},
  description={\textcolor{catAE}{\textbf{[Aerospace]}}},
}

\newglossaryentry{Engine-Display}{ sort={Engine Display},
  name={Engine Display},
  description={\textcolor{catAE}{\textbf{[Aerospace]}}},
}

\newglossaryentry{System-Display}{ sort={System Display},
  name={System Display},
  description={\textcolor{catAE}{\textbf{[Aerospace]}}},
}

\newglossaryentry{Caution-Warning-System}{ sort={Caution Warning System},
  name={Caution Warning System},
  description={\textcolor{catAE}{\textbf{[Aerospace]}}},
}

\newglossaryentry{Master-Warning}{ sort={Master Warning},
  name={Master Warning},
  description={\textcolor{catAE}{\textbf{[Aerospace]}}},
}

\newglossaryentry{EICAS}{ sort={EICAS},
  name={EICAS},
  description={\textcolor{catAE}{\textbf{[Aerospace]}}},
}

\newglossaryentry{ECAM}{ sort={ECAM},
  name={ECAM},
  description={\textcolor{catAE}{\textbf{[Aerospace]}}},
}

\newglossaryentry{Data-Bus}{ sort={Data Bus},
  name={Data Bus},
  description={\textcolor{catAE}{\textbf{[Aerospace]}}},
}

\newglossaryentry{ARINC-429}{ sort={ARINC 429},
  name={ARINC 429},
  description={\textcolor{catAE}{\textbf{[Aerospace]}}},
}

\newglossaryentry{MIL-STD-1553}{ sort={MIL-STD-1553},
  name={MIL-STD-1553},
  description={\textcolor{catAE}{\textbf{[Aerospace]}}},
}

\newglossaryentry{Aircraft-Wiring}{ sort={Aircraft Wiring},
  name={Aircraft Wiring},
  description={\textcolor{catAE}{\textbf{[Aerospace]}}},
}

\newglossaryentry{Circuit-Breaker}{ sort={Circuit Breaker},
  name={Circuit Breaker},
  description={\textcolor{catAE}{\textbf{[Aerospace]}}},
}

\newglossaryentry{Bus-Bar}{ sort={Bus Bar},
  name={Bus Bar},
  description={\textcolor{catAE}{\textbf{[Aerospace]}}},
}

\newglossaryentry{Generator}{ sort={Generator},
  name={Generator},
  description={\textcolor{catAE}{\textbf{[Aerospace]}}},
}

\newglossaryentry{Alternator}{ sort={Alternator},
  name={Alternator},
  description={\textcolor{catAE}{\textbf{[Aerospace]}}},
}


\newglossaryentry{APU}{ sort={APU},
  name={APU},
  description={\textcolor{catAE}{\textbf{[Aerospace]}}},
}

\newglossaryentry{Inverter}{ sort={Inverter},
  name={Inverter},
  description={\textcolor{catAE}{\textbf{[Aerospace]}}},
}

\newglossaryentry{Transformer}{ sort={Transformer},
  name={Transformer},
  description={\textcolor{catAE}{\textbf{[Aerospace]}}},
}

\newglossaryentry{Rectifier}{ sort={Rectifier},
  name={Rectifier},
  description={\textcolor{catAE}{\textbf{[Aerospace]}}},
}

\newglossaryentry{Voltage-Regulator}{ sort={Voltage Regulator},
  name={Voltage Regulator},
  description={\textcolor{catAE}{\textbf{[Aerospace]}}},
}

\newglossaryentry{Oxygen-System}{ sort={Oxygen System},
  name={Oxygen System},
  description={\textcolor{catAE}{\textbf{[Aerospace]}}},
}

\newglossaryentry{Pressurization-System}{ sort={Pressurization System},
  name={Pressurization System},
  description={\textcolor{catAE}{\textbf{[Aerospace]}}},
}

\newglossaryentry{Air-Conditioning}{ sort={Air Conditioning},
  name={Air Conditioning},
  description={\textcolor{catAE}{\textbf{[Aerospace]}}},
}

\newglossaryentry{Bleed-Air}{ sort={Bleed Air},
  name={Bleed Air},
  description={\textcolor{catAE}{\textbf{[Aerospace]}}},
}

\newglossaryentry{ECS-Pack}{ sort={ECS Pack},
  name={ECS Pack},
  description={\textcolor{catAE}{\textbf{[Aerospace]}}},
}

\newglossaryentry{Ventilation}{ sort={Ventilation},
  name={Ventilation},
  description={\textcolor{catAE}{\textbf{[Aerospace]}}},
}

\newglossaryentry{Rotor-Brake}{ sort={Rotor Brake},
  name={Rotor Brake},
  description={\textcolor{catAE}{\textbf{[Aerospace]}}},
}

\newglossaryentry{Rotor-Head}{ sort={Rotor Head},
  name={Rotor Head},
  description={\textcolor{catAE}{\textbf{[Aerospace]}}},
}

\newglossaryentry{Swashplate}{ sort={Swashplate},
  name={Swashplate},
  description={\textcolor{catAE}{\textbf{[Aerospace]}}},
}

\newglossaryentry{Collective-Pitch}{ sort={Collective Pitch},
  name={Collective Pitch},
  description={\textcolor{catAE}{\textbf{[Aerospace]}}},
}

\newglossaryentry{Cyclic-Pitch}{ sort={Cyclic Pitch},
  name={Cyclic Pitch},
  description={\textcolor{catAE}{\textbf{[Aerospace]}}},
}

\newglossaryentry{Lead-Lag-Hinge}{ sort={Lead-Lag Hinge},
  name={Lead-Lag Hinge},
  description={\textcolor{catAE}{\textbf{[Aerospace]}}},
}

\newglossaryentry{Flapping-Hinge}{ sort={Flapping Hinge},
  name={Flapping Hinge},
  description={\textcolor{catAE}{\textbf{[Aerospace]}}},
}

\newglossaryentry{Feathering-Hinge}{ sort={Feathering Hinge},
  name={Feathering Hinge},
  description={\textcolor{catAE}{\textbf{[Aerospace]}}},
}

\newglossaryentry{Fully-Articulated-Rotor}{ sort={Fully Articulated Rotor},
  name={Fully Articulated Rotor},
  description={\textcolor{catAE}{\textbf{[Aerospace]}}},
}

\newglossaryentry{Semi-Rigid-Rotor}{ sort={Semi-Rigid Rotor},
  name={Semi-Rigid Rotor},
  description={\textcolor{catAE}{\textbf{[Aerospace]}}},
}

\newglossaryentry{Rigid-Rotor}{ sort={Rigid Rotor},
  name={Rigid Rotor},
  description={\textcolor{catAE}{\textbf{[Aerospace]}}},
}

\newglossaryentry{Tail-Rotor}{ sort={Tail Rotor},
  name={Tail Rotor},
  description={\textcolor{catAE}{\textbf{[Aerospace]}}},
}

\newglossaryentry{Fenestron}{ sort={Fenestron},
  name={Fenestron},
  description={\textcolor{catAE}{\textbf{[Aerospace]}}},
}

\newglossaryentry{NOTAR}{ sort={NOTAR},
  name={NOTAR},
  description={\textcolor{catAE}{\textbf{[Aerospace]}}},
}

\newglossaryentry{Main-Rotor}{ sort={Main Rotor},
  name={Main Rotor},
  description={\textcolor{catAE}{\textbf{[Aerospace]}}},
}

\newglossaryentry{Tip-Speed}{ sort={Tip Speed},
  name={Tip Speed},
  description={\textcolor{catAE}{\textbf{[Aerospace]}}},
}

\newglossaryentry{Hover}{ sort={Hover},
  name={Hover},
  description={\textcolor{catAE}{\textbf{[Aerospace]}}},
}

\newglossaryentry{Autorotation-Landing}{ sort={Autorotation Landing},
  name={Autorotation Landing},
  description={\textcolor{catAE}{\textbf{[Aerospace]}}},
}

\newglossaryentry{Ground-Resonance}{ sort={Ground Resonance},
  name={Ground Resonance},
  description={\textcolor{catAE}{\textbf{[Aerospace]}}},
}

\newglossaryentry{Retreating-Blade-Stall}{ sort={Retreating Blade Stall},
  name={Retreating Blade Stall},
  description={\textcolor{catAE}{\textbf{[Aerospace]}}},
}

\newglossaryentry{Dissymmetry-of-Lift}{ sort={Dissymmetry of Lift},
  name={Dissymmetry of Lift},
  description={\textcolor{catAE}{\textbf{[Aerospace]}}},
}

\newglossaryentry{Translational-Lift}{ sort={Translational Lift},
  name={Translational Lift},
  description={\textcolor{catAE}{\textbf{[Aerospace]}}},
}

\newglossaryentry{Effective-Translational-Lift}{ sort={Effective Translational Lift},
  name={Effective Translational Lift},
  description={\textcolor{catAE}{\textbf{[Aerospace]}}},
}

\newglossaryentry{Ground-Effect-Helicopter}{ sort={Ground Effect Helicopter},
  name={Ground Effect Helicopter},
  description={\textcolor{catAE}{\textbf{[Aerospace]}}},
}

\newglossaryentry{Vortex-Ring-State}{ sort={Vortex Ring State},
  name={Vortex Ring State},
  description={\textcolor{catAE}{\textbf{[Aerospace]}}},
}

\newglossaryentry{Settling-With-Power}{ sort={Settling With Power},
  name={Settling With Power},
  description={\textcolor{catAE}{\textbf{[Aerospace]}}},
}

\newglossaryentry{Power-Required-Curve}{ sort={Power Required Curve},
  name={Power Required Curve},
  description={\textcolor{catAE}{\textbf{[Aerospace]}}},
}

\newglossaryentry{Power-Available-Curve}{ sort={Power Available Curve},
  name={Power Available Curve},
  description={\textcolor{catAE}{\textbf{[Aerospace]}}},
}

\newglossaryentry{Height-Velocity-Diagram}{ sort={Height-Velocity Diagram},
  name={Height-Velocity Diagram},
  description={\textcolor{catAE}{\textbf{[Aerospace]}}},
}

\newglossaryentry{H-V-Curve}{ sort={H-V Curve},
  name={H-V Curve},
  description={\textcolor{catAE}{\textbf{[Aerospace]}}},
}

\newglossaryentry{Cat-A-Takeoff}{ sort={Cat A Takeoff},
  name={Cat A Takeoff},
  description={\textcolor{catAE}{\textbf{[Aerospace]}}},
}

\newglossaryentry{Cat-B-Takeoff}{ sort={Cat B Takeoff},
  name={Cat B Takeoff},
  description={\textcolor{catAE}{\textbf{[Aerospace]}}},
}

\newglossaryentry{Stall-Warning-System}{ sort={Stall Warning System},
  name={Stall Warning System},
  description={\textcolor{catAE}{\textbf{[Aerospace]}}},
}

\newglossaryentry{Stick-Shaker}{ sort={Stick Shaker},
  name={Stick Shaker},
  description={\textcolor{catAE}{\textbf{[Aerospace]}}},
}

\newglossaryentry{Yaw-Damper}{ sort={Yaw Damper},
  name={Yaw Damper},
  description={\textcolor{catAE}{\textbf{[Aerospace]}}},
}

\newglossaryentry{Autothrottle}{ sort={Autothrottle},
  name={Autothrottle},
  description={\textcolor{catAE}{\textbf{[Aerospace]}}},
}

\newglossaryentry{Flight-Director}{ sort={Flight Director},
  name={Flight Director},
  description={\textcolor{catAE}{\textbf{[Aerospace]}}},
}

\newglossaryentry{TCAS}{ sort={TCAS},
  name={TCAS},
  description={\textcolor{catAE}{\textbf{[Aerospace]}}},
}

\newglossaryentry{GPWS}{ sort={GPWS},
  name={GPWS},
  description={\textcolor{catAE}{\textbf{[Aerospace]}}},
}

\newglossaryentry{Radar-Altimeter}{ sort={Radar Altimeter},
  name={Radar Altimeter},
  description={\textcolor{catAE}{\textbf{[Aerospace]}}},
}


\newglossaryentry{Rocket-Engine}{ sort={Rocket Engine},
  name={Rocket Engine},
  description={\textcolor{catRM}{\textbf{[Rockets and Missiles]}}},
}

\newglossaryentry{Solid-Rocket-Motor}{ sort={Solid Rocket Motor},
  name={Solid Rocket Motor},
  description={\textcolor{catRM}{\textbf{[Rockets and Missiles]}}},
}

\newglossaryentry{Liquid-Rocket-Engine}{ sort={Liquid Rocket Engine},
  name={Liquid Rocket Engine},
  description={\textcolor{catRM}{\textbf{[Rockets and Missiles]}}},
}

\newglossaryentry{Hybrid-Rocket}{ sort={Hybrid Rocket},
  name={Hybrid Rocket},
  description={\textcolor{catRM}{\textbf{[Rockets and Missiles]}}},
}

\newglossaryentry{Specific-Impulse}{ sort={Specific Impulse},
  name={Specific Impulse},
  description={\textcolor{catRM}{\textbf{[Rockets and Missiles]}}},
}

\newglossaryentry{Total-Impulse}{ sort={Total Impulse},
  name={Total Impulse},
  description={\textcolor{catRM}{\textbf{[Rockets and Missiles]}}},
}

\newglossaryentry{Thrust}{ sort={Thrust},
  name={Thrust},
  description={\textcolor{catRM}{\textbf{[Rockets and Missiles]}}},
}

\newglossaryentry{Thrust-to-Weight-Ratio}{ sort={Thrust-to-Weight Ratio},
  name={Thrust-to-Weight Ratio},
  description={\textcolor{catRM}{\textbf{[Rockets and Missiles]}}},
}

\newglossaryentry{Propellant}{ sort={Propellant},
  name={Propellant},
  description={\textcolor{catRM}{\textbf{[Rockets and Missiles]}}},
}

\newglossaryentry{Oxidizer}{ sort={Oxidizer},
  name={Oxidizer},
  description={\textcolor{catRM}{\textbf{[Rockets and Missiles]}}},
}

\newglossaryentry{Bipropellant}{ sort={Bipropellant},
  name={Bipropellant},
  description={\textcolor{catRM}{\textbf{[Rockets and Missiles]}}},
}

\newglossaryentry{Monopropellant}{ sort={Monopropellant},
  name={Monopropellant},
  description={\textcolor{catRM}{\textbf{[Rockets and Missiles]}}},
}

\newglossaryentry{Hypergolic}{ sort={Hypergolic},
  name={Hypergolic},
  description={\textcolor{catRM}{\textbf{[Rockets and Missiles]}}},
}

\newglossaryentry{Cryogenic-Propellant}{ sort={Cryogenic Propellant},
  name={Cryogenic Propellant},
  description={\textcolor{catRM}{\textbf{[Rockets and Missiles]}}},
}

\newglossaryentry{RP-1}{ sort={RP-1},
  name={RP-1},
  description={\textcolor{catRM}{\textbf{[Rockets and Missiles]}}},
}

\newglossaryentry{Liquid-Hydrogen}{ sort={Liquid Hydrogen},
  name={Liquid Hydrogen},
  description={\textcolor{catRM}{\textbf{[Rockets and Missiles]}}},
}

\newglossaryentry{Liquid-Oxygen}{ sort={Liquid Oxygen},
  name={Liquid Oxygen},
  description={\textcolor{catRM}{\textbf{[Rockets and Missiles]}}},
}

\newglossaryentry{Nitrogen-Tetroxide}{ sort={Nitrogen Tetroxide},
  name={Nitrogen Tetroxide},
  description={\textcolor{catRM}{\textbf{[Rockets and Missiles]}}},
}

\newglossaryentry{Hydrazine}{ sort={Hydrazine},
  name={Hydrazine},
  description={\textcolor{catRM}{\textbf{[Rockets and Missiles]}}},
}

\newglossaryentry{MMH}{ sort={MMH},
  name={MMH},
  description={\textcolor{catRM}{\textbf{[Rockets and Missiles]}}},
}

\newglossaryentry{UDMH}{ sort={UDMH},
  name={UDMH},
  description={\textcolor{catRM}{\textbf{[Rockets and Missiles]}}},
}

\newglossaryentry{Propellant-Tank}{ sort={Propellant Tank},
  name={Propellant Tank},
  description={\textcolor{catRM}{\textbf{[Rockets and Missiles]}}},
}

\newglossaryentry{Tank-Pressurization}{ sort={Tank Pressurization},
  name={Tank Pressurization},
  description={\textcolor{catRM}{\textbf{[Rockets and Missiles]}}},
}

\newglossaryentry{Combustion-Chamber}{ sort={Combustion Chamber},
  name={Combustion Chamber},
  description={\textcolor{catRM}{\textbf{[Rockets and Missiles]}}},
}

\newglossaryentry{Injector}{ sort={Injector},
  name={Injector},
  description={\textcolor{catRM}{\textbf{[Rockets and Missiles]}}},
}

\newglossaryentry{Nozzle}{ sort={Nozzle},
  name={Nozzle},
  description={\textcolor{catRM}{\textbf{[Rockets and Missiles]}}},
}

\newglossaryentry{De-Laval-Nozzle}{ sort={De Laval Nozzle},
  name={De Laval Nozzle},
  description={\textcolor{catRM}{\textbf{[Rockets and Missiles]}}},
}

\newglossaryentry{Throat}{ sort={Throat},
  name={Throat},
  description={\textcolor{catRM}{\textbf{[Rockets and Missiles]}}},
}

\newglossaryentry{Expansion-Ratio}{ sort={Expansion Ratio},
  name={Expansion Ratio},
  description={\textcolor{catRM}{\textbf{[Rockets and Missiles]}}},
}

\newglossaryentry{Nozzle-Contour}{ sort={Nozzle Contour},
  name={Nozzle Contour},
  description={\textcolor{catRM}{\textbf{[Rockets and Missiles]}}},
}

\newglossaryentry{Rao-Nozzle}{ sort={Rao Nozzle},
  name={Rao Nozzle},
  description={\textcolor{catRM}{\textbf{[Rockets and Missiles]}}},
}

\newglossaryentry{Bell-Nozzle}{ sort={Bell Nozzle},
  name={Bell Nozzle},
  description={\textcolor{catRM}{\textbf{[Rockets and Missiles]}}},
}

\newglossaryentry{Plug-Nozzle}{ sort={Plug Nozzle},
  name={Plug Nozzle},
  description={\textcolor{catRM}{\textbf{[Rockets and Missiles]}}},
}

\newglossaryentry{Aerospike-Nozzle}{ sort={Aerospike Nozzle},
  name={Aerospike Nozzle},
  description={\textcolor{catRM}{\textbf{[Rockets and Missiles]}}},
}

\newglossaryentry{Regenerative-Cooling}{ sort={Regenerative Cooling},
  name={Regenerative Cooling},
  description={\textcolor{catRM}{\textbf{[Rockets and Missiles]}}},
}

\newglossaryentry{Film-Cooling}{ sort={Film Cooling},
  name={Film Cooling},
  description={\textcolor{catRM}{\textbf{[Rockets and Missiles]}}},
}

\newglossaryentry{Ablative-Cooling}{ sort={Ablative Cooling},
  name={Ablative Cooling},
  description={\textcolor{catRM}{\textbf{[Rockets and Missiles]}}},
}

\newglossaryentry{Radiative-Cooling}{ sort={Radiative Cooling},
  name={Radiative Cooling},
  description={\textcolor{catRM}{\textbf{[Rockets and Missiles]}}},
}

\newglossaryentry{Igniter}{ sort={Igniter},
  name={Igniter},
  description={\textcolor{catRM}{\textbf{[Rockets and Missiles]}}},
}

\newglossaryentry{Pyrotechnic-Igniter}{ sort={Pyrotechnic Igniter},
  name={Pyrotechnic Igniter},
  description={\textcolor{catRM}{\textbf{[Rockets and Missiles]}}},
}

\newglossaryentry{Spark-Igniter}{ sort={Spark Igniter},
  name={Spark Igniter},
  description={\textcolor{catRM}{\textbf{[Rockets and Missiles]}}},
}

\newglossaryentry{Catalytic-Igniter}{ sort={Catalytic Igniter},
  name={Catalytic Igniter},
  description={\textcolor{catRM}{\textbf{[Rockets and Missiles]}}},
}

\newglossaryentry{Thrust-Chamber}{ sort={Thrust Chamber},
  name={Thrust Chamber},
  description={\textcolor{catRM}{\textbf{[Rockets and Missiles]}}},
}

\newglossaryentry{Engine-Cycle}{ sort={Engine Cycle},
  name={Engine Cycle},
  description={\textcolor{catRM}{\textbf{[Rockets and Missiles]}}},
}

\newglossaryentry{Pressure-Fed-Cycle}{ sort={Pressure Fed Cycle},
  name={Pressure Fed Cycle},
  description={\textcolor{catRM}{\textbf{[Rockets and Missiles]}}},
}

\newglossaryentry{Gas-Generator-Cycle}{ sort={Gas Generator Cycle},
  name={Gas Generator Cycle},
  description={\textcolor{catRM}{\textbf{[Rockets and Missiles]}}},
}

\newglossaryentry{Staged-Combustion-Cycle}{ sort={Staged Combustion Cycle},
  name={Staged Combustion Cycle},
  description={\textcolor{catRM}{\textbf{[Rockets and Missiles]}}},
}



\newglossaryentry{Expander-Cycle}{ sort={Expander Cycle},
  name={Expander Cycle},
  description={\textcolor{catRM}{\textbf{[Rockets and Missiles]}}},
}

\newglossaryentry{Tap-Off-Cycle}{ sort={Tap-Off Cycle},
  name={Tap-Off Cycle},
  description={\textcolor{catRM}{\textbf{[Rockets and Missiles]}}},
}

\newglossaryentry{Electric-Pump-Cycle}{ sort={Electric Pump Cycle},
  name={Electric Pump Cycle},
  description={\textcolor{catRM}{\textbf{[Rockets and Missiles]}}},
}

\newglossaryentry{Turbopump}{ sort={Turbopump},
  name={Turbopump},
  description={\textcolor{catRM}{\textbf{[Rockets and Missiles]}}},
}

\newglossaryentry{Preburner}{ sort={Preburner},
  name={Preburner},
  description={\textcolor{catRM}{\textbf{[Rockets and Missiles]}}},
}

\newglossaryentry{Main-Combustion-Chamber}{ sort={Main Combustion Chamber},
  name={Main Combustion Chamber},
  description={\textcolor{catRM}{\textbf{[Rockets and Missiles]}}},
}

\newglossaryentry{Thrust-Vector-Control}{ sort={Thrust Vector Control},
  name={Thrust Vector Control},
  description={\textcolor{catRM}{\textbf{[Rockets and Missiles]}}},
}


\newglossaryentry{Vernier-Engine}{ sort={Vernier Engine},
  name={Vernier Engine},
  description={\textcolor{catRM}{\textbf{[Rockets and Missiles]}}},
}


\newglossaryentry{Cold-Gas-Thruster}{ sort={Cold Gas Thruster},
  name={Cold Gas Thruster},
  description={\textcolor{catRM}{\textbf{[Rockets and Missiles]}}},
}

\newglossaryentry{Launch-Vehicle}{ sort={Launch Vehicle},
  name={Launch Vehicle},
  description={\textcolor{catRM}{\textbf{[Rockets and Missiles]}}},
}

\newglossaryentry{Expendable-Launch-Vehicle}{ sort={Expendable Launch Vehicle},
  name={Expendable Launch Vehicle},
  description={\textcolor{catRM}{\textbf{[Rockets and Missiles]}}},
}

\newglossaryentry{Reusable-Launch-Vehicle}{ sort={Reusable Launch Vehicle},
  name={Reusable Launch Vehicle},
  description={\textcolor{catRM}{\textbf{[Rockets and Missiles]}}},
}

\newglossaryentry{Multi-Stage-Rocket}{ sort={Multi-Stage Rocket},
  name={Multi-Stage Rocket},
  description={\textcolor{catRM}{\textbf{[Rockets and Missiles]}}},
}

\newglossaryentry{First-Stage}{ sort={First Stage},
  name={First Stage},
  description={\textcolor{catRM}{\textbf{[Rockets and Missiles]}}},
}

\newglossaryentry{Second-Stage}{ sort={Second Stage},
  name={Second Stage},
  description={\textcolor{catRM}{\textbf{[Rockets and Missiles]}}},
}

\newglossaryentry{Upper-Stage}{ sort={Upper Stage},
  name={Upper Stage},
  description={\textcolor{catRM}{\textbf{[Rockets and Missiles]}}},
}

\newglossaryentry{Strap-On-Booster}{ sort={Strap-On Booster},
  name={Strap-On Booster},
  description={\textcolor{catRM}{\textbf{[Rockets and Missiles]}}},
}

\newglossaryentry{Payload-Fairing}{ sort={Payload Fairing},
  name={Payload Fairing},
  description={\textcolor{catRM}{\textbf{[Rockets and Missiles]}}},
}

\newglossaryentry{Interstage}{ sort={Interstage},
  name={Interstage},
  description={\textcolor{catRM}{\textbf{[Rockets and Missiles]}}},
}

\newglossaryentry{Stage-Separation}{ sort={Stage Separation},
  name={Stage Separation},
  description={\textcolor{catRM}{\textbf{[Rockets and Missiles]}}},
}

\newglossaryentry{Fairing-Separation}{ sort={Fairing Separation},
  name={Fairing Separation},
  description={\textcolor{catRM}{\textbf{[Rockets and Missiles]}}},
}


\newglossaryentry{Launch-Tower}{ sort={Launch Tower},
  name={Launch Tower},
  description={\textcolor{catRM}{\textbf{[Rockets and Missiles]}}},
}

\newglossaryentry{Launch-Rail}{ sort={Launch Rail},
  name={Launch Rail},
  description={\textcolor{catRM}{\textbf{[Rockets and Missiles]}}},
}

\newglossaryentry{Launch-Tube}{ sort={Launch Tube},
  name={Launch Tube},
  description={\textcolor{catRM}{\textbf{[Rockets and Missiles]}}},
}

\newglossaryentry{Mobile-Launcher}{ sort={Mobile Launcher},
  name={Mobile Launcher},
  description={\textcolor{catRM}{\textbf{[Rockets and Missiles]}}},
}

\newglossaryentry{Silo}{ sort={Silo},
  name={Silo},
  description={\textcolor{catRM}{\textbf{[Rockets and Missiles]}}},
}


\newglossaryentry{T-0}{ sort={T-0},
  name={T-0},
  description={\textcolor{catRM}{\textbf{[Rockets and Missiles]}}},
}

\newglossaryentry{Lift-Off}{ sort={Lift-Off},
  name={Lift-Off},
  description={\textcolor{catRM}{\textbf{[Rockets and Missiles]}}},
}

\newglossaryentry{Max-Q}{ sort={Max Q},
  name={Max Q},
  description={\textcolor{catRM}{\textbf{[Rockets and Missiles]}}},
}

\newglossaryentry{MECO}{ sort={MECO},
  name={MECO},
  description={\textcolor{catRM}{\textbf{[Rockets and Missiles]}}},
}

\newglossaryentry{SECO}{ sort={SECO},
  name={SECO},
  description={\textcolor{catRM}{\textbf{[Rockets and Missiles]}}},
}

\newglossaryentry{Insertion}{ sort={Insertion},
  name={Insertion},
  description={\textcolor{catRM}{\textbf{[Rockets and Missiles]}}},
}

\newglossaryentry{Injection-Burn}{ sort={Injection Burn},
  name={Injection Burn},
  description={\textcolor{catRM}{\textbf{[Rockets and Missiles]}}},
}

\newglossaryentry{Circularization-Burn}{ sort={Circularization Burn},
  name={Circularization Burn},
  description={\textcolor{catRM}{\textbf{[Rockets and Missiles]}}},
}

\newglossaryentry{Deorbit-Burn}{ sort={Deorbit Burn},
  name={Deorbit Burn},
  description={\textcolor{catRM}{\textbf{[Rockets and Missiles]}}},
}

\newglossaryentry{Retro-Rocket}{ sort={Retro Rocket},
  name={Retro Rocket},
  description={\textcolor{catRM}{\textbf{[Rockets and Missiles]}}},
}

\newglossaryentry{Ullage-Motor}{ sort={Ullage Motor},
  name={Ullage Motor},
  description={\textcolor{catRM}{\textbf{[Rockets and Missiles]}}},
}

\newglossaryentry{Spin-Motor}{ sort={Spin Motor},
  name={Spin Motor},
  description={\textcolor{catRM}{\textbf{[Rockets and Missiles]}}},
}

\newglossaryentry{Separation-Motor}{ sort={Separation Motor},
  name={Separation Motor},
  description={\textcolor{catRM}{\textbf{[Rockets and Missiles]}}},
}

\newglossaryentry{IRFNA}{ sort={IRFNA},
  name={IRFNA},
  description={\textcolor{catRM}{\textbf{[Rockets and Missiles]}}},
}

\newglossaryentry{LOX}{ sort={LOX},
  name={LOX},
  description={\textcolor{catRM}{\textbf{[Rockets and Missiles]}}},
}

\newglossaryentry{LH2}{ sort={LH2},
  name={LH2},
  description={\textcolor{catRM}{\textbf{[Rockets and Missiles]}}},
}

\newglossaryentry{LNG}{ sort={LNG},
  name={LNG},
  description={\textcolor{catRM}{\textbf{[Rockets and Missiles]}}},
}

\newglossaryentry{Propellant-Mass-Fraction}{ sort={Propellant Mass Fraction},
  name={Propellant Mass Fraction},
  description={\textcolor{catRM}{\textbf{[Rockets and Missiles]}}},
}

\newglossaryentry{Dry-Mass}{ sort={Dry Mass},
  name={Dry Mass},
  description={\textcolor{catRM}{\textbf{[Rockets and Missiles]}}},
}

\newglossaryentry{Wet-Mass}{ sort={Wet Mass},
  name={Wet Mass},
  description={\textcolor{catRM}{\textbf{[Rockets and Missiles]}}},
}

\newglossaryentry{Payload-Mass-Fraction}{ sort={Payload Mass Fraction},
  name={Payload Mass Fraction},
  description={\textcolor{catRM}{\textbf{[Rockets and Missiles]}}},
}

\newglossaryentry{Tsiolkovsky-Equation}{ sort={Tsiolkovsky Equation},
  name={Tsiolkovsky Equation},
  description={\textcolor{catRM}{\textbf{[Rockets and Missiles]}}},
}

\newglossaryentry{Characteristic-Velocity}{ sort={Characteristic Velocity},
  name={Characteristic Velocity},
  description={\textcolor{catRM}{\textbf{[Rockets and Missiles]}}},
}

\newglossaryentry{Exhaust-Velocity}{ sort={Exhaust Velocity},
  name={Exhaust Velocity},
  description={\textcolor{catRM}{\textbf{[Rockets and Missiles]}}},
}

\newglossaryentry{Effective-Exhaust-Velocity}{ sort={Effective Exhaust Velocity},
  name={Effective Exhaust Velocity},
  description={\textcolor{catRM}{\textbf{[Rockets and Missiles]}}},
}

\newglossaryentry{Mass-Flow-Rate}{ sort={Mass Flow Rate},
  name={Mass Flow Rate},
  description={\textcolor{catRM}{\textbf{[Rockets and Missiles]}}},
}

\newglossaryentry{Chamber-Pressure}{ sort={Chamber Pressure},
  name={Chamber Pressure},
  description={\textcolor{catRM}{\textbf{[Rockets and Missiles]}}},
}

\newglossaryentry{Expansion-Ratio-Effect}{ sort={Expansion Ratio Effect},
  name={Expansion Ratio Effect},
  description={\textcolor{catRM}{\textbf{[Rockets and Missiles]}}},
}

\newglossaryentry{Optimal-Expansion}{ sort={Optimal Expansion},
  name={Optimal Expansion},
  description={\textcolor{catRM}{\textbf{[Rockets and Missiles]}}},
}

\newglossaryentry{Overexpansion}{ sort={Overexpansion},
  name={Overexpansion},
  description={\textcolor{catRM}{\textbf{[Rockets and Missiles]}}},
}

\newglossaryentry{Underexpansion}{ sort={Underexpansion},
  name={Underexpansion},
  description={\textcolor{catRM}{\textbf{[Rockets and Missiles]}}},
}

\newglossaryentry{Flow-Separation-in-Nozzle}{ sort={Flow Separation in Nozzle},
  name={Flow Separation in Nozzle},
  description={\textcolor{catRM}{\textbf{[Rockets and Missiles]}}},
}

\newglossaryentry{Side-Load}{ sort={Side Load},
  name={Side Load},
  description={\textcolor{catRM}{\textbf{[Rockets and Missiles]}}},
}

\newglossaryentry{Nozzle-Extension}{ sort={Nozzle Extension},
  name={Nozzle Extension},
  description={\textcolor{catRM}{\textbf{[Rockets and Missiles]}}},
}

\newglossaryentry{Dual-Bell-Nozzle}{ sort={Dual Bell Nozzle},
  name={Dual Bell Nozzle},
  description={\textcolor{catRM}{\textbf{[Rockets and Missiles]}}},
}

\newglossaryentry{Thrust-Chamber-Liner}{ sort={Thrust Chamber Liner},
  name={Thrust Chamber Liner},
  description={\textcolor{catRM}{\textbf{[Rockets and Missiles]}}},
}

\newglossaryentry{Combustion-Stability}{ sort={Combustion Stability},
  name={Combustion Stability},
  description={\textcolor{catRM}{\textbf{[Rockets and Missiles]}}},
}

\newglossaryentry{Combustion-Instability}{ sort={Combustion Instability},
  name={Combustion Instability},
  description={\textcolor{catRM}{\textbf{[Rockets and Missiles]}}},
}

\newglossaryentry{Chugging}{ sort={Chugging},
  name={Chugging},
  description={\textcolor{catRM}{\textbf{[Rockets and Missiles]}}},
}

\newglossaryentry{Buzzing}{ sort={Buzzing},
  name={Buzzing},
  description={\textcolor{catRM}{\textbf{[Rockets and Missiles]}}},
}

\newglossaryentry{Screaming}{ sort={Screaming},
  name={Screaming},
  description={\textcolor{catRM}{\textbf{[Rockets and Missiles]}}},
}

\newglossaryentry{Pogo-Oscillation}{ sort={Pogo Oscillation},
  name={Pogo Oscillation},
  description={\textcolor{catRM}{\textbf{[Rockets and Missiles]}}},
}

\newglossaryentry{Acoustic-Cavity}{ sort={Acoustic Cavity},
  name={Acoustic Cavity},
  description={\textcolor{catRM}{\textbf{[Rockets and Missiles]}}},
}

\newglossaryentry{Baffle}{ sort={Baffle},
  name={Baffle},
  description={\textcolor{catRM}{\textbf{[Rockets and Missiles]}}},
}

\newglossaryentry{Resonator}{ sort={Resonator},
  name={Resonator},
  description={\textcolor{catRM}{\textbf{[Rockets and Missiles]}}},
}

\newglossaryentry{Injector-Pattern}{ sort={Injector Pattern},
  name={Injector Pattern},
  description={\textcolor{catRM}{\textbf{[Rockets and Missiles]}}},
}

\newglossaryentry{Showerhead-Injector}{ sort={Showerhead Injector},
  name={Showerhead Injector},
  description={\textcolor{catRM}{\textbf{[Rockets and Missiles]}}},
}

\newglossaryentry{Impinging-Injector}{ sort={Impinging Injector},
  name={Impinging Injector},
  description={\textcolor{catRM}{\textbf{[Rockets and Missiles]}}},
}

\newglossaryentry{Coaxial-Injector}{ sort={Coaxial Injector},
  name={Coaxial Injector},
  description={\textcolor{catRM}{\textbf{[Rockets and Missiles]}}},
}

\newglossaryentry{Pintle-Injector}{ sort={Pintle Injector},
  name={Pintle Injector},
  description={\textcolor{catRM}{\textbf{[Rockets and Missiles]}}},
}

\newglossaryentry{Shear-Injector}{ sort={Shear Injector},
  name={Shear Injector},
  description={\textcolor{catRM}{\textbf{[Rockets and Missiles]}}},
}

\newglossaryentry{Atomization}{ sort={Atomization},
  name={Atomization},
  description={\textcolor{catRM}{\textbf{[Rockets and Missiles]}}},
}

\newglossaryentry{Droplet-Combustion}{ sort={Droplet Combustion},
  name={Droplet Combustion},
  description={\textcolor{catRM}{\textbf{[Rockets and Missiles]}}},
}

\newglossaryentry{Combustion-Efficiency}{ sort={Combustion Efficiency},
  name={Combustion Efficiency},
  description={\textcolor{catRM}{\textbf{[Rockets and Missiles]}}},
}

\newglossaryentry{C-Star-Efficiency}{ sort={C Star Efficiency},
  name={C Star Efficiency},
  description={\textcolor{catRM}{\textbf{[Rockets and Missiles]}}},
}

\newglossaryentry{Thrust-Coefficient}{ sort={Thrust Coefficient},
  name={Thrust Coefficient},
  description={\textcolor{catRM}{\textbf{[Rockets and Missiles]}}},
}

\newglossaryentry{Specific-Impulse-Efficiency}{ sort={Specific Impulse Efficiency},
  name={Specific Impulse Efficiency},
  description={\textcolor{catRM}{\textbf{[Rockets and Missiles]}}},
}

\newglossaryentry{Grain}{ sort={Grain},
  name={Grain},
  description={\textcolor{catRM}{\textbf{[Rockets and Missiles]}}},
}

\newglossaryentry{Grain-Configuration}{ sort={Grain Configuration},
  name={Grain Configuration},
  description={\textcolor{catRM}{\textbf{[Rockets and Missiles]}}},
}

\newglossaryentry{Burning-Surface}{ sort={Burning Surface},
  name={Burning Surface},
  description={\textcolor{catRM}{\textbf{[Rockets and Missiles]}}},
}

\newglossaryentry{Neutral-Burning}{ sort={Neutral Burning},
  name={Neutral Burning},
  description={\textcolor{catRM}{\textbf{[Rockets and Missiles]}}},
}

\newglossaryentry{Progressive-Burning}{ sort={Progressive Burning},
  name={Progressive Burning},
  description={\textcolor{catRM}{\textbf{[Rockets and Missiles]}}},
}

\newglossaryentry{Regressive-Burning}{ sort={Regressive Burning},
  name={Regressive Burning},
  description={\textcolor{catRM}{\textbf{[Rockets and Missiles]}}},
}

\newglossaryentry{End-Burning}{ sort={End Burning},
  name={End Burning},
  description={\textcolor{catRM}{\textbf{[Rockets and Missiles]}}},
}

\newglossaryentry{Internal-Burning}{ sort={Internal Burning},
  name={Internal Burning},
  description={\textcolor{catRM}{\textbf{[Rockets and Missiles]}}},
}

\newglossaryentry{Cigarette-Burning}{ sort={Cigarette Burning},
  name={Cigarette Burning},
  description={\textcolor{catRM}{\textbf{[Rockets and Missiles]}}},
}

\newglossaryentry{Star-Grain}{ sort={Star Grain},
  name={Star Grain},
  description={\textcolor{catRM}{\textbf{[Rockets and Missiles]}}},
}

\newglossaryentry{Wagon-Wheel-Grain}{ sort={Wagon Wheel Grain},
  name={Wagon Wheel Grain},
  description={\textcolor{catRM}{\textbf{[Rockets and Missiles]}}},
}

\newglossaryentry{Slotted-Tube-Grain}{ sort={Slotted Tube Grain},
  name={Slotted Tube Grain},
  description={\textcolor{catRM}{\textbf{[Rockets and Missiles]}}},
}

\newglossaryentry{Inhibitor}{ sort={Inhibitor},
  name={Inhibitor},
  description={\textcolor{catRM}{\textbf{[Rockets and Missiles]}}},
}

\newglossaryentry{Propellant-Liner}{ sort={Propellant Liner},
  name={Propellant Liner},
  description={\textcolor{catRM}{\textbf{[Rockets and Missiles]}}},
}

\newglossaryentry{Case}{ sort={Case},
  name={Case},
  description={\textcolor{catRM}{\textbf{[Rockets and Missiles]}}},
}

\newglossaryentry{Motor-Case}{ sort={Motor Case},
  name={Motor Case},
  description={\textcolor{catRM}{\textbf{[Rockets and Missiles]}}},
}

\newglossaryentry{Insulation}{ sort={Insulation},
  name={Insulation},
  description={\textcolor{catRM}{\textbf{[Rockets and Missiles]}}},
}

\newglossaryentry{Nozzle-Throat-Insert}{ sort={Nozzle Throat Insert},
  name={Nozzle Throat Insert},
  description={\textcolor{catRM}{\textbf{[Rockets and Missiles]}}},
}

\newglossaryentry{Erosion}{ sort={Erosion},
  name={Erosion},
  description={\textcolor{catRM}{\textbf{[Rockets and Missiles]}}},
}

\newglossaryentry{Graphite-Nozzle}{ sort={Graphite Nozzle},
  name={Graphite Nozzle},
  description={\textcolor{catRM}{\textbf{[Rockets and Missiles]}}},
}

\newglossaryentry{Carbon-Carbon-Nozzle}{ sort={Carbon-Carbon Nozzle},
  name={Carbon-Carbon Nozzle},
  description={\textcolor{catRM}{\textbf{[Rockets and Missiles]}}},
}

\newglossaryentry{Metal-Nozzle}{ sort={Metal Nozzle},
  name={Metal Nozzle},
  description={\textcolor{catRM}{\textbf{[Rockets and Missiles]}}},
}

\newglossaryentry{Nozzle-Closure}{ sort={Nozzle Closure},
  name={Nozzle Closure},
  description={\textcolor{catRM}{\textbf{[Rockets and Missiles]}}},
}

\newglossaryentry{Igniter-Grain}{ sort={Igniter Grain},
  name={Igniter Grain},
  description={\textcolor{catRM}{\textbf{[Rockets and Missiles]}}},
}

\newglossaryentry{Pyrogen-Igniter}{ sort={Pyrogen Igniter},
  name={Pyrogen Igniter},
  description={\textcolor{catRM}{\textbf{[Rockets and Missiles]}}},
}

\newglossaryentry{Safe-and-Arm-Device}{ sort={Safe and Arm Device},
  name={Safe and Arm Device},
  description={\textcolor{catRM}{\textbf{[Rockets and Missiles]}}},
}

\newglossaryentry{Thrust-Termination}{ sort={Thrust Termination},
  name={Thrust Termination},
  description={\textcolor{catRM}{\textbf{[Rockets and Missiles]}}},
}

\newglossaryentry{Destruct-System}{ sort={Destruct System},
  name={Destruct System},
  description={\textcolor{catRM}{\textbf{[Rockets and Missiles]}}},
}


\newglossaryentry{Range-Safety}{ sort={Range Safety},
  name={Range Safety},
  description={\textcolor{catRM}{\textbf{[Rockets and Missiles]}}},
}

\newglossaryentry{Missile}{ sort={Missile},
  name={Missile},
  description={\textcolor{catRM}{\textbf{[Rockets and Missiles]}}},
}

\newglossaryentry{Ballistic-Missile}{ sort={Ballistic Missile},
  name={Ballistic Missile},
  description={\textcolor{catRM}{\textbf{[Rockets and Missiles]}}},
}

\newglossaryentry{Cruise-Missile}{ sort={Cruise Missile},
  name={Cruise Missile},
  description={\textcolor{catRM}{\textbf{[Rockets and Missiles]}}},
}

\newglossaryentry{ICBM}{ sort={ICBM},
  name={ICBM},
  description={\textcolor{catRM}{\textbf{[Rockets and Missiles]}}},
}

\newglossaryentry{IRBM}{ sort={IRBM},
  name={IRBM},
  description={\textcolor{catRM}{\textbf{[Rockets and Missiles]}}},
}

\newglossaryentry{SRBM}{ sort={SRBM},
  name={SRBM},
  description={\textcolor{catRM}{\textbf{[Rockets and Missiles]}}},
}

\newglossaryentry{SLBM}{ sort={SLBM},
  name={SLBM},
  description={\textcolor{catRM}{\textbf{[Rockets and Missiles]}}},
}

\newglossaryentry{SAM}{ sort={SAM},
  name={SAM},
  description={\textcolor{catRM}{\textbf{[Rockets and Missiles]}}},
}

\newglossaryentry{AAM}{ sort={AAM},
  name={AAM},
  description={\textcolor{catRM}{\textbf{[Rockets and Missiles]}}},
}

\newglossaryentry{ASM}{ sort={ASM},
  name={ASM},
  description={\textcolor{catRM}{\textbf{[Rockets and Missiles]}}},
}

\newglossaryentry{ATGM}{ sort={ATGM},
  name={ATGM},
  description={\textcolor{catRM}{\textbf{[Rockets and Missiles]}}},
}

\newglossaryentry{Anti-Ship-Missile}{ sort={Anti-Ship Missile},
  name={Anti-Ship Missile},
  description={\textcolor{catRM}{\textbf{[Rockets and Missiles]}}},
}

\newglossaryentry{Hypersonic-Missile}{ sort={Hypersonic Missile},
  name={Hypersonic Missile},
  description={\textcolor{catRM}{\textbf{[Rockets and Missiles]}}},
}

\newglossaryentry{Hypersonic-Glide-Vehicle}{ sort={Hypersonic Glide Vehicle},
  name={Hypersonic Glide Vehicle},
  description={\textcolor{catRM}{\textbf{[Rockets and Missiles]}}},
}

\newglossaryentry{MIRV}{ sort={MIRV},
  name={MIRV},
  description={\textcolor{catRM}{\textbf{[Rockets and Missiles]}}},
}

\newglossaryentry{Reentry-Vehicle}{ sort={Reentry Vehicle},
  name={Reentry Vehicle},
  description={\textcolor{catRM}{\textbf{[Rockets and Missiles]}}},
}

\newglossaryentry{Warhead}{ sort={Warhead},
  name={Warhead},
  description={\textcolor{catRM}{\textbf{[Rockets and Missiles]}}},
}

\newglossaryentry{Conventional-Warhead}{ sort={Conventional Warhead},
  name={Conventional Warhead},
  description={\textcolor{catRM}{\textbf{[Rockets and Missiles]}}},
}

\newglossaryentry{Nuclear-Warhead}{ sort={Nuclear Warhead},
  name={Nuclear Warhead},
  description={\textcolor{catRM}{\textbf{[Rockets and Missiles]}}},
}

\newglossaryentry{Guidance-System}{ sort={Guidance System},
  name={Guidance System},
  description={\textcolor{catRM}{\textbf{[Rockets and Missiles]}}},
}

\newglossaryentry{Inertial-Guidance}{ sort={Inertial Guidance},
  name={Inertial Guidance},
  description={\textcolor{catRM}{\textbf{[Rockets and Missiles]}}},
}

\newglossaryentry{GPS-Guidance}{ sort={GPS Guidance},
  name={GPS Guidance},
  description={\textcolor{catRM}{\textbf{[Rockets and Missiles]}}},
}

\newglossaryentry{Terrain-Contour-Matching}{ sort={Terrain Contour Matching},
  name={Terrain Contour Matching},
  description={\textcolor{catRM}{\textbf{[Rockets and Missiles]}}},
}

\newglossaryentry{Scene-Matching}{ sort={Scene Matching},
  name={Scene Matching},
  description={\textcolor{catRM}{\textbf{[Rockets and Missiles]}}},
}

\newglossaryentry{Active-Radar-Homing}{ sort={Active Radar Homing},
  name={Active Radar Homing},
  description={\textcolor{catRM}{\textbf{[Rockets and Missiles]}}},
}

\newglossaryentry{Semi-Active-Radar-Homing}{ sort={Semi-Active Radar Homing},
  name={Semi-Active Radar Homing},
  description={\textcolor{catRM}{\textbf{[Rockets and Missiles]}}},
}

\newglossaryentry{Passive-Radar-Homing}{ sort={Passive Radar Homing},
  name={Passive Radar Homing},
  description={\textcolor{catRM}{\textbf{[Rockets and Missiles]}}},
}

\newglossaryentry{Infrared-Homing}{ sort={Infrared Homing},
  name={Infrared Homing},
  description={\textcolor{catRM}{\textbf{[Rockets and Missiles]}}},
}

\newglossaryentry{Laser-Homing}{ sort={Laser Homing},
  name={Laser Homing},
  description={\textcolor{catRM}{\textbf{[Rockets and Missiles]}}},
}

\newglossaryentry{Command-Guidance}{ sort={Command Guidance},
  name={Command Guidance},
  description={\textcolor{catRM}{\textbf{[Rockets and Missiles]}}},
}

\newglossaryentry{Beam-Riding}{ sort={Beam Riding},
  name={Beam Riding},
  description={\textcolor{catRM}{\textbf{[Rockets and Missiles]}}},
}

\newglossaryentry{Wire-Guidance}{ sort={Wire Guidance},
  name={Wire Guidance},
  description={\textcolor{catRM}{\textbf{[Rockets and Missiles]}}},
}

\newglossaryentry{Fiber-Optic-Guidance}{ sort={Fiber Optic Guidance},
  name={Fiber Optic Guidance},
  description={\textcolor{catRM}{\textbf{[Rockets and Missiles]}}},
}

\newglossaryentry{Terminal-Guidance}{ sort={Terminal Guidance},
  name={Terminal Guidance},
  description={\textcolor{catRM}{\textbf{[Rockets and Missiles]}}},
}

\newglossaryentry{Midcourse-Guidance}{ sort={Midcourse Guidance},
  name={Midcourse Guidance},
  description={\textcolor{catRM}{\textbf{[Rockets and Missiles]}}},
}

\newglossaryentry{Boost-Phase}{ sort={Boost Phase},
  name={Boost Phase},
  description={\textcolor{catRM}{\textbf{[Rockets and Missiles]}}},
}

\newglossaryentry{Coast-Phase}{ sort={Coast Phase},
  name={Coast Phase},
  description={\textcolor{catRM}{\textbf{[Rockets and Missiles]}}},
}

\newglossaryentry{Reentry-Phase}{ sort={Reentry Phase},
  name={Reentry Phase},
  description={\textcolor{catRM}{\textbf{[Rockets and Missiles]}}},
}

\newglossaryentry{Trajectory}{ sort={Trajectory},
  name={Trajectory},
  description={\textcolor{catRM}{\textbf{[Rockets and Missiles]}}},
}

\newglossaryentry{Ballistic-Trajectory}{ sort={Ballistic Trajectory},
  name={Ballistic Trajectory},
  description={\textcolor{catRM}{\textbf{[Rockets and Missiles]}}},
}

\newglossaryentry{Lifting-Trajectory}{ sort={Lifting Trajectory},
  name={Lifting Trajectory},
  description={\textcolor{catRM}{\textbf{[Rockets and Missiles]}}},
}

\newglossaryentry{Depressed-Trajectory}{ sort={Depressed Trajectory},
  name={Depressed Trajectory},
  description={\textcolor{catRM}{\textbf{[Rockets and Missiles]}}},
}

\newglossaryentry{Lofted-Trajectory}{ sort={Lofted Trajectory},
  name={Lofted Trajectory},
  description={\textcolor{catRM}{\textbf{[Rockets and Missiles]}}},
}

\newglossaryentry{Minimum-Energy-Trajectory}{ sort={Minimum Energy Trajectory},
  name={Minimum Energy Trajectory},
  description={\textcolor{catRM}{\textbf{[Rockets and Missiles]}}},
}

\newglossaryentry{Maximum-Range-Trajectory}{ sort={Maximum Range Trajectory},
  name={Maximum Range Trajectory},
  description={\textcolor{catRM}{\textbf{[Rockets and Missiles]}}},
}

\newglossaryentry{Launch-Azimuth}{ sort={Launch Azimuth},
  name={Launch Azimuth},
  description={\textcolor{catRM}{\textbf{[Rockets and Missiles]}}},
}

\newglossaryentry{Impact-Point}{ sort={Impact Point},
  name={Impact Point},
  description={\textcolor{catRM}{\textbf{[Rockets and Missiles]}}},
}

\newglossaryentry{CEP}{ sort={CEP},
  name={CEP},
  description={\textcolor{catRM}{\textbf{[Rockets and Missiles]}}},
}

\newglossaryentry{Kill-Probability}{ sort={Kill Probability},
  name={Kill Probability},
  description={\textcolor{catRM}{\textbf{[Rockets and Missiles]}}},
}

\newglossaryentry{Lethality}{ sort={Lethality},
  name={Lethality},
  description={\textcolor{catRM}{\textbf{[Rockets and Missiles]}}},
}

\newglossaryentry{Fragmentation-Warhead}{ sort={Fragmentation Warhead},
  name={Fragmentation Warhead},
  description={\textcolor{catRM}{\textbf{[Rockets and Missiles]}}},
}

\newglossaryentry{Shaped-Charge}{ sort={Shaped Charge},
  name={Shaped Charge},
  description={\textcolor{catRM}{\textbf{[Rockets and Missiles]}}},
}

\newglossaryentry{Kinetic-Kill-Vehicle}{ sort={Kinetic Kill Vehicle},
  name={Kinetic Kill Vehicle},
  description={\textcolor{catRM}{\textbf{[Rockets and Missiles]}}},
}

\newglossaryentry{Decoy}{ sort={Decoy},
  name={Decoy},
  description={\textcolor{catRM}{\textbf{[Rockets and Missiles]}}},
}

\newglossaryentry{Countermeasure}{ sort={Countermeasure},
  name={Countermeasure},
  description={\textcolor{catRM}{\textbf{[Rockets and Missiles]}}},
}

\newglossaryentry{Chaff}{ sort={Chaff},
  name={Chaff},
  description={\textcolor{catRM}{\textbf{[Rockets and Missiles]}}},
}

\newglossaryentry{Flare}{ sort={Flare},
  name={Flare},
  description={\textcolor{catRM}{\textbf{[Rockets and Missiles]}}},
}

\newglossaryentry{Electronic-Countermeasure}{ sort={Electronic Countermeasure},
  name={Electronic Countermeasure},
  description={\textcolor{catRM}{\textbf{[Rockets and Missiles]}}},
}

\newglossaryentry{Penetration-Aid}{ sort={Penetration Aid},
  name={Penetration Aid},
  description={\textcolor{catRM}{\textbf{[Rockets and Missiles]}}},
}

\newglossaryentry{Anti-Ballistic-Missile}{ sort={Anti-Ballistic Missile},
  name={Anti-Ballistic Missile},
  description={\textcolor{catRM}{\textbf{[Rockets and Missiles]}}},
}

\newglossaryentry{Interceptor}{ sort={Interceptor},
  name={Interceptor},
  description={\textcolor{catRM}{\textbf{[Rockets and Missiles]}}},
}

\newglossaryentry{Hit-to-Kill}{ sort={Hit-to-Kill},
  name={Hit-to-Kill},
  description={\textcolor{catRM}{\textbf{[Rockets and Missiles]}}},
}

\newglossaryentry{Proximity-Fuse}{ sort={Proximity Fuse},
  name={Proximity Fuse},
  description={\textcolor{catRM}{\textbf{[Rockets and Missiles]}}},
}

\newglossaryentry{Contact-Fuse}{ sort={Contact Fuse},
  name={Contact Fuse},
  description={\textcolor{catRM}{\textbf{[Rockets and Missiles]}}},
}

\newglossaryentry{Time-Fuse}{ sort={Time Fuse},
  name={Time Fuse},
  description={\textcolor{catRM}{\textbf{[Rockets and Missiles]}}},
}

\newglossaryentry{Airburst}{ sort={Airburst},
  name={Airburst},
  description={\textcolor{catRM}{\textbf{[Rockets and Missiles]}}},
}

\newglossaryentry{Terminal-Phase}{ sort={Terminal Phase},
  name={Terminal Phase},
  description={\textcolor{catRM}{\textbf{[Rockets and Missiles]}}},
}

\newglossaryentry{Launch-Detection}{ sort={Launch Detection},
  name={Launch Detection},
  description={\textcolor{catRM}{\textbf{[Rockets and Missiles]}}},
}

\newglossaryentry{Tracking-Radar}{ sort={Tracking Radar},
  name={Tracking Radar},
  description={\textcolor{catRM}{\textbf{[Rockets and Missiles]}}},
}

\newglossaryentry{Fire-Control-Radar}{ sort={Fire Control Radar},
  name={Fire Control Radar},
  description={\textcolor{catRM}{\textbf{[Rockets and Missiles]}}},
}

\newglossaryentry{Phased-Array-Radar}{ sort={Phased Array Radar},
  name={Phased Array Radar},
  description={\textcolor{catRM}{\textbf{[Rockets and Missiles]}}},
}

\newglossaryentry{Early-Warning-Radar}{ sort={Early Warning Radar},
  name={Early Warning Radar},
  description={\textcolor{catRM}{\textbf{[Rockets and Missiles]}}},
}

\newglossaryentry{Over-the-Horizon-Radar}{ sort={Over-the-Horizon Radar},
  name={Over-the-Horizon Radar},
  description={\textcolor{catRM}{\textbf{[Rockets and Missiles]}}},
}

\newglossaryentry{Satellite-Warning-System}{ sort={Satellite Warning System},
  name={Satellite Warning System},
  description={\textcolor{catRM}{\textbf{[Rockets and Missiles]}}},
}

\newglossaryentry{Ballistic-Missile-Defense}{ sort={Ballistic Missile Defense},
  name={Ballistic Missile Defense},
  description={\textcolor{catRM}{\textbf{[Rockets and Missiles]}}},
}

\newglossaryentry{THAAD}{ sort={THAAD},
  name={THAAD},
  description={\textcolor{catRM}{\textbf{[Rockets and Missiles]}}},
}

\newglossaryentry{Patriot}{ sort={Patriot},
  name={Patriot},
  description={\textcolor{catRM}{\textbf{[Rockets and Missiles]}}},
}

\newglossaryentry{Aegis}{ sort={Aegis},
  name={Aegis},
  description={\textcolor{catRM}{\textbf{[Rockets and Missiles]}}},
}

\newglossaryentry{Ground-Based-Interceptor}{ sort={Ground-Based Interceptor},
  name={Ground-Based Interceptor},
  description={\textcolor{catRM}{\textbf{[Rockets and Missiles]}}},
}

\newglossaryentry{Iron-Dome}{ sort={Iron Dome},
  name={Iron Dome},
  description={\textcolor{catRM}{\textbf{[Rockets and Missiles]}}},
}

\newglossaryentry{David-Sling}{ sort={David Sling},
  name={David Sling},
  description={\textcolor{catRM}{\textbf{[Rockets and Missiles]}}},
}

\newglossaryentry{Arrow}{ sort={Arrow},
  name={Arrow},
  description={\textcolor{catRM}{\textbf{[Rockets and Missiles]}}},
}

\newglossaryentry{S-400}{ sort={S-400},
  name={S-400},
  description={\textcolor{catRM}{\textbf{[Rockets and Missiles]}}},
}

\newglossaryentry{S-500}{ sort={S-500},
  name={S-500},
  description={\textcolor{catRM}{\textbf{[Rockets and Missiles]}}},
}

\newglossaryentry{Ramjet-Missile}{ sort={Ramjet Missile},
  name={Ramjet Missile},
  description={\textcolor{catRM}{\textbf{[Rockets and Missiles]}}},
}

\newglossaryentry{Scramjet-Missile}{ sort={Scramjet Missile},
  name={Scramjet Missile},
  description={\textcolor{catRM}{\textbf{[Rockets and Missiles]}}},
}

\newglossaryentry{Solid-Rocket-Booster}{ sort={Solid Rocket Booster},
  name={Solid Rocket Booster},
  description={\textcolor{catRM}{\textbf{[Rockets and Missiles]}}},
}

\newglossaryentry{Dual-Thrust-Motor}{ sort={Dual Thrust Motor},
  name={Dual Thrust Motor},
  description={\textcolor{catRM}{\textbf{[Rockets and Missiles]}}},
}

\newglossaryentry{Boost-Sustain-Motor}{ sort={Boost-Sustain Motor},
  name={Boost-Sustain Motor},
  description={\textcolor{catRM}{\textbf{[Rockets and Missiles]}}},
}

\newglossaryentry{Thrust-Profile}{ sort={Thrust Profile},
  name={Thrust Profile},
  description={\textcolor{catRM}{\textbf{[Rockets and Missiles]}}},
}

\newglossaryentry{Flame-Deflector}{ sort={Flame Deflector},
  name={Flame Deflector},
  description={\textcolor{catRM}{\textbf{[Rockets and Missiles]}}},
}

\newglossaryentry{Exhaust-Plume}{ sort={Exhaust Plume},
  name={Exhaust Plume},
  description={\textcolor{catRM}{\textbf{[Rockets and Missiles]}}},
}

\newglossaryentry{Mach-Diamond}{ sort={Mach Diamond},
  name={Mach Diamond},
  description={\textcolor{catRM}{\textbf{[Rockets and Missiles]}}},
}

\newglossaryentry{Sounding-Rocket}{ sort={Sounding Rocket},
  name={Sounding Rocket},
  description={\textcolor{catRM}{\textbf{[Rockets and Missiles]}}},
}

\newglossaryentry{Meteorological-Rocket}{ sort={Meteorological Rocket},
  name={Meteorological Rocket},
  description={\textcolor{catRM}{\textbf{[Rockets and Missiles]}}},
}

\newglossaryentry{Target-Drone}{ sort={Target Drone},
  name={Target Drone},
  description={\textcolor{catRM}{\textbf{[Rockets and Missiles]}}},
}

\newglossaryentry{Burn-Rate}{ sort={Burn Rate},
  name={Burn Rate},
  description={\textcolor{catRM}{\textbf{[Rockets and Missiles]}}},
}

\newglossaryentry{Strand-Burn-Rate}{ sort={Strand Burn Rate},
  name={Strand Burn Rate},
  description={\textcolor{catRM}{\textbf{[Rockets and Missiles]}}},
}

\newglossaryentry{Exhaust-Trail}{ sort={Exhaust Trail},
  name={Exhaust Trail},
  description={\textcolor{catRM}{\textbf{[Rockets and Missiles]}}},
}

\newglossaryentry{Liquid-Injection-Thrust-Vectoring}{ sort={Liquid Injection Thrust Vectoring},
  name={Liquid Injection Thrust Vectoring},
  description={\textcolor{catRM}{\textbf{[Rockets and Missiles]}}},
}

\newglossaryentry{Jet-Vanes}{ sort={Jet Vanes},
  name={Jet Vanes},
  description={\textcolor{catRM}{\textbf{[Rockets and Missiles]}}},
}

\newglossaryentry{Jet-Tab}{ sort={Jet Tab},
  name={Jet Tab},
  description={\textcolor{catRM}{\textbf{[Rockets and Missiles]}}},
}

\newglossaryentry{Flexible-Bearing}{ sort={Flexible Bearing},
  name={Flexible Bearing},
  description={\textcolor{catRM}{\textbf{[Rockets and Missiles]}}},
}

\newglossaryentry{Hot-Gas-Valve}{ sort={Hot Gas Valve},
  name={Hot Gas Valve},
  description={\textcolor{catRM}{\textbf{[Rockets and Missiles]}}},
}

\newglossaryentry{Secondary-Injection}{ sort={Secondary Injection},
  name={Secondary Injection},
  description={\textcolor{catRM}{\textbf{[Rockets and Missiles]}}},
}

\newglossaryentry{Main-Propellant-Valve}{ sort={Main Propellant Valve},
  name={Main Propellant Valve},
  description={\textcolor{catRM}{\textbf{[Rockets and Missiles]}}},
}

\newglossaryentry{Fill-Valve}{ sort={Fill Valve},
  name={Fill Valve},
  description={\textcolor{catRM}{\textbf{[Rockets and Missiles]}}},
}

\newglossaryentry{Drain-Valve}{ sort={Drain Valve},
  name={Drain Valve},
  description={\textcolor{catRM}{\textbf{[Rockets and Missiles]}}},
}

\newglossaryentry{Vent-Valve}{ sort={Vent Valve},
  name={Vent Valve},
  description={\textcolor{catRM}{\textbf{[Rockets and Missiles]}}},
}

\newglossaryentry{Relief-Valve}{ sort={Relief Valve},
  name={Relief Valve},
  description={\textcolor{catRM}{\textbf{[Rockets and Missiles]}}},
}

\newglossaryentry{Check-Valve}{ sort={Check Valve},
  name={Check Valve},
  description={\textcolor{catRM}{\textbf{[Rockets and Missiles]}}},
}

\newglossaryentry{Solenoid-Valve}{ sort={Solenoid Valve},
  name={Solenoid Valve},
  description={\textcolor{catRM}{\textbf{[Rockets and Missiles]}}},
}

\newglossaryentry{Poppet-Valve}{ sort={Poppet Valve},
  name={Poppet Valve},
  description={\textcolor{catRM}{\textbf{[Rockets and Missiles]}}},
}

\newglossaryentry{Ball-Valve}{ sort={Ball Valve},
  name={Ball Valve},
  description={\textcolor{catRM}{\textbf{[Rockets and Missiles]}}},
}

\newglossaryentry{Butterfly-Valve}{ sort={Butterfly Valve},
  name={Butterfly Valve},
  description={\textcolor{catRM}{\textbf{[Rockets and Missiles]}}},
}

\newglossaryentry{Pneumatic-Valve}{ sort={Pneumatic Valve},
  name={Pneumatic Valve},
  description={\textcolor{catRM}{\textbf{[Rockets and Missiles]}}},
}

\newglossaryentry{Pyro-Valve}{ sort={Pyro Valve},
  name={Pyro Valve},
  description={\textcolor{catRM}{\textbf{[Rockets and Missiles]}}},
}

\newglossaryentry{Propellant-Line}{ sort={Propellant Line},
  name={Propellant Line},
  description={\textcolor{catRM}{\textbf{[Rockets and Missiles]}}},
}

\newglossaryentry{Suction-Line}{ sort={Suction Line},
  name={Suction Line},
  description={\textcolor{catRM}{\textbf{[Rockets and Missiles]}}},
}

\newglossaryentry{Discharge-Line}{ sort={Discharge Line},
  name={Discharge Line},
  description={\textcolor{catRM}{\textbf{[Rockets and Missiles]}}},
}

\newglossaryentry{Helium-Pressurization}{ sort={Helium Pressurization},
  name={Helium Pressurization},
  description={\textcolor{catRM}{\textbf{[Rockets and Missiles]}}},
}

\newglossaryentry{Nitrogen-Pressurization}{ sort={Nitrogen Pressurization},
  name={Nitrogen Pressurization},
  description={\textcolor{catRM}{\textbf{[Rockets and Missiles]}}},
}

\newglossaryentry{Autogenous-Pressurization}{ sort={Autogenous Pressurization},
  name={Autogenous Pressurization},
  description={\textcolor{catRM}{\textbf{[Rockets and Missiles]}}},
}

\newglossaryentry{Accumulator}{ sort={Accumulator},
  name={Accumulator},
  description={\textcolor{catRM}{\textbf{[Rockets and Missiles]}}},
}

\newglossaryentry{Pressure-Regulator}{ sort={Pressure Regulator},
  name={Pressure Regulator},
  description={\textcolor{catRM}{\textbf{[Rockets and Missiles]}}},
}

\newglossaryentry{Gas-Generator}{ sort={Gas Generator},
  name={Gas Generator},
  description={\textcolor{catRM}{\textbf{[Rockets and Missiles]}}},
}

\newglossaryentry{Heat-Exchanger}{ sort={Heat Exchanger},
  name={Heat Exchanger},
  description={\textcolor{catRM}{\textbf{[Rockets and Missiles]}}},
}

\newglossaryentry{Gas-Cooler}{ sort={Gas Cooler},
  name={Gas Cooler},
  description={\textcolor{catRM}{\textbf{[Rockets and Missiles]}}},
}

\newglossaryentry{Filter}{ sort={Filter},
  name={Filter},
  description={\textcolor{catRM}{\textbf{[Rockets and Missiles]}}},
}

\newglossaryentry{Propellant-Management}{ sort={Propellant Management},
  name={Propellant Management},
  description={\textcolor{catRM}{\textbf{[Rockets and Missiles]}}},
}

\newglossaryentry{Propellant-Settling}{ sort={Propellant Settling},
  name={Propellant Settling},
  description={\textcolor{catRM}{\textbf{[Rockets and Missiles]}}},
}

\newglossaryentry{Fin-Stabilized-Rocket}{ sort={Fin-Stabilized Rocket},
  name={Fin-Stabilized Rocket},
  description={\textcolor{catRM}{\textbf{[Rockets and Missiles]}}},
}

\newglossaryentry{Spin-Stabilized-Rocket}{ sort={Spin-Stabilized Rocket},
  name={Spin-Stabilized Rocket},
  description={\textcolor{catRM}{\textbf{[Rockets and Missiles]}}},
}

\newglossaryentry{Tail-Fins}{ sort={Tail Fins},
  name={Tail Fins},
  description={\textcolor{catRM}{\textbf{[Rockets and Missiles]}}},
}

\newglossaryentry{Canard-Control}{ sort={Canard Control},
  name={Canard Control},
  description={\textcolor{catRM}{\textbf{[Rockets and Missiles]}}},
}

\newglossaryentry{Drone}{ sort={Drone},
  name={Drone},
  description={\textcolor{catD}{\textbf{[Drones]}}},
}

\newglossaryentry{UAV}{ sort={UAV},
  name={UAV},
  description={\textcolor{catD}{\textbf{[Drones]}}},
}

\newglossaryentry{UAS}{ sort={UAS},
  name={UAS},
  description={\textcolor{catD}{\textbf{[Drones]}}},
}

\newglossaryentry{RPAS}{ sort={RPAS},
  name={RPAS},
  description={\textcolor{catD}{\textbf{[Drones]}}},
}

\newglossaryentry{Multirotor}{ sort={Multirotor},
  name={Multirotor},
  description={\textcolor{catD}{\textbf{[Drones]}}},
}

\newglossaryentry{Quadcopter}{ sort={Quadcopter},
  name={Quadcopter},
  description={\textcolor{catD}{\textbf{[Drones]}}},
}

\newglossaryentry{Hexacopter}{ sort={Hexacopter},
  name={Hexacopter},
  description={\textcolor{catD}{\textbf{[Drones]}}},
}

\newglossaryentry{Octocopter}{ sort={Octocopter},
  name={Octocopter},
  description={\textcolor{catD}{\textbf{[Drones]}}},
}

\newglossaryentry{Fixed-Wing-Drone}{ sort={Fixed-Wing Drone},
  name={Fixed-Wing Drone},
  description={\textcolor{catD}{\textbf{[Drones]}}},
}

\newglossaryentry{VTOL-Drone}{ sort={VTOL Drone},
  name={VTOL Drone},
  description={\textcolor{catD}{\textbf{[Drones]}}},
}

\newglossaryentry{Hybrid-Drone}{ sort={Hybrid Drone},
  name={Hybrid Drone},
  description={\textcolor{catD}{\textbf{[Drones]}}},
}

\newglossaryentry{Nano-Drone}{ sort={Nano Drone},
  name={Nano Drone},
  description={\textcolor{catD}{\textbf{[Drones]}}},
}

\newglossaryentry{Micro-Drone}{ sort={Micro Drone},
  name={Micro Drone},
  description={\textcolor{catD}{\textbf{[Drones]}}},
}

\newglossaryentry{Tactical-Drone}{ sort={Tactical Drone},
  name={Tactical Drone},
  description={\textcolor{catD}{\textbf{[Drones]}}},
}

\newglossaryentry{MALE-Drone}{ sort={MALE Drone},
  name={MALE Drone},
  description={\textcolor{catD}{\textbf{[Drones]}}},
}

\newglossaryentry{HALE-Drone}{ sort={HALE Drone},
  name={HALE Drone},
  description={\textcolor{catD}{\textbf{[Drones]}}},
}

\newglossaryentry{Combat-Drone}{ sort={Combat Drone},
  name={Combat Drone},
  description={\textcolor{catD}{\textbf{[Drones]}}},
}

\newglossaryentry{Reconnaissance-Drone}{ sort={Reconnaissance Drone},
  name={Reconnaissance Drone},
  description={\textcolor{catD}{\textbf{[Drones]}}},
}

\newglossaryentry{Swarm-Drone}{ sort={Swarm Drone},
  name={Swarm Drone},
  description={\textcolor{catD}{\textbf{[Drones]}}},
}

\newglossaryentry{Delivery-Drone}{ sort={Delivery Drone},
  name={Delivery Drone},
  description={\textcolor{catD}{\textbf{[Drones]}}},
}

\newglossaryentry{Agriculture-Drone}{ sort={Agriculture Drone},
  name={Agriculture Drone},
  description={\textcolor{catD}{\textbf{[Drones]}}},
}

\newglossaryentry{Surveying-Drone}{ sort={Surveying Drone},
  name={Surveying Drone},
  description={\textcolor{catD}{\textbf{[Drones]}}},
}

\newglossaryentry{Inspection-Drone}{ sort={Inspection Drone},
  name={Inspection Drone},
  description={\textcolor{catD}{\textbf{[Drones]}}},
}

\newglossaryentry{Search-and-Rescue-Drone}{ sort={Search and Rescue Drone},
  name={Search and Rescue Drone},
  description={\textcolor{catD}{\textbf{[Drones]}}},
}

\newglossaryentry{Racing-Drone}{ sort={Racing Drone},
  name={Racing Drone},
  description={\textcolor{catD}{\textbf{[Drones]}}},
}

\newglossaryentry{FPV-Drone}{ sort={FPV Drone},
  name={FPV Drone},
  description={\textcolor{catD}{\textbf{[Drones]}}},
}

\newglossaryentry{Cinematic-Drone}{ sort={Cinematic Drone},
  name={Cinematic Drone},
  description={\textcolor{catD}{\textbf{[Drones]}}},
}

\newglossaryentry{Submersible-Drone}{ sort={Submersible Drone},
  name={Submersible Drone},
  description={\textcolor{catD}{\textbf{[Drones]}}},
}

\newglossaryentry{Autonomous-Drone}{ sort={Autonomous Drone},
  name={Autonomous Drone},
  description={\textcolor{catD}{\textbf{[Drones]}}},
}

\newglossaryentry{Tethered-Drone}{ sort={Tethered Drone},
  name={Tethered Drone},
  description={\textcolor{catD}{\textbf{[Drones]}}},
}

\newglossaryentry{Solar-Drone}{ sort={Solar Drone},
  name={Solar Drone},
  description={\textcolor{catD}{\textbf{[Drones]}}},
}

\newglossaryentry{Flight-Controller}{ sort={Flight Controller},
  name={Flight Controller},
  description={\textcolor{catD}{\textbf{[Drones]}}},
}

\newglossaryentry{IMU}{ sort={IMU},
  name={IMU},
  description={\textcolor{catD}{\textbf{[Drones]}}},
}

\newglossaryentry{Accelerometer}{ sort={Accelerometer},
  name={Accelerometer},
  description={\textcolor{catD}{\textbf{[Drones]}}},
}



\newglossaryentry{Barometer}{ sort={Barometer},
  name={Barometer},
  description={\textcolor{catD}{\textbf{[Drones]}}},
}

\newglossaryentry{GPS-Module}{ sort={GPS Module},
  name={GPS Module},
  description={\textcolor{catD}{\textbf{[Drones]}}},
}

\newglossaryentry{GLONASS}{ sort={GLONASS},
  name={GLONASS},
  description={\textcolor{catD}{\textbf{[Drones]}}},
}

\newglossaryentry{Compass}{ sort={Compass},
  name={Compass},
  description={\textcolor{catD}{\textbf{[Drones]}}},
}

\newglossaryentry{ESC}{ sort={ESC},
  name={ESC},
  description={\textcolor{catD}{\textbf{[Drones]}}},
}

\newglossaryentry{Motor}{ sort={Motor},
  name={Motor},
  description={\textcolor{catD}{\textbf{[Drones]}}},
}

\newglossaryentry{Brushless-Motor}{ sort={Brushless Motor},
  name={Brushless Motor},
  description={\textcolor{catD}{\textbf{[Drones]}}},
}

\newglossaryentry{Propeller}{ sort={Propeller},
  name={Propeller},
  description={\textcolor{catD}{\textbf{[Drones]}}},
}

\newglossaryentry{Pitch-Propeller}{ sort={Pitch Propeller},
  name={Pitch Propeller},
  description={\textcolor{catD}{\textbf{[Drones]}}},
}

\newglossaryentry{Yaw-Propeller}{ sort={Yaw Propeller},
  name={Yaw Propeller},
  description={\textcolor{catD}{\textbf{[Drones]}}},
}

\newglossaryentry{LiPo-Battery}{ sort={LiPo Battery},
  name={LiPo Battery},
  description={\textcolor{catD}{\textbf{[Drones]}}},
}

\newglossaryentry{Battery-Capacity}{ sort={Battery Capacity},
  name={Battery Capacity},
  description={\textcolor{catD}{\textbf{[Drones]}}},
}

\newglossaryentry{S-Rating}{ sort={S Rating},
  name={S Rating},
  description={\textcolor{catD}{\textbf{[Drones]}}},
}

\newglossaryentry{C-Rating}{ sort={C Rating},
  name={C Rating},
  description={\textcolor{catD}{\textbf{[Drones]}}},
}

\newglossaryentry{Power-Distribution-Board}{ sort={Power Distribution Board},
  name={Power Distribution Board},
  description={\textcolor{catD}{\textbf{[Drones]}}},
}

\newglossaryentry{PDB}{ sort={PDB},
  name={PDB},
  description={\textcolor{catD}{\textbf{[Drones]}}},
}

\newglossaryentry{BEC}{ sort={BEC},
  name={BEC},
  description={\textcolor{catD}{\textbf{[Drones]}}},
}


\newglossaryentry{Receiver}{ sort={Receiver},
  name={Receiver},
  description={\textcolor{catD}{\textbf{[Drones]}}},
}

\newglossaryentry{Transmitter}{ sort={Transmitter},
  name={Transmitter},
  description={\textcolor{catD}{\textbf{[Drones]}}},
}

\newglossaryentry{Radio-Link}{ sort={Radio Link},
  name={Radio Link},
  description={\textcolor{catD}{\textbf{[Drones]}}},
}

\newglossaryentry{Control-Link}{ sort={Control Link},
  name={Control Link},
  description={\textcolor{catD}{\textbf{[Drones]}}},
}

\newglossaryentry{Video-Link}{ sort={Video Link},
  name={Video Link},
  description={\textcolor{catD}{\textbf{[Drones]}}},
}

\newglossaryentry{Data-Link}{ sort={Data Link},
  name={Data Link},
  description={\textcolor{catD}{\textbf{[Drones]}}},
}

\newglossaryentry{Frequency-Band}{ sort={Frequency Band},
  name={Frequency Band},
  description={\textcolor{catD}{\textbf{[Drones]}}},
}

\newglossaryentry{24-GHz}{ sort={2.4 GHz},
  name={2.4 GHz},
  description={\textcolor{catD}{\textbf{[Drones]}}},
}

\newglossaryentry{58-GHz}{ sort={5.8 GHz},
  name={5.8 GHz},
  description={\textcolor{catD}{\textbf{[Drones]}}},
}

\newglossaryentry{900-MHz}{ sort={900 MHz},
  name={900 MHz},
  description={\textcolor{catD}{\textbf{[Drones]}}},
}

\newglossaryentry{433-MHz}{ sort={433 MHz},
  name={433 MHz},
  description={\textcolor{catD}{\textbf{[Drones]}}},
}

\newglossaryentry{FHSS}{ sort={FHSS},
  name={FHSS},
  description={\textcolor{catD}{\textbf{[Drones]}}},
}

\newglossaryentry{DSSS}{ sort={DSSS},
  name={DSSS},
  description={\textcolor{catD}{\textbf{[Drones]}}},
}

\newglossaryentry{LBT}{ sort={LBT},
  name={LBT},
  description={\textcolor{catD}{\textbf{[Drones]}}},
}

\newglossaryentry{RSSI}{ sort={RSSI},
  name={RSSI},
  description={\textcolor{catD}{\textbf{[Drones]}}},
}


\newglossaryentry{Omnidirectional-Antenna}{ sort={Omnidirectional Antenna},
  name={Omnidirectional Antenna},
  description={\textcolor{catD}{\textbf{[Drones]}}},
}

\newglossaryentry{Directional-Antenna}{ sort={Directional Antenna},
  name={Directional Antenna},
  description={\textcolor{catD}{\textbf{[Drones]}}},
}

\newglossaryentry{Patch-Antenna}{ sort={Patch Antenna},
  name={Patch Antenna},
  description={\textcolor{catD}{\textbf{[Drones]}}},
}

\newglossaryentry{Yagi-Antenna}{ sort={Yagi Antenna},
  name={Yagi Antenna},
  description={\textcolor{catD}{\textbf{[Drones]}}},
}

\newglossaryentry{Gimbal}{ sort={Gimbal},
  name={Gimbal},
  description={\textcolor{catD}{\textbf{[Drones]}}},
}

\newglossaryentry{3-Axis-Gimbal}{ sort={3-Axis Gimbal},
  name={3-Axis Gimbal},
  description={\textcolor{catD}{\textbf{[Drones]}}},
}

\newglossaryentry{2-Axis-Gimbal}{ sort={2-Axis Gimbal},
  name={2-Axis Gimbal},
  description={\textcolor{catD}{\textbf{[Drones]}}},
}

\newglossaryentry{PID-Tuning}{ sort={PID Tuning},
  name={PID Tuning},
  description={\textcolor{catD}{\textbf{[Drones]}}},
}

\newglossaryentry{Autotune}{ sort={Autotune},
  name={Autotune},
  description={\textcolor{catD}{\textbf{[Drones]}}},
}

\newglossaryentry{Rate-Mode}{ sort={Rate Mode},
  name={Rate Mode},
  description={\textcolor{catD}{\textbf{[Drones]}}},
}

\newglossaryentry{Angle-Mode}{ sort={Angle Mode},
  name={Angle Mode},
  description={\textcolor{catD}{\textbf{[Drones]}}},
}

\newglossaryentry{Horizon-Mode}{ sort={Horizon Mode},
  name={Horizon Mode},
  description={\textcolor{catD}{\textbf{[Drones]}}},
}

\newglossaryentry{Altitude-Hold}{ sort={Altitude Hold},
  name={Altitude Hold},
  description={\textcolor{catD}{\textbf{[Drones]}}},
}

\newglossaryentry{Position-Hold}{ sort={Position Hold},
  name={Position Hold},
  description={\textcolor{catD}{\textbf{[Drones]}}},
}

\newglossaryentry{Return-to-Home}{ sort={Return to Home},
  name={Return to Home},
  description={\textcolor{catD}{\textbf{[Drones]}}},
}

\newglossaryentry{RTH}{ sort={RTH},
  name={RTH},
  description={\textcolor{catD}{\textbf{[Drones]}}},
}

\newglossaryentry{Loiter}{ sort={Loiter},
  name={Loiter},
  description={\textcolor{catD}{\textbf{[Drones]}}},
}

\newglossaryentry{Waypoint-Navigation}{ sort={Waypoint Navigation},
  name={Waypoint Navigation},
  description={\textcolor{catD}{\textbf{[Drones]}}},
}

\newglossaryentry{Follow-Me}{ sort={Follow Me},
  name={Follow Me},
  description={\textcolor{catD}{\textbf{[Drones]}}},
}

\newglossaryentry{Orbit-Mode}{ sort={Orbit Mode},
  name={Orbit Mode},
  description={\textcolor{catD}{\textbf{[Drones]}}},
}

\newglossaryentry{Mission-Planning}{ sort={Mission Planning},
  name={Mission Planning},
  description={\textcolor{catD}{\textbf{[Drones]}}},
}

\newglossaryentry{Geofence}{ sort={Geofence},
  name={Geofence},
  description={\textcolor{catD}{\textbf{[Drones]}}},
}

\newglossaryentry{Low-Battery-Warning}{ sort={Low Battery Warning},
  name={Low Battery Warning},
  description={\textcolor{catD}{\textbf{[Drones]}}},
}

\newglossaryentry{Collision-Avoidance}{ sort={Collision Avoidance},
  name={Collision Avoidance},
  description={\textcolor{catD}{\textbf{[Drones]}}},
}

\newglossaryentry{Obstacle-Detection}{ sort={Obstacle Detection},
  name={Obstacle Detection},
  description={\textcolor{catD}{\textbf{[Drones]}}},
}

\newglossaryentry{Emergency-Stop}{ sort={Emergency Stop},
  name={Emergency Stop},
  description={\textcolor{catD}{\textbf{[Drones]}}},
}

\newglossaryentry{Failsafe}{ sort={Failsafe},
  name={Failsafe},
  description={\textcolor{catD}{\textbf{[Drones]}}},
}

\newglossaryentry{Flight-Termination-System}{ sort={Flight Termination System},
  name={Flight Termination System},
  description={\textcolor{catD}{\textbf{[Drones]}}},
}

\newglossaryentry{Parachute-System}{ sort={Parachute System},
  name={Parachute System},
  description={\textcolor{catD}{\textbf{[Drones]}}},
}

\newglossaryentry{VLOS}{ sort={VLOS},
  name={VLOS},
  description={\textcolor{catD}{\textbf{[Drones]}}},
}

\newglossaryentry{BVLOS}{ sort={BVLOS},
  name={BVLOS},
  description={\textcolor{catD}{\textbf{[Drones]}}},
}

\newglossaryentry{EVLOS}{ sort={EVLOS},
  name={EVLOS},
  description={\textcolor{catD}{\textbf{[Drones]}}},
}

\newglossaryentry{Remote-ID}{ sort={Remote ID},
  name={Remote ID},
  description={\textcolor{catD}{\textbf{[Drones]}}},
}

\newglossaryentry{No-Fly-Zone}{ sort={No-Fly Zone},
  name={No-Fly Zone},
  description={\textcolor{catD}{\textbf{[Drones]}}},
}

\newglossaryentry{Airspace}{ sort={Airspace},
  name={Airspace},
  description={\textcolor{catD}{\textbf{[Drones]}}},
}

\newglossaryentry{UTM}{ sort={UTM},
  name={UTM},
  description={\textcolor{catD}{\textbf{[Drones]}}},
}

\newglossaryentry{NOTAM}{ sort={NOTAM},
  name={NOTAM},
  description={\textcolor{catD}{\textbf{[Drones]}}},
}

\newglossaryentry{Flight-Authorization}{ sort={Flight Authorization},
  name={Flight Authorization},
  description={\textcolor{catD}{\textbf{[Drones]}}},
}

\newglossaryentry{LAANC}{ sort={LAANC},
  name={LAANC},
  description={\textcolor{catD}{\textbf{[Drones]}}},
}

\newglossaryentry{Part-107}{ sort={Part 107},
  name={Part 107},
  description={\textcolor{catD}{\textbf{[Drones]}}},
}

\newglossaryentry{Drone-Registration}{ sort={Drone Registration},
  name={Drone Registration},
  description={\textcolor{catD}{\textbf{[Drones]}}},
}

\newglossaryentry{Operator-License}{ sort={Operator License},
  name={Operator License},
  description={\textcolor{catD}{\textbf{[Drones]}}},
}

\newglossaryentry{Remote-Pilot-Certificate}{ sort={Remote Pilot Certificate},
  name={Remote Pilot Certificate},
  description={\textcolor{catD}{\textbf{[Drones]}}},
}

\newglossaryentry{Aerial-Photography}{ sort={Aerial Photography},
  name={Aerial Photography},
  description={\textcolor{catD}{\textbf{[Drones]}}},
}

\newglossaryentry{Precision-Agriculture}{ sort={Precision Agriculture},
  name={Precision Agriculture},
  description={\textcolor{catD}{\textbf{[Drones]}}},
}

\newglossaryentry{Crop-Monitoring}{ sort={Crop Monitoring},
  name={Crop Monitoring},
  description={\textcolor{catD}{\textbf{[Drones]}}},
}

\newglossaryentry{Mapping}{ sort={Mapping},
  name={Mapping},
  description={\textcolor{catD}{\textbf{[Drones]}}},
}

\newglossaryentry{Surveying}{ sort={Surveying},
  name={Surveying},
  description={\textcolor{catD}{\textbf{[Drones]}}},
}

\newglossaryentry{Photogrammetry}{ sort={Photogrammetry},
  name={Photogrammetry},
  description={\textcolor{catD}{\textbf{[Drones]}}},
}

\newglossaryentry{Infrastructure-Inspection}{ sort={Infrastructure Inspection},
  name={Infrastructure Inspection},
  description={\textcolor{catD}{\textbf{[Drones]}}},
}

\newglossaryentry{Search-and-Rescue}{ sort={Search and Rescue},
  name={Search and Rescue},
  description={\textcolor{catD}{\textbf{[Drones]}}},
}

\newglossaryentry{Disaster-Response}{ sort={Disaster Response},
  name={Disaster Response},
  description={\textcolor{catD}{\textbf{[Drones]}}},
}

\newglossaryentry{Wildlife-Monitoring}{ sort={Wildlife Monitoring},
  name={Wildlife Monitoring},
  description={\textcolor{catD}{\textbf{[Drones]}}},
}

\newglossaryentry{Parcel-Delivery}{ sort={Parcel Delivery},
  name={Parcel Delivery},
  description={\textcolor{catD}{\textbf{[Drones]}}},
}

\newglossaryentry{Environmental-Monitoring}{ sort={Environmental Monitoring},
  name={Environmental Monitoring},
  description={\textcolor{catD}{\textbf{[Drones]}}},
}

\newglossaryentry{Mission-Planner}{ sort={Mission Planner},
  name={Mission Planner},
  description={\textcolor{catD}{\textbf{[Drones]}}},
}

\newglossaryentry{QGroundControl}{ sort={QGroundControl},
  name={QGroundControl},
  description={\textcolor{catD}{\textbf{[Drones]}}},
}

\newglossaryentry{PX4}{ sort={PX4},
  name={PX4},
  description={\textcolor{catD}{\textbf{[Drones]}}},
}

\newglossaryentry{ArduPilot}{ sort={ArduPilot},
  name={ArduPilot},
  description={\textcolor{catD}{\textbf{[Drones]}}},
}

\newglossaryentry{MAVLink}{ sort={MAVLink},
  name={MAVLink},
  description={\textcolor{catD}{\textbf{[Drones]}}},
}

\newglossaryentry{MAVSDK}{ sort={MAVSDK},
  name={MAVSDK},
  description={\textcolor{catD}{\textbf{[Drones]}}},
}

\newglossaryentry{Gazebo}{ sort={Gazebo},
  name={Gazebo},
  description={\textcolor{catD}{\textbf{[Drones]}}},
}

\newglossaryentry{AirSim}{ sort={AirSim},
  name={AirSim},
  description={\textcolor{catD}{\textbf{[Drones]}}},
}

\newglossaryentry{SITL}{ sort={SITL},
  name={SITL},
  description={\textcolor{catD}{\textbf{[Drones]}}},
}

\newglossaryentry{HITL}{ sort={HITL},
  name={HITL},
  description={\textcolor{catD}{\textbf{[Drones]}}},
}

\newglossaryentry{Digital-Twin}{ sort={Digital Twin},
  name={Digital Twin},
  description={\textcolor{catD}{\textbf{[Drones]}}},
}

\newglossaryentry{ROS}{ sort={ROS},
  name={ROS},
  description={\textcolor{catD}{\textbf{[Drones]}}},
}

\newglossaryentry{DroneKit}{ sort={DroneKit},
  name={DroneKit},
  description={\textcolor{catD}{\textbf{[Drones]}}},
}

\newglossaryentry{Camera-Payload}{ sort={Camera Payload},
  name={Camera Payload},
  description={\textcolor{catD}{\textbf{[Drones]}}},
}

\newglossaryentry{LiDAR-Payload}{ sort={LiDAR Payload},
  name={LiDAR Payload},
  description={\textcolor{catD}{\textbf{[Drones]}}},
}

\newglossaryentry{Sprayer}{ sort={Sprayer},
  name={Sprayer},
  description={\textcolor{catD}{\textbf{[Drones]}}},
}

\newglossaryentry{Cargo-Pod}{ sort={Cargo Pod},
  name={Cargo Pod},
  description={\textcolor{catD}{\textbf{[Drones]}}},
}

\newglossaryentry{Loudspeaker}{ sort={Loudspeaker},
  name={Loudspeaker},
  description={\textcolor{catD}{\textbf{[Drones]}}},
}

\newglossaryentry{Spotlight}{ sort={Spotlight},
  name={Spotlight},
  description={\textcolor{catD}{\textbf{[Drones]}}},
}

\newglossaryentry{Winch}{ sort={Winch},
  name={Winch},
  description={\textcolor{catD}{\textbf{[Drones]}}},
}

\newglossaryentry{Delivery-Mechanism}{ sort={Delivery Mechanism},
  name={Delivery Mechanism},
  description={\textcolor{catD}{\textbf{[Drones]}}},
}

\newglossaryentry{Scientific-Instrument}{ sort={Scientific Instrument},
  name={Scientific Instrument},
  description={\textcolor{catD}{\textbf{[Drones]}}},
}

\newglossaryentry{Hyperspectral-Sensor}{ sort={Hyperspectral Sensor},
  name={Hyperspectral Sensor},
  description={\textcolor{catD}{\textbf{[Drones]}}},
}

\newglossaryentry{Thermal-Camera}{ sort={Thermal Camera},
  name={Thermal Camera},
  description={\textcolor{catD}{\textbf{[Drones]}}},
}

\newglossaryentry{Multispectral-Sensor}{ sort={Multispectral Sensor},
  name={Multispectral Sensor},
  description={\textcolor{catD}{\textbf{[Drones]}}},
}

\newglossaryentry{Flight-Time}{ sort={Flight Time},
  name={Flight Time},
  description={\textcolor{catD}{\textbf{[Drones]}}},
}

\newglossaryentry{Hover-Time}{ sort={Hover Time},
  name={Hover Time},
  description={\textcolor{catD}{\textbf{[Drones]}}},
}

\newglossaryentry{Endurance}{ sort={Endurance},
  name={Endurance},
  description={\textcolor{catD}{\textbf{[Drones]}}},
}

\newglossaryentry{Range}{ sort={Range},
  name={Range},
  description={\textcolor{catD}{\textbf{[Drones]}}},
}

\newglossaryentry{Payload-Capacity}{ sort={Payload Capacity},
  name={Payload Capacity},
  description={\textcolor{catD}{\textbf{[Drones]}}},
}

\newglossaryentry{MTOM}{ sort={MTOM},
  name={MTOM},
  description={\textcolor{catD}{\textbf{[Drones]}}},
}

\newglossaryentry{MTOW}{ sort={MTOW},
  name={MTOW},
  description={\textcolor{catD}{\textbf{[Drones]}}},
}

\newglossaryentry{Service-Ceiling}{ sort={Service Ceiling},
  name={Service Ceiling},
  description={\textcolor{catD}{\textbf{[Drones]}}},
}

\newglossaryentry{Maximum-Speed}{ sort={Maximum Speed},
  name={Maximum Speed},
  description={\textcolor{catD}{\textbf{[Drones]}}},
}

\newglossaryentry{Cruise-Speed}{ sort={Cruise Speed},
  name={Cruise Speed},
  description={\textcolor{catD}{\textbf{[Drones]}}},
}

\newglossaryentry{Wind-Resistance}{ sort={Wind Resistance},
  name={Wind Resistance},
  description={\textcolor{catD}{\textbf{[Drones]}}},
}

\newglossaryentry{IP-Rating}{ sort={IP Rating},
  name={IP Rating},
  description={\textcolor{catD}{\textbf{[Drones]}}},
}

\newglossaryentry{Drone-Swarm}{ sort={Drone Swarm},
  name={Drone Swarm},
  description={\textcolor{catD}{\textbf{[Drones]}}},
}

\newglossaryentry{Swarm-Intelligence}{ sort={Swarm Intelligence},
  name={Swarm Intelligence},
  description={\textcolor{catD}{\textbf{[Drones]}}},
}

\newglossaryentry{Formation-Flight}{ sort={Formation Flight},
  name={Formation Flight},
  description={\textcolor{catD}{\textbf{[Drones]}}},
}

\newglossaryentry{Cooperative-Navigation}{ sort={Cooperative Navigation},
  name={Cooperative Navigation},
  description={\textcolor{catD}{\textbf{[Drones]}}},
}

\newglossaryentry{Distributed-Control}{ sort={Distributed Control},
  name={Distributed Control},
  description={\textcolor{catD}{\textbf{[Drones]}}},
}

\newglossaryentry{Mesh-Network}{ sort={Mesh Network},
  name={Mesh Network},
  description={\textcolor{catD}{\textbf{[Drones]}}},
}

\newglossaryentry{Multi-Agent-System}{ sort={Multi-Agent System},
  name={Multi-Agent System},
  description={\textcolor{catD}{\textbf{[Drones]}}},
}

\newglossaryentry{Consensus-Algorithm}{ sort={Consensus Algorithm},
  name={Consensus Algorithm},
  description={\textcolor{catD}{\textbf{[Drones]}}},
}

\newglossaryentry{Swarm-Coordination}{ sort={Swarm Coordination},
  name={Swarm Coordination},
  description={\textcolor{catD}{\textbf{[Drones]}}},
}

\newglossaryentry{Collision-Avoidance-Swarm}{ sort={Collision Avoidance Swarm},
  name={Collision Avoidance Swarm},
  description={\textcolor{catD}{\textbf{[Drones]}}},
}

\newglossaryentry{FAA}{ sort={FAA},
  name={FAA},
  description={\textcolor{catD}{\textbf{[Drones]}}},
}

\newglossaryentry{EASA}{ sort={EASA},
  name={EASA},
  description={\textcolor{catD}{\textbf{[Drones]}}},
}

\newglossaryentry{ICAO}{ sort={ICAO},
  name={ICAO},
  description={\textcolor{catD}{\textbf{[Drones]}}},
}

\newglossaryentry{ASTM-International}{ sort={ASTM International},
  name={ASTM International},
  description={\textcolor{catD}{\textbf{[Drones]}}},
}

\newglossaryentry{RTCA}{ sort={RTCA},
  name={RTCA},
  description={\textcolor{catD}{\textbf{[Drones]}}},
}

\newglossaryentry{Dronecode-Foundation}{ sort={Dronecode Foundation},
  name={Dronecode Foundation},
  description={\textcolor{catD}{\textbf{[Drones]}}},
}

\newglossaryentry{PX4-Community}{ sort={PX4 Community},
  name={PX4 Community},
  description={\textcolor{catD}{\textbf{[Drones]}}},
}

\newglossaryentry{ArduPilot-Project}{ sort={ArduPilot Project},
  name={ArduPilot Project},
  description={\textcolor{catD}{\textbf{[Drones]}}},
}

\newglossaryentry{Flying-Weight}{ sort={Flying Weight},
  name={Flying Weight},
  description={\textcolor{catD}{\textbf{[Drones]}}},
}

\newglossaryentry{Frame}{ sort={Frame},
  name={Frame},
  description={\textcolor{catD}{\textbf{[Drones]}}},
}

\newglossaryentry{Arm}{ sort={Arm},
  name={Arm},
  description={\textcolor{catD}{\textbf{[Drones]}}},
}

\newglossaryentry{Landing-Skid}{ sort={Landing Skid},
  name={Landing Skid},
  description={\textcolor{catD}{\textbf{[Drones]}}},
}

\newglossaryentry{Motor-Mount}{ sort={Motor Mount},
  name={Motor Mount},
  description={\textcolor{catD}{\textbf{[Drones]}}},
}

\newglossaryentry{Propeller-Guard}{ sort={Propeller Guard},
  name={Propeller Guard},
  description={\textcolor{catD}{\textbf{[Drones]}}},
}

\newglossaryentry{Camera-Mount}{ sort={Camera Mount},
  name={Camera Mount},
  description={\textcolor{catD}{\textbf{[Drones]}}},
}

\newglossaryentry{Vibration-Damping}{ sort={Vibration Damping},
  name={Vibration Damping},
  description={\textcolor{catD}{\textbf{[Drones]}}},
}

\newglossaryentry{Anti-Vibration-Mount}{ sort={Anti-Vibration Mount},
  name={Anti-Vibration Mount},
  description={\textcolor{catD}{\textbf{[Drones]}}},
}

\newglossaryentry{Foam-Damping}{ sort={Foam Damping},
  name={Foam Damping},
  description={\textcolor{catD}{\textbf{[Drones]}}},
}

\newglossaryentry{GPS-Mast}{ sort={GPS Mast},
  name={GPS Mast},
  description={\textcolor{catD}{\textbf{[Drones]}}},
}

\newglossaryentry{Compass-Mount}{ sort={Compass Mount},
  name={Compass Mount},
  description={\textcolor{catD}{\textbf{[Drones]}}},
}

\newglossaryentry{LED-Indicator}{ sort={LED Indicator},
  name={LED Indicator},
  description={\textcolor{catD}{\textbf{[Drones]}}},
}

\newglossaryentry{Buzzer}{ sort={Buzzer},
  name={Buzzer},
  description={\textcolor{catD}{\textbf{[Drones]}}},
}

\newglossaryentry{Current-Sensor}{ sort={Current Sensor},
  name={Current Sensor},
  description={\textcolor{catD}{\textbf{[Drones]}}},
}

\newglossaryentry{Voltage-Sensor}{ sort={Voltage Sensor},
  name={Voltage Sensor},
  description={\textcolor{catD}{\textbf{[Drones]}}},
}

\newglossaryentry{Power-Module}{ sort={Power Module},
  name={Power Module},
  description={\textcolor{catD}{\textbf{[Drones]}}},
}

\newglossaryentry{Pixhawk}{ sort={Pixhawk},
  name={Pixhawk},
  description={\textcolor{catD}{\textbf{[Drones]}}},
}

\newglossaryentry{Cube}{ sort={Cube},
  name={Cube},
  description={\textcolor{catD}{\textbf{[Drones]}}},
}

\newglossaryentry{Kakute}{ sort={Kakute},
  name={Kakute},
  description={\textcolor{catD}{\textbf{[Drones]}}},
}

\newglossaryentry{F4-Processor}{ sort={F4 Processor},
  name={F4 Processor},
  description={\textcolor{catD}{\textbf{[Drones]}}},
}

\newglossaryentry{F7-Processor}{ sort={F7 Processor},
  name={F7 Processor},
  description={\textcolor{catD}{\textbf{[Drones]}}},
}

\newglossaryentry{H7-Processor}{ sort={H7 Processor},
  name={H7 Processor},
  description={\textcolor{catD}{\textbf{[Drones]}}},
}

\newglossaryentry{Betaflight}{ sort={Betaflight},
  name={Betaflight},
  description={\textcolor{catD}{\textbf{[Drones]}}},
}

\newglossaryentry{INAV}{ sort={INAV},
  name={INAV},
  description={\textcolor{catD}{\textbf{[Drones]}}},
}

\newglossaryentry{Cleanflight}{ sort={Cleanflight},
  name={Cleanflight},
  description={\textcolor{catD}{\textbf{[Drones]}}},
}

\newglossaryentry{Kiss}{ sort={Kiss},
  name={Kiss},
  description={\textcolor{catD}{\textbf{[Drones]}}},
}

\newglossaryentry{ButterFlight}{ sort={ButterFlight},
  name={ButterFlight},
  description={\textcolor{catD}{\textbf{[Drones]}}},
}

\newglossaryentry{EmuFlight}{ sort={EmuFlight},
  name={EmuFlight},
  description={\textcolor{catD}{\textbf{[Drones]}}},
}

\newglossaryentry{Raceflight}{ sort={Raceflight},
  name={Raceflight},
  description={\textcolor{catD}{\textbf{[Drones]}}},
}

\newglossaryentry{ExpressLRS}{ sort={ExpressLRS},
  name={ExpressLRS},
  description={\textcolor{catD}{\textbf{[Drones]}}},
}

\newglossaryentry{Crossfire}{ sort={Crossfire},
  name={Crossfire},
  description={\textcolor{catD}{\textbf{[Drones]}}},
}

\newglossaryentry{Tracer}{ sort={Tracer},
  name={Tracer},
  description={\textcolor{catD}{\textbf{[Drones]}}},
}

\newglossaryentry{Ghost}{ sort={Ghost},
  name={Ghost},
  description={\textcolor{catD}{\textbf{[Drones]}}},
}

\newglossaryentry{DSMX}{ sort={DSMX},
  name={DSMX},
  description={\textcolor{catD}{\textbf{[Drones]}}},
}

\newglossaryentry{SBUS}{ sort={SBUS},
  name={SBUS},
  description={\textcolor{catD}{\textbf{[Drones]}}},
}

\newglossaryentry{CRSF}{ sort={CRSF},
  name={CRSF},
  description={\textcolor{catD}{\textbf{[Drones]}}},
}

\newglossaryentry{PWM}{ sort={PWM},
  name={PWM},
  description={\textcolor{catD}{\textbf{[Drones]}}},
}

\newglossaryentry{DShot}{ sort={DShot},
  name={DShot},
  description={\textcolor{catD}{\textbf{[Drones]}}},
}

\newglossaryentry{Multishot}{ sort={Multishot},
  name={Multishot},
  description={\textcolor{catD}{\textbf{[Drones]}}},
}

\newglossaryentry{Oneshot}{ sort={Oneshot},
  name={Oneshot},
  description={\textcolor{catD}{\textbf{[Drones]}}},
}

\newglossaryentry{Bidirectional-DShot}{ sort={Bidirectional DShot},
  name={Bidirectional DShot},
  description={\textcolor{catD}{\textbf{[Drones]}}},
}

\newglossaryentry{RPM-Filtering}{ sort={RPM Filtering},
  name={RPM Filtering},
  description={\textcolor{catD}{\textbf{[Drones]}}},
}

\newglossaryentry{Dynamic-Filter}{ sort={Dynamic Filter},
  name={Dynamic Filter},
  description={\textcolor{catD}{\textbf{[Drones]}}},
}

\newglossaryentry{Notch-Filter}{ sort={Notch Filter},
  name={Notch Filter},
  description={\textcolor{catD}{\textbf{[Drones]}}},
}

\newglossaryentry{Low-Pass-Filter}{ sort={Low Pass Filter},
  name={Low Pass Filter},
  description={\textcolor{catD}{\textbf{[Drones]}}},
}

\newglossaryentry{Gyro-Sampling}{ sort={Gyro Sampling},
  name={Gyro Sampling},
  description={\textcolor{catD}{\textbf{[Drones]}}},
}

\newglossaryentry{PID-Loop-Frequency}{ sort={PID Loop Frequency},
  name={PID Loop Frequency},
  description={\textcolor{catD}{\textbf{[Drones]}}},
}

\newglossaryentry{Desync}{ sort={Desync},
  name={Desync},
  description={\textcolor{catD}{\textbf{[Drones]}}},
}

\newglossaryentry{Motor-Timing}{ sort={Motor Timing},
  name={Motor Timing},
  description={\textcolor{catD}{\textbf{[Drones]}}},
}

\newglossaryentry{Demag-Compensation}{ sort={Demag Compensation},
  name={Demag Compensation},
  description={\textcolor{catD}{\textbf{[Drones]}}},
}

\newglossaryentry{Throttle-Scale}{ sort={Throttle Scale},
  name={Throttle Scale},
  description={\textcolor{catD}{\textbf{[Drones]}}},
}

\newglossaryentry{Throttle-Limit}{ sort={Throttle Limit},
  name={Throttle Limit},
  description={\textcolor{catD}{\textbf{[Drones]}}},
}

\newglossaryentry{Motor-Output}{ sort={Motor Output},
  name={Motor Output},
  description={\textcolor{catD}{\textbf{[Drones]}}},
}

\newglossaryentry{Direction-Lock}{ sort={Direction Lock},
  name={Direction Lock},
  description={\textcolor{catD}{\textbf{[Drones]}}},
}

\newglossaryentry{Resource-Remapping}{ sort={Resource Remapping},
  name={Resource Remapping},
  description={\textcolor{catD}{\textbf{[Drones]}}},
}

\newglossaryentry{Serial-Port}{ sort={Serial Port},
  name={Serial Port},
  description={\textcolor{catD}{\textbf{[Drones]}}},
}

\newglossaryentry{I2C}{ sort={I2C},
  name={I2C},
  description={\textcolor{catD}{\textbf{[Drones]}}},
}

\newglossaryentry{SPI}{ sort={SPI},
  name={SPI},
  description={\textcolor{catD}{\textbf{[Drones]}}},
}

\newglossaryentry{UART}{ sort={UART},
  name={UART},
  description={\textcolor{catD}{\textbf{[Drones]}}},
}

\newglossaryentry{CAN-Bus}{ sort={CAN Bus},
  name={CAN Bus},
  description={\textcolor{catD}{\textbf{[Drones]}}},
}

\newglossaryentry{SD-Card-Logging}{ sort={SD Card Logging},
  name={SD Card Logging},
  description={\textcolor{catD}{\textbf{[Drones]}}},
}

\newglossaryentry{Blackbox}{ sort={Blackbox},
  name={Blackbox},
  description={\textcolor{catD}{\textbf{[Drones]}}},
}

\newglossaryentry{OSD}{ sort={OSD},
  name={OSD},
  description={\textcolor{catD}{\textbf{[Drones]}}},
}

\newglossaryentry{Betaflight-OSD}{ sort={Betaflight OSD},
  name={Betaflight OSD},
  description={\textcolor{catD}{\textbf{[Drones]}}},
}

\newglossaryentry{DJI-OSD}{ sort={DJI OSD},
  name={DJI OSD},
  description={\textcolor{catD}{\textbf{[Drones]}}},
}

\newglossaryentry{FPV-Camera}{ sort={FPV Camera},
  name={FPV Camera},
  description={\textcolor{catD}{\textbf{[Drones]}}},
}

\newglossaryentry{VTX}{ sort={VTX},
  name={VTX},
  description={\textcolor{catD}{\textbf{[Drones]}}},
}

\newglossaryentry{Video-Receiver}{ sort={Video Receiver},
  name={Video Receiver},
  description={\textcolor{catD}{\textbf{[Drones]}}},
}

\newglossaryentry{FPV-Goggles}{ sort={FPV Goggles},
  name={FPV Goggles},
  description={\textcolor{catD}{\textbf{[Drones]}}},
}

\newglossaryentry{FPV-Monitor}{ sort={FPV Monitor},
  name={FPV Monitor},
  description={\textcolor{catD}{\textbf{[Drones]}}},
}

\newglossaryentry{Diversity-Receiver}{ sort={Diversity Receiver},
  name={Diversity Receiver},
  description={\textcolor{catD}{\textbf{[Drones]}}},
}

\newglossaryentry{Antenna-Polarization}{ sort={Antenna Polarization},
  name={Antenna Polarization},
  description={\textcolor{catD}{\textbf{[Drones]}}},
}

\newglossaryentry{Linear-Polarization}{ sort={Linear Polarization},
  name={Linear Polarization},
  description={\textcolor{catD}{\textbf{[Drones]}}},
}

\newglossaryentry{Circular-Polarization}{ sort={Circular Polarization},
  name={Circular Polarization},
  description={\textcolor{catD}{\textbf{[Drones]}}},
}

\newglossaryentry{RHCP}{ sort={RHCP},
  name={RHCP},
  description={\textcolor{catD}{\textbf{[Drones]}}},
}

\newglossaryentry{LHCP}{ sort={LHCP},
  name={LHCP},
  description={\textcolor{catD}{\textbf{[Drones]}}},
}

\newglossaryentry{SMA-Connector}{ sort={SMA Connector},
  name={SMA Connector},
  description={\textcolor{catD}{\textbf{[Drones]}}},
}

\newglossaryentry{RP-SMA}{ sort={RP-SMA},
  name={RP-SMA},
  description={\textcolor{catD}{\textbf{[Drones]}}},
}

\newglossaryentry{MMCX-Connector}{ sort={MMCX Connector},
  name={MMCX Connector},
  description={\textcolor{catD}{\textbf{[Drones]}}},
}

\newglossaryentry{UFL-Connector}{ sort={UFL Connector},
  name={UFL Connector},
  description={\textcolor{catD}{\textbf{[Drones]}}},
}

\newglossaryentry{IPX-Connector}{ sort={IPX Connector},
  name={IPX Connector},
  description={\textcolor{catD}{\textbf{[Drones]}}},
}

\newglossaryentry{LQ}{ sort={LQ},
  name={LQ},
  description={\textcolor{catD}{\textbf{[Drones]}}},
}

\newglossaryentry{SNR}{ sort={SNR},
  name={SNR},
  description={\textcolor{catD}{\textbf{[Drones]}}},
}

\newglossaryentry{Latency}{ sort={Latency},
  name={Latency},
  description={\textcolor{catD}{\textbf{[Drones]}}},
}

\newglossaryentry{End-to-End-Latency}{ sort={End to End Latency},
  name={End to End Latency},
  description={\textcolor{catD}{\textbf{[Drones]}}},
}

\newglossaryentry{Frame-Rate}{ sort={Frame Rate},
  name={Frame Rate},
  description={\textcolor{catD}{\textbf{[Drones]}}},
}

\newglossaryentry{Resolution}{ sort={Resolution},
  name={Resolution},
  description={\textcolor{catD}{\textbf{[Drones]}}},
}

\newglossaryentry{Bitrate}{ sort={Bitrate},
  name={Bitrate},
  description={\textcolor{catD}{\textbf{[Drones]}}},
}

\newglossaryentry{Codec}{ sort={Codec},
  name={Codec},
  description={\textcolor{catD}{\textbf{[Drones]}}},
}

\newglossaryentry{H264}{ sort={H.264},
  name={H.264},
  description={\textcolor{catD}{\textbf{[Drones]}}},
}

\newglossaryentry{H265}{ sort={H.265},
  name={H.265},
  description={\textcolor{catD}{\textbf{[Drones]}}},
}

\newglossaryentry{Mavlink-Protocol}{ sort={Mavlink Protocol},
  name={Mavlink Protocol},
  description={\textcolor{catD}{\textbf{[Drones]}}},
}

\newglossaryentry{Mavlink-2}{ sort={Mavlink 2},
  name={Mavlink 2},
  description={\textcolor{catD}{\textbf{[Drones]}}},
}

\newglossaryentry{Mavlink-Inspector}{ sort={Mavlink Inspector},
  name={Mavlink Inspector},
  description={\textcolor{catD}{\textbf{[Drones]}}},
}

\newglossaryentry{Parameter}{ sort={Parameter},
  name={Parameter},
  description={\textcolor{catD}{\textbf{[Drones]}}},
}

\newglossaryentry{Mission}{ sort={Mission},
  name={Mission},
  description={\textcolor{catD}{\textbf{[Drones]}}},
}

\newglossaryentry{Flight-Mode}{ sort={Flight Mode},
  name={Flight Mode},
  description={\textcolor{catD}{\textbf{[Drones]}}},
}

\newglossaryentry{Arming}{ sort={Arming},
  name={Arming},
  description={\textcolor{catD}{\textbf{[Drones]}}},
}

\newglossaryentry{Disarming}{ sort={Disarming},
  name={Disarming},
  description={\textcolor{catD}{\textbf{[Drones]}}},
}

\newglossaryentry{Pre-Arm-Check}{ sort={Pre-Arm Check},
  name={Pre-Arm Check},
  description={\textcolor{catD}{\textbf{[Drones]}}},
}

\newglossaryentry{Calibration}{ sort={Calibration},
  name={Calibration},
  description={\textcolor{catD}{\textbf{[Drones]}}},
}

\newglossaryentry{Level-Calibration}{ sort={Level Calibration},
  name={Level Calibration},
  description={\textcolor{catD}{\textbf{[Drones]}}},
}

\newglossaryentry{Compass-Calibration}{ sort={Compass Calibration},
  name={Compass Calibration},
  description={\textcolor{catD}{\textbf{[Drones]}}},
}

\newglossaryentry{ESC-Calibration}{ sort={ESC Calibration},
  name={ESC Calibration},
  description={\textcolor{catD}{\textbf{[Drones]}}},
}

\newglossaryentry{Radio-Calibration}{ sort={Radio Calibration},
  name={Radio Calibration},
  description={\textcolor{catD}{\textbf{[Drones]}}},
}

\newglossaryentry{Throttle-Calibration}{ sort={Throttle Calibration},
  name={Throttle Calibration},
  description={\textcolor{catD}{\textbf{[Drones]}}},
}

\newglossaryentry{Safety-Switch}{ sort={Safety Switch},
  name={Safety Switch},
  description={\textcolor{catD}{\textbf{[Drones]}}},
}

\newglossaryentry{Beacon}{ sort={Beacon},
  name={Beacon},
  description={\textcolor{catD}{\textbf{[Drones]}}},
}

\newglossaryentry{Return-to-Launch}{ sort={Return to Launch},
  name={Return to Launch},
  description={\textcolor{catD}{\textbf{[Drones]}}},
}

\newglossaryentry{Land-Immediately}{ sort={Land Immediately},
  name={Land Immediately},
  description={\textcolor{catD}{\textbf{[Drones]}}},
}

\newglossaryentry{Geofence-Action}{ sort={Geofence Action},
  name={Geofence Action},
  description={\textcolor{catD}{\textbf{[Drones]}}},
}

\newglossaryentry{Battery-Failsafe}{ sort={Battery Failsafe},
  name={Battery Failsafe},
  description={\textcolor{catD}{\textbf{[Drones]}}},
}

\newglossaryentry{RC-Failsafe}{ sort={RC Failsafe},
  name={RC Failsafe},
  description={\textcolor{catD}{\textbf{[Drones]}}},
}

\newglossaryentry{GPS-Failsafe}{ sort={GPS Failsafe},
  name={GPS Failsafe},
  description={\textcolor{catD}{\textbf{[Drones]}}},
}

\newglossaryentry{Compass-Error}{ sort={Compass Error},
  name={Compass Error},
  description={\textcolor{catD}{\textbf{[Drones]}}},
}

\newglossaryentry{Gyro-Calibration}{ sort={Gyro Calibration},
  name={Gyro Calibration},
  description={\textcolor{catD}{\textbf{[Drones]}}},
}

\newglossaryentry{Accel-Calibration}{ sort={Accel Calibration},
  name={Accel Calibration},
  description={\textcolor{catD}{\textbf{[Drones]}}},
}

\newglossaryentry{Space-Science}{ sort={Space Science},
  name={Space Science},
  description={\textcolor{catSS}{\textbf{[Space Science]}}},
}

\newglossaryentry{Astronomy}{ sort={Astronomy},
  name={Astronomy},
  description={\textcolor{catSS}{\textbf{[Space Science]}}},
}

\newglossaryentry{Astrophysics}{ sort={Astrophysics},
  name={Astrophysics},
  description={\textcolor{catSS}{\textbf{[Space Science]}}},
}

\newglossaryentry{Cosmology}{ sort={Cosmology},
  name={Cosmology},
  description={\textcolor{catSS}{\textbf{[Space Science]}}},
}

\newglossaryentry{Planetary-Science}{ sort={Planetary Science},
  name={Planetary Science},
  description={\textcolor{catSS}{\textbf{[Space Science]}}},
}

\newglossaryentry{Exoplanet}{ sort={Exoplanet},
  name={Exoplanet},
  description={\textcolor{catSS}{\textbf{[Space Science]}}},
}

\newglossaryentry{Solar-System}{ sort={Solar System},
  name={Solar System},
  description={\textcolor{catSS}{\textbf{[Space Science]}}},
}

\newglossaryentry{Star}{ sort={Star},
  name={Star},
  description={\textcolor{catSS}{\textbf{[Space Science]}}},
}

\newglossaryentry{Planet}{ sort={Planet},
  name={Planet},
  description={\textcolor{catSS}{\textbf{[Space Science]}}},
}

\newglossaryentry{Moon}{ sort={Moon},
  name={Moon},
  description={\textcolor{catSS}{\textbf{[Space Science]}}},
}

\newglossaryentry{Asteroid}{ sort={Asteroid},
  name={Asteroid},
  description={\textcolor{catSS}{\textbf{[Space Science]}}},
}

\newglossaryentry{Comet}{ sort={Comet},
  name={Comet},
  description={\textcolor{catSS}{\textbf{[Space Science]}}},
}

\newglossaryentry{Meteor}{ sort={Meteor},
  name={Meteor},
  description={\textcolor{catSS}{\textbf{[Space Science]}}},
}

\newglossaryentry{Meteorite}{ sort={Meteorite},
  name={Meteorite},
  description={\textcolor{catSS}{\textbf{[Space Science]}}},
}

\newglossaryentry{Meteoroid}{ sort={Meteoroid},
  name={Meteoroid},
  description={\textcolor{catSS}{\textbf{[Space Science]}}},
}

\newglossaryentry{Asteroid-Belt}{ sort={Asteroid Belt},
  name={Asteroid Belt},
  description={\textcolor{catSS}{\textbf{[Space Science]}}},
}

\newglossaryentry{Kuiper-Belt}{ sort={Kuiper Belt},
  name={Kuiper Belt},
  description={\textcolor{catSS}{\textbf{[Space Science]}}},
}

\newglossaryentry{Oort-Cloud}{ sort={Oort Cloud},
  name={Oort Cloud},
  description={\textcolor{catSS}{\textbf{[Space Science]}}},
}

\newglossaryentry{Dwarf-Planet}{ sort={Dwarf Planet},
  name={Dwarf Planet},
  description={\textcolor{catSS}{\textbf{[Space Science]}}},
}

\newglossaryentry{Pluto}{ sort={Pluto},
  name={Pluto},
  description={\textcolor{catSS}{\textbf{[Space Science]}}},
}

\newglossaryentry{Ceres}{ sort={Ceres},
  name={Ceres},
  description={\textcolor{catSS}{\textbf{[Space Science]}}},
}

\newglossaryentry{Sun}{ sort={Sun},
  name={Sun},
  description={\textcolor{catSS}{\textbf{[Space Science]}}},
}

\newglossaryentry{Solar-Flare}{ sort={Solar Flare},
  name={Solar Flare},
  description={\textcolor{catSS}{\textbf{[Space Science]}}},
}

\newglossaryentry{Coronal-Mass-Ejection}{ sort={Coronal Mass Ejection},
  name={Coronal Mass Ejection},
  description={\textcolor{catSS}{\textbf{[Space Science]}}},
}

\newglossaryentry{Solar-Wind}{ sort={Solar Wind},
  name={Solar Wind},
  description={\textcolor{catSS}{\textbf{[Space Science]}}},
}

\newglossaryentry{Heliosphere}{ sort={Heliosphere},
  name={Heliosphere},
  description={\textcolor{catSS}{\textbf{[Space Science]}}},
}

\newglossaryentry{Magnetosphere}{ sort={Magnetosphere},
  name={Magnetosphere},
  description={\textcolor{catSS}{\textbf{[Space Science]}}},
}

\newglossaryentry{Van-Allen-Belt}{ sort={Van Allen Belt},
  name={Van Allen Belt},
  description={\textcolor{catSS}{\textbf{[Space Science]}}},
}

\newglossaryentry{Aurora}{ sort={Aurora},
  name={Aurora},
  description={\textcolor{catSS}{\textbf{[Space Science]}}},
}

\newglossaryentry{Northern-Lights}{ sort={Northern Lights},
  name={Northern Lights},
  description={\textcolor{catSS}{\textbf{[Space Science]}}},
}

\newglossaryentry{Southern-Lights}{ sort={Southern Lights},
  name={Southern Lights},
  description={\textcolor{catSS}{\textbf{[Space Science]}}},
}

\newglossaryentry{Ionosphere}{ sort={Ionosphere},
  name={Ionosphere},
  description={\textcolor{catSS}{\textbf{[Space Science]}}},
}

\newglossaryentry{Atmosphere}{ sort={Atmosphere},
  name={Atmosphere},
  description={\textcolor{catSS}{\textbf{[Space Science]}}},
}

\newglossaryentry{Exosphere}{ sort={Exosphere},
  name={Exosphere},
  description={\textcolor{catSS}{\textbf{[Space Science]}}},
}

\newglossaryentry{Thermosphere}{ sort={Thermosphere},
  name={Thermosphere},
  description={\textcolor{catSS}{\textbf{[Space Science]}}},
}

\newglossaryentry{Mesosphere}{ sort={Mesosphere},
  name={Mesosphere},
  description={\textcolor{catSS}{\textbf{[Space Science]}}},
}

\newglossaryentry{Stratosphere}{ sort={Stratosphere},
  name={Stratosphere},
  description={\textcolor{catSS}{\textbf{[Space Science]}}},
}

\newglossaryentry{Troposphere}{ sort={Troposphere},
  name={Troposphere},
  description={\textcolor{catSS}{\textbf{[Space Science]}}},
}

\newglossaryentry{Ozone-Layer}{ sort={Ozone Layer},
  name={Ozone Layer},
  description={\textcolor{catSS}{\textbf{[Space Science]}}},
}

\newglossaryentry{Greenhouse-Effect}{ sort={Greenhouse Effect},
  name={Greenhouse Effect},
  description={\textcolor{catSS}{\textbf{[Space Science]}}},
}

\newglossaryentry{Carbon-Dioxide}{ sort={Carbon Dioxide},
  name={Carbon Dioxide},
  description={\textcolor{catSS}{\textbf{[Space Science]}}},
}

\newglossaryentry{Orbital-Mechanics}{ sort={Orbital Mechanics},
  name={Orbital Mechanics},
  description={\textcolor{catSS}{\textbf{[Space Science]}}},
}

\newglossaryentry{Celestial-Mechanics}{ sort={Celestial Mechanics},
  name={Celestial Mechanics},
  description={\textcolor{catSS}{\textbf{[Space Science]}}},
}

\newglossaryentry{Keplers-Laws}{ sort={Kepler's Laws},
  name={Kepler's Laws},
  description={\textcolor{catSS}{\textbf{[Space Science]}}},
}

\newglossaryentry{Keplers-First-Law}{ sort={Kepler's First Law},
  name={Kepler's First Law},
  description={\textcolor{catSS}{\textbf{[Space Science]}}},
}

\newglossaryentry{Keplers-Second-Law}{ sort={Kepler's Second Law},
  name={Kepler's Second Law},
  description={\textcolor{catSS}{\textbf{[Space Science]}}},
}

\newglossaryentry{Keplers-Third-Law}{ sort={Kepler's Third Law},
  name={Kepler's Third Law},
  description={\textcolor{catSS}{\textbf{[Space Science]}}},
}

\newglossaryentry{Newtons-Law-of-Gravitation}{ sort={Newton's Law of Gravitation},
  name={Newton's Law of Gravitation},
  description={\textcolor{catSS}{\textbf{[Space Science]}}},
}

\newglossaryentry{Gravitational-Constant}{ sort={Gravitational Constant},
  name={Gravitational Constant},
  description={\textcolor{catSS}{\textbf{[Space Science]}}},
}

\newglossaryentry{Gravity}{ sort={Gravity},
  name={Gravity},
  description={\textcolor{catSS}{\textbf{[Space Science]}}},
}

\newglossaryentry{Microgravity}{ sort={Microgravity},
  name={Microgravity},
  description={\textcolor{catSS}{\textbf{[Space Science]}}},
}

\newglossaryentry{Weightlessness}{ sort={Weightlessness},
  name={Weightlessness},
  description={\textcolor{catSS}{\textbf{[Space Science]}}},
}

\newglossaryentry{Free-Fall}{ sort={Free Fall},
  name={Free Fall},
  description={\textcolor{catSS}{\textbf{[Space Science]}}},
}

\newglossaryentry{Escape-Velocity}{ sort={Escape Velocity},
  name={Escape Velocity},
  description={\textcolor{catSS}{\textbf{[Space Science]}}},
}

\newglossaryentry{Orbital-Velocity}{ sort={Orbital Velocity},
  name={Orbital Velocity},
  description={\textcolor{catSS}{\textbf{[Space Science]}}},
}

\newglossaryentry{Circular-Velocity}{ sort={Circular Velocity},
  name={Circular Velocity},
  description={\textcolor{catSS}{\textbf{[Space Science]}}},
}

\newglossaryentry{Hohmann-Transfer-Orbit}{ sort={Hohmann Transfer Orbit},
  name={Hohmann Transfer Orbit},
  description={\textcolor{catSS}{\textbf{[Space Science]}}},
}

\newglossaryentry{Bi-Elliptic-Transfer}{ sort={Bi-Elliptic Transfer},
  name={Bi-Elliptic Transfer},
  description={\textcolor{catSS}{\textbf{[Space Science]}}},
}

\newglossaryentry{Gravity-Assist}{ sort={Gravity Assist},
  name={Gravity Assist},
  description={\textcolor{catSS}{\textbf{[Space Science]}}},
}

\newglossaryentry{Swing-By}{ sort={Swing-By},
  name={Swing-By},
  description={\textcolor{catSS}{\textbf{[Space Science]}}},
}

\newglossaryentry{Slingshot}{ sort={Slingshot},
  name={Slingshot},
  description={\textcolor{catSS}{\textbf{[Space Science]}}},
}

\newglossaryentry{Lagrange-Point}{ sort={Lagrange Point},
  name={Lagrange Point},
  description={\textcolor{catSS}{\textbf{[Space Science]}}},
}

\newglossaryentry{L1}{ sort={L1},
  name={L1},
  description={\textcolor{catSS}{\textbf{[Space Science]}}},
}

\newglossaryentry{L2}{ sort={L2},
  name={L2},
  description={\textcolor{catSS}{\textbf{[Space Science]}}},
}

\newglossaryentry{Lissajous-Orbit}{ sort={Lissajous Orbit},
  name={Lissajous Orbit},
  description={\textcolor{catSS}{\textbf{[Space Science]}}},
}

\newglossaryentry{Halo-Orbit}{ sort={Halo Orbit},
  name={Halo Orbit},
  description={\textcolor{catSS}{\textbf{[Space Science]}}},
}

\newglossaryentry{Patched-Conics}{ sort={Patched Conics},
  name={Patched Conics},
  description={\textcolor{catSS}{\textbf{[Space Science]}}},
}

\newglossaryentry{Sphere-of-Influence}{ sort={Sphere of Influence},
  name={Sphere of Influence},
  description={\textcolor{catSS}{\textbf{[Space Science]}}},
}

\newglossaryentry{Launch-Window}{ sort={Launch Window},
  name={Launch Window},
  description={\textcolor{catSS}{\textbf{[Space Science]}}},
}

\newglossaryentry{Interplanetary-Transfer}{ sort={Interplanetary Transfer},
  name={Interplanetary Transfer},
  description={\textcolor{catSS}{\textbf{[Space Science]}}},
}

\newglossaryentry{Earth-Orbit}{ sort={Earth Orbit},
  name={Earth Orbit},
  description={\textcolor{catSS}{\textbf{[Space Science]}}},
}

\newglossaryentry{Low-Earth-Orbit}{ sort={Low Earth Orbit},
  name={Low Earth Orbit},
  description={\textcolor{catSS}{\textbf{[Space Science]}}},
}

\newglossaryentry{Medium-Earth-Orbit}{ sort={Medium Earth Orbit},
  name={Medium Earth Orbit},
  description={\textcolor{catSS}{\textbf{[Space Science]}}},
}

\newglossaryentry{Geosynchronous-Orbit}{ sort={Geosynchronous Orbit},
  name={Geosynchronous Orbit},
  description={\textcolor{catSS}{\textbf{[Space Science]}}},
}

\newglossaryentry{Geostationary-Orbit}{ sort={Geostationary Orbit},
  name={Geostationary Orbit},
  description={\textcolor{catSS}{\textbf{[Space Science]}}},
}

\newglossaryentry{Geostationary-Transfer-Orbit}{ sort={Geostationary Transfer Orbit},
  name={Geostationary Transfer Orbit},
  description={\textcolor{catSS}{\textbf{[Space Science]}}},
}

\newglossaryentry{Highly-Elliptical-Orbit}{ sort={Highly Elliptical Orbit},
  name={Highly Elliptical Orbit},
  description={\textcolor{catSS}{\textbf{[Space Science]}}},
}

\newglossaryentry{Molniya-Orbit}{ sort={Molniya Orbit},
  name={Molniya Orbit},
  description={\textcolor{catSS}{\textbf{[Space Science]}}},
}

\newglossaryentry{Tundra-Orbit}{ sort={Tundra Orbit},
  name={Tundra Orbit},
  description={\textcolor{catSS}{\textbf{[Space Science]}}},
}

\newglossaryentry{Polar-Orbit}{ sort={Polar Orbit},
  name={Polar Orbit},
  description={\textcolor{catSS}{\textbf{[Space Science]}}},
}

\newglossaryentry{Sun-Synchronous-Orbit}{ sort={Sun-Synchronous Orbit},
  name={Sun-Synchronous Orbit},
  description={\textcolor{catSS}{\textbf{[Space Science]}}},
}

\newglossaryentry{Dawn-Dusk-Orbit}{ sort={Dawn-Dusk Orbit},
  name={Dawn-Dusk Orbit},
  description={\textcolor{catSS}{\textbf{[Space Science]}}},
}

\newglossaryentry{Frozen-Orbit}{ sort={Frozen Orbit},
  name={Frozen Orbit},
  description={\textcolor{catSS}{\textbf{[Space Science]}}},
}

\newglossaryentry{Graveyard-Orbit}{ sort={Graveyard Orbit},
  name={Graveyard Orbit},
  description={\textcolor{catSS}{\textbf{[Space Science]}}},
}

\newglossaryentry{Parking-Orbit}{ sort={Parking Orbit},
  name={Parking Orbit},
  description={\textcolor{catSS}{\textbf{[Space Science]}}},
}

\newglossaryentry{Phasing-Orbit}{ sort={Phasing Orbit},
  name={Phasing Orbit},
  description={\textcolor{catSS}{\textbf{[Space Science]}}},
}

\newglossaryentry{Co-orbital}{ sort={Co-orbital},
  name={Co-orbital},
  description={\textcolor{catSS}{\textbf{[Space Science]}}},
}

\newglossaryentry{Retrograde-Orbit}{ sort={Retrograde Orbit},
  name={Retrograde Orbit},
  description={\textcolor{catSS}{\textbf{[Space Science]}}},
}

\newglossaryentry{Prograde-Orbit}{ sort={Prograde Orbit},
  name={Prograde Orbit},
  description={\textcolor{catSS}{\textbf{[Space Science]}}},
}

\newglossaryentry{Ascending-Node}{ sort={Ascending Node},
  name={Ascending Node},
  description={\textcolor{catSS}{\textbf{[Space Science]}}},
}

\newglossaryentry{Descending-Node}{ sort={Descending Node},
  name={Descending Node},
  description={\textcolor{catSS}{\textbf{[Space Science]}}},
}

\newglossaryentry{Line-of-Nodes}{ sort={Line of Nodes},
  name={Line of Nodes},
  description={\textcolor{catSS}{\textbf{[Space Science]}}},
}

\newglossaryentry{Argument-of-Perigee}{ sort={Argument of Perigee},
  name={Argument of Perigee},
  description={\textcolor{catSS}{\textbf{[Space Science]}}},
}

\newglossaryentry{Right-Ascension-of-Ascending-Node}{ sort={Right Ascension of Ascending Node},
  name={Right Ascension of Ascending Node},
  description={\textcolor{catSS}{\textbf{[Space Science]}}},
}

\newglossaryentry{True-Anomaly}{ sort={True Anomaly},
  name={True Anomaly},
  description={\textcolor{catSS}{\textbf{[Space Science]}}},
}

\newglossaryentry{Mean-Anomaly}{ sort={Mean Anomaly},
  name={Mean Anomaly},
  description={\textcolor{catSS}{\textbf{[Space Science]}}},
}

\newglossaryentry{Eccentric-Anomaly}{ sort={Eccentric Anomaly},
  name={Eccentric Anomaly},
  description={\textcolor{catSS}{\textbf{[Space Science]}}},
}

\newglossaryentry{Semi-Major-Axis}{ sort={Semi-Major Axis},
  name={Semi-Major Axis},
  description={\textcolor{catSS}{\textbf{[Space Science]}}},
}

\newglossaryentry{Semi-Minor-Axis}{ sort={Semi-Minor Axis},
  name={Semi-Minor Axis},
  description={\textcolor{catSS}{\textbf{[Space Science]}}},
}

\newglossaryentry{Eccentricity}{ sort={Eccentricity},
  name={Eccentricity},
  description={\textcolor{catSS}{\textbf{[Space Science]}}},
}

\newglossaryentry{Inclination}{ sort={Inclination},
  name={Inclination},
  description={\textcolor{catSS}{\textbf{[Space Science]}}},
}

\newglossaryentry{Periapsis}{ sort={Periapsis},
  name={Periapsis},
  description={\textcolor{catSS}{\textbf{[Space Science]}}},
}

\newglossaryentry{Apoapsis}{ sort={Apoapsis},
  name={Apoapsis},
  description={\textcolor{catSS}{\textbf{[Space Science]}}},
}

\newglossaryentry{Perigee}{ sort={Perigee},
  name={Perigee},
  description={\textcolor{catSS}{\textbf{[Space Science]}}},
}

\newglossaryentry{Apogee}{ sort={Apogee},
  name={Apogee},
  description={\textcolor{catSS}{\textbf{[Space Science]}}},
}

\newglossaryentry{Perihelion}{ sort={Perihelion},
  name={Perihelion},
  description={\textcolor{catSS}{\textbf{[Space Science]}}},
}

\newglossaryentry{Aphelion}{ sort={Aphelion},
  name={Aphelion},
  description={\textcolor{catSS}{\textbf{[Space Science]}}},
}

\newglossaryentry{Orbital-Period}{ sort={Orbital Period},
  name={Orbital Period},
  description={\textcolor{catSS}{\textbf{[Space Science]}}},
}

\newglossaryentry{Synodic-Period}{ sort={Synodic Period},
  name={Synodic Period},
  description={\textcolor{catSS}{\textbf{[Space Science]}}},
}

\newglossaryentry{Sidereal-Period}{ sort={Sidereal Period},
  name={Sidereal Period},
  description={\textcolor{catSS}{\textbf{[Space Science]}}},
}

\newglossaryentry{Nodal-Precession}{ sort={Nodal Precession},
  name={Nodal Precession},
  description={\textcolor{catSS}{\textbf{[Space Science]}}},
}

\newglossaryentry{Apsidal-Precession}{ sort={Apsidal Precession},
  name={Apsidal Precession},
  description={\textcolor{catSS}{\textbf{[Space Science]}}},
}

\newglossaryentry{Perturbation}{ sort={Perturbation},
  name={Perturbation},
  description={\textcolor{catSS}{\textbf{[Space Science]}}},
}

\newglossaryentry{J2-Perturbation}{ sort={J2 Perturbation},
  name={J2 Perturbation},
  description={\textcolor{catSS}{\textbf{[Space Science]}}},
}

\newglossaryentry{Third-Body-Perturbation}{ sort={Third Body Perturbation},
  name={Third Body Perturbation},
  description={\textcolor{catSS}{\textbf{[Space Science]}}},
}

\newglossaryentry{Drag-Perturbation}{ sort={Drag Perturbation},
  name={Drag Perturbation},
  description={\textcolor{catSS}{\textbf{[Space Science]}}},
}

\newglossaryentry{Radiation-Pressure}{ sort={Radiation Pressure},
  name={Radiation Pressure},
  description={\textcolor{catSS}{\textbf{[Space Science]}}},
}

\newglossaryentry{Poynting-Robertson-Effect}{ sort={Poynting-Robertson Effect},
  name={Poynting-Robertson Effect},
  description={\textcolor{catSS}{\textbf{[Space Science]}}},
}

\newglossaryentry{Yarkovsky-Effect}{ sort={Yarkovsky Effect},
  name={Yarkovsky Effect},
  description={\textcolor{catSS}{\textbf{[Space Science]}}},
}

\newglossaryentry{Roche-Limit}{ sort={Roche Limit},
  name={Roche Limit},
  description={\textcolor{catSS}{\textbf{[Space Science]}}},
}

\newglossaryentry{Hill-Sphere}{ sort={Hill Sphere},
  name={Hill Sphere},
  description={\textcolor{catSS}{\textbf{[Space Science]}}},
}

\newglossaryentry{Tidal-Force}{ sort={Tidal Force},
  name={Tidal Force},
  description={\textcolor{catSS}{\textbf{[Space Science]}}},
}

\newglossaryentry{Tidal-Locking}{ sort={Tidal Locking},
  name={Tidal Locking},
  description={\textcolor{catSS}{\textbf{[Space Science]}}},
}

\newglossaryentry{Synchronous-Rotation}{ sort={Synchronous Rotation},
  name={Synchronous Rotation},
  description={\textcolor{catSS}{\textbf{[Space Science]}}},
}

\newglossaryentry{Spacecraft}{ sort={Spacecraft},
  name={Spacecraft},
  description={\textcolor{catSS}{\textbf{[Space Science]}}},
}

\newglossaryentry{Satellite}{ sort={Satellite},
  name={Satellite},
  description={\textcolor{catSS}{\textbf{[Space Science]}}},
}

\newglossaryentry{Space-Station}{ sort={Space Station},
  name={Space Station},
  description={\textcolor{catSS}{\textbf{[Space Science]}}},
}

\newglossaryentry{ISS}{ sort={ISS},
  name={ISS},
  description={\textcolor{catSS}{\textbf{[Space Science]}}},
}

\newglossaryentry{Capsule}{ sort={Capsule},
  name={Capsule},
  description={\textcolor{catSS}{\textbf{[Space Science]}}},
}

\newglossaryentry{Space-Probe}{ sort={Space Probe},
  name={Space Probe},
  description={\textcolor{catSS}{\textbf{[Space Science]}}},
}

\newglossaryentry{Lander}{ sort={Lander},
  name={Lander},
  description={\textcolor{catSS}{\textbf{[Space Science]}}},
}

\newglossaryentry{Rover}{ sort={Rover},
  name={Rover},
  description={\textcolor{catSS}{\textbf{[Space Science]}}},
}

\newglossaryentry{Orbiter}{ sort={Orbiter},
  name={Orbiter},
  description={\textcolor{catSS}{\textbf{[Space Science]}}},
}

\newglossaryentry{Flyby-Spacecraft}{ sort={Flyby Spacecraft},
  name={Flyby Spacecraft},
  description={\textcolor{catSS}{\textbf{[Space Science]}}},
}

\newglossaryentry{Sample-Return-Mission}{ sort={Sample Return Mission},
  name={Sample Return Mission},
  description={\textcolor{catSS}{\textbf{[Space Science]}}},
}

\newglossaryentry{Crewed-Spacecraft}{ sort={Crewed Spacecraft},
  name={Crewed Spacecraft},
  description={\textcolor{catSS}{\textbf{[Space Science]}}},
}

\newglossaryentry{Space-Suit}{ sort={Space Suit},
  name={Space Suit},
  description={\textcolor{catSS}{\textbf{[Space Science]}}},
}

\newglossaryentry{EVA-Suit}{ sort={EVA Suit},
  name={EVA Suit},
  description={\textcolor{catSS}{\textbf{[Space Science]}}},
}

\newglossaryentry{Life-Support-System}{ sort={Life Support System},
  name={Life Support System},
  description={\textcolor{catSS}{\textbf{[Space Science]}}},
}

\newglossaryentry{ECLSS}{ sort={ECLSS},
  name={ECLSS},
  description={\textcolor{catSS}{\textbf{[Space Science]}}},
}

\newglossaryentry{Habitat-Module}{ sort={Habitat Module},
  name={Habitat Module},
  description={\textcolor{catSS}{\textbf{[Space Science]}}},
}

\newglossaryentry{Airlock}{ sort={Airlock},
  name={Airlock},
  description={\textcolor{catSS}{\textbf{[Space Science]}}},
}

\newglossaryentry{Docking-Mechanism}{ sort={Docking Mechanism},
  name={Docking Mechanism},
  description={\textcolor{catSS}{\textbf{[Space Science]}}},
}

\newglossaryentry{Androgynous-Docking}{ sort={Androgynous Docking},
  name={Androgynous Docking},
  description={\textcolor{catSS}{\textbf{[Space Science]}}},
}

\newglossaryentry{Probe-and-Drogue}{ sort={Probe and Drogue},
  name={Probe and Drogue},
  description={\textcolor{catSS}{\textbf{[Space Science]}}},
}

\newglossaryentry{Capture}{ sort={Capture},
  name={Capture},
  description={\textcolor{catSS}{\textbf{[Space Science]}}},
}

\newglossaryentry{Soft-Capture}{ sort={Soft Capture},
  name={Soft Capture},
  description={\textcolor{catSS}{\textbf{[Space Science]}}},
}

\newglossaryentry{Hard-Capture}{ sort={Hard Capture},
  name={Hard Capture},
  description={\textcolor{catSS}{\textbf{[Space Science]}}},
}

\newglossaryentry{Berthing}{ sort={Berthing},
  name={Berthing},
  description={\textcolor{catSS}{\textbf{[Space Science]}}},
}

\newglossaryentry{Robotic-Arm}{ sort={Robotic Arm},
  name={Robotic Arm},
  description={\textcolor{catSS}{\textbf{[Space Science]}}},
}

\newglossaryentry{Canadarm}{ sort={Canadarm},
  name={Canadarm},
  description={\textcolor{catSS}{\textbf{[Space Science]}}},
}

\newglossaryentry{Remote-Manipulator-System}{ sort={Remote Manipulator System},
  name={Remote Manipulator System},
  description={\textcolor{catSS}{\textbf{[Space Science]}}},
}

\newglossaryentry{Space-Tether}{ sort={Space Tether},
  name={Space Tether},
  description={\textcolor{catSS}{\textbf{[Space Science]}}},
}

\newglossaryentry{Electrodynamic-Tether}{ sort={Electrodynamic Tether},
  name={Electrodynamic Tether},
  description={\textcolor{catSS}{\textbf{[Space Science]}}},
}

\newglossaryentry{Solar-Sail}{ sort={Solar Sail},
  name={Solar Sail},
  description={\textcolor{catSS}{\textbf{[Space Science]}}},
}

\newglossaryentry{Solar-Panel}{ sort={Solar Panel},
  name={Solar Panel},
  description={\textcolor{catSS}{\textbf{[Space Science]}}},
}

\newglossaryentry{Photovoltaic-Cell}{ sort={Photovoltaic Cell},
  name={Photovoltaic Cell},
  description={\textcolor{catSS}{\textbf{[Space Science]}}},
}

\newglossaryentry{Solar-Array}{ sort={Solar Array},
  name={Solar Array},
  description={\textcolor{catSS}{\textbf{[Space Science]}}},
}

\newglossaryentry{RTG}{ sort={RTG},
  name={RTG},
  description={\textcolor{catSS}{\textbf{[Space Science]}}},
}

\newglossaryentry{Radioisotope-Thermoelectric-Generator}{ sort={Radioisotope Thermoelectric Generator},
  name={Radioisotope Thermoelectric Generator},
  description={\textcolor{catSS}{\textbf{[Space Science]}}},
}

\newglossaryentry{Nuclear-Power-in-Space}{ sort={Nuclear Power in Space},
  name={Nuclear Power in Space},
  description={\textcolor{catSS}{\textbf{[Space Science]}}},
}

\newglossaryentry{Battery}{ sort={Battery},
  name={Battery},
  description={\textcolor{catSS}{\textbf{[Space Science]}}},
}

\newglossaryentry{Fuel-Cell}{ sort={Fuel Cell},
  name={Fuel Cell},
  description={\textcolor{catSS}{\textbf{[Space Science]}}},
}

\newglossaryentry{Attitude-Control-System}{ sort={Attitude Control System},
  name={Attitude Control System},
  description={\textcolor{catSS}{\textbf{[Space Science]}}},
}

\newglossaryentry{Reaction-Wheel}{ sort={Reaction Wheel},
  name={Reaction Wheel},
  description={\textcolor{catSS}{\textbf{[Space Science]}}},
}

\newglossaryentry{Control-Moment-Gyro}{ sort={Control Moment Gyro},
  name={Control Moment Gyro},
  description={\textcolor{catSS}{\textbf{[Space Science]}}},
}

\newglossaryentry{Thruster}{ sort={Thruster},
  name={Thruster},
  description={\textcolor{catSS}{\textbf{[Space Science]}}},
}

\newglossaryentry{Reaction-Control-System}{ sort={Reaction Control System},
  name={Reaction Control System},
  description={\textcolor{catSS}{\textbf{[Space Science]}}},
}

\newglossaryentry{Star-Tracker}{ sort={Star Tracker},
  name={Star Tracker},
  description={\textcolor{catSS}{\textbf{[Space Science]}}},
}

\newglossaryentry{Sun-Sensor}{ sort={Sun Sensor},
  name={Sun Sensor},
  description={\textcolor{catSS}{\textbf{[Space Science]}}},
}

\newglossaryentry{Horizon-Sensor}{ sort={Horizon Sensor},
  name={Horizon Sensor},
  description={\textcolor{catSS}{\textbf{[Space Science]}}},
}

\newglossaryentry{Earth-Sensor}{ sort={Earth Sensor},
  name={Earth Sensor},
  description={\textcolor{catSS}{\textbf{[Space Science]}}},
}

\newglossaryentry{Magnetometer}{ sort={Magnetometer},
  name={Magnetometer},
  description={\textcolor{catSS}{\textbf{[Space Science]}}},
}

\newglossaryentry{Gyroscope}{ sort={Gyroscope},
  name={Gyroscope},
  description={\textcolor{catSS}{\textbf{[Space Science]}}},
}

\newglossaryentry{Inertial-Measurement-Unit}{ sort={Inertial Measurement Unit},
  name={Inertial Measurement Unit},
  description={\textcolor{catSS}{\textbf{[Space Science]}}},
}

\newglossaryentry{Attitude-Determination}{ sort={Attitude Determination},
  name={Attitude Determination},
  description={\textcolor{catSS}{\textbf{[Space Science]}}},
}

\newglossaryentry{Attitude-Propagation}{ sort={Attitude Propagation},
  name={Attitude Propagation},
  description={\textcolor{catSS}{\textbf{[Space Science]}}},
}

\newglossaryentry{Despin}{ sort={Despin},
  name={Despin},
  description={\textcolor{catSS}{\textbf{[Space Science]}}},
}

\newglossaryentry{Spin-Stabilization}{ sort={Spin Stabilization},
  name={Spin Stabilization},
  description={\textcolor{catSS}{\textbf{[Space Science]}}},
}

\newglossaryentry{Three-Axis-Stabilization}{ sort={Three-Axis Stabilization},
  name={Three-Axis Stabilization},
  description={\textcolor{catSS}{\textbf{[Space Science]}}},
}

\newglossaryentry{Gravity-Gradient-Stabilization}{ sort={Gravity Gradient Stabilization},
  name={Gravity Gradient Stabilization},
  description={\textcolor{catSS}{\textbf{[Space Science]}}},
}

\newglossaryentry{Magnetic-Torquer}{ sort={Magnetic Torquer},
  name={Magnetic Torquer},
  description={\textcolor{catSS}{\textbf{[Space Science]}}},
}

\newglossaryentry{Magnetorquer}{ sort={Magnetorquer},
  name={Magnetorquer},
  description={\textcolor{catSS}{\textbf{[Space Science]}}},
}

\newglossaryentry{Thermal-Control-System}{ sort={Thermal Control System},
  name={Thermal Control System},
  description={\textcolor{catSS}{\textbf{[Space Science]}}},
}

\newglossaryentry{Passive-Thermal-Control}{ sort={Passive Thermal Control},
  name={Passive Thermal Control},
  description={\textcolor{catSS}{\textbf{[Space Science]}}},
}

\newglossaryentry{Active-Thermal-Control}{ sort={Active Thermal Control},
  name={Active Thermal Control},
  description={\textcolor{catSS}{\textbf{[Space Science]}}},
}

\newglossaryentry{Multi-Layer-Insulation}{ sort={Multi-Layer Insulation},
  name={Multi-Layer Insulation},
  description={\textcolor{catSS}{\textbf{[Space Science]}}},
}

\newglossaryentry{Radiator}{ sort={Radiator},
  name={Radiator},
  description={\textcolor{catSS}{\textbf{[Space Science]}}},
}

\newglossaryentry{Heat-Pipe}{ sort={Heat Pipe},
  name={Heat Pipe},
  description={\textcolor{catSS}{\textbf{[Space Science]}}},
}

\newglossaryentry{Louver}{ sort={Louver},
  name={Louver},
  description={\textcolor{catSS}{\textbf{[Space Science]}}},
}

\newglossaryentry{Phase-Change-Material}{ sort={Phase Change Material},
  name={Phase Change Material},
  description={\textcolor{catSS}{\textbf{[Space Science]}}},
}

\newglossaryentry{Thermal-Blanket}{ sort={Thermal Blanket},
  name={Thermal Blanket},
  description={\textcolor{catSS}{\textbf{[Space Science]}}},
}

\newglossaryentry{Cold-Plate}{ sort={Cold Plate},
  name={Cold Plate},
  description={\textcolor{catSS}{\textbf{[Space Science]}}},
}

\newglossaryentry{Telemetry}{ sort={Telemetry},
  name={Telemetry},
  description={\textcolor{catSS}{\textbf{[Space Science]}}},
}

\newglossaryentry{Telecommand}{ sort={Telecommand},
  name={Telecommand},
  description={\textcolor{catSS}{\textbf{[Space Science]}}},
}

\newglossaryentry{Ground-Station}{ sort={Ground Station},
  name={Ground Station},
  description={\textcolor{catSS}{\textbf{[Space Science]}}},
}

\newglossaryentry{Antenna}{ sort={Antenna},
  name={Antenna},
  description={\textcolor{catSS}{\textbf{[Space Science]}}},
}

\newglossaryentry{High-Gain-Antenna}{ sort={High Gain Antenna},
  name={High Gain Antenna},
  description={\textcolor{catSS}{\textbf{[Space Science]}}},
}

\newglossaryentry{Low-Gain-Antenna}{ sort={Low Gain Antenna},
  name={Low Gain Antenna},
  description={\textcolor{catSS}{\textbf{[Space Science]}}},
}

\newglossaryentry{Parabolic-Antenna}{ sort={Parabolic Antenna},
  name={Parabolic Antenna},
  description={\textcolor{catSS}{\textbf{[Space Science]}}},
}

\newglossaryentry{Phased-Array}{ sort={Phased Array},
  name={Phased Array},
  description={\textcolor{catSS}{\textbf{[Space Science]}}},
}

\newglossaryentry{Transponder}{ sort={Transponder},
  name={Transponder},
  description={\textcolor{catSS}{\textbf{[Space Science]}}},
}

\newglossaryentry{RF-Communication}{ sort={RF Communication},
  name={RF Communication},
  description={\textcolor{catSS}{\textbf{[Space Science]}}},
}

\newglossaryentry{S-Band}{ sort={S Band},
  name={S Band},
  description={\textcolor{catSS}{\textbf{[Space Science]}}},
}

\newglossaryentry{X-Band}{ sort={X Band},
  name={X Band},
  description={\textcolor{catSS}{\textbf{[Space Science]}}},
}

\newglossaryentry{Ka-Band}{ sort={Ka Band},
  name={Ka Band},
  description={\textcolor{catSS}{\textbf{[Space Science]}}},
}

\newglossaryentry{Deep-Space-Network}{ sort={Deep Space Network},
  name={Deep Space Network},
  description={\textcolor{catSS}{\textbf{[Space Science]}}},
}

\newglossaryentry{DSN}{ sort={DSN},
  name={DSN},
  description={\textcolor{catSS}{\textbf{[Space Science]}}},
}

\newglossaryentry{Data-Rate}{ sort={Data Rate},
  name={Data Rate},
  description={\textcolor{catSS}{\textbf{[Space Science]}}},
}

\newglossaryentry{Compression}{ sort={Compression},
  name={Compression},
  description={\textcolor{catSS}{\textbf{[Space Science]}}},
}

\newglossaryentry{Error-Correction}{ sort={Error Correction},
  name={Error Correction},
  description={\textcolor{catSS}{\textbf{[Space Science]}}},
}

\newglossaryentry{Reed-Solomon-Code}{ sort={Reed-Solomon Code},
  name={Reed-Solomon Code},
  description={\textcolor{catSS}{\textbf{[Space Science]}}},
}

\newglossaryentry{Convolutional-Code}{ sort={Convolutional Code},
  name={Convolutional Code},
  description={\textcolor{catSS}{\textbf{[Space Science]}}},
}

\newglossaryentry{Turbo-Code}{ sort={Turbo Code},
  name={Turbo Code},
  description={\textcolor{catSS}{\textbf{[Space Science]}}},
}

\newglossaryentry{LDPC}{ sort={LDPC},
  name={LDPC},
  description={\textcolor{catSS}{\textbf{[Space Science]}}},
}

\newglossaryentry{Spacecraft-Bus}{ sort={Spacecraft Bus},
  name={Spacecraft Bus},
  description={\textcolor{catSS}{\textbf{[Space Science]}}},
}

\newglossaryentry{Structure}{ sort={Structure},
  name={Structure},
  description={\textcolor{catSS}{\textbf{[Space Science]}}},
}

\newglossaryentry{Honeycomb-Panel}{ sort={Honeycomb Panel},
  name={Honeycomb Panel},
  description={\textcolor{catSS}{\textbf{[Space Science]}}},
}

\newglossaryentry{Aluminum-Honeycomb}{ sort={Aluminum Honeycomb},
  name={Aluminum Honeycomb},
  description={\textcolor{catSS}{\textbf{[Space Science]}}},
}

\newglossaryentry{Carbon-Fiber-Structure}{ sort={Carbon Fiber Structure},
  name={Carbon Fiber Structure},
  description={\textcolor{catSS}{\textbf{[Space Science]}}},
}

\newglossaryentry{Deployable-Structure}{ sort={Deployable Structure},
  name={Deployable Structure},
  description={\textcolor{catSS}{\textbf{[Space Science]}}},
}

\newglossaryentry{Solar-Array-Deployment}{ sort={Solar Array Deployment},
  name={Solar Array Deployment},
  description={\textcolor{catSS}{\textbf{[Space Science]}}},
}

\newglossaryentry{Antenna-Deployment}{ sort={Antenna Deployment},
  name={Antenna Deployment},
  description={\textcolor{catSS}{\textbf{[Space Science]}}},
}

\newglossaryentry{Boom}{ sort={Boom},
  name={Boom},
  description={\textcolor{catSS}{\textbf{[Space Science]}}},
}

\newglossaryentry{Mast}{ sort={Mast},
  name={Mast},
  description={\textcolor{catSS}{\textbf{[Space Science]}}},
}

\newglossaryentry{Mechanism}{ sort={Mechanism},
  name={Mechanism},
  description={\textcolor{catSS}{\textbf{[Space Science]}}},
}

\newglossaryentry{Pyrotechnic-Device}{ sort={Pyrotechnic Device},
  name={Pyrotechnic Device},
  description={\textcolor{catSS}{\textbf{[Space Science]}}},
}

\newglossaryentry{Frangible-Nut}{ sort={Frangible Nut},
  name={Frangible Nut},
  description={\textcolor{catSS}{\textbf{[Space Science]}}},
}

\newglossaryentry{Cable-Cutter}{ sort={Cable Cutter},
  name={Cable Cutter},
  description={\textcolor{catSS}{\textbf{[Space Science]}}},
}

\newglossaryentry{Pin-Puller}{ sort={Pin Puller},
  name={Pin Puller},
  description={\textcolor{catSS}{\textbf{[Space Science]}}},
}

\newglossaryentry{Separation-Nut}{ sort={Separation Nut},
  name={Separation Nut},
  description={\textcolor{catSS}{\textbf{[Space Science]}}},
}

\newglossaryentry{Launch-Adapter}{ sort={Launch Adapter},
  name={Launch Adapter},
  description={\textcolor{catSS}{\textbf{[Space Science]}}},
}

\newglossaryentry{Payload-Adapter}{ sort={Payload Adapter},
  name={Payload Adapter},
  description={\textcolor{catSS}{\textbf{[Space Science]}}},
}

\newglossaryentry{Clamp-Band}{ sort={Clamp Band},
  name={Clamp Band},
  description={\textcolor{catSS}{\textbf{[Space Science]}}},
}

\newglossaryentry{Marman-Clamp}{ sort={Marman Clamp},
  name={Marman Clamp},
  description={\textcolor{catSS}{\textbf{[Space Science]}}},
}

\newglossaryentry{Lightband}{ sort={Lightband},
  name={Lightband},
  description={\textcolor{catSS}{\textbf{[Space Science]}}},
}

\newglossaryentry{Separation-System}{ sort={Separation System},
  name={Separation System},
  description={\textcolor{catSS}{\textbf{[Space Science]}}},
}

\newglossaryentry{Yo-Yo-Deploy}{ sort={Yo-Yo Deploy},
  name={Yo-Yo Deploy},
  description={\textcolor{catSS}{\textbf{[Space Science]}}},
}

\newglossaryentry{Spin-Table}{ sort={Spin Table},
  name={Spin Table},
  description={\textcolor{catSS}{\textbf{[Space Science]}}},
}

\newglossaryentry{Balance}{ sort={Balance},
  name={Balance},
  description={\textcolor{catSS}{\textbf{[Space Science]}}},
}

\newglossaryentry{Center-of-Mass}{ sort={Center of Mass},
  name={Center of Mass},
  description={\textcolor{catSS}{\textbf{[Space Science]}}},
}

\newglossaryentry{Margin-of-Safety}{ sort={Margin of Safety},
  name={Margin of Safety},
  description={\textcolor{catSS}{\textbf{[Space Science]}}},
}

\newglossaryentry{Qualification-Model}{ sort={Qualification Model},
  name={Qualification Model},
  description={\textcolor{catSS}{\textbf{[Space Science]}}},
}

\newglossaryentry{Flight-Model}{ sort={Flight Model},
  name={Flight Model},
  description={\textcolor{catSS}{\textbf{[Space Science]}}},
}

\newglossaryentry{Engineering-Model}{ sort={Engineering Model},
  name={Engineering Model},
  description={\textcolor{catSS}{\textbf{[Space Science]}}},
}

\newglossaryentry{Structural-Model}{ sort={Structural Model},
  name={Structural Model},
  description={\textcolor{catSS}{\textbf{[Space Science]}}},
}

\newglossaryentry{Thermal-Model}{ sort={Thermal Model},
  name={Thermal Model},
  description={\textcolor{catSS}{\textbf{[Space Science]}}},
}

\newglossaryentry{Protoflight-Model}{ sort={Protoflight Model},
  name={Protoflight Model},
  description={\textcolor{catSS}{\textbf{[Space Science]}}},
}

\newglossaryentry{Clean-Room}{ sort={Clean Room},
  name={Clean Room},
  description={\textcolor{catSS}{\textbf{[Space Science]}}},
}

\newglossaryentry{Cleanliness}{ sort={Cleanliness},
  name={Cleanliness},
  description={\textcolor{catSS}{\textbf{[Space Science]}}},
}

\newglossaryentry{Outgassing}{ sort={Outgassing},
  name={Outgassing},
  description={\textcolor{catSS}{\textbf{[Space Science]}}},
}

\newglossaryentry{Vacuum}{ sort={Vacuum},
  name={Vacuum},
  description={\textcolor{catSS}{\textbf{[Space Science]}}},
}

\newglossaryentry{Thermal-Vacuum-Test}{ sort={Thermal Vacuum Test},
  name={Thermal Vacuum Test},
  description={\textcolor{catSS}{\textbf{[Space Science]}}},
}

\newglossaryentry{Vibration-Test}{ sort={Vibration Test},
  name={Vibration Test},
  description={\textcolor{catSS}{\textbf{[Space Science]}}},
}

\newglossaryentry{Acoustic-Test}{ sort={Acoustic Test},
  name={Acoustic Test},
  description={\textcolor{catSS}{\textbf{[Space Science]}}},
}

\newglossaryentry{Shock-Test}{ sort={Shock Test},
  name={Shock Test},
  description={\textcolor{catSS}{\textbf{[Space Science]}}},
}

\newglossaryentry{EMIEMC-Test}{ sort={EMI/EMC Test},
  name={EMI/EMC Test},
  description={\textcolor{catSS}{\textbf{[Space Science]}}},
}

\newglossaryentry{Mass-Properties-Measurement}{ sort={Mass Properties Measurement},
  name={Mass Properties Measurement},
  description={\textcolor{catSS}{\textbf{[Space Science]}}},
}

\newglossaryentry{Spin-Balance}{ sort={Spin Balance},
  name={Spin Balance},
  description={\textcolor{catSS}{\textbf{[Space Science]}}},
}

\newglossaryentry{Alignment}{ sort={Alignment},
  name={Alignment},
  description={\textcolor{catSS}{\textbf{[Space Science]}}},
}

\newglossaryentry{Integration}{ sort={Integration},
  name={Integration},
  description={\textcolor{catSS}{\textbf{[Space Science]}}},
}

\newglossaryentry{Test}{ sort={Test},
  name={Test},
  description={\textcolor{catSS}{\textbf{[Space Science]}}},
}

\newglossaryentry{Launch-Campaign}{ sort={Launch Campaign},
  name={Launch Campaign},
  description={\textcolor{catSS}{\textbf{[Space Science]}}},
}

\newglossaryentry{Launch-Site}{ sort={Launch Site},
  name={Launch Site},
  description={\textcolor{catSS}{\textbf{[Space Science]}}},
}

\newglossaryentry{Cape-Canaveral}{ sort={Cape Canaveral},
  name={Cape Canaveral},
  description={\textcolor{catSS}{\textbf{[Space Science]}}},
}

\newglossaryentry{Kennedy-Space-Center}{ sort={Kennedy Space Center},
  name={Kennedy Space Center},
  description={\textcolor{catSS}{\textbf{[Space Science]}}},
}

\newglossaryentry{Vandenberg-SFB}{ sort={Vandenberg SFB},
  name={Vandenberg SFB},
  description={\textcolor{catSS}{\textbf{[Space Science]}}},
}

\newglossaryentry{Wallops-Flight-Facility}{ sort={Wallops Flight Facility},
  name={Wallops Flight Facility},
  description={\textcolor{catSS}{\textbf{[Space Science]}}},
}

\newglossaryentry{Baikonur}{ sort={Baikonur},
  name={Baikonur},
  description={\textcolor{catSS}{\textbf{[Space Science]}}},
}

\newglossaryentry{Kourou}{ sort={Kourou},
  name={Kourou},
  description={\textcolor{catSS}{\textbf{[Space Science]}}},
}

\newglossaryentry{Spaceport}{ sort={Spaceport},
  name={Spaceport},
  description={\textcolor{catSS}{\textbf{[Space Science]}}},
}

\newglossaryentry{Launch-Pad}{ sort={Launch Pad},
  name={Launch Pad},
  description={\textcolor{catSS}{\textbf{[Space Science]}}},
}

\newglossaryentry{Launch-Complex}{ sort={Launch Complex},
  name={Launch Complex},
  description={\textcolor{catSS}{\textbf{[Space Science]}}},
}

\newglossaryentry{Integration-Facility}{ sort={Integration Facility},
  name={Integration Facility},
  description={\textcolor{catSS}{\textbf{[Space Science]}}},
}

\newglossaryentry{Payload-Processing}{ sort={Payload Processing},
  name={Payload Processing},
  description={\textcolor{catSS}{\textbf{[Space Science]}}},
}

\newglossaryentry{Propellant-Loading}{ sort={Propellant Loading},
  name={Propellant Loading},
  description={\textcolor{catSS}{\textbf{[Space Science]}}},
}

\newglossaryentry{Countdown}{ sort={Countdown},
  name={Countdown},
  description={\textcolor{catSS}{\textbf{[Space Science]}}},
}

\newglossaryentry{Big-Bang}{ sort={Big Bang},
  name={Big Bang},
  description={\textcolor{catSS}{\textbf{[Space Science]}}},
}

\newglossaryentry{Dark-Matter}{ sort={Dark Matter},
  name={Dark Matter},
  description={\textcolor{catSS}{\textbf{[Space Science]}}},
}

\newglossaryentry{Dark-Energy}{ sort={Dark Energy},
  name={Dark Energy},
  description={\textcolor{catSS}{\textbf{[Space Science]}}},
}

\newglossaryentry{Black-Hole}{ sort={Black Hole},
  name={Black Hole},
  description={\textcolor{catSS}{\textbf{[Space Science]}}},
}

\newglossaryentry{Neutron-Star}{ sort={Neutron Star},
  name={Neutron Star},
  description={\textcolor{catSS}{\textbf{[Space Science]}}},
}

\newglossaryentry{White-Dwarf}{ sort={White Dwarf},
  name={White Dwarf},
  description={\textcolor{catSS}{\textbf{[Space Science]}}},
}

\newglossaryentry{Supernova}{ sort={Supernova},
  name={Supernova},
  description={\textcolor{catSS}{\textbf{[Space Science]}}},
}

\newglossaryentry{Nebula}{ sort={Nebula},
  name={Nebula},
  description={\textcolor{catSS}{\textbf{[Space Science]}}},
}

\newglossaryentry{Galaxy}{ sort={Galaxy},
  name={Galaxy},
  description={\textcolor{catSS}{\textbf{[Space Science]}}},
}

\newglossaryentry{Milky-Way}{ sort={Milky Way},
  name={Milky Way},
  description={\textcolor{catSS}{\textbf{[Space Science]}}},
}

\newglossaryentry{Light-Year}{ sort={Light-Year},
  name={Light-Year},
  description={\textcolor{catSS}{\textbf{[Space Science]}}},
}

\newglossaryentry{Astronomical-Unit}{ sort={Astronomical Unit},
  name={Astronomical Unit},
  description={\textcolor{catSS}{\textbf{[Space Science]}}},
}

\newglossaryentry{Parsec}{ sort={Parsec},
  name={Parsec},
  description={\textcolor{catSS}{\textbf{[Space Science]}}},
}

\newglossaryentry{Redshift}{ sort={Redshift},
  name={Redshift},
  description={\textcolor{catSS}{\textbf{[Space Science]}}},
}

\newglossaryentry{Cosmic-Microwave-Background}{ sort={Cosmic Microwave Background},
  name={Cosmic Microwave Background},
  description={\textcolor{catSS}{\textbf{[Space Science]}}},
}

\newglossaryentry{Terrestrial-Planet}{ sort={Terrestrial Planet},
  name={Terrestrial Planet},
  description={\textcolor{catSS}{\textbf{[Space Science]}}},
}

\newglossaryentry{Gas-Giant}{ sort={Gas Giant},
  name={Gas Giant},
  description={\textcolor{catSS}{\textbf{[Space Science]}}},
}

\newglossaryentry{Impact-Crater}{ sort={Impact Crater},
  name={Impact Crater},
  description={\textcolor{catSS}{\textbf{[Space Science]}}},
}

\newglossaryentry{Regolith}{ sort={Regolith},
  name={Regolith},
  description={\textcolor{catSS}{\textbf{[Space Science]}}},
}

\newglossaryentry{Cubesat}{ sort={Cubesat},
  name={Cubesat},
  description={\textcolor{catSS}{\textbf{[Space Science]}}},
}

\newglossaryentry{SmallSat}{ sort={SmallSat},
  name={SmallSat},
  description={\textcolor{catSS}{\textbf{[Space Science]}}},
}

\newglossaryentry{Space-Tourism}{ sort={Space Tourism},
  name={Space Tourism},
  description={\textcolor{catSS}{\textbf{[Space Science]}}},
}

\newglossaryentry{Planetary-Protection}{ sort={Planetary Protection},
  name={Planetary Protection},
  description={\textcolor{catSS}{\textbf{[Space Science]}}},
}



\begin{document}

\frontmatter
\begin{titlepage}
\centering
\vspace*{3cm}
{\Huge\bfseries \booktitle\par}
\vspace{1cm}
{\Large Five chapters covering\par}
\vspace{0.5cm}
{\large
\textcolor{catAD}{\textbf{Aerodynamics}} \quad
\textcolor{catAE}{\textbf{Aerospace}} \\
\textcolor{catRM}{\textbf{Rockets and Missiles}} \\
\textcolor{catD}{\textbf{Drones}} \quad
\textcolor{catSS}{\textbf{Space Science}}
\par}
\vfill
{\large\today\par}
\end{titlepage}

\tableofcontents
\clearpage
\mainmatter

\clearpage
\chapter*{Aerodynamics}
\addcontentsline{toc}{chapter}{Aerodynamics}
\markboth{Aerodynamics}{Aerodynamics}
\clearpage
\begin{multicols}{2}
\small
The study of air in motion and the forces it creates on solid bodies. From boundary layers to shock waves, airfoils to wind tunnels.\par
\vspace{0.5cm}
\glsadd{Aerodynamics}\dictentry{Aerodynamics}{Aerodynamics is the branch of fluid dynamics concerned with the motion of air and the forces acting on bodies immersed in it. It studies how air molecules interact with surfaces, producing lift, drag, and moments that govern flight. The discipline draws on conservation laws of mass, momentum, and energy to predict flow behavior. Aerodynamics is foundational to the design of aircraft, missiles, automobiles, and any vehicle moving through the atmosphere.}{catAD}
\glsadd{Aerodynamic-Center}\dictentry{Aerodynamic Center}{The Aerodynamic Center is the point on an airfoil about which the pitching moment coefficient remains constant regardless of angle of attack. For subsonic flow it typically lies near the quarter-chord point, shifting aft as the flow becomes supersonic. Locating the aerodynamic center simplifies stability analysis because moment changes due to angle of attack are eliminated. Aircraft designers use this point to set the required static margin for longitudinal stability.}{catAD}
\glsadd{Aerodynamic-Force}\dictentry{Aerodynamic Force}{Aerodynamic Force is the net force exerted on a body by the surrounding airflow. It is conventionally resolved into three perpendicular components: lift perpendicular to the freestream, drag parallel to it, and side force along the lateral axis. These components arise from the integrated pressure and shear stress distributions over the body surface. Accurate prediction of aerodynamic forces is essential for sizing wings, control surfaces, and structural members.}{catAD}
\glsadd{Aerodynamic-Heating}\dictentry{Aerodynamic Heating}{Aerodynamic Heating is the thermal energy generated when a body compresses and shears the air flowing around it at high speed. At supersonic and hypersonic velocities the stagnation temperature can exceed thousands of kelvins, demanding robust thermal protection. Heating intensity depends on Mach number, nose radius, and boundary layer state. Engineers must account for aerodynamic heating when selecting materials for high-speed aircraft and re-entry vehicles.}{catAD}
\glsadd{Aeroelasticity}\dictentry{Aeroelasticity}{Aeroelasticity is the study of the coupled interaction between aerodynamic loads and the elastic deformation of a structure. When airflow forces deform a wing, the changed shape alters the aerodynamic forces in return, creating a feedback loop. This coupling can produce dangerous phenomena such as flutter, divergence, and control surface reversal. Understanding aeroelastic behavior is critical for ensuring the structural integrity of lightweight, high-aspect-ratio wings on modern aircraft.}{catAD}
\glsadd{Aileron}\dictentry{Aileron}{An Aileron is a hinged control surface located on the trailing edge of each wing, typically near the wingtips. Deflecting the left and right ailerons in opposite directions creates a differential lift that rolls the aircraft about its longitudinal axis. Aileron design must balance roll responsiveness against adverse yaw, the tendency to yaw opposite the intended turn. Many aircraft supplement or replace conventional ailerons with spoilers or differential tail surfaces.}{catAD}
\glsadd{Airfoil}\dictentry{Airfoil}{An Airfoil is the two-dimensional cross-sectional shape of a wing, blade, or other surface designed to interact with a moving fluid to produce a useful force. Its geometry is defined by parameters such as chord, camber, and thickness distribution. The pressure difference between upper and lower surfaces generates lift according to Bernoulli's principle and circulation theory. Airfoil selection directly influences the lift-to-drag ratio, stall characteristics, and overall efficiency of the aerodynamic surface.}{catAD}
\glsadd{Angle-of-Attack}\dictentry{Angle of Attack}{The Angle of Attack is the angle between the chord line of an airfoil and the oncoming relative wind. Increasing the angle of attack raises the lift coefficient up to a critical value, beyond which the airflow separates and the wing stalls. Pilots control angle of attack through pitch inputs, and many aircraft carry dedicated sensors to provide this vital data. Proper management of angle of attack is essential for safe takeoff, maneuvering, and landing.}{catAD}
\glsadd{Angle-of-Incidence}\dictentry{Angle of Incidence}{The Angle of Incidence is the fixed angle between the wing chord line and the longitudinal axis of the fuselage. It is set during design to ensure that the wing operates at a favorable angle of attack during cruise. By tilting the wing upward relative to the fuselage, designers can reduce fuselage drag at cruise speed. This geometric parameter differs from the angle of attack, which varies continuously with flight conditions.}{catAD}
\glsadd{Anhedral}\dictentry{Anhedral}{Anhedral is the downward angular offset of a wing from the horizontal, measured at the root. It reduces lateral stability by allowing the low wing of a banked aircraft to generate more lift, tending to increase the bank angle. Designers apply anhedral to highly swept or high-wing aircraft that already possess excessive natural lateral stability. The result is improved maneuverability and more responsive roll behavior.}{catAD}
\glsadd{Aspect-Ratio}\dictentry{Aspect Ratio}{Aspect Ratio is the ratio of a wing's span squared to its planform area, equivalently span divided by mean chord. A high aspect ratio reduces induced drag, making long, narrow wings efficient for cruise, as seen on gliders. Low aspect ratio wings favor structural strength and rapid roll response, common in fighter aircraft. The choice of aspect ratio represents a key trade-off between aerodynamic efficiency and structural weight.}{catAD}
\glsadd{Attached-Flow}\dictentry{Attached Flow}{Attached Flow describes air that remains in continuous contact with a surface as it moves along it. In this state the boundary layer stays thin and orderly, producing predictable lift and minimal pressure drag. Attached flow is the design condition for wings and control surfaces during normal flight. If the flow separates, lift drops sharply and drag rises, leading to stall.}{catAD}
\glsadd{Autorotation}\dictentry{Autorotation}{Autorotation is a flight condition in which a helicopter rotor is driven entirely by the upward flow of air through the disc rather than by engine power. It occurs during a powered descent when the aircraft sinks fast enough for the relative wind to spin the blades. Pilots use autorotation as an emergency landing technique following engine failure. The stored rotational energy is converted into a flare just before touchdown.}{catAD}
\glsadd{Boundary-Layer}\dictentry{Boundary Layer}{The Boundary Layer is the thin region of fluid adjacent to a surface where velocity increases from zero at the wall to the freestream value. Viscous effects are concentrated in this layer, producing skin friction drag. The boundary layer may be laminar, turbulent, or transitional, each with different drag and heat transfer characteristics. Managing boundary layer behavior is central to reducing drag and delaying flow separation.}{catAD}
\glsadd{Boundary-Layer-Separation}\dictentry{Boundary Layer Separation}{Boundary Layer Separation occurs when the boundary layer detaches from the surface due to an adverse pressure gradient that the slow-moving fluid cannot overcome. The separated region produces large pressure drag and unsteady wake formation. On wings, separation leads to loss of lift and the onset of stall. Techniques such as vortex generators and boundary layer suction are employed to delay separation.}{catAD}
\glsadd{Boundary-Layer-Transition}\dictentry{Boundary Layer Transition}{Boundary Layer Transition is the process by which a laminar boundary layer becomes turbulent. Triggered by surface roughness, pressure gradients, or freestream disturbances, transition increases skin friction but also energizes the layer, delaying separation. Predicting transition location is essential for accurate drag estimation. Laminar flow control technologies aim to delay transition and reduce friction drag on aircraft surfaces.}{catAD}
\glsadd{Camber}\dictentry{Camber}{Camber is the measure of curvature of the mean line of an airfoil, describing how much it deviates from a straight chord. Positive camber increases the lift coefficient at a given angle of attack and shifts the zero-lift angle to a negative value. Highly cambered sections are used on takeoff and landing flaps to boost low-speed lift. Symmetrical airfoils have zero camber and produce no lift at zero angle of attack.}{catAD}
\glsadd{Center-of-Pressure}\dictentry{Center of Pressure}{The Center of Pressure is the point on an airfoil where the total aerodynamic force effectively acts. Unlike the aerodynamic center, the center of pressure moves fore and aft with changing angle of attack. Its migration complicates stability analysis, which is why the aerodynamic center is often preferred. Knowledge of the center of pressure is still needed for structural load calculations.}{catAD}
\glsadd{Chord}\dictentry{Chord}{Chord is the straight-line distance from the leading edge to the trailing edge of an airfoil. It serves as the characteristic length in Reynolds number calculations and non-dimensional aerodynamic coefficients. Chord may vary along the span in tapered or swept wings, in which case the mean aerodynamic chord is used. Accurate chord measurement is essential for wind tunnel testing and computational fluid dynamics validation.}{catAD}
\glsadd{Chord-Line}\dictentry{Chord Line}{The Chord Line is the straight line connecting the leading edge to the trailing edge of an airfoil. It serves as the geometric reference from which the angle of attack and camber are measured. The chord line is distinct from the camber line, which follows the mean curvature of the section. Aerodynamic properties such as pitching moment are often referenced to the chord line.}{catAD}
\glsadd{Circulation}\dictentry{Circulation}{Circulation is the line integral of the velocity vector around a closed loop in a fluid. In aerodynamic theory, circulation around an airfoil is directly proportional to the lift it generates, as stated by the Kutta-Joukowski theorem. Circulation arises from the asymmetry of flow over a cambered surface and the enforcement of the Kutta condition. This concept unifies potential flow theory with practical lift prediction.}{catAD}
\glsadd{Compressible-Flow}\dictentry{Compressible Flow}{Compressible Flow is fluid motion in which density changes are significant and cannot be ignored. It becomes important at Mach numbers above roughly 0.3, where shock waves and expansion fans appear. Compressible flow analysis governs the design of jet engines, supersonic inlets, and re-entry bodies. The governing equations must couple the energy equation with continuity and momentum to account for density variations.}{catAD}
\glsadd{Critical-Angle-of-Attack}\dictentry{Critical Angle of Attack}{The Critical Angle of Attack is the angle at which the airflow over a wing separates to the point of producing maximum lift coefficient, beyond which lift decreases sharply. This marks the onset of aerodynamic stall. The critical angle is a fixed property of the airfoil section and is not significantly affected by airspeed. Pilots must maintain awareness of this limit to avoid loss of control during low-speed flight.}{catAD}
\glsadd{Critical-Mach-Number}\dictentry{Critical Mach Number}{The Critical Mach Number is the freestream Mach number at which sonic flow first appears somewhere on the airframe, typically at the thickest point of the wing. Above this speed, local supersonic pockets and shock waves begin to form, dramatically increasing drag. The critical Mach number marks the onset of transonic flight regime. Swept wings and thin airfoil sections raise the critical Mach number, extending efficient subsonic cruise.}{catAD}
\glsadd{Dihedral}\dictentry{Dihedral}{Dihedral is the upward angular offset of a wing measured from the horizontal at the root. When a gust rolls the aircraft, the lower wing develops a higher effective angle of attack and generates more lift, producing a restoring rolling moment. This geometric feature provides passive lateral stability. Designers tune the amount of dihedral to balance stability with maneuverability requirements.}{catAD}
\glsadd{Downwash}\dictentry{Downwash}{Downwash is the downward deflection of air caused by a wing generating lift during flight. The wing pushes air downward per Newton's third law, reducing the effective angle of attack behind it. Downwash affects tail surfaces and trailing aircraft, influencing stability and wake turbulence. Understanding downwash is critical for formation flying and tail surface design.}{catAD}
\glsadd{Drag}\dictentry{Drag}{Drag is the component of aerodynamic force that opposes the motion of a body through the air. It arises from skin friction, pressure differences, and induced effects related to lift production. Reducing drag is a primary objective in aircraft design because it directly affects fuel consumption and range. Drag is non-dimensionalized as a drag coefficient to allow comparison across different scales and speeds.}{catAD}
\glsadd{Drag-Coefficient}\dictentry{Drag Coefficient}{The Drag Coefficient is a dimensionless quantity that relates the drag force on a body to the dynamic pressure and reference area. It is defined as drag divided by the product of dynamic pressure and wing area. A lower drag coefficient indicates a more aerodynamically efficient shape. Engineers track drag coefficient across Mach numbers and angles of attack to characterize the full performance envelope.}{catAD}
\glsadd{Dynamic-Stall}\dictentry{Dynamic Stall}{Dynamic Stall is a transient aerodynamic phenomenon that occurs when an airfoil rapidly changes its angle of attack beyond the static stall angle. A strong vortex sheds from the leading edge, producing a temporary surge in lift before a dramatic drop. The process is hysteretic, with different behavior during pitch-up and pitch-down. Dynamic stall is critical in helicopter rotor blade aerodynamics and in fighter aircraft performing aggressive maneuvers.}{catAD}
\glsadd{Dynamic-Stability}\dictentry{Dynamic Stability}{Dynamic Stability refers to the time-history behavior of an aircraft's motion following a disturbance. An aircraft is dynamically stable if its oscillations decay over time, even if it is statically unstable. The analysis considers coupled modes such as phugoid, short period, Dutch roll, and spiral. Dynamic stability is assessed through eigenvalue analysis of the linearized equations of motion.}{catAD}
\glsadd{Dutch-Roll}\dictentry{Dutch Roll}{Dutch Roll is a coupled lateral-directional oscillation combining yaw and roll motions, resembling the skating motion of Dutch ice skaters. It arises when directional stability is strong relative to lateral stability. The mode is typically lightly damped and can cause passenger discomfort. Yaw dampers are commonly installed on airliners to suppress Dutch Roll oscillations.}{catAD}
\glsadd{Elevator}\dictentry{Elevator}{The Elevator is a hinged control surface on the horizontal tail that controls the aircraft's pitch attitude. Deflecting the elevator changes the tail lift, creating a pitching moment about the lateral axis. Pilots use the elevator to climb, descend, and maintain trim. Some aircraft use an all-moving stabilator instead, which avoids compressibility effects at high speeds.}{catAD}
\glsadd{Elevon}\dictentry{Elevon}{An Elevon is a combined elevator and aileron control surface found on tailless or delta-wing aircraft. It is mounted on the trailing edge of the wing and can deflect symmetrically for pitch control or differentially for roll control. Elevons simplify the control system by merging two functions into one surface. They are used on aircraft such as the Concorde and many flying-wing designs.}{catAD}
\glsadd{Expansion-Wave}\dictentry{Expansion Wave}{An Expansion Wave is a continuous compression-free region in supersonic flow that occurs when the flow turns around a convex corner. As the air expands, its pressure, temperature, and density decrease while velocity increases. Unlike a shock wave, an expansion wave is isentropic and does not increase entropy. Expansion fans are a key feature in supersonic nozzle and external aerodynamic analysis.}{catAD}
\glsadd{Flap}\dictentry{Flap}{A Flap is a high-lift device attached to the trailing edge of a wing. When extended, it increases the wing's camber and sometimes its area, raising the maximum lift coefficient for takeoff and landing. Common types include plain, split, slotted, and Fowler flaps. Flap deployment also increases drag, which is useful for descent and speed management.}{catAD}
\glsadd{Flow-Separation}\dictentry{Flow Separation}{Flow Separation occurs when the boundary layer detaches from a surface, causing loss of lift and increased drag. It happens when airflow cannot follow the surface contour due to adverse pressure gradients. Separation leads to stall conditions and reduced effectiveness of control surfaces. Engineers use vortex generators and careful airfoil design to delay separation and improve performance.}{catAD}
\glsadd{Form-Drag}\dictentry{Form Drag}{Form Drag, also called pressure drag, results from the pressure difference between the front and rear of a body immersed in a flow. It depends primarily on the shape and frontal area of the object. Streamlining reduces form drag by allowing the flow to close smoothly behind the body. Form drag is a major component of parasite drag on non-lifting bodies such as fuselages and landing gear.}{catAD}
\glsadd{Friction-Drag}\dictentry{Friction Drag}{Friction Drag is the component of drag caused by shear stresses acting along the surface of a body in contact with a viscous fluid. It depends on the wetted area and the state of the boundary layer, being higher for turbulent flows. Laminar flow surfaces experience significantly less friction drag than turbulent ones. Surface polishing and laminar flow airfoils are techniques used to minimize friction drag.}{catAD}
\glsadd{Hypersonic-Flow}\dictentry{Hypersonic Flow}{Hypersonic Flow refers to airflow at Mach numbers greater than five, where extreme compression heats the gas enough to cause real-gas effects such as dissociation and ionization. Viscous interaction, chemical non-equilibrium, and intense aerodynamic heating dominate the flow physics. Hypersonic regimes govern the re-entry of spacecraft and the design of scramjet-powered vehicles. Accurate prediction requires coupled thermochemical models.}{catAD}
\glsadd{Incompressible-Flow}\dictentry{Incompressible Flow}{Incompressible Flow is an idealization in which fluid density is assumed constant throughout the flow field. This approximation is valid for low-speed aerodynamics, typically below Mach 0.3. Under this assumption the governing equations simplify considerably, allowing analytical and semi-analytical solutions. Incompressible flow theory underpins much of classical airfoil and wing analysis.}{catAD}
\glsadd{Induced-Drag}\dictentry{Induced Drag}{Induced Drag is the drag penalty that accompanies lift production. It arises from the finite span of a wing, which generates tip vortices that redirect the lift vector backward. Induced drag decreases with increasing aspect ratio and is proportional to the square of the lift coefficient. It dominates the total drag at low speeds and high angles of attack, such as during climb and landing.}{catAD}
\glsadd{Induced-Velocity}\dictentry{Induced Velocity}{Induced Velocity is the component of airflow velocity produced by the bound and trailing vortex system of a lifting wing. It manifests primarily as a downward velocity, or downwash, in the wake. This downwash reduces the effective angle of attack seen by each wing section. In momentum theory, induced velocity is used to calculate the induced drag directly from the downwash at the lifting line.}{catAD}
\glsadd{Interference-Drag}\dictentry{Interference Drag}{Interference Drag is the additional drag produced when the flow fields of two or more aircraft components interact in unexpected ways. Wing-fuselage junctions and pylon-nacelle intersections are common sources. The merging boundary layers and pressure fields can cause local separation and increased losses. Fairings and careful geometric blending are used to minimize interference effects.}{catAD}
\glsadd{Inviscid-Flow}\dictentry{Inviscid Flow}{Inviscid Flow is a theoretical idealization in which fluid viscosity is neglected entirely. Under this assumption there is no skin friction and the boundary layer does not form. Inviscid theory, such as potential flow, provides useful predictions for lift and pressure distributions away from surfaces. It forms the outer-boundary condition for boundary layer methods in coupled viscous-inviscid analyses.}{catAD}
\glsadd{Irrotational-Flow}\dictentry{Irrotational Flow}{Irrotational Flow is a flow field in which the vorticity is zero at every point, meaning fluid particles do not rotate about their own axes. Such flows can be described by a velocity potential that satisfies Laplace's equation. Irrotational, inviscid, incompressible flow forms the basis of classical potential flow theory. This framework enables analytical solutions for flows around airfoils and bodies.}{catAD}
\glsadd{Kutta-Condition}\dictentry{Kutta Condition}{The Kutta Condition is a boundary condition applied at the trailing edge of an airfoil stating that the flow must leave smoothly. It eliminates the infinite velocity that potential theory would otherwise predict at a sharp trailing edge. The condition effectively fixes the circulation and hence the lift of the airfoil. Without it, potential flow would predict zero lift for any body.}{catAD}
\glsadd{Kutta-Joukowski-Theorem}\dictentry{Kutta-Joukowski Theorem}{The Kutta-Joukowski Theorem states that the lift per unit span on a two-dimensional body in an inviscid, incompressible flow equals the product of freestream density, velocity, and circulation. It is one of the most important results in classical aerodynamics. The theorem directly links the concept of circulation to measurable lift force. It remains a cornerstone for validating computational and experimental methods.}{catAD}
\glsadd{Laminar-Flow}\dictentry{Laminar Flow}{Laminar Flow is a regime of fluid motion in which adjacent layers slide smoothly over one another without mixing. Laminar boundary layers have low skin friction compared to turbulent ones but are more susceptible to separation under adverse pressure gradients. Achieving extensive laminar flow on aircraft surfaces is a goal of drag reduction research. Surface quality, pressure gradient, and freestream turbulence all influence the extent of laminar flow.}{catAD}
\glsadd{Lateral-Stability}\dictentry{Lateral Stability}{Lateral Stability is the tendency of an aircraft to resist and recover from rolling disturbances about its longitudinal axis. It is influenced by wing dihedral, sweep, and the vertical tail. A laterally stable aircraft will return to wings-level flight after a bank disturbance. Excessive lateral stability can make the aircraft sluggish in roll, so designers carefully balance this property.}{catAD}
\glsadd{Leading-Edge}\dictentry{Leading Edge}{The Leading Edge is the foremost part of an airfoil or wing where the oncoming airflow first makes contact. Its radius and shape influence the pressure distribution, stall behavior, and sensitivity to angle of attack. Sharp leading edges favor high-speed flight, while rounded leading edges improve low-speed performance. Leading edge devices such as slats and slots are used to enhance maximum lift.}{catAD}
\glsadd{Lift}\dictentry{Lift}{Lift is the aerodynamic force component perpendicular to the freestream that supports the weight of an aircraft in level flight. It is generated by the pressure difference between upper and lower wing surfaces. Lift depends on airspeed, air density, wing area, and the lift coefficient. Understanding and maximizing lift while minimizing its associated drag is a central challenge in aerodynamic design.}{catAD}
\glsadd{Lift-Coefficient}\dictentry{Lift Coefficient}{The Lift Coefficient is a dimensionless number that relates lift to dynamic pressure and wing area. It is defined as lift divided by the product of dynamic pressure and reference area. The lift coefficient varies with angle of attack, Mach number, and Reynolds number. Plotting lift coefficient against angle of attack reveals the stall characteristics of an airfoil or wing.}{catAD}
\glsadd{Longitudinal-Stability}\dictentry{Longitudinal Stability}{Longitudinal Stability is the tendency of an aircraft to resist and recover from pitch disturbances about its lateral axis. It requires that the center of gravity lie ahead of the neutral point. A longitudinally stable aircraft generates a restoring pitching moment when displaced from its trim angle of attack. Static margin, the distance between CG and neutral point, quantifies this stability.}{catAD}
\glsadd{Mach-Angle}\dictentry{Mach Angle}{The Mach Angle is the half-angle of the cone formed by the shock wave from a point source in supersonic flow. It is calculated as the inverse sine of one divided by the Mach number. A higher Mach number produces a narrower Mach cone with a smaller angle. This geometric relationship is fundamental in supersonic aerodynamics and sonic boom prediction.}{catAD}
\glsadd{Mach-Cone}\dictentry{Mach Cone}{The Mach Cone is the conical wavefront produced by a small disturbance moving at supersonic speed through a fluid. The cone's apex is at the disturbance source and its half-angle is the Mach angle. Flow properties change abruptly across the cone surface. The Mach cone concept explains why supersonic aircraft produce shock waves that reach the ground as sonic booms.}{catAD}
\glsadd{Mean-Aerodynamic-Chord}\dictentry{Mean Aerodynamic Chord}{The Mean Aerodynamic Chord is the chord of an equivalent rectangular wing that would produce the same pitching moment characteristics as the actual wing. It is used as the reference length for longitudinal stability and control calculations. The MAC accounts for spanwise variations in chord, twist, and airfoil section. Center of gravity positions are often expressed as a percentage of MAC.}{catAD}
\glsadd{Momentum-Thickness}\dictentry{Momentum Thickness}{Momentum Thickness is an integral boundary layer parameter that quantifies the momentum deficit caused by viscous retardation of the flow near a surface. It is directly related to the skin friction drag experienced by the body. Momentum thickness grows along the surface as the boundary layer develops. It is one of several thickness scales used to characterize boundary layer profiles.}{catAD}
\glsadd{Navier-Stokes-Equations}\dictentry{Navier-Stokes Equations}{The Navier-Stokes Equations are the fundamental partial differential equations governing viscous fluid motion. They express conservation of momentum with viscous stress terms included. These equations describe phenomena from boundary layer behavior to turbulence, but no general analytical solution exists. Computational fluid dynamics solves them numerically for practical engineering problems.}{catAD}
\glsadd{Neutral-Point}\dictentry{Neutral Point}{The Neutral Point is the longitudinal position where the pitching moment coefficient is independent of angle of attack. When the center of gravity coincides with the neutral point, the aircraft has zero static margin and is neutrally stable. Moving the CG ahead of this point provides static stability. The neutral point shifts aft as Mach number increases into the supersonic regime.}{catAD}
\glsadd{Normal-Shock}\dictentry{Normal Shock}{A Normal Shock is a standing wave perpendicular to the direction of supersonic flow. Across the shock, pressure, temperature, and density rise abruptly while the flow decelerates to subsonic speed. Entropy increases across the shock, making it an irreversible process. Normal shocks appear in supersonic inlets and at the throat of converging-diverging nozzles.}{catAD}
\glsadd{Nozzle}\dictentry{Nozzle}{A Nozzle is a duct with a varying cross-section that accelerates a fluid to high velocity by converting pressure energy into kinetic energy. In aerodynamics, wind tunnel nozzles shape and accelerate the test section flow to the desired Mach number. The convergent-divergent (de Laval) nozzle design enables supersonic flow by choking at the throat and expanding through the divergent section. Nozzle contour design controls flow uniformity, turbulence levels, and boundary layer growth in experimental facilities.}{catAD}
\glsadd{Oblique-Shock}\dictentry{Oblique Shock}{An Oblique Shock is a shock wave that forms at an angle to the direction of supersonic flow, typically from a concave corner or wedge. The flow decelerates and turns toward the surface across the shock. Oblique shocks are weaker than normal shocks and produce less entropy increase. They are a fundamental feature of supersonic inlet and external compression design.}{catAD}
\glsadd{Parasite-Drag}\dictentry{Parasite Drag}{Parasite Drag is the sum of all drag components that are not associated with lift production, including skin friction, form drag, and interference drag. It increases with the square of airspeed and is a major factor in determining cruise efficiency. Streamlining, surface finish, and component integration all affect parasite drag. Reducing parasite drag is a key objective in transport aircraft design.}{catAD}
\glsadd{Phugoid}\dictentry{Phugoid}{Phugoid is a long-period longitudinal oscillation in which the aircraft trades kinetic energy for potential energy, oscillating in speed and altitude while the angle of attack remains nearly constant. The period can range from tens of seconds in light aircraft to minutes in large transports. Phugoid motion is typically lightly damped. Autopilot systems often include a phugoid suppression function.}{catAD}
\glsadd{Pitching-Moment}\dictentry{Pitching Moment}{Pitching Moment is the aerodynamic moment acting about an aircraft's lateral axis, tending to rotate the nose up or down. It arises from the chordwise distribution of pressure and is referenced to a specific point such as the aerodynamic center. A nose-up pitching moment is conventionally positive. Pilots control pitching moment with the elevator and horizontal tail.}{catAD}
\glsadd{Pitot-Tube}\dictentry{Pitot Tube}{A Pitot Tube is an L-shaped probe that faces the oncoming flow to measure stagnation pressure. Combined with a static pressure port, it provides the dynamic pressure from which airspeed is calculated. Pitot tubes are standard instruments on all aircraft for indicated airspeed measurement. Blockage by ice or debris is a recognized hazard that can lead to erroneous speed readings.}{catAD}
\glsadd{Prandtl-Meyer-Expansion}\dictentry{Prandtl-Meyer Expansion}{Prandtl-Meyer Expansion describes the isentropic acceleration of supersonic flow around a convex corner through a centered expansion fan. The flow velocity increases while pressure and temperature drop across the fan. The expansion angle is a function of the upstream and downstream Mach numbers through the Prandtl-Meyer function. This mechanism is essential for analyzing supersonic external flows and nozzle contours.}{catAD}
\glsadd{Pressure-Drag}\dictentry{Pressure Drag}{Pressure Drag is the component of drag arising from the integrated pressure distribution over a body. It depends on the body shape and the extent of flow separation at the rear. Streamlined bodies with gentle rear closures experience low pressure drag. For bluff bodies pressure drag dominates due to massive separation.}{catAD}
\glsadd{Propeller}\dictentry{Propeller}{A Propeller is a rotating set of airfoil-shaped blades that converts engine power into thrust by accelerating a stream of air rearward. Blade element theory describes how each section generates lift with a forward thrust component. Pitch angle, rotational speed, and blade shape determine the air mass moved and resulting thrust. Propellers are most efficient at subsonic speeds, which is why they dominate general aviation and UAV propulsion. Advances in composite materials and computational blade design continue to improve propulsive efficiency.}{catAD}
\glsadd{Reattachment}\dictentry{Reattachment}{Reattachment is the process by which a previously separated boundary layer re-establishes contact with the surface downstream. It can occur when the flow encounters a favorable pressure gradient or a geometric feature that redirects the streamlines. Reattachment reduces the separated wake size and lowers pressure drag. It plays a role in dynamic stall recovery and in the design of diffuser passages.}{catAD}
\glsadd{Relative-Wind}\dictentry{Relative Wind}{Relative Wind is the direction and speed of the airflow as perceived by an airfoil or aircraft in motion. It is equal and opposite to the velocity of the aircraft through the air. The angle between the relative wind and the chord line defines the angle of attack. Understanding relative wind is essential for analyzing all aerodynamic forces and moments.}{catAD}
\glsadd{Reynolds-Number}\dictentry{Reynolds Number}{The Reynolds Number is a dimensionless parameter representing the ratio of inertial forces to viscous forces in a flow. It determines the transition from laminar to turbulent flow and scales boundary layer behavior. Reynolds number is computed using a characteristic length such as chord and the freestream velocity. Matching Reynolds number between wind tunnel models and full-scale aircraft is critical for accurate data correlation.}{catAD}
\glsadd{Rolling-Moment}\dictentry{Rolling Moment}{Rolling Moment is the aerodynamic moment about an aircraft's longitudinal axis, tending to rotate one wing up and the other down. It is generated by asymmetric lift distributions, primarily through aileron deflection. Rolling moment coefficient is non-dimensionalized by dynamic pressure, wing area, and wingspan. Adequate rolling moment authority is essential for control and maneuvering.}{catAD}
\glsadd{Rotor}\dictentry{Rotor}{A Rotor is a rotating assembly of blades that generates both lift and propulsive force on rotary-wing aircraft. Each blade functions as an airfoil experiencing a complex, time-varying aerodynamic environment. Rotor dynamics involve flapping, lead-lag, and torsional degrees of freedom. The rotor is the defining component of helicopters and autogyros.}{catAD}
\glsadd{Rudder}\dictentry{Rudder}{The Rudder is a vertical control surface mounted on the tail that controls yaw about the vertical axis. Deflecting the rudder creates a side force on the tail, producing a yawing moment. Pilots use the rudder for coordination in turns, crosswind landings, and engine out scenarios. The rudder works in conjunction with ailerons for effective turning.}{catAD}
\glsadd{Schlieren-Photography}\dictentry{Schlieren Photography}{Schlieren Photography is an optical visualization technique that reveals density gradients in transparent media. A collimated light source passes through the test region and is focused onto a knife edge, where refracted rays are either blocked or passed. This produces an image where brightness corresponds to the first derivative of density. Schlieren systems are widely used in supersonic wind tunnels to photograph shock waves.}{catAD}
\glsadd{Separated-Flow}\dictentry{Separated Flow}{Separated Flow is a region where the boundary layer has detached from the surface, creating a recirculating or dead-air zone. Pressure in the separated region is relatively uniform and lower than the freestream, causing high pressure drag. Separated flow often exhibits unsteady behavior with vortex shedding. It is a primary contributor to stall on wings and performance loss in diffusers.}{catAD}
\glsadd{Shock-Wave}\dictentry{Shock Wave}{A Shock Wave is an extremely thin discontinuity across which flow properties change abruptly. Across a shock, pressure, temperature, and density rise sharply while velocity drops. Shocks are dissipative structures that increase entropy. They are omnipresent in supersonic and hypersonic flight, significantly affecting drag, heating, and engine performance.}{catAD}
\glsadd{Short-Period-Mode}\dictentry{Short Period Mode}{The Short Period Mode is a heavily damped, high-frequency longitudinal oscillation primarily involving changes in angle of attack with relatively constant speed. It is the dominant short-term response of an aircraft to pitch disturbances. A well-damped short period mode is essential for good handling qualities. Its frequency and damping are key criteria in flying qualities specifications.}{catAD}
\glsadd{Skin-Friction}\dictentry{Skin Friction}{Skin Friction is the shear stress exerted by a viscous fluid flowing over a surface, arising from the velocity gradient at the wall. It is a direct consequence of the no-slip condition. Skin friction contributes significantly to total drag, especially on streamlined bodies with large wetted areas. Reducing skin friction through laminar flow control is an active area of research.}{catAD}
\glsadd{Skin-Friction-Drag}\dictentry{Skin Friction Drag}{Skin Friction Drag is the drag force resulting from the integrated wall shear stress over the surface of a body. It is directly proportional to the wetted area and depends strongly on whether the boundary layer is laminar or turbulent. For streamlined bodies such as wings, skin friction drag constitutes the majority of total drag. Surface coatings and laminar flow technologies aim to reduce this component.}{catAD}
\glsadd{Slat}\dictentry{Slat}{A Slat is a high-lift device attached to the leading edge of a wing. When deployed, it creates a slot that channels high-energy air over the upper surface, delaying boundary layer separation. This increases the maximum lift coefficient and lowers the stall speed. Slats are essential for safe takeoff and landing at the lower speeds required at airports with short runways.}{catAD}
\glsadd{Sonic-Boom}\dictentry{Sonic Boom}{A Sonic Boom is the impulsive noise generated when the shock waves from a supersonic vehicle coalesce and reach the ground. It appears as a double pressure rise resembling an explosion. The intensity depends on vehicle size, shape, altitude, and speed. Sonic booms have limited the operational envelope of supersonic transports over populated areas.}{catAD}
\glsadd{Spoiler}\dictentry{Spoiler}{A Spoiler is a flat or slightly curved panel mounted on the upper wing surface that can be raised into the airflow. It disrupts the smooth flow, reducing lift and increasing drag on that wing panel. Spoilers are used symmetrically as speed brakes and differentially for roll control. They offer rapid lift dumping, useful during landing rollout.}{catAD}
\glsadd{Stall}\dictentry{Stall}{Stall is the aerodynamic condition in which a wing exceeds its critical angle of attack and the airflow separates, causing a sudden loss of lift. It results from the inability of the boundary layer to remain attached under a strong adverse pressure gradient. Stall recovery requires reducing the angle of attack below the critical value. Stall warning and prevention systems are mandated on all civil aircraft.}{catAD}
\glsadd{Static-Margin}\dictentry{Static Margin}{Static Margin is the longitudinal distance between the neutral point and the center of gravity, expressed as a percentage of mean aerodynamic chord. A positive static margin means the CG is ahead of the neutral point, providing inherent static stability. Larger static margins increase stability but may reduce maneuverability and increase trim drag.}{catAD}
\glsadd{Static-Port}\dictentry{Static Port}{A Static Port is a flush-mounted opening on the aircraft skin that samples the ambient static pressure of the surrounding air. It feeds the altimeter, vertical speed indicator, and airspeed indicator. Static ports are carefully located to minimize errors from local flow acceleration. Alternate static sources are provided for use if the primary ports become blocked.}{catAD}
\glsadd{Static-Stability}\dictentry{Static Stability}{Static Stability describes the initial tendency of an aircraft to return toward its equilibrium position after a disturbance. If the aircraft tends to return, it is statically stable; if it moves further away, it is statically unstable. Static stability does not guarantee that oscillations will decay, which is a dynamic stability property. The degree of static stability is measured by the static margin.}{catAD}
\glsadd{Steady-Flow}\dictentry{Steady Flow}{Steady Flow is a condition in which all fluid properties at any given point in the flow field remain constant over time. In steady flow, streamlines, pathlines, and streaklines all coincide. This simplification greatly reduces the complexity of aerodynamic analysis. Most cruise-condition analyses assume steady flow as a first approximation.}{catAD}
\glsadd{Subsonic-Flow}\dictentry{Subsonic Flow}{Subsonic Flow is airflow in which the Mach number remains below one everywhere. Pressure disturbances propagate faster than the flow and communicate upstream, allowing the fluid to adjust smoothly around bodies. No shock waves form in purely subsonic flow. Subsonic aerodynamics governs the majority of general aviation and commercial aircraft operations.}{catAD}
\glsadd{Supersonic-Flow}\dictentry{Supersonic Flow}{Supersonic Flow is airflow in which the Mach number exceeds one everywhere. Pressure disturbances cannot propagate upstream, so the fluid cannot anticipate the presence of obstacles. Shock waves and expansion fans dominate the flow structure. Supersonic aerodynamics governs the design of military fighters, supersonic transports, and missile systems.}{catAD}
\glsadd{Sweep-Angle}\dictentry{Sweep Angle}{Sweep Angle is the angle between the wing's quarter-chord line and the lateral axis of the aircraft. Sweeping the wing delays the onset of compressibility drag rise by reducing the effective Mach number normal to the leading edge. Higher sweep angles allow higher cruise speeds but increase structural weight and low-speed handling challenges.}{catAD}
\glsadd{Thickness-Ratio}\dictentry{Thickness Ratio}{Thickness Ratio is the ratio of the maximum thickness of an airfoil to its chord length. It directly influences the lift capability, drag, and structural depth available for fuel and landing gear. Thicker airfoils provide more volume but produce higher form drag and lower critical Mach numbers. The choice of thickness ratio balances aerodynamic and structural requirements.}{catAD}
\glsadd{Thrust}\dictentry{Thrust}{Thrust is the propulsive force that moves an aircraft forward through the air. It is generated by accelerating a mass of air rearward, producing a reaction force in the forward direction per Newton's third law. In jet engines, thrust comes from compressing and heating air then expanding it through a nozzle. In propeller-driven aircraft, the rotating blades accelerate air to produce thrust more efficiently at lower speeds. Thrust must overcome drag for level flight and exceed weight for climb.}{catAD}
\glsadd{Thrust-Vectoring}\dictentry{Thrust Vectoring}{Thrust Vectoring is the ability to redirect the exhaust nozzle of a jet engine to produce a thrust component perpendicular to the aircraft's centerline. It augments or replaces conventional control surfaces for pitch, yaw, and roll. Thrust vectoring enhances maneuverability at low speeds and high angles of attack where aerodynamic surfaces lose effectiveness. It is employed on advanced fighters such as the F-22 and Su-57.}{catAD}
\glsadd{Trailing-Edge}\dictentry{Trailing Edge}{The Trailing Edge is the aft-most part of an airfoil where the upper and lower surfaces meet. Its geometry affects the pressure recovery, wake characteristics, and control surface hinge moments. A sharp trailing edge satisfies the Kutta condition in subsonic flow. Control surfaces such as ailerons, elevators, and flaps are typically hinged at the trailing edge.}{catAD}
\glsadd{Transonic-Flow}\dictentry{Transonic Flow}{Transonic Flow exists in the Mach number range from approximately 0.8 to 1.2, where regions of both subsonic and supersonic flow coexist on the same body. Shock waves form on the upper surface of the wing and move aft as speed increases. This regime is characterized by rapid changes in drag and pitching moment. Transonic aerodynamics presents some of the most complex design challenges in aviation.}{catAD}
\glsadd{Trim-Drag}\dictentry{Trim Drag}{Trim Drag is the incremental drag incurred when control surfaces are deflected to maintain a trimmed flight condition. It arises because deflected surfaces produce additional induced and pressure drag. Minimizing trim drag requires careful placement of the center of gravity relative to the aerodynamic center. Fuel management systems in modern airliners shift fuel to optimize CG position and reduce trim drag.}{catAD}
\glsadd{Turbofan}\dictentry{Turbofan}{A Turbofan diverts a portion of incoming air around the combustion core through a front fan. Bypass air provides thrust more efficiently, reducing fuel consumption and noise. High-bypass turbofans power most modern commercial airliners with bypass ratios exceeding 10:1. Turbofans dominate subsonic aviation due to superior fuel economy and reduced emissions.}{catAD}
\glsadd{Turbojet}\dictentry{Turbojet}{A Turbojet is a jet engine where all incoming air passes through the combustion core. Air is compressed, mixed with fuel, ignited, and expelled at high velocity for thrust. Turbojets were the first practical jet engines powering early military and commercial aircraft. They are efficient at high speeds but produce more noise and fuel consumption than turbofans.}{catAD}
\glsadd{Turbulent-Flow}\dictentry{Turbulent Flow}{Turbulent Flow is a regime of fluid motion characterized by chaotic, irregular fluctuations in velocity and pressure. Eddies of many sizes interact nonlinearly, transferring energy from large scales to small scales in a process known as the energy cascade. Turbulence increases mixing, heat transfer, and aerodynamic drag compared to laminar flow. It is the dominant flow regime for most practical aircraft and is notoriously difficult to predict from first principles.}{catAD}
\glsadd{Unsteady-Flow}\dictentry{Unsteady Flow}{Unsteady Flow describes any fluid motion where the velocity, pressure, or other flow properties at a given point change over time. Unsteadiness can arise from oscillating bodies, gusting freestream conditions, or transient phenomena such as starting vortices. Many real aerodynamic problems, including flutter and vortex shedding, are inherently unsteady. Solving unsteady flows demands time-accurate computation or measurement, making it more costly than steady analysis.}{catAD}
\glsadd{Upwash}\dictentry{Upwash}{Upwash is the upward component of airflow ahead of and above a lifting wing. The pressure field created by the wing deflects the incoming stream upward before it reaches the leading edge, effectively increasing the local angle of attack. Upwash influences the aerodynamic performance of surfaces located ahead of the wing, such as canards or propeller blades. It also plays a role in ground-effect vehicles, where the ground modifies the upwash pattern.}{catAD}
\glsadd{Vortex}\dictentry{Vortex}{A Vortex is a rotating region of fluid whose axis of rotation is either a line or a point. In aerodynamics, wingtip vortices form as high-pressure air beneath the wing curls around the tip toward the low-pressure region above. These vortices are a primary source of induced drag and persist in the atmosphere as trailing wake vortices. Understanding vortex behavior is essential for safe aircraft separation distances and for designing vortex-controlled devices.}{catAD}
\glsadd{Vortex-Generator}\dictentry{Vortex Generator}{A Vortex Generator is a small, angled vane mounted on a surface to energize the boundary layer. By producing streamwise vortices, it mixes high-energy freestream air into the slow-moving boundary layer, delaying flow separation. Vortex generators are widely used on wings, engine nacelles, and fuselages to improve performance at high angles of attack. They are simple, passive devices that require no moving parts or external power.}{catAD}
\glsadd{Wake}\dictentry{Wake}{The Wake is the region of disturbed, lower-energy flow trailing behind a body moving through a fluid. It contains deficit momentum, turbulence, and vorticity shed from the body surfaces. Wake properties directly affect the drag of the body and the aerodynamic loads on downstream structures. Aircraft wake turbulence, generated by large transport aircraft, is a critical safety concern for following aircraft during takeoff and landing.}{catAD}
\glsadd{Wall-Shear-Stress}\dictentry{Wall Shear Stress}{Wall Shear Stress is the tangential frictional force per unit area exerted by a flowing fluid on an adjacent solid surface. It arises from the velocity gradient within the boundary layer at the wall and is a direct measure of skin-friction drag. Accurate prediction of wall shear stress is essential for estimating total drag and for designing efficient aerodynamic surfaces. In boundary-layer control and transition research, wall shear stress measurements serve as key diagnostic quantities.}{catAD}
\glsadd{Washout}\dictentry{Washout}{Washout is an intentional aerodynamic twist that reduces the angle of attack at the wing tip relative to the root. This design feature ensures that the wing root stalls before the tip, preserving aileron effectiveness during the early stages of a stall. Washout improves lateral control and overall stall behavior, enhancing flight safety. It is achieved structurally during wing manufacture and is standard on most modern aircraft wings.}{catAD}
\glsadd{Wave-Drag}\dictentry{Wave Drag}{Wave Drag is the component of drag arising from the formation of shock waves at transonic and supersonic speeds. As local airflow over the wing exceeds the speed of sound, normal shocks form and cause abrupt pressure rises and boundary layer separation. Wave drag increases sharply near Mach 1 and is a primary obstacle to supersonic cruise efficiency. Design features such as wing sweep, area ruling, and supercritical airfoils reduce wave drag.}{catAD}
\glsadd{Wing}\dictentry{Wing}{The Wing is the primary aerodynamic lifting surface of an aircraft, generating the force that sustains flight. Its shape, planform, and airfoil sections are designed to produce lift efficiently while minimizing drag. Wings also serve as structural mounts for engines, landing gear, fuel tanks, and control surfaces. The design of the wing largely determines an aircraft's performance envelope, including speed, range, and maneuverability.}{catAD}
\glsadd{Wing-Area}\dictentry{Wing Area}{Wing Area is the planform surface area of the wing, a critical parameter in aerodynamic force calculations. Lift and drag are both proportional to wing area, so it directly influences wing loading and overall aircraft performance. The reference wing area typically includes the area of the fuselage portion that lies within the wing span. Correct definition of wing area is essential for consistent aerodynamic coefficient comparisons.}{catAD}
\glsadd{Wing-Span}\dictentry{Wing Span}{Wing Span is the straight-line distance from one wingtip to the other. It is a key geometric parameter that affects induced drag, structural weight, and roll performance. Longer spans reduce induced drag by distributing lift over a greater width, improving fuel efficiency. However, increased span adds structural weight and may be limited by hangar size and airport gate constraints.}{catAD}
\glsadd{Winglet}\dictentry{Winglet}{A Winglet is a vertical or near-vertical extension at the wingtip designed to reduce induced drag. By partially blocking the high-pressure airflow from curling around the tip, winglets weaken the wingtip vortex and improve the effective lift distribution. The result is lower fuel burn and extended range for the same fuel load. Modern winglet designs, including blended and split-tip variants, continue to evolve for incremental performance gains.}{catAD}
\glsadd{Wing-Root}\dictentry{Wing Root}{The Wing Root is the junction where the wing attaches to the fuselage. It experiences the highest bending moments during flight and must be structurally reinforced accordingly. Aerodynamically, the root region often has a thicker airfoil section to accommodate structural loads and internal systems. The wing-fuselage fairing at the root is shaped to minimize interference drag and smooth the airflow transition.}{catAD}
\glsadd{Wing-Tip}\dictentry{Wing Tip}{The Wing Tip is the outermost section of the wing, where the wingtip vortex originates. Tip design significantly affects induced drag and overall aerodynamic efficiency. Many modern aircraft employ winglets or raked tips to mitigate vortex strength and reduce drag. The tip is also a common location for navigation lights and static discharge wicks.}{catAD}
\glsadd{Yawing-Moment}\dictentry{Yawing Moment}{The Yawing Moment is the aerodynamic moment acting about the vertical axis of an aircraft, tending to rotate the nose left or right. It is generated by asymmetric lift or drag distributions, such as those caused by sideslip or rudder deflection. Yawing moment must be controlled to maintain directional stability and coordinated flight. Excessive or uncontrolled yawing can lead to dangerous conditions such as spin entry.}{catAD}
\glsadd{Adverse-Yaw}\dictentry{Adverse Yaw}{Adverse Yaw is the tendency of an aircraft to yaw in the direction opposite to the intended roll when ailerons are deflected. It occurs because the downgoing aileron generates more drag than the upgoing aileron, creating a net yawing moment. Adverse yaw is most pronounced at low speeds and high angles of attack. It is counteracted using differential aileron deflection, rudder input, or devices such as Frise ailerons.}{catAD}
\glsadd{Airbrake}\dictentry{Airbrake}{An Airbrake is a deployable device that increases aerodynamic drag to decelerate an aircraft during flight. Unlike spoilers, airbrakes are typically symmetrical panels that open into the airflow on both sides of the fuselage or wings. They are used during descent, approach, and landing rollout to control speed without reducing engine thrust. Some military aircraft use airbrakes to manage energy during combat maneuvers.}{catAD}
\glsadd{Boundary-Layer-Control}\dictentry{Boundary Layer Control}{Boundary Layer Control encompasses techniques to modify the state or behavior of the boundary layer on an aerodynamic surface. Methods include suction through porous surfaces, blowing high-energy air over the surface, and the use of vortex generators. The goal is to delay separation, reduce drag, or increase maximum lift. Effective boundary layer control enables higher angles of attack and improved aerodynamic efficiency.}{catAD}
\glsadd{Choked-Flow}\dictentry{Choked Flow}{Choked Flow occurs in a duct or nozzle when the flow velocity reaches the speed of sound at the throat, limiting the maximum mass flow rate. Further reduction of downstream pressure does not increase the mass flow through the choked section. This phenomenon is fundamental in the design of rocket nozzles, jet engine inlets, and supersonic wind tunnels. Understanding choked flow is critical for sizing propulsion system components.}{catAD}
\glsadd{Computational-Fluid-Dynamics}\dictentry{Computational Fluid Dynamics}{Computational Fluid Dynamics is a branch of fluid mechanics that uses numerical methods to simulate and analyze fluid flow, heat transfer, and associated phenomena. CFD solves the governing equations of fluid motion on discretized grids using algorithms such as finite volume, finite element, or finite difference methods. It enables engineers to predict aerodynamic performance before physical prototypes are built. Modern CFD tools can simulate complex three-dimensional flows including turbulence and chemical reactions.}{catAD}
\glsadd{Contra-Rotating-Propeller}\dictentry{Contra-Rotating Propeller}{A Contra-Rotating Propeller consists of two coaxial propellers rotating in opposite directions on the same shaft. The aft propeller recovers rotational energy lost by the forward propeller, increasing overall propulsive efficiency. This configuration also cancels torque effects, eliminating the need for a tail rotor on helicopters. Contra-rotating designs are used in high-performance turboprop aircraft and advanced marine propulsion systems.}{catAD}
\glsadd{Diffuser}\dictentry{Diffuser}{A Diffuser is a duct with an expanding cross-sectional area that decelerates fluid flow and converts kinetic energy into pressure rise. In jet engines, diffusers slow incoming air before the compressor, improving its efficiency. In supersonic inlets, diffuser geometry manages shock wave patterns to achieve efficient pressure recovery. Poorly designed diffusers can cause flow separation and pressure losses that degrade engine performance.}{catAD}
\glsadd{Displacement-Thickness}\dictentry{Displacement Thickness}{Displacement Thickness is the distance by which the external potential flow is displaced outward due to the presence of the boundary layer. It represents the equivalent thickness of a layer of inviscid fluid that would carry the same mass deficit as the real boundary layer. Displacement thickness is used to correct for boundary layer effects in inviscid flow calculations. It provides a useful measure of the overall influence of viscosity on the flow field.}{catAD}
\glsadd{Divergence}\dictentry{Divergence}{Divergence is an aeroelastic instability in which a structural deformation of the wing increases the aerodynamic load, which in turn increases the deformation in a self-reinforcing manner. It occurs when the dynamic pressure exceeds a critical value at which the aerodynamic stiffness overwhelms the structural stiffness. Divergence can lead to catastrophic structural failure if not prevented by adequate wing stiffness or design limits. It is a primary concern in the design of high-speed and lightweight aircraft.}{catAD}
\glsadd{Euler-Equations}\dictentry{Euler Equations}{The Euler Equations describe the motion of an inviscid, non-heat-conducting fluid. They represent conservation of mass, momentum, and energy and form the basis for many aerodynamic flow models. When viscosity is negligible, such as in parts of the flow far from surfaces, the Euler equations provide accurate predictions. They are simpler to solve than the full Navier-Stokes equations but cannot predict boundary layers or skin friction.}{catAD}
\glsadd{Finite-Element-Method}\dictentry{Finite Element Method}{The Finite Element Method is a numerical technique for obtaining approximate solutions to boundary value problems governed by partial differential equations. The domain is divided into small elements, and the solution is approximated within each element using simple shape functions. FEM is widely used in structural analysis, heat transfer, and fluid mechanics. In aerodynamics, it is particularly valuable for solving complex structural and coupled fluid-structure interaction problems.}{catAD}
\glsadd{Finite-Volume-Method}\dictentry{Finite Volume Method}{The Finite Volume Method solves partial differential equations by discretizing the domain into control volumes and integrating governing equations over each. This ensures fluxes of conserved quantities are exactly balanced across cell interfaces, making FVM well suited for computational fluid dynamics where conservation is paramount. FVM handles complex geometries with unstructured meshes and is the basis of most commercial and research CFD codes.}{catAD}
\glsadd{Force-Balance}\dictentry{Force Balance}{A Force Balance is an instrumented device that measures the aerodynamic forces and moments acting on a wind tunnel model. It typically uses strain gauges or piezoelectric sensors to resolve forces into lift, drag, and side force components, along with pitching, rolling, and yawing moments. Force balance data provide the direct aerodynamic coefficients needed for aircraft design and performance prediction. Precision calibration and model alignment are essential for accurate measurements.}{catAD}
\glsadd{Grid}\dictentry{Grid}{A Grid, in computational fluid dynamics, is the discretized representation of the fluid domain into cells or elements where the governing equations are solved. Grid quality, including cell size, shape, and alignment with flow features, directly affects solution accuracy and computational cost. Structured grids follow regular connectivity patterns, while unstructured grids offer greater geometric flexibility. Adaptive grid refinement allows resolution to be concentrated in regions of high gradients such as shock waves.}{catAD}
\glsadd{LES}\dictentry{LES}{Large Eddy Simulation is a turbulence modeling approach that directly resolves large-scale turbulent structures while modeling the effects of smaller, more universal eddies. LES provides greater fidelity than Reynolds-Averaged methods for unsteady flows, capturing phenomena such as vortex shedding and acoustic noise generation. It requires finer grids and smaller time steps than RANS, making it computationally expensive. LES is increasingly used for complex aerodynamic and propulsion applications.}{catAD}
\glsadd{Mach-Number}\dictentry{Mach Number}{The Mach Number is the ratio of a flow's velocity to the local speed of sound. It is the fundamental dimensionless parameter governing compressibility effects in aerodynamics. Flows below Mach 0.3 are generally treated as incompressible, while transonic, supersonic, and hypersonic regimes begin at approximately Mach 0.8, Mach 1.2, and Mach 5 respectively. The Mach number determines the nature of shock waves and expansion fans around a body.}{catAD}
\glsadd{Mesh}\dictentry{Mesh}{A Mesh is the collection of cells, elements, or nodes that form the discretized representation of a fluid volume in computational simulation. Mesh generation is a critical step in CFD, as mesh quality influences accuracy, convergence, and computational expense. Meshes may be structured, unstructured, or hybrid depending on the geometry and solver requirements. Boundary layer meshes with highly stretched cells near walls are essential for resolving viscous effects.}{catAD}
\glsadd{Particle-Image-Velocimetry}\dictentry{Particle Image Velocimetry}{Particle Image Velocimetry is an optical measurement technique that captures instantaneous velocity fields in a fluid flow. Seeding particles are illuminated by a laser sheet, and their positions are recorded by a camera at two closely spaced times. Cross-correlation analysis of image pairs yields the displacement and thus the velocity of the fluid. PIV provides full-field, non-intrusive measurements that are invaluable for validating computational fluid dynamics results.}{catAD}
\glsadd{Ramjet}\dictentry{Ramjet}{A Ramjet is an air-breathing jet engine that uses the forward motion of the vehicle to compress incoming air without a mechanical compressor. Air is decelerated in an inlet diffuser, mixed with fuel, and burned in a combustion chamber. The resulting high-velocity exhaust produces thrust. Ramjets are efficient at supersonic speeds but cannot start from rest, requiring a booster to reach their operational velocity.}{catAD}
\glsadd{RANS}\dictentry{RANS}{Reynolds-Averaged Navier-Stokes is a turbulence modeling approach that decomposes the flow into mean and fluctuating components and models the effect of turbulence on the mean flow through Reynolds stresses. RANS is the most widely used method in industrial CFD due to its computational efficiency. Various turbulence models, such as k-epsilon and k-omega, close the RANS equations by relating Reynolds stresses to mean flow quantities. RANS is well suited for steady, attached flows but less accurate for massively separated or unsteady flows.}{catAD}
\glsadd{DNS}\dictentry{DNS}{Direct Numerical Simulation solves the full Navier-Stokes equations without any turbulence model, resolving all scales of motion down to the smallest dissipative eddies. DNS provides the most accurate representation of turbulent flow but requires extremely fine grids and small time steps, limiting it to low Reynolds numbers and simple geometries. It serves as a benchmark for validating turbulence models and as a research tool for studying fundamental turbulence physics. Advances in supercomputing continue to extend DNS accessibility.}{catAD}
\glsadd{Scramjet}\dictentry{Scramjet}{A Scramjet is a supersonic combustion ramjet in which the airflow through the engine remains supersonic throughout. Unlike a conventional ramjet, which decelerates the flow to subsonic speeds before combustion, a scramjet burns fuel in a supersonic stream. This eliminates the total pressure losses associated with deceleration and enables efficient propulsion at hypersonic speeds. Scramjet technology is being developed for hypersonic vehicles capable of sustained flight above Mach 5.}{catAD}
\glsadd{Smoke-Visualization}\dictentry{Smoke Visualization}{Smoke Visualization is an experimental technique in which smoke is introduced into a airflow to reveal flow patterns around a model. The smoke streaklines show the paths of fluid particles, making features such as separation, vortices, and transition visible. It is a simple, qualitative method widely used in low-speed wind tunnel testing and for educational demonstrations. Modern variants use laser sheet illumination to create planar visualizations of three-dimensional flow structures.}{catAD}
\glsadd{Static-Test}\dictentry{Static Test}{A Static Test is a structural test in which predetermined loads are applied to an aircraft component or structure without dynamic excitation. The purpose is to verify that the structure can withstand design limit and ultimate loads without failure or excessive deformation. Static tests are performed on complete airframes and individual components such as wing boxes and fuselage sections. Results validate structural analysis methods and confirm compliance with airworthiness standards.}{catAD}
\glsadd{Test-Section}\dictentry{Test Section}{The Test Section is the portion of a wind tunnel where the model is placed and exposed to a controlled, uniform airflow. It is designed to produce flow of known velocity, turbulence level, and temperature with minimal spatial variation. Test section geometry, including shape and size, affects wall interference and blockage effects that must be accounted for in data reduction. Modern test sections may include optical access windows for flow visualization and laser-based measurements.}{catAD}
\glsadd{Transitional-Flow}\dictentry{Transitional Flow}{Transitional Flow is the regime between fully laminar and fully turbulent flow in which the boundary layer undergoes a gradual change from orderly to chaotic motion. Transition is influenced by surface roughness, pressure gradient, freestream turbulence, and Reynolds number. Predicting the transition location is critical because it determines the extent of laminar and turbulent boundary layers, which have very different drag and heat transfer characteristics. Laminar flow control technologies aim to delay transition for drag reduction.}{catAD}
\glsadd{Turbulence-Model}\dictentry{Turbulence Model}{A Turbulence Model is a set of additional equations used in computational fluid dynamics to account for the effects of unresolved turbulent motions on the mean flow. Common models include the k-epsilon, k-omega SST, and Spalart-Allmaras families, each with different strengths and limitations. The choice of turbulence model significantly affects predicted flow separation, mixing, and heat transfer. No single model performs well for all flow conditions, so model selection requires engineering judgment.}{catAD}
\glsadd{Vortex-Shedding}\dictentry{Vortex Shedding}{Vortex Shedding is the periodic formation and detachment of alternating vortices from the sides of a bluff body in a flow. The resulting von Karman vortex street produces oscillating lift and drag forces that can induce structural vibrations. Vortex shedding frequency is characterized by the Strouhal number, which depends on body geometry and Reynolds number. Excessive vortex-induced vibrations have caused structural failures in bridges, chimneys, and offshore structures.}{catAD}
\glsadd{Wind-Tunnel}\dictentry{Wind Tunnel}{A Wind Tunnel is an experimental facility that generates a controlled stream of air flowing past a stationary model to measure aerodynamic forces, pressures, and flow patterns. Wind tunnels range from small subsonic units for education to large transonic and supersonic facilities for full-scale aircraft testing. They remain indispensable for validating computational predictions and for exploring flow phenomena that are difficult to simulate. Data from wind tunnel tests directly inform aircraft design and certification.}{catAD}
\glsadd{Aerothermal-Heating}\dictentry{Aerothermal Heating}{Aerothermal Heating is the thermal energy transferred to a body surface by the combined effects of convective heat transfer and compressive heating of the air in high-speed flow. At hypersonic velocities, the stagnation temperature of the flow can exceed several thousand degrees, causing severe heating that must be managed by thermal protection systems. Aerothermal heating analysis couples fluid dynamics with heat transfer to predict surface temperatures. It is a critical design driver for re-entry vehicles and hypersonic aircraft.}{catAD}
\glsadd{Ablation}\dictentry{Ablation}{Ablation is the removal of surface material through vaporization, melting, or chemical erosion as a mechanism for protecting a structure from extreme heat. Ablative thermal protection systems absorb heat by undergoing phase change and carrying energy away with the removed material. This approach is widely used on re-entry capsules, rocket nozzles, and leading edges exposed to hypersonic flow. Ablation rates and effectiveness depend on material properties and the severity of the thermal environment.}{catAD}
\glsadd{Re-entry}\dictentry{Re-entry}{Re-entry is the phase of flight in which a spacecraft returns from space and enters the Earth's atmosphere. The vehicle encounters rapidly increasing air density at hypersonic speeds, generating intense aerodynamic heating and deceleration forces. Heat shields and thermal protection systems are essential to prevent structural failure during this phase. Proper trajectory management is critical to limit peak heating rates and ensure the vehicle reaches its intended landing site safely.}{catAD}
\glsadd{Shock-Layer}\dictentry{Shock Layer}{The Shock Layer is the region of compressed, heated gas between the bow shock wave and the surface of a body traveling at supersonic or hypersonic speed. The shock wave drastically increases pressure, temperature, and density of the air, creating a hostile thermal environment for the vehicle. Shock layer properties vary strongly with Mach number, vehicle geometry, and altitude. Accurate prediction of shock layer conditions is essential for designing thermal protection systems and estimating aerodynamic loads.}{catAD}
\glsadd{Thermal-Protection-System}\dictentry{Thermal Protection System}{A Thermal Protection System is the ensemble of materials and structural components designed to protect a spacecraft from extreme heat during atmospheric entry or high-speed flight. TPS technologies include ablative materials, reusable ceramic tiles, and metallic heat shields, each suited to different thermal environments. The TPS must withstand temperatures from hundreds to thousands of degrees while maintaining structural integrity. It is one of the most safety-critical systems on any vehicle experiencing aerothermal heating.}{catAD}
\glsadd{Pulsed-Detonation-Engine}\dictentry{Pulsed Detonation Engine}{A Pulsed Detonation Engine is a propulsion device that generates thrust through repeated detonation waves rather than conventional deflagration. Each detonation cycle compresses and heats the fuel-air mixture almost instantaneously, achieving higher thermodynamic efficiency than a Brayton cycle engine. PDEs have the potential to simplify engine construction by eliminating the need for complex turbomachinery. Although still largely experimental, PDE technology promises significant performance gains for future aerospace propulsion systems.}{catAD}
\glsadd{Turbine}\dictentry{Turbine}{A Turbine is a rotary mechanical device that extracts energy from a moving fluid and converts it into useful work. In aerospace, turbines are central to jet engines, where they extract energy from hot combustion gases to drive the compressor and accessories. Turbine blade design must balance aerodynamic efficiency with extreme thermal and centrifugal loading. Advances in blade cooling and high-temperature materials continue to raise turbine inlet temperatures, improving engine cycle efficiency.}{catAD}
\glsadd{Compressor}\dictentry{Compressor}{A Compressor is a device that increases the pressure of a gas by doing work on it, typically through rotating blades. In gas turbine engines, the compressor draws in ambient air and raises its pressure before combustion, a critical step in the thermodynamic cycle. Axial compressors achieve high flow rates with multiple blade stages, while centrifugal compressors offer simplicity and compactness. Compressor efficiency directly affects overall engine fuel consumption and thrust output.}{catAD}
\glsadd{Spoileron}\dictentry{Spoileron}{A Spoileron is a spoiler surface deployed differentially to produce a rolling moment while simultaneously increasing drag. Unlike conventional ailerons, spoilerons primarily reduce lift on one wing rather than increasing it on the other. They are commonly used on aircraft where wing flexibility or supersonic flight limits conventional aileron effectiveness. Spoilerons also serve as speed brakes when deployed symmetrically.}{catAD}
\glsadd{Strake}\dictentry{Strake}{A Strake is a narrow, elongated aerodynamic surface extending from the fuselage or wing root, typically at a high angle of attack. Strakes generate powerful vortices that energize the airflow over the main wing, delaying stall and improving high-angle-of-attack handling. They are commonly found on fighter aircraft and high-performance designs where aggressive maneuvering is required. Strakes also contribute to directional stability and can house sensors or antennas.}{catAD}
\glsadd{Canard}\dictentry{Canard}{A Canard is a small forewing positioned ahead of the main wing on an aircraft. It provides pitch control and, in some configurations, contributes to lift. Because the canard operates in clean air before the main wing, it can provide effective control authority at high angles of attack. Canard configurations offer benefits in stall characteristics and short-field performance but require careful design to manage pitch stability.}{catAD}
\glsadd{Blended-Wing-Body}\dictentry{Blended Wing Body}{A Blended Wing Body is an aircraft configuration in which the wing and fuselage are smoothly integrated into a single lifting surface. This design reduces wetted area and interference drag compared to conventional tube-and-wing configurations, offering significant fuel efficiency improvements. The wide interior can accommodate passengers, cargo, or fuel in a non-cylindrical pressurized structure. BWB designs present challenges in structural pressurization, emergency evacuation, and airport compatibility.}{catAD}
\glsadd{Flying-Wing}\dictentry{Flying Wing}{A Flying Wing is an aircraft configuration in which the entire structure serves as the lifting surface, with no distinct fuselage or tail. This arrangement minimizes wetted area and parasite drag, offering the potential for very high aerodynamic efficiency. Control and stability are achieved through split drag rudders, elevons, and wing shaping. Flying wings present challenges in directional stability, payload volume, and pitch control that require innovative solutions.}{catAD}
\glsadd{Delta-Wing}\dictentry{Delta Wing}{A Delta Wing is a triangular planform wing with a highly swept leading edge and a straight or nearly straight trailing edge. The delta configuration offers a large internal volume for fuel and structure, high structural strength, and favorable characteristics at supersonic speeds. At low speeds, delta wings generate lift through leading-edge vortices rather than conventional attached flow. The trade-off is high induced drag and poor low-speed performance without additional surfaces such as canards.}{catAD}
\glsadd{Swept-Wing}\dictentry{Swept Wing}{A Swept Wing is a wing angled rearward from the root, reducing the effective component of freestream velocity perpendicular to the leading edge. This delays the onset of compressibility effects and raises the critical Mach number, enabling efficient transonic and supersonic flight. Swept wings are standard on most modern commercial and military aircraft. The downside is a tendency for tip stall and reduced low-speed lift, which must be managed through design features like washout and leading-edge devices.}{catAD}
\glsadd{Forward-Swept-Wing}\dictentry{Forward Swept Wing}{A Forward Swept Wing is a wing angled forward from the root, offering several aerodynamic advantages over conventional aft-swept designs. Inbound flow along the span migrates toward the root, preventing tip stall and preserving aileron effectiveness at high angles of attack. Forward-swept wings also allow a thinner, more efficient airfoil and lower trim drag. The configuration requires careful structural design to resist aeroelastic divergence, which is more severe than for aft-swept wings.}{catAD}
\glsadd{Ogive}\dictentry{Ogive}{An Ogive is a pointed, curved shape formed by the intersection of circular arcs, widely used for nose cones and projectile tips. The smooth curvature of an ogive minimizes wave drag at supersonic speeds compared to simpler conical shapes. Tangent ogives blend smoothly into the cylindrical body, while secant ogives have a discontinuity at the junction. Ogive geometry is a standard consideration in the design of missiles, rockets, and high-speed aircraft noses.}{catAD}
\glsadd{Streamline}\dictentry{Streamline}{A Streamline is a curve that is everywhere tangent to the instantaneous velocity vector of the flow. In steady flow, streamlines represent the actual paths followed by fluid particles. They provide a powerful visual tool for understanding flow patterns, with closely spaced streamlines indicating high velocity and widely spaced streamlines indicating low velocity. Streamlines cannot cross each other in steady flow, and their shape reveals regions of acceleration, deceleration, and separation.}{catAD}
\glsadd{Pathline}\dictentry{Pathline}{A Pathline is the actual trajectory traced by an individual fluid particle over a period of time. In steady flow, pathlines coincide with streamlines, but in unsteady flow they diverge because the velocity field changes as the particle moves. Pathlines are recorded experimentally using particle tracking methods such as particle image velocimetry. Understanding pathlines is essential for analyzing mixing, pollutant transport, and unsteady aerodynamic phenomena.}{catAD}
\glsadd{Streakline}\dictentry{Streakline}{A Streakline is the locus of all fluid particles that have passed through a fixed point in space at various earlier times. In steady flow, streaklines, streamlines, and pathlines are identical, but they differ in unsteady flow. Streaklines are the most commonly visualized flow feature in experiments using smoke or dye injection. They reveal the instantaneous shape of the flow field as seen from the injection point.}{catAD}
\glsadd{Continuity-Equation}\dictentry{Continuity Equation}{The Continuity Equation expresses the principle of conservation of mass in fluid flow. It states that the rate of mass entering a control volume must equal the rate of mass leaving plus any accumulation within the volume. For an incompressible fluid, the continuity equation reduces to the requirement that the velocity field is divergence-free. It is one of the fundamental governing equations used in all branches of fluid mechanics and aerodynamics.}{catAD}
\glsadd{Momentum-Equation}\dictentry{Momentum Equation}{The Momentum Equation applies Newton's second law to fluid motion, relating the forces acting on a fluid element to its acceleration. In integral form, it accounts for pressure forces, viscous stresses, body forces, and momentum flux through control surfaces. The momentum equation is one of the three governing equations of fluid dynamics, alongside continuity and energy. It forms the basis for analyzing aerodynamic forces on bodies and for computing flow fields in computational fluid dynamics.}{catAD}
\glsadd{Energy-Equation}\dictentry{Energy Equation}{The Energy Equation expresses the conservation of total energy in a fluid flow, accounting for internal energy, kinetic energy, heat transfer, and work done by surface and body forces. In compressible flow, the energy equation is essential for determining temperature changes associated with compression and expansion. Combined with the continuity and momentum equations, it forms a complete system for solving thermally active flows. It is particularly important in high-speed aerodynamics where compressibility and heat transfer are significant.}{catAD}
\glsadd{Bernoulli-Equation}\dictentry{Bernoulli Equation}{The Bernoulli Equation states that the sum of static pressure, dynamic pressure, and hydrostatic pressure is constant along a streamline in steady, inviscid flow. It is a simplified form of the momentum equation applicable when viscous effects and heat transfer are negligible. Bernoulli's principle explains the fundamental relationship between flow velocity and pressure that underpins lift generation. While limited to incompressible, irrotational flow in its simplest form, it remains a foundational tool in aerodynamic analysis.}{catAD}
\glsadd{Camber-Line}\dictentry{Camber Line}{The Camber Line is the locus of points equidistant between the upper and lower surfaces of an airfoil, running from leading edge to trailing edge. It defines the mean curvature of the airfoil and is a primary determinant of the airfoil's zero-lift angle of attack and maximum lift coefficient. A symmetric airfoil has a straight camber line, while a cambered airfoil has a curved one. Camber is a fundamental design parameter for tailoring airfoil performance to specific flight conditions.}{catAD}
\glsadd{Mean-Line}\dictentry{Mean Line}{Mean Line is an alternative term for the camber line of an airfoil, describing the curve midway between the upper and lower surfaces. In thin airfoil theory, the mean line is used to calculate lift and pitching moment distributions independently of thickness. Aerodynamic design often focuses on optimizing the mean line shape to achieve desired lift characteristics while thickness is chosen for structural requirements. The mean line concept simplifies airfoil analysis by separating camber effects from thickness effects.}{catAD}
\glsadd{Pressure-Coefficient}\dictentry{Pressure Coefficient}{The Pressure Coefficient is a dimensionless number that normalizes the local static pressure difference from freestream by the freestream dynamic pressure. A value of one indicates a stagnation point, while negative values indicate regions where local pressure is below freestream. Pressure coefficient distributions over airfoil and body surfaces are fundamental data for aerodynamic analysis and design. They reveal the loading distribution and are used to compute lift, drag, and moment coefficients.}{catAD}
\glsadd{Stagnation-Point}\dictentry{Stagnation Point}{The Stagnation Point is the location on a body surface where the local flow velocity is zero and the pressure reaches its maximum value. At this point, all kinetic energy of the incoming flow has been converted to pressure energy. The stagnation point location varies with angle of attack and body geometry, and its movement influences surface pressure distributions. Stagnation point heating is also the most severe thermal environment on a hypersonic vehicle.}{catAD}
\glsadd{Stagnation-Pressure}\dictentry{Stagnation Pressure}{Stagnation Pressure, also known as total pressure, is the pressure that would be measured at a point where the flow is brought to rest isentropically. It represents the total energy content of the flow per unit volume. Stagnation pressure decreases across shock waves and in the presence of viscous losses, making it a useful indicator of flow quality and efficiency. In propulsion systems, stagnation pressure recovery in the inlet is a key performance metric.}{catAD}
\glsadd{Static-Pressure}\dictentry{Static Pressure}{Static Pressure is the thermodynamic pressure of the fluid as measured by an instrument moving with the flow. It is the pressure that would be sensed by a pressure tap flush with a surface parallel to the flow. Static pressure decreases as flow velocity increases, as described by Bernoulli's equation in incompressible flow. It is one of the fundamental properties used to characterize a fluid flow, alongside dynamic and stagnation pressure.}{catAD}
\glsadd{Total-Pressure}\dictentry{Total Pressure}{Total Pressure is the sum of the static pressure and dynamic pressure in a flow. In incompressible flow, it is constant along a streamline in the absence of viscous losses. Total pressure represents the complete pressure potential of the flow and is a direct measure of the mechanical energy available. Loss of total pressure across shocks, diffusers, or viscous regions indicates irreversible energy dissipation that degrades performance.}{catAD}
\glsadd{Aerodynamic-Twist}\dictentry{Aerodynamic Twist}{Aerodynamic Twist is the variation in airfoil section properties, typically camber or zero-lift angle, along the wingspan. It affects the local lift characteristics at each spanwise station and is used to tailor the spanwise lift distribution. Aerodynamic twist works in conjunction with geometric twist to optimize performance and stall behavior. It is an important consideration in wing design, particularly for achieving near-elliptical lift distributions.}{catAD}
\glsadd{Geometric-Twist}\dictentry{Geometric Twist}{Geometric Twist is the physical change in the angle of incidence of the wing sections along the span. It is achieved by structurally rotating the wing so that the angle of attack varies from root to tip. Geometric twist is the primary mechanism for controlling stall progression and optimizing the spanwise lift distribution. Most aircraft wings incorporate some degree of geometric twist, usually washout, to improve stall characteristics and reduce induced drag.}{catAD}
\glsadd{Wash-In}\dictentry{Wash-In}{Wash-In is a wing twist that increases the angle of attack at the wing tip relative to the root. This configuration causes the tip to generate more lift per unit area than the root, which can be used to tailor the spanwise load distribution. Wash-in is the opposite of the more common washout and is less frequently used because it promotes tip-first stall. It may be employed in specific designs where tip loading or control characteristics are desired.}{catAD}
\glsadd{Oswald-Span-Efficiency}\dictentry{Oswald Span Efficiency}{The Oswald Span Efficiency Factor is a dimensionless measure of how closely the actual spanwise lift distribution approaches the ideal elliptical distribution that minimizes induced drag. A value of one indicates a perfectly elliptical distribution with minimum induced drag, while lower values indicate less efficient distributions. The factor typically ranges from 0.7 to 0.95 for practical aircraft configurations. It is used in preliminary design to estimate induced drag from the lift coefficient, aspect ratio, and wing planform.}{catAD}
\glsadd{Elliptical-Wing}\dictentry{Elliptical Wing}{An Elliptical Wing has a planform shape in which the local chord varies elliptically along the span, producing a uniform downwash distribution. This distribution minimizes induced drag for a given lift and span, representing the theoretically optimal spanwise loading. The elliptical wing was famously used on the Supermarine Spitfire. In practice, near-elliptical lift distributions can be achieved with tapered wings and appropriate twist, avoiding the structural and manufacturing complexity of a true elliptical planform.}{catAD}
\glsadd{Taper-Ratio}\dictentry{Taper Ratio}{The Taper Ratio is the ratio of the wing tip chord to the root chord. It influences the spanwise lift distribution, structural weight, and stall characteristics of the wing. A taper ratio of one indicates a rectangular wing, while zero represents a pure triangular planform. Typical values range from 0.2 to 0.5, and the selection involves trade-offs between aerodynamic efficiency, structural efficiency, and manufacturing considerations.}{catAD}
\glsadd{Wing-Loading}\dictentry{Wing Loading}{Wing Loading is the ratio of aircraft weight to wing area, expressed in units of force per unit area. It is a fundamental parameter that influences stall speed, climb rate, maneuver performance, and turbulence response. Higher wing loading results in higher cruise speeds and smoother ride in turbulence but requires higher approach and takeoff speeds. Wing loading is a key parameter in aircraft sizing and is used to compare aircraft of different sizes and types.}{catAD}
\glsadd{Load-Factor}\dictentry{Load Factor}{Load Factor is the ratio of the aerodynamic lift acting on an aircraft to its weight, expressed in units of gravitational acceleration g. In straight-and-level flight the load factor is one, while in a banked turn it increases to sustain the centripetal acceleration. Structural strength requirements and human tolerance limits define the maximum allowable load factor for each aircraft category. Load factor is a central concept in structural design and flight envelope definition.}{catAD}
\glsadd{V-n-Diagram}\dictentry{V-n Diagram}{The V-n Diagram is a graph of airspeed versus load factor that defines the complete flight envelope of an aircraft. It shows the structural limits of the airframe, including maximum positive and negative load factors, stall boundaries, and maneuver limits. Every point within the diagram represents a combination of speed and load factor that the aircraft can safely sustain. The V-n diagram is a fundamental document in aircraft structural design and certification.}{catAD}
\glsadd{Maneuverability}\dictentry{Maneuverability}{Maneuverability is the ability of an aircraft to change its flight path, attitude, and velocity rapidly and precisely. It is determined by the aircraft's thrust-to-weight ratio, wing loading, control surface effectiveness, and structural load factor limits. High maneuverability is essential for military fighter aircraft engaged in air combat. Trade-offs exist between maneuverability and other performance parameters such as range, payload, and stealth.}{catAD}
\glsadd{Agility}\dictentry{Agility}{Agility is the measure of an aircraft's capability to rapidly change its attitude and direction in three dimensions. While maneuverability concerns the ability to sustain turns and load factors, agility emphasizes the speed of response and the rate of attitude change. Agility depends on control power, rotational inertia, and the responsiveness of the flight control system. It is a critical performance metric for modern fighter aircraft engaged in close-range combat.}{catAD}
\glsadd{Turning-Radius}\dictentry{Turning Radius}{Turning Radius is the radius of the circular path traced by an aircraft during a level, constant-speed turn. It depends on true airspeed and bank angle, with tighter turns achievable at lower speeds or higher bank angles. The minimum turning radius is limited by structural load factor, stall speed, and engine performance. Pilots and mission planners use turning radius to assess terrain clearance and tactical maneuvering capability.}{catAD}
\glsadd{Rate-of-Turn}\dictentry{Rate of Turn}{Rate of Turn is the angular velocity at which an aircraft changes its heading during a turn, measured in degrees per second. It is a function of true airspeed, bank angle, and load factor. Maximum rate of turn occurs at minimum airspeed and maximum allowable bank angle. Rate of turn is an important performance parameter for tactical aircraft, air traffic control procedures, and flight training.}{catAD}
\glsadd{Coordinated-Turn}\dictentry{Coordinated Turn}{A Coordinated Turn is a maneuver in which the aircraft turns with no sideslip, meaning the nose tracks smoothly along the curved flight path without yawing relative to the relative wind. Aerodynamic forces are balanced so that lift acts perpendicular to the flight path while centripetal acceleration is provided by the horizontal component of lift. Coordination requires proper use of ailerons and rudder to maintain zero sideslip. Coordinated turns are more efficient and comfortable than uncoordinated ones.}{catAD}
\glsadd{Sideslip}\dictentry{Sideslip}{Sideslip is the angle between the aircraft's longitudinal axis and the direction of the relative wind in the horizontal plane. It is caused by uncoordinated rudder input, crosswind, or asymmetric thrust. Sideslip generates a lateral force on the fuselage and vertical tail, which provides directional stability. Controlled sideslip is used in crosswind landings and single-engine operations, while uncontrolled sideslip can lead to adverse handling qualities.}{catAD}
\glsadd{Ground-Effect}\dictentry{Ground Effect}{Ground Effect is the reduction in induced drag that occurs when an aircraft flies close to a surface, typically within one wingspan of the ground. The ground interferes with the development of wingtip vortices, reducing the downwash and the associated induced drag. Ground effect enables shorter takeoff and landing distances and can significantly increase payload capability for low-altitude flight. ekranoplans, or ground-effect vehicles, exploit this phenomenon for high-speed transport over water.}{catAD}
\glsadd{Reynolds-Analogy}\dictentry{Reynolds Analogy}{The Reynolds Analogy is a relationship that connects skin friction drag to convective heat transfer in a boundary layer. It states that the Stanton number is proportional to the friction coefficient, implying that momentum and heat transfer occur at similar rates. The analogy is most accurate for gases with a Prandtl number near one and for turbulent boundary layers. It provides a useful estimate of heat transfer rates when only aerodynamic data are available.}{catAD}
\glsadd{Prandtl-Number}\dictentry{Prandtl Number}{The Prandtl Number is a dimensionless parameter that represents the ratio of momentum diffusivity to thermal diffusivity in a fluid. It governs the relative thickness of the velocity and thermal boundary layers. For air, the Prandtl number is approximately 0.71, indicating that momentum diffuses slightly faster than heat. It is a key parameter in convective heat transfer correlations and in the analysis of aerothermal heating environments.}{catAD}
\glsadd{Mach-Tuck}\dictentry{Mach Tuck}{Mach Tuck is a nose-down pitching moment that develops as an aircraft enters the transonic speed range. It is caused by the aftward movement of the shock wave and the resulting rearward shift of the center of pressure on the wing. Mach tuck can lead to a rapid, uncontrollable dive if not counteracted by trim systems or automatic pitch control. Understanding and managing mach tuck is essential for safe transonic and supersonic flight operations.}{catAD}
\glsadd{Shock-Stall}\dictentry{Shock Stall}{Shock Stall is a loss of lift and increase in drag caused by shock-induced boundary layer separation on a wing at transonic speeds. As the local supersonic flow terminates in a normal shock, the resulting adverse pressure gradient causes the boundary layer to separate downstream of the shock. This separation destroys the pressure recovery on the upper surface and can lead to a sudden loss of lift. Shock stall is a critical consideration in the design of transonic airfoils and wings.}{catAD}
\glsadd{Buffeting}\dictentry{Buffeting}{Buffeting is the structural vibration caused by unsteady, separated airflow impinging on aircraft surfaces. It is commonly experienced near stall, in transonic flight, or behind large separated flow regions. Buffeting serves as a tactile warning to pilots of impending aerodynamic limits such as stall or shock-induced separation. Persistent buffeting can cause structural fatigue and must be considered in airframe design and operational limits.}{catAD}
\glsadd{Flutter}\dictentry{Flutter}{Flutter is a dynamic aeroelastic instability in which aerodynamic, inertial, and elastic forces interact to produce self-excited oscillations of a structure. Once the flutter speed is exceeded, oscillations grow rapidly and can lead to catastrophic structural failure within seconds. Flutter analysis is a mandatory part of aircraft certification, and flight testing must demonstrate adequate margin above the design flutter speed. The phenomenon involves complex coupling between multiple structural modes and aerodynamic forces.}{catAD}
\glsadd{Control-Reversal}\dictentry{Control Reversal}{Control Reversal is a phenomenon in which the aerodynamic response to a control surface deflection is opposite to the intended effect, occurring at high speeds due to aeroelastic deformation of the wing. As the wing twists in response to the control surface load, the change in local angle of attack can overwhelm the direct control effect. Control reversal sets an effective upper speed limit for conventional control systems. Composite wing structures and active control systems are used to raise or eliminate the reversal speed.}{catAD}
\glsadd{Gust-Load}\dictentry{Gust Load}{Gust Load is the transient aerodynamic force induced on an aircraft by a sudden change in the atmospheric velocity field, such as an updraft or downdraft. The severity of gust loads depends on gust velocity, aircraft speed, weight, and wing characteristics. Aircraft structures must be designed to withstand the most severe gust loads expected in service without damage. Gust load alleviation systems use sensors and control surfaces to detect and counteract gust-induced loads in real time.}{catAD}
\glsadd{Maneuvering-Load}\dictentry{Maneuvering Load}{Maneuvering load is additional aerodynamic load caused by pilot- or autopilot-induced maneuvers. These loads arise because the aircraft must accelerate beyond straight-and-level flight, increasing load factor. The magnitude depends on maneuver severity, speed, and configuration. Designers must ensure the airframe withstands limit and ultimate loads per certification standards.}{catAD}
\glsadd{Lift-Distribution}\dictentry{Lift Distribution}{Lift distribution describes how lift varies along a wing span in flight. The ideal elliptical distribution minimizes induced drag, as predicted by Prandtl's lifting line theory. In practice, wing taper, twist, and sweep produce deviations. The distribution determines trailing vortex strength and location. Optimizing it is critical for efficient wing design.}{catAD}
\glsadd{Trailing-Vortex}\dictentry{Trailing Vortex}{A trailing vortex is a rotating tube of air forming at a lifting wingtip and extending downstream in the wake. Higher pressure beneath the wing drives flow around the tip to the lower-pressure upper surface, creating spiraling motion. Vortex strength is proportional to lift and inversely proportional to speed and wingspan. Trailing vortices are the primary source of induced drag and can hazard following aircraft.}{catAD}
\glsadd{Starting-Vortex}\dictentry{Starting Vortex}{The starting vortex is a counter-clockwise vortex shed from the trailing edge of an airfoil when it begins moving or abruptly changes angle of attack. By the Kelvin circulation theorem, the sudden generation of bound circulation is balanced by shedding a vortex of opposite sign. This vortex moves downstream and becomes part of the wake. The starting vortex explains how circulation develops on a wing to produce lift.}{catAD}
\glsadd{Horseshoe-Vortex}\dictentry{Horseshoe Vortex}{The horseshoe vortex models the vortex system around a finite wing as a bound vortex along the span with two trailing vortices from the wingtips. This satisfies the Helmholtz theorems by forming a continuous vortex line in the fluid. The bound vortex represents lift-generating circulation, while trailing vortices account for finite span effects. The model is foundational in lifting line theory and explains induced drag generation.}{catAD}
\glsadd{Panel-Method}\dictentry{Panel Method}{A panel method computes potential flow around complex bodies by representing surfaces as panels carrying singularities. Strengths are set by enforcing flow-tangency at control points. The method yields velocity and pressure over the surface. It is widely used in preliminary aircraft design for efficiency, though limited to inviscid flow.}{catAD}
\glsadd{Vortex-Lattice-Method}\dictentry{Vortex Lattice Method}{The vortex lattice method models lifting surfaces as a lattice of vortex rings or horseshoe vortices. Each element induces velocities on all others, and circulation strengths are solved by enforcing no-penetration at collocation points. It captures three-dimensional effects including induced drag, downwash, and lift distribution. More accurate than lifting line theory for complex planforms, it remains a standard tool in aerodynamic design.}{catAD}
\glsadd{Source-Panel-Method}\dictentry{Source Panel Method}{The source panel method covers a body surface with panels carrying source or sink distributions of uniform strength. Source strengths are set by requiring zero normal velocity on the body surface. It is effective for modeling non-lifting bodies such as fuselages and nacelles. The method produces accurate pressure distributions for inviscid, incompressible flow around symmetric shapes. Vortex elements must be added for lifting surfaces.}{catAD}
\glsadd{Doublet}\dictentry{Doublet}{A doublet is a potential flow element from a source and sink of equal strength placed infinitesimally close together. The velocity field decays as the inverse cube of distance. Doublets combined with uniform flow model cylinders and spheres. In 3D, a doublet produces flow around a sphere. They build blocks for complex potential flow solutions.}{catAD}
\glsadd{Rankine-Body}\dictentry{Rankine Body}{A Rankine body is a closed streamlined shape obtained by superimposing a source in a uniform free stream. The body is blunt at the nose with a stagnation point, tapering smoothly to a closed trailing edge along the dividing streamline. Its pressure distribution can be computed analytically from potential flow theory. Though no real body matches exactly, it serves as a useful model for blunt-nosed body aerodynamics.}{catAD}
\glsadd{Joukowski-Airfoil}\dictentry{Joukowski Airfoil}{The Joukowski airfoil is generated by conformal mapping of a circle in the complex plane. Offsetting the center produces airfoils with rounded leading and sharp trailing edges. The sharp edge enforces the Kutta condition, determining circulation and lift. The transformation provides closed-form solutions. It was among the first models to predict finite airfoil lift.}{catAD}
\glsadd{Thin-Airfoil-Theory}\dictentry{Thin Airfoil Theory}{Thin airfoil theory analyzes airfoils of small thickness and camber at low angles using linearized equations. It replaces the airfoil with a camber line distributing vorticity to satisfy the Kutta condition. Lift coefficient varies linearly with angle of attack, with the aerodynamic center at quarter-chord. The theory also predicts the zero-lift angle for cambered airfoils. It remains essential for understanding fundamental airfoil behavior.}{catAD}
\glsadd{Lifting-Line-Theory}\dictentry{Lifting Line Theory}{Prandtl's lifting line theory models a finite wing as a bound vortex with trailing vortices. It relates local circulation to angle of attack and downwash from the trailing sheet. An elliptical distribution minimizes induced drag. The theory explains why high-aspect-ratio wings are efficient. Despite simplifications, it remains a foundational framework.}{catAD}
\glsadd{Lifting-Surface-Theory}\dictentry{Lifting Surface Theory}{Lifting surface theory distributes vorticity over the entire wing planform rather than along a single line. It handles arbitrary planforms including swept, tapered, and delta shapes with greater accuracy. The theory accounts for both chordwise and spanwise loading variations. It forms the basis for numerical methods like the vortex lattice method. It is valuable for analyzing low-aspect-ratio and highly swept wings.}{catAD}
\glsadd{Sweep-Theory}\dictentry{Sweep Theory}{Sweep theory analyzes swept wings by resolving velocity into normal and parallel components. Only the normal component affects pressure and lift, reducing effective aerodynamic performance. This explains why swept wings delay compressibility effects. Spanwise flow can cause boundary layer problems. The theory is fundamental to transonic wing design.}{catAD}
\glsadd{Slender-Body-Theory}\dictentry{Slender Body Theory}{Slender body theory predicts forces on long, thin bodies at small angles of attack. It assumes cross-sections change slowly and perturbations are dominated by the crossflow plane. The method decomposes flow into two-dimensional crossflow and longitudinal components. It is useful for missiles, fuselages, and slender configurations. The theory provides analytical insight where classical wing theory cannot apply.}{catAD}
\glsadd{Area-Rule}\dictentry{Area Rule}{The area rule, discovered by Richard Whitcomb at NACA, states that an aircraft's total cross-sectional area should vary smoothly to minimize transonic wave drag. Sharp area changes cause strong shocks and drag rise near Mach 1. Practical application involves adding fairings to create a coke-bottle shape. The rule revolutionized transonic design for military and civilian aircraft. It remains a guiding principle for transonic performance.}{catAD}
\glsadd{Supercritical-Airfoil}\dictentry{Supercritical Airfoil}{A supercritical airfoil delays strong shock waves and drag rise at transonic speeds. Its flattened upper surface weakens the local supersonic region while a cambered aft section recovers lift. The design enables efficient cruise at higher subsonic Mach numbers. Pioneered by Richard Whitcomb at NASA, these airfoils are now standard on modern aircraft. They enable fuel savings through higher cruise speeds or altitudes.}{catAD}
\glsadd{Shock-Free-Airfoil}\dictentry{Shock-Free Airfoil}{A shock-free airfoil avoids shock formation at its design cruise condition. The pressure distribution is tailored to prevent abrupt supersonic-to-subsonic deceleration. Computational methods achieve near-zero entropy rise at the design point. It represents the ideal of wave-drag-free transonic flight. Off-design conditions still produce shocks.}{catAD}
\glsadd{Natural-Laminar-Flow}\dictentry{Natural Laminar Flow}{Natural laminar flow shapes the airfoil pressure distribution to maintain a laminar boundary layer over much of the chord. A favorable pressure gradient over the forward section delays transition to turbulence. Success requires careful control of surface roughness and manufacturing imperfections. When achieved, it significantly reduces skin friction drag. The approach is applied on general aviation aircraft and some airliner wing sections.}{catAD}
\glsadd{Laminar-Flow-Control}\dictentry{Laminar Flow Control}{Laminar flow control uses active techniques to maintain laminar boundary layers beyond what shaping alone achieves. Methods include suction through porous surfaces or slots removing low-momentum wall fluid. This stabilizes the boundary layer and reduces skin friction drag. Demonstrated on the F-16XL, practical use must address contamination and weight.}{catAD}
\glsadd{Riblet}\dictentry{Riblet}{Riblets are small longitudinal grooves reducing turbulent skin friction drag. Sized at tens of micrometers, comparable to the viscous sublayer, they inhibit spanwise vortex motion. Drag reductions of five to ten percent have been measured. Riblet films are evaluated for aircraft fuselages and nacelles. The technology is passive but requires clean surfaces.}{catAD}
\glsadd{Dimple}\dictentry{Dimple}{A dimple is a surface indentation that reduces pressure drag by promoting earlier transition and delaying separation. The golf ball is the classic example, where turbulated flow stays attached longer, reducing wake size. Engineered dimple patterns are studied for bluff bodies and aircraft components. They transfer momentum from outer flow to the near-wall region.}{catAD}
\glsadd{Base-Drag}\dictentry{Base Drag}{Base drag results from low-pressure recirculating flow behind a blunt trailing edge or body base. Flow separation at a sharp corner creates a wake with pressure well below freestream, producing a rearward force. It is a major drag contributor on projectiles, rockets, and afterbodies. The magnitude depends on base diameter, Mach number, and boundary layer state. Reducing base drag is a key objective in boattail and base bleed design.}{catAD}
\glsadd{Boattail}\dictentry{Boattail}{A boattail is a tapered contour on the aft body reducing base drag by gradually decreasing area before the base. The convergence shrinks the separated wake and raises base pressure. Angles are limited to about fifteen degrees to avoid surface separation. Common on projectiles, rockets, and nacelles, effectiveness depends on closure angle and Mach number.}{catAD}
\glsadd{Bump}\dictentry{Bump}{A bump is a localized protuberance on a smooth surface that alters the local flow field. Bumps can be intentional, like vortex generators, or unintentional, like antenna housings. Well-designed bumps can energize the boundary layer and delay separation. Unfaired bumps increase drag and may cause local shock formation at transonic speeds. Streamlining surface bumps is important for drag reduction on production aircraft.}{catAD}
\glsadd{Fence}\dictentry{Fence}{A fence is a thin vertical plate mounted on an aerodynamic surface to block spanwise boundary layer flow. On swept wings, crossflow can drive the boundary layer toward the tips, causing separation and loss of aileron effectiveness. Fences interrupt this migration and maintain attached flow over outboard sections. They are located where separation is most likely. Fences are a simple, lightweight solution for improving swept-wing handling.}{catAD}
\glsadd{Wing-Fence}\dictentry{Wing Fence}{A wing fence is mounted on the upper surface of a swept or tapered wing to prevent spanwise boundary layer migration. Extending from the leading edge aft, it keeps flow chordwise. It is critical on swept wings where crossflow causes tip stall. The fence divides the wing into sections with better pressure gradients. Classic swept-wing aircraft use them for low-speed improvement.}{catAD}
\glsadd{Dogtooth}\dictentry{Dogtooth}{A dogtooth is a small triangular notch at the leading edge of a wing, typically at the inboard-outboard junction. It generates a vortex that energizes the upper-surface boundary layer and delays separation at high angles. It serves a similar purpose to a leading edge vortex generator but is simpler. The feature appeared on aircraft like the Messerschmitt Me 262. It improves maximum lift and stall characteristics without movable devices.}{catAD}
\glsadd{Leading-Edge-Extension}\dictentry{Leading Edge Extension}{A leading edge extension, or LEX, extends forward from the wing root along the fuselage. At high angles it generates a powerful vortex sweeping over the wing, delaying stall. This vortex lift mechanism enhances agility and high-angle performance. The LEX also contributes lift through its induced low-pressure field. It is a defining feature of modern fighters.}{catAD}
\glsadd{LEX-Fence}\dictentry{LEX Fence}{A LEX fence is a small vertical plate on a leading edge extension controlling the generated vortex. Without it, the vortex can wander spanwise and interact poorly with wing and tail surfaces. The fence stabilizes the core position and delays breakdown. It is typically a short chordwise plate near the trailing edge. The F/A-18 uses LEX fences for tail interaction control.}{catAD}
\glsadd{Strakes}\dictentry{Strakes}{Strakes are sharp-edged, low-aspect-ratio surfaces at the wing root generating controlled vortices at high angles. Like leading edge extensions but more pronounced, they maintain attached flow and provide vortex lift. Strakes improve post-stall handling on fighters, missiles, and racing vehicles. Geometry is optimized to balance lift benefits against drag penalties.}{catAD}
\glsadd{Chine}\dictentry{Chine}{A chine is a sharp longitudinal edge along a fuselage side, creating a cross-section break. At angle of attack it generates a vortex contributing lift, like a delta wing leading edge. Chined fuselages appear on stealth aircraft and high-speed boats. The chine controls forebody vortex systems affecting directional stability. The SR-71 features prominent chines.}{catAD}
\glsadd{Belly-Flap}\dictentry{Belly Flap}{A belly flap is mounted on the underside of a fuselage or nacelle to generate drag or modify lift for deceleration or trim. Some belly flaps function as speed brakes, increasing pressure drag when deployed. They can also adjust pitching moment during descent or approach. Typically hydraulically actuated, they may deploy symmetrically or independently. Their lower-surface location keeps disrupted airflow away from wings and control surfaces.}{catAD}
\glsadd{Deceleron}\dictentry{Deceleron}{A deceleron is a split aileron whose upper and lower halves deflect independently, acting as both aileron and airbrake. Symmetric opening creates drag for deceleration. Differential deflection produces a rolling moment like a conventional aileron. This dual functionality saves weight and complexity. The deceleron was used on the Vought F-8 Crusader to improve roll authority and deceleration during approach and combat.}{catAD}
\glsadd{Flapperon}\dictentry{Flapperon}{A flapperon combines the functions of a flap and aileron on a single trailing-edge surface. Symmetric deflection increases camber and lift for takeoff and landing, while differential deflection provides roll control. This reduces the number of surfaces needed on the wing. Flapperons are common on light aircraft and some military designs. The mechanical design must accommodate both functions' range of motion and load requirements.}{catAD}
\glsadd{Taileron}\dictentry{Taileron}{A taileron is a horizontal tail deflected differentially to provide roll control in addition to normal pitch. Both sides together act as an elevator; opposite deflections create a rolling moment. Taileron systems are found on aircraft without conventional ailerons, such as delta-wing designs. They provide an alternative roll mechanism while reducing distinct control surfaces. The system requires mixing of pilot roll and pitch inputs.}{catAD}
\glsadd{Stabilator}\dictentry{Stabilator}{A stabilator is a fully movable horizontal tail pivoting about a spanwise axis, replacing the fixed stabilizer and movable elevator. The entire surface moves, avoiding hinge moment limitations. Stabilators are common on high-performance military aircraft and some jet transports. The surface is balanced aerodynamically or with mass weights to prevent flutter. Trim tabs or differential deflection fine-tune control forces at low speeds.}{catAD}
\glsadd{V-Tail}\dictentry{V-Tail}{A V-tail mounts two surfaces in a V arrangement, combining horizontal and vertical stabilizer functions. A ruddervator mechanism mixes control inputs for pitch and yaw. The V-tail reduces surfaces from three to two, lowering interference drag and wetted area. However, it introduces mixing complexity and pitch-yaw coupling. V-tails were used on the Beechcraft Bonanza and the Lockheed F-117.}{catAD}
\glsadd{Cruciform-Tail}\dictentry{Cruciform Tail}{A cruciform tail has the horizontal stabilizer partway up the vertical fin, forming a cross shape. This keeps the horizontal surface above the wing wake and exhaust, improving effectiveness. It is a compromise between the low tail's simplicity and the T-tail's aerodynamic advantages. Cruciform tails are used on military transports, bombers, and some airliners. The surface intersection must be carefully faired to minimize interference drag.}{catAD}
\glsadd{T-Tail}\dictentry{T-Tail}{A T-tail places the horizontal stabilizer at the top of the vertical fin, above the wing wake and exhaust. This provides cleaner airflow for consistent pitch control effectiveness. The configuration can also reduce fuselage bending loads. However, in deep stall the surface can become blanketed by the stalled wing wake, risking loss of pitch control. This requires careful stall testing and may necessitate a stick-pusher system.}{catAD}
\glsadd{H-Tail}\dictentry{H-Tail}{An H-tail uses two vertical fins on the horizontal stabilizer, forming an H shape. This provides redundant directional stability where a single fin is impractical. Twin fins can be shorter than one of equivalent effectiveness, reducing height. H-tails appear on multi-engine transports and business jets. Design must account for surface interference.}{catAD}
\glsadd{Twin-Tail}\dictentry{Twin Tail}{A twin tail uses two separate vertical fins instead of a single central stabilizer. This reduces overall height, aiding carrier operations and hangar storage. Twin tails improve directional control at high angles where a single fin might be blanketed. They are common on fighters like the F-15 Eagle and F-22 Raptor. Fin spacing and sizing must be designed to avoid adverse aerodynamic interactions.}{catAD}
\glsadd{Endplate}\dictentry{Endplate}{An endplate is a vertical plate at a wingtip that reduces induced drag by partially blocking the tip vortex. It prevents high-pressure air from flowing around the tip, effectively increasing aspect ratio. Endplates are common on racing cars and sailing vessels where span is limited. On aircraft they have been explored as alternatives to winglets. Optimal geometry depends on conditions and the trade-off between drag reduction and weight.}{catAD}
\glsadd{Hoerner-Tip}\dictentry{Hoerner Tip}{The Hoerner tip is a rounded, drooped wingtip by Sighard Hoerner to reduce induced drag. It encourages the tip vortex to form where it effectively increases span. The rounded contour smooths the pressure distribution, weakening the vortex. Simpler and lighter than winglets, it provides meaningful drag reduction. Widely used on general aviation aircraft.}{catAD}
\glsadd{Raked-Wingtip}\dictentry{Raked Wingtip}{A raked wingtip has a swept trailing edge extending effective wingspan without adding much physical span. The shape shifts the tip vortex outward and reduces its strength, lowering induced drag. Unlike winglets, raked tips add less structural weight and interference drag. The Boeing 787 Dreamliner uses raked wingtips as a key efficiency feature. They are most effective on long-range cruise aircraft.}{catAD}
\glsadd{Tip-Fence}\dictentry{Tip Fence}{A tip fence is a small vertical surface at the wingtip extending above and below the wing plane. It functions as a combined endplate and mini-winglet, blocking spanwise flow and weakening the tip vortex. Compact and lightweight, it is attractive where full winglets are impractical. Used on military and research aircraft, effectiveness depends on geometry.}{catAD}
\glsadd{Wingtip-Vortex}\dictentry{Wingtip Vortex}{A wingtip vortex is a concentrated rotating flow at the tip of a lifting wing from the surface pressure difference. High-pressure air curls around the tip creating a spiraling wake. Vortices are the primary cause of induced drag, an irreducible cost of finite-span lift. Strength increases with lift and decreases with speed and wingspan.}{catAD}
\glsadd{Contrail}\dictentry{Contrail}{A contrail, or condensation trail, is a line of water droplets or ice crystals forming behind an aircraft at high altitudes. Exhaust vapor mixes with cold air, condensing if temperature is low enough. Persistence depends on humidity: dry air dissipates them quickly while humid air lets them spread. Persistent contrails trap infrared radiation, impacting climate.}{catAD}
\glsadd{Vapor-Cone}\dictentry{Vapor Cone}{A vapor cone is visible water vapor forming around an aircraft at transonic speeds. Local supersonic flow causes temperature drops below the dew point, triggering condensation into a cone-shaped cloud. Vapor cones appear mainly during low-altitude high-speed passes in humid air. The phenomenon demonstrates transonic effects and the Prandtl-Glauert singularity.}{catAD}
\glsadd{Shock-Diamond}\dictentry{Shock Diamond}{Shock diamonds are repeating bright patterns in supersonic engine exhaust at off-design conditions. They form when exhaust pressure differs from ambient, creating shocks and expansion fans reflecting off the jet boundary. Each intersection increases density, scattering light into visible diamonds. Spacing relates to pressure ratio and Mach number. Common on military aircraft and rockets.}{catAD}
\glsadd{Mach-Wave}\dictentry{Mach Wave}{A Mach wave is a weak pressure disturbance propagating at the Mach angle in supersonic flow. Unlike a shock, no finite property change occurs across it. The Mach angle is the arcsine of the reciprocal of Mach number. Each small disturbance on a body generates a Mach wave within the Mach cone. Mach waves collectively define the boundary of upstream influence in supersonic flow.}{catAD}
\glsadd{Bow-Shock}\dictentry{Bow Shock}{A bow shock is a detached, curved shock wave ahead of a blunt body in supersonic flow. The body's bluntness prevents sharp turning, so the shock stands off at finite distance from the nose. Behind it, flow is subsonic near the stagnation point and accelerates to supersonic around the body. Standoff distance depends on bluntness and Mach number. Bow shocks are encountered on re-entering spacecraft and blunt supersonic objects.}{catAD}
\glsadd{Attached-Shock}\dictentry{Attached Shock}{An attached shock is an oblique shock anchored to the sharp leading edge of a body in supersonic flow. Its wave angle is set by the deflection angle and freestream Mach number. Flow behind it remains supersonic and follows the body surface. Attached shocks occur on sharp-nosed wedges, cones, and pointed inlets at sufficient Mach numbers. Transition from attached to detached happens below a geometry-dependent critical Mach number.}{catAD}
\glsadd{Lambda-Shock}\dictentry{Lambda Shock}{A lambda shock resembles the Greek letter lambda in transonic flow over airfoils. A normal shock on the upper surface bifurcates near the surface into oblique branches. The structure results from shock-boundary layer interaction causing local separation. This reduces wave drag versus a single strong shock. Lambda shocks are observed on supercritical airfoils.}{catAD}
\glsadd{Fish-Tail-Shock}\dictentry{Fish-Tail Shock}{A fish-tail shock is a pattern resembling a fish tail on overexpanded convergent-divergent nozzles. The exhaust jet recompresses through oblique shocks matching ambient pressure. Shocks intersect at the centerline and reflect from the jet boundary, creating a repeating structure. Visible in schlieren photography and sometimes in the plume itself, it indicates significant deviation from ideal expansion and wasted propulsive efficiency.}{catAD}
\glsadd{Detached-Shock}\dictentry{Detached Shock}{A detached shock stands at finite distance from a body rather than attaching to a sharp feature. It occurs when the body is too blunt or Mach number too low for an attached shock. Subsonic flow exists between shock and body near the nose, weakening toward the flanks. Standoff increases with lower Mach and greater bluntness. It produces high stagnation temperatures.}{catAD}
\glsadd{Expansion-Fan}\dictentry{Expansion Fan}{An expansion fan is a region of continuous isentropic turning when supersonic flow meets a convex corner. Unlike a shock, properties change smoothly with velocity increasing and pressure decreasing. The fan spreads outward from the corner at angles set by the local Mach number. Total turning is computed via the Prandtl-Meyer function. Expansion fans are fundamental in supersonic flow analysis for nozzles, wings, and external corners.}{catAD}
\glsadd{Prandtl-Meyer-Function}\dictentry{Prandtl-Meyer Function}{The Prandtl-Meyer function gives the maximum isentropic turning angle for a given Mach number in supersonic flow. Derived from the characteristic equations, it depends on Mach number and the ratio of specific heats. The function increases monotonically, meaning higher-speed flows undergo greater expansion. It computes flow properties across expansion fans and around convex corners. It is one of the most important tools in supersonic aerodynamics.}{catAD}
\glsadd{Method-of-Characteristics}\dictentry{Method of Characteristics}{The method of characteristics solves supersonic flow PDEs by reducing them to ODEs along characteristic lines. These lines, oriented at the Mach angle, carry disturbances. The technique solves 2D and axisymmetric problems including expansion fans, shocks, and nozzle design. It provides exact inviscid solutions and benchmarks computational methods.}{catAD}
\glsadd{Supersonic-Wind-Tunnel}\dictentry{Supersonic Wind Tunnel}{A supersonic wind tunnel produces controlled airflow at Mach 1.2 to 5.0 for aerodynamic testing. Flow accelerates through a converging-diverging nozzle to supersonic speeds in the test section. Continuous tunnels use closed circuits with compressors and heat exchangers. Blow-down tunnels store high-pressure air for shorter, cheaper test runs. These tunnels study shock waves, wave drag, and aircraft behavior above Mach 1.}{catAD}
\glsadd{Hypersonic-Wind-Tunnel}\dictentry{Hypersonic Wind Tunnel}{A hypersonic wind tunnel produces airflow above Mach 5.0, simulating re-entry conditions. Very high stagnation pressures and temperatures are needed, often with heating to prevent gas liquefaction. Facilities may be blow-down, shock-driven, or arc-heated. Real gas effects like dissociation complicate design. Critical for thermal protection and hypersonic research.}{catAD}
\glsadd{Shock-Tunnel}\dictentry{Shock Tunnel}{A shock tunnel uses a shock wave in a driver tube to produce high-enthalpy flow for milliseconds. Pressurized driver gas is suddenly released, heating and accelerating the test gas. Shock tunnels simulate extreme hypersonic enthalpies without continuous heating. Short test time is offset by high data acquisition rates and unique conditions. They are widely used for hypersonic research and scramjet testing.}{catAD}
\glsadd{Arc-Heated-Wind-Tunnel}\dictentry{Arc-Heated Wind Tunnel}{An arc-heated wind tunnel uses an electric arc to heat gas to extreme temperatures, simulating re-entry conditions. The high-current arc creates plasma-like flow through dissociation and ionization. Heated gas expands through a nozzle to desired conditions. Arc facilities achieve enthalpies far beyond combustion tunnels. They are primarily used for thermal protection material testing.}{catAD}
\glsadd{Cryogenic-Wind-Tunnel}\dictentry{Cryogenic Wind Tunnel}{A cryogenic wind tunnel uses nitrogen cooled to around 100 Kelvin for high Reynolds numbers in moderate facilities. Cooling increases density and decreases viscosity, both raising Reynolds number. NASA's Transonic Cryogenic Tunnel pioneered full-scale transonic testing. The approach requires specialized materials and insulation. It provides highly accurate transonic data.}{catAD}
\glsadd{Indraft-Wind-Tunnel}\dictentry{Indraft Wind Tunnel}{An indraft wind tunnel draws room air through the test section via a downstream exhaust fan. The room serves as the settling chamber, providing calm air that enters through a contraction cone. Indraft tunnels are simple and inexpensive, popular for education and low-speed research. They are sensitive to room temperature, drafts, and background turbulence. Flow quality depends on the room environment, requiring careful control.}{catAD}
\glsadd{Open-Circuit-Tunnel}\dictentry{Open Circuit Tunnel}{An open circuit tunnel draws air from and exhausts it to the atmosphere without recirculation. A fan pushes or pulls air through a contraction, test section, and diffuser. The design is straightforward to build and maintain with no return ductwork. It suits low-speed and educational applications with modest flow quality needs. The main disadvantage is sensitivity to ambient conditions and higher operating costs.}{catAD}
\glsadd{Closed-Circuit-Tunnel}\dictentry{Closed Circuit Tunnel}{A closed circuit tunnel recirculates air through a loop of test section, diffuser, corners, settling chamber, contraction, and fan. Recirculation provides superior flow quality with minimal disturbance. The closed loop enables independent temperature and humidity control while reducing noise. Return ducts add complexity but offset by lower power requirements. Most modern research tunnels use closed circuit designs for performance and efficiency.}{catAD}
\glsadd{Low-Speed-Tunnel}\dictentry{Low Speed Tunnel}{A low speed tunnel operates below Mach 0.3 where compressibility is negligible. It tests aircraft takeoff and landing, ground vehicle aerodynamics, wind loads, and sports equipment. Flow is driven by axial or centrifugal fans at velocities from a few to about 100 meters per second. Low speed tunnels are the most common type due to simplicity and low cost. They remain essential for high-lift, landing gear, and propeller studies.}{catAD}
\glsadd{High-Speed-Tunnel}\dictentry{High Speed Tunnel}{A high speed tunnel operates in the transonic to low supersonic range, from Mach 0.5 to about 5.0. It requires converging-diverging nozzles and test sections with perforated or slotted walls for transonic interference control. These tunnels study shock formation, wave drag, and buffet. They may be continuous, blow-down, or intermittent. Data are critical for transport, fighter, and missile design.}{catAD}
\glsadd{Wall-Interference}\dictentry{Wall Interference}{Wall interference is flow distortion around a model caused by tunnel walls constraining the flow differently from free air. In subsonic flow, walls artificially accelerate or decelerate the flow. Transonic interference is particularly complex due to shock reflections. Corrections are essential for accurate free-air data. Modern transonic tunnels use adaptive or porous walls to minimize this effect.}{catAD}
\glsadd{Blockage}\dictentry{Blockage}{Blockage is the ratio of model frontal area to tunnel cross-sectional area. High blockage accelerates flow around the model, distorting pressures and altering Reynolds number. It also affects shock location and strength in transonic tests. Operators typically limit blockage to a few percent. Blockage corrections are applied to test data to account for confinement effects.}{catAD}
\glsadd{Turbulence-Intensity}\dictentry{Turbulence Intensity}{Turbulence intensity measures velocity fluctuation level as the ratio of root-mean-square fluctuation to mean velocity. It characterizes wind tunnel flow quality, with lower values indicating smoother conditions. High intensity triggers premature transition and alters drag measurements. Research tunnels achieve below one percent, industrial tunnels several percent. Turbulence is controlled by upstream honeycomb straighteners and screens.}{catAD}
\glsadd{Flow-Angularity}\dictentry{Flow Angularity}{Flow angularity is the deviation of local flow direction from the tunnel axis. Even small misalignments cause significant lift and moment errors, especially for high-aspect wings. It arises from imperfections in the contraction, turning vanes, or model alignment. It is measured using calibration probes traversed across the test section. Correcting for flow angularity is standard in data reduction.}{catAD}
\glsadd{Sting-Mount}\dictentry{Sting Mount}{A sting mount holds a model from behind, extending from the base through the wake. It minimizes interference because it contacts the model only at the base where flow is already disturbed. This makes it preferred for aircraft and missile models needing accurate forebody data. The sting must be strong yet slender enough to minimize wake interference. Integrated sting balances measure forces and moments simultaneously.}{catAD}
\glsadd{Strut-Mount}\dictentry{Strut Mount}{A strut mount supports a model from below, above, or the side using rigid struts. The strut contacts the model away from the base, potentially interfering with surface flow. However, it allows base flow access and is necessary when rear stings are impractical. Struts are faired to minimize interference with corrections applied. Strut mounts are common for automotive, building, and store separation testing.}{catAD}
\glsadd{Magnetic-Suspension}\dictentry{Magnetic Suspension}{Magnetic suspension holds a model using electromagnetic forces with no physical contact. Ferromagnetic elements in the model interact with external magnets maintaining position. This eliminates all support interference for purest aerodynamic data. Systems are complex and expensive, requiring sophisticated control. Used where zero interference is critical, such as precision drag measurement.}{catAD}
\glsadd{Pressure-Sensitive-Paint}\dictentry{Pressure Sensitive Paint}{Pressure sensitive paint fluoresces under ultraviolet light with intensity varying with local air pressure. Luminescent molecules are quenched by oxygen, whose partial pressure is proportional to total pressure. A camera records the pattern, providing continuous pressure mapping. This replaces hundreds of discrete pressure taps with a single global measurement. The technique has revolutionized aerodynamic testing with full-field pressure data.}{catAD}
\glsadd{Temperature-Sensitive-Paint}\dictentry{Temperature Sensitive Paint}{Temperature sensitive paint has luminescent properties that change predictably with surface temperature. Phosphor or dye molecules emit light whose intensity or lifetime varies with temperature. When excited, the paint reveals the spatial temperature distribution. It provides direct surface temperature data on complex geometries where infrared viewing may be obstructed. It is particularly useful for studying hypersonic heating patterns.}{catAD}
\glsadd{Infrared-Thermography}\dictentry{Infrared Thermography}{Infrared thermography uses a camera to detect thermal radiation and map surface temperatures. Every object above absolute zero emits radiation proportional to temperature and emissivity. It identifies transition locations, hot spots, and heating distributions in testing. Calibration and emissivity knowledge are needed. It provides full-field data with high resolution and no intrusiveness.}{catAD}
\glsadd{Optical-Strain-Measurement}\dictentry{Optical Strain Measurement}{Optical strain measurement uses cameras and image correlation to measure deformation without contact. A speckle pattern on the surface is tracked as the structure deforms, computing full-field strain. This eliminates discrete strain gauges, providing continuous data over the visible surface. It is widely used in aerospace structural testing. The method offers high accuracy for complex strain patterns.}{catAD}
\glsadd{Accelerometer}\dictentry{Accelerometer}{An Accelerometer is a sensor that measures linear acceleration along three orthogonal axes using a micro-mechanical proof mass. In a drone it detects gravitational acceleration to determine tilt angle and measures vibration and dynamic motion during flight. The accelerometer data is fused with gyroscope readings by the flight controller to maintain accurate attitude estimation. High-quality accelerometers with low noise and stable bias are essential for precise altitude hold and smooth flight performance.}{catAD}
\glsadd{Strain-Gauge}\dictentry{Strain Gauge}{A strain gauge is thin foil bonded to a surface whose resistance changes when deformed. Structural stretching or compression changes gauge dimensions, producing resistance change proportional to strain. Wheatstone bridge circuits amplify signals and compensate for temperature. Strain gauges are the most widely used deformation sensors in aerospace. Data inform stress analysis and fatigue prediction.}{catAD}
\glsadd{Thermocouple}\dictentry{Thermocouple}{A thermocouple uses two dissimilar metals joined at one end, generating voltage via the Seebeck effect proportional to temperature difference. It provides simple, robust measurement over wide ranges. Widely used for surface and gas temperatures in aerodynamic testing. Types cover below freezing to over two thousand degrees. Small size, fast response, and reliability make them indispensable.}{catAD}
\glsadd{Balance}\dictentry{Balance}{Balance refers to the process of adjusting the mass distribution of a spacecraft so that its center of mass aligns with the intended reference point. Proper balance is essential for stable flight, accurate thrust vectoring, and proper deployment dynamics. Balancing involves adding or relocating small masses to correct asymmetries. An unbalanced spacecraft can exhibit unwanted tumbling or off-nominal trajectory behavior.}{catAD}
\glsadd{External-Balance}\dictentry{External Balance}{An external balance is mounted outside the test section, connected through walls via linkages. It measures forces and moments transmitted from the model to its support. External balances are accessible for calibration and can be made stiff for large loads. However, wall-penetrating linkages introduce friction and interference. They are common for half-model testing and models with limited internal space.}{catAD}
\glsadd{Internal-Balance}\dictentry{Internal Balance}{An internal balance is a sealed chamber built into the leading edge of a control surface, with a flexible seal allowing a portion of the surface to deflect aerodynamically within the chamber. Air pressure acting on the balanced portion reduces the hinge moment, lowering the force required from the pilot. The sealed design prevents external airflow from entering the balance cavity, maintaining predictable force reduction across a range of speeds. Internal balances are common on high-performance and transport aircraft.}{catAD}
\glsadd{Half-Model-Testing}\dictentry{Half Model Testing}{Half model testing mounts a half-span model against a symmetry plane, typically a splitter plate or wall. The wall reflects the flow to simulate the full aircraft's symmetry. This allows larger scales and higher Reynolds numbers within a given tunnel. The splitter plate boundary layer must not interfere with the model. Half model testing is particularly useful for wing and nacelle studies.}{catAD}
\glsadd{Full-Model-Testing}\dictentry{Full Model Testing}{Full model testing evaluates a complete aircraft model including all major components. This captures aerodynamic interactions lost in isolated or half-model tests. Full model tests provide the most accurate prediction of integrated aerodynamic characteristics. Model size is limited by tunnel dimensions and blockage. It is the definitive wind tunnel step before flight.}{catAD}
\glsadd{Dynamic-Testing}\dictentry{Dynamic Testing}{Dynamic testing measures forces during time-varying motion such as oscillation, pitching, or plunging. It determines stability derivatives, flutter boundaries, and turbulence response. The model is driven by actuators while force and displacement are recorded at high rates. It captures damping derivatives and frequency-dependent loads that static testing cannot. Results are critical for flight simulation, control design, and flutter clearance.}{catAD}
\glsadd{Static-Testing}\dictentry{Static Testing}{Static testing measures steady-state forces and moments at fixed angles of attack and sideslip. It provides baseline characteristics: lift curves, drag polars, and moment coefficients. Static tests are the most common form of wind tunnel testing. The model is stepped through an angle matrix while the balance records forces. Data are used for performance, stability, and flight simulation.}{catAD}
\glsadd{Oscillatory-Testing}\dictentry{Oscillatory Testing}{Oscillatory testing forces the model through sinusoidal motions in pitch, yaw, or roll. Phase lag and amplitude ratio between motion and forces yield unsteady aerodynamic derivatives. These derivatives predict flutter, limit cycle oscillations, and dynamic stability. The technique requires precise excitation and high-bandwidth force measurement. Data are important for flexible aircraft and validating computational unsteady methods.}{catAD}
\glsadd{Grid-Turbulence}\dictentry{Grid Turbulence}{Grid turbulence is nearly homogeneous isotropic flow generated by passing uniform air through a mesh. The grid creates vortices that decay downstream with well-defined statistics. It is a canonical flow for studying energy cascade, decay rates, and length scale growth. In tunnels, grids produce controlled disturbance levels for model response testing. Grid turbulence is one of the most studied flows in fluid mechanics.}{catAD}
\glsadd{Hot-Wire-Anemometry}\dictentry{Hot Wire Anemometry}{Hot wire anemometry uses a thin heated wire to determine velocity from convective heat transfer. Faster flow cools the wire more, changing the power needed to maintain its temperature. The technique provides very high frequency response for turbulent fluctuations up to hundreds of kilohertz. Probes can measure single or multi-component velocities. Despite fragility, it remains one of the most precise tools for turbulent flow studies.}{catAD}
\glsadd{Laser-Doppler-Velocimetry}\dictentry{Laser Doppler Velocimetry}{Laser Doppler velocimetry measures local velocity without inserting a probe. Two crossed laser beams create interference fringes, and seed particles scatter light at a frequency proportional to velocity. The Doppler shift is detected and processed for instantaneous velocity. The technique is non-intrusive, highly accurate, and measures very low to very high velocities. It is widely used in wind tunnels, combustion, and biomedical flows.}{catAD}
\glsadd{Phase-Doppler-Anemometry}\dictentry{Phase Doppler Anemometry}{Phase Doppler anemometry extends laser Doppler velocimetry to measure particle velocity and size simultaneously. Multiple detectors measure phase differences for diameter while velocity comes from Doppler frequency. Invaluable for sprays, atomization, and two-phase flows. Used in fuel injection, meteorology, and pharmaceutical research. Non-intrusive and accurate point measurements.}{catAD}
\glsadd{Pressure-Scanner}\dictentry{Pressure Scanner}{A pressure scanner measures static pressure at many points via a multiplexed transducer. Ports connect to pressure taps with rapid channel switching. Modern scanners achieve Pascal-level accuracy across hundreds of channels in seconds. They reduce instrumentation complexity versus individual transducers. Standard in wind tunnel testing for detailed pressure measurements.}{catAD}
\glsadd{Wind-Tunnel-Fan}\dictentry{Wind Tunnel Fan}{The wind tunnel fan is the primary mechanism driving airflow through the tunnel circuit. Usually axial or centrifugal, it overcomes total circuit pressure losses. Power depends on dynamic pressure, test section area, and circuit efficiency. Variable-pitch or speed fans allow flow adjustment over wide ranges. The fan must produce uniform flow with minimal noise and vibration.}{catAD}
\glsadd{Honeycomb-Straightener}\dictentry{Honeycomb Straightener}{A honeycomb straightener is a bundle of small tubes in a hexagonal pattern in the settling chamber. It removes lateral velocity components and swirl, producing uniform axial flow before contraction. The length-to-cell-diameter ratio determines effectiveness. It also breaks large eddies into smaller ones that decay faster. Honeycomb straighteners are standard in high-quality wind tunnel circuits.}{catAD}
\glsadd{Settling-Chamber}\dictentry{Settling Chamber}{The settling chamber is the wide region between flow conditioning devices and the contraction nozzle. Its large cross-section ensures low velocities for flow to settle and become uniform. It houses screens and honeycomb elements that reduce turbulence and eliminate angularity. Length and diameter provide sufficient residence time for conditioning. A well-designed settling chamber is essential for high-quality, low-turbulence test flow.}{catAD}
\glsadd{Contraction-Cone}\dictentry{Contraction Cone}{A Contraction Cone is the converging section of a wind tunnel that accelerates airflow from the settling chamber to the test section. By smoothly reducing cross-sectional area, it increases velocity while lowering turbulence intensity. The contraction ratio, typically 4:1 to 16:1, determines the speed gain and flow quality. It is essential for producing the precise, uniform airflow needed for accurate aerodynamic measurements.}{catAD}
\glsadd{Diffuser-Section}\dictentry{Diffuser Section}{A Diffuser Section is the expanding duct downstream of a wind tunnel test section that decelerates airflow and recovers static pressure. It gradually increases area to convert kinetic energy back into pressure energy. A well-designed diffuser minimizes pressure losses and prevents flow separation. It is critical for reducing the power required to operate closed-circuit wind tunnels.}{catAD}
\glsadd{Corner-Vanes}\dictentry{Corner Vanes}{Corner Vanes are curved guide surfaces installed in the bends of closed-circuit wind tunnels to redirect airflow with minimal energy loss. They split the flow into guided channels that follow smooth curvatures, preventing separation at corners. Without them, flow distortion and pressure losses would degrade test-section quality. These vanes are essential for maintaining efficient, uniform circulation throughout the tunnel.}{catAD}
\glsadd{Turning-Vanes}\dictentry{Turning Vanes}{Turning Vanes are aerodynamic guide surfaces placed inside duct bends to redirect airflow around corners efficiently. They prevent boundary-layer separation and reduce total pressure losses by ensuring smooth flow paths. Turning vanes are found in wind tunnels, HVAC systems, and aircraft inlet ducts. Their presence significantly improves system efficiency by minimizing turbulence and recirculation.}{catAD}
\glsadd{Screens}\dictentry{Screens}{Fine-mesh Screens are grids placed in a wind tunnel settling chamber to reduce turbulence and improve flow uniformity. They break large turbulent eddies into smaller ones that dissipate more quickly. Multiple screens with progressively finer mesh are often installed in series for optimal results. The small pressure drop they cause is far outweighed by the improved flow quality for precision testing.}{catAD}
\glsadd{Acoustic-Testing}\dictentry{Acoustic Testing}{Acoustic Testing is a specialized wind tunnel methodology focused on measuring noise generated by aerodynamic flows over aircraft components. It requires anechoic or treated test sections with extremely low background noise to isolate specific sources. Microphone arrays and beamforming techniques map noise distributions across airframes and engines. This testing is essential for meeting stringent noise certification requirements.}{catAD}
\glsadd{Aeroacoustics}\dictentry{Aeroacoustics}{Aeroacoustics is the study of noise generation by turbulent flow and aerodynamic forces interacting with surfaces. It encompasses jet mixing noise, airframe self-noise, propeller acoustics, and fan tone generation. The field combines computational fluid dynamics with acoustic modeling to predict and reduce aircraft noise. It is critical for certification, community noise reduction, and passenger comfort.}{catAD}
\glsadd{Jet-Noise}\dictentry{Jet Noise}{Jet Noise is sound produced by turbulent mixing between a high-velocity engine exhaust plume and ambient air. It is generated primarily by large-scale eddies in the shear layer and increases with the eighth power of jet velocity. Jet noise dominates during takeoff when engines operate at maximum thrust. Mitigation strategies include chevron nozzles, mixers, and higher bypass-ratio engine designs.}{catAD}
\glsadd{Airframe-Noise}\dictentry{Airframe Noise}{Airframe Noise is aerodynamic noise from non-propulsive aircraft components such as landing gear, flaps, slats, and wing trailing edges. It becomes dominant during approach and landing when engines operate at low thrust. As modern engines have grown quieter, airframe noise has become a major challenge for meeting noise regulations. Solutions include streamlined gear fairings and serrated trailing edges.}{catAD}
\glsadd{Propeller-Noise}\dictentry{Propeller Noise}{Propeller Noise is sound generated by rotating blades due to periodic blade loading, thickness effects, and tip vortex interactions. It consists of tonal components at the blade passage frequency and broadband noise from turbulence. It is a primary noise concern for turboprop aircraft near populated areas. Advanced blade designs with swept tips and optimized distributions reduce noise levels.}{catAD}
\glsadd{Rotor-Noise}\dictentry{Rotor Noise}{Rotor Noise is the acoustic signature of helicopter main and tail rotors, including blade-vortex interaction and high-speed impulsive noise. It is particularly challenging because it is inherently tonal and low-frequency, making it difficult to attenuate. This is a critical consideration for urban air mobility and military operations near communities. Active and passive blade design strategies are employed to reduce rotor noise.}{catAD}
\glsadd{Fan-Noise}\dictentry{Fan Noise}{Fan Noise is sound generated by the inlet fan stage of turbofan engines, consisting of tonal noise at the blade passage frequency and broadband turbulence noise. The fan is typically the loudest component during approach. Noise propagates both forward through the inlet and rearward through the bypass duct. Acoustic nacelle liners and optimized blade geometry are the primary reduction methods.}{catAD}
\glsadd{Combustion-Noise}\dictentry{Combustion Noise}{Combustion Noise is generated by unsteady heat release and turbulent mixing inside a gas turbine combustor. It produces both direct noise and indirect noise through entropy waves traveling downstream. Modern lean-burn combustors reduce NOx but can increase combustion noise. Careful acoustic design of the combustor geometry is required to manage this trade-off.}{catAD}
\glsadd{Noise-Certification}\dictentry{Noise Certification}{Noise Certification is the regulatory process by which aircraft demonstrate compliance with maximum noise limits set by ICAO, FAA, and EASA. It defines standardized measurement procedures and noise limits for takeoff, approach, and sideline positions. Aircraft must pass before entering commercial service. These standards have progressively tightened, driving continuous noise reduction development.}{catAD}
\glsadd{Noise-Reduction}\dictentry{Noise Reduction}{Noise Reduction encompasses techniques to decrease aircraft noise, including nacelle liners, chevron nozzles, airframe shielding, and optimized flight procedures. Methods range from passive acoustic treatments to active noise control systems. They are combined for maximum effect across all flight phases. Reduction efforts are driven by certification standards and community noise concerns around airports.}{catAD}
\glsadd{Chevron-Nozzle}\dictentry{Chevron Nozzle}{A Chevron Nozzle has a serrated or sawtooth trailing edge that promotes mixing between hot exhaust and ambient air. The chevrons break up large-scale turbulent structures and shift noise energy to higher frequencies. Higher-frequency sound is more easily attenuated by the atmosphere, reducing perceived noise. Chevrons are now standard on many modern turbofan engines.}{catAD}
\glsadd{Hush-Kit}\dictentry{Hush Kit}{A Hush Kit is a set of acoustic modifications applied to older engines to reduce noise emissions. Components typically include acoustic liners, exhaust mixers, and nozzle treatments. Hush kits allow aging aircraft to meet modern noise certification standards without full engine replacement. They represent a cost-effective solution for fleet noise compliance.}{catAD}
\glsadd{Emission-Index}\dictentry{Emission Index}{An Emission Index measures the mass of a pollutant produced per unit mass of fuel burned, expressed in grams per kilogram of fuel. Separate indices track NOx, CO, unburned hydrocarbons, and soot from aircraft engines. Emission indices are measured during certification and used in atmospheric impact modeling. Lower values indicate cleaner-burning engines and fuels.}{catAD}
\glsadd{Carbon-Footprint}\dictentry{Carbon Footprint}{A Carbon Footprint is the total CO2 emissions from an aircraft or aviation operation over a defined period or mission. It includes direct fuel combustion emissions and lifecycle emissions from fuel production. Aviation accounts for roughly two to three percent of global CO2 output. Reduction strategies include sustainable fuels, aerodynamic improvements, and optimized operations.}{catAD}
\glsadd{Sustainable-Aviation-Fuel}\dictentry{Sustainable Aviation Fuel}{Sustainable aviation fuel is a liquid fuel derived from non-petroleum sources such as waste oils, agricultural residues, and synthesized from captured carbon dioxide and green hydrogen. SAF can be blended with conventional jet fuel and used in existing engines with little or no modification. It significantly reduces lifecycle carbon emissions compared to fossil-based jet fuel. SAF adoption is central to aviation's decarbonization strategy.}{catAD}
\glsadd{Hydrogen-Propulsion}\dictentry{Hydrogen Propulsion}{Hydrogen Propulsion uses hydrogen as aircraft fuel, either burned in modified turbines or converted to electricity via fuel cells. It produces zero CO2 emissions at the point of use and can reduce contrail formation. However, it requires large cryogenic storage tanks and new fueling infrastructure. It is a leading candidate for deep decarbonization of long-range aviation.}{catAD}
\glsadd{Electric-Propulsion}\dictentry{Electric Propulsion}{Electric Propulsion uses electric motors powered by batteries or fuel cells to drive propellers or fans. It enables distributed propulsion architectures with multiple small motors along wings. Electric propulsion produces zero local emissions and significantly reduces noise. It is promising for urban air mobility and short-range regional aircraft.}{catAD}
\glsadd{Hybrid-Propulsion}\dictentry{Hybrid Propulsion}{Hybrid Propulsion combines a thermal engine with an electric motor and battery system. The thermal engine operates at its most efficient cruise condition while the electric motor provides peak power for takeoff and climb. This reduces fuel consumption while extending range beyond pure battery aircraft. It is an emerging architecture for regional and urban air mobility.}{catAD}
\glsadd{Solar-Powered-Flight}\dictentry{Solar Powered Flight}{Solar Powered Flight uses photovoltaic cells on aircraft surfaces to convert sunlight into electrical energy for motors and batteries. Solar aircraft can achieve sustained high-altitude flight for days or months, limited primarily by energy storage and daylight cycles. These platforms suit communications relay, earth observation, and atmospheric research. Challenges include energy storage density and structural weight minimization.}{catAD}
\glsadd{High-Altitude-Flight}\dictentry{High Altitude Flight}{High Altitude Flight refers to operations above fifty thousand feet, where low air density demands high aerodynamic efficiency and specialized propulsion. Vehicles include stratospheric drones, research aircraft, and high-altitude balloons. The thin atmosphere reduces drag but also limits engine performance and control authority. High altitude platforms serve communications, surveillance, and scientific research.}{catAD}
\glsadd{Stratospheric-Flight}\dictentry{Stratospheric Flight}{Stratospheric Flight occurs within the stratosphere, roughly thirty-five thousand to one hundred sixty thousand feet. This layer offers stable conditions, minimal weather, and low turbulence. Stratospheric platforms can remain airborne for extended durations, ideal for telecommunications and earth observation. Low air density poses challenges for propulsion efficiency and aerodynamic lift.}{catAD}
\glsadd{Suborbital-Flight}\dictentry{Suborbital Flight}{Suborbital Flight follows a ballistic trajectory reaching space above the Kármán line but lacking velocity for orbit. The vehicle ascends under power, experiences brief weightlessness, then returns on a ballistic arc. It is used for sounding rockets, space tourism, and hypersonic research. It provides microgravity access at lower cost than orbital missions.}{catAD}
\glsadd{Atmospheric-Re-entry}\dictentry{Atmospheric Re-entry}{Atmospheric Re-entry is a spacecraft's return through a planetary atmosphere from space. It involves hypersonic velocities generating intense heating and high deceleration loads. Thermal protection systems and trajectory design are essential to prevent structural failure. Successful re-entry balances heat management with controlled deceleration for a safe landing.}{catAD}
\glsadd{Re-entry-Corridor}\dictentry{Re-entry Corridor}{A Re-entry Corridor is the narrow range of flight path angles for safe atmospheric entry. Entering too steeply causes excessive heating and structural loads. Entering too shallowly causes the vehicle to skip off the atmosphere. The corridor is defined by thermal, structural, and navigation constraints.}{catAD}
\glsadd{Re-entry-Angle}\dictentry{Re-entry Angle}{A Re-entry Angle is the angle between a spacecraft's flight path and the local horizontal at atmospheric entry. A steep angle produces high peak heating and rapid deceleration. A shallow angle reduces peak loads but risks atmospheric skip-out. Selecting the optimal angle is critical for vehicle survival.}{catAD}
\glsadd{Skip-Re-entry}\dictentry{Skip Re-entry}{Skip Re-entry involves a vehicle descending into the upper atmosphere, lifting back out, then re-entering a second time. This distributes heating over a longer period, reducing peak temperatures. It also extends cross-range capability and total flight distance. It was proposed for lunar missions and is used in some modern capsule designs.}{catAD}
\glsadd{Ballistic-Re-entry}\dictentry{Ballistic Re-entry}{Ballistic Re-entry follows a purely gravitational path with no aerodynamic lift, relying solely on atmospheric drag. This produces very high peak deceleration and heating concentrated in a short time. It was used by early capsules including Mercury and Vostok. The trade-off is limited trajectory control and extreme peak thermal conditions.}{catAD}
\glsadd{Lifting-Re-entry}\dictentry{Lifting Re-entry}{Lifting Re-entry uses aerodynamic lift to control descent and reduce peak deceleration. By modulating lift-to-drag ratio, the vehicle spreads heating over a longer period. It also enables cross-range maneuvering to reach distant landing sites. This approach was used by the Space Shuttle and is planned for Mars return vehicles.}{catAD}
\glsadd{Re-entry-Velocity}\dictentry{Re-entry Velocity}{Re-entry Velocity is a spacecraft's speed at atmospheric entry, ranging from about seven point eight km per second for low Earth orbit to eleven point two for escape velocity. Higher velocities produce dramatically greater kinetic energy dissipated as heat. It is the primary driver of thermal protection system design. Interplanetary returns face the most severe velocity challenges.}{catAD}
\glsadd{Re-entry-Heating-Rate}\dictentry{Re-entry Heating Rate}{Re-entry Heating Rate is thermal energy transferred to a spacecraft surface per unit area per unit time, measured in watts per square centimeter. It combines convective and radiative heating components. Peak heating occurs near maximum deceleration. Accurate prediction is essential for thermal protection system engineering.}{catAD}
\glsadd{Peak-Heating}\dictentry{Peak Heating}{Peak Heating is the maximum instantaneous heat flux a spacecraft experiences during re-entry. It occurs where velocity and atmospheric density combine to produce the highest convective and radiative heat transfer. Peak heating drives the selection and sizing of thermal protection materials. Managing it is the central challenge of re-entry thermal design.}{catAD}
\glsadd{Heat-Flux}\dictentry{Heat Flux}{Heat Flux is the rate of thermal energy transfer per unit area at a surface, measured in watts per square meter. During re-entry it combines convective, radiative, and catalytic heating components. Heat flux varies over the vehicle surface with stagnation points experiencing the highest values. Accurate prediction is fundamental to thermal protection engineering.}{catAD}
\glsadd{Convective-Heating}\dictentry{Convective Heating}{Convective Heating transfers thermal energy from hot boundary-layer gases to a vehicle surface through molecular conduction and turbulent mixing. It dominates at moderate re-entry velocities below about ten km per second. Laminar convective heating is lower and more predictable than turbulent. Controlling boundary layer transition is essential for managing convective heat loads.}{catAD}
\glsadd{Radiative-Heating}\dictentry{Radiative Heating}{Radiative Heating transfers thermal energy by radiation from the high-temperature shock layer ahead of the vehicle. It becomes dominant at hypersonic velocities above ten km per second. It is particularly severe for Mars return and planetary entry at very high speeds. Accurate prediction requires detailed models of high-temperature gas radiation properties.}{catAD}
\glsadd{Catalytic-Heating}\dictentry{Catalytic Heating}{Catalytic Heating occurs when atomic oxygen and nitrogen recombine on a vehicle surface, releasing chemical energy as additional heat. The amount depends on the surface material's recombination efficiency. High-catalicity surfaces experience significantly elevated total heat flux. Low-catalicity coatings are an important strategy for reducing thermal loads.}{catAD}
\glsadd{Thermal-Protection-Material}\dictentry{Thermal Protection Material}{Thermal Protection Materials withstand extreme heat loads during re-entry or hypersonic flight. Options include ablative composites, reusable ceramic tiles, and carbon-carbon composites, each suited to specific temperatures and durations. Material selection depends on peak heating, total heat load, and reusability needs. Research focuses on lighter, more durable materials for next-generation vehicles.}{catAD}
\glsadd{Ablative-TPS}\dictentry{Ablative TPS}{An Ablative TPS manages heat by sacrificing material through controlled surface ablation. The material chars, melts, and vaporizes, carrying energy away from the structure. Ablators are lightweight and effective for single-use high-heat-flux applications like ICBM vehicles and planetary probes. The drawback is that each system is consumed and cannot be reused.}{catAD}
\glsadd{Reusable-TPS}\dictentry{Reusable TPS}{A Reusable TPS survives multiple re-entry cycles using ceramic tiles, fibrous blankets, or metallic panels. These systems radiate absorbed heat back to space between flights. Reusable TPS was developed for the Space Shuttle and is essential for viable reusable launch vehicles. The challenge is maintaining integrity through repeated thermal cycling.}{catAD}
\glsadd{Insulation}\dictentry{Insulation}{Insulation in a solid rocket motor is a layer of thermally protective material bonded to the inside of the motor case or applied to the grain surface. It prevents the case from reaching critical temperatures during the burn and protects bond lines from hot combustion gases. Insulation materials ablate and char during operation, forming a protective layer. Proper insulation design is essential for case integrity and crew safety in launch vehicles.}{catAD}
\glsadd{Thermal-Barrier-Coating}\dictentry{Thermal Barrier Coating}{A Thermal Barrier Coating is a thin ceramic layer applied to metallic engine components to reduce heat transfer. It enables turbine blades to operate at gas temperatures exceeding the metal's melting point. The coating creates a steep temperature gradient across its thickness. It is essential for achieving high engine efficiency and component durability.}{catAD}
\glsadd{Hot-Structure}\dictentry{Hot Structure}{A Hot Structure is a design approach where the primary load-bearing structure is permitted to reach high temperatures during flight. By eliminating heavy insulation, it reduces vehicle mass. It requires materials with excellent high-temperature strength such as superalloys or ceramic matrix composites. Hot structures are used in hypersonic vehicles and rocket nozzle extensions.}{catAD}
\glsadd{Cold-Structure}\dictentry{Cold Structure}{A Cold Structure uses external insulation to protect the load-bearing structure, keeping it at low operating temperatures. This allows conventional materials like aluminum with well-understood properties. The trade-off is added insulation weight and attachment hardware. Cold structures are common in re-entry backshells and lower-temperature hypersonic vehicle regions.}{catAD}
\glsadd{Nose-Cap}\dictentry{Nose Cap}{A Nose Cap is the forward-most component of a re-entry vehicle, positioned at the stagnation point of highest heating. It experiences the most intense temperatures and pressures during atmospheric entry. Nose caps are typically reinforced carbon-carbon or ablative materials. Their design is one of the most demanding thermal protection challenges.}{catAD}
\glsadd{Leading-Edge-TPS}\dictentry{Leading Edge TPS}{Leading Edge TPS protects the sharp leading edges of wings and fins on re-entry and hypersonic vehicles. These edges experience intense localized heating from shock-boundary layer interactions. The system must withstand high heat fluxes while preserving the aerodynamic contour. Reinforced carbon-carbon is the typical material for reusable vehicle leading edges.}{catAD}
\glsadd{Tiles}\dictentry{Tiles}{Tiles are rigid ceramic insulating blocks used as reusable thermal protection on vehicles like the Space Shuttle. They are made from high-purity silica fibers, giving extremely low density and excellent insulation. Each tile has a borosilicate glass coating for waterproofing and emissivity control. They are individually shaped and bonded to the structure.}{catAD}
\glsadd{Reinforced-Carbon-Carbon}\dictentry{Reinforced Carbon Carbon}{Reinforced Carbon Carbon is a composite of carbon fibers in a carbon matrix, withstanding temperatures up to three thousand degrees Fahrenheit. It maintains strength where metals melt and ceramics degrade. It is used for nose caps and wing leading edges on re-entry vehicles. An oxidation-resistant coating prevents degradation in atomic oxygen environments.}{catAD}
\glsadd{Flexible-TPS}\dictentry{Flexible TPS}{Flexible TPS is a lightweight, conformable material made from woven ceramic or polymer fibers. Unlike rigid tiles, it can be draped over complex curved surfaces, simplifying installation and maintenance. It suits areas with moderate heating on re-entry vehicles. Flexible TPS reduces weight and improves maintainability compared to rigid tile systems.}{catAD}
\glsadd{Inflatable-TPS}\dictentry{Inflatable TPS}{An Inflatable TPS is a deployable system using a flexible aeroshell that inflates before entry to provide drag and thermal shielding. It creates a much larger drag area than a rigid shield of equivalent mass, reducing peak heating. This technology is valuable for delivering payloads to planetary surfaces. It is a promising approach for Mars cargo and crew missions.}{catAD}
\glsadd{Aeroshell}\dictentry{Aeroshell}{An Aeroshell is the protective outer shell of a spacecraft providing aerodynamic deceleration and thermal protection during entry. It consists of a forward heat shield and an aft backshell enclosing the payload. Aeroshells are used for planetary entry where parachutes alone cannot decelerate sufficiently. The design balances drag area, thermal protection mass, and structural integrity.}{catAD}
\glsadd{Heat-Shield}\dictentry{Heat Shield}{A heat shield is a thermal protection system that dissipates the intense heat generated during atmospheric entry. It is typically made of ablative materials such as phenolic resin or reinforced carbon-carbon that char and carry heat away. Heat shields protect spacecraft and crew from temperatures exceeding several thousand degrees. Their design depends on entry velocity, vehicle geometry, and mission reusability requirements.}{catAD}
\glsadd{Buoyancy}\dictentry{Buoyancy}{Buoyancy is the upward force a fluid exerts on an immersed body, equal to the weight of displaced fluid per Archimedes' principle. In aerodynamics it underpins lighter-than-air vehicles like balloons and airships. The lifting force equals the density difference between lifting gas and ambient air times displaced volume and gravity. It enables long-endurance, low-speed flight for surveillance and communications.}{catAD}
\glsadd{Clean-Configuration}\dictentry{Clean Configuration}{Clean Configuration is the aircraft state with all high-lift devices and landing gear fully retracted. This minimizes parasitic drag and allows maximum speed and fuel efficiency. Clean configuration is used during cruise, climb, and high-speed flight. It contrasts with dirty configurations where devices are extended for takeoff and landing.}{catAD}
\glsadd{Decalage}\dictentry{Decalage}{Decalage is the angular difference between the chord lines of upper and lower wings on a biplane. Positive decalage gives the upper wing a higher incidence, generating more lift. It is used to tailor stall characteristics so the lower wing stalls first, preserving aileron effectiveness. It is a fundamental parameter for biplane longitudinal stability.}{catAD}
\glsadd{Disk-Loading}\dictentry{Disk Loading}{Disk Loading is the average pressure change across a propeller or rotor disk, defined as thrust divided by disk area. Low disk loading characterizes large helicopter rotors and high-lift propellers. High disk loading typifies small high-speed propellers and ducted fans. It directly affects induced velocity, efficiency, and noise.}{catAD}
\glsadd{Enstrophy}\dictentry{Enstrophy}{Enstrophy measures the integrated square of vorticity magnitude in a flow, representing rotational motion intensity. It is directly related to the rate of kinetic energy dissipation into heat through viscous action. High enstrophy indicates strong, concentrated vortical activity such as turbulent mixing. It is fundamental for understanding turbulence cascade processes.}{catAD}
\glsadd{Keel-Effect}\dictentry{Keel Effect}{The Keel Effect is the directional stability provided by lateral area aft of the center of gravity, analogous to a ship's keel. In aircraft it is primarily the vertical tail and aft fuselage generating restoring side forces during sideslip. It produces weathercock stability, aligning the nose into the relative wind. Adequate keel area is essential for preventing Dutch roll.}{catAD}
\glsadd{Rogallo-Wing}\dictentry{Rogallo Wing}{A Rogallo Wing is a flexible delta-shaped lifting surface of two conical fabric panels held taut by a rigid frame. Invented by Francis Rogallo, it folds for storage and deploys rapidly. It is used in hang gliders, parawings, and spacecraft recovery systems where low weight matters. The design offers simplicity and reliability for soaring and capsule recovery.}{catAD}
\glsadd{Viscosity}\dictentry{Viscosity}{Viscosity measures a fluid's internal resistance to shear deformation, arising from intermolecular forces and momentum transfer between layers. In aerodynamics it governs boundary layer behavior, skin friction drag, and laminar-to-turbulent transition. It also controls flow separation affecting lift and drag. Understanding viscosity is fundamental to predicting aerodynamic performance.}{catAD}
\glsadd{Biplane}\dictentry{Biplane}{A biplane is a fixed-wing aircraft with two stacked wings connected by struts and bracing wires. The dual-wing arrangement generates greater lift at lower speeds than a single equivalent wing. Biplanes dominated early aviation from the Wright Flyer through World War I. The structural bracing creates substantial parasitic drag, leading to their replacement by monoplanes as engine power advanced.}{catAD}
\glsadd{Coanda-Effect}\dictentry{Coanda Effect}{The Coanda Effect is the tendency of a fluid jet to follow a nearby curved surface rather than continuing straight. When moving air encounters a convex surface, the flow attaches and bends around the contour due to boundary layer pressure differences. This phenomenon is exploited in boundary layer control, circulation control wings, and thrust vectoring nozzles. Engineers must account for it on curved aerodynamic surfaces.}{catAD}
\glsadd{Fineness-Ratio}\dictentry{Fineness Ratio}{Fineness Ratio is the ratio of an aerodynamic body's maximum diameter to its overall length. A higher ratio produces a more slender form, generally reducing wave drag at supersonic speeds but increasing friction drag. For fuselages, the optimal ratio balances structural volume against aerodynamic efficiency. Typical subsonic fuselages have ratios between 8 and 12, while supersonic designs may use 15 or higher.}{catAD}
\glsadd{Lift-to-Drag-Ratio}\dictentry{Lift-to-Drag Ratio}{The lift-to-drag ratio is a dimensionless measure of aerodynamic efficiency, defined as lift force divided by drag force. Higher values indicate more efficient performance, generating more lift for a given drag level. Commercial airliners cruise at ratios around 15 to 20, while high-performance gliders can exceed 40. Maximizing this ratio is central to wing design, affecting fuel economy, range, and payload.}{catAD}
\glsadd{Moment-Coefficient}\dictentry{Moment Coefficient}{The moment coefficient is a dimensionless parameter quantifying aerodynamic moment about a reference point on an airfoil or aircraft. It is calculated by dividing moment by the product of dynamic pressure, reference area, and a reference length. These coefficients vary with angle of attack, Mach number, and Reynolds number, forming the core of stability analysis. Designers use them for control surface sizing.}{catAD}
\glsadd{Shock-Boundary-Layer-Interaction}\dictentry{Shock-Boundary Layer Interaction}{Shock-Boundary Layer Interaction occurs when a shock wave impinges on or forms within a boundary layer at supersonic or transonic speeds. The abrupt pressure rise can cause separation, increasing drag and causing buffet. This is significant in transonic wing design where local supersonic pockets terminate in normal shocks. Managing this interaction through airfoil shaping is critical near Mach 1.}{catAD}
\glsadd{Stability-Derivative}\dictentry{Stability Derivative}{Stability derivatives are partial derivatives of aerodynamic forces and moments with respect to flow parameter changes such as angle of attack and angular rates. These coefficients quantify how loads respond to small perturbations from trimmed conditions. They form the core of linearized equations of motion in flight dynamics and control design. Positive derivatives in certain combinations indicate natural stability.}{catAD}
\glsadd{Aerodynamic-Balancing}\dictentry{Aerodynamic Balancing}{Aerodynamic balancing positions a control surface hinge line so a portion extends forward of the pivot, creating a tab effect. The forward segment experiences loads counteracting hinge moments on the main portion. This reduces pilot control force or actuator requirements. It is fundamental to designing ailerons, elevators, and rudders with acceptable control forces throughout the flight envelope.}{catAD}
\glsadd{Aileron-Trim}\dictentry{Aileron Trim}{Aileron trim adjusts the neutral aileron position to eliminate sustained lateral inputs needed to maintain wings-level flight. Trim devices such as tabs relieve pilot forces caused by asymmetric loading, fuel imbalance, or crosswinds. Proper trim reduces pilot fatigue during extended cruise. Most modern aircraft incorporate electric or hydraulic trim systems integrated with the autopilot.}{catAD}
\glsadd{Body-Force}\dictentry{Body Force}{A body force is an external force acting uniformly throughout a fluid volume rather than at its surface. Gravity and electromagnetic forces are common examples in aerodynamics, driving natural convection effects. In computational fluid dynamics, body force terms are incorporated into the Navier-Stokes equations. Unlike surface forces such as pressure and shear, body forces scale with the fluid element's mass.}{catAD}
\glsadd{Boundary-Layer-Suction}\dictentry{Boundary Layer Suction}{Boundary layer suction removes low-energy fluid from the boundary layer through porous surfaces or slots. By extracting slow-moving wall fluid, it delays separation and can maintain laminar flow over extended areas. The technique has been applied on research aircraft, high-lift systems, and turbine blades. The energy cost must be weighed against drag reduction benefits for each application.}{catAD}
\glsadd{Chordwise-Distribution}\dictentry{Chordwise Distribution}{Chordwise distribution describes how aerodynamic properties vary from leading edge to trailing edge of an airfoil. The pressure distribution determines peak suction location, pressure gradient extent, and likelihood of transition or separation. Designers shape profiles to achieve desired loadings such as flat-top or rear-loaded distributions. Wind tunnel measurements and computational tools characterize and optimize these.}{catAD}
\glsadd{Control-Volume}\dictentry{Control Volume}{A control volume is a fixed or moving region through which fluid flows, used for applying conservation laws. Mass, momentum, and energy balances are formulated across boundaries to derive integral flow relationships. This analysis underpins fundamental aerodynamic equations including thrust calculations and Bernoulli derivations. The approach simplifies complex fields by focusing on aggregate quantities at boundaries.}{catAD}
\glsadd{Convergent-Divergent-Duct}\dictentry{Convergent-Divergent Duct}{A convergent-divergent duct first narrows to a throat then expands, enabling acceleration of compressible fluid from subsonic to supersonic velocities. The convergent section accelerates subsonic flow, while the divergent section further accelerates flow reaching sonic conditions at the throat. This geometry is essential in supersonic nozzles, rocket engines, and intake diffusers. The area ratio determines exit Mach number.}{catAD}
\glsadd{Cruise-Efficiency}\dictentry{Cruise Efficiency}{Cruise efficiency relates useful propulsive work to total energy expended during cruise flight. It encompasses aerodynamic efficiency, propulsive efficiency, and structural weight optimization. Maximizing it reduces fuel consumption per unit distance, the primary driver of operating costs. Modern aircraft achieve efficiencies above 80 percent through advanced airframes, high-bypass engines, and optimized altitudes.}{catAD}
\glsadd{Direct-Simulation-Monte-Carlo}\dictentry{Direct Simulation Monte Carlo}{Direct Simulation Monte Carlo is a computational method for modeling rarefied gas dynamics by tracking representative molecule trajectories and collisions. Unlike Navier-Stokes solvers, it explicitly simulates molecular motion for flows where mean free path is comparable to characteristic length. It is widely used for hypersonic re-entry and spacecraft plume modeling. Accuracy improves with particle count but demands substantial resources.}{catAD}
\glsadd{Drag-Bucket}\dictentry{Drag Bucket}{The drag bucket is a region in a laminar flow airfoil's drag polar where section drag remains near its minimum over a moderate lift coefficient range. This low-drag plateau persists while laminar flow prevails across the surface. Supercritical and natural laminar airfoils are designed to provide wide drag buckets. Once outside, transition and separation cause a sharp drag increase.}{catAD}
\glsadd{Drag-Divergence}\dictentry{Drag Divergence}{Drag divergence is the rapid aerodynamic drag increase as an aircraft approaches and exceeds the speed of sound, driven by shock wave formation introducing wave drag. The drag divergence Mach number marks where drag rises steeply, typically by 20 counts. Designers employ sweep, area ruling, and supercritical airfoils to push this number higher. It is a key parameter in transonic aircraft design.}{catAD}
\glsadd{Effective-Aspect-Ratio}\dictentry{Effective Aspect Ratio}{Effective aspect ratio is the aspect ratio of a finite wing producing the same induced drag as the actual configuration, accounting for sweep, end plates, or winglet effects. It provides a single parameter characterizing lift-dependent drag of complex planforms. A higher value corresponds to lower induced drag at a given lift coefficient. Engineers use it to compare designs without detailed vortex lattice calculations.}{catAD}
\glsadd{Elastic-Axis}\dictentry{Elastic Axis}{The elastic axis of a wing is the locus where applied transverse force produces bending without torsional deformation. This structural reference line is critical in aeroelastic analysis because the distance between the aerodynamic center and elastic axis determines flutter-driving torsional moments. If the aerodynamic center lies aft of it, the wing tends to be torsionally stable. Precise knowledge is essential for flutter prediction.}{catAD}
\glsadd{Equivalent-Body-of-Revolution}\dictentry{Equivalent Body of Revolution}{An equivalent body of revolution is a hypothetical axisymmetric shape producing the same wave drag as a given non-axisymmetric body in supersonic flow. This allows engineers to estimate wave drag of complex shapes using simpler one-dimensional area distribution analysis. The cross-sectional area distribution along the longitudinal axis governs wave drag. Designers use this to guide area ruling for minimum drag.}{catAD}
\glsadd{Fin-Flutter}\dictentry{Fin Flutter}{Fin flutter is a self-excited aeroelastic oscillation of an aircraft vertical tail occurring when aerodynamic forces, elasticity, and inertial effects couple to produce unstable vibrations. The onset is defined by a critical dynamic pressure above which oscillations grow, potentially causing failure in seconds. Mass distribution, stiffness, and component positions influence the boundary. Designers address it through stiffening or mass balancing.}{catAD}
\glsadd{Flow-visualization}\dictentry{Flow visualization}{Flow visualization encompasses techniques rendering fluid flow patterns visible to observers. Physical methods include smoke injection, oil-film techniques, schlieren imaging, and particle image velocimetry. Computational visualization transforms simulation data into contour plots, streamlines, and vector fields. These techniques are indispensable for validating designs, identifying separation zones and vortices, and communicating phenomena.}{catAD}
\glsadd{Freestream}\dictentry{Freestream}{The freestream refers to undisturbed flow conditions far upstream of an aerodynamic body, where velocity, pressure, temperature, and density are uniform. Freestream conditions define the baseline reference from which all aerodynamic coefficients are measured. In wind tunnel testing, achieving uniform freestream with minimal turbulence is critical. Flight tests similarly reference ambient atmospheric conditions as the freestream state.}{catAD}
\glsadd{Hypersonic-Boundary-Layer}\dictentry{Hypersonic Boundary Layer}{A hypersonic boundary layer is the thin viscous region adjacent to a surface where local Mach number exceeds approximately five, characterized by extreme compressibility and high-temperature gas effects. Aerodynamic heating causes molecular vibration, dissociation, and ionization of air species. These real-gas effects profoundly alter velocity and thermal profiles. Accurate prediction is essential for thermal protection on re-entry vehicles.}{catAD}
\glsadd{Induced-Drag-Factor}\dictentry{Induced Drag Factor}{The induced drag factor is a dimensionless coefficient in the induced drag equation accounting for actual spanwise lift distribution relative to ideal elliptical loading. For perfect elliptical loading, the factor equals one, producing theoretical minimum induced drag. Non-elliptical distributions yield values greater than one. Wing planform shape, twist, and non-planar configurations all influence this factor during optimization.}{catAD}
\glsadd{Interference-Coefficient}\dictentry{Interference Coefficient}{An interference coefficient quantifies additional aerodynamic force on one component due to another in a multi-body configuration. For example, wing-fuselage interaction modifies local flow fields beyond isolated component behavior. These coefficients are determined through wind tunnel testing of complete versus isolated configurations. Understanding interference is critical for accurate drag, lift, and moment prediction with external stores.}{catAD}
\glsadd{Isentropic-Flow}\dictentry{Isentropic Flow}{Isentropic flow is an idealized process where fluid entropy remains constant, implying adiabatic and reversible conditions. Stagnation properties are preserved along streamlines, and relationships between local Mach number and thermodynamic variables follow well-defined equations. These relations are fundamental to compressible aerodynamics and used in nozzle design and performance estimation. Real flows deviate due to viscosity and shocks.}{catAD}
\glsadd{Laminar-Bucket}\dictentry{Laminar Bucket}{The laminar bucket is a pronounced low-drag region in a laminar flow airfoil's drag versus lift curve where section drag remains near its minimum. Within this bucket, laminar flow prevails across most of the surface. The bucket is bounded by lift coefficients where transition or separation causes sudden drag increase. Airfoil design optimizing these buckets is central to low-drag performance for gliders and fuel-efficient airliners.}{catAD}
\glsadd{Laminar-Separation-Bubble}\dictentry{Laminar Separation Bubble}{A laminar separation bubble forms when a laminar boundary layer separates due to an adverse pressure gradient and transitions to turbulent flow downstream. The separated shear layer undergoes rapid transition, producing a recirculating region. These bubbles increase profile drag and alter effective aerodynamic shape at low Reynolds numbers. They are common on unmanned vehicles, helicopter blades, and wind turbine airfoils.}{catAD}
\glsadd{Lifting-Body}\dictentry{Lifting Body}{A lifting body is an aircraft or spacecraft where the fuselage itself generates significant lift, reducing the need for conventional wings. The body is contoured with shaped cross sections and camber for useful lift-to-drag ratios. Lifting body designs were explored extensively during early re-entry research, including NASA's M2-F2 and HL-10. Modern interest has revived for reusable launch vehicle and hypersonic transport concepts.}{catAD}
\glsadd{Mach-Disk}\dictentry{Mach Disk}{A Mach disk is a normal shock wave forming in the exhaust plume of a supersonic nozzle at overexpanded conditions, where exit pressure is significantly below ambient. Flow passes through converging oblique shocks and terminates at the disk, decelerating subsonically. Disk location and strength depend on nozzle pressure ratio and overexpansion degree. Mach disks are commonly observed in rocket exhaust plumes.}{catAD}
\glsadd{Mangler-Transform}\dictentry{Mangler Transform}{The Mangler transform is a coordinate transformation converting three-dimensional axisymmetric boundary layer equations into equivalent two-dimensional form. By rescaling coordinates appropriately, the transformed equations admit solutions using established 2D methods. This simplifies analysis of boundary layers on bodies of revolution such as fuselages and nose cones. The transform is valid for thin boundary layers relative to body radius.}{catAD}
\glsadd{Mass-Flow-Ratio}\dictentry{Mass Flow Ratio}{The mass flow ratio is the ratio of actual mass flow captured by an inlet to the flow captured by a freestream tube of the same area. A ratio of unity indicates perfect capture with no spillage. At supersonic speeds, the ratio depends strongly on inlet geometry, shock structure, and angle of attack. Inlet design aims to maximize this ratio across the operating envelope.}{catAD}
\glsadd{Mean-Camber-Line}\dictentry{Mean Camber Line}{The mean camber line is the locus of points equidistant between an airfoil's upper and lower surfaces, perpendicular to the chord. This curve defines the fundamental camber shape, determining zero-lift angle and pitching moment. A symmetric airfoil has its mean camber line along the chord, while a cambered one generates lift at zero angle of attack. Design begins with specifying this distribution and adding thickness.}{catAD}
\glsadd{Minimum-Drag}\dictentry{Minimum Drag}{Minimum drag is the lowest total aerodynamic drag achievable at a given flight condition, representing the bottom of the drag curve. It results from optimal balance between parasitic components and induced drag at a specific lift coefficient. For an airfoil, minimum drag occurs near zero lift where the boundary layer is thinnest. Aircraft typically cruise near the lift coefficient for maximum lift-to-drag ratio.}{catAD}
\glsadd{Mixed-Boundary-Layer}\dictentry{Mixed Boundary Layer}{A mixed boundary layer contains both laminar and turbulent regions, typically with laminar flow near the leading edge transitioning downstream. The transition location depends on turbulence, roughness, pressure gradient, and Reynolds number. In many applications, the boundary layer is transitional over significant surface portions. Predicting transition and resulting characteristics is essential for drag estimation and heat transfer analysis.}{catAD}
\glsadd{Net-Thrust}\dictentry{Net Thrust}{Net thrust is the useful propulsive force from a jet or rocket engine after accounting for momentum drag from ingested inlet air. It equals gross thrust minus the additive drag from the freestream ram effect. For subsonic turbojets and turbofans, net thrust decreases with increasing speed due to growing momentum drag. Net thrust is the figure of merit for aircraft performance and mission analysis.}{catAD}
\glsadd{Non-Inertial-Reference-Frame}\dictentry{Non-Inertial Reference Frame}{A non-inertial reference frame is a coordinate system accelerating relative to an inertial frame, requiring fictitious forces like Coriolis and centrifugal forces. In aerodynamics, aircraft-fixed frames are non-inertial during maneuvering, and equations of motion must include these apparent forces. Meteorological analyses also employ rotating Earth-fixed frames. Proper treatment is essential for accurate flight dynamics modeling and navigation.}{catAD}
\glsadd{Normal-Force}\dictentry{Normal Force}{Normal force is the aerodynamic force component acting perpendicular to the chord line or longitudinal axis of an aircraft or airfoil. It is the primary contributor to lift at small angles of attack and relates directly to surface pressure distribution. The normal force coefficient is used in stability analyses to separate perpendicular and axial contributions. For missiles at high angles it is often more significant than lift-drag decomposition.}{catAD}
\glsadd{Oblique-Flow}\dictentry{Oblique Flow}{Oblique flow refers to fluid approaching a surface at an angle to the principal reference direction, such as crossflow over a swept wing. The component perpendicular to the leading edge drives local aerodynamic forces, while the lateral component modifies boundary layer behavior. Oblique conditions are pervasive as aircraft frequently operate at sideslip angles. Understanding it is essential for swept wing and wing-body junction design.}{catAD}
\glsadd{Open-Jet-Tunnel}\dictentry{Open Jet Tunnel}{An open jet wind tunnel discharges airflow as a free jet into a large room rather than confining it within walls. Models are placed within the jet, and the absence of wall interference eliminates complex correction procedures. These tunnels are useful for powered model testing and propeller noise measurements. The trade-off is susceptibility to jet boundary instabilities and entrainment effects.}{catAD}
\glsadd{Pitching-Moment-Coefficient}\dictentry{Pitching Moment Coefficient}{The pitching moment coefficient is a dimensionless parameter quantifying aerodynamic moment about the lateral pitch axis, normalized by dynamic pressure, reference area, and chord. It determines an aircraft's tendency to nose up or down and is central to longitudinal stability. A negative slope with angle of attack indicates stability for conventional configurations. Pilots must trim this moment to zero for steady flight.}{catAD}
\glsadd{Planform}\dictentry{Planform}{The planform is a wing's shape as viewed from above, defined by leading edge, trailing edge, and tip outlines. Common shapes include rectangular, tapered, elliptical, swept, and delta, each with distinct aerodynamic trade-offs. The planform directly determines spanwise lift distribution and induced drag. Modern computational tools optimize planforms using inverse design and evolutionary algorithms.}{catAD}
\glsadd{Pressure-Recovery}\dictentry{Pressure Recovery}{Pressure recovery is the process by which static pressure increases as flow decelerates, converting kinetic energy back into pressure. In subsonic diffusers and inlets, it reflects efficiency in converting dynamic pressure into static rise for the engine. Higher recovery means higher-pressure air reaching the engine, improving thrust. At supersonic speeds, pressure recovery is severely impacted by shock losses.}{catAD}
\glsadd{Pressure-Surface}\dictentry{Pressure Surface}{The pressure surface of an airfoil is the side where local static pressure exceeds freestream pressure, producing a force into the surface. On a positively cambered airfoil at positive angle, the lower surface is the pressure surface while the upper experiences suction. The net difference between suction and pressure distributions produces lift. Pressure surface design influences pitching moment and high-speed performance.}{catAD}
\glsadd{Ramp-Angle}\dictentry{Ramp Angle}{Ramp angle is the angle of an inclined surface relative to oncoming flow, commonly used for supersonic inlet ramps and control surface deflections. In supersonic inlets, oblique shocks from ramp surfaces progressively decelerate flow before a final normal shock. The ramp angle determines each shock's strength and must be selected for the design Mach number. It influences aerodynamic effectiveness and hinge moments.}{catAD}
\glsadd{Rankine-Hugoniot-Relations}\dictentry{Rankine-Hugoniot Relations}{The Rankine-Hugoniot relations are fundamental equations describing property jumps across a shock wave in compressible flow. Derived from conservation of mass, momentum, and energy, they predict that normal shock flow always decelerates from supersonic to subsonic with increased pressure, density, and entropy. They form the analytical basis for all shock calculations in compressible aerodynamics and are essential for nozzle and inlet analysis.}{catAD}
\glsadd{Reduced-Frequency}\dictentry{Reduced Frequency}{Reduced frequency is a dimensionless parameter measuring flow unsteadiness, defined as the ratio of characteristic flow time to motion time. For an oscillating airfoil, it equals angular frequency times half-chord divided by freestream velocity. At values near zero, flow is effectively quasi-steady. As it increases, unsteady effects like delayed pressure response become significant, requiring unsteady aerodynamic analysis.}{catAD}
\glsadd{Reference-Area}\dictentry{Reference Area}{Reference area is the characteristic area used to non-dimensionalize aerodynamic coefficients. For wings, planform area is conventional; for fuselages, maximum cross-sectional area is used. Full aircraft coefficients typically reference wing planform area. The choice does not affect physical forces but is critical for consistent coefficient reporting and comparing data from different sources.}{catAD}
\glsadd{Reynolds-Number-Scaling}\dictentry{Reynolds Number Scaling}{Reynolds number scaling adjusts aerodynamic data for differences between model-scale wind tunnel tests and full-scale flight. Since Reynolds number governs boundary layer transition, turbulence, and separation, coefficients at low values may differ substantially from flight. Correction methods include empirical criteria and computational extrapolation. Accurate scaling is essential for translating wind tunnel results into reliable predictions.}{catAD}
\glsadd{Roll-Rate}\dictentry{Roll Rate}{Roll rate is the angular velocity of an aircraft about its longitudinal axis, measured in degrees or radians per second. It is produced by asymmetric lift from ailerons and is a primary measure of lateral maneuverability. Achievable roll rate depends on aileron effectiveness, roll damping, moment of inertia, and control authority. Excessive roll rate can cause pilot-induced oscillation or structural loads.}{catAD}
\glsadd{Roughness-Element}\dictentry{Roughness Element}{A roughness element is a small surface protrusion that disrupts the laminar boundary layer and can trigger premature transition. Engineered elements are placed on surfaces to study transition or provide controlled trips. Natural roughness from manufacturing, insects, or ice similarly affects behavior. Element height, spacing, and shape relative to boundary layer thickness determine influence on transition and skin friction.}{catAD}
\glsadd{Schlieren-System}\dictentry{Schlieren System}{A schlieren system is an optical visualization technique rendering density gradients visible by exploiting light refraction through varying refractive index. The setup uses collimated light, a test section, and a knife-edge at the focal point to selectively block deviated rays. Schlieren imaging is indispensable for visualizing shock waves, expansion fans, and boundary layers in compressible flow research.}{catAD}
\glsadd{Scoop}\dictentry{Scoop}{An aerodynamic scoop is an inlet designed to capture and direct airflow into a system such as an engine, heat exchanger, or cooling duct. The geometry is shaped for efficient pressure recovery while minimizing drag and flow distortion. Internal design must account for boundary layer displacement, spillage drag, and downstream pressure recovery. Scoops are common on aircraft for oil cooling and race cars for cooling.}{catAD}
\glsadd{Section-Drag-Coefficient}\dictentry{Section Drag Coefficient}{The section drag coefficient is a dimensionless measure of two-dimensional drag per unit span of an airfoil, referenced to dynamic pressure and chord. It quantifies combined skin friction, pressure, and wave drag effects. This coefficient plotted against section lift coefficient forms the two-dimensional drag polar, the fundamental building block for aircraft drag buildup.}{catAD}
\glsadd{Section-Lift-Coefficient}\dictentry{Section Lift Coefficient}{The section lift coefficient is a dimensionless parameter characterizing lift per unit span of a two-dimensional airfoil, normalized by dynamic pressure and chord. It depends on airfoil shape, camber, and angle of attack, varying linearly in attached flow before reaching maximum at stall. The curve slope is approximately two pi per radian for thin airfoils. These coefficients are fundamental to wing aerodynamic analysis.}{catAD}
\glsadd{Separation-Bubble}\dictentry{Separation Bubble}{A separation bubble is a localized recirculating flow region formed when a boundary layer separates and reattaches downstream. The bubble creates a virtual aerodynamic surface altering the pressure distribution and can trigger early transition. Small leading-edge bubbles may be benign and promote transition to a more separation-resistant layer. Large or bursting bubbles cause abrupt lift loss and increased drag.}{catAD}
\glsadd{Shock-Expansin-Theory}\dictentry{Shock Expansin Theory}{Shock-expansion theory computes pressure distributions on supersonic surfaces by applying oblique shock relations at corners and Prandtl-Meyer expansion fans along convex surfaces. It constructs the complete pressure and Mach number field by tracing characteristic lines. The method provides more accurate results than linearized theory for moderate compressibility effects. It is used for preliminary supersonic airfoil and nozzle design.}{catAD}
\glsadd{Slat-Track}\dictentry{Slat Track}{A slat track is the structural guide rail along which a leading-edge slat translates between stowed and deployed positions. The track defines the deployment path, extension distance, and angle relative to the wing's leading edge. Tracks use rollers or slide bearings for smooth operation under aerodynamic loads. The geometry is optimized to maximize lift increment and delay stall while minimizing retracted drag.}{catAD}
\glsadd{Slipstream}\dictentry{Slipstream}{The slipstream is the accelerated, rotating air column passing through and behind a propeller during operation. It has higher velocity and lower pressure than ambient air, impinging on tail surfaces and fuselage. In single-engine aircraft, the spiraling slipstream creates asymmetric airflow producing a yawing moment. Understanding slipstream behavior is essential for tail surface sizing and propeller-wing interaction analysis.}{catAD}
\glsadd{Span-Loading}\dictentry{Span Loading}{Span loading is the distribution of aerodynamic lift along the wingspan, typically expressed as lift per unit span. It directly determines induced drag, with elliptical distribution providing minimum drag for a given total lift and span. Practical designs approach elliptical loading through taper, washout twist, and flap scheduling. These distributions are measured in flight tests or estimated computationally.}{catAD}
\glsadd{Spanwise-Flow}\dictentry{Spanwise Flow}{Spanwise flow is the velocity component running along the wingspan direction, perpendicular to chordwise flow. On swept wings, boundary layer spanwise momentum can lead to premature tip separation known as tip stall. It is also significant on delta wings where leading-edge vortices drive flow across the upper surface. Managing it through fences, vortex generators, and twist is essential for uniform loading.}{catAD}
\glsadd{Speed-Brake}\dictentry{Speed Brake}{A speed brake is a movable panel that deploys to increase aerodynamic drag and reduce airspeed without significantly altering lift. Unlike spoilers which primarily reduce lift, speed brakes maximize drag for rapid deceleration during descent or rollout. They are typically on fuselage or wing surfaces, actuated hydraulically. They are valuable in military aircraft for energy management and commercial aircraft for glideslope control.}{catAD}
\glsadd{Subsonic-Leading-Edge}\dictentry{Subsonic Leading Edge}{A subsonic leading edge is a rounded leading edge whose local flow remains below sonic even when freestream Mach number is transonic. The rounded shape allows smooth acceleration without a detached shock. Subsonic leading edges are characteristic of conventional airfoils for subsonic cruise. At supersonic freestream conditions, they generate an attached bow shock weaker than a blunt nose detached shock.}{catAD}
\glsadd{Supersonic-Leading-Edge}\dictentry{Supersonic Leading Edge}{A supersonic leading edge is a sharp leading edge whose local flow exceeds sonic speed, resulting in an attached oblique shock. The sharp geometry forces the flow through an oblique shock rather than a detached bow shock, significantly reducing wave drag. These leading edges are defining features of thin, highly swept wings on supersonic aircraft. The sweep must exceed the Mach cone angle for the shock to remain attached.}{catAD}
\glsadd{Supersonic-Area-Rule}\dictentry{Supersonic Area Rule}{The supersonic area rule states that wave drag depends on cross-sectional area distribution perpendicular to flow, weighted by an influence function differing from the subsonic case. Application leads to distinctive fuselage shaping such as necking near wing junctions. The rule enables rapid wave drag estimation from complex configurations. It accounts for Mach line interference between components in supersonic flow.}{catAD}
\glsadd{Surface-Pressure-Distribution}\dictentry{Surface Pressure Distribution}{The surface pressure distribution describes how static pressure varies over an aerodynamic body's surface as a function of position. It is presented as pressure coefficient plots or contour maps. The distribution determines integrated forces, identifies pressure gradient regions, and reveals shocks, separation, and transition. Computational fluid dynamics and wind tunnel measurements are the primary tools for obtaining this data.}{catAD}
\glsadd{Sweep-Effect}\dictentry{Sweep Effect}{The sweep effect describes aerodynamic changes from orienting a wing at an angle to oncoming flow, reducing the normal velocity component. Sweeping delays compressibility effects and increases drag divergence Mach number. However, it reduces lift curve slope, promotes spanwise flow and tip stall, and introduces structural coupling. Sweep angle is one of the most influential planform parameters in transonic and supersonic design.}{catAD}
\glsadd{Terminal-Velocity}\dictentry{Terminal Velocity}{Terminal velocity is the constant speed reached when gravitational force is balanced by aerodynamic drag during descent. Once achieved, net force is zero and descent is steady without further acceleration. The value depends on mass, frontal area, drag coefficient, and fluid density. Terminal velocity calculations are important for parachute design, atmospheric entry analysis, and high-altitude instrumentation.}{catAD}
\glsadd{Thermal-Boundary-Layer}\dictentry{Thermal Boundary Layer}{The thermal boundary layer is the thin region adjacent to a heated or cooled surface where fluid temperature transitions from surface to freestream values. It develops with the velocity boundary layer, but thickness depends on thermal diffusivity relative to momentum diffusivity, characterized by the Prandtl number. Accurate prediction is critical for aerodynamic heating loads on re-entry vehicles and heat exchanger design.}{catAD}
\glsadd{Transonic-Dip}\dictentry{Transonic Dip}{The transonic dip is the phenomenon where flutter speed of a wing or control surface decreases in the transonic range before recovering at supersonic speeds. Shock waves alter unsteady pressure distributions, reducing flutter margins. The dip's magnitude depends on thickness, mode coupling, and Reynolds number. Identifying and mitigating it is critical for aeroelastic certification of transonic aircraft.}{catAD}
\glsadd{Two-Dimensional-Flow}\dictentry{Two-Dimensional Flow}{Two-dimensional flow is idealized fluid motion where all properties vary in only two spatial dimensions with no spanwise variation. This applies to infinite-span wings and two-dimensional nozzles where end effects are negligible. The analysis simplifies governing equations, allowing closed-form solutions for problems like flat plate boundary layers. Real three-dimensional effects such as tip vortices are absent.}{catAD}
\glsadd{Velocity-Potential}\dictentry{Velocity Potential}{The velocity potential is a scalar function whose gradient equals the velocity vector in irrotational flow, providing a powerful tool for potential flow analysis. In incompressible flow, it satisfies Laplace's equation, admitting solutions through superposition and conformal mapping. The formulation is the foundation of panel methods and vortex lattice methods for inviscid analysis. It breaks down for rotational or viscous flows.}{catAD}
\glsadd{Viscous-Drag}\dictentry{Viscous Drag}{Viscous drag is the aerodynamic drag component from shear stresses on a surface due to fluid viscosity. It consists primarily of skin friction from tangential stress and form drag from separation-induced pressure imbalance. It dominates at subsonic speeds for streamlined bodies and depends on boundary layer state, surface roughness, and Reynolds number. Reducing it through laminar flow control is a primary objective.}{catAD}
\glsadd{Volume-Coefficient}\dictentry{Volume Coefficient}{The volume coefficient is a dimensionless parameter relating an aerodynamic body's volume to its frontal or reference area, characterizing bluffness. A lower coefficient for a given enclosed volume implies a more slender, efficient shape. The parameter is useful for parametric studies and comparing aerodynamic efficiency of different geometries during preliminary design phases.}{catAD}
\glsadd{Vortex-Core}\dictentry{Vortex Core}{The vortex core is the central region of a concentrated vortex where rotational velocity reaches its maximum and static pressure its minimum. In the Rankine model, the core exhibits solid-body rotation surrounded by a potential vortex. Real cores are viscous regions with finite pressure minima that can cause structural loads and hazards for trailing aircraft. Understanding core structure is important for wake turbulence studies.}{catAD}
\glsadd{Wind-Shear}\dictentry{Wind Shear}{Wind shear is a rapid change in wind speed and/or direction over a short atmospheric distance, producing sudden aerodynamic force changes. It is particularly dangerous during takeoff and landing when aircraft are at low altitude near stall speed. Microbursts, frontal boundaries, and terrain effects are common sources. Modern airliners are equipped with predictive detection systems and pilots train in recovery procedures.}{catAD}
\glsadd{Zero-Lift-Angle}\dictentry{Zero-Lift Angle}{The zero-lift angle of attack is the angle where an airfoil produces zero net lift, negative for cambered sections and zero for symmetric ones. It depends on camber and twist distribution, determining trim conditions. The zero-lift angle sets geometric incidence of wings and tails for desired cruise lift coefficients. It is essential for longitudinal stability analysis and flight condition optimization.}{catAD}
\glsadd{Zero-Lift-Drag}\dictentry{Zero-Lift Drag}{Zero-lift drag is total aerodynamic drag when producing no lift, comprising skin friction, form drag from volume displacement, and interference drag. It represents the parasitic drag floor below which total drag cannot fall. It is estimated by summing component contributions using empirical buildup methods. Minimizing it through streamlining and surface finish is a continuous design and maintenance objective.}{catAD}
\glsadd{Aileron-Lock}\dictentry{Aileron Lock}{An aileron lock is a mechanical device securing ailerons in neutral position when parked or towed. It prevents wind-induced deflections that could damage the control system or create flutter. Locks are typically brightly colored with Remove Before Flight tags. Some aircraft incorporate internal hydraulic locks within actuators serving the same protective function.}{catAD}
\glsadd{Adverse-Pressure-Gradient}\dictentry{Adverse Pressure Gradient}{An adverse pressure gradient occurs when static pressure increases in the flow direction, causing boundary layer fluid to decelerate. As low-energy wall fluid slows, it may stop and reverse, causing separation. Adverse gradients are encountered on aft airfoil portions, diffusers, and behind shocks. Managing them through shaping, vortex generators, and suction is a central aerodynamic design challenge.}{catAD}
\glsadd{Boundary-Layer-Thickness}\dictentry{Boundary Layer Thickness}{Boundary layer thickness is the distance from a surface to where local velocity reaches 99 percent of freestream velocity. It grows along the surface as viscous effects diffuse outward, with rate depending on whether the layer is laminar or turbulent. A thicker layer produces greater displacement effects and higher drag. This fundamental length scale is used in Reynolds number calculations.}{catAD}
\glsadd{Centerline-Pressure}\dictentry{Centerline Pressure}{Centerline pressure is the static pressure measured along a body's geometric centerline or flow field core, used as a reference for internal flows. In wind tunnel testing, centerline surveys provide flow quality calibration data. For nozzles and inlets, the distribution reveals shock train and expansion wave locations. Accurate measurement is essential for validating computational internal flow predictions.}{catAD}
\glsadd{Characteristic-Mach-Angle}\dictentry{Characteristic Mach Angle}{The characteristic Mach angle is the angle between a Mach wave and flow direction, given by the arcsine of the reciprocal of local Mach number. It defines the cone of influence within which disturbances propagate in supersonic flow. As Mach number increases, the angle decreases. The Mach angle is fundamental to the method of characteristics and supersonic inlet and nozzle design.}{catAD}
\glsadd{Compressibility-Correction}\dictentry{Compressibility Correction}{Compressibility correction refers to adjustments applied to incompressible aerodynamic data to account for density changes at higher Mach numbers. The Prandtl-Glauert rule provides a first-order correction increasing lift slope by one over the square root of one minus Mach number squared. More advanced corrections include the Karman-Tsien rule. These allow designers to extrapolate low-speed data into compressible regimes.}{catAD}
\glsadd{Couette-Flow}\dictentry{Couette Flow}{Couette flow describes viscous fluid motion between two parallel plates where one moves relative to the other. The steady, incompressible velocity profile is linear, varying from zero at the stationary wall to plate velocity at the moving wall. This canonical flow provides exact Navier-Stokes solutions and serves as a fundamental model for shear-driven flows and lubrication theory.}{catAD}
\glsadd{Entrance-Length}\dictentry{Entrance Length}{Entrance length is the distance from a duct inlet that flow requires to develop from uniform inlet conditions to a fully developed profile. In laminar pipe flow, it scales as approximately 0.06 times Reynolds number times diameter, while turbulent lengths are shorter. Within this region, boundary layer grows and core velocity accelerates. Knowledge is important for instrument positioning and heat exchanger design.}{catAD}
\glsadd{Favorable-Pressure-Gradient}\dictentry{Favorable Pressure Gradient}{A favorable pressure gradient occurs when static pressure decreases in the flow direction, causing boundary layer fluid to accelerate and gain kinetic energy. This stabilizes the layer, delays transition, and resists separation. Favorable gradients exist on forward airfoil portions and in converging nozzles. Designers seek to maintain them over as much surface as possible to minimize drag.}{catAD}
\glsadd{Hagen-Poiseuille-Flow}\dictentry{Hagen-Poiseuille Flow}{Hagen-Poiseuille flow describes steady, fully developed laminar flow through a constant-diameter circular pipe. The velocity profile is parabolic with maximum at centerline and zero at walls, and flow rate is proportional to pressure gradient and the fourth power of radius. This canonical flow provides exact Navier-Stokes solutions and is widely used to determine fluid viscosity from pipe experiments.}{catAD}
\glsadd{Hydraulic-Diameter}\dictentry{Hydraulic Diameter}{Hydraulic diameter is defined as four times cross-sectional area divided by wetted perimeter, extending circular pipe correlations to non-circular geometries. For circular pipes it equals the actual diameter; for rectangular or annular channels it provides an equivalent for Reynolds number calculations. This concept allows engineers to use circular pipe results for non-circular ducts in heat exchangers and hydraulic circuits.}{catAD}
\glsadd{Kinematic-Viscosity}\dictentry{Kinematic Viscosity}{Kinematic viscosity is the ratio of dynamic viscosity to density, representing momentum diffusion rate under viscous forces alone. It appears in the Reynolds number definition and governs boundary layer growth and vorticity diffusion. Air at standard conditions has a kinematic viscosity of approximately 1.5 times 10 to the negative fifth power square meters per second. It varies strongly with temperature.}{catAD}
\glsadd{Laminar-Turbulent-Transition}\dictentry{Laminar-Turbulent Transition}{Laminar-turbulent transition is the breakdown of organized laminar boundary layer flow into chaotic turbulence, typically triggered by Tollmien-Schlichting waves or bypass mechanisms. The transition location depends on Reynolds number, pressure gradient, roughness, and freestream turbulence. Transition causes sudden increases in skin friction and heat transfer. Techniques like laminar flow control and surface suction aim to delay it.}{catAD}
\glsadd{Linearized-Theory}\dictentry{Linearized Theory}{Linearized theory simplifies governing flow equations by assuming small perturbations from a uniform state, neglecting nonlinear terms. For subsonic and supersonic flows, this yields the linearized potential equation admitting closed-form solutions for thin airfoils and slender bodies. It provides fundamental results like the thin airfoil theorem. While limited to small angles, it offers invaluable physical insight.}{catAD}
\glsadd{Local-Mach-Number}\dictentry{Local Mach Number}{The local Mach number is the ratio of local flow velocity to local speed of sound at a given point. Unlike freestream Mach number, it varies spatially due to acceleration around curved surfaces and across shocks. On a transonic airfoil, it can exceed unity even at subsonic freestream, creating local supersonic pockets. It determines local pressure, temperature, and density.}{catAD}
\glsadd{No-Slip-Condition}\dictentry{No-Slip Condition}{The no-slip condition requires fluid velocity at a solid surface to equal the surface velocity. At a stationary wall, fluid velocity is zero, creating the velocity gradient driving viscous shear and boundary layer formation. This arises from molecular interactions and is universally applicable for Newtonian fluids at continuum scales. It is the physical origin of all skin friction drag.}{catAD}
\glsadd{Potential-Flow}\dictentry{Potential Flow}{Potential flow is an idealized model assuming irrotational, inviscid fluid motion, allowing velocity to be expressed as the gradient of a scalar potential. This reduces the governing equation to Laplace's equation, admitting solutions through superposition and conformal mapping. It provides accurate pressure predictions on streamlined bodies and forms the basis of panel methods. It fails for viscous and rotational flows.}{catAD}
\glsadd{Stokes-Flow}\dictentry{Stokes Flow}{Stokes flow, or creeping flow, describes fluid motion at very low Reynolds numbers where viscous forces completely dominate inertial forces. The nonlinear terms are neglected, leaving a linear system admitting exact solutions for many geometries. It is characterized by smooth, reversible streamlines and no turbulence. This regime is relevant to microorganism motion, sedimentation, and microfluidic devices.}{catAD}
\glsadd{Stream-Function}\dictentry{Stream Function}{The stream function is a scalar defined for two-dimensional incompressible flows where constant-value lines are streamlines with no flow crossing. Velocity components are obtained by differentiating it, automatically satisfying continuity. In irrotational flow, it satisfies Laplace's equation, enabling complex potential theory. It provides an elegant framework for analyzing two-dimensional patterns.}{catAD}
\glsadd{Taylor-Couette-Flow}\dictentry{Taylor-Couette Flow}{Taylor-Couette flow is fluid motion between two concentric rotating cylinders. At low speeds the flow is purely azimuthal, but as the Taylor number exceeds a critical value, centrifugal instability produces counter-rotating toroidal Taylor vortices. Further increases lead to wavy vortices and eventually turbulence. This canonical problem is fundamental to understanding centrifugal instability and rotating machinery.}{catAD}
\glsadd{Thermal-Conductivity}\dictentry{Thermal Conductivity}{Thermal conductivity quantifies heat conduction rate through a substance by molecular diffusion, defined as heat flux to temperature gradient. In aerodynamics, it determines heat transfer rates, particularly important in hypersonic heating. Air at standard conditions has thermal conductivity of approximately 0.026 watts per meter per kelvin. It appears in the energy equation and Prandtl number.}{catAD}
\glsadd{Turbulent-Reynolds-Stress}\dictentry{Turbulent Reynolds Stress}{Turbulent Reynolds stress is the additional apparent stress in time-averaged equations arising from velocity fluctuation correlations. These stresses represent momentum transport by turbulent eddies and are typically much larger than viscous stresses. The Reynolds stress tensor has six independent components requiring modeling to close turbulence equations, which is the central challenge of turbulence modeling.}{catAD}
\glsadd{Velocity-Profile}\dictentry{Velocity Profile}{A velocity profile describes how fluid velocity varies with position, typically distance normal to a wall or across a duct cross section. Laminar profiles are smooth and predictable, while turbulent profiles are fuller with steep near-wall gradients and logarithmic outer regions. Profiles are measured using pitot probes or laser velocimetry. The shape determines displacement thickness, momentum thickness, and drag characteristics.}{catAD}
\glsadd{Vorticity}\dictentry{Vorticity}{Vorticity is a vector quantity describing local rotational motion of fluid elements, defined as the curl of the velocity field. In irrotational flow it is zero everywhere, while viscous flows generate it at surfaces by the no-slip condition. Vorticity governs boundary layer formation, vortex shedding, and the Biot-Savart law relating vortex elements to induced velocities. The vorticity transport equation is fundamental to vortex methods.}{catAD}
\end{multicols}

\clearpage
\chapter*{Aerospace}
\addcontentsline{toc}{chapter}{Aerospace}
\markboth{Aerospace}{Aerospace}
\clearpage
\begin{multicols}{2}
\small
Aircraft and spacecraft systems: avionics, propulsion (non-rocket), flight controls, airframe structures, and the engineered systems that enable flight.\par
\vspace{0.5cm}
\glsadd{Aerospace-Engineering}\dictentry{Aerospace Engineering}{Aerospace Engineering is the primary field covering the design, development, and testing of aircraft and spacecraft. It encompasses two closely related disciplines: aeronautics for atmospheric flight and astronautics for space travel. Engineers apply principles of aerodynamics, structural mechanics, propulsion, and control systems. The discipline drives innovation in commercial aviation, defense, satellite technology, and space exploration.}{catAE}
\glsadd{Aircraft}\dictentry{Aircraft}{An Aircraft is any vehicle designed to travel through the atmosphere, generating lift through aerodynamic or lighter-than-air means. Aircraft range from small unmanned drones to massive commercial airliners and military fighters. They are classified by method of lift generation, including fixed-wing, rotary-wing, and lighter-than-air designs. Modern aircraft rely on advanced materials, propulsion, and avionics for safe, efficient flight.}{catAE}
\glsadd{Airplane}\dictentry{Airplane}{An Airplane is a powered fixed-wing aircraft that generates lift through forward motion of its wings. It consists of a fuselage, wings, empennage, and one or more engines mounted on the airframe. Airplanes are the most common form of air transportation, carrying passengers and cargo worldwide. They range from single-engine light aircraft to wide-body jets capable of intercontinental travel.}{catAE}
\glsadd{Helicopter}\dictentry{Helicopter}{A Helicopter is a rotary-wing aircraft with one or more horizontal rotors providing both lift and thrust. Unlike fixed-wing aircraft, helicopters can hover, take off vertically, and maneuver in any direction. The main rotor is powered by a turbine or piston engine, while a tail rotor counters torque. Helicopters are widely used in emergency medical services, law enforcement, and military operations.}{catAE}
\glsadd{Rotorcraft}\dictentry{Rotorcraft}{Rotorcraft is a class of heavier-than-air craft using rotating blades to generate lift and thrust. This category includes helicopters, gyrocopters, and tiltrotor vehicles, all relying on rotor aerodynamics. Rotorcraft offer unique capabilities such as vertical takeoff, hovering, and low-speed maneuvering. They serve critical roles where conventional runways are impractical.}{catAE}
\glsadd{Tiltrotor}\dictentry{Tiltrotor}{A Tiltrotor combines vertical takeoff and landing capability with efficient forward flight. Rotating nacelles transition from vertical helicopter mode to horizontal airplane mode during flight. In helicopter mode rotors provide lift, while in airplane mode they function as propellers. The V-22 Osprey is the most prominent example, serving the United States military.}{catAE}
\glsadd{Gyrocopter}\dictentry{Gyrocopter}{A Gyrocopter is a rotorcraft with an unpowered rotor spinning freely in autorotation. Forward flight comes from a separate engine-driven propeller, while the spinning rotor creates lift. Gyrocopters are mechanically simpler and more affordable than helicopters. They can operate from short fields and are inherently safe since autorotation maintains lift if the engine fails.}{catAE}
\glsadd{Glider}\dictentry{Glider}{A Glider is a fixed-wing aircraft without an engine that relies on rising air currents for sustained flight. Pilots launch from tow planes or winches, then seek thermals and ridge lift to extend duration. Modern gliders feature high aspect ratio wings, achieving glide ratios exceeding 60:1. They are used for pilot training, sport aviation, and atmospheric research.}{catAE}
\glsadd{UAV}\dictentry{UAV}{UAV stands for Unmanned Aerial Vehicle and refers to the physical aircraft component of an unmanned system. It is the airborne platform that carries propulsion, sensors, navigation, and communication hardware. A UAV operates without a pilot on board, relying on onboard computers or a remote operator for guidance. The term is widely used in military, commercial, and recreational contexts to describe any pilotless flying machine.}{catAE}
\glsadd{UAS}\dictentry{UAS}{UAS stands for Unmanned Aircraft System and encompasses the complete operational package beyond just the aircraft. It includes the UAV itself, the ground control station, the data link communications, and the support equipment needed for operation. The UAS framework recognizes that an unmanned aircraft cannot function without its supporting infrastructure of controllers, telemetry, and software. This holistic term is preferred in regulatory and professional contexts.}{catAE}
\glsadd{eVTOL}\dictentry{eVTOL}{An Electric Vertical Takeoff and Landing aircraft is a battery-powered vehicle capable of vertical liftoff. eVTOL designs typically use multiple electric motors driving rotors for lift and propulsion. These vehicles aim to enable urban air mobility through quiet, emission-free transportation. Several manufacturers are developing eVTOL platforms with certification expected in coming years.}{catAE}
\glsadd{Airship}\dictentry{Airship}{An Airship is a powered, steerable lighter-than-air craft using buoyant gas for lift and engines for propulsion. The gas, typically helium, provides sufficient lift to carry the envelope, structure, and payload. Airships consume far less fuel than heavier-than-air aircraft and can stay airborne for extended durations. They are used for tourism, surveillance, and research where endurance and low fuel consumption matter.}{catAE}
\glsadd{Balloon}\dictentry{Balloon}{A Balloon is an unpowered lighter-than-air craft that rises and descends by controlling buoyancy. Hot air balloons use heated air for lift, while gas balloons fill the envelope with helium. Balloons have been used for scientific observation and recreational flight since the eighteenth century. Weather balloons routinely ascend to the stratosphere carrying atmospheric instruments.}{catAE}
\glsadd{Spacecraft}\dictentry{Spacecraft}{A spacecraft is a vehicle designed to operate beyond Earth's atmosphere and travel through the vacuum of space. Spacecraft range from uncrewed probes and satellites to crewed capsules and space stations. They carry their own life support, power, thermal control, and communication systems to survive the harsh space environment. Spacecraft are launched by rockets and use propulsion for orbit insertion, maneuvering, and attitude control.}{catAE}
\glsadd{Satellite}\dictentry{Satellite}{A satellite is any object placed into orbit around a celestial body, whether natural or artificial. Artificial satellites are manufactured objects launched for communication, navigation, Earth observation, and research. They range from CubeSats of a few kilograms to large geostationary platforms weighing several tons.}{catAE}
\glsadd{Launch-Vehicle}\dictentry{Launch Vehicle}{A Launch Vehicle is a rocket designed to transport payloads from Earth's surface to space, including orbit, suborbital trajectory, or escape velocity. Launch vehicles range from small sounding rockets to massive heavy-lift systems delivering tens of tons to low Earth orbit. They are typically multi-stage to optimize mass efficiency by discarding empty structure during ascent.}{catAE}
\glsadd{Sounding-Rocket}\dictentry{Sounding Rocket}{A Sounding Rocket is a suborbital rocket carrying scientific instruments into the upper atmosphere or near-space for brief measurement periods. Following a ballistic trajectory, they typically reach altitudes between fifty and a thousand kilometers. Sounding rockets provide a cost-effective means of conducting microgravity, atmospheric chemistry, and space physics experiments. They have been instrumental in advancing understanding of the upper atmosphere.}{catAE}
\glsadd{Spaceplane}\dictentry{Spaceplane}{A Spaceplane is a vehicle capable of operating in both atmosphere and space within a single mission. It combines aerodynamic surfaces for atmospheric flight with reaction control for space maneuvering. Spaceplanes use wings for gliding descent and heat-resistant structures for re-entry. Historic examples include the Space Shuttle, with modern designs like Dream Chaser continuing development.}{catAE}
\glsadd{Hypersonic-Vehicle}\dictentry{Hypersonic Vehicle}{A Hypersonic Vehicle is an aircraft or missile capable of sustained flight above Mach 5. At these speeds aerodynamic heating becomes extreme, requiring specialized thermal protection materials. Hypersonic flight enables rapid global transportation and advanced military strike capabilities. Research into scramjet engines and thermal management continues to advance practical hypersonic systems.}{catAE}
\glsadd{Flight-Mechanics}\dictentry{Flight Mechanics}{Flight Mechanics is the study of forces acting on an aircraft and the resulting motion through the atmosphere. It analyzes lift, drag, thrust, and gravity to determine performance, stability, and control characteristics. Engineers use flight mechanics to predict takeoff distances, cruise efficiency, and maneuvering capability. The discipline forms the foundation for aircraft design, simulation, and pilot training.}{catAE}
\glsadd{Sideslip-Angle}\dictentry{Sideslip Angle}{The Sideslip Angle is the angle between the aircraft's longitudinal axis and the oncoming relative wind. A positive sideslip occurs when the nose points to one side while the aircraft moves along a different path. Pilots and autopilot systems monitor sideslip to coordinate turns and prevent dangerous conditions. The sideslip indicator helps pilots maintain coordinated flight during all maneuvers.}{catAE}
\glsadd{Flight-Path}\dictentry{Flight Path}{The Flight Path is the trajectory traced by an aircraft's center of gravity through the atmosphere. It is defined by ground track, altitude profile, and vertical speed along the planned route. Flight path analysis considers wind, atmospheric conditions, and aircraft performance for optimal efficiency. Modern avionics display the flight path in real time to monitor adherence to the planned route.}{catAE}
\glsadd{Glide-Ratio}\dictentry{Glide Ratio}{The Glide Ratio is the horizontal distance an aircraft travels per unit of altitude lost in unpowered flight. It equals the lift-to-drag ratio and measures aerodynamic efficiency. High-performance gliders achieve ratios exceeding 60:1, while airliners typically glide around 17:1. Understanding glide ratio is essential for emergency descent planning and evaluating unpowered range.}{catAE}
\glsadd{Ceiling}\dictentry{Ceiling}{The Ceiling is the maximum altitude at which an aircraft can sustain level flight. Absolute ceiling is the highest achievable altitude, while service ceiling is where climb rate drops to 100 feet per minute. Ceiling is limited by decreasing air density, which reduces engine power and aerodynamic lift. High-altitude aircraft like the U-2 operate near their ceiling to maximize performance.}{catAE}
\glsadd{Stall-Recovery}\dictentry{Stall Recovery}{Stall Recovery encompasses the techniques pilots use to regain controlled flight after an aerodynamic stall. A stall occurs when the wing exceeds its critical angle of attack and loses sufficient lift. Recovery requires reducing the angle of attack by pushing the nose down, then smoothly reapplying power. Pilot training emphasizes stall recognition and recovery as fundamental skills for safe flight.}{catAE}
\glsadd{Spin-Recovery}\dictentry{Spin Recovery}{Spin Recovery involves procedures for exiting an autorotative spin where the aircraft descends in a helical path. A spin results from a stalled wing combined with yawing motion, causing asymmetric lift. Standard recovery involves opposite rudder, forward control input, then neutralizing controls as rotation stops. Spin recovery techniques vary by aircraft type and are required for certification.}{catAE}
\glsadd{Fuselage}\dictentry{Fuselage}{The Fuselage is the central body of an aircraft housing the cockpit, crew, passengers, and cargo. It is designed to minimize aerodynamic drag while providing structural integrity for attachments. Fuselage construction typically uses semi-monocoque design with frames, stringers, and stressed skin. The shape and size directly impact capacity, drag characteristics, and overall performance.}{catAE}
\glsadd{Wing-Spar}\dictentry{Wing Spar}{The Wing Spar is the principal spanwise structural member carrying bending loads from lift and weight. Spars run from the wing root to the tip, typically constructed from aluminum or composite materials. Most wings have one or two main spars along with secondary spars to distribute loads. The spar is the primary load path, and its design determines wing strength and fatigue life.}{catAE}
\glsadd{Rib}\dictentry{Rib}{A Rib is a chordwise structural member maintaining the airfoil shape of a wing between spars and skin. Ribs transfer aerodynamic loads from the skin to the spars and provide attachment points. They are typically stamped from aluminum or molded from composite materials to match the airfoil contour. Rib spacing affects the wing's resistance to torsion and ability to maintain shape under load.}{catAE}
\glsadd{Stringer}\dictentry{Stringer}{A Stringer is a thin longitudinal reinforcing element attached to the inner surface of aircraft skin. Stringers increase the skin's resistance to buckling under compressive and shear loads. They work with frames and bulkheads to create the semi-monocoque structure of modern aircraft. Stringer spacing and cross-section optimize the strength-to-weight ratio of the airframe.}{catAE}
\glsadd{Bulkhead}\dictentry{Bulkhead}{A Bulkhead is a transverse structural partition carrying concentrated loads and maintaining fuselage shape. Bulkheads are placed at major attachment points such as landing gear mounts and pressure boundaries. They transfer loads from concentrated sources into the surrounding skin and stringer structure. Pressurized bulkheads maintain the pressure differential essential for high-altitude flight.}{catAE}
\glsadd{Longeron}\dictentry{Longeron}{A Longeron is a major longitudinal structural member running along the length of the fuselage. Longerons carry primary bending loads from flight forces, landing impacts, and pressurization. They are heavy-duty extrusions attached to frames at regular intervals. Together with stringers and skin, longerons form the structural backbone of the fuselage.}{catAE}
\glsadd{Empennage}\dictentry{Empennage}{The Empennage is the tail assembly consisting of horizontal and vertical surfaces providing stability and control. It includes the vertical stabilizer with rudder for directional control and horizontal stabilizer with elevator for pitch. The empennage counteracts the natural tendency to yaw, pitch, or deviate from stable flight. Its design balances aerodynamic effectiveness with structural weight and drag.}{catAE}
\glsadd{Vertical-Stabilizer}\dictentry{Vertical Stabilizer}{The Vertical Stabilizer is the fixed vertical surface providing directional stability to the aircraft. It prevents uncommanded yaw by generating a restoring side force when heading deviates. The rudder, a movable trailing edge surface, allows the pilot to control yaw actively. Sizing requires analysis of engine-out conditions, crosswind landings, and spinning characteristics.}{catAE}
\glsadd{Horizontal-Stabilizer}\dictentry{Horizontal Stabilizer}{The Horizontal Stabilizer is the fixed horizontal surface providing longitudinal stability. It generates a downward force to counter the nose-down pitching moment from wing and fuselage. The elevator on the trailing edge allows the pilot to control pitch attitude. Some designs use a fully movable stabilator for more effective pitch control at all speeds.}{catAE}
\glsadd{Fin}\dictentry{Fin}{The Fin is another term for the vertical stabilizer, the fixed vertical tail surface of an aircraft. It provides directional stability and resists uncommanded yawing motion. The term is commonly used in aviation and appears in design literature as a synonym for vertical tail. The fin works with the rudder to maintain heading control during all flight phases.}{catAE}
\glsadd{Skin}\dictentry{Skin}{The Skin is the outer covering transferring aerodynamic loads to the internal aircraft framework. In semi-monocoque construction, skin carries shear and torsional loads while forming the aerodynamic surface. Aircraft skin is typically aluminum alloy, titanium, or composite laminates. Surface smoothness directly affects drag, fuel efficiency, and aerodynamic performance.}{catAE}
\glsadd{Honeycomb-Panel}\dictentry{Honeycomb Panel}{A honeycomb panel is a sandwich structure consisting of a lightweight honeycomb core bonded between two thin face sheets. This construction provides exceptional strength and stiffness at very low weight. Honeycomb panels are used extensively in spacecraft structures, including equipment panels, antenna reflectors, and solar array substrates. The honeycomb geometry efficiently distributes loads while minimizing material mass.}{catAE}
\glsadd{Monocoque}\dictentry{Monocoque}{Monocoque is a structural design where the outer skin carries all or most loads without internal bracing. The name comes from French meaning single shell, reflecting the load-bearing exterior concept. Early aircraft and modern race car chassis use monocoque construction for lightweight smooth exteriors. Pure monocoque structures are limited because skin alone cannot resist large concentrated loads.}{catAE}
\glsadd{Semi-Monocoque}\dictentry{Semi-Monocoque}{Semi-Monocoque is a stressed-skin design where loads are shared between the outer skin and internal framework. Internal members include stringers, longerons, frames, and bulkheads reinforcing the skin. This is the standard construction for modern aircraft fuselages and wings due to its strength-to-weight ratio. Semi-monocoque structures provide damage tolerance through redundant load paths.}{catAE}
\glsadd{Aluminum-Alloy}\dictentry{Aluminum Alloy}{Aluminum Alloy is a lightweight metallic material used extensively in aircraft structures. Common aerospace alloys such as 2024 and 7075 offer high strength, fatigue resistance, and ease of manufacturing. Aluminum alloys can be formed and machined into complex structural shapes for airframe components. Although composites are growing, aluminum remains dominant due to proven reliability and cost.}{catAE}
\glsadd{Titanium-Alloy}\dictentry{Titanium Alloy}{Titanium Alloy is a high-strength, corrosion-resistant metal for critical high-temperature components. Aerospace titanium maintains properties at temperatures where aluminum alloys lose strength. It is used in engine components, landing gear, fasteners, and structural elements. Despite higher cost, titanium's unique properties make it indispensable in demanding applications.}{catAE}
\glsadd{Carbon-Fiber}\dictentry{Carbon Fiber}{Carbon Fiber is a strong, stiff composite material made from thin filaments of carbon atoms. When embedded in polymer matrix, carbon fiber reinforced plastic achieves strength-to-weight ratios exceeding aluminum. It is used extensively in modern aircraft wings, fuselages, and control surfaces. Composites reduce weight, improve fuel efficiency, and enable complex aerodynamic shapes.}{catAE}
\glsadd{Fiberglass}\dictentry{Fiberglass}{Fiberglass consists of glass fibers embedded in a polymer resin matrix, forming a lightweight composite. It offers good strength, electrical insulation, and resistance to corrosion and moisture. In aerospace, fiberglass is used for radomes, interior panels, fairings, and secondary structures. While not as strong as carbon fiber, fiberglass is less expensive for non-critical applications.}{catAE}
\glsadd{Kevlar}\dictentry{Kevlar}{Kevlar is a strong, heat-resistant synthetic fiber belonging to the aramid polymer family. It provides exceptional tensile strength at very low weight for impact-resistant applications. In aerospace, Kevlar is used for fuel tanks, pressure vessels, and crash-resistant structures. Kevlar composites absorb energy effectively and resist penetration in critical components.}{catAE}
\glsadd{Composite-Material}\dictentry{Composite Material}{A Composite Material combines two or more constituents with different properties for superior performance. In aerospace, composites typically consist of high-strength fibers in a polymer, metal, or ceramic matrix. They offer high strength-to-weight ratios, corrosion resistance, and tailorable directional properties. Modern aircraft increasingly use composites to reduce weight and lower maintenance requirements.}{catAE}
\glsadd{Ceramic-Matrix-Composite}\dictentry{Ceramic Matrix Composite}{A Ceramic Matrix Composite has a ceramic matrix reinforced by fibers to improve toughness and thermal resistance. Unlike monolithic ceramics, CMCs withstand thermal shock without catastrophic brittle failure. They are used in high-temperature jet engine sections and thermal protection systems. CMCs enable higher engine operating temperatures, directly improving thermodynamic efficiency.}{catAE}
\glsadd{Superalloy}\dictentry{Superalloy}{A Superalloy maintains high strength, creep resistance, and oxidation resistance at elevated temperatures. Nickel-based, cobalt-based, and iron-based superalloys serve in the hottest jet engine sections. These alloys withstand temperatures exceeding 1,000 degrees Celsius while resisting fatigue and corrosion. Superalloys are critical for turbine blades, combustion chambers, and exhaust components.}{catAE}
\glsadd{Shape-Memory-Alloy}\dictentry{Shape Memory Alloy}{A Shape Memory Alloy returns to a predetermined shape when heated after being deformed. This property results from a reversible phase transformation between crystal structures. In aerospace, shape memory alloys are used for actuators, morphing structures, and self-deploying antennas. These materials enable compact mechanisms without traditional motors or hydraulic actuators.}{catAE}
\glsadd{Ablative-Material}\dictentry{Ablative Material}{An Ablative Material erodes in a controlled manner to absorb extreme heat during re-entry. As the material burns away, it carries thermal energy from the underlying structure. Ablative heat shields are used on re-entry capsules, rocket nozzles, and solid motor cases. This thermal protection approach is simple, reliable, and proven since early spaceflights.}{catAE}
\glsadd{Hydraulic-System}\dictentry{Hydraulic System}{A Hydraulic System uses pressurized fluid to actuate flight controls, landing gear, and mechanical systems. Hydraulic fluid is pumped through lines to actuators converting pressure into linear or rotary motion. Hydraulic systems provide the high force and precise control needed for large control surfaces. Redundant circuits ensure critical controls remain operational during partial system failures.}{catAE}
\glsadd{Pneumatic-System}\dictentry{Pneumatic System}{A Pneumatic System uses compressed air to power aircraft subsystems including engine starting and environmental controls. Bleed air from jet engines is typically the source, distributed through valves and ducts. Pneumatic systems also power anti-ice boots and windshield wipers on commercial aircraft. The simplicity and reliability of pneumatic systems make them valuable for many functions.}{catAE}
\glsadd{Electrical-System}\dictentry{Electrical System}{The Electrical System provides power to operate avionics, lighting, environmental controls, and subsystems. Power is generated by engine-driven generators and supplemented by batteries and auxiliary power units. Modern aircraft use variable-frequency AC systems along with DC buses for different equipment. Electrical redundancy through multiple generators ensures critical systems remain powered.}{catAE}
\glsadd{Fuel-System}\dictentry{Fuel System}{The Fuel System stores, manages, and delivers fuel from tanks to engines in correct quantity and pressure. It includes tanks, pumps, valves, filters, and distribution lines for safe operation in all conditions. Fuel system design accounts for center of gravity management and fuel jettison capability. The system provides quantity indication, temperature monitoring, and contamination detection.}{catAE}
\glsadd{Environmental-Control-System}\dictentry{Environmental Control System}{The Environmental Control System manages cabin pressure, temperature, humidity, and air quality. It uses engine bleed air or electric compressors to provide conditioned air at appropriate conditions. The ECS includes air filtration, ozone conversion, and cooling for a habitable environment. Reliable environmental control is essential for crew performance and passenger safety.}{catAE}
\glsadd{Cabin-Pressurization}\dictentry{Cabin Pressurization}{Cabin Pressurization maintains safe pressure altitude inside the aircraft during high-altitude flight. Without pressurization, passengers would experience hypoxia at cruising altitudes above 10,000 feet. Controlled outflow valves regulate cabin altitude, typically equivalent to 6,000 to 8,000 feet. The fuselage must withstand the pressure differential between cabin and thin outside atmosphere.}{catAE}
\glsadd{Landing-Gear}\dictentry{Landing Gear}{Landing Gear is the undercarriage supporting the aircraft on the ground, absorbing landing loads. It consists of wheels, tires, brakes, shock absorbers, and retraction mechanisms. Landing gear must withstand high impact forces while remaining lightweight for flight. Configurations include tricycle, taildragger, and multi-wheel designs based on aircraft requirements.}{catAE}
\glsadd{Braking-System}\dictentry{Braking System}{The Braking System decelerates the aircraft during landing rollout and rejected takeoffs. Carbon or steel disc brakes are actuated by hydraulic pressure with anti-skid systems. Brake energy absorption is critical as landing at high speed generates enormous thermal energy. Brake temperature monitoring alerts pilots to overheating that could cause tire failure.}{catAE}
\glsadd{Anti-Skid-System}\dictentry{Anti-Skid System}{The Anti-Skid System prevents wheel lockup during braking by modulating brake pressure based on slip. It functions like automotive anti-lock brakes but is designed for aircraft loads and speeds. Sensors measure wheel rotational speed, and the system reduces pressure when skid is detected. Anti-skid maintains directional control and braking effectiveness on contaminated runways.}{catAE}
\glsadd{Ice-Protection-System}\dictentry{Ice Protection System}{The Ice Protection System prevents or removes ice on wings, tail, engines, and sensors. Methods include thermal anti-ice using hot bleed air, electrothermal de-ice boots, and chemical systems. Ice disrupts airflow, reduces lift, increases drag, and can damage engine components. Reliable ice protection is essential for safe flight in hazardous icing conditions.}{catAE}
\glsadd{Flight-Computer}\dictentry{Flight Computer}{The Flight Computer is the central processing unit integrating data from aircraft sensors and avionics. It calculates navigation solutions, performance parameters, and provides inputs to autopilot systems. Modern flight computers use redundant processors and voting logic for data integrity. The flight computer coordinates information flow between all electronic systems aboard the aircraft.}{catAE}
\glsadd{Flight-Management-System}\dictentry{Flight Management System}{The Flight Management System automates in-flight tasks related to navigation and performance. It combines GPS, inertial navigation, and air data for precise positioning and route guidance. The FMS computes optimal climb profiles, cruise altitudes, and fuel predictions. Pilots program the FMS to reduce workload and enable efficient operations across the air network.}{catAE}
\glsadd{Autopilot}\dictentry{Autopilot}{An Autopilot automatically controls the aircraft's flight path without direct pilot input. It uses sensors to measure attitude, heading, altitude, and airspeed, then commands control surfaces. Autopilot modes include altitude hold, heading select, approach coupling, and full FMS guidance. Modern autopilot reduces workload, improves precision, and enhances safety during all flight phases.}{catAE}
\glsadd{Fly-by-Wire}\dictentry{Fly-by-Wire}{Fly-by-Wire electronically transmits pilot inputs from the cockpit to control surface actuators. Electrical signals travel through redundant wiring and flight control computers to command surfaces. FBW enables relaxed static stability designs that improve maneuverability and fuel efficiency. The computers also provide envelope protection, preventing exceeding safe structural limits.}{catAE}
\glsadd{Glass-Cockpit}\dictentry{Glass Cockpit}{A Glass Cockpit features electronic multi-function displays instead of traditional analog instruments. Glass cockpits present flight data, navigation, engine parameters, and system status on digital screens. They improve situational awareness by integrating and prioritizing information visually. The transition to glass cockpits has been a major advancement in aviation safety.}{catAE}
\glsadd{Head-Up-Display}\dictentry{Head-Up Display}{A Head-Up Display projects critical flight information onto a transparent screen in the pilot's line of sight. The HUD allows monitoring airspeed, altitude, and heading while maintaining visual contact outside. This reduces the need to scan between instruments and the external environment during critical phases. HUDs are especially valuable during low-visibility landings with synthetic vision cues.}{catAE}
\glsadd{Multi-Function-Display}\dictentry{Multi-Function Display}{A Multi-Function Display is a configurable screen presenting various types of flight information. Pilots can switch between navigation maps, weather radar, engine data, and system status pages. MFDs provide flexibility by allowing customization based on the current phase of flight. Multiple MFDs work together to present a comprehensive picture of aircraft conditions.}{catAE}
\glsadd{Air-Data-Computer}\dictentry{Air Data Computer}{The Air Data Computer processes pitot and static port inputs to compute essential flight parameters. It calculates barometric altitude, airspeed, Mach number, vertical speed, and temperature. ADC corrections account for instrument, compressibility, and position errors for accuracy. Modern air data computers integrate with the flight management system and autopilot.}{catAE}
\glsadd{GPS}\dictentry{GPS}{GPS, the Global Positioning System, provides precise three-dimensional position and time via satellites. The constellation transmits timing signals that receivers process to determine location. In aviation, GPS enables precision approaches, direct routing, and enhanced situational awareness. Augmentation systems like WAAS improve accuracy for instrument flight operations.}{catAE}
\glsadd{GNSS}\dictentry{GNSS}{GNSS refers to the collective of satellite navigation constellations including GPS, GLONASS, Galileo, and BeiDou. Using multiple constellations improves availability, accuracy, and integrity of positioning data. GNSS receivers track satellites across constellations to compute a robust navigation solution. Multiple GNSS integration represents the future of global air navigation with greater redundancy.}{catAE}
\glsadd{INS}\dictentry{INS}{An Inertial Navigation System uses accelerometers and gyroscopes to track position and velocity. By measuring acceleration and rotation, the INS computes position changes from a known starting point. INS provides continuous data independent of external signals, immune to jamming or signal loss. Modern strapdown INS uses laser or fiber optic gyroscopes for aviation and spaceflight accuracy.}{catAE}
\glsadd{Dead-Reckoning}\dictentry{Dead Reckoning}{Dead Reckoning estimates current position from a known starting point, speed, time, and direction. It is one of the oldest navigation techniques used by mariners and aviators. Dead reckoning accumulates errors over time, requiring periodic position fixes. In modern aviation it serves as backup when electronic systems are unavailable.}{catAE}
\glsadd{Waypoint}\dictentry{Waypoint}{A Waypoint is a predetermined geographic coordinate defining a point along the planned path of an aircraft. Each is specified by latitude and longitude with an alphanumeric identifier for referencing. Pilots and flight management systems use these points to build flight plans and manage airway routing. Spacing between consecutive waypoints affects fuel efficiency, traffic separation, and ATC compliance.}{catAE}
\glsadd{Great-Circle-Route}\dictentry{Great Circle Route}{A Great Circle Route is the shortest path between two points on a sphere. On flat maps these routes appear curved but represent minimum surface distance. Airlines use great circle routing to minimize fuel consumption and flight time. Actual paths may deviate due to winds, traffic restrictions, and oceanic procedures.}{catAE}
\glsadd{Magnetic-Heading}\dictentry{Magnetic Heading}{Magnetic Heading is the direction an aircraft points relative to magnetic north. The Earth's magnetic field does not align with true geographic north, creating variation. Pilots convert between magnetic and true headings using local magnetic variation. Magnetic heading remains essential for compass-based navigation and instrument flying.}{catAE}
\glsadd{True-Heading}\dictentry{True Heading}{True Heading is the direction an aircraft points relative to true geographic north. It is calculated by applying magnetic variation to the magnetic heading obtained from a compass. Flight plans and charts use this as the standard reference for routing and navigation. Modern avionics systems automatically compute it from GPS and inertial data.}{catAE}
\glsadd{Indicated-Airspeed}\dictentry{Indicated Airspeed}{Indicated Airspeed is the speed read directly from the airspeed indicator via pitot-static pressure. IAS does not account for instrument, position, or compressibility errors. Despite limitations, IAS is the primary reference as aerodynamic limits are based on it. Airspeed indicators are calibrated to display IAS for all flight phases.}{catAE}
\glsadd{Calibrated-Airspeed}\dictentry{Calibrated Airspeed}{Calibrated Airspeed corrects Indicated Airspeed for instrument and position errors. CAS provides more accurate dynamic pressure representation for precise calculations. The difference between IAS and CAS increases at higher speeds and altitudes. Pilots use CAS for performance-critical phases like takeoff and landing.}{catAE}
\glsadd{True-Airspeed}\dictentry{True Airspeed}{True Airspeed is the actual speed of the aircraft relative to the undisturbed air mass. TAS is higher than indicated airspeed because air density decreases with altitude. TAS is used for navigation, fuel planning, and position reporting to ATC. Flight management systems continuously compute and display TAS from air density measurements.}{catAE}
\glsadd{Ground-Speed}\dictentry{Ground Speed}{Ground Speed is the actual speed of an aircraft relative to the ground surface below. It equals true airspeed adjusted for wind component along the flight path. GPS and inertial navigation systems provide accurate measurement independent of air data sensors. This value determines travel time between waypoints and is the primary reference for navigation and scheduling.}{catAE}
\glsadd{Turboprop}\dictentry{Turboprop}{A Turboprop is a jet engine where a gas turbine drives a propeller through a reduction gearbox. Most thrust comes from the propeller, with small contribution from jet exhaust. Turboprops excel in regional aviation, military transport, and utility operations. They offer better fuel efficiency than turbofans at lower speeds and shorter routes.}{catAE}
\glsadd{Turboshaft}\dictentry{Turboshaft}{A Turboshaft is a gas turbine optimized to deliver shaft power rather than jet thrust. Nearly all turbine energy drives an output shaft for rotors, propellers, or equipment. Turboshafts provide high power-to-weight ratios and are the standard helicopter powerplant. Their compact size also suits auxiliary power units and marine applications.}{catAE}
\glsadd{Pulsejet}\dictentry{Pulsejet}{A Pulsejet operates through intermittent combustion cycles rather than continuous burn. Valves or acoustic resonance control air intake and exhaust, creating rapid thrust pulses. Pulsejets are mechanically simple but produce high noise and lower efficiency than turbojets. The German V-1 flying bomb of World War II was the most famous pulsejet weapon.}{catAE}
\glsadd{Combustion-Chamber}\dictentry{Combustion Chamber}{The Combustion Chamber is the enclosed volume within a rocket engine where propellants mix and burn at high temperature and pressure. Its size and shape determine the residence time of propellants and the completeness of combustion. Chamber walls experience extreme thermal loads and require active cooling or thermal protection. The combustion chamber connects directly to the nozzle, forming the core of the thrust chamber assembly.}{catAE}
\glsadd{Afterburner}\dictentry{Afterburner}{An Afterburner is a secondary combustion section temporarily increasing jet engine thrust. Additional fuel is injected into hot exhaust and reignited for a significant thrust boost. Afterburners consume fuel at extremely high rates, limiting use to short durations. Military fighters use afterburners for supersonic flight and high-performance combat maneuvers.}{catAE}
\glsadd{Specific-Impulse}\dictentry{Specific Impulse}{Specific Impulse (Isp) is a fundamental measure of rocket propellant efficiency, expressed in seconds. It represents thrust produced per unit weight flow rate of propellant, with higher values indicating greater momentum extracted per kilogram of fuel. Isp varies between vacuum and sea-level conditions due to ambient pressure effects on nozzle expansion. Typical values range from about 200 seconds for solid motors to over 450 seconds for liquid hydrogen engines, directly determining delta-v capability through the Tsiolkovsky equation.}{catAE}
\glsadd{Thrust-Vector-Control}\dictentry{Thrust Vector Control}{Thrust Vector Control is the system that steers a rocket by redirecting the direction of the engine thrust vector. Methods include gimbaling the entire engine, using secondary injection into the exhaust, or employing deflection vanes. TVC is essential during atmospheric flight for guidance and during powered flight for trajectory corrections.}{catAE}
\glsadd{Payload}\dictentry{Payload}{The Payload is the cargo carried by a launch vehicle or aircraft to accomplish its mission objectives. For rockets, this includes satellites, probes, crew capsules, and scientific instruments. Capacity determines vehicle class and influences design, staging, and trajectory optimization. Maximizing the fraction of total vehicle mass devoted to useful cargo is a primary goal in aerospace vehicle design.}{catAE}
\glsadd{Fairing}\dictentry{Fairing}{The Fairing is an aerodynamic shell enclosing and protecting the payload during launch. It shields the payload from heating, acoustic vibration, and debris during atmospheric ascent. After exiting dense atmosphere, the fairing is jettisoned to reduce weight. Design balances aerodynamic performance with structural strength to protect sensitive equipment.}{catAE}
\glsadd{Booster}\dictentry{Booster}{A Booster is the first stage providing initial thrust to lift the rocket off the pad. Boosters generate the greatest thrust because they must overcome the full vehicle weight. Many vehicles use strap-on solid or liquid boosters for additional launch capability. Booster recovery and reuse have dramatically reduced the cost of access to space.}{catAE}
\glsadd{Upper-Stage}\dictentry{Upper Stage}{The Upper Stage is the final propelled stage of a launch vehicle, responsible for precise orbit insertion and any required plane changes. Upper stages operate entirely in vacuum and use high-specific-impulse propellant combinations such as liquid hydrogen and liquid oxygen. Many upper stages are designed to restart multiple times for complex mission profiles.}{catAE}
\glsadd{Delta-v}\dictentry{Delta-v}{Delta-v represents the total velocity change capacity available for spacecraft maneuvers. It is determined by the Tsiolkovsky rocket equation relating specific impulse and mass ratio. Delta-v budgets account for launch, orbit insertion, transfers, and attitude corrections. Sufficient delta-v enables reaching target orbit, course corrections, and safe return.}{catAE}
\glsadd{Escape-Velocity}\dictentry{Escape Velocity}{Escape Velocity is the minimum speed needed to break free from a body's gravitational pull without further propulsion. On Earth, this is approximately 11.2 kilometers per second from the surface. The value depends on the body's mass and radius. Escape velocity is critical for rocket mission design, determining energy requirements for leaving Earth orbit. It also determines whether a body can retain an atmosphere long-term.}{catAE}
\glsadd{Orbital-Velocity}\dictentry{Orbital Velocity}{Orbital Velocity is the speed needed to maintain a stable orbit around a central body at a given altitude. A satellite in low Earth orbit at 200 kilometers must travel about 7.8 kilometers per second. Required velocity decreases at higher altitudes as gravity weakens. Orbital velocity balances gravitational acceleration with centripetal acceleration for curved motion. It is a fundamental parameter in spacecraft operations and mission planning.}{catAE}
\glsadd{Staging}\dictentry{Staging}{Staging is discarding empty rocket stages during ascent to reduce mass for subsequent stages. Each stage carries its own engines and propellant, igniting after the previous stage is jettisoned. Staging dramatically improves efficiency by allowing smaller engines to accelerate less mass. Stage number and size are optimized based on target orbit, payload, and performance requirements.}{catAE}
\glsadd{Gravity-Turn}\dictentry{Gravity Turn}{A Gravity Turn uses gravity to gradually pitch the launch vehicle from vertical to horizontal during ascent. After initial vertical rise, the vehicle performs a slight pitch-over to begin the turn. Gravity naturally bends the trajectory toward horizontal, reducing steering forces and loads. The gravity turn is the standard ascent profile for most launch vehicles.}{catAE}
\glsadd{Hot-Staging}\dictentry{Hot Staging}{Hot staging is a rocket staging technique in which the upper stage engine ignites while still attached to the lower stage. This eliminates the need for a coast phase and ensures propellant settles under thrust before ignition. It improves payload capacity by reducing gravity losses during the staging transition. The method was first used in Soviet rockets and is now employed by modern launch vehicles such as SpaceX Starship.}{catAE}
\glsadd{Spacecraft-Bus}\dictentry{Spacecraft Bus}{The spacecraft bus is the primary structural and functional platform that supports and integrates all subsystems of a spacecraft. It provides the structure, power, thermal control, propulsion, and communication capabilities needed for mission operations. The bus is designed to be compatible with various payloads and can often be adapted for different missions. It serves as the foundation upon which the entire spacecraft is built and operated.}{catAE}
\glsadd{Solar-Panel}\dictentry{Solar Panel}{A solar panel converts sunlight into electricity using the photovoltaic effect in semiconductor materials. Spacecraft solar panels consist of photovoltaic cells on rigid or flexible substrates protected by cover glass against radiation. They provide primary power for most satellites and space stations, operating continuously in sunlit orbit portions. Efficiency, radiation tolerance, and power-to-weight ratio are critical design parameters.}{catAE}
\glsadd{Reaction-Wheel}\dictentry{Reaction Wheel}{A reaction wheel is a spinning flywheel that controls spacecraft orientation without propellant. Changing the wheel's spin speed transfers angular momentum between wheel and spacecraft, causing rotation in the opposite direction. Three or more wheels on different axes provide full three-axis control. They are primary actuators for most satellites and observatories, though accumulated momentum must periodically be desaturated.}{catAE}
\glsadd{Star-Tracker}\dictentry{Star Tracker}{A star tracker determines spacecraft attitude by photographing the star field and comparing patterns to an onboard catalog. Modern star trackers achieve accuracy better than one arcsecond by identifying multiple stars and computing three-axis orientation. They are among the most accurate attitude determination instruments and operate autonomously.}{catAE}
\glsadd{Sun-Sensor}\dictentry{Sun Sensor}{A sun sensor measures the angle of incoming sunlight relative to a spacecraft reference axis to determine attitude. Coarse sensors use photodiodes while fine sensors employ optical systems for higher accuracy. They provide critical references for initialization after launch, solar panel orientation, and safe-pointing during anomalies. Sun sensors operate autonomously without complex processing, making them a robust backup to star trackers.}{catAE}
\glsadd{Magnetorquer}\dictentry{Magnetorquer}{A magnetorquer is an electromagnetic coil that generates a magnetic dipole moment when current flows, producing torque against a planetary magnetic field. Magnetorquers provide attitude control in low Earth orbit where the geomagnetic field is strong enough for useful torque. They serve as primary actuators for simple spacecraft and momentum management in systems using reaction wheels.}{catAE}
\glsadd{Docking-Port}\dictentry{Docking Port}{A docking port is a mechanical and electrical interface that enables two spacecraft to connect in orbit. It features alignment guides, latching mechanisms, and data and power transfer links. Docking ports allow crew transfer, cargo resupply, and assembly of modular space stations. Standardized designs such as the International Docking Adapter enable compatibility across different spacecraft.}{catAE}
\glsadd{Airlock}\dictentry{Airlock}{An airlock is a chamber enabling passage between environments of different pressures, such as between a pressurized spacecraft and the vacuum of space. It features two independent seals that allow one side to depressurize while the other remains habitable. Airlocks are essential for conducting spacewalks and transferring equipment between interior and exterior. They minimize atmosphere loss and reduce preparation time for extravehicular activity.}{catAE}
\glsadd{Orbit}\dictentry{Orbit}{An orbit is the curved path of an object around a celestial body, maintained by the balance of gravitational pull and the object's velocity. Orbits are described by parameters such as semi-major axis, eccentricity, and inclination. They can be circular, elliptical, parabolic, or hyperbolic depending on the object's energy. Understanding orbits is fundamental to mission planning and satellite operations.}{catAE}
\glsadd{Elliptical-Orbit}\dictentry{Elliptical Orbit}{An elliptical orbit is an orbital path with eccentricity between zero and one, causing the orbiting body to vary its distance from the central body. The point closest to the primary is called periapsis and the farthest is apoapsis. Kepler's laws govern the timing and speed variations along the path. Elliptical orbits are commonly used for transfer maneuvers and highly optimized satellite constellations.}{catAE}
\glsadd{Circular-Orbit}\dictentry{Circular Orbit}{A circular orbit is an orbital path with an eccentricity of exactly zero, maintaining a constant altitude above the central body. The orbital velocity is uniform throughout the path, simplifying calculations and operations. Many communication and Earth observation satellites use circular orbits for consistent coverage. Maintaining a true circular orbit requires precise injection and occasional station-keeping.}{catAE}
\glsadd{Geostationary-Orbit}\dictentry{Geostationary Orbit}{A Geostationary Orbit, or GEO, is circular at about 35,786 kilometers directly above Earth's equator with zero inclination. A GEO satellite appears stationary relative to the ground, ideal for continuous communications, broadcasting, and weather monitoring. Proposed by Arthur C. Clarke in 1945, it is now crowded and requires periodic station-keeping. The limited orbital slots are managed through international coordination.}{catAE}
\glsadd{Low-Earth-Orbit}\dictentry{Low Earth Orbit}{Low Earth Orbit, or LEO, ranges from roughly 160 to 2,000 kilometers altitude. Objects travel at about 7.8 kilometers per second, completing an orbit in 90 to 120 minutes. The ISS, many Earth-observation satellites, and communication constellations operate here. LEO experiences atmospheric drag causing gradual decay without reboost. It is the most crowded orbital regime, making debris management increasingly important.}{catAE}
\glsadd{Medium-Earth-Orbit}\dictentry{Medium Earth Orbit}{Medium Earth Orbit, or MEO, ranges from about 2,000 to 35,786 kilometers altitude. Navigation constellations including GPS, GLONASS, and Galileo are its most prominent application. MEO satellites have orbital periods of several hours, providing global coverage with small constellations. This regime is above most atmospheric drag, allowing longer orbital lifetimes. It is also used for some communication and Earth-observation systems.}{catAE}
\glsadd{Polar-Orbit}\dictentry{Polar Orbit}{A Polar Orbit has approximately 90-degree inclination, passing over or near both geographic poles. Satellites travel north-south as Earth rotates beneath, covering the entire surface. Typically at 600 to 800 kilometers, these orbits serve Earth-observation, reconnaissance, and weather monitoring. Complete global coverage makes them essential for mapping and environmental monitoring. Launching requires more energy due to no rotational assist.}{catAE}
\glsadd{Sun-Synchronous-Orbit}\dictentry{Sun-Synchronous Orbit}{A Sun-Synchronous Orbit is a near-polar orbit maintaining a constant angle relative to the Sun, ensuring consistent local solar time passes over any point. This provides uniform illumination for Earth-observation and remote sensing. The orbit achieves this through controlled nodal precession from Earth's oblateness. Typical altitudes are 600 to 800 kilometers. It is widely used for weather, land-monitoring, and scientific missions.}{catAE}
\glsadd{Transfer-Orbit}\dictentry{Transfer Orbit}{A transfer orbit is an intermediate trajectory used to move a spacecraft between two different orbits. It bridges the initial and target orbits using one or more propulsive maneuvers. The choice of transfer orbit depends on fuel efficiency, time constraints, and mission objectives. Common examples include Hohmann transfers, bi-elliptic transfers, and phasing orbits.}{catAE}
\glsadd{Hohmann-Transfer}\dictentry{Hohmann Transfer}{A Hohmann transfer is the most fuel-efficient two-impulse maneuver for moving between two coplanar circular orbits. It uses an elliptical path that is tangent to both the departure and arrival orbits. The first burn raises the apogee to the target altitude and the second circularizes at the destination. It is widely used for orbital rendezvous and interplanetary trajectory design.}{catAE}
\glsadd{Orbital-Inclination}\dictentry{Orbital Inclination}{Orbital inclination is the angle between the orbital plane and the Earth's equatorial plane, measured at the ascending node. An inclination of zero degrees indicates an equatorial orbit, while ninety degrees indicates a polar orbit. The inclination determines the latitudinal range a satellite can observe directly. Launch sites near the equator offer the most flexibility for reaching a wide range of inclinations.}{catAE}
\glsadd{Eccentricity}\dictentry{Eccentricity}{Eccentricity is a dimensionless parameter measuring an orbit's deviation from a circle. Zero represents a circular orbit, while values between zero and one describe ellipses. One corresponds to parabolic escape, and greater than one to hyperbolic flyby orbits. It determines the ratio between minimum and maximum orbital distances from the central body. One of six classical orbital elements, it is key to orbit classification and mission design.}{catAE}
\glsadd{Orbital-Period}\dictentry{Orbital Period}{Orbital period is the time required to complete one full revolution around a central body. For circular orbits it depends only on the semi-major axis and central body mass, as described by Kepler's third law. Low Earth orbit satellites have periods near ninety minutes, while geosynchronous satellites orbit in one sidereal day. The period governs revisit times, communication windows, and ground track patterns.}{catAE}
\glsadd{Guidance}\dictentry{Guidance}{Guidance is the process of determining the desired trajectory and steering commands for a vehicle during flight. It involves computing the optimal path from the current state to the target destination in real time. Modern guidance systems use onboard computers and inertial sensors to execute pre-programmed or adaptive algorithms. Effective guidance is critical for orbit insertion, docking, and precision landing.}{catAE}
\glsadd{Navigation}\dictentry{Navigation}{Navigation is the process of determining a vehicle's current position, velocity, and orientation relative to a reference frame. It uses data from sensors such as inertial measurement units, star trackers, GPS receivers, and radar. Navigation solutions are updated continuously to correct for drift and measurement errors. Accurate navigation is essential for trajectory control, rendezvous operations, and surface landing.}{catAE}
\glsadd{Attitude-Control}\dictentry{Attitude Control}{Attitude control is the management of a spacecraft's angular orientation in three-dimensional space. It uses actuators such as reaction wheels, control moment gyroscopes, thrusters, and magnetorquers to apply torque. The control system maintains or adjusts the pointing direction of instruments, antennas, and solar panels. Precise attitude control is fundamental to mission success for imaging, communications, and scientific observation.}{catAE}
\glsadd{Attitude-Determination}\dictentry{Attitude Determination}{Attitude determination calculates a spacecraft's orientation using star trackers, sun sensors, magnetometers, and gyroscopes. Raw measurements are processed through least-squares or Kalman filter algorithms to compute a three-axis solution. Accurate determination is essential for pointing instruments, antennas, and solar panels at their targets.}{catAE}
\glsadd{PID-Controller}\dictentry{PID Controller}{A PID controller is a feedback control loop that calculates a correction based on proportional, integral, and derivative terms of the error signal. The proportional term responds to current error, the integral eliminates accumulated offset, and the derivative dampens oscillations. It is the most widely used control algorithm in aerospace and industrial systems. Its simplicity makes it suitable for a broad range of dynamic systems.}{catAE}
\glsadd{Kalman-Filter}\dictentry{Kalman Filter}{A Kalman filter is a recursive algorithm that estimates the state of a dynamic system from a series of noisy measurements. It combines a predictive model with observation data to produce an optimal estimate in a least-squares sense. The filter continuously updates its estimate and covariance as new data arrives. It is fundamental to navigation, guidance, and sensor fusion in aerospace applications.}{catAE}
\glsadd{State-Estimation}\dictentry{State Estimation}{State estimation is the mathematical process of inferring the full state of a system from partial and noisy sensor observations. It uses models of the system dynamics along with measurement data to produce the best possible estimate. Techniques include Kalman filtering, extended Kalman filtering, and particle filtering. Accurate state estimation is essential for control, health monitoring, and decision-making.}{catAE}
\glsadd{Sensor-Fusion}\dictentry{Sensor Fusion}{Sensor fusion is the technique of combining data from multiple sensors to produce information that is more accurate and reliable than any single source alone. It reduces uncertainty by exploiting complementary strengths and compensating for individual sensor weaknesses. Common approaches include weighted averaging, Bayesian estimation, and Kalman filtering. It is critical for navigation, perception, and situational awareness in aerospace systems.}{catAE}
\glsadd{Control-Law}\dictentry{Control Law}{A control law is a mathematical equation or algorithm that maps system states and reference inputs to control outputs. It defines how the controller responds to errors and disturbances to achieve the desired behavior. Control laws can be linear or nonlinear, time-invariant or adaptive. The design of the control law determines system stability, performance, and robustness.}{catAE}
\glsadd{Microgravity}\dictentry{Microgravity}{Microgravity is the condition of apparent weightlessness experienced during free fall or orbit. Despite the name, gravity in orbit is only slightly weaker; the effect comes from continuous free fall. It enables unique experiments in fluid physics, materials science, and biology. Research aboard the ISS prepares for long-duration missions beyond Earth orbit. Prolonged exposure causes bone density loss and muscle atrophy in astronauts.}{catAE}
\glsadd{Space-Vacuum}\dictentry{Space Vacuum}{Space vacuum refers to the extremely low ambient pressure found beyond Earth's atmosphere, where particle density is negligible. In this environment, heat transfer occurs only through radiation, and materials can outgas or sublimate. The vacuum poses challenges for thermal management, lubrication, and electrical discharge. Understanding space vacuum conditions is essential for spacecraft design and material selection.}{catAE}
\glsadd{Space-Weather}\dictentry{Space Weather}{Space weather describes the dynamic environmental conditions in near-Earth space driven by solar activity. It includes solar flares, coronal mass ejections, geomagnetic storms, and energetic particle events. Space weather can disrupt satellite operations, communications, navigation systems, and power grids. Monitoring and forecasting space weather is critical for protecting space-based infrastructure and crew.}{catAE}
\glsadd{Solar-Wind}\dictentry{Solar Wind}{The Solar Wind is a continuous stream of charged particles, mainly protons and electrons, flowing from the Sun's corona at 300 to 800 kilometers per second. It carries the Sun's magnetic field into interplanetary space and shapes planetary magnetospheres. Solar wind intensity variations drive space weather effects on satellites and power systems. It extends to the heliopause where it meets the interstellar medium.}{catAE}
\glsadd{Van-Allen-Belts}\dictentry{Van Allen Belts}{The Van Allen belts are two toroidal zones of energetic charged particles trapped by Earth's magnetic field. The inner belt contains mainly protons and the outer belt primarily electrons. These radiation zones pose hazards to spacecraft electronics and astronaut health. Satellite orbits must be designed to minimize exposure or include appropriate radiation shielding.}{catAE}
\glsadd{Cosmic-Ray}\dictentry{Cosmic Ray}{Cosmic rays are high-energy particles originating from outside the solar system, including protons, heavy nuclei, and electrons. They travel at relativistic speeds and can penetrate spacecraft hulls, causing single-event effects and long-term degradation. Galactic cosmic rays are modulated by the solar cycle and Earth's magnetosphere. Shielding against cosmic rays remains one of the major challenges for long-duration human spaceflight.}{catAE}
\glsadd{Space-Plasma}\dictentry{Space Plasma}{Space plasma is a hot, ionized gas composed of free electrons, ions, and neutral particles found in the magnetosphere, solar wind, and interstellar medium. Plasma behaves collectively under electromagnetic forces and exhibits phenomena such as magnetic reconnection. It affects spacecraft charging, communications, and navigation signals. Understanding plasma physics is essential for predicting space environment effects on missions.}{catAE}
\glsadd{Space-Debris}\dictentry{Space Debris}{Space debris consists of non-functional artificial objects in Earth orbit, including spent rocket stages, defunct satellites, and fragmentation remnants. The debris population is growing and poses collision risks to operational spacecraft and the International Space Station. Even small particles cause significant damage at orbital velocities. Active debris removal and mitigation guidelines are critical for the sustainability of space activities.}{catAE}
\glsadd{Meteoroid}\dictentry{Meteoroid}{A Meteoroid is a small rocky or metallic body traveling through interplanetary space, typically smaller than an asteroid but larger than a molecule. They originate from asteroid collisions, cometary debris, or planetary impact ejecta. When entering Earth's atmosphere, they produce visible meteors or shooting stars. Sizes range from sand grains to objects about one meter across. They are among the most common small bodies near Earth.}{catAE}
\glsadd{Micrometeoroid}\dictentry{Micrometeoroid}{A micrometeoroid is an extremely small particle in space, typically less than one millimeter, traveling at velocities of several kilometers per second. Despite their tiny size, they can damage spacecraft surfaces, optics, and solar panels through hypervelocity impacts. MMOD shielding uses layered materials to break up and absorb particle energy. Monitoring the micrometeoroid environment supports risk assessment for long-duration missions.}{catAE}
\glsadd{Additive-Manufacturing}\dictentry{Additive Manufacturing}{Additive manufacturing builds three-dimensional objects by depositing material layer by layer based on a digital model. In aerospace, it is used to produce complex structural components, engine parts, and tooling with reduced weight and material waste. Processes include selective laser melting, electron beam melting, and fused deposition modeling. Additive manufacturing enables rapid prototyping and on-demand production of flight hardware.}{catAE}
\glsadd{CNC-Machining}\dictentry{CNC Machining}{CNC machining is a subtractive manufacturing process in which computer-controlled tools remove material from a solid block to create precise parts. It is widely used in aerospace for producing structural components, engine parts, and fixtures from metals and composites. CNC offers high accuracy, repeatability, and the ability to work with a wide range of materials. It remains the backbone of precision manufacturing for flight-critical hardware.}{catAE}
\glsadd{Friction-Stir-Welding}\dictentry{Friction Stir Welding}{Friction stir welding is a solid-state joining process that uses a rotating tool to generate frictional heat and plasticize material without melting it. The tool traverses along the joint line, forging the two pieces together. It produces joints with superior mechanical properties compared to fusion welding, especially in aluminum alloys. It is used extensively in aerospace structural assembly, including fuel tanks and fuselage panels.}{catAE}
\glsadd{Electron-Beam-Welding}\dictentry{Electron Beam Welding}{Electron beam welding uses a focused beam of high-velocity electrons in a vacuum to join metals. The kinetic energy of the electrons is converted to heat upon impact, creating deep, narrow welds with minimal heat-affected zones. It is capable of joining dissimilar metals and thick sections with high precision. The process is commonly used in aerospace for turbine blades, landing gear, and structural assemblies.}{catAE}
\glsadd{Autoclave-Curing}\dictentry{Autoclave Curing}{Autoclave curing is a composite manufacturing process that subjects parts to elevated temperature and pressure inside a sealed vessel. The combination of heat and pressure consolidates laminate layers, removes voids, and cures the resin matrix. It produces high-quality, low-porosity composite structures with excellent mechanical properties. Autoclave-cured composites are used extensively in aircraft fuselages, wings, and control surfaces.}{catAE}
\glsadd{Resin-Transfer-Molding}\dictentry{Resin Transfer Molding}{Resin transfer molding is a closed-mold composite manufacturing process in which dry fiber preforms are placed in a mold and resin is injected under pressure. The resin flows through the fibers, wetting them and filling the mold cavity before curing. It offers good surface finish on both sides and dimensional accuracy. RTM is used for producing aerospace structural components with consistent quality and moderate production rates.}{catAE}
\glsadd{Filament-Winding}\dictentry{Filament Winding}{Filament winding is a composite manufacturing process in which resin-impregnated fibers are wound under tension around a rotating mandrel. The winding pattern and fiber angle are controlled to optimize structural performance for the intended loading. After curing, the mandrel is removed to leave a hollow composite structure. It is widely used for pressure vessels, rocket motor cases, and cylindrical aerospace structures.}{catAE}
\glsadd{Fatigue-Test}\dictentry{Fatigue Test}{A fatigue test evaluates the durability of a material or component by subjecting it to repeated cyclic loading over many cycles. The test determines the fatigue life, or the number of cycles a part can withstand before failure. In aerospace, fatigue testing is critical for ensuring structural integrity of airframes and engine components over their service life. Results guide inspection intervals and retirement criteria.}{catAE}
\glsadd{Flutter-Test}\dictentry{Flutter Test}{A flutter test is an experimental procedure to determine the speed at which aeroelastic flutter occurs in an aircraft structure. Flutter is a self-excited oscillation caused by the coupling of aerodynamic, elastic, and inertial forces. The test involves flying the aircraft at progressively higher speeds while monitoring structural response. Flutter boundaries define the safe operating envelope and validate aerodynamic and structural models.}{catAE}
\glsadd{Vibration-Test}\dictentry{Vibration Test}{A vibration test subjects a spacecraft or component to the mechanical vibration environment expected during launch. The test simulates the random and sinusoidal vibrations produced by the rocket engines and aerodynamic forces. It verifies structural integrity and the functionality of all components under launch loads. Vibration testing is performed on all flight hardware to ensure survival through the launch phase.}{catAE}
\glsadd{Thermal-Vacuum-Test}\dictentry{Thermal Vacuum Test}{A thermal vacuum test simulates the combined thermal and vacuum conditions a spacecraft will experience in space. The test article is placed in a vacuum chamber and subjected to extreme hot and cold temperature cycles. This verifies that all components function correctly under realistic space conditions. Thermal vacuum testing is typically the most comprehensive environmental test in a spacecraft program.}{catAE}
\glsadd{Structural-Test}\dictentry{Structural Test}{A structural test applies loads to an aerospace component or assembly to verify that it meets design strength and stiffness requirements. Tests include static proof loading, ultimate strength testing, and deflection measurements. Results validate analytical models and confirm structural margins of safety. Structural testing is required before first flight and at intervals throughout the production program.}{catAE}
\glsadd{Flight-Test}\dictentry{Flight Test}{A flight test is an in-flight evaluation of an aircraft or spacecraft to verify performance, stability, and handling qualities under real operating conditions. It progresses through phases from envelope expansion to systems integration and certification testing. Flight tests are conducted by specialized test pilots and engineers using instrumented aircraft. The data collected validates design assumptions and ensures compliance with standards.}{catAE}
\glsadd{Qualification-Test}\dictentry{Qualification Test}{A qualification test demonstrates that a design meets all specified performance and environmental requirements under worst-case conditions. It subjects hardware to conditions more severe than expected in service, including thermal, vibration, shock, and electromagnetic environments. Passing qualification tests grants approval for the design to enter production. It is a key milestone in aerospace development programs.}{catAE}
\glsadd{Acceptance-Test}\dictentry{Acceptance Test}{An acceptance test verifies that each production unit meets its individual specifications before delivery. Unlike qualification testing, acceptance tests are performed on every unit under nominal conditions. They confirm manufacturing quality, functional performance, and dimensional compliance. Acceptance testing ensures traceability and confidence in every piece of flight hardware.}{catAE}
\glsadd{Redundancy}\dictentry{Redundancy}{Redundancy is the inclusion of backup components or systems that can take over if the primary system fails. In aerospace, redundancy is implemented at hardware, software, and functional levels to increase reliability and safety. Common configurations include dual and triple modular redundancy with voting logic. It is a fundamental design principle for flight-critical systems in both aircraft and spacecraft.}{catAE}
\glsadd{Fail-Safe-Design}\dictentry{Fail-Safe Design}{Fail-safe design is an engineering philosophy that ensures a system will fail in a predictable and manageable way rather than catastrophically. It incorporates features such as load paths, containment, and graceful degradation to protect crew and equipment. Fail-safe principles are applied in airframe design, propulsion systems, and flight control architecture. They form a core requirement of airworthiness certification.}{catAE}
\glsadd{Fault-Tolerance}\dictentry{Fault Tolerance}{Fault tolerance is the ability of a system to continue operating correctly even when one or more components have failed. It is achieved through redundancy, error detection, and automatic reconfiguration. In aerospace, fault-tolerant design is essential for flight control computers, navigation systems, and propulsion controllers. It enables missions to proceed safely despite hardware or software anomalies.}{catAE}
\glsadd{MTBF}\dictentry{MTBF}{MTBF is the mean time between failures, a statistical measure of a system's reliability expressed in operating hours. It represents the average expected time a non-repairable component will function before failure. MTBF is calculated from reliability testing data and used to predict maintenance needs and spare parts requirements. Higher MTBF values indicate greater reliability and lower lifecycle cost.}{catAE}
\glsadd{Reliability}\dictentry{Reliability}{Reliability is the probability that a system or component will perform its intended function without failure for a specified period under stated conditions. It is quantified through reliability testing, field data analysis, and predictive modeling. In aerospace, reliability requirements are among the most stringent of any industry. Achieving high reliability drives design choices in materials, redundancy, manufacturing, and testing.}{catAE}
\glsadd{Risk-Assessment}\dictentry{Risk Assessment}{Risk assessment is the systematic process of identifying, analyzing, and evaluating potential hazards and their consequences. It combines probability estimation with impact severity to prioritize risks and guide mitigation efforts. In aerospace, risk assessment informs design decisions, test planning, and operational procedures. It is a continuous activity performed throughout the lifecycle of a program.}{catAE}
\glsadd{Hazard-Analysis}\dictentry{Hazard Analysis}{Hazard analysis is a structured examination of a system to identify conditions or events that could lead to harm to people, equipment, or the environment. Techniques include fault tree analysis, failure modes and effects analysis, and hazard and operability studies. The analysis determines root causes and evaluates the effectiveness of safeguards. It is a mandatory step in aerospace safety certification.}{catAE}
\glsadd{Flight-Safety}\dictentry{Flight Safety}{Flight safety encompasses the policies, procedures, and engineering measures that protect aircraft, crew, passengers, and the public during flight operations. It includes hazard identification, risk mitigation, accident investigation, and safety management systems. Flight safety is governed by regulatory authorities and enforced through certification and operational standards. A strong safety culture is fundamental to the aerospace industry.}{catAE}
\glsadd{Range-Safety}\dictentry{Range Safety}{Range Safety encompasses procedures, systems, and personnel dedicated to protecting the public during rocket and missile launch operations. It includes trajectory analysis, hazard area calculation, real-time tracking, and the authority to terminate flight. Range safety officers monitor every launch and can activate destruct systems if the vehicle leaves its approved corridor. These measures are mandated by national and international regulations.}{catAE}
\glsadd{Airworthiness}\dictentry{Airworthiness}{Airworthiness is the condition of an aircraft that meets its design and safety standards and is safe to fly. It is established through type certification, production conformity, and ongoing maintenance and inspection. Airworthiness directives are issued to address known safety issues. Maintaining airworthiness is a legal requirement for all aircraft in commercial and private operations.}{catAE}
\glsadd{Type-Certificate}\dictentry{Type Certificate}{A type certificate is an official document issued by an aviation authority confirming that an aircraft design meets all applicable airworthiness and environmental standards. It covers the complete design including structure, systems, and performance. Once issued, the certificate permits mass production and operational use of the type. Modifications require supplemental type certificates or amendments.}{catAE}
\glsadd{Supplemental-Type-Certificate}\dictentry{Supplemental Type Certificate}{A supplemental type certificate is an approval issued for major modifications or repairs to an already type-certificated aircraft design. It certifies that the modification does not compromise the airworthiness of the original design. STCs enable equipment upgrades, engine changes, and structural modifications throughout an aircraft's service life. They are issued by national aviation authorities after thorough design review.}{catAE}
\glsadd{FAA}\dictentry{FAA}{The Federal Aviation Administration (FAA) is the U.S. agency regulating civil aviation including unmanned aircraft. It establishes registration requirements, pilot certification, airspace rules, and operational limitations. The FAA manages the LAANC system and Part 107 framework. It issues waivers for beyond-visual-line-of-sight and other non-standard operations.}{catAE}
\glsadd{EASA}\dictentry{EASA}{The European Union Aviation Safety Agency (EASA) regulates civil aviation safety across EU states including drones. EASA established the Open, Specific, and Certified category framework calibrated by risk. It defines drone class markings C0 through C4, pilot competency requirements, and uniform operational rules. Its framework is increasingly adopted as an international reference model.}{catAE}
\glsadd{ICAO}\dictentry{ICAO}{The International Civil Aviation Organization (ICAO) is a UN agency setting global aviation standards including unmanned aircraft integration. ICAO publishes Annex 13 for drone safety and develops the UTM framework used by national regulators. It coordinates harmonized regulations for cross-border operations. ICAO standards form the baseline for national authorities like the FAA and EASA.}{catAE}
\glsadd{DO-178C}\dictentry{DO-178C}{DO-178C is the Software Considerations in Airborne Systems and Equipment Certification standard. It defines the processes, activities, and objectives for developing airborne software that meets airworthiness requirements. The standard specifies five software levels based on the severity of failure conditions. Compliance with DO-178C is mandatory for all commercial aviation software and is widely adopted in defense and space.}{catAE}
\glsadd{DO-254}\dictentry{DO-254}{DO-254 is the Design Assurance Guidance for Airborne Electronic Hardware. It provides a structured framework for the design, development, and verification of complex electronic hardware in aircraft. The standard defines design assurance levels corresponding to the criticality of hardware functions. Compliance with DO-254 is required for programmable logic devices, ASICs, and custom electronic components in certified aircraft.}{catAE}
\glsadd{Urban-Air-Mobility}\dictentry{Urban Air Mobility}{Urban air mobility is a concept for air transportation services operating within and between cities at low altitudes. It envisions networks of electric vertical takeoff and landing aircraft carrying passengers and cargo on short urban routes. UAM aims to reduce ground traffic congestion and provide rapid point-to-point transit. It requires advances in autonomous flight, battery technology, and airspace management.}{catAE}
\glsadd{Electric-Aircraft}\dictentry{Electric Aircraft}{An electric aircraft is powered by electric motors that draw energy from batteries, fuel cells, or other electrical storage systems. Electric propulsion offers reduced noise, lower emissions, and simplified maintenance compared to conventional engines. Current battery energy density limits range and payload, driving research into hybrid architectures. Electric aircraft are emerging for short-range regional flights and urban air mobility.}{catAE}
\glsadd{Hybrid-Electric-Propulsion}\dictentry{Hybrid-Electric Propulsion}{Hybrid-electric propulsion combines conventional thermal engines with electric motors and energy storage systems. The thermal engine can drive a generator to power electric motors or provide direct thrust, while batteries handle peak loads. This architecture reduces fuel consumption and emissions while extending range beyond pure electric capability. It is a key enabling technology for next-generation regional aircraft.}{catAE}
\glsadd{Hydrogen-Aircraft}\dictentry{Hydrogen Aircraft}{A hydrogen aircraft uses hydrogen as its primary energy source, either through fuel cells producing electricity or direct combustion in modified turbine engines. Hydrogen offers high specific energy and produces only water vapor when used in fuel cells. The main challenges include onboard storage at cryogenic temperatures and airframe redesign for large fuel tanks. Hydrogen aircraft represent a pathway toward zero-emission long-range aviation.}{catAE}
\glsadd{Autonomous-Flight}\dictentry{Autonomous Flight}{Autonomous flight refers to aircraft operations conducted without direct human pilot intervention, relying on onboard artificial intelligence and control systems. It encompasses capabilities from automated takeoff and landing to fully autonomous mission execution. These systems perceive the environment, make decisions, and execute control actions using sensors and algorithms. It promises to enhance safety and enable new operational concepts.}{catAE}
\glsadd{Digital-Twin}\dictentry{Digital Twin}{A Digital Twin is a real-time virtual replica of a physical drone mirroring its state and environmental conditions. It receives telemetry and uses simulation models to predict behavior and detect anomalies. Digital twins enable predictive maintenance and mission rehearsal in virtual environments. They are becoming central to fleet management and autonomous system certification.}{catAE}
\glsadd{Morphing-Wing}\dictentry{Morphing Wing}{A morphing wing is an aircraft wing capable of changing its shape in flight to optimize aerodynamic performance for different conditions. Morphing mechanisms can alter sweep, span, camber, and twist using flexible structures, actuators, and smart materials. The goal is to improve fuel efficiency, extend the flight envelope, and reduce mechanical complexity. Morphing wing technology draws on advances in materials science and adaptive control.}{catAE}
\glsadd{Reusable-Launch-Vehicle}\dictentry{Reusable Launch Vehicle}{A Reusable Launch Vehicle is a rocket system designed to recover and refly one or more stages, dramatically reducing the cost of space access. Reusability requires robust thermal protection, landing systems, and inspection and refurbishment infrastructure. The Falcon 9 first stage has demonstrated reliable reuse across dozens of missions.}{catAE}
\glsadd{Cockpit}\dictentry{Cockpit}{The cockpit is the compartment at the front of an aircraft from which the pilot controls the flight. It contains the instrument panel, flight controls, communication radios, navigation displays, and avionics systems. The cockpit layout is designed for ergonomic efficiency and situational awareness. In modern aircraft, glass cockpit displays have replaced traditional analog instruments with integrated digital screens.}{catAE}
\glsadd{Canopy}\dictentry{Canopy}{The canopy is the transparent enclosure over the cockpit that provides the pilot with external visibility while maintaining a pressurized environment. It is typically made of acrylic or polycarbonate laminates designed to withstand bird strikes, pressure loads, and optical distortion. The canopy may be fixed, hinged, sliding, or removable for egress. Its design balances visibility, structural strength, weight, and emergency escape requirements.}{catAE}
\glsadd{Nose-Gear}\dictentry{Nose Gear}{The nose gear is the forward landing gear assembly of an aircraft, consisting of one or more wheels mounted on a strut. It supports the front portion of the aircraft on the ground and provides steering capability during taxi. Nose gear assemblies include shock absorbers, retraction mechanisms, and steering actuators. Proper design and maintenance are essential for safe ground handling and landing operations.}{catAE}
\glsadd{Main-Gear}\dictentry{Main Gear}{The main gear is the primary landing gear assembly that supports the majority of an aircraft's weight during ground operations. It is typically located near the center of gravity and consists of one or more wheel bogies on shock-absorbing struts. The main gear absorbs landing impact loads and houses the primary braking system. Its design must accommodate high loads and provide directional stability during rollout.}{catAE}
\glsadd{Tailwheel}\dictentry{Tailwheel}{A tailwheel is a small wheel or caster located at the rear of a conventional landing gear configuration. It supports the tail section while the aircraft is on the ground and allows steering through differential braking or linkage to the rudder pedals. Tailwheel landing gear is common on bush planes and aerobatic aircraft due to its lighter weight and lower drag. It requires more pilot skill compared to tricycle gear.}{catAE}
\glsadd{Tricycle-Gear}\dictentry{Tricycle Gear}{Tricycle gear is a landing gear configuration with a steerable nose wheel at the front and two main gear units behind the center of gravity. It provides excellent ground visibility, simpler ground handling, and resistance to ground looping. Tricycle gear is the standard configuration for modern commercial and military aircraft. It facilitates safer takeoff and landing operations, especially for less experienced pilots.}{catAE}
\glsadd{Brake}\dictentry{Brake}{An aircraft brake is a friction-based device that converts kinetic energy into heat to decelerate the wheels during landing and taxi. Disc brakes made of steel or carbon are the most common type, actuated by hydraulic pressure from the pilot's toe pedals. Anti-skid systems prevent wheel lockup to maintain directional control and maximize braking efficiency. Brake design must handle extreme thermal loads under all conditions.}{catAE}
\glsadd{Parking-Brake}\dictentry{Parking Brake}{A parking brake is a mechanical or hydraulic device that locks the aircraft's wheel brakes to prevent movement while parked. It is typically set by the pilot before shutdown and released before taxi. The parking brake ensures the aircraft remains stationary on sloped surfaces and in windy conditions. It is a simple but essential safety system for ground operations.}{catAE}
\glsadd{Control-Column}\dictentry{Control Column}{The control column is the primary flight control input device in the cockpit, used to command pitch and roll of the aircraft. In fixed-wing aircraft, it is a column with a wheel or stick at the top connected to flight control surfaces through cables, pushrods, or fly-by-wire systems. The control column provides tactile feedback and allows the pilot to apply precise inputs. It has been standard since the earliest days of powered flight.}{catAE}
\glsadd{Yoke}\dictentry{Yoke}{A yoke is a wheel-shaped control device mounted on the control column of many fixed-wing aircraft. Rotating the yoke commands roll, while pushing or pulling it commands pitch. The yoke provides a natural and ergonomic interface for sustained two-axis control. It is commonly found in general aviation and commercial transport aircraft.}{catAE}
\glsadd{Cyclic}\dictentry{Cyclic}{The cyclic is the primary flight control for a helicopter, used to tilt the main rotor disk and control the aircraft's direction of horizontal movement. Moving the cyclic changes the pitch of the rotor blades cyclically as they rotate, creating a thrust vector in the desired direction. It is mounted between the pilot's legs or to the side of the seat. Mastering the cyclic is fundamental to helicopter piloting.}{catAE}
\glsadd{Collective}\dictentry{Collective}{The collective is a helicopter control that simultaneously changes the pitch angle of all main rotor blades by the same amount. Raising the collective increases lift to climb, while lowering it decreases lift to descend. It is typically located to the left of the pilot's seat and connected to the engine throttle for coordinated power management. The collective provides the primary means of altitude control in helicopters.}{catAE}
\glsadd{Throttle}\dictentry{Throttle}{A throttle is the control lever or mechanism that regulates the power output of an engine by adjusting fuel flow, air intake, or propeller pitch. In aircraft, the throttle is typically located on the center console or overhead panel for easy pilot access. Precise throttle management is essential for maintaining desired airspeed, altitude, and fuel efficiency. Modern engines use digital control to optimize performance automatically.}{catAE}
\glsadd{Mixture-Control}\dictentry{Mixture Control}{The mixture control is a cockpit lever that adjusts the ratio of fuel to air entering a piston aircraft engine. Leaning the mixture reduces fuel flow at altitude where air density is lower, improving efficiency and reducing engine wear. Rich mixture provides more power for takeoff and cooling during high-power operation. Proper mixture management is a critical skill for piston-engine aircraft pilots.}{catAE}
\glsadd{Propeller-Control}\dictentry{Propeller Control}{Propeller control refers to the system that regulates the rotational speed and blade pitch of a constant-speed propeller. The pilot sets a desired RPM with the propeller control lever, and a governor adjusts blade pitch automatically to maintain that speed under varying flight conditions. This optimizes propeller efficiency across the entire flight envelope. Constant-speed propellers provide better performance than fixed-pitch designs.}{catAE}
\glsadd{Trim-Tab}\dictentry{Trim Tab}{A trim tab is a small adjustable surface attached to a primary flight control surface that reduces the pilot's control forces. By deflecting the tab, an aerodynamic force is created that holds the control surface in the desired position without continuous pilot input. Trim tabs can be manually adjusted from the cockpit or automatically controlled by the autopilot. They are essential for maintaining stable flight with minimal workload.}{catAE}
\glsadd{Balance-Tab}\dictentry{Balance Tab}{A balance tab is a small hinged surface attached to a primary control surface that moves in the opposite direction to reduce aerodynamic forces required to deflect the control surface. Unlike a trim tab, it is not independently adjustable and functions automatically as part of the control surface assembly. Balance tabs reduce control forces during all phases of flight. They are found on aircraft with large control surfaces.}{catAE}
\glsadd{Anti-Servo-Tab}\dictentry{Anti-Servo Tab}{An anti-servo tab is a control surface attached to the trailing edge of a primary control surface that moves in the same direction as the control surface deflection. It increases the aerodynamic force opposing the pilot's input, providing artificial feel and enhancing stability. By resisting excessive control movement, the anti-servo tab prevents overcontrol and helps maintain precise handling qualities. These tabs are essential on light aircraft where direct mechanical control forces would otherwise be too light at certain airspeeds.}{catAE}
\glsadd{Mass-Balance}\dictentry{Mass Balance}{Mass balance involves attaching weights to control surfaces or their linkage systems to counterbalance the aerodynamic and inertial forces acting ahead of the hinge line. This redistribution of mass ensures the center of gravity of the control surface is at or near the hinge axis. Proper mass balancing prevents flutter, a dangerous self-excited oscillation that can destroy control surfaces at high speed. Every control surface on an aircraft must be mass balanced within strict tolerances specified by certification standards.}{catAE}
\glsadd{Aerodynamic-Balance}\dictentry{Aerodynamic Balance}{An aerodynamic balance is a portion of a control surface extending forward of the hinge line, allowing air pressure to create a moment opposing the pilot's deflection force. This reduces the stick force needed to move the control surface. The designer must carefully size the balance area because too large a balance can make controls overly sensitive or cause control reversal. Aerodynamic balancing is one of several methods used to achieve appropriate control force gradients.}{catAE}
\glsadd{Horn-Balance}\dictentry{Horn Balance}{A horn balance is a physical extension of the control surface that projects forward of the hinge line in the shape of a horn or paddle. When the surface deflects, airflow over the horn generates an aerodynamic force that assists the pilot, reducing required control forces. Horn balances are widely used on ailerons and rudders due to their simplicity and effectiveness. They must be carefully designed to avoid excessive sensitivity at high speeds or adverse flutter characteristics.}{catAE}
\glsadd{Boosted-Controls}\dictentry{Boosted Controls}{Boosted controls use hydraulic or electro-mechanical actuators to augment the pilot's manual input, reducing the physical effort required to move primary flight control surfaces. A mechanical linkage still connects the pilot's controls to the surface, but the booster provides most of the force needed for deflection. This technology is essential on large aircraft where pure mechanical advantage cannot provide adequate control forces. Boosted systems include feedback mechanisms that preserve a sense of aerodynamic load for the pilot.}{catAE}
\glsadd{Power-Controls}\dictentry{Power Controls}{Power controls are fully powered flight control systems in which actuators driven by hydraulic or electric power move the control surfaces with no direct mechanical connection to the pilot's controls. Commands are transmitted electrically or via fly-by-wire signals to the actuators. This design eliminates heavy mechanical linkages and allows flight envelope protection logic to be integrated. Power controls are standard on modern airliners and military aircraft, forming the foundation for fly-by-wire architectures.}{catAE}
\glsadd{Control-Horn}\dictentry{Control Horn}{A control horn is a rigid arm or bracket attached to a control surface that provides an attachment point for the control linkage, such as pushrods or cables. It extends the moment arm between the linkage attachment point and the hinge line, giving the pilot mechanical advantage for surface deflection. The geometry of the control horn determines the ratio of pilot input to surface travel, directly affecting control sensitivity. Control horns are fundamental components in all mechanically linked flight control systems.}{catAE}
\glsadd{Pushrod}\dictentry{Pushrod}{A pushrod is a rigid rod used in aircraft flight control systems to transmit mechanical force from the cockpit controls to the control surfaces. It operates primarily in compression and tension along its length, transferring the pilot's inputs through bellcranks and linkages. Pushrods are typically made of steel or aluminum tubing to provide high strength with low weight. They are a key element of direct mechanical control systems found on light aircraft and as part of larger linkage networks on heavier aircraft.}{catAE}
\glsadd{Cable}\dictentry{Cable}{A steel wire cable is used in aircraft control systems to transmit tensile forces over long distances between the cockpit and control surfaces. Cables run through fairleads and over pulleys to follow the aircraft structure, allowing control inputs to be routed around obstacles. Cable systems are lightweight, reliable, and have been a standard control method since the earliest days of aviation. Proper cable tension is critical for precise control response, and cables are regularly inspected as part of maintenance programs.}{catAE}
\glsadd{Bellerank}\dictentry{Bellerank}{A bellerank, also known as a bellcrank, is a pivoting lever used in aircraft control linkages to change the direction of a mechanical force or motion. Shaped like an L or V with a pivot at the vertex, it translates a push or pull input from one cable or pushrod into a perpendicular output to another. Bellcranks are essential for routing control movements around structural members where straight-line connections are not feasible. They are fabricated from aluminum or steel with specific arm ratios to achieve desired deflection ratios.}{catAE}
\glsadd{Quadrant}\dictentry{Quadrant}{A control quadrant is a sector-shaped lever mounted in the cockpit that the pilot moves through an arc to command engine power, propeller pitch, or other adjustable parameters. The curved guide path allows large angular inputs within a compact space, and detent positions mark key settings such as idle, cruise, and full power. Quadrants provide tactile and visual references for the current setting and have been a standard cockpit control element since the early days of aviation.}{catAE}
\glsadd{Linkage}\dictentry{Linkage}{A control linkage is the complete mechanical system of rods, cables, bellcranks, pulleys, and hinges that transmits the pilot's inputs from the cockpit controls to the aircraft's control surfaces. The linkage must be designed with minimal free play, appropriate friction, and adequate strength for precise and reliable control response. Redundant load paths are incorporated so that failure of a single component does not result in loss of control. Design and maintenance are governed by strict certification standards.}{catAE}
\glsadd{Control-Surface}\dictentry{Control Surface}{A control surface is a movable aerodynamic surface attached to the wing, tail, or fuselage that the pilot deflects to change the aircraft's attitude or flight path. By altering airflow, it generates moments about the center of gravity, producing pitch, roll, or yaw. Primary surfaces include elevators, ailerons, and the rudder, while secondary surfaces include flaps, trim tabs, and spoilers. Effectiveness depends on airspeed, surface area, and distance from the center of gravity.}{catAE}
\glsadd{Primary-Control}\dictentry{Primary Control}{Primary control surfaces are the main aerodynamic surfaces that allow the pilot to control rotation about all three axes: pitch through the elevator, roll through the ailerons, and yaw through the rudder. These surfaces must remain functional throughout the flight envelope and are designed with redundancy and structural margins far exceeding normal loads. Sizing and placement directly determine handling qualities, stability, and maneuverability. Loss of any primary control generally constitutes an emergency.}{catAE}
\glsadd{Secondary-Control}\dictentry{Secondary Control}{Secondary control surfaces are auxiliary aerodynamic devices that modify performance characteristics but are not used for direct attitude control. They include flaps for increasing lift during takeoff and landing, trim tabs for relieving steady control forces, and spoilers for reducing lift or increasing drag. While not essential for basic flight, secondary controls expand the operational envelope and reduce pilot workload. Many are operated automatically by flight control computers in modern aircraft.}{catAE}
\glsadd{High-Lift-Device}\dictentry{High-Lift Device}{A high-lift device is any aerodynamic mechanism, typically a flap or slat, that increases the maximum lift coefficient of a wing beyond its clean configuration value. By changing the wing's camber, area, or both, high-lift devices allow safe flight at lower speeds for takeoff and landing. Transport aircraft may increase lift coefficient by 50 to 100 percent through combined deployment of leading-edge slats and trailing-edge flaps. The complexity and weight are justified by the dramatic reduction in approach and takeoff speeds they enable.}{catAE}
\glsadd{Leading-Edge-Device}\dictentry{Leading Edge Device}{A leading edge device is an aerodynamic surface on or near the front edge of a wing that can be extended to increase lift at low speeds. Common types include slats, which create a slot gap allowing high-energy air to flow over the upper surface, and Krueger flaps, which increase camber. Leading-edge devices delay flow separation and stall to higher angles of attack, critical for maintaining control during approach and landing. They are typically retracted during cruise to minimize drag.}{catAE}
\glsadd{Trailing-Edge-Device}\dictentry{Trailing Edge Device}{A trailing edge device is an aerodynamic surface on the rear edge of a wing that can be deflected downward to increase lift and drag for low-speed flight. Flaps of various types, including plain, split, slotted, and Fowler designs, are the most common trailing edge devices. By increasing camber and sometimes chord area, they allow the aircraft to maintain lift at lower airspeeds. They are essential for reducing takeoff and landing distances and are deployed in stages depending on flight phase.}{catAE}
\glsadd{Fowler-Flap}\dictentry{Fowler Flap}{A Fowler flap is a trailing-edge flap that moves aft and downward on tracks, simultaneously increasing both camber and planform area. This dual effect produces a larger lift coefficient increase than camber-only designs, making Fowler flaps particularly effective for transport aircraft. The extended position increases wing area, reducing induced load per unit area and improving low-speed performance. Typically driven by screw jacks, Fowler flaps may incorporate multiple slots for additional lift augmentation.}{catAE}
\glsadd{Split-Flap}\dictentry{Split Flap}{A split flap is a trailing-edge device consisting of a hinged panel on the lower surface of the wing that deflects downward while the upper surface remains fixed. When deployed, it creates a large region of separated low-pressure flow behind it, increasing both lift and drag significantly. The high drag characteristic is useful during approach and landing for steep descent angles without increasing airspeed. Split flaps are mechanically simpler than slotted or Fowler types but less efficient at producing lift.}{catAE}
\glsadd{Slotted-Flap}\dictentry{Slotted Flap}{A slotted flap is a trailing-edge device that, when deflected, creates a gap between the flap and the wing's trailing edge. High-energy air from the lower surface flows through this slot and re-energizes the boundary layer on the flap's upper surface, delaying separation and allowing greater deflection angles. This slot effect produces substantially more lift than a plain flap of comparable size. Slotted flaps are widely used on general aviation and commercial aircraft and may incorporate multiple slots.}{catAE}
\glsadd{Krueger-Flap}\dictentry{Krueger Flap}{A Krueger flap is a leading-edge high-lift device consisting of a hinged panel that folds out from the lower surface of the wing's leading edge to increase camber. Unlike slats, Krueger flaps do not create a slot gap, but their deployment changes the leading-edge radius and curvature to improve airflow attachment at high angles of attack. They are most commonly found on inboard sections of transport aircraft wings where mechanical simplicity is advantageous.}{catAE}
\glsadd{Handley-Page-Slat}\dictentry{Handley Page Slat}{A Handley Page slat is a leading-edge device that extends forward from the wing's leading edge, creating a slot gap through which high-pressure air from below flows to the upper surface. This energized airflow delays boundary-layer separation, allowing higher angles of attack before stalling. The slat may extend automatically under aerodynamic loads or be actuated mechanically. Handley Page slats have been a fundamental high-lift technology since the 1920s and remain widely used on transport and military aircraft.}{catAE}
\glsadd{Blown-Flap}\dictentry{Blown Flap}{A blown flap is a high-lift device that uses compressed air, typically bled from engines, blown over the flap surfaces to control the boundary layer and delay flow separation. By energizing the airflow, blown flaps allow much greater flap deflection angles without stall, producing significantly more lift than passive designs. The increased lift reduces takeoff and landing distances, enabling operation from shorter runways. Blown flap systems were used on several STOL aircraft.}{catAE}
\glsadd{Jet-Flap}\dictentry{Jet Flap}{A jet flap is a high-lift system that directs high-velocity engine exhaust over the trailing-edge flaps, dramatically increasing lift through both increased circulation and direct momentum transfer. The jet blast energizes the boundary layer over the flap surfaces, preventing separation at extreme deflection angles. Jet flaps can nearly double maximum lift coefficient compared to conventional flaps but impose significant weight and complexity penalties. The concept has been explored in several experimental STOL aircraft.}{catAE}
\glsadd{Circulation-Control-Wing}\dictentry{Circulation Control Wing}{A circulation control wing employs continuous blowing of compressed air over a rounded trailing edge to control circulation and lift. The Coanda effect causes the blown jet to follow the curved trailing edge, shifting the rear stagnation point and modifying lift without mechanical moving parts. By varying the blowing rate, lift can be adjusted over a wide range. This technology offers simpler, lighter control systems with faster response than conventional hinged surfaces.}{catAE}
\glsadd{Active-Flow-Control}\dictentry{Active Flow Control}{Active flow control encompasses techniques using energy input from actuators to manipulate airflow over aircraft surfaces in real time. Methods include synthetic jet actuators, pulsed blowing, plasma actuators, and micro-electro-mechanical systems that alter boundary layer characteristics. The goal is to delay separation, reduce drag, or increase lift without traditional mechanical devices. Active flow control is an active research area promising more efficient aircraft with fewer moving parts.}{catAE}
\glsadd{Load-Alleviation}\dictentry{Load Alleviation}{Load alleviation is an active flight control technique that uses sensors and control computers to detect and counteract excessive structural loads in real time. When a load event such as a gust or hard maneuver is detected, the system commands small control surface deflections to reduce the incremental load before it reaches the structure. This allows lighter airframe design because it need not withstand every possible load condition. Load alleviation systems are standard on modern transport and military aircraft.}{catAE}
\glsadd{Gust-Alleviation}\dictentry{Gust Alleviation}{Gust alleviation targets structural loads and passenger discomfort from atmospheric turbulence. Accelerometers and air data sensors detect gust onset, and the flight control system commands rapid control surface deflections counteracting gust-induced lift before the entire aircraft responds. The result is a smoother ride and reduced fatigue loads on the airframe. Gust alleviation is increasingly important for composite structures, which are more sensitive to repeated load cycles.}{catAE}
\glsadd{Mode-Control-Panel}\dictentry{Mode Control Panel}{The mode control panel is the primary cockpit interface through which the flight crew selects and manages the autopilot and autothrottle operating modes. Located on the glare shield for easy access, it provides buttons, knobs, and display windows for selecting heading, altitude, vertical speed, airspeed, and navigation modes. The MCP is the central element of the automated flight control system, allowing transition between manual and automated flight with minimal workload. Its design and layout are standardized across aircraft types.}{catAE}
\glsadd{Navigation-Display}\dictentry{Navigation Display}{The navigation display is a cockpit instrument presenting the aircraft's position, planned route, nearby waypoints, navaids, and traffic information in a consolidated graphical format. It typically shows a moving map or arc view with the planned track overlaid, along with data blocks for distance, ground speed, wind, and estimated time of arrival. The navigation display is a primary reference for route monitoring and situational awareness. Modern versions integrate GPS, inertial reference, radio navigation, and traffic surveillance systems.}{catAE}
\glsadd{Primary-Flight-Display}\dictentry{Primary Flight Display}{The primary flight display is the main instrument in a modern glass cockpit, presenting airspeed, altitude, attitude, vertical speed, heading, and flight director command bars on a single electronic screen. It replaces traditional mechanical instruments with a customizable digital display. The PFD enhances situational awareness by integrating data from air data computers, inertial reference units, and radio navigation. It is the pilot's primary reference for controlling flight path and attitude.}{catAE}
\glsadd{Engine-Display}\dictentry{Engine Display}{An engine display is a cockpit screen presenting real-time engine performance parameters including thrust, temperatures, pressures, fuel flow, and vibration levels. On multi-engine aircraft, each engine's data is displayed side by side for easy comparison, allowing quick anomaly detection. Modern engine displays integrate into the electronic instrument system and include caution and advisory messages when parameters exceed normal ranges. They are essential for monitoring engine health during all flight phases.}{catAE}
\glsadd{System-Display}\dictentry{System Display}{A system display is a cockpit screen providing the flight crew with status and diagnostic information about hydraulic, electrical, pneumatic, fuel, and environmental control systems. When a malfunction occurs, the display presents the affected system's schematic with the failed component highlighted, along with corrective action procedures. This integration of monitoring and crew alerting significantly reduces pilot workload during abnormal situations.}{catAE}
\glsadd{Caution-Warning-System}\dictentry{Caution Warning System}{A caution warning system continuously checks aircraft systems for abnormal conditions and alerts the crew through visual and auditory indications. Caution messages indicate conditions requiring awareness but not immediate action, while warnings indicate conditions requiring immediate corrective action. The system monitors dozens of parameters simultaneously and prioritizes alerts to prevent information overload during cascading failures.}{catAE}
\glsadd{Master-Warning}\dictentry{Master Warning}{The master warning is a prominent visual and auditory alert in the cockpit indicating a condition requiring immediate crew action. It is typically a large illuminated button on the instrument panel that activates when any critical system parameter exceeds safe limits. The master warning draws the pilot's attention to the specific system alert displayed on monitoring screens. It serves as the top-level alert in the hierarchical crew alerting architecture used in modern transport aircraft.}{catAE}
\glsadd{EICAS}\dictentry{EICAS}{Engine Indicating and Crew Alerting System is a computerized monitoring system on Boeing aircraft integrating engine performance data display with caution and warning functions. EICAS presents engine parameters such as N1 speed, EGT, fuel flow, and oil pressure while monitoring all aircraft systems for abnormal conditions. When a malfunction is detected, the system displays the appropriate caution or warning message along with the affected system's status. EICAS reduces crew workload by consolidating multiple monitoring functions.}{catAE}
\glsadd{ECAM}\dictentry{ECAM}{Electronic Centralized Aircraft Monitor is an Airbus flight deck system integrating engine and system monitoring with automated crew alerting and electronic checklists. When a failure occurs, ECAM automatically presents the affected system's schematic, identifies the failed component, and displays recommended corrective actions in priority order. The system replaces the need for paper checklists during emergencies, reducing response time and procedural errors. ECAM is a core element of the Airbus philosophy of flight deck automation.}{catAE}
\glsadd{Data-Bus}\dictentry{Data Bus}{A data bus is a shared communication pathway that allows avionics components and Line Replaceable Units to exchange digital information using standardized protocols. By using a common bus architecture, aircraft reduce individual wiring connections, saving weight and complexity. The data bus must ensure high reliability, real-time performance, and resistance to electromagnetic interference. Modern aircraft may use several bus types optimized for different data rates and criticality levels.}{catAE}
\glsadd{ARINC-429}\dictentry{ARINC 429}{ARINC 429 is a widely adopted data communication standard in commercial aviation defining electrical characteristics, data format, and protocol for serial digital information transfer between avionics components. It uses a twisted-pair, simplex transmission scheme where each transmitter sends data on its own bus and receivers listen passively. The standard specifies a 32-bit word format with labels, source and destination identifiers, and parity checking. ARINC 429 has been the backbone of commercial avionics data communication for decades.}{catAE}
\glsadd{MIL-STD-1553}\dictentry{MIL-STD-1553}{MIL-STD-1553 is a U.S. military standard defining mechanical, electrical, and functional characteristics of a serial data communication bus for military and aerospace systems. It uses a dual-redundant twisted-pair bus with time-division multiplexing, where a bus controller manages all data transfers and remote terminals respond to commands. The standard supports high reliability through built-in error detection, redundant buses, and deterministic timing. It is the primary data bus standard for most Western military aircraft.}{catAE}
\glsadd{Aircraft-Wiring}\dictentry{Aircraft Wiring}{Aircraft wiring refers to the network of electrical cables and wire harnesses distributing power, carrying control signals, and transmitting data throughout an aircraft. Every wire must meet stringent standards for insulation, conductor material, shielding, and flame resistance to operate reliably in the harsh aerospace environment of vibration, temperature extremes, and low pressure. Wire routing must avoid electromagnetic interference sources and maintain separation from hydraulic and fuel lines.}{catAE}
\glsadd{Circuit-Breaker}\dictentry{Circuit Breaker}{A circuit breaker is a resettable overcurrent protection device that automatically opens the electrical circuit when current exceeds a predetermined safe level, preventing damage from overload or short circuits. Unlike a fuse, a circuit breaker can be reset after the fault condition is cleared. In aircraft, circuit breakers are organized on dedicated panels accessible to the flight crew, allowing isolation of faulty circuits in flight. Proper rating and maintenance of circuit breakers is essential for preventing electrical fires.}{catAE}
\glsadd{Bus-Bar}\dictentry{Bus Bar}{A bus bar is a low-impedance electrical conductor, typically a flat strip of copper or aluminum, used to distribute power from a source to multiple downstream circuits. In aircraft, bus bars are housed in power distribution panels and serve as the central connection point for generators, batteries, and circuit breakers. Multiple bus bars may separate essential and non-essential loads, ensuring critical systems remain powered even if a main bus fails.}{catAE}
\glsadd{Generator}\dictentry{Generator}{An aircraft generator converts mechanical power from the engine or auxiliary power unit into electrical power for the aircraft's systems. Aircraft generators typically produce alternating current at a regulated voltage and frequency, with the engine-driven generator being the primary source during flight. Redundant generators on multi-engine aircraft ensure loss of one engine does not result in loss of electrical power. Modern brushless generators with integrated voltage regulators offer high reliability and low maintenance.}{catAE}
\glsadd{Alternator}\dictentry{Alternator}{An alternator is an alternating current generator commonly used in light aircraft and as a component of the aircraft electrical system. It converts engine rotation into AC electrical power, which is rectified to DC if needed for the aircraft's electrical bus. Alternators are preferred over DC generators because they are lighter, more efficient, and produce a smoother output with less electrical noise. In modern aircraft, the terms generator and alternator are often used interchangeably.}{catAE}
\glsadd{Battery}\dictentry{Battery}{A battery stores electrical energy as chemical energy and releases it as direct current when connected to a load. Spacecraft batteries charge during sunlit orbit portions and discharge during eclipses when solar panels are in shadow. Lithium-ion batteries are now standard due to high energy density and long cycle life. Battery management systems control charge and discharge rates to maximize lifespan through extreme temperature swings.}{catAE}
\glsadd{APU}\dictentry{APU}{The Auxiliary Power Unit is a small gas turbine engine, typically in the aircraft's tail section, that provides electrical power and compressed air for ground operations without running the main engines. The APU drives a generator for electrical power and supplies bleed air for air conditioning, cabin pressurization, and engine start. It allows the aircraft to be self-sufficient at airports without ground power or air start carts. The APU is normally started before pushback and shut down after the main engines are running.}{catAE}
\glsadd{Inverter}\dictentry{Inverter}{An inverter converts direct current from the aircraft's DC electrical system or battery into alternating current at the required voltage and frequency. Aircraft AC-powered instruments and motors require the stable AC supply that the inverter provides. Static inverters using semiconductor switching elements have replaced older motor-generator types due to higher reliability, lower weight, and no moving parts. They are critical in aircraft with DC primary systems that must also support AC loads.}{catAE}
\glsadd{Transformer}\dictentry{Transformer}{A transformer is a passive electromagnetic device that changes the voltage level of an alternating current power supply while maintaining the same frequency. In aircraft electrical systems, transformers step voltages up or down as required by different loads. They are also used for electrical isolation between circuits, preventing fault current propagation. Transformer rectifier units combine a transformer with a rectifier to convert AC to DC at a different voltage level, serving as key components in aircraft power distribution.}{catAE}
\glsadd{Rectifier}\dictentry{Rectifier}{A rectifier is an electrical device that converts alternating current into direct current, allowing AC-powered sources to supply DC loads in the aircraft electrical system. Semiconductor diodes or thyristors are used as the rectifying elements, offering high efficiency and rapid response. Rectifier units in aircraft often include transformers and filters to produce a clean, regulated DC output from the AC generator supply. They are essential for charging batteries, powering DC avionics, and supplying the DC bus from AC-producing generators.}{catAE}
\glsadd{Voltage-Regulator}\dictentry{Voltage Regulator}{A voltage regulator automatically maintains a constant output voltage despite variations in input voltage, load current, or temperature. In aircraft systems, it ensures all equipment receives power within specified voltage tolerance, preventing damage from overvoltage or malfunction from undervoltage. Modern solid-state regulators are integrated into the generator control unit and respond in milliseconds. Reliable voltage regulation is essential for sensitive avionics and electrical equipment longevity.}{catAE}
\glsadd{Oxygen-System}\dictentry{Oxygen System}{An aircraft oxygen system provides supplemental breathable oxygen to flight crew and passengers at altitudes where ambient air pressure is too low to sustain consciousness. Systems range from simple chemical oxygen generators for passengers to high-pressure gaseous systems for crew members requiring reliable supply for extended durations. The system must deliver oxygen at the correct flow rate and pressure for the cabin altitude. Oxygen systems are critical safety equipment for high-altitude flight and decompression emergencies.}{catAE}
\glsadd{Pressurization-System}\dictentry{Pressurization System}{An aircraft pressurization system maintains a safe cabin altitude by controlling the pressure difference between the interior and outside atmosphere during flight. Compressed air from engine bleed is regulated and supplied to the cabin, while outflow valves control air discharge to maintain desired pressure. The system must handle conditions from ground level to maximum altitude while providing automatic and manual control modes. Pressurization failure at high altitude is a serious emergency requiring immediate descent.}{catAE}
\glsadd{Air-Conditioning}\dictentry{Air Conditioning}{The aircraft air conditioning system controls temperature, humidity, and cleanliness of air supplied to the cabin and flight deck. Air from engine bleed or electric compressors is cooled by air cycle machines or vapor cycle systems and distributed through environmental control system ducting. The system must maintain comfortable conditions despite extreme outside temperatures ranging from minus 60 degrees Celsius at cruise to plus 50 degrees on hot desert runways. Air conditioning directly affects passenger comfort and crew performance.}{catAE}
\glsadd{Bleed-Air}\dictentry{Bleed Air}{Bleed air is compressed air extracted from the compressor stage of a gas turbine engine before the combustion chamber. This hot, high-pressure air powers numerous aircraft systems including cabin pressurization, air conditioning, wing and engine anti-icing, and pneumatic engine starting. The bleed air must be cooled and regulated before use due to temperatures that can exceed 200 degrees Celsius. Some modern aircraft designs are moving toward electric compressor systems to eliminate the efficiency penalty of bleed air extraction.}{catAE}
\glsadd{ECS-Pack}\dictentry{ECS Pack}{An Environmental Control System pack is a self-contained air cycle machine that conditions bleed air from the engines to the temperature and pressure required for the cabin. Each pack typically includes a compressor, turbine, and heat exchanger that cool and regulate the air supply. Most transport aircraft are equipped with two or three packs for redundancy and sufficient capacity. The ECS pack is the core component of the environmental control system, directly affecting passenger comfort and system reliability.}{catAE}
\glsadd{Ventilation}\dictentry{Ventilation}{Cabin ventilation refers to the continuous circulation and replacement of air within the aircraft cabin to maintain air quality, prevent buildup of carbon dioxide and contaminants, and ensure passenger comfort. Fresh air enters from the environmental control system while stale air is exhausted overboard through outflow valves. The system must provide adequate air changes per hour while distributing air evenly to avoid hot or cold spots. Proper ventilation is essential for preventing the spread of airborne illnesses during long flights.}{catAE}
\glsadd{Rotor-Brake}\dictentry{Rotor Brake}{A rotor brake is a mechanical braking device installed on the main rotor mast to stop rotor rotation after engine shutdown. Without it, aerodynamic drag alone may take several minutes to stop the rotor, which is unacceptable in emergencies. The rotor brake uses hydraulic pressure or mechanical force to apply friction to the rotor shaft, bringing blades to a controlled stop. It is also used as an anti-rotation device to prevent windmilling in high winds while parked.}{catAE}
\glsadd{Rotor-Head}\dictentry{Rotor Head}{The rotor head is the mechanical assembly at the top of the main rotor mast connecting blades to the helicopter's drive system and providing hinge mechanisms for blade movement. It must transmit enormous centrifugal and aerodynamic forces from the blades to the airframe while allowing complex motions required for flight. The rotor head design determines handling characteristics, vibration levels, and maintenance requirements. It is one of the most highly stressed and critical components in any helicopter.}{catAE}
\glsadd{Swashplate}\dictentry{Swashplate}{A swashplate is a mechanical device of two annular plates separated by bearings that translates non-rotating control inputs into rotating blade pitch changes on the main rotor. The lower non-rotating plate connects to control rods from the cockpit, while the upper rotating plate turns with the mast and connects to blade pitch change links. By tilting the swashplate, the pilot changes blade pitch cyclically around the disk, directing the thrust vector for directional control.}{catAE}
\glsadd{Collective-Pitch}\dictentry{Collective Pitch}{Collective pitch is a helicopter flight control that simultaneously changes the angle of attack of all main rotor blades by the same amount. Increasing collective pitch raises total rotor thrust, causing climb, while decreasing it reduces thrust for descent. The collective lever on the left side of the pilot's seat connects to the swashplate to raise or lower the entire plate without tilting it. Managing collective pitch in coordination with engine power is one of the most fundamental helicopter piloting skills.}{catAE}
\glsadd{Cyclic-Pitch}\dictentry{Cyclic Pitch}{Cyclic pitch is a helicopter flight control that varies each blade's pitch angle as it rotates around the disk, creating periodic lift variation that tilts the rotor thrust vector in the desired direction. The pilot's cyclic stick tilts the swashplate, commanding higher pitch on one side and lower on the opposite. This asymmetry causes the helicopter to move in the direction the disk is tilted. Cyclic control is the primary method for directing horizontal flight path and attitude.}{catAE}
\glsadd{Lead-Lag-Hinge}\dictentry{Lead-Lag Hinge}{A lead-lag hinge, also called a drag hinge, is a hinge on a helicopter rotor blade allowing the blade to pivot forward and aft in the plane of rotation. This accommodates the Coriolis effect, which causes the blade's rotational speed to change as it flaps up and down during each revolution. Without it, resulting internal stresses would be transmitted to the rotor hub and could cause structural failure. Dampers are typically installed at lead-lag hinges to prevent ground resonance.}{catAE}
\glsadd{Flapping-Hinge}\dictentry{Flapping Hinge}{A flapping hinge is a hinge on a helicopter rotor blade allowing upward and downward movement relative to the rotor hub. This freedom is essential because it allows advancing and retreating blades to produce equal average lift by flapping to different heights during each revolution. Without it, the helicopter would experience a strong rolling moment due to the airspeed difference between advancing and retreating sides of the disk. The flapping hinge is a key feature distinguishing articulated rotors from rigid designs.}{catAE}
\glsadd{Feathering-Hinge}\dictentry{Feathering Hinge}{A feathering hinge allows a helicopter rotor blade to rotate about its spanwise axis, changing the blade's angle of attack. Controlled by collective and cyclic pitch through the swashplate mechanism, it provides the primary means for controlling lift and direction. The feathering hinge allows smooth pitch changes while transmitting enormous centrifugal forces from spinning blades to the rotor hub. All helicopter rotor systems incorporate feathering, whether through a physical hinge or flexible blade root sections.}{catAE}
\glsadd{Fully-Articulated-Rotor}\dictentry{Fully Articulated Rotor}{A fully articulated rotor attaches each blade to the hub by three separate hinges: a flapping hinge for vertical motion, a lead-lag hinge for in-plane motion, and a feathering hinge for pitch changes. This maximum individual blade freedom reduces stresses transmitted to the hub and allows independent blade response to aerodynamic and inertial forces. Fully articulated rotors are found on medium and large transport helicopters where their complexity is justified by superior vibration and load management.}{catAE}
\glsadd{Semi-Rigid-Rotor}\dictentry{Semi-Rigid Rotor}{A semi-rigid rotor is a helicopter rotor system where blades are rigidly attached to each other but connected to the mast through a teetering hinge allowing the entire disk to tilt as a unit. The two-bladed teetering design allows one blade to rise while the other descends, accommodating dissymmetry of lift without individual flapping hinges. This simplification reduces moving parts and maintenance compared to fully articulated systems. Semi-rigid rotors are characteristic of many light helicopters, including the iconic Bell series.}{catAE}
\glsadd{Rigid-Rotor}\dictentry{Rigid Rotor}{A rigid rotor attaches blades rigidly to the hub with no mechanical hinges, relying on controlled flexibility of blade roots and hub structure to accommodate necessary blade motions. This eliminates complex hinge mechanisms, bearings, and dampers, resulting in lower weight, reduced maintenance, and improved responsiveness. Rigid rotors require sophisticated blade design and materials to withstand stresses otherwise absorbed by hinges. They are used on some modern light helicopters and military tiltrotor aircraft.}{catAE}
\glsadd{Tail-Rotor}\dictentry{Tail Rotor}{The tail rotor is a small rotor mounted vertically on the tail boom of a conventional helicopter that produces a horizontal thrust to counteract the torque reaction of the main rotor. Without the tail rotor, the helicopter fuselage would spin in the opposite direction of the main rotor due to Newton's third law. The tail rotor also provides yaw control, as the pilot adjusts its thrust to turn the helicopter's nose left or right. Tail rotors consume a significant portion of engine power and add noise and complexity to the helicopter design.}{catAE}
\glsadd{Fenestron}\dictentry{Fenestron}{A Fenestron is a ducted tail rotor system consisting of multiple small blades mounted within a circular shroud in the helicopter's tail fin. The duct protects blades from foreign object damage and significantly reduces noise compared to an open tail rotor. The higher blade count running at lower tip speed, combined with duct acoustic treatment, makes the Fenestron substantially quieter, advantageous for operations over populated areas. Developed by Sud Aviation, it is used extensively on Airbus Helicopters products.}{catAE}
\glsadd{NOTAR}\dictentry{NOTAR}{NOTAR is an acronym for No Tail Rotor, an anti-torque system replacing the conventional tail rotor with alternative methods of counteracting main rotor torque. The system uses direct jet thrusters at the tail boom end and a circulation control system blowing air over the tail boom surface to generate side force through the Coanda effect. NOTAR eliminates noise, vibration, and safety hazards associated with conventional tail rotors. Developed by McDonnell Douglas, it is used on the MD 520N and MD 902.}{catAE}
\glsadd{Main-Rotor}\dictentry{Main Rotor}{The main rotor is the large, horizontally rotating wing system of a helicopter providing both lift and thrust for directional control. Each blade acts as a spinning airfoil, generating forces the pilot modulates through collective and cyclic pitch changes. The main rotor operates at relatively low rotational speeds, with tip speeds typically 200 to 220 meters per second. Its design is the most critical factor in determining a helicopter's performance, handling qualities, and vibration characteristics.}{catAE}
\glsadd{Tip-Speed}\dictentry{Tip Speed}{Rotor tip speed is the linear velocity of the outermost point of a rotor blade as it rotates, typically expressed in meters per second or knots. It is determined by the product of rotational speed and blade radius, and is a critical design parameter because it approaches the speed of sound at the tips. If tip speed becomes too high, compressibility effects at the advancing tip and retreating blade stall limit forward speed. Designers must balance tip speed to optimize hover efficiency while avoiding aerodynamic penalties.}{catAE}
\glsadd{Hover}\dictentry{Hover}{Hover is a flight condition where the helicopter maintains a fixed position relative to the ground with zero velocity. In hover, the main rotor must generate thrust equal to the aircraft's weight while the anti-torque system counters rotor torque. Hovering requires the most engine power of any helicopter flight regime because there is no translational lift. The ability to hover is the defining characteristic of rotary-wing aircraft, enabling vertical takeoff, precision placement, and rescue from confined areas.}{catAE}
\glsadd{Autorotation-Landing}\dictentry{Autorotation Landing}{Autorotation landing is a technique for safely landing a helicopter after engine failure by using upward airflow through the rotor disk to maintain rotation and control. During descent, the pilot adjusts collective pitch to keep rotor speed constant, storing kinetic energy in the spinning blades. Just before touchdown, the pilot flares by pulling collective pitch up, converting rotor energy into a momentary reduction of descent rate for a gentle landing. Autorotation training is mandatory in all helicopter pilot certification programs.}{catAE}
\glsadd{Ground-Resonance}\dictentry{Ground Resonance}{Ground resonance is a destructive self-excited oscillation that can occur when a helicopter is on the ground or in very low hover, involving coupling between rotor out-of-track blade motion and airframe natural vibration modes. When blades are not tracking properly, the resulting imbalance creates periodic forces that excite the airframe, which feeds energy back into the rotor. If insufficiently damped, oscillations grow rapidly and can destroy the helicopter in seconds. Lead-lag dampers on articulated rotors are the primary defense.}{catAE}
\glsadd{Retreating-Blade-Stall}\dictentry{Retreating Blade Stall}{Retreating blade stall is an aerodynamic phenomenon limiting helicopter forward speed, occurring when the blade on the retreating side of the disk reaches its critical angle of attack and stalls. As forward speed increases, the retreating blade experiences progressively lower relative airflow, requiring higher pitch angles until it exceeds the stall angle. Onset causes vibration, a pitching or rolling tendency, and loss of lift on the affected side. Pilots must recognize symptoms and reduce airspeed or collective pitch to recover.}{catAE}
\glsadd{Dissymmetry-of-Lift}\dictentry{Dissymmetry of Lift}{Dissymmetry of lift is the fundamental condition in forward helicopter flight where the advancing blade side produces more lift than the retreating side due to the difference in relative airflow velocity. On the advancing side, blade rotational velocity adds to forward speed, while on the retreating side it subtracts. This imbalance must be compensated through blade flapping, cyclic pitch changes, or rotor system design. It is a primary factor limiting helicopter forward speed and shaping rotor design.}{catAE}
\glsadd{Translational-Lift}\dictentry{Translational Lift}{Translational lift is the increase in rotor efficiency and lift as a helicopter transitions from hover to forward flight. Forward motion causes the rotor to operate in relatively undisturbed air rather than its own recirculated downwash, increasing effective angle of attack and reducing induced velocity. This additional lift allows the helicopter to climb or maintain altitude with less power than required in hover. Translational lift becomes significant above approximately 15 to 20 knots and continues to increase with forward speed.}{catAE}
\glsadd{Effective-Translational-Lift}\dictentry{Effective Translational Lift}{Effective translational lift is the component of translational lift that develops as the helicopter gains forward speed and the rotor begins to fly into undisturbed air rather than its own downwash. The transition is most pronounced between 15 and 30 knots, where the pilot feels a distinct improvement in performance. Unlike induced translational lift, this specifically refers to the benefit gained as the rotor disk moves clear of the vortex ring and recirculating flow patterns of hover.}{catAE}
\glsadd{Ground-Effect-Helicopter}\dictentry{Ground Effect Helicopter}{Ground effect in helicopters refers to increased lift and reduced power requirement when the rotor operates within approximately one rotor diameter of the ground. The ground interferes with the rotor's downwash pattern, reducing induced velocity and increasing net thrust for a given power input. This allows helicopters to hover at lower power settings near the ground than at higher altitudes. Ground effect is most beneficial during hover taxi, takeoff, and landing, diminishing rapidly above approximately one rotor diameter.}{catAE}
\glsadd{Vortex-Ring-State}\dictentry{Vortex Ring State}{Vortex ring state is a dangerous condition occurring when a helicopter descends into its own rotor downwash, causing airflow to reverse over portions of the disk and form a toroidal vortex ring. This results in rapid, uncommanded descent with high vibration, loss of control effectiveness, and increased power requirements. It typically develops during steep, slow descents with low airspeed and power applied. Recovery requires nose-down cyclic to gain forward airspeed and reduction of collective pitch.}{catAE}
\glsadd{Settling-With-Power}\dictentry{Settling With Power}{Settling with power is a condition where a helicopter experiences continued descent that cannot be arrested by applying additional collective pitch, because increased power actually worsens the condition. It describes the same phenomenon as vortex ring state, emphasizing the pilot's perception that adding power makes the descent worse. It typically occurs in high gross weight, high density altitude situations during steep approaches. Recovery requires reducing collective pitch and applying forward cyclic.}{catAE}
\glsadd{Power-Required-Curve}\dictentry{Power Required Curve}{The power required curve is a graphical representation of total power needed to maintain level flight at each airspeed, plotted with airspeed on the horizontal axis and power on the vertical axis. The curve is typically U-shaped, with high power at hover due to induced losses, a minimum at intermediate cruise speed, and increasing power at high speed due to profile and parasite drag. Understanding this curve is essential for helicopter performance planning, determining range speed, endurance speed, and maximum speed capability.}{catAE}
\glsadd{Power-Available-Curve}\dictentry{Power Available Curve}{The power available curve shows the maximum power the engine can deliver to the rotor at each airspeed, accounting for installation losses and corrections for altitude and temperature. Plotted on the same axes as the power required curve, the intersection determines maximum level flight speed. If the power available falls below power required at any speed, the helicopter cannot sustain flight there. The gap between curves represents excess power available for climbing or accelerating.}{catAE}
\glsadd{Height-Velocity-Diagram}\dictentry{Height-Velocity Diagram}{The height-velocity diagram is a helicopter flight envelope chart defining combinations of altitude and airspeed from which a safe autorotation landing can be made following engine failure. It shows two shaded regions, a low-speed high-altitude zone and a low-altitude high-speed zone, between which autorotation is possible. Flying within the shaded regions, the dead man's curve, may prevent successful autorotation. Pilots study this diagram and plan profiles to avoid critical regions.}{catAE}
\glsadd{H-V-Curve}\dictentry{H-V Curve}{The H-V curve is a common abbreviation for the height-velocity diagram, a fundamental performance chart every helicopter pilot must understand. It defines the safe envelope for engine-out autorotation by showing the minimum altitude required for successful autorotation at each airspeed, and conversely the minimum speed needed at each altitude. The curve is unique to each helicopter type and varies with gross weight and density altitude. Pilots use it to plan approaches and departures keeping the helicopter within recoverable conditions.}{catAE}
\glsadd{Cat-A-Takeoff}\dictentry{Cat A Takeoff}{A Category A takeoff is a helicopter takeoff profile ensuring the helicopter can safely continue the takeoff or land safely if an engine fails at any point during departure. The profile requires climbing to a predetermined decision speed and altitude before transitioning to forward flight, with the ability to either continue or execute a controlled descent to a safe area. Cat A procedures are mandatory at heliports with obstacles or airports with complex departure paths. The additional performance margin reduces payload capability.}{catAE}
\glsadd{Cat-B-Takeoff}\dictentry{Cat B Takeoff}{A Category B takeoff assumes the engine will not fail during departure, allowing a simpler, more direct climb to cruise altitude. Unlike Cat A, a Cat B takeoff does not require the ability to safely land or continue following an engine failure at any point. This lower performance requirement allows more payload or operation from higher-elevation heliports. Cat B takeoffs are appropriate in unobstructed areas where the helicopter can safely autorotate to a suitable landing area if necessary.}{catAE}
\glsadd{Stall-Warning-System}\dictentry{Stall Warning System}{A stall warning system monitors the wing's angle of attack and alerts the pilot when approaching the critical angle at which stall will occur. The system typically activates several degrees before the actual stall, providing time for corrective action. On transport aircraft, stall warning is a certified, redundant system essential for safe flight. Modern systems use multiple angle of attack sensors and voting logic to ensure reliable operation even with sensor failures.}{catAE}
\glsadd{Stick-Shaker}\dictentry{Stick Shaker}{A stick shaker is a stall warning device on the pilot's control column that physically vibrates the stick when approaching the critical angle of attack. The vibration is produced by an electric motor with eccentric weights, creating a distinct tactile and auditory alert impossible to ignore. It activates several degrees before the actual stall, giving time to recover. This mechanical warning is required on all transport category aircraft and provides the most direct stall alert available.}{catAE}
\glsadd{Yaw-Damper}\dictentry{Yaw Damper}{A yaw damper is an automatic flight control system that detects and dampens Dutch roll oscillations in swept-wing aircraft. A rate gyro senses yaw velocity and commands small, automatic rudder deflections out of phase with the oscillation to suppress it. On transport aircraft, the yaw damper operates continuously and is essential for passenger comfort and structural fatigue prevention. Some aircraft have primary and standby yaw dampers for redundancy, as loss can result in uncomfortable and potentially dangerous oscillations.}{catAE}
\glsadd{Autothrottle}\dictentry{Autothrottle}{An autothrottle automatically adjusts engine thrust to maintain a target airspeed, Mach number, or flight path angle without manual pilot input. Air data inputs compute required thrust and command fuel control or throttle actuators to achieve the desired setting. Autothrottle reduces pilot workload during cruise and ensures precise speed control during approach and landing. It integrates with the autopilot for coordinated speed and path management, and is standard on modern transport aircraft.}{catAE}
\glsadd{Flight-Director}\dictentry{Flight Director}{A flight director displays command guidance cues telling the pilot how to fly the aircraft to follow a selected flight path or maintain desired speed and altitude. It computes necessary control inputs based on selected autopilot and navigation modes and presents them as bars or symbols on the primary flight display. The pilot manually flies to match the command bars, combining automation guidance with manual control. The flight director is particularly valuable during instrument approaches and complex maneuvers.}{catAE}
\glsadd{TCAS}\dictentry{TCAS}{Traffic Alert and Collision Avoidance System is an airborne system monitoring airspace for transponder-equipped traffic and providing advisories to prevent mid-air collisions. TCAS interrogates nearby transponders to determine range, bearing, and altitude, tracking relative motion to predict conflicts. When collision risk is detected, TCAS issues traffic advisories and, if the threat continues, resolution advisories commanding climb or descent. TCAS is mandated on all large commercial aircraft worldwide.}{catAE}
\glsadd{GPWS}\dictentry{GPWS}{Ground Proximity Warning System is an aircraft safety system using radar altimeter data, barometric altitude, and rate of change to detect dangerous flight conditions near terrain. GPWS provides distinctive aural warnings such as terrain, pull up, and excessive descent rate alerts when the aircraft is in danger of controlled flight into the ground. The system uses predictive algorithms to issue warnings several seconds before a predicted terrain impact. GPWS has prevented numerous accidents and is mandated on commercial IFR aircraft.}{catAE}
\glsadd{Radar-Altimeter}\dictentry{Radar Altimeter}{A radar altimeter measures absolute height above the terrain directly below by transmitting a radio signal downward and timing the reflected return. Unlike a barometric altimeter measuring pressure altitude above sea level, it provides actual ground clearance critical during low-level flight, approach, and landing. The reading is displayed on the primary flight display and feeds the ground proximity warning system and autoland systems. Operating range is from 2,500 feet to ground level.}{catAE}
\glsadd{Transponder}\dictentry{Transponder}{A transponder is a communication device aboard a spacecraft that receives a signal on one frequency and retransmits it on another. It amplifies and conditions the signal to ensure reliable data delivery over vast distances. Transponders are essential for satellite communication, telemetry, and command links. They enable continuous two-way communication between ground stations and orbiting or deep-space vehicles.}{catAE}
\glsadd{Abort}\dictentry{Abort}{A planned or emergency termination of a flight or launch sequence. An abort occurs when a pilot or automated system determines that continuing the mission is unsafe and initiates procedures to land or separate from the vehicle. Aborts can range from immediate emergency ejections to controlled return-to-launch-site scenarios. The ability to safely abort at various stages of flight is a fundamental requirement in aerospace vehicle design.}{catAE}
\glsadd{Aerobraking}\dictentry{Aerobraking}{A spaceflight maneuver that uses atmospheric drag to reduce a spacecraft's orbital velocity. By repeatedly dipping into a planet's upper atmosphere at periapsis, the vehicle gradually lowers its orbit without expending propellant. Aerobraking requires careful thermal management and precise navigation to avoid excessive heating or structural loads. This technique has been used successfully on missions to Mars and Venus to circularize orbits.}{catAE}
\glsadd{Ejection-Seat}\dictentry{Ejection Seat}{An emergency escape system that launches a crew member out of an aircraft using an explosive charge or rocket motor. The system typically fires the canopy, extracts the occupant, and deploys a parachute in a carefully timed sequence. Ejection seats are designed for use at altitudes and speeds within a specific safe envelope. Modern variants can accommodate zero-zero conditions, meaning zero altitude and zero airspeed.}{catAE}
\glsadd{Instrument-Landing-System}\dictentry{Instrument Landing System}{A ground-based radio navigation system that provides precision lateral and vertical guidance to aircraft during approach and landing. The ILS transmits two signals: a localizer for horizontal alignment with the runway and a glideslope for the correct descent path. Pilots follow these signals on cockpit instruments to land safely in low-visibility conditions. ILS remains one of the most widely used precision approach systems worldwide.}{catAE}
\glsadd{JATO}\dictentry{JATO}{Jet-Assisted Takeoff, a method of using small rocket boosters to augment an aircraft's engines during takeoff. JATO bottles are mounted on the fuselage and ignited to provide extra thrust for short field operations or heavy payloads. The rockets burn for only a few seconds but produce enormous additional force. This technique was commonly used on military transports and is now largely replaced by improved engine technology.}{catAE}
\glsadd{Parachute}\dictentry{Parachute}{A fabric device designed to create aerodynamic drag and slow the descent of a body through the atmosphere. The canopy catches air and decelerates the payload or person to a survivable landing speed. Parachutes come in various configurations including round, ram-air, and drogue types for different applications. They are essential safety equipment for aircraft, spacecraft recovery, skydiving, and military operations.}{catAE}
\glsadd{Traffic-Collision-Avoidance-System}\dictentry{Traffic Collision Avoidance System}{An onboard avionics system that monitors surrounding air traffic and provides resolution advisories to pilots to prevent mid-air collisions. TCAS interrogates transponders of nearby aircraft and calculates their relative positions, altitudes, and closing rates. If a potential conflict is detected, the system issues visual and auditory alerts with specific climbing or descending instructions. TCAS is mandated on most commercial aircraft and operates independently of ground-based air traffic control.}{catAE}
\glsadd{Terrain-Awareness-and-Warning-System}\dictentry{Terrain Awareness and Warning System}{An avionics system that alerts pilots when an aircraft is in immediate danger of flying into terrain. TAWS uses a terrain database and GPS position data to compare the aircraft's flight path against surrounding topography. The system provides visual and aural warnings such as terrain ahead and pull up commands. TAWS has significantly reduced controlled-flight-into-terrain accidents since its widespread adoption.}{catAE}
\glsadd{Barometric-Altimeter}\dictentry{Barometric Altimeter}{An instrument that measures altitude above a reference point by comparing atmospheric pressure at the aircraft's position with a known pressure setting. As an aircraft climbs, decreasing ambient pressure is converted to a height reading via an aneroid barometer. Pilots calibrate the altimeter to local barometric pressure for accurate readings during different phases of flight. This instrument is a fundamental component of the basic six-pack of flight instruments.}{catAE}
\glsadd{Capstan-Winch}\dictentry{Capstan Winch}{A mechanical device used on aircraft and ships to provide controlled pulling force through a rotating drum. The operator wraps a rope around the capstan and controls the tension by varying friction. In aerospace applications, capstan winches are used for tasks such as tensioning cables and operating cargo handling systems. The design allows for fine control and the ability to hold loads without additional locking mechanisms.}{catAE}
\glsadd{Cockpit-Voice-Recorder}\dictentry{Cockpit Voice Recorder}{An onboard device that records audio in the flight deck, typically capturing pilot conversations, radio transmissions, and ambient sounds. The CVR stores the last two hours of audio on a solid-state memory unit encased in a crash-survivable housing. It is a critical tool for accident investigation, providing investigators with context leading up to an incident. CVRs are required on virtually all commercial and many general aviation aircraft.}{catAE}
\glsadd{Directional-Gyroscope}\dictentry{Directional Gyroscope}{A gyroscopic instrument that provides the pilot with a stable heading reference, replacing the traditional magnetic compass in most modern cockpits. The device uses a spinning gyroscope to maintain its orientation regardless of aircraft movements. It must be periodically reset to match the magnetic compass due to precession drift. The directional gyro offers more stable and accurate readings than a standard compass, especially during turns.}{catAE}
\glsadd{Electronic-Flight-Instrument-System}\dictentry{Electronic Flight Instrument System}{An integrated suite of digital displays that present flight data, engine parameters, and navigation information on electronic screens in the cockpit. EFIS replaced traditional analog instruments with compact cathode ray tube or liquid crystal displays. The system typically includes primary flight displays and multi-function displays that pilots can customize. EFIS is now standard on most modern commercial and military aircraft.}{catAE}
\glsadd{Flight-Data-Recorder}\dictentry{Flight Data Recorder}{An onboard device that continuously records aircraft flight parameters such as airspeed, altitude, heading, and control surface positions. The FDR stores data on crash-protected solid-state memory, typically retaining at least 25 hours of flight information. Combined with the cockpit voice recorder, it provides investigators with a comprehensive picture of an aircraft's performance before an incident. FDRs are mandatory on commercial airliners worldwide.}{catAE}
\glsadd{Flight-Envelope}\dictentry{Flight Envelope}{The complete range of conditions under which an aircraft is designed and certified to operate safely. This includes parameters such as airspeed, altitude, load factor, and Mach number. Exceeding the flight envelope can lead to structural failure, loss of control, or other dangerous conditions. Pilots must understand and respect these limits, which are typically depicted on the aircraft's airspeed indicator.}{catAE}
\glsadd{Flight-Level}\dictentry{Flight Level}{A standard altitude measurement used in aviation, expressed in hundreds of feet based on a standard pressure setting of 29.92 inches of mercury. Flight Level 350, for example, corresponds to 35,000 feet under standard atmospheric conditions. Using a common pressure reference ensures vertical separation between aircraft regardless of local barometric variations. Flight levels are assigned by air traffic control above the transition altitude.}{catAE}
\glsadd{Gyroscopic-Precession}\dictentry{Gyroscopic Precession}{A phenomenon in which a force applied to a spinning gyroscope produces a rotational effect 90 degrees later in the direction of rotation. When a torque is applied about one axis, the gyroscope tilts about an axis perpendicular to both the applied force and the spin axis. This property is fundamental to the operation of attitude indicators, heading gyros, and turn coordinators. Understanding precession is essential for accurate interpretation of gyroscopic flight instruments.}{catAE}
\glsadd{Heading-Indicator}\dictentry{Heading Indicator}{A gyroscopic flight instrument that displays the aircraft's magnetic heading on a rotary compass card. It is more stable and easier to read than a magnetic compass, especially during turns and acceleration. The heading indicator must be periodically realigned with the magnetic compass to correct for gyroscopic drift. It is one of the primary instruments used for navigation and maintaining course in the cockpit.}{catAE}
\glsadd{Horizontal-Situation-Indicator}\dictentry{Horizontal Situation Indicator}{A cockpit instrument that combines heading information with navigation guidance on a single display. The HSI shows the aircraft's heading, bearing to a selected navigational aid, and deviation from a course or approach path. It integrates the functions of a heading indicator and a VOR indicator into one intuitive instrument. The HSI greatly reduces pilot workload during navigation and instrument approaches.}{catAE}
\glsadd{Machmeter}\dictentry{Machmeter}{An instrument that displays the aircraft's speed as a ratio of true airspeed to the local speed of sound, known as the Mach number. The machmeter accounts for variations in temperature and pressure with altitude to provide an accurate reading. Operating near or above Mach 1 introduces compressibility effects that affect aircraft performance and control. The machmeter is essential for high-altitude and high-speed flight operations.}{catAE}
\glsadd{Maneuvering-Speed}\dictentry{Maneuvering Speed}{The maximum airspeed at which the aircraft can safely perform full control deflections without exceeding structural design limits. At this speed, the aircraft will stall before structural components are overloaded. Maneuvering speed varies with aircraft weight and is lower than the maximum structural cruising speed. Pilots reduce speed to this value when encountering turbulence or performing aggressive maneuvers.}{catAE}
\glsadd{Maximum-Certified-Altitude}\dictentry{Maximum Certified Altitude}{The highest altitude at which an aircraft has been tested and approved for operation by the certifying authority. This limit is determined by factors such as pressurization system capability, oxygen supply, and engine performance. Operating above this altitude could compromise structural integrity or crew and passenger safety. The certified altitude is published in the aircraft's operating limitations.}{catAE}
\glsadd{Maximum-Landing-Weight}\dictentry{Maximum Landing Weight}{The greatest weight at which an aircraft is certified to land safely. This limit is typically lower than the maximum takeoff weight because the airframe must absorb landing loads without damage. Exceeding this weight can result in structural failure or an inability to stop within available runway length. Pilots plan fuel loads and cargo to ensure the aircraft is at or below this weight upon arrival.}{catAE}
\glsadd{Maximum-Takeoff-Weight}\dictentry{Maximum Takeoff Weight}{The heaviest weight at which an aircraft is permitted to begin its takeoff roll. This limit considers structural strength, engine thrust, wing lift capability, and braking performance. Exceeding the maximum takeoff weight reduces climb performance and increases required runway length. Airlines calculate this weight for each flight based on payload, fuel requirements, and environmental conditions.}{catAE}
\glsadd{Minimum-Control-Speed}\dictentry{Minimum Control Speed}{The lowest airspeed at which an aircraft can maintain controlled flight following an engine failure or other asymmetric power condition. Operating below this speed after losing an engine may result in loss of directional control. Minimum control speed varies with aircraft configuration, weight, and altitude. Pilots must know this value for all critical phases of flight, especially during takeoff.}{catAE}
\glsadd{Operating-Limit-Speed}\dictentry{Operating Limit Speed}{A set of maximum and minimum airspeeds defined by the aircraft manufacturer that must not be exceeded during normal operations. These limits protect the airframe from structural damage, flutter, and aerodynamic instability. Common operating limits include Vno, Vne, and Vfe for different configurations. Pilots must remain within these speeds to ensure safe flight.}{catAE}
\glsadd{Radius-of-Action}\dictentry{Radius of Action}{The maximum distance an aircraft can travel from its base and return safely with adequate fuel reserves. This value depends on aircraft type, payload, fuel capacity, and prevailing wind conditions. Radius of action is critical for military mission planning and search-and-rescue operations. It determines the operational reach of a force without requiring forward refueling.}{catAE}
\glsadd{Rate-of-Climb}\dictentry{Rate of Climb}{The speed at which an aircraft gains altitude, typically measured in feet per minute. This value depends on available engine power, aircraft weight, and aerodynamic efficiency. The rate of climb decreases with altitude until reaching the absolute ceiling, where climb capability reaches zero. Pilots monitor this value during takeoff and climb to ensure performance meets minimum requirements.}{catAE}
\glsadd{Vne}\dictentry{Vne}{Never-exceed speed, the maximum airspeed at which the aircraft can be operated without risk of structural damage or loss of control. Exceeding this speed may cause flutter, control reversal, or catastrophic structural failure. Vne is typically marked with a red line on the airspeed indicator. Pilots must never intentionally fly at or above this speed under any circumstances.}{catAE}
\glsadd{Vno}\dictentry{Vno}{Maximum structural cruising speed, the highest speed at which the aircraft can be operated in smooth air without exceeding design stress limits. Turbulence encountered at speeds above Vno could cause structural damage. Vno is marked with a green arc upper limit on the airspeed indicator. It represents the normal maximum operating speed for most flight conditions.}{catAE}
\glsadd{Vfe}\dictentry{Vfe}{Maximum flap extended speed, the highest airspeed at which the flaps can be safely extended and operated. Exceeding this speed with flaps down can cause structural damage to the flap mechanism and surfaces. Vfe is typically marked with a white arc upper limit on the airspeed indicator. Pilots must ensure they are below this speed before extending flaps for approach and landing.}{catAE}
\glsadd{Vle}\dictentry{Vle}{Maximum landing gear extended speed, the highest airspeed at which the landing gear can be safely lowered and maintained in the extended position. Flying faster than Vle with the gear down risks structural damage to the gear doors, struts, and mechanisms due to aerodynamic loads. This speed is lower than Vno and is published in the aircraft's operating manual. Pilots must decelerate below this value before extending the landing gear.}{catAE}
\glsadd{Vmca}\dictentry{Vmca}{Minimum control speed in the air, the lowest airspeed at which an aircraft can maintain directional control after an engine failure on a multi-engine airplane. This speed accounts for the asymmetric thrust conditions created by the inoperative engine. Vmca is typically higher with the critical engine feathered and all drag-producing devices retracted. Pilots must fly above this speed during all critical phases of single-engine operation.}{catAE}
\glsadd{Weather-Radar}\dictentry{Weather Radar}{An onboard radar system that detects precipitation and turbulence in the atmosphere ahead of the aircraft. The system transmits microwave pulses that reflect off water droplets and ice particles, creating a display of storm intensity and location. Color coding indicates precipitation severity, allowing pilots to navigate around hazardous weather. Modern weather radar systems can also detect wind shear and ground clutter.}{catAE}
\glsadd{Attitude-Heading-Reference-System}\dictentry{Attitude Heading Reference System}{An electronic system that provides pitch, roll, and heading information using gyroscopic and accelerometer sensors. AHRS replaces traditional mechanical gyros with solid-state devices that are more reliable and require less maintenance. The system outputs attitude data to cockpit displays and autopilot systems. AHRS is a key component of glass cockpit architectures in modern aircraft.}{catAE}
\glsadd{Automatic-Direction-Finder}\dictentry{Automatic Direction Finder}{An avionics system that receives signals from non-directional beacons and displays the bearing to the station. ADF uses a loop antenna to determine the direction of incoming radio signals and indicates relative or magnetic bearing. Pilots use ADF for position fixing, tracking to or from navigational aids, and instrument approaches. While largely supplemented by GPS, ADF remains a backup navigation tool.}{catAE}
\glsadd{Control-Loading}\dictentry{Control Loading}{The artificial feel system that provides pilots with realistic resistance forces on flight controls. In fly-by-wire aircraft, control loading is simulated electronically rather than generated mechanically. This feedback helps pilots judge the amount of input required and prevents over-controlling the aircraft. Control loading systems are calibrated to provide consistent and appropriate forces across different flight conditions.}{catAE}
\glsadd{Ditching}\dictentry{Ditching}{A controlled emergency landing of an aircraft on water. Ditching requires careful management of airspeed, descent rate, and aircraft attitude to maximize survivability upon water impact. The aircraft must be properly configured with flaps and landing gear as recommended for the specific type. Pilots receive specialized training in ditching procedures, and aircraft are designed with flotation features to aid post-landing survival.}{catAE}
\glsadd{Emergency-Locator-Transmitter}\dictentry{Emergency Locator Transmitter}{An onboard device that automatically or manually transmits distress signals on designated frequencies to aid search and rescue operations. ELTs typically activate upon detecting the sudden deceleration of a crash and broadcast a signal that can be detected by satellites and ground stations. Modern 406 MHz ELTs provide more accurate location data than older analog models. ELTs are required equipment on most certified aircraft.}{catAE}
\glsadd{Engine-Indicating-and-Crew-Alerting-System}\dictentry{Engine Indicating and Crew Alerting System}{An integrated avionics system that monitors engine parameters and alerts the crew to abnormal or emergency conditions. EICAS displays information such as engine thrust settings, temperatures, fuel flow, and oil pressure on cockpit screens. The system prioritizes alerts to ensure pilots receive the most critical information first. EICAS reduces crew workload and improves response time to in-flight anomalies.}{catAE}
\glsadd{Fuel-Management}\dictentry{Fuel Management}{The planning, monitoring, and control of fuel consumption throughout all phases of a flight. Effective fuel management ensures sufficient fuel is available for the planned route, reserves, and any diversions. Pilots track fuel balance between tanks, monitor burn rates, and adjust flight profiles as needed. Poor fuel management has been a contributing factor in numerous aviation accidents.}{catAE}
\glsadd{Glide-Slope}\dictentry{Glide Slope}{The vertical guidance component of an Instrument Landing System that provides a defined descent path to the runway. The glideslope signal typically guides aircraft along a three-degree glide path from approximately 2,000 feet above ground level to the touchdown zone. Deviations from the correct path are displayed on cockpit instruments as needle deflections. The glideslope works in conjunction with the localizer for precision approaches.}{catAE}
\glsadd{Ground-Proximity-Warning-System}\dictentry{Ground Proximity Warning System}{An avionics system that alerts pilots when the aircraft is dangerously close to terrain or water. GPWS uses radar altimeter data and barometric information to calculate closure rates with the ground. The system provides voice warnings such as terrain and minimums followed by appropriate pilot actions. Enhanced GPWS integrates GPS and terrain databases for greater accuracy and earlier warnings.}{catAE}
\glsadd{Landing-Distance-Available}\dictentry{Landing Distance Available}{The length of runway available for an aircraft to land and come to a complete stop, measured from the runway threshold to the end of the paved surface. LDA may be shorter than the total runway length if displaced thresholds or obstacles are present. Pilots must ensure that the computed landing distance does not exceed the LDA under existing conditions. This value is critical for performance planning and go-around decision-making.}{catAE}
\glsadd{Load-Manifest}\dictentry{Load Manifest}{A document that details the weight and distribution of all cargo, passengers, and fuel aboard an aircraft before departure. The manifest ensures the aircraft remains within its center of gravity limits and maximum weight restrictions. Pilots and dispatchers use this information to calculate performance data such as takeoff and landing distances. An accurate load manifest is essential for safe flight operations.}{catAE}
\glsadd{Navigation-Fix}\dictentry{Navigation Fix}{A specific geographical point used in air navigation, determined by the intersection of bearings from navigational aids or by GPS coordinates. Fixes serve as waypoints for flight planning, reporting positions, and defining airways and approach procedures. Pilots identify fixes using visual references, radio navigation instruments, or satellite-based systems. Accurate identification of fixes is fundamental to safe and orderly air traffic flow.}{catAE}
\glsadd{Obstacle-Clearance}\dictentry{Obstacle Clearance}{The vertical distance maintained between an aircraft and objects on the ground during departure, approach, and landing procedures. Regulatory minimum obstacle clearance altitudes ensure aircraft can safely clear terrain, buildings, and other obstacles along their flight path. Pilots and flight planners must account for obstacle clearance when selecting altitudes and designing instrument procedures. Insufficient obstacle clearance has been a factor in numerous accidents.}{catAE}
\glsadd{Runway-Visual-Range}\dictentry{Runway Visual Range}{The maximum distance a pilot can see down a runway, measured in meters or feet using instrument-based sensors. RVR accounts for the effects of runway lighting intensity and atmospheric visibility conditions. It is measured at touchdown, midpoint, and rollout zones of the runway. RVR values are critical for determining whether conditions permit landing under instrument flight rules.}{catAE}
\glsadd{Stall-Warning}\dictentry{Stall Warning}{An alert system that notifies the pilot when the aircraft is approaching an aerodynamic stall. The warning typically activates five to ten knots above the actual stall speed using sensors that detect changes in airflow over the wing. Alerts can be aural warnings, stick shakers, or visual indicators. Early stall warning gives pilots time to reduce angle of attack and prevent an unrecoverable stall.}{catAE}
\glsadd{Taxiway}\dictentry{Taxiway}{A paved surface at an airport designed for aircraft movement between runways, terminals, hangars, and other facilities. Taxiways are marked with yellow centerlines and edge markings and are designated by letter or alphanumeric identifiers. Pilots follow taxiway instructions from ground control to navigate the airport surface safely. Proper taxiway use prevents incursions onto active runways and maintains orderly ground operations.}{catAE}
\glsadd{Vertical-Speed-Indicator}\dictentry{Vertical Speed Indicator}{An instrument that displays the rate at which an aircraft is climbing or descending, measured in feet per minute. The VSI measures changes in atmospheric pressure to determine vertical movement. A positive reading indicates a climb while a negative reading indicates descent. The instrument is essential for maintaining assigned altitudes and executing standard instrument departures and arrivals.}{catAE}
\glsadd{Wind-Component}\dictentry{Wind Component}{The portion of wind speed that acts parallel to the runway, affecting an aircraft's takeoff and landing performance. A headwind component reduces ground roll while a tailwind component increases it. Pilots calculate wind components to determine required takeoff and landing distances and to adjust approach speeds. Wind component calculations are essential for safe and efficient flight operations in varying weather conditions.}{catAE}
\glsadd{Crosswind-Component}\dictentry{Crosswind Component}{The portion of wind speed that acts perpendicular to the runway, causing lateral drift during takeoff and landing. Exceeding an aircraft's crosswind capability can result in loss of directional control. Pilots use crosswind correction techniques such as crabbing and side-slipping to maintain runway alignment. Each aircraft has a maximum demonstrated crosswind component published in its flight manual.}{catAE}
\glsadd{Go-Around}\dictentry{Go-Around}{A pilot-initiated maneuver to abort a landing approach and climb back to a safe altitude. This may be executed due to an unstabilized approach, runway obstruction, or air traffic control instructions. The procedure involves applying maximum available power, retracting flaps to the takeoff setting, and climbing on the published go-around path. Go-arounds are a normal and essential safety procedure in everyday aviation operations.}{catAE}
\glsadd{Missed-Approach}\dictentry{Missed Approach}{A published procedure that an aircraft follows when a landing cannot be completed at the published minimums. The missed approach point is a defined location on an instrument approach where the pilot must decide to land or go around. Following this procedure ensures obstacle clearance and safe separation from terrain and other traffic. Missed approach procedures are published on instrument approach charts and briefed before every approach.}{catAE}
\glsadd{Holding-Pattern}\dictentry{Holding Pattern}{A predefined flight path that aircraft follow while waiting for clearance to proceed, typically depicted as an oval racetrack pattern. Standard holds are made up of two 180-degree turns connected by two straight legs of equal length. Aircraft fly the pattern at specified altitudes and speeds while remaining within protected airspace. Holding patterns are used for traffic management, weather delays, and sequencing into busy airports.}{catAE}
\glsadd{Traffic-Pattern}\dictentry{Traffic Pattern}{A standardized rectangular flight path that aircraft follow in the vicinity of an airport for takeoff and landing operations. The pattern consists of upwind, crosswind, downwind, base, and final legs at specified altitudes and directions. Standard patterns are flown to the left, though right-hand patterns may be designated where terrain or noise abatement requires. Adherence to the traffic pattern ensures orderly and safe operations at non-towered airports.}{catAE}
\glsadd{Final-Approach}\dictentry{Final Approach}{The last segment of an aircraft's flight path leading to the runway threshold for landing. On a visual approach, the pilot aligns with the runway and configures the aircraft for landing. In an instrument approach, the final approach segment follows a defined course and descent gradient to the landing area. Proper configuration, speed, and alignment on final are critical for a safe touchdown.}{catAE}
\glsadd{Base-Leg}\dictentry{Base Leg}{The portion of a standard traffic pattern flown perpendicular to the runway between the downwind leg and final approach. During the base leg, the pilot begins configuring the aircraft for landing by reducing power and extending flaps. The turn from base to final must be coordinated to avoid excessive bank angles at low altitude. Proper timing of the base leg determines the aircraft's spacing and alignment on final approach.}{catAE}
\glsadd{Downwind-Leg}\dictentry{Downwind Leg}{The segment of a standard traffic pattern flown parallel to the runway in the opposite direction of landing. Aircraft maintain pattern altitude and speed while positioned one to two miles from the runway. Pilots perform pre-landing checklists and configure the aircraft during this phase. The downwind leg provides the framework for proper spacing and sequencing with other traffic in the pattern.}{catAE}
\glsadd{Straight-In-Approach}\dictentry{Straight-In Approach}{An approach to landing where the aircraft does not enter the standard traffic pattern but instead flies directly to the final approach segment. Straight-in approaches are common for instrument approaches or when directed by air traffic control. They require careful planning to ensure proper configuration and spacing. While efficient, straight-in approaches may conflict with traffic in the standard pattern if not properly coordinated.}{catAE}
\glsadd{Circling-Approach}\dictentry{Circling Approach}{An instrument approach procedure that allows a pilot to circle around the airport to land on a runway different from the one aligned with the initial instrument approach segment. The pilot descends to the circling minimums and then visually maneuvers to the landing runway while maintaining obstacle clearance. Circling approaches require high situational awareness and are limited by visibility and cloud ceiling conditions. They are used when a straight-in approach to the desired runway is not available.}{catAE}
\glsadd{Visual-Approach}\dictentry{Visual Approach}{An approach to landing conducted under visual meteorological conditions where the pilot maintains visual reference to the ground and surrounding terrain. The pilot may be cleared by air traffic control to join the traffic pattern or fly a straight-in approach. Visual approaches allow more flexibility and efficiency than instrument approaches. They require sufficient visibility and ceiling to safely identify the airport and runway environment.}{catAE}
\glsadd{Instrument-Approach}\dictentry{Instrument Approach}{A published procedure that allows pilots to descend from cruise altitude to a point where a landing can be made using cockpit instruments alone. Instrument approaches are designed using ground-based radio aids or satellite navigation and include specific courses, altitudes, and minimum descent altitudes. The procedure enables landings in weather conditions too poor for visual flight. Standardized instrument approaches are essential for maintaining safe air traffic flow in low-visibility conditions.}{catAE}
\glsadd{Precision-Approach}\dictentry{Precision Approach}{An instrument approach that provides both lateral and vertical guidance to the runway, allowing descent along a defined glide path. The most common precision approaches are ILS, GBAS Landing System, and PAR. These procedures offer lower minimums than non-precision approaches, permitting landings in worse weather conditions. Precision approaches provide the highest level of guidance accuracy and safety for instrument landings.}{catAE}
\glsadd{Non-Precision-Approach}\dictentry{Non-Precision Approach}{An instrument approach that provides only lateral course guidance to the runway without vertical path guidance. Examples include VOR, NDB, and RNAV approaches where the pilot descends stepwise to published altitudes. Non-precision approaches typically have higher minimums than precision approaches. Pilots must carefully manage descent rates and altitudes to avoid premature descent below the minimum descent altitude.}{catAE}
\glsadd{Decision-Altitude}\dictentry{Decision Altitude}{The altitude at which a pilot on a precision approach must have the required visual references to continue the landing or execute a missed approach. Below DA, the pilot must have the runway environment in sight to proceed. If visual references are not acquired, a mandatory go-around must be initiated. DA is published on approach charts and accounts for obstacle clearance and signal accuracy.}{catAE}
\glsadd{Minimum-Descent-Altitude}\dictentry{Minimum Descent Altitude}{The lowest altitude to which a pilot may descend on a non-precision approach without having the runway environment in sight. Unlike decision altitude, MDA is a hard floor that the pilot may level off at and fly to the missed approach point. If the runway is not visible by that point, a go-around is mandatory. MDA provides obstacle clearance during the final segment of non-precision approaches.}{catAE}
\glsadd{Localizer}\dictentry{Localizer}{The lateral guidance component of an Instrument Landing System that transmits a signal aligned with the runway centerline. The localizer operates in the VHF frequency range and provides left-right guidance to pilots during approach. Deviation from the centerline is displayed on the cockpit instrument as a needle deflection. The localizer works with the glideslope to provide full precision guidance for landing in low-visibility conditions.}{catAE}
\glsadd{Glideslope-Antenna}\dictentry{Glideslope Antenna}{The ground-based antenna that transmits the vertical guidance signal for an ILS glideslope. Typically located beside the runway at approximately 750 feet from the threshold, it broadcasts a signal that defines the correct descent angle, usually around three degrees. The antenna transmits two overlapping signals at different modulation frequencies, creating a null zone along the glide path. Aircraft receivers interpret these signals to display vertical deviation on cockpit instruments.}{catAE}
\glsadd{Marker-Beacon}\dictentry{Marker Beacon}{An electronic device that transmits a directional signal upward to identify specific points along an instrument approach path. Marker beacons transmit on 75 MHz and are identified by different audio tones and cockpit lights corresponding to outer, middle, and inner markers. They help pilots verify their position along the approach and confirm the correct distance to the runway. Marker beacons are being gradually supplemented by GPS-based position reporting.}{catAE}
\glsadd{Distance-Measuring-Equipment}\dictentry{Distance Measuring Equipment}{An airborne and ground-based radio system that provides continuous slant-range distance from an aircraft to a ground station. DME operates by measuring the round-trip time of a UHF signal transmitted between the aircraft and the ground facility. The distance is displayed in nautical miles on cockpit instruments. DME is commonly paired with VOR stations and is a fundamental tool for en route navigation and approach procedures.}{catAE}
\glsadd{VOR}\dictentry{VOR}{VHF Omnidirectional Range, a ground-based radio navigation aid that provides bearing information to and from the station. VOR transmits two signals that create a phase difference proportional to the aircraft's radial position relative to the station. Pilots use VOR radials for course tracking, position fixing, and airway navigation. VOR has been the backbone of the airway navigation system for decades and remains widely used worldwide.}{catAE}
\glsadd{VORTAC}\dictentry{VORTAC}{A combined navigation facility that co-locates a VOR with TACAN, providing both civilian and military users with bearing and distance information. VORTAC offers the civilian VOR function for commercial aviation while TACAN serves military aircraft with a combined bearing and distance capability. This co-location ensures interoperability between civil and military navigation systems. VORTAC stations are a key component of the airway structure across many countries.}{catAE}
\glsadd{Non-Directional-Beacon}\dictentry{Non-Directional Beacon}{A low or medium frequency radio beacon that transmits an omnidirectional signal for aircraft navigation. NDB signals are automatically direction-finder-equipped aircraft to determine bearing to the beacon. NDBs are relatively simple and inexpensive to maintain, making them useful in remote areas and developing regions. Despite their limitations in accuracy and range, NDBs continue to serve as backup navigation aids and support approach procedures.}{catAE}
\glsadd{ADF}\dictentry{ADF}{Automatic Direction Finder, an airborne radio receiver that determines the bearing to a non-directional beacon or commercial AM broadcast station. The ADF uses a loop antenna or sense antenna to detect the direction of incoming radio signals. It displays bearing information on a relative or magnetic indicator that the pilot interprets for navigation. While largely supplemented by GPS and more modern systems, ADF remains a certified navigation backup.}{catAE}
\glsadd{Traffic-Information-Service}\dictentry{Traffic Information Service}{A free broadcast service provided by air traffic control that advises pilots of nearby radar-detected aircraft. TIS enhances pilot awareness of surrounding traffic beyond what is provided by standard ATC instructions. The service transmits bearing, range, altitude, and trend information to equipped aircraft cockpit displays. TIS works in conjunction with Mode C and Mode S transponders to deliver position data to pilots.}{catAE}
\glsadd{Safety-Management-System}\dictentry{Safety Management System}{A systematic approach to managing safety that includes organizational policies, procedures, and practices for identifying hazards and mitigating risks. SMS is required by most aviation regulatory authorities for airlines, airports, and maintenance organizations. The framework includes safety risk management, safety assurance, safety policy, and safety promotion. SMS represents a shift from reactive to proactive safety management in aviation.}{catAE}
\glsadd{Airworthiness-Directive}\dictentry{Airworthiness Directive}{A legally enforceable regulation issued by an aviation authority that requires correction of an unsafe condition in an aircraft, engine, or component. ADs are issued when a design defect, manufacturing flaw, or operational issue is discovered that affects safety. Compliance is mandatory and must be completed within specified timeframes. ADs ensure that all affected aircraft are modified or inspected to maintain safe operation.}{catAE}
\glsadd{Production-Certificate}\dictentry{Production Certificate}{An authorization issued by an aviation authority that permits a manufacturer to produce aircraft or aircraft components under approved quality standards. The certificate requires the manufacturer to maintain a production inspection system and comply with all applicable design specifications. Holders must undergo regular audits to retain their certification. A production certificate is essential for selling aircraft or parts in regulated aviation markets.}{catAE}
\glsadd{Deicing-System}\dictentry{Deicing System}{An onboard system designed to remove ice that has already accumulated on aircraft surfaces such as wings, tail, and engine inlets. Deicing systems typically use heated air, electric heating elements, or fluid application to break the bond between ice and the airframe. Effective deicing is critical because ice disrupts airflow over wings and reduces lift. These systems are essential for safe operation in icing conditions encountered during flight.}{catAE}
\glsadd{Anti-Icing-System}\dictentry{Anti-Icing System}{An onboard system that prevents ice from forming on critical aircraft surfaces during flight through continuous heating or fluid application. Anti-icing differs from deicing in that it acts before ice accumulation occurs rather than removing existing ice. Common methods include thermal anti-ice using engine bleed air and electrothermal systems. Anti-icing is essential for aircraft certification for flight into known icing conditions.}{catAE}
\glsadd{Fire-Suppression}\dictentry{Fire Suppression}{An onboard system designed to detect and extinguish fires in the aircraft's engines, cargo holds, and avionics bays. The system typically uses halon or other clean extinguishing agents discharged through overhead nozzles. Fire detection uses thermal sensors or continuous-loop detection wire that triggers alarms and suppression activation. Fire suppression systems are mandatory on all commercial aircraft and are regularly inspected and serviced.}{catAE}
\glsadd{Hydraulic-Pump}\dictentry{Hydraulic Pump}{A mechanical or electrically driven device that pressurizes hydraulic fluid to power aircraft systems such as flight controls, landing gear, and brakes. Hydraulic pumps are typically driven by the engine through accessory gearboxes or electric motors. Multiple independent hydraulic systems provide redundancy for critical flight functions. High-pressure hydraulic systems allow powerful and precise actuation of large control surfaces with relatively small actuator sizes.}{catAE}
\glsadd{Thermal-Anti-Ice}\dictentry{Thermal Anti-Ice}{A method of preventing ice formation on aircraft surfaces by routing hot engine bleed air through internal passages in the leading edges of wings and engine inlets. The heated air warms the surface above the freezing point, preventing supercooled water droplets from freezing on contact. Thermal anti-ice systems are effective but consume engine bleed air, which reduces overall engine efficiency. They are the primary anti-ice method on most commercial jet aircraft.}{catAE}
\glsadd{Windshield-Heating}\dictentry{Windshield Heating}{An electrical or bleed-air system that heats the cockpit windshields to prevent ice formation, fogging, and to maintain structural integrity. Electrically heated windshields use embedded conductive elements to generate warmth across the glass surface. Some systems also provide rain repellent by evaporating water on contact with the heated surface. Proper windshield heating is essential for pilot visibility and safety during all weather conditions.}{catAE}
\glsadd{Ground-Handling}\dictentry{Ground Handling}{The range of services and procedures performed on an aircraft while it is on the ground between flights. This includes refueling, baggage and cargo loading, catering, cleaning, and pushback. Ground handling can be performed by airline staff or outsourced to specialized ground handling companies. Efficient ground handling minimizes turnaround time and reduces delays while maintaining safety standards.}{catAE}
\glsadd{Pushback}\dictentry{Pushback}{The process of moving an aircraft backward from its gate or parking position using a specialized tug vehicle. Most commercial aircraft cannot reverse under their own power and require a pushback tractor connected to the nose gear. The pushback driver follows signals from a ground marshal or wing walkers to safely maneuver the aircraft. After pushback is complete, the aircraft starts its engines and taxis to the runway under its own power.}{catAE}
\glsadd{Turnaround}\dictentry{Turnaround}{The entire process of servicing an aircraft between arrival and its next departure, including unloading, cleaning, refueling, maintenance checks, and loading. Turnaround time is a critical metric for airline efficiency and profitability. A well-coordinated turnaround minimizes ground time and maximizes aircraft utilization. Airlines use detailed checklists and specialized teams to complete turnarounds within tight schedules.}{catAE}
\glsadd{Line-Maintenance}\dictentry{Line Maintenance}{Routine maintenance activities performed on an aircraft between flights or during overnight stops to keep it in an airworthy condition. Line maintenance includes daily inspections, minor repairs, component replacements, and servicing tasks such as fluid replenishment. These tasks are typically performed by certified maintenance technicians at the airport. Line maintenance is the first line of defense in ensuring ongoing aircraft safety and reliability.}{catAE}
\glsadd{Heavy-Maintenance}\dictentry{Heavy Maintenance}{Comprehensive, scheduled maintenance events that involve extensive inspection, repair, and overhaul of an aircraft's structure, systems, and components. Heavy maintenance checks are performed in specialized hangars over several weeks or months. These checks involve detailed inspections of the airframe, engines, avionics, and all mechanical systems. Heavy maintenance is essential for extending the aircraft's operational life and maintaining airworthiness certification.}{catAE}
\glsadd{Letter-Check}\dictentry{Letter Check}{A scheduled maintenance inspection category designated by a letter, such as A, B, C, or D, with increasing levels of scope and complexity. Each letter check has specific tasks, intervals, and requirements defined by the aircraft manufacturer and regulatory authority. The checks progress from relatively brief line-level inspections to major structural overhauls. Letter checks form the backbone of an aircraft's maintenance program.}{catAE}
\glsadd{A-Check}\dictentry{A-Check}{The most frequent and least intensive scheduled maintenance check in the letter check system, typically performed every 400 to 600 flight hours. An A-check involves visual inspections of key systems, fluid level checks, and servicing of routine items. The check is usually completed overnight or during a short ground period. A-checks ensure the aircraft remains in safe operating condition between more comprehensive maintenance events.}{catAE}
\glsadd{C-Check}\dictentry{C-Check}{A comprehensive scheduled maintenance inspection performed approximately every 20 to 24 months that requires the aircraft to be out of service for several days. The C-check involves detailed inspection of the airframe, flight controls, engines, and all major systems. Panels and access points are opened for thorough examination that goes beyond visual inspection. C-checks are performed in hangars and require significant manpower and engineering support.}{catAE}
\glsadd{D-Check}\dictentry{D-Check}{The most extensive scheduled maintenance event in the letter check system, performed approximately every six to ten years. A D-check involves complete disassembly and detailed inspection of virtually every component of the aircraft. The airframe is stripped, inspected for corrosion and fatigue, and repainted. The check can take several months to complete and may cost millions of dollars. D-checks are essential for extending an aircraft's operational life to its maximum potential.}{catAE}
\glsadd{Airframe}\dictentry{Airframe}{The main structural body of an aircraft, excluding the engines, avionics, and powerplant. The airframe includes the fuselage, wings, tail surfaces, and all structural components that bear flight loads. It is designed and manufactured to withstand the stresses of flight, pressurization, and landing throughout the aircraft's operational life. Airframe integrity is paramount to safe flight and is the primary focus of structural maintenance programs.}{catAE}
\glsadd{Fuselage-Station}\dictentry{Fuselage Station}{A longitudinal measurement system used to locate structural components and systems along the length of the aircraft fuselage. Fuselage stations are typically measured in inches from a reference point called the datum, usually located at or near the nose. These measurements are essential for weight and balance calculations and for locating structural reference points during maintenance. Every component in the fuselage has an associated fuselage station number.}{catAE}
\glsadd{Waterline}\dictentry{Waterline}{A vertical measurement system used in aircraft design and maintenance that references distances above or below a horizontal datum line. Waterlines indicate the height of components and structural features relative to this reference plane. Along with fuselage stations and butt lines, waterlines form the three-dimensional coordinate system used to locate any point on the aircraft. Waterline measurements are critical for weight and balance and structural repair procedures.}{catAE}
\glsadd{Butt-Line}\dictentry{Butt Line}{A lateral measurement system used to locate components and reference points along the wingspan of an aircraft. Butt lines are measured in inches from the aircraft centerline, with left and right designations. Together with fuselage stations and waterlines, butt lines provide a complete three-dimensional coordinate system for the airframe. These measurements are essential for weight and balance calculations and structural repair documentation.}{catAE}
\glsadd{Aircraft-Classification-Number}\dictentry{Aircraft Classification Number}{A number assigned to an aircraft to indicate its relative impact on airport pavement surfaces during landing and takeoff operations. The ACN is calculated based on the aircraft's weight, landing gear configuration, and tire pressure. It is used in conjunction with the pavement classification number to determine whether an airport's runway can support a particular aircraft. The ACN-PCN system provides a standardized method for evaluating pavement compatibility.}{catAE}
\glsadd{Pavement-Classification-Number}\dictentry{Pavement Classification Number}{A number that indicates the load-carrying capacity of an airport's pavement for aircraft operations. The PCN is determined through engineering analysis of the pavement structure and is published in the airport's aeronautical information. An aircraft can operate on a pavement as long as its ACN does not exceed the reported PCN. This system allows airports to communicate their pavement capabilities to operators worldwide using a standardized format.}{catAE}
\glsadd{Runway-Strength}\dictentry{Runway Strength}{A measure of an airport runway's ability to support aircraft loads without structural damage. Runway strength is expressed as a pavement classification number and depends on the pavement thickness, subgrade conditions, and construction quality. The strength determines which aircraft types and maximum weights can safely use the runway. Runway strength is tested periodically and reported to ensure operators can plan operations within safe limits.}{catAE}
\glsadd{FOD}\dictentry{FOD}{Foreign Object Debris, any material or object found on the airport surface that could damage aircraft or engines. FOD includes items such as loose hardware, pavement fragments, wildlife remains, and discarded equipment. Ingestion of FOD by jet engines can cause catastrophic damage. Airports conduct regular FOD walks and inspections to minimize this hazard. FOD prevention is a critical safety concern across all airport operations.}{catAE}
\glsadd{RAM}\dictentry{RAM}{Radio Absorbing Material, a substance designed to absorb electromagnetic radiation rather than reflect it. In aerospace applications, RAM is used to reduce the radar cross-section of aircraft for stealth purposes. The material works by converting incoming radar energy into heat through dielectric or magnetic absorption. RAM coatings and structures are carefully engineered to be effective across specific frequency ranges while minimizing weight and aerodynamic impact.}{catAE}
\glsadd{Windshear}\dictentry{Windshear}{A sudden and significant change in wind speed or direction over a short distance in the atmosphere. Windshear can be caused by temperature inversions, thunderstorm downdrafts, or terrain effects. It poses a serious threat to aircraft during takeoff and landing, where it can cause sudden loss of airspeed and altitude. Modern windshear detection systems and pilot training have significantly improved safety in windshear conditions.}{catAE}
\glsadd{Microburst}\dictentry{Microburst}{A localized, intense downdraft that strikes the ground and spreads outward in all directions, creating severe windshear. Microbursts are typically associated with thunderstorms and produce dangerous wind shifts that can affect aircraft during the critical approach and landing phase. The phenomenon can produce wind changes of up to 100 knots over very short distances. Microburst detection systems and pilot training programs have greatly reduced the risk of microburst-related accidents.}{catAE}
\glsadd{Clear-Air-Turbulence}\dictentry{Clear Air Turbulence}{Turbulence encountered in cloud-free regions of the atmosphere, often associated with jet streams, mountain waves, or frontal boundaries. CAT cannot be detected by conventional weather radar because it occurs in the absence of moisture and precipitation. It can cause sudden and violent aircraft movements, posing risks to passenger safety and structural integrity. CAT is often forecast using numerical weather prediction models and pilot reports.}{catAE}
\glsadd{Lee-Wave}\dictentry{Lee Wave}{A type of atmospheric gravity wave formed when stable air flows over a mountain range. The air oscillates up and down on the downwind side of the mountains, creating a series of waves that can extend far downstream. Lee waves can generate severe turbulence, particularly in the rotor zones near the surface. Pilots should avoid flying in the lee of mountains in strong winds to prevent encountering these turbulence-producing waves.}{catAE}
\glsadd{Mountain-Wave}\dictentry{Mountain Wave}{A standing atmospheric wave system that forms when strong winds blow perpendicular to a mountain range. The wave pattern extends from the mountain crest through the entire troposphere and can reach into the stratosphere. Mountain waves produce turbulence, severe updrafts and downdrafts, and can generate lenticular clouds as visible indicators. These waves can affect aircraft performance and safety over distances far greater than the mountains themselves.}{catAE}
\glsadd{Wake-Turbulence}\dictentry{Wake Turbulence}{The disturbance in the atmosphere created behind an aircraft as it passes through the air, consisting primarily of wingtip vortices. Wake turbulence is strongest behind large, heavy aircraft and can persist for several minutes. Following aircraft can encounter these vortices, which can cause uncommanded roll and structural stress. Air traffic controllers apply separation standards to ensure following aircraft avoid dangerous wake turbulence encounters.}{catAE}
\glsadd{Jet-Blast}\dictentry{Jet Blast}{The high-velocity exhaust stream from a jet engine that extends behind the aircraft and can cause significant damage to objects and other aircraft on the ground. Jet blast can reach temperatures exceeding 200 degrees and velocities over 100 knots near the runway. Airports establish jet blast barriers and minimum distance requirements to protect ground personnel and equipment. Jet blast is a major consideration during taxi, takeoff, and parking operations.}{catAE}
\glsadd{Propwash}\dictentry{Propwash}{The disturbed airflow behind a propeller-driven aircraft created by the spinning blades. Propwash produces a spiraling pattern of air that can affect other aircraft and ground personnel in the vicinity. The wash extends well behind the propeller and can cause damage to loose objects on the ground. Pilots and ground crews must be aware of propwash hazards when operating near running propeller aircraft.}{catAE}
\glsadd{Slipstream-Effect}\dictentry{Slipstream Effect}{The acceleration of air behind a propeller or rotor that increases the velocity of airflow over portions of the aircraft's structure. The slipstream can affect tail surface effectiveness, cause asymmetric forces on single-engine aircraft, and influence handling characteristics. Pilots must account for slipstream effects during takeoff, climb, and when operating at low airspeeds with high power settings. Understanding slipstream behavior is essential for safe single-engine operation.}{catAE}
\glsadd{Asymmetric-Thrust}\dictentry{Asymmetric Thrust}{A condition in multi-engine aircraft where one engine produces significantly less thrust than the other, creating a yawing moment. Asymmetric thrust can result from engine failure, fuel imbalance, or malfunctioning engine controls. The pilot must apply corrective rudder and maintain directional control using available control inputs. Managing asymmetric thrust is a critical skill for multi-engine pilots and is emphasized in simulator training.}{catAE}
\glsadd{Critical-Engine}\dictentry{Critical Engine}{The engine on a multi-engine aircraft whose failure would have the most adverse effect on flight performance due to its position and the P-factor effect. The critical engine is typically the one that, when inoperative, produces the greatest yawing moment. Pilots must identify the critical engine and understand the control inputs required to maintain safe flight if it fails. Knowledge of the critical engine is essential for multi-engine flight training and operations.}{catAE}
\glsadd{Vmcg}\dictentry{Vmcg}{Minimum control speed on the ground, the lowest airspeed at which an aircraft can safely continue a takeoff after an engine failure. At this speed, the rudder and other directional controls can maintain runway alignment despite the asymmetric thrust from the failed engine. Vmcg is typically higher than the corresponding in-air minimum control speed due to reduced control effectiveness on the ground. Pilots must ensure takeoff speeds are above Vmcg to ensure safety.}{catAE}
\glsadd{Balanced-Field-Length}\dictentry{Balanced Field Length}{A takeoff performance calculation that accounts for both the takeoff distance required and the accelerate-stop distance required. The balanced field length is the runway length at which the aircraft can either continue the takeoff or abort and stop safely within the available distance. It represents the most critical condition for takeoff safety. Pilots must verify that the available runway meets or exceeds the balanced field length for current conditions.}{catAE}
\glsadd{Takeoff-Distance-Required}\dictentry{Takeoff Distance Required}{The total ground distance an aircraft needs to accelerate to takeoff speed, lift off, and clear a 35-foot obstacle at the end of the runway. TDR depends on aircraft weight, runway slope, atmospheric conditions, and engine performance. It must be compared against the takeoff distance available on the intended runway. Accurate calculation of TDR is essential for ensuring safe takeoff operations under all conditions.}{catAE}
\glsadd{Accelerate-Stop-Distance}\dictentry{Accelerate-Stop Distance}{The total distance required for an aircraft to accelerate from brake release to the decision speed and then decelerate to a complete stop. This distance accounts for the engine failure recognition time, pilot reaction time, and braking performance. It is one of the critical performance calculations for takeoff planning. The available runway must be at least as long as the computed accelerate-stop distance for safe operation.}{catAE}
\glsadd{Climb-Gradient}\dictentry{Climb Gradient}{The rate of altitude gain per unit of horizontal distance traveled, expressed as a percentage. A three-percent climb gradient means the aircraft gains three feet in altitude for every one hundred feet of horizontal distance. Climb gradients are specified in instrument departure procedures to ensure obstacle clearance. Pilots must achieve minimum climb gradients with one engine inoperative to meet regulatory safety requirements during departure.}{catAE}
\end{multicols}

\clearpage
\chapter*{Rockets and Missiles}
\addcontentsline{toc}{chapter}{Rockets and Missiles}
\markboth{Rockets and Missiles}{Rockets and Missiles}
\clearpage
\begin{multicols}{2}
\small
Rocket propulsion, launch vehicle design, and guided missile systems. From the Tsiolkovsky equation to terminal guidance and hypersonic glide vehicles.\par
\vspace{0.5cm}
\glsadd{Rocket-Engine}\dictentry{Rocket Engine}{A Rocket Engine is a propulsion device that generates thrust by expelling mass at high velocity through a nozzle. It operates on Newton's Third Law: as hot gas accelerates rearward, an equal reaction force pushes the vehicle forward. Rocket engines function in the vacuum of space where no atmospheric oxygen is available. They power launch vehicles, spacecraft maneuvering systems, and missiles across atmospheric and orbital regimes.}{catRM}
\glsadd{Solid-Rocket-Motor}\dictentry{Solid Rocket Motor}{A Solid Rocket Motor is a rocket engine that uses solid chemical propellant cast directly into the motor casing. The propellant grain is ignited and burns from its exposed surfaces outward, requiring no pumps or complex plumbing. Solid motors offer high reliability, simple construction, and the ability to remain in a ready state for extended periods.}{catRM}
\glsadd{Liquid-Rocket-Engine}\dictentry{Liquid Rocket Engine}{A Liquid Rocket Engine uses separately stored liquid fuel and oxidizer pumped into a combustion chamber for burning. Turbopumps or pressure feeds deliver propellants at precise flow rates, enabling throttle control and restart capability. Liquid engines achieve higher specific impulse than solid motors, making them efficient for orbital insertion and deep-space maneuvers.}{catRM}
\glsadd{Hybrid-Rocket}\dictentry{Hybrid Rocket}{A Hybrid Rocket combines a solid fuel grain with a liquid or gaseous oxidizer, blending advantages of both solid and liquid systems. The oxidizer is injected into the solid grain, where combustion occurs at the interface. Hybrid rockets offer throttle capability and safer handling than fully solid motors, since the fuel cannot ignite without the oxidizer.}{catRM}
\glsadd{Total-Impulse}\dictentry{Total Impulse}{Total Impulse is the integral of thrust over the entire burn duration of a rocket motor. It represents the total momentum change delivered by the motor and is measured in Newton-seconds or pound-force-seconds. Total impulse determines the overall capability of a rocket motor class and is used to classify model and amateur rocket motors. It directly correlates with the maximum velocity change the motor can impart to a given vehicle mass.}{catRM}
\glsadd{Thrust-to-Weight-Ratio}\dictentry{Thrust-to-Weight Ratio}{Thrust-to-Weight Ratio is the ratio of engine thrust to the total weight of the vehicle at a given moment. A ratio greater than one is required for the rocket to lift off the pad, with typical values ranging from 1.2 to 1.5. As propellant burns and the vehicle becomes lighter, this ratio increases, causing acceleration to grow throughout the burn. Engineers must balance this ratio to stay within structural and payload acceleration limits.}{catRM}
\glsadd{Propellant}\dictentry{Propellant}{Propellant is the chemical substance or combination of substances burned in a rocket engine to produce thrust. It includes both the fuel, which is oxidized, and the oxidizer, which supplies the oxygen for combustion. Propellant selection determines specific impulse, storage requirements, and engine complexity. The choice between solid, liquid, and hybrid propellants is a primary driver of overall rocket design and mission capability.}{catRM}
\glsadd{Oxidizer}\dictentry{Oxidizer}{An Oxidizer is the component of a rocket propellant that supplies oxygen for the combustion of fuel. Since rockets must operate in the vacuum of space, they carry their own oxidizer rather than relying on atmospheric oxygen. Common oxidizers include liquid oxygen, nitrogen tetroxide, and hydrogen peroxide. The oxidizer typically makes up the majority of propellant mass, and its chemical properties strongly influence engine performance and safety.}{catRM}
\glsadd{Bipropellant}\dictentry{Bipropellant}{A Bipropellant system uses two separate propellants, a fuel and an oxidizer, stored in independent tanks and mixed in the combustion chamber. This separation allows the use of high-energy propellant combinations that cannot be pre-mixed safely. Bipropellant engines are the standard for orbital launch vehicles because they deliver the highest performance.}{catRM}
\glsadd{Monopropellant}\dictentry{Monopropellant}{A Monopropellant system uses a single chemical that decomposes exothermically when passed over a catalyst, producing hot gas for thrust. Hydrazine is the most common monopropellant, offering reliable operation without an ignition system. These engines produce relatively low thrust and specific impulse but are valued for simplicity and reliability. They are widely used for spacecraft attitude control and orbital maintenance maneuvers.}{catRM}
\glsadd{Hypergolic}\dictentry{Hypergolic}{Hypergolic propellants are fuel and oxidizer combinations that ignite spontaneously upon contact, eliminating the need for an ignition system. This characteristic makes them extremely reliable for applications requiring reliable restart, such as orbital maneuvering and reaction control. The most common pairing is nitrogen tetroxide with a hydrazine-based fuel.}{catRM}
\glsadd{Cryogenic-Propellant}\dictentry{Cryogenic Propellant}{Cryogenic Propellants are rocket propellants stored at extremely low temperatures to maintain them in liquid form. Liquid hydrogen and liquid oxygen are the most widely used pair, offering the highest chemical performance of any practical combination. Storage requires insulated tanks and continuous thermal management to prevent boil-off.}{catRM}
\glsadd{RP-1}\dictentry{RP-1}{RP-1 is a highly refined form of kerosene used as rocket fuel in liquid-propellant engines. It is stable at room temperature, dense, and relatively inexpensive compared to cryogenic fuels. RP-1 is typically burned with liquid oxygen and is the fuel of choice for first-stage engines such as the F-1 and Merlin. Its high density allows smaller tanks and contributes to a compact, high-thrust vehicle configuration.}{catRM}
\glsadd{Liquid-Hydrogen}\dictentry{Liquid Hydrogen}{Liquid Hydrogen is a cryogenic fuel stored at minus 253 degrees Celsius, used in high-performance rocket engines. When combined with liquid oxygen, it achieves the highest specific impulse of any chemical propellant combination. Its extremely low density requires large tank volumes, which is acceptable in upper stages where mass efficiency outweighs volume constraints.}{catRM}
\glsadd{Liquid-Oxygen}\dictentry{Liquid Oxygen}{Liquid Oxygen is a cryogenic oxidizer stored at minus 183 degrees Celsius, used in virtually all high-performance liquid rocket engines. It is non-toxic, abundant, and relatively inexpensive to produce through air separation. Liquid oxygen is paired with fuels such as liquid hydrogen, RP-1, and methane across a wide range of engines. Its handling requires strict contamination control since liquid oxygen reacts vigorously with many materials.}{catRM}
\glsadd{Nitrogen-Tetroxide}\dictentry{Nitrogen Tetroxide}{Nitrogen Tetroxide is a storable liquid oxidizer widely used in hypergolic rocket engines. It is self-pressurizing and remains liquid across a wide temperature range, making it suitable for long-duration space missions. When contacted with hydrazine-based fuels, it ignites spontaneously without any ignition source.}{catRM}
\glsadd{Hydrazine}\dictentry{Hydrazine}{Hydrazine is a versatile chemical used as both a monopropellant and a bipropellant fuel in rocket engines. As a monopropellant, it decomposes over an iridium catalyst to produce hot nitrogen and hydrogen gases. When used as a fuel with an oxidizer like nitrogen tetroxide, it provides hypergolic ignition and reliable restart.}{catRM}
\glsadd{MMH}\dictentry{MMH}{MMH, or Monomethylhydrazine, is a hypergolic fuel commonly used in spacecraft reaction control systems and upper-stage engines. It ignites on contact with oxidizers such as nitrogen tetroxide, providing reliable and repeatable combustion. MMH offers better thermal stability and lower freezing point than pure hydrazine, making it suitable for long-duration space missions.}{catRM}
\glsadd{UDMH}\dictentry{UDMH}{UDMH, or Unsymmetrical Dimethylhydrazine, is a hypergolic fuel used extensively in military and orbital launch vehicle engines. It combines with nitrogen tetroxide to produce reliable hypergolic ignition for storable propulsion systems. UDMH offers good performance and can be stored for extended periods, although it is toxic and suspected of being a carcinogen. It was the primary fuel for engines of the Titan and Proton launch vehicle families.}{catRM}
\glsadd{Propellant-Tank}\dictentry{Propellant Tank}{A Propellant Tank is a pressure vessel designed to store liquid or solid propellant aboard a rocket vehicle. Tank design must balance structural integrity, mass, thermal management, and propellant settleability in microgravity. Common materials include aluminum alloys, stainless steel, and composite overwrapped structures.}{catRM}
\glsadd{Tank-Pressurization}\dictentry{Tank Pressurization}{Tank Pressurization applies gas pressure above the liquid propellant to force it toward the engine. Methods include pressurant gas from separate bottles, autogenous pressurization using heated engine exhaust, and electric pump-fed systems. Stable tank pressure prevents pump cavitation and maintains consistent engine performance. The pressurization system adds mass but eliminates the need for extremely heavy tank structures.}{catRM}
\glsadd{Injector}\dictentry{Injector}{The Injector is the component of a rocket engine that delivers and mixes fuel and oxidizer into the combustion chamber. It must produce uniform spray patterns, maintain proper mixing ratios, and prevent combustion instability. Injector design varies from simple showerhead patterns to complex impinging and coaxial element configurations. The injector face also serves as a critical thermal boundary, often requiring regenerative cooling.}{catRM}
\glsadd{De-Laval-Nozzle}\dictentry{De Laval Nozzle}{The De Laval Nozzle is a convergent-divergent nozzle design that accelerates gas from subsonic to supersonic velocity. It was developed by Gustaf de Laval for steam turbines and later adopted for rocket engines. The convergent section accelerates flow to Mach 1 at the throat, and the divergent section expands it further to supersonic speeds.}{catRM}
\glsadd{Throat}\dictentry{Throat}{The Throat is the narrowest cross-section of a rocket nozzle, located between the convergent and divergent sections. All propellant combustion products must pass through this constricted area, where the flow reaches Mach 1 and becomes choked. The throat experiences the highest heat flux in the entire engine and requires intensive cooling.}{catRM}
\glsadd{Expansion-Ratio}\dictentry{Expansion Ratio}{The Expansion Ratio is the ratio of the nozzle exit area to the throat area, and it governs how much the exhaust gas is expanded. A higher expansion ratio extracts more energy from the gas, increasing specific impulse in vacuum but reducing performance at sea level. Sea-level engines typically have expansion ratios of 10 to 20, while vacuum-optimized engines may exceed 200. This parameter is a key trade-off in multi-stage rocket design.}{catRM}
\glsadd{Nozzle-Contour}\dictentry{Nozzle Contour}{The Nozzle Contour is the geometric profile of the nozzle wall from throat to exit. This shape determines how exhaust gases expand and directly affects thrust efficiency and flow uniformity. A well-designed contour minimizes flow separation and maximizes the axial component of thrust. Various contour shapes such as conical, bell, and parabolic offer different trade-offs between length, weight, and performance.}{catRM}
\glsadd{Rao-Nozzle}\dictentry{Rao Nozzle}{The Rao Nozzle is a bell-shaped nozzle contour optimized to achieve maximum thrust for a given nozzle length. Developed by G.V.Rao, it uses a 15-degree half-angle conical section followed by a curved bell that redirects flow axially. This design recovers most of the performance of an ideal parabolic nozzle in a shorter, lighter package. Rao contours are the standard for modern high-performance rocket engines.}{catRM}
\glsadd{Bell-Nozzle}\dictentry{Bell Nozzle}{A Bell Nozzle is a curved nozzle contour designed to turn exhaust flow back toward the axial direction after the initial expansion. The bell shape provides higher thrust efficiency than a simple conical nozzle while remaining shorter than an ideal parabolic contour. Bell nozzles are the most common type used on operational rocket engines.}{catRM}
\glsadd{Plug-Nozzle}\dictentry{Plug Nozzle}{A Plug Nozzle is an altitude-compensating nozzle design where the exhaust expands against an external plug surface rather than an outer wall. As ambient pressure changes with altitude, the flow automatically adjusts its expansion, maintaining near-optimal performance across a range of conditions. This self-compensation eliminates the fixed-geometry penalty of conventional nozzles.}{catRM}
\glsadd{Aerospike-Nozzle}\dictentry{Aerospike Nozzle}{The Aerospike Nozzle is an altitude-compensating design that uses an external aerodynamic surface to confine and direct the exhaust. It provides efficient thrust from sea level through vacuum by allowing ambient atmosphere to act as the outer nozzle boundary. The aerospike can be shorter and lighter than an equivalent bell nozzle while delivering consistent performance.}{catRM}
\glsadd{Regenerative-Cooling}\dictentry{Regenerative Cooling}{Regenerative Cooling routes one propellant, typically the fuel, through channels in the chamber and nozzle walls before injecting it into the combustion chamber. The propellant absorbs heat from the hot gas side and is simultaneously preheated, improving overall efficiency. This technique enables sustained operation at extreme chamber pressures.}{catRM}
\glsadd{Film-Cooling}\dictentry{Film Cooling}{Film Cooling injects a thin layer of cooler fluid along the inner wall of the combustion chamber. This film acts as a thermal barrier between the hot combustion gases and the chamber wall. It is typically applied near the injector face where heat flux is highest. While film cooling protects the structure, it slightly reduces combustion efficiency since the film does not fully participate in the primary combustion.}{catRM}
\glsadd{Ablative-Cooling}\dictentry{Ablative Cooling}{Ablative Cooling protects a rocket combustion chamber by lining the inner wall with a material that chars, melts, or sublimes as heat is absorbed. The phase change and mass loss carry energy away from the structure, keeping wall temperature within safe limits. Ablative chambers are simpler and lighter than regeneratively cooled designs but cannot sustain operation for extended durations.}{catRM}
\glsadd{Radiative-Cooling}\dictentry{Radiative Cooling}{Radiative Cooling radiates absorbed heat directly from the engine structure into space. This method is most effective in vacuum and for nozzle sections not directly cooled by propellant flow. Materials with high emissivity and melting points, such as niobium alloys, are used for radiatively cooled nozzle extensions. The RL10 engine nozzle extension is a notable example of this approach.}{catRM}
\glsadd{Igniter}\dictentry{Igniter}{An Igniter is a device that initiates combustion of propellants in a rocket engine combustion chamber. It must deliver sufficient energy to raise the propellant mixture above its auto-ignition temperature. Igniter types range from simple pyrotechnic charges to electrical spark systems and catalytic devices. Reliable ignition is critical for mission success, particularly for engines that must restart in orbit.}{catRM}
\glsadd{Pyrotechnic-Igniter}\dictentry{Pyrotechnic Igniter}{A Pyrotechnic Igniter uses a combustible chemical charge to generate hot gas and flame for initiating engine combustion. The charge is typically activated electrically and burns briefly to provide the necessary heat and radical species. Pyrotechnic igniters are simple, reliable, and widely used in both solid and liquid rocket motors. Their limitation is that they are single-use devices and cannot provide repeated ignition for restart missions.}{catRM}
\glsadd{Spark-Igniter}\dictentry{Spark Igniter}{A Spark Igniter uses an electric spark to ignite propellants in a rocket combustion chamber. The spark provides localized energy input that raises the local mixture above its ignition temperature. Spark igniters can be cycled multiple times, making them ideal for engines requiring in-flight restart capability. They are commonly used on upper-stage engines such as the RL10 and Vacuum Merlin.}{catRM}
\glsadd{Catalytic-Igniter}\dictentry{Catalytic Igniter}{A Catalytic Igniter uses a catalytic material to initiate decomposition or reaction of propellants without an external energy source. For monopropellant engines, passing hydrazine over a heated catalyst bed starts the exothermic decomposition reaction. Catalytic igniters are inherently simple and reliable because they have no moving parts or electrical requirements. They are the standard ignition method for spacecraft reaction control thrusters.}{catRM}
\glsadd{Thrust-Chamber}\dictentry{Thrust Chamber}{The Thrust Chamber is the integrated assembly comprising the combustion chamber, injector, and nozzle of a rocket engine. It is the core hardware where propellants are converted into high-velocity exhaust and thrust. The thrust chamber must withstand extreme temperatures, pressures, and vibration loads simultaneously. Its design is the primary determinant of engine performance, weight, and reliability.}{catRM}
\glsadd{Engine-Cycle}\dictentry{Engine Cycle}{The Engine Cycle describes how a liquid rocket engine delivers propellants from the tanks to the combustion chamber and manages the energy needed to drive the turbopumps. Different cycles trade complexity, performance, and reliability in distinct ways. Common cycles include pressure fed, gas generator, staged combustion, expander, and electric pump. The choice of cycle is one of the most consequential decisions in engine design.}{catRM}
\glsadd{Pressure-Fed-Cycle}\dictentry{Pressure Fed Cycle}{A Pressure Fed Cycle propels propellants to the combustion chamber by pressurizing the tanks with an inert gas such as helium. This eliminates the need for turbopumps and their associated plumbing, greatly simplifying the engine. The trade-off is that tanks must be heavier to withstand internal pressure, and chamber pressure is limited. Pressure fed engines are used in spacecraft propulsion systems and smaller launch vehicle upper stages.}{catRM}
\glsadd{Gas-Generator-Cycle}\dictentry{Gas Generator Cycle}{A Gas Generator Cycle burns a small portion of propellant in a separate preburner to generate hot gas that drives the turbopump turbine. The turbine exhaust is dumped overboard or into the nozzle at low pressure, representing a loss of potential performance. This cycle offers relative simplicity and is the basis for many engines including the Merlin and RS-27. It provides a practical balance of performance and reliability.}{catRM}
\glsadd{Staged-Combustion-Cycle}\dictentry{Staged Combustion Cycle}{A Staged Combustion Cycle burns propellant in a preburner at fuel-rich or oxidizer-rich conditions, and the gas drives the turbopump before entering the main combustion chamber for full burning. This recovers the energy of the turbine exhaust, yielding higher performance than a gas generator cycle. The complexity and high pressures involved make this cycle more challenging to develop.}{catRM}
\glsadd{Full-Flow-Staged-Combustion}\dictentry{Full Flow Staged Combustion}{Full Flow Staged Combustion uses two preburners, one fuel-rich and one oxidizer-rich, to drive the fuel and oxidizer turbopumps respectively. Every propellant molecule passes through a turbine before entering the main chamber, maximizing energy extraction. This cycle achieves the highest theoretical performance for a chemical rocket engine and reduces turbine temperature requirements. The SpaceX Raptor is the first operational full-flow engine.}{catRM}
\glsadd{Expander-Cycle}\dictentry{Expander Cycle}{An Expander Cycle uses heat absorbed by the propellant from the combustion chamber and nozzle walls to drive the turbopump turbine. The warmed propellant is then injected into the combustion chamber for final burning. This cycle provides high efficiency with fewer moving parts than staged combustion. It is limited in thrust by the amount of heat that can be absorbed, making it ideal for upper-stage engines.}{catRM}
\glsadd{Tap-Off-Cycle}\dictentry{Tap-Off Cycle}{A Tap-Off Cycle extracts hot combustion gas directly from the main combustion chamber to drive the turbopump turbine, then exhausts it overboard. This eliminates the need for a separate preburner, reducing engine complexity and part count. The tapped gas is drawn from a specific region of the chamber where conditions are suitable for turbine operation. The J-2 engine used on the Saturn V upper stages was a notable tap-off cycle engine.}{catRM}
\glsadd{Electric-Pump-Cycle}\dictentry{Electric Pump Cycle}{An Electric Pump Cycle uses electric motors and batteries to drive the propellant turbopumps instead of a gas-driven turbine. This eliminates the preburner, turbine, and hot gas plumbing entirely. It provides independent pump control and enables restart without a separate ignition system. The trade-off is battery mass, limiting this cycle to smaller engines. The Rocket Lab Rutherford pioneered this cycle for orbital use.}{catRM}
\glsadd{Turbopump}\dictentry{Turbopump}{A Turbopump is a high-speed rotating machine comprising a turbine and one or more centrifugal pumps that deliver propellants from the tanks to the combustion chamber at high pressure. It allows engines to operate at high chamber pressures, directly improving performance. Turbopumps operate at tens of thousands of revolutions per minute under extreme conditions. Their reliability is often the single greatest factor in overall engine success.}{catRM}
\glsadd{Preburner}\dictentry{Preburner}{A Preburner is a small combustion chamber that burns propellant at either fuel-rich or oxidizer-rich conditions to produce hot gas for driving the turbopump turbine. In a gas generator cycle, the preburner is the gas generator. In a staged combustion cycle, the preburner operates at much higher pressures and its exhaust is routed to the main chamber.}{catRM}
\glsadd{Main-Combustion-Chamber}\dictentry{Main Combustion Chamber}{The Main Combustion Chamber is the primary volume in a staged combustion or full-flow engine where the preburner exhausts mix with remaining propellants for complete combustion. It operates at the highest pressure in the engine system, often exceeding 200 bar in modern engines. The main chamber is where the majority of thrust-producing energy is released.}{catRM}
\glsadd{Gimbal}\dictentry{Gimbal}{A Gimbal is a swivel mount that allows a rocket engine to pivot and change the direction of its thrust vector. By tilting the engine on two axes, the vehicle can be steered without aerodynamic control surfaces. Gimbals are actuated hydraulically, pneumatically, or electromechanically in response to guidance commands. They are the primary method of thrust vector control for most liquid-propellant launch vehicles.}{catRM}
\glsadd{Vernier-Engine}\dictentry{Vernier Engine}{A Vernier Engine is a small auxiliary rocket engine used for fine trajectory adjustments and precise orbit insertion. It produces significantly lower thrust than the main engine, allowing delicate control of velocity increments. Vernier engines may be integral to the main engine or mounted separately on the vehicle. They are particularly valued in upper stages and spacecraft where precision orbit delivery is critical.}{catRM}
\glsadd{Reaction-Control-System}\dictentry{Reaction Control System}{A reaction control system is a network of small thrusters and propellant tanks controlling spacecraft attitude and small translational maneuvers. RCS thrusters are arranged around the spacecraft to provide forces and torques in all directions. They use hypergolic propellants that ignite on contact for reliable instantaneous response. The RCS is critical for launch vehicle control, orbital rendezvous, docking, and reentry attitude maintenance.}{catRM}
\glsadd{Cold-Gas-Thruster}\dictentry{Cold Gas Thruster}{A Cold Gas Thruster produces thrust by releasing compressed gas through a simple nozzle, without combustion. Common propellants include nitrogen, helium, and refrigerants stored at high pressure. Cold gas systems are the simplest and safest form of propulsion, with no ignition or combustion hardware required. They are used for attitude control on small satellites andCubeSats where thrust requirements are minimal.}{catRM}
\glsadd{Expendable-Launch-Vehicle}\dictentry{Expendable Launch Vehicle}{An Expendable Launch Vehicle is a rocket designed for a single flight, with all stages consumed or discarded during the mission. This approach simplifies design since components do not need to survive re-entry or be refurbished. Expendable vehicles dominated the launch industry for decades and remain in use for high-energy missions. However, their per-launch cost is higher since each flight requires an entirely new vehicle.}{catRM}
\glsadd{Multi-Stage-Rocket}\dictentry{Multi-Stage Rocket}{A Multi-Stage Rocket is a launch vehicle composed of two or more stages that are sequentially discarded as their propellant is exhausted. Each stage ignites after the previous one separates, shedding dead weight and continuing to accelerate efficiently. This staging approach dramatically improves payload capability compared to a single-stage design.}{catRM}
\glsadd{First-Stage}\dictentry{First Stage}{The First Stage is the lowermost stage of a multi-stage rocket, providing the initial thrust for liftoff and atmospheric ascent. It is typically the largest and most powerful stage, generating millions of pounds of thrust to overcome Earth's gravity and atmospheric drag. First stages burn dense propellants such as RP-1 or methane and may be recovered for reuse. They operate from sea level through the thickest portion of the atmosphere.}{catRM}
\glsadd{Second-Stage}\dictentry{Second Stage}{The Second Stage ignites after the first stage separates, continuing to accelerate the vehicle toward orbital velocity. It operates at higher altitudes where atmospheric pressure is reduced, allowing a more efficient nozzle expansion ratio. The second stage is typically smaller and lighter than the first stage since it carries less propellant and operates against reduced aerodynamic drag. It often performs the primary orbital insertion burn.}{catRM}
\glsadd{Strap-On-Booster}\dictentry{Strap-On Booster}{A Strap-On Booster is an auxiliary rocket motor attached to the side of a launch vehicle's first stage to provide additional liftoff thrust. Boosters are typically jettisoned before the core stage reaches orbit, reducing dead weight. They allow a single core vehicle to serve a wide range of payload requirements by adding or removing boosters. Common examples include the Space Shuttle SRBs and the boosters used on Ariane 5.}{catRM}
\glsadd{Payload-Fairing}\dictentry{Payload Fairing}{The Payload Fairing is the aerodynamic nose cone structure that encloses and protects the payload during ascent through the atmosphere. It shields the payload from aerodynamic heating, acoustic loads, and dynamic pressures experienced during the first minutes of flight. The fairing is typically constructed from lightweight composite materials and is jettisoned once the vehicle reaches sufficient altitude where aerodynamic forces are negligible.}{catRM}
\glsadd{Interstage}\dictentry{Interstage}{The Interstage is the structural section connecting two stages of a multi-stage rocket. It must withstand the compressive loads from the upper stage engine thrust while maintaining alignment between stages. Interstages are typically lightweight shell structures made from aluminum or composite materials. In some designs the interstage houses avionics, separation mechanisms, or even the upper stage engine itself in a hot-staging configuration.}{catRM}
\glsadd{Stage-Separation}\dictentry{Stage Separation}{Stage Separation is the process of disconnecting and distancing an expended rocket stage from the vehicle during flight. It involves pyrotechnic bolt severance, spring or thruster-powered push-off, and positive ignition of the next stage engine. The separation event must be precisely timed and occur with minimal disturbance to the remaining vehicle. Improper separation can cause tumbling, collision, or loss of mission control.}{catRM}
\glsadd{Fairing-Separation}\dictentry{Fairing Separation}{Fairing Separation is the jettisoning of the payload fairing once the vehicle has climbed above the dense atmosphere. The fairing is split into halves or petals and pushed away by springs or small thrusters. Fairing separation typically occurs between 100 and 200 kilometers altitude. Separating the fairing as early as possible reduces the dead weight carried during the remaining ascent to orbit.}{catRM}
\glsadd{Launch-Pad}\dictentry{Launch Pad}{A launch pad is the platform structure from which a rocket or spacecraft is launched. It provides support for the vehicle during pre-launch operations, including fueling, checkout, and crew access. The pad includes hold-down mechanisms, deluge systems for sound suppression, and lightning protection. After ignition, the launch pad must safely manage the extreme acoustic, thermal, and vibration loads produced by the rocket engines.}{catRM}
\glsadd{Launch-Tower}\dictentry{Launch Tower}{A Launch Tower is a vertical structure at the launch pad that supports the rocket before liftoff and provides access for crew, propellant, and umbilical connections. It houses swing arms that retract at launch and may incorporate a crew escape system for crewed missions. The tower must be designed to survive the blast environment of a pad abort. It is a visible and iconic element of the launch complex.}{catRM}
\glsadd{Launch-Rail}\dictentry{Launch Rail}{A Launch Rail is a guided track that directs the rocket during the initial moments of flight until it achieves sufficient velocity for aerodynamic stability. It is commonly used for sounding rockets and small launch vehicles that lack active thrust vector control at liftoff. The rail ensures the vehicle departs at the correct orientation and prevents tip-off during the critical low-speed phase.}{catRM}
\glsadd{Launch-Tube}\dictentry{Launch Tube}{A Launch Tube is an enclosed duct through which a missile or rocket is expelled using its own propulsion or an auxiliary launch charge. The tube protects the missile during initial acceleration and provides directional guidance until it clears the launcher. Launch tubes are commonly used for submarine-launched ballistic missiles and man-portable anti-tank weapons.}{catRM}
\glsadd{Mobile-Launcher}\dictentry{Mobile Launcher}{A Mobile Launcher is a transportable platform that can carry, fuel, and launch a rocket from different locations. It provides strategic flexibility by avoiding the need for fixed infrastructure. Mobile launchers are used extensively for military ballistic missiles and are being adopted for commercial small launch vehicles. They require self-contained propellant loading, power, and communication systems to operate independently at remote sites.}{catRM}
\glsadd{Silo}\dictentry{Silo}{A Silo is a hardened underground facility designed to store and launch ballistic missiles. It protects the missile from the effects of a nearby nuclear detonation and conceals its readiness state. Silo-based ICBMs can be launched within minutes of receiving a launch order. While mobile launchers have supplanted silos in some missile forces, silos remain a cornerstone of nuclear deterrence for multiple nations.}{catRM}
\glsadd{Countdown}\dictentry{Countdown}{A countdown is the structured timeline of events leading up to a rocket launch, typically measured in hours and minutes before liftoff. It includes milestones such as propellant loading, system checkouts, weather assessments, and final go/no-go decisions. The countdown synchronizes the activities of hundreds of personnel and dozens of systems. If a critical anomaly is detected, the countdown is held or scrubbed, and the sequence may be reset.}{catRM}
\glsadd{T-0}\dictentry{T-0}{T-0 is the designated moment of engine ignition or vehicle liftoff in a launch countdown. All time references before this point are negative (T-minus) and after are positive (T-plus). T-0 is precisely calculated based on orbital mechanics requirements, launch window constraints, and the vehicle's ascent profile. The accuracy of timing at T-0 is critical to achieving the correct orbital insertion point.}{catRM}
\glsadd{Lift-Off}\dictentry{Lift-Off}{Lift-Off is the moment when the rocket leaves the launch pad under its own power. It occurs shortly after engine ignition once thrust exceeds the vehicle weight and the hold-down clamps release. Lift-off marks the transition from ground operations to flight. The initial moments after lift-off are among the most dynamic, as the vehicle accelerates through the dense atmosphere at increasing speed.}{catRM}
\glsadd{Max-Q}\dictentry{Max Q}{Max Q is the point during ascent where the rocket experiences the maximum aerodynamic dynamic pressure. It occurs when the product of air density and velocity squared reaches its peak, typically at an altitude of 10 to 15 kilometers. The vehicle structure must be designed to withstand the peak loads at Max Q, which is often a critical design condition. Engines are sometimes throttled briefly to reduce loads at this point.}{catRM}
\glsadd{MECO}\dictentry{MECO}{MECO, or Main Engine Cutoff, is the event where the primary propulsion engines shut down after completing their burn. It typically marks the end of the first-stage or core-stage powered flight phase. MECO is followed by stage separation and, in most multi-stage vehicles, ignition of the next stage engine. The timing of MECO directly determines the velocity and trajectory state passed to the subsequent stage.}{catRM}
\glsadd{SECO}\dictentry{SECO}{SECO, or Second Engine Cutoff, marks the shutdown of the second-stage engine after it has completed its burn. At this point the vehicle is typically in or near its target orbit. SECO may be followed by coast phases and subsequent upper-stage burns for orbit circularization or transfer trajectory insertion. The precision of SECO timing is critical to achieving the desired orbital parameters.}{catRM}
\glsadd{Insertion}\dictentry{Insertion}{Insertion is the process of placing a spacecraft or payload into its intended orbit around Earth or another celestial body. It involves precise timing, direction, and magnitude of the final engine burn. Orbital insertion requires matching the target orbit's altitude, inclination, and eccentricity within tight tolerances. Successful insertion is the culminating objective of the launch vehicle's mission.}{catRM}
\glsadd{Injection-Burn}\dictentry{Injection Burn}{An Injection Burn is an engine firing that places a spacecraft onto a transfer trajectory, such as a Hohmann transfer orbit to another planet. The burn is calculated to change the spacecraft's velocity so that its new orbit intersects the target body's orbit at the correct time. Injection burns require precise navigation and timing to achieve the desired encounter conditions.}{catRM}
\glsadd{Circularization-Burn}\dictentry{Circularization Burn}{A Circularization Burn is an engine firing performed to transform an elliptical orbit into a circular one. It is typically executed at apoapsis, the highest point of the orbit, to raise the periapsis altitude. This maneuver is essential for operational satellites that require a constant altitude for communications or Earth observation. The magnitude of the burn is determined by the difference between the current and target orbit velocities.}{catRM}
\glsadd{Deorbit-Burn}\dictentry{Deorbit Burn}{A Deorbit Burn is an engine firing performed in the direction opposite to the spacecraft's velocity to reduce its orbital speed. This lowers the periapsis into the atmosphere, causing drag to eventually remove the spacecraft from orbit. The burn is timed so that the resulting entry interface occurs at the desired landing site. Deorbit burns are a standard procedure for spacecraft returning from low Earth orbit.}{catRM}
\glsadd{Retro-Rocket}\dictentry{Retro Rocket}{A Retro Rocket is a rocket motor that fires in the direction opposite to the vehicle's motion, producing a decelerating force. Retro rockets are used for deorbit maneuvers, landing, and trajectory corrections that reduce velocity. They must be oriented to produce thrust against the direction of travel. Retro rockets were famously used on the Apollo lunar module for descent and on Soyuz for Earth reentry braking.}{catRM}
\glsadd{Ullage-Motor}\dictentry{Ullage Motor}{An Ullage Motor is a small solid or cold-gas rocket that fires briefly before the main engine ignites in zero gravity. Its purpose is to settle the liquid propellants at the bottom of their tanks by creating a small acceleration. Without ullage settling, propellants float freely in microgravity and cannot be reliably drawn to the engine inlets. Ullage motors are essential for upper-stage restarts during coast phases.}{catRM}
\glsadd{Spin-Motor}\dictentry{Spin Motor}{A Spin Motor is a small rocket motor mounted tangentially to a rocket stage to induce rotation about its longitudinal axis. The spin provides gyroscopic stability, keeping the vehicle pointed in the correct direction without active attitude control. Spin stabilization was widely used in early sounding rockets and solid-fuel ballistic missiles. It remains relevant for some small launch vehicles and spinning upper stages.}{catRM}
\glsadd{Separation-Motor}\dictentry{Separation Motor}{A Separation Motor is a small solid rocket motor used to push rocket stages apart during stage separation. It provides a positive, reliable separation force independent of spring or pneumatic systems. Separation motors are typically mounted on the discarded stage and fire outward to ensure safe distancing. They are simple, reliable, and widely used across multiple launch vehicle families.}{catRM}
\glsadd{IRFNA}\dictentry{IRFNA}{IRFNA, or Inhibited Red Fuming Nitric Acid, is a storable liquid oxidizer used in hypergolic rocket propulsion systems. The inhibition refers to the addition of fluoride compounds that reduce the oxidizer's corrosiveness to metals. IRFNA ignites spontaneously when it contacts fuel, providing reliable hypergolic ignition. It is used in several national ballistic missile and launch vehicle programs where storable propellants are required.}{catRM}
\glsadd{LOX}\dictentry{LOX}{LOX, or Liquid Oxygen, is the most widely used oxidizer in modern high-performance rocket engines. It is produced by liquefying and fractionally distilling atmospheric air at temperatures below minus 183 degrees Celsius. LOX is dense, non-toxic, and inexpensive, making it the oxidizer of choice for launch vehicle first and upper stages. It is typically paired with liquid hydrogen, RP-1, or methane.}{catRM}
\glsadd{LH2}\dictentry{LH2}{LH2, or Liquid Hydrogen, is a cryogenic fuel stored at minus 253 degrees Celsius that provides the highest specific impulse of any practical chemical rocket propellant. It is colorless, odorless, and has an extremely low density, requiring large insulated tanks. LH2 paired with LOX is the standard propellant for upper stages on vehicles such as the Delta IV, Atlas V, and Space Launch System.}{catRM}
\glsadd{LNG}\dictentry{LNG}{LNG, or Liquefied Natural Gas, is a rocket fuel composed primarily of methane liquefied at minus 162 degrees Celsius. LNG offers a middle ground between RP-1 and liquid hydrogen in terms of performance and handling complexity. Methane burns cleanly, reducing coking and enabling simpler engine reusability.}{catRM}
\glsadd{Propellant-Mass-Fraction}\dictentry{Propellant Mass Fraction}{Propellant Mass Fraction is the ratio of propellant mass to the total initial mass of a rocket stage. It is a fundamental measure of how mass-efficient a stage design is, with values typically ranging from 0.85 to 0.95 for well-designed stages. A higher propellant mass fraction means less structural overhead and more delta-v available for the mission. Improving this fraction is one of the primary challenges of rocket structural design.}{catRM}
\glsadd{Dry-Mass}\dictentry{Dry Mass}{Dry Mass is the total mass of a rocket stage or vehicle excluding all propellant. It includes the structure, engines, avionics, wiring, plumbing, and any payload. Minimizing dry mass while maintaining structural integrity is the central challenge of aerospace structural engineering. Even small reductions in dry mass translate directly into significant increases in payload capacity or mission delta-v.}{catRM}
\glsadd{Wet-Mass}\dictentry{Wet Mass}{Wet Mass is the total mass of a rocket stage or vehicle including all propellants. It represents the initial mass at engine ignition and is the starting point for Tsiolkovsky equation calculations. The difference between wet mass and dry mass determines the mass ratio, which directly governs the achievable delta-v. Wet mass is the primary input for trajectory analysis and structural load calculations.}{catRM}
\glsadd{Payload-Mass-Fraction}\dictentry{Payload Mass Fraction}{Payload Mass Fraction is the ratio of the payload mass to the total initial mass of a launch vehicle. It represents the fraction of the total vehicle mass that is actually useful payload delivered to orbit. Typical payload mass fractions range from one to five percent for chemical launch vehicles, illustrating the extreme mass penalty of space launch. Maximizing this fraction is the ultimate objective of all launch vehicle design optimization.}{catRM}
\glsadd{Tsiolkovsky-Equation}\dictentry{Tsiolkovsky Equation}{The Tsiolkovsky Equation, also known as the Rocket Equation, relates the achievable change in velocity to the exhaust velocity and mass ratio. Delta-v equals exhaust velocity multiplied by the natural logarithm of initial mass over final mass. It reveals the exponential relationship between payload fraction and velocity change, explaining why staging is necessary for orbital flight.}{catRM}
\glsadd{Characteristic-Velocity}\dictentry{Characteristic Velocity}{Characteristic Velocity (c*) is a figure of merit describing the performance of a rocket propellant combination inside the combustion chamber, independent of nozzle geometry. It is calculated from the chamber pressure, mass flow rate, and throat area. Higher c* values indicate more energetic propellant combinations. Engineers use it to compare propellant performance and diagnose combustion quality.}{catRM}
\glsadd{Exhaust-Velocity}\dictentry{Exhaust Velocity}{Exhaust Velocity is the speed at which combustion gases exit the rocket nozzle, measured at the nozzle exit plane. It depends on chamber temperature, molecular weight of the exhaust, and the expansion ratio of the nozzle. Exhaust velocity directly determines the specific impulse and thus the fuel efficiency of the engine. Maximizing exhaust velocity is a central goal in nozzle design.}{catRM}
\glsadd{Effective-Exhaust-Velocity}\dictentry{Effective Exhaust Velocity}{Effective Exhaust Velocity accounts for the pressure imbalance at the nozzle exit by combining the momentum thrust and pressure thrust into a single equivalent velocity. It equals the actual exhaust velocity plus the exit pressure times exit area divided by mass flow rate. This quantity allows thrust to be expressed as mass flow rate times effective exhaust velocity, simplifying performance calculations for engines operating at varying altitudes.}{catRM}
\glsadd{Mass-Flow-Rate}\dictentry{Mass Flow Rate}{Mass Flow Rate is the mass of propellant consumed per unit time, typically expressed in kilograms per second. It is determined by the injector design, chamber pressure, and propellant density. Mass flow rate combined with exhaust velocity gives the thrust produced by the engine. Accurate control of mass flow rate is essential for maintaining stable and predictable engine performance.}{catRM}
\glsadd{Chamber-Pressure}\dictentry{Chamber Pressure}{Chamber Pressure is the internal pressure within the combustion chamber of a rocket engine during operation. Higher chamber pressure generally improves combustion efficiency and allows smaller nozzle throat dimensions. The choice of operating pressure affects engine weight, complexity, and the structural demands on the combustion chamber and turbopump systems.}{catRM}
\glsadd{Expansion-Ratio-Effect}\dictentry{Expansion Ratio Effect}{Expansion Ratio Effect refers to how the ratio of nozzle exit area to throat area influences engine performance at different altitudes. A larger expansion ratio extracts more energy from the exhaust at high altitude but causes flow separation at low altitude. This trade-off drives the design of altitude-compensating nozzles and determines whether an engine is optimized for sea level or vacuum operation.}{catRM}
\glsadd{Optimal-Expansion}\dictentry{Optimal Expansion}{Optimal Expansion occurs when the pressure of the exhaust gases at the nozzle exit exactly matches the surrounding ambient pressure. Under this condition all available thrust is directed axially and no energy is wasted on expansion or compression of the exhaust plume. No single fixed nozzle can achieve optimal expansion at all altitudes because ambient pressure changes during ascent, motivating altitude-compensating nozzle designs.}{catRM}
\glsadd{Overexpansion}\dictentry{Overexpansion}{Overexpansion happens when the nozzle exit pressure falls below the ambient pressure, causing the exhaust plume to be compressed by the atmosphere after leaving the nozzle. This reduces thrust efficiency and can cause flow separation inside the nozzle at severe levels. Overexpansion is common for sea-level engines at altitude and for upper-stage engines operating briefly at low altitude during launch.}{catRM}
\glsadd{Underexpansion}\dictentry{Underexpansion}{Underexpansion occurs when the nozzle exit pressure exceeds ambient pressure, meaning the exhaust gases continue to expand after leaving the nozzle. The resulting outward-spreading plume wastes energy that could have been converted to additional thrust. Underexpansion is typical for sea-level engines because their nozzles are designed with modest expansion ratios to avoid overexpansion at low altitude.}{catRM}
\glsadd{Flow-Separation-in-Nozzle}\dictentry{Flow Separation in Nozzle}{Flow Separation in Nozzle is the detachment of the exhaust gas flow from the inner wall of a convergent-divergent nozzle. It occurs when the nozzle is overexpanded beyond a critical pressure ratio, creating localized shock waves that push the boundary layer off the wall. Separation causes asymmetric side loads and reduced performance. Managing separation is critical for nozzle structural integrity and predictable thrust.}{catRM}
\glsadd{Side-Load}\dictentry{Side Load}{Side Load refers to lateral forces exerted on a rocket engine nozzle during operation, primarily caused by asymmetric flow separation. These forces can exceed the primary axial thrust load during certain transient conditions such as engine start and shutdown. Side loads are a major design driver for nozzle structural margins and can cause catastrophic nozzle failure if not adequately accounted for.}{catRM}
\glsadd{Nozzle-Extension}\dictentry{Nozzle Extension}{Nozzle Extension is an add-on section attached to the divergent portion of a nozzle to increase the expansion ratio for high-altitude or vacuum operation. It allows a sea-level engine to achieve better performance when operating in vacuum conditions. Some designs use extendable nozzles that deploy after launch, combining sea-level efficiency with optimized vacuum thrust. The extension must handle extreme thermal and structural loads during firing.}{catRM}
\glsadd{Dual-Bell-Nozzle}\dictentry{Dual Bell Nozzle}{Dual Bell Nozzle is an altitude-compensating nozzle with two distinct bell contours separated by a pressure break. At low altitude the exhaust flows in the smaller lower-expansion contour, while at high altitude pressure buildup causes the flow to attach to the larger upper contour. This provides near-optimal expansion across a wide altitude range without moving parts. It offers a passive alternative to complex adjustable nozzles.}{catRM}
\glsadd{Thrust-Chamber-Liner}\dictentry{Thrust Chamber Liner}{Thrust Chamber Liner is a thin wall element placed inside the thrust chamber to protect the outer structural wall from the extreme heat of combustion. It is typically made from regeneratively cooled channels or ablative materials. The liner absorbs thermal loads and extends the operational life of the chamber. Its design directly affects engine reusability and maintenance intervals.}{catRM}
\glsadd{Combustion-Stability}\dictentry{Combustion Stability}{Combustion Stability describes the condition in which the combustion process inside a rocket engine proceeds without sustained pressure oscillations. A stable engine maintains steady chamber pressure and uniform heat release over time. Achieving stability requires careful design of the injector, chamber geometry, and propellant feed system. Stability is typically validated through extensive hot-fire testing.}{catRM}
\glsadd{Combustion-Instability}\dictentry{Combustion Instability}{Combustion Instability refers to self-sustained pressure oscillations in the combustion chamber driven by coupling between heat release, acoustics, and propellant flow. Instabilities can cause destructive vibrations, localized hot spots, and engine failure within milliseconds. They are classified by frequency as low-frequency (chugging), intermediate-frequency (buzzing), or high-frequency (screaming) modes.}{catRM}
\glsadd{Chugging}\dictentry{Chugging}{Chugging is a low-frequency combustion instability, typically below 300 Hz, caused by coupling between the propellant feed system and the combustion process. The oscillation creates visible pulsation of the engine thrust and exhaust plume. Chugging is usually controlled by increasing injector pressure drop or modifying feed system compliance. It is the most common instability encountered during engine development.}{catRM}
\glsadd{Buzzing}\dictentry{Buzzing}{Buzzing is an intermediate-frequency combustion instability occurring in the range of roughly 300 to 1000 Hz. It arises from acoustic resonance in the combustion chamber coupled with the combustion process. Buzzing can produce significant structural vibration and is more difficult to diagnose than chugging due to its complex wave patterns. Suppression typically requires baffle modifications or acoustic damping.}{catRM}
\glsadd{Screaming}\dictentry{Screaming}{Screaming is a high-frequency combustion instability, generally above 1000 Hz, involving transverse acoustic modes in the combustion chamber. It produces extremely intense pressure oscillations that can destroy an engine in fractions of a second. Screaming instabilities are among the most dangerous and difficult to predict, often requiring extensive high-frequency instrumentation and iterative design to resolve.}{catRM}
\glsadd{Pogo-Oscillation}\dictentry{Pogo Oscillation}{Pogo Oscillation is a longitudinal vibration of the entire rocket caused by feedback between engine thrust fluctuations and the structural dynamics of the vehicle. As the rocket oscillates, propellant flow varies, which changes thrust, which reinforces the oscillation. Pogo can damage the vehicle structure and payload. Suppression devices such as accumulators are installed in propellant lines to break the feedback loop.}{catRM}
\glsadd{Acoustic-Cavity}\dictentry{Acoustic Cavity}{Acoustic Cavity is a tuned resonant chamber incorporated into the combustion chamber or injector to absorb acoustic energy at specific frequencies. It works by creating destructive interference with the pressure waves driving instability. Acoustic cavities are a passive suppression method requiring no moving parts. They must be precisely tuned to the problematic frequency modes identified during testing.}{catRM}
\glsadd{Baffle}\dictentry{Baffle}{Baffle is a physical barrier mounted on the face of an injector to disrupt tangential and radial acoustic modes in the combustion chamber. By breaking up the wave patterns, baffles prevent the constructive feedback that sustains high-frequency instabilities. Baffles must withstand the full thermal environment of the combustion zone. They are one of the most effective and widely used stability augmentation devices in rocket engine design.}{catRM}
\glsadd{Resonator}\dictentry{Resonator}{Resonator is a tuned acoustic device installed in or near the combustion chamber to damp pressure oscillations at targeted frequencies. Unlike baffles, resonators absorb energy rather than physically blocking wave propagation. They are typically quarter-wave or Helmholtz-type cavities designed to match the frequency of observed instability modes. Resonators add minimal weight and no moving parts, making them an attractive stability solution.}{catRM}
\glsadd{Injector-Pattern}\dictentry{Injector Pattern}{Injector Pattern is the geometric arrangement of fuel and oxidizer injection elements on the injector face. The pattern determines how propellants are distributed and mixed within the combustion chamber. Common patterns include self-impinging, unlike-doublet, and coaxial arrangements. The pattern directly affects combustion efficiency, chamber temperature uniformity, and stability characteristics.}{catRM}
\glsadd{Showerhead-Injector}\dictentry{Showerhead Injector}{Showerhead Injector is a simple injector design featuring a large number of small, uniformly spaced holes across the injector face. Each hole injects propellant as a single stream without intentional impingement. Mixing relies on turbulent diffusion in the combustion chamber rather than on interstream collision. The design is mechanically robust but typically produces lower combustion efficiency than more sophisticated injector types.}{catRM}
\glsadd{Impinging-Injector}\dictentry{Impinging Injector}{Impinging Injector creates mixing by directing propellant streams to collide with each other at predetermined angles. The impact atomizes the propellants into fine droplets and promotes rapid mixing. Common configurations include unlike doublets, triplet, and quintlet impinging patterns. This injector type provides good atomization and mixing but requires precise manufacturing to maintain correct impingement geometry.}{catRM}
\glsadd{Coaxial-Injector}\dictentry{Coaxial Injector}{Coaxial Injector places the fuel and oxidizer in concentric annular streams that mix through shear forces at their interface. This design is widely used in large engines because it scales well and provides good mixing at high flow rates. The shear between the two streams atomizes the propellants without requiring direct impingement. Coaxial injectors are favored for engines using liquid oxygen and hydrocarbon or hydrogen fuels.}{catRM}
\glsadd{Pintle-Injector}\dictentry{Pintle Injector}{Pintle Injector uses a central post or pintle to distribute one propellant radially while the other flows axially around it. The mixing occurs through the collision of these orthogonal streams. Pintle injectors allow deep throttling because the flow distribution remains stable across a wide range of operating conditions. This capability makes them ideal for landing engines and variable-thrust applications.}{catRM}
\glsadd{Shear-Injector}\dictentry{Shear Injector}{Shear Injector atomizes propellants primarily through the velocity difference between adjacent streams of fuel and oxidizer. The shear forces at the stream interface break the liquids into fine droplets without requiring direct collision between jets. This approach produces uniform atomization and reduces the risk of injector face erosion. Shear injectors are common in engines using cryogenic propellants with large velocity differentials.}{catRM}
\glsadd{Atomization}\dictentry{Atomization}{Atomization is the process of breaking bulk liquid propellant into a fine spray of small droplets. Smaller droplets have a larger surface-area-to-volume ratio, which dramatically increases the evaporation and combustion rate. Methods include pressure atomization through small orifices, impinging stream atomization, and shear atomization. Effective atomization is essential for achieving high combustion efficiency and stable operation.}{catRM}
\glsadd{Droplet-Combustion}\dictentry{Droplet Combustion}{Droplet Combustion is the burning of individual propellant droplets within the combustion chamber. Each droplet evaporates at its surface, and the resulting vapor reacts with the oxidizer in a diffusion flame around it. The combustion rate depends on droplet diameter, ambient temperature, and local oxidizer concentration. Understanding droplet combustion is key to modeling overall engine performance and predicting combustion efficiency.}{catRM}
\glsadd{Combustion-Efficiency}\dictentry{Combustion Efficiency}{Combustion Efficiency measures how completely the chemical energy of the propellants is released during combustion. It is defined as the ratio of actual heat release to the theoretical maximum based on propellant mass and heating value. Factors affecting efficiency include mixing quality, droplet size, and residence time in the chamber. Typical well-designed engines achieve combustion efficiencies above 95 percent.}{catRM}
\glsadd{C-Star-Efficiency}\dictentry{C Star Efficiency}{C Star Efficiency is a measure of combustion quality expressed as the ratio of the achieved characteristic velocity to the theoretical c* for the given propellant combination. It isolates the performance of the injector and chamber from the nozzle. A value of 100 percent represents ideal combustion. C star efficiency is a standard metric used during engine development testing to evaluate injector performance.}{catRM}
\glsadd{Thrust-Coefficient}\dictentry{Thrust Coefficient}{Thrust Coefficient is a dimensionless parameter that quantifies how effectively the nozzle converts chamber pressure into thrust. It depends on the ratio of specific heats of the exhaust gas, the nozzle expansion ratio, and the ambient pressure. A higher thrust coefficient means more thrust for a given chamber pressure and throat area. Engine designers optimize the nozzle contour to maximize this coefficient for the intended operating altitude.}{catRM}
\glsadd{Specific-Impulse-Efficiency}\dictentry{Specific Impulse Efficiency}{Specific Impulse Efficiency is the ratio of the actually achieved specific impulse to the theoretical maximum for the propellant combination and operating conditions. It captures all real-world losses including incomplete combustion, non-ideal expansion, and boundary layer effects. This metric provides a single number reflecting the overall engine performance quality. It is widely used in trade studies comparing different engine designs.}{catRM}
\glsadd{Grain}\dictentry{Grain}{Grain is a preformed mass of solid propellant cast into a specific geometric shape inside the rocket motor case. The grain's geometry determines the burning surface area profile and therefore the thrust-time curve of the motor. Grains are typically made from composite propellants containing an oxidizer, fuel binder, and metal additive. The grain design is the primary mechanism for tailoring the motor's performance characteristics.}{catRM}
\glsadd{Grain-Configuration}\dictentry{Grain Configuration}{Grain Configuration refers to the cross-sectional geometry of a solid propellant grain, including features such as bore shape, number of perforations, and slot geometry. The configuration controls how the burning surface evolves over time, which determines whether thrust is neutral, progressive, or regressive. Designers select configurations based on mission thrust requirements and structural constraints.}{catRM}
\glsadd{Burning-Surface}\dictentry{Burning Surface}{Burning Surface is the total area of exposed propellant that is actively combusting at any given instant. The burning surface area multiplied by the linear burn rate gives the volumetric burn rate, which determines the mass flow rate and thus the thrust. As the grain burns, the surface area changes depending on grain geometry. Controlling this evolution is the fundamental challenge of solid motor grain design.}{catRM}
\glsadd{Neutral-Burning}\dictentry{Neutral Burning}{Neutral Burning describes a grain design in which the total burning surface area remains approximately constant throughout the burn. This produces a nearly constant thrust level from ignition to burnout. Neutral grains are achieved by balancing the increase in surface area from inward-burning surfaces with the decrease from other surfaces. The cylindrical bore grain is the simplest example of neutral burning geometry.}{catRM}
\glsadd{Progressive-Burning}\dictentry{Progressive Burning}{Progressive Burning is a grain design where the burning surface area increases as the propellant burns, producing rising thrust over time. This profile is useful for applications requiring maximum thrust later in the burn. Progressive burning is typically achieved with multi-perforated grains or star-shaped bores that open up as burning proceeds. The increasing thrust must be accommodated by vehicle structural margins.}{catRM}
\glsadd{Regressive-Burning}\dictentry{Regressive Burning}{Regressive Burning is a grain configuration where the burning surface area decreases during the burn, producing declining thrust. This profile can be useful for matching the reducing aerodynamic loads as a vehicle ascends. Regressive burning is commonly achieved with external-burning tubular grains where the outer diameter shrinks over time. Care must be taken to ensure residual propellant does not create thrust tail-off problems.}{catRM}
\glsadd{End-Burning}\dictentry{End Burning}{End Burning is a grain configuration where combustion proceeds linearly from one face of the grain toward the other, like a cigarette. The burning surface area equals the cross-sectional area of the grain and remains constant, producing neutral thrust. This simple geometry maximizes the volumetric loading of propellant in the case. End-burning grains are used in sustain motors and gas generators where steady output is required.}{catRM}
\glsadd{Internal-Burning}\dictentry{Internal Burning}{Internal Burning describes a grain design where combustion occurs from passages or perforations inside the grain, with the outer surface bonded to the motor case. This approach protects the case from hot combustion gases and eliminates the need for internal insulation on bonded surfaces. Internal burning allows complex grain geometries to produce a wide range of thrust profiles. Most modern solid rocket motors use internal burning grains.}{catRM}
\glsadd{Cigarette-Burning}\dictentry{Cigarette Burning}{Cigarette Burning is another term for end burning of a solid propellant grain, where combustion progresses axially from one end to the other at a constant rate. The burning surface is simply the circular cross-section of the grain. This produces a neutral thrust profile with maximum propellant loading fraction. Cigarette-burning grains are preferred in applications where simplicity and long burn duration are priorities.}{catRM}
\glsadd{Star-Grain}\dictentry{Star Grain}{Star Grain is a solid propellant grain with a star-shaped internal bore cross-section. As burning proceeds the star points erode, increasing the burning surface area. This geometry can be tailored to produce progressive, neutral, or regressive thrust profiles depending on the star dimensions. Star grains are one of the most widely used configurations in tactical and strategic solid motors.}{catRM}
\glsadd{Wagon-Wheel-Grain}\dictentry{Wagon Wheel Grain}{Wagon Wheel Grain is a solid propellant grain with a cross-section resembling a wagon wheel, featuring a central bore connected to radial slots. The slots increase the initial burning surface area and can be shaped to control the thrust profile over time. This configuration provides designers with additional parameters to match specific thrust requirements. Wagon wheel grains are used in motors requiring complex thrust histories.}{catRM}
\glsadd{Slotted-Tube-Grain}\dictentry{Slotted Tube Grain}{Slotted Tube Grain is a tubular propellant grain with axial slots cut into its surface to increase and control the burning area. The slots allow designers to modify the thrust-time curve beyond what a simple tubular grain can achieve. As the slots burn outward, additional surface area is exposed, typically producing a progressive component. This grain type offers a practical compromise between geometric complexity and manufacturing simplicity.}{catRM}
\glsadd{Inhibitor}\dictentry{Inhibitor}{Inhibitor is a coating or liner applied to selected surfaces of a solid propellant grain to prevent those surfaces from burning. By controlling which surfaces burn, inhibitors allow designers to achieve the desired grain geometry and thrust profile. Common inhibitor materials include rubber-based compounds and reinforced polymers. The inhibitor must adhere reliably throughout storage and withstand the thermal environment during ignition.}{catRM}
\glsadd{Propellant-Liner}\dictentry{Propellant Liner}{Propellant Liner is a thin layer of compliant material applied between the solid propellant grain and the motor case to create a strong adhesive bond. The liner must accommodate differential thermal expansion between the propellant and case material. It also provides a chemical barrier preventing ingredient migration between propellant and case. Liner reliability is critical because bond failure can expose the case to combustion gases.}{catRM}
\glsadd{Case}\dictentry{Case}{Case is the structural pressure vessel that contains the solid propellant grain and withstands the high internal pressures generated during combustion. It is typically made from high-strength steel, titanium, or filament-wound composite materials. The case must be as lightweight as possible while maintaining adequate safety margins. Case design and material selection are major drivers of the motor's overall mass fraction and performance.}{catRM}
\glsadd{Motor-Case}\dictentry{Motor Case}{Motor Case is the outer shell of a solid rocket motor that serves as both the pressure vessel and the primary structure. It transfers aerodynamic and thrust loads to the vehicle airframe. The case is often the heaviest single component of the motor, so minimizing its mass is a key design objective. Modern composite cases wound from carbon or glass fiber offer significant weight savings over metallic designs.}{catRM}
\glsadd{Nozzle-Throat-Insert}\dictentry{Nozzle Throat Insert}{Nozzle Throat Insert is a component made from a high-temperature erosion-resistant material such as graphite or tungsten, installed at the narrowest point of the nozzle. The throat sees the most severe thermal and erosive environment in the engine. The insert maintains a consistent throat area throughout the burn, critical for stable thrust and chamber pressure. Throat area growth due to erosion is a key performance degradation mechanism.}{catRM}
\glsadd{Erosion}\dictentry{Erosion}{Erosion is the progressive loss of material from the nozzle throat and other hot-gas-exposed surfaces during motor firing. It is caused by the combined effects of extreme temperature, high-velocity gas flow, and chemical attack from combustion products. Erosion increases the throat area, which reduces chamber pressure and thrust over time. Controlling erosion is critical for maintaining predicted performance, especially in long-duration burns.}{catRM}
\glsadd{Graphite-Nozzle}\dictentry{Graphite Nozzle}{Graphite Nozzle uses graphite as the primary material for the throat and divergent sections due to its excellent high-temperature strength and thermal shock resistance. Graphite sublimes rather than melting, allowing it to withstand the extreme throat environment. However, graphite is susceptible to oxidative erosion from water vapor and combustion products. It is commonly used in solid rocket motor nozzles for tactical and strategic missiles.}{catRM}
\glsadd{Carbon-Carbon-Nozzle}\dictentry{Carbon-Carbon Nozzle}{Carbon-Carbon Nozzle is constructed from a carbon fiber reinforced carbon matrix composite. This material retains its strength at extremely high temperatures and has a low erosion rate, making it suitable for the most demanding nozzle applications. Carbon-carbon nozzles are used in solid rocket motors for intercontinental ballistic missiles and space launch systems. The material is expensive and requires specialized manufacturing processes.}{catRM}
\glsadd{Metal-Nozzle}\dictentry{Metal Nozzle}{Metal Nozzle is a nozzle constructed primarily from high-temperature metal alloys such as columbium, molybdenum, or nickel-based superalloys. Metals offer good machinability and ductility, which simplifies fabrication and provides damage tolerance. However, metals erode faster than carbon-based materials and require thermal protection coatings for long-duration firings. Metal nozzles are common in liquid rocket engines and shorter-burn solid motors.}{catRM}
\glsadd{Nozzle-Closure}\dictentry{Nozzle Closure}{Nozzle Closure is a protective seal installed in the nozzle exit of a solid rocket motor before ignition. It prevents moisture, contaminants, and foreign objects from entering the motor during storage and transport. The closure must fragment cleanly and be expelled during ignition without damaging the nozzle or creating debris. It is typically held in place by frangible mounts designed to break under initial ignition pressure.}{catRM}
\glsadd{Igniter-Grain}\dictentry{Igniter Grain}{Igniter Grain is a small solid propellant charge used to initiate combustion in the main motor grain. It generates hot gas and particles that heat the main grain surface above its ignition temperature. The igniter grain is itself ignited by an electric squib or pyrotechnic device. Its design must ensure uniform and rapid ignition of the full burning surface to avoid asymmetric pressurization.}{catRM}
\glsadd{Pyrogen-Igniter}\dictentry{Pyrogen Igniter}{Pyrogen Igniter is a pyrotechnic device that produces a large volume of hot gas and incandescent particles to ignite a rocket motor. It consists of a pyrotechnic composition housed in a small casing with an electrical initiator. The hot gas stream is directed onto the grain surface to ensure rapid and uniform ignition. Pyrogen igniters are used in both solid rocket motors and as igniters for liquid engine gas generators.}{catRM}
\glsadd{Safe-and-Arm-Device}\dictentry{Safe and Arm Device}{Safe and Arm Device is a mechanical or electro-mechanical safety mechanism that prevents inadvertent ignition of a rocket motor or missile. In the safe position the device creates a physical interruption in the ignition circuit. It is switched to the armed position only after specific authorization and arming conditions are met. This device is a critical safety element throughout manufacturing, transportation, storage, and pre-launch operations.}{catRM}
\glsadd{Thrust-Termination}\dictentry{Thrust Termination}{Thrust Termination is a system that provides a means to end the thrust of a solid rocket motor before the propellant is fully consumed. It typically works by opening ports in the forward end of the case to vent pressure and halt the combustion process. Thrust termination is essential for precise orbit insertion in upper-stage solid motors and for range safety compliance. The system must act reliably within milliseconds of the termination command.}{catRM}
\glsadd{Destruct-System}\dictentry{Destruct System}{Destruct System is an explosive device installed on a rocket or missile that fragments the vehicle when commanded by range safety officers. It is activated if the vehicle deviates from its planned trajectory and threatens populated areas. The system uses linear shaped charges to cut through the motor case and structural members. It is a last-resort safety measure that ensures public protection during launch operations.}{catRM}
\glsadd{Flight-Termination-System}\dictentry{Flight Termination System}{Flight Termination System is the integrated collection of sensors, receivers, command decoders, and explosive devices that enable range safety officers to destroy a malfunctioning vehicle in flight. It includes redundant command receivers and independent power supplies for reliability. The system meets stringent safety requirements set by launch range authorities. A properly functioning FTS ensures that no debris reaches inhabited areas.}{catRM}
\glsadd{Missile}\dictentry{Missile}{Missile is a self-propelled weapon system that carries a warhead to a designated target using rocket or jet propulsion and a guidance system. Missiles are classified by launch platform, target type, and range. They may use solid, liquid, or air-breathing propulsion depending on the mission requirements. Modern missiles integrate advanced guidance, navigation, and warhead technologies to achieve high accuracy and lethality.}{catRM}
\glsadd{Ballistic-Missile}\dictentry{Ballistic Missile}{Ballistic Missile is a missile that follows a predominantly unpowered ballistic trajectory after its initial boost phase. The warhead travels on a parabolic arc through the atmosphere or space before reentering and striking the target. The trajectory is determined at launch and cannot be significantly altered during the coast phase. Ballistic missiles range from short-range tactical weapons to intercontinental strategic systems.}{catRM}
\glsadd{Cruise-Missile}\dictentry{Cruise Missile}{Cruise Missile is a self-navigating missile that maintains powered flight throughout most of its trajectory, using wings and an air-breathing engine to fly at low altitude. It uses inertial navigation, GPS, and terrain matching to achieve high accuracy over long distances. Cruise missiles can be launched from land, sea, or air platforms. Their low-altitude flight profile makes them difficult to detect and intercept.}{catRM}
\glsadd{ICBM}\dictentry{ICBM}{ICBM is an Intercontinental Ballistic Missile with a range exceeding 5500 kilometers, designed to deliver nuclear or conventional warheads across continents. It is launched from a fixed silo or mobile launcher and follows a ballistic trajectory through space before reentering the atmosphere. ICBMs are a primary component of strategic nuclear deterrent forces. Modern ICBMs may carry multiple independently targetable warheads.}{catRM}
\glsadd{IRBM}\dictentry{IRBM}{IRBM is an Intermediate-Range Ballistic Missile with a typical range between 3000 and 5500 kilometers. It fills the gap between theater-range missiles and intercontinental systems. IRBMs are typically deployed on mobile launchers for increased survivability. These missiles played a significant role in Cold War nuclear strategy and remain in the arsenals of several nations.}{catRM}
\glsadd{SRBM}\dictentry{SRBM}{SRBM is a Short-Range Ballistic Missile with a range typically under 1000 kilometers. It provides tactical strike capability for battlefield and theater operations. SRBMs are highly mobile and can be rapidly deployed and fired, making them effective in time-sensitive targeting. Their short flight time makes defense against them particularly challenging.}{catRM}
\glsadd{SLBM}\dictentry{SLBM}{SLBM is a Submarine-Launched Ballistic Missile designed to be fired from submerged nuclear-powered ballistic missile submarines. SLBMs provide a survivable second-strike capability because submarines are extremely difficult to track. They use solid-propellant rocket motors and are ejected from underwater tubes before ignition. SLBMs are a critical leg of the nuclear triad for major nuclear powers.}{catRM}
\glsadd{SAM}\dictentry{SAM}{SAM is a Surface-to-Air Missile designed to intercept and destroy aircraft, helicopters, cruise missiles, and other aerial threats. SAMs are launched from ground-based or shipboard launchers and use radar, infrared, or command guidance to reach their targets. They range from short-range point defense systems to long-range area defense interceptors. SAM systems are a cornerstone of integrated air defense networks worldwide.}{catRM}
\glsadd{AAM}\dictentry{AAM}{AAM is an Air-to-Air Missile launched from an aircraft to destroy another aerial target. AAMs use infrared homing for short-range engagements and active or semi-active radar homing for beyond-visual-range intercepts. Modern AAMs feature high off-boresight capability and advanced counter-countermeasure systems. They are the primary offensive weapon of fighter and interceptor aircraft.}{catRM}
\glsadd{ASM}\dictentry{ASM}{ASM is an Air-to-Surface Missile designed to strike ground or naval targets from an aircraft. ASMs may use radar, infrared, GPS, or laser guidance to achieve precision strikes against fixed and moving targets. They carry a variety of warhead types including penetrating, fragmentation, and blast charges. ASMs extend the strike range of combat aircraft beyond the reach of bombs and unguided rockets.}{catRM}
\glsadd{ATGM}\dictentry{ATGM}{ATGM is an Anti-Tank Guided Missile designed to defeat armored vehicles from a stand-off range. ATGMs use wire, laser, infrared, or radar guidance to steer toward the target with high accuracy. They typically carry a shaped-charge warhead capable of penetrating heavy armor. Modern ATGMs are man-portable or vehicle-mounted and are effective against the most protected targets.}{catRM}
\glsadd{Anti-Ship-Missile}\dictentry{Anti-Ship Missile}{Anti-Ship Missile is a missile specifically designed to engage and destroy naval vessels. It typically uses a low-altitude sea-skimming flight profile to evade ship defenses and active radar or infrared homing for terminal guidance. Anti-ship missiles carry large warheads designed to cripple or sink warships. They are deployed from ships, submarines, aircraft, and coastal launch sites.}{catRM}
\glsadd{Hypersonic-Missile}\dictentry{Hypersonic Missile}{Hypersonic Missile is a missile capable of sustained flight at speeds exceeding Mach 5. At these velocities the missile can reach any target in the world within minutes and its extreme speed and maneuverability make interception extremely difficult. Hypersonic missiles use either scramjet air-breathing engines or boost-glide configurations. They represent a transformative development in strategic and tactical strike capabilities.}{catRM}
\glsadd{Hypersonic-Glide-Vehicle}\dictentry{Hypersonic Glide Vehicle}{Hypersonic Glide Vehicle is an unpowered maneuverable warhead that is boosted to hypersonic speed by a rocket and then glides through the upper atmosphere to its target at speeds above Mach 5. Unlike a ballistic reentry vehicle it can perform significant cross-range and down-range maneuvers during flight. This maneuverability combined with extreme speed makes it very challenging for missile defense systems to track and intercept.}{catRM}
\glsadd{MIRV}\dictentry{MIRV}{MIRV stands for Multiple Independently Targetable Reentry Vehicle. A single missile carries several warheads, each on its own reentry vehicle, that can be aimed at separate targets. After the boost phase the missile dispenses individual warheads along different trajectories. MIRV technology greatly increases the destructive capability per missile and complicates defensive countermeasures.}{catRM}
\glsadd{Reentry-Vehicle}\dictentry{Reentry Vehicle}{Reentry Vehicle is the portion of a ballistic missile that carries the warhead through the intense heating and mechanical loads of atmospheric reentry. It is shaped for aerodynamic stability and protected by a heat shield made from ablative or reusable thermal protection materials. The reentry vehicle must survive temperatures exceeding several thousand degrees while delivering its payload to the target with precision.}{catRM}
\glsadd{Warhead}\dictentry{Warhead}{Warhead is the payload section of a missile or torpedo that delivers the destructive effect against the target. It may contain high explosives, nuclear material, fragmentation elements, or specialized submunitions. The warhead is designed to survive the forces of launch, flight, and impact before detonating at the appropriate time and location. Warhead design balances lethality with safety and reliability requirements.}{catRM}
\glsadd{Conventional-Warhead}\dictentry{Conventional Warhead}{Conventional Warhead is a warhead that uses non-nuclear explosive materials to produce blast, fragmentation, or penetrating effects. Types include blast-fragmentation, shaped-charge, and kinetic energy warheads. Conventional warheads are used in the majority of tactical and strategic missile systems. Their design focuses on maximizing effectiveness against specific target types while minimizing collateral damage.}{catRM}
\glsadd{Nuclear-Warhead}\dictentry{Nuclear Warhead}{Nuclear Warhead is a warhead that derives its destructive force from nuclear fission, fusion, or a combination of both reactions. A single nuclear warhead can destroy an entire city, making it orders of magnitude more powerful than any conventional explosive. Nuclear warheads are deployed on ballistic missiles, cruise missiles, and gravity bombs. Their existence shapes global strategic deterrence and arms control policy.}{catRM}
\glsadd{Guidance-System}\dictentry{Guidance System}{Guidance System is the integrated set of sensors, computers, and actuators that directs a missile toward its target. It determines the missile's current position, computes the required trajectory, and issues steering commands to the control surfaces or thrust vectors. Guidance methods include inertial, GPS, terminal homing, and command guidance. The choice of guidance system determines accuracy, range, and resistance to countermeasures.}{catRM}
\glsadd{Inertial-Guidance}\dictentry{Inertial Guidance}{Inertial Guidance is a self-contained navigation method using onboard accelerometers and gyroscopes to measure the missile's acceleration and rotation. A computer integrates these measurements over time to compute position and velocity without external signals. This makes inertial guidance immune to jamming and interference. However, sensor errors accumulate over time, so inertial systems are often augmented with GPS or other corrections.}{catRM}
\glsadd{GPS-Guidance}\dictentry{GPS Guidance}{GPS Guidance uses signals from the Global Positioning System satellite constellation to determine the missile's precise position and velocity in real time. A GPS receiver on the missile compares satellite signals to compute a three-dimensional position fix. This provides highly accurate guidance that does not degrade over time. GPS guidance is vulnerable to jamming and spoofing, so it is typically combined with inertial navigation.}{catRM}
\glsadd{Terrain-Contour-Matching}\dictentry{Terrain Contour Matching}{Terrain Contour Matching is a guidance method that compares a radar altimeter profile of the terrain below the missile with a stored digital terrain map. The missile matches the observed terrain profile to the map to determine its position and apply corrections. This technique provides high accuracy without relying on external navigation signals. It is particularly effective for cruise missiles flying at low altitude over well-mapped terrain.}{catRM}
\glsadd{Scene-Matching}\dictentry{Scene Matching}{Scene Matching is a terminal guidance technique that compares an onboard camera image of the target area with a preloaded reference image. The missile identifies corresponding features to determine its position relative to the target and applies final course corrections. Scene matching provides very high accuracy for precision strike applications. It requires good visibility and pre-mission intelligence to create the reference imagery.}{catRM}
\glsadd{Active-Radar-Homing}\dictentry{Active Radar Homing}{Active Radar Homing is a missile guidance method where the missile carries both a radar transmitter and receiver. In the terminal phase the missile emits radar pulses and tracks the reflections from the target autonomously. This provides fire-and-forget capability because the launch platform does not need to illuminate the target after launch. Active radar homing is used in modern beyond-visual-range air-to-air missiles.}{catRM}
\glsadd{Semi-Active-Radar-Homing}\dictentry{Semi-Active Radar Homing}{Semi-Active Radar Homing is a guidance method where the launch platform illuminates the target with radar and the missile homes on the reflected energy. The missile carries only a receiver, which reduces its size and cost compared to active radar systems. However, the launch aircraft must maintain radar lock on the target until impact. Semi-active homing remains widely used for surface-to-air and air-to-air missile systems.}{catRM}
\glsadd{Passive-Radar-Homing}\dictentry{Passive Radar Homing}{Passive Radar Homing is a missile guidance method that tracks electromagnetic emissions from the target itself, such as its own radar transmissions. The missile does not emit any energy, making it virtually undetectable by the target. This provides a significant stealth advantage and can be effective against ships and aircraft with active radar. The limitation is dependence on the target continuously emitting radar signals.}{catRM}
\glsadd{Infrared-Homing}\dictentry{Infrared Homing}{Infrared Homing is a guidance method where the missile tracks the thermal radiation emitted by the target, typically from engine exhaust or aerodynamic heating. The missile carries an infrared seeker that locks onto the heat signature and guides autonomously. Modern infrared seekers are highly sensitive and can engage at long range. Infrared homing provides fire-and-forget capability but can be affected by countermeasure flares and weather.}{catRM}
\glsadd{Laser-Homing}\dictentry{Laser Homing}{Laser Homing is a guidance method where the missile tracks reflected laser energy from the target. A designator on the launch platform or a ground team illuminates the target with a coded laser beam, and the missile's seeker locks onto the reflected energy. This provides very high accuracy, often within a meter of the aim point. Laser homing requires line-of-sight to the target and can be degraded by smoke, dust, or fog.}{catRM}
\glsadd{Command-Guidance}\dictentry{Command Guidance}{Command Guidance is a missile guidance method where an external station tracks both the missile and the target and transmits steering commands to the missile via a radio link. The missile follows these commands to intercept the target. This approach keeps the guidance electronics on the ground, simplifying the missile design. However, it requires continuous tracking and command transmission, which limits the number of simultaneous engagements.}{catRM}
\glsadd{Beam-Riding}\dictentry{Beam Riding}{Beam Riding is a guidance method where the missile attempts to stay within a radar or laser beam directed at the target. The missile senses its position relative to the beam center and applies corrections to remain on the beam axis. The operator or fire control system keeps the beam pointed at the target throughout the engagement. Beam riding is simple to implement but accuracy degrades at longer ranges as the beam spreads.}{catRM}
\glsadd{Wire-Guidance}\dictentry{Wire Guidance}{Wire Guidance is a method where steering commands are transmitted from the launch platform to the missile through a thin wire unspooled during flight. The wire carries electrical signals encoding guidance commands. This method is resistant to electronic countermeasures because the link is physical rather than radio-based. Wire guidance is commonly used in anti-tank missiles where the operator maintains line of sight to the target.}{catRM}
\glsadd{Fiber-Optic-Guidance}\dictentry{Fiber Optic Guidance}{Fiber Optic Guidance is an advanced form of wire guidance using a thin fiber optic cable instead of an electrical wire to transmit commands and return video from the missile's seeker. The fiber optic link provides much higher bandwidth, enabling the operator to see the missile's viewpoint in real time. This allows mid-course target updates and man-in-the-loop control. Fiber optic guidance is used in modern anti-tank and multi-role missile systems.}{catRM}
\glsadd{Terminal-Guidance}\dictentry{Terminal Guidance}{Terminal Guidance is the final phase of a missile's flight in which it acquires and homes in on the specific target. The missile switches from midcourse navigation to its onboard seeker for precise target tracking. Terminal guidance requires the highest accuracy and is where the missile is most vulnerable to countermeasures. The seeker and algorithms used in this phase largely determine the missile's single-shot kill probability.}{catRM}
\glsadd{Midcourse-Guidance}\dictentry{Midcourse Guidance}{Midcourse Guidance is the intermediate phase of a missile's flight between the initial boost and terminal homing phases. During midcourse the missile follows a planned trajectory using inertial navigation, GPS, or command updates from the launch platform. This phase covers the longest portion of the flight distance. Efficient midcourse guidance conserves energy for the terminal phase and enables the missile to reach distant targets.}{catRM}
\glsadd{Boost-Phase}\dictentry{Boost Phase}{The Boost Phase is the initial powered segment of a missile flight when the propulsion system generates thrust. The vehicle accelerates from rest, shedding mass as propellant is consumed and sometimes jettisoning spent stages. This phase determines the velocity and trajectory parameters governing all subsequent flight. For intercontinental missiles, the boost phase lasts only minutes but produces the brightest infrared signature.}{catRM}
\glsadd{Coast-Phase}\dictentry{Coast Phase}{The Coast Phase is the unpowered segment of missile flight after propulsion shuts down. The vehicle travels under gravity and existing momentum, following a ballistic arc if no guidance corrections are applied. The coast phase can account for the majority of a ballistic missile's total flight time on long-range trajectories. Minimizing perturbations during this phase is essential for maintaining accuracy at the target.}{catRM}
\glsadd{Reentry-Phase}\dictentry{Reentry Phase}{The Reentry Phase is the final ballistic trajectory segment where the vehicle reenters the atmosphere at extreme velocity. Aerodynamic heating and deceleration forces reach their peak, requiring robust thermal protection systems. For intercontinental missiles, reentry speeds exceed Mach 20, making thermal management a critical engineering challenge. The reentry phase is also where atmospheric interaction can degrade accuracy if the vehicle is not properly designed.}{catRM}
\glsadd{Trajectory}\dictentry{Trajectory}{A Trajectory is the three-dimensional path followed by a missile or rocket from launch to impact or orbit insertion. Its shape is determined by thrust profile, guidance commands, gravitational forces, and aerodynamic effects. Trajectory design balances range, payload capacity, and survivability requirements. Different missions produce different shapes, from depressed profiles for fast attack to lofted arcs for maximum altitude.}{catRM}
\glsadd{Ballistic-Trajectory}\dictentry{Ballistic Trajectory}{A Ballistic Trajectory is the parabolic path a projectile follows after propulsion ceases, governed entirely by gravity and initial velocity. Once a missile enters its ballistic phase, no further thrust alters its course unless an active reentry vehicle maneuvers. Engineers calculate range and impact point using well-established equations of motion. This type of flight defines intermediate and intercontinental range ballistic missiles.}{catRM}
\glsadd{Lifting-Trajectory}\dictentry{Lifting Trajectory}{A Lifting Trajectory uses aerodynamic lift to extend range or modify a reentry vehicle's path rather than following a purely ballistic arc. The vehicle glides or maneuvers within the atmosphere to reach distant or evasive targets. This approach can significantly increase effective range without additional propulsion. Lifting trajectories complicate interception because the vehicle can change direction during atmospheric flight.}{catRM}
\glsadd{Depressed-Trajectory}\dictentry{Depressed Trajectory}{A Depressed Trajectory is flatter and lower than the standard minimum-energy trajectory for a given range. By launching at a lower angle and higher velocity, the missile spends more time at lower altitudes. This reduces adversary warning time because the missile stays below the radar horizon longer. The tradeoff is higher fuel expenditure and increased structural loads during boost.}{catRM}
\glsadd{Lofted-Trajectory}\dictentry{Lofted Trajectory}{A Lofted Trajectory arcs significantly higher than the minimum-energy trajectory for a given range. The missile launches at a steeper angle, reaching extreme altitudes before descending toward its target. Lofted trajectories test missile performance under high reentry heating conditions or reduce the horizontal range of test flights. The steep descent angle also changes the geometry of interception for defense systems.}{catRM}
\glsadd{Minimum-Energy-Trajectory}\dictentry{Minimum Energy Trajectory}{A Minimum Energy Trajectory requires the least propellant to deliver a payload to a specified range. For any target distance, an optimal launch angle and burn profile minimize total energy expenditure. Engineers use minimum energy trajectories as the baseline for performance calculations and range planning. Departures from this solution in either direction reduce achievable range or require additional propellant.}{catRM}
\glsadd{Maximum-Range-Trajectory}\dictentry{Maximum Range Trajectory}{A Maximum Range Trajectory allows a missile to reach the farthest possible distance with its available propellant. This path is typically close to the minimum energy solution but may incorporate modifications for atmospheric losses and Earth rotation effects. Achieving maximum range is a primary design objective for intercontinental ballistic missiles. Range margin provides flexibility in choosing launch sites and targets.}{catRM}
\glsadd{Launch-Azimuth}\dictentry{Launch Azimuth}{Launch Azimuth is the compass direction, measured clockwise from true north, in which a missile is fired. It determines the vehicle's ground track and, combined with trajectory profile, defines the impact point or orbital plane. For intercontinental missiles, the azimuth is chosen to avoid overflying neutral territories and optimize trajectory geometry. Small azimuth errors translate directly into cross-range errors at the target.}{catRM}
\glsadd{Impact-Point}\dictentry{Impact Point}{The Impact Point is the geographic location where a missile warhead or reentry vehicle strikes the Earth's surface. Predicting it requires accurate knowledge of trajectory, atmospheric conditions, and any terminal maneuvers by the vehicle. For guided missiles, the impact point is continuously refined by the guidance system throughout flight. The distance between the intended and actual impact points defines the missile's accuracy.}{catRM}
\glsadd{CEP}\dictentry{CEP}{CEP, or Circular Error Probable, measures accuracy as the radius of a circle centered on the target within which fifty percent of warheads fall. A smaller CEP indicates higher precision, with modern strategic warheads achieving CEPs in tens of meters. CEP is used to assess weapon system effectiveness and determine how many warheads are needed to destroy a target. Improving CEP remains a continuous goal in missile guidance technology.}{catRM}
\glsadd{Kill-Probability}\dictentry{Kill Probability}{Kill Probability is the statistical likelihood that a missile or interceptor will successfully destroy its target upon detonation. It depends on accuracy, warhead yield, target hardness, and the effectiveness of any target defenses. Kill probability is a central metric in nuclear strategy and missile defense planning, directly influencing force structure requirements. Multiple warheads or repeated engagements may raise cumulative kill probability against hardened targets.}{catRM}
\glsadd{Lethality}\dictentry{Lethality}{Lethality describes a warhead's capacity to cause the desired level of damage to a target. It is a function of yield, blast overpressure, thermal radiation, fragmentation pattern, and target vulnerability to each effect. Lethality analysis ensures warheads meet mission requirements against specific target types. Maximizing lethality while minimizing weapon weight is a fundamental challenge in warhead engineering.}{catRM}
\glsadd{Fragmentation-Warhead}\dictentry{Fragmentation Warhead}{A Fragmentation Warhead scatters thousands of high-velocity metal fragments upon detonation, creating a lethal shrapnel cloud. The casing is pre-scored to break into fragments optimized for the intended target type. These warheads are widely used in surface-to-air and air-to-air missiles where direct hits are unlikely. The effective kill radius depends on fragment mass, velocity, and pattern density.}{catRM}
\glsadd{Shaped-Charge}\dictentry{Shaped Charge}{A Shaped Charge uses a concave metal liner focused by explosive detonation into a high-velocity metal jet. This jet penetrates thick armor or reinforced concrete by concentrating enormous energy on a small area. Shaped charges are used in anti-tank missiles, bunker-busting warheads, and armor-piercing projectiles. The Monroe effect powering them was discovered in the nineteenth century but not practically applied until the mid-twentieth century.}{catRM}
\glsadd{Kinetic-Kill-Vehicle}\dictentry{Kinetic Kill Vehicle}{A Kinetic Kill Vehicle destroys its target through pure kinetic energy of direct impact without an explosive warhead. The interceptor must achieve extremely precise guidance to collide with a target moving at thousands of meters per second. This approach simplifies interceptor design by eliminating warhead complexity but demands extraordinary sensor and guidance accuracy. The concept is central to several modern ballistic missile defense architectures.}{catRM}
\glsadd{Decoy}\dictentry{Decoy}{A Decoy is a device deployed during missile flight to confuse enemy sensors and interceptors. Decoys replicate the radar, infrared, or visual signatures of actual warheads to force defenders to waste resources on false targets. In strategic missile systems, decoys are essential for ensuring some warheads reach their targets. Their effectiveness depends on the sophistication of the defense system's discrimination capabilities.}{catRM}
\glsadd{Countermeasure}\dictentry{Countermeasure}{A Countermeasure is any technique or tactic employed to negate or degrade an enemy's weapons or sensors. In missile contexts, countermeasures range from electronic jamming to physical decoys and evasive maneuvers. Countermeasure development is an ongoing competition with defense systems, as each new countermeasure prompts defensive improvements. Modern missile designs integrate multiple countermeasure layers to defeat increasingly capable defenses.}{catRM}
\glsadd{Chaff}\dictentry{Chaff}{Chaff consists of thin aluminum or metallic-coated fiber strips dispensed to create a radar-reflective cloud. When a radar-guided missile encounters chaff, it may lock onto the false target instead of the actual vehicle. Chaff has been a standard electronic countermeasure since World War Two and remains effective against many radar systems. Modern dispensers can release precisely metered quantities to create decoy clouds of optimal size.}{catRM}
\glsadd{Flare}\dictentry{Flare}{A Flare is an infrared countermeasure that emits intense heat to divert heat-seeking missiles. When an infrared-guided missile approaches, flares are dispensed to present a more attractive heat source that the missile tracks instead of the engine. Flares are a standard defensive system on military aircraft and some ground vehicles. Advanced infrared seekers now employ counter-countermeasure techniques to distinguish flares from actual targets.}{catRM}
\glsadd{Electronic-Countermeasure}\dictentry{Electronic Countermeasure}{An Electronic Countermeasure uses electromagnetic energy to degrade or deny the effectiveness of enemy radar, communications, or guidance systems. Techniques include noise jamming, deception jamming, and electronic deception, each targeting different aspects of adversary systems. Electronic countermeasures are critical for integrated defense suppression during strike missions. Their effectiveness depends on knowledge of the enemy's operating frequencies and signal characteristics.}{catRM}
\glsadd{Penetration-Aid}\dictentry{Penetration Aid}{A Penetration Aid is any device or technique that helps a missile warhead survive interception. These include decoys, chaff, flares, jamming devices, and maneuvering reentry vehicles, often used in combination. Strategic missile designers integrate penetration aid suites to maximize the probability of at least one warhead reaching its target. The competition between penetration aids and defense discrimination is a central dynamic in strategic stability.}{catRM}
\glsadd{Anti-Ballistic-Missile}\dictentry{Anti-Ballistic Missile}{An Anti-Ballistic Missile is designed to intercept and destroy incoming ballistic missile warheads. These systems use radar and satellite tracking to compute intercept trajectories, guiding the interceptor to a collision or proximity detonation. Anti-ballistic missiles must operate under extreme time pressure, as engagements may last only minutes. Their development has profound implications for strategic nuclear deterrence and international security.}{catRM}
\glsadd{Interceptor}\dictentry{Interceptor}{An Interceptor is a missile designed to destroy incoming threats including ballistic missiles, cruise missiles, and aircraft. Interceptors rely on advanced sensors, guidance algorithms, and high maneuverability for precise timing. Modern families include terminal defense interceptors for point defense and midcourse interceptors for area defense. Effectiveness is measured by single-shot kill probability against representative threats.}{catRM}
\glsadd{Hit-to-Kill}\dictentry{Hit-to-Kill}{Hit-to-Kill destroys the target through direct kinetic impact without a warhead. At closing velocities of several kilometers per second, the energy released equals a substantial explosive detonation. This approach eliminates interceptor warhead complexity and weight. Hit-to-Kill requires extraordinary guidance precision, with the kill vehicle making final adjustments in the terminal seconds of flight.}{catRM}
\glsadd{Proximity-Fuse}\dictentry{Proximity Fuse}{A Proximity Fuse detects the distance to a target and triggers detonation within a preset lethal range. It dramatically increases kill probability compared to contact fuses because the warhead need not physically strike the target. Proximity fuses are standard in surface-to-air and air-to-air missiles where direct impact is unlikely. Modern versions use radar or laser sensing for reliable detection in complex environments.}{catRM}
\glsadd{Contact-Fuse}\dictentry{Contact Fuse}{A Contact Fuse triggers the warhead at the moment of physical impact with the target. It is the simplest and most reliable fuse type, requiring no electronics or sensing apparatus. Contact fuses are used when direct hits are expected or when penetration before detonation is desired. Their limitation is that they provide no kill probability enhancement for near misses.}{catRM}
\glsadd{Time-Fuse}\dictentry{Time Fuse}{A Time Fuse triggers the warhead after a precisely measured interval from launch or another reference event. It is used when weapons must detonate at specific trajectory points, such as for airburst effects. Modern electronic time fuses achieve microsecond accuracy for precise control of burst altitude. They are important for fragmentation weapons where optimal burst height maximizes the lethal area.}{catRM}
\glsadd{Airburst}\dictentry{Airburst}{An Airburst is a detonation at predetermined altitude above the ground rather than at the surface. Airbursts maximize the area affected by blast overpressure or fragmentation since the explosion is not absorbed by the ground. The optimal burst height depends on weapon yield, target type, and desired damage criteria. For nuclear weapons, airbursting can significantly increase the effective blast radius compared to ground detonation.}{catRM}
\glsadd{Terminal-Phase}\dictentry{Terminal Phase}{The Terminal Phase is the interceptor's final flight segment where onboard sensors acquire the target and execute terminal guidance maneuvers. The interceptor must rapidly compute and execute last-second corrections to close the remaining distance. This phase is the most demanding portion of interception, requiring sensors and actuators with extremely fast response times. Success or failure of the entire engagement is determined in these final moments.}{catRM}
\glsadd{Launch-Detection}\dictentry{Launch Detection}{Launch Detection identifies a missile launch through infrared satellite sensors, radar, or other surveillance platforms. Early detection is the first link in the ballistic missile defense chain, providing initial tracking data for intercept solutions. Space-based infrared sensors can detect a boosting rocket's heat signature within seconds of ignition. Rapid and reliable detection is essential for providing adequate warning time to decision-makers and defensive systems.}{catRM}
\glsadd{Tracking-Radar}\dictentry{Tracking Radar}{Tracking Radar continuously follows a specific target, providing precise measurements of position, velocity, and trajectory. These radars use narrow beamwidths and high update rates to maintain locks on fast-moving missiles or aircraft. The data feeds fire control computers that compute intercept solutions. Modern tracking radars employ monopulse or track-while-scan techniques for high accuracy in cluttered environments.}{catRM}
\glsadd{Fire-Control-Radar}\dictentry{Fire Control Radar}{Fire Control Radar is integrated with weapons platforms to provide targeting data for directing missile or gun engagements. It tracks targets, computes intercept points, and provides guidance commands to weapons in flight. Fire control radars must balance accuracy, range, and resistance to electronic countermeasures. They form the critical link between detection and engagement in all modern air defense systems.}{catRM}
\glsadd{Phased-Array-Radar}\dictentry{Phased Array Radar}{A Phased Array Radar uses antenna elements whose signal phases are electronically controlled to steer the beam without physical movement. This allows rapid beam switching, simultaneous multi-target tracking, and concurrent search and track functions. Phased array radars are the backbone of modern ballistic missile defense networks. Their electronic steering provides far greater flexibility and response speed than mechanically scanned systems.}{catRM}
\glsadd{Early-Warning-Radar}\dictentry{Early Warning Radar}{Early Warning Radar detects incoming missiles or aircraft at the greatest possible distance to maximize warning time. These radars typically operate at lower frequencies for better range, trading some tracking precision for early detection. They form the outer layer of a nation's air and missile defense architecture. Data from early warning radars is passed to tracking and fire control systems for engagement.}{catRM}
\glsadd{Over-the-Horizon-Radar}\dictentry{Over-the-Horizon Radar}{Over-the-Horizon Radar uses radio wave reflection off the ionosphere to detect targets beyond conventional radar line of sight. By bouncing signals off the ionosphere back to Earth's surface, these radars can detect targets thousands of kilometers away. They provide valuable early warning against low-flying cruise missiles hidden by Earth's curvature. The technique is sensitive to ionospheric conditions that vary with time of day and solar activity.}{catRM}
\glsadd{Satellite-Warning-System}\dictentry{Satellite Warning System}{A Satellite Warning System uses space-based infrared sensors to detect missile launches from orbit, providing global early warning coverage. These systems detect launches within seconds and relay information to ground stations for processing. Space-based detection is the only means of achieving timely warning for submarine-launched ballistic missiles over open ocean. Several nations maintain satellite warning constellations as critical strategic defense components.}{catRM}
\glsadd{Ballistic-Missile-Defense}\dictentry{Ballistic Missile Defense}{Ballistic Missile Defense is the coordinated system of sensors, command and control, and interceptors designed to destroy incoming ballistic missiles before they reach targets. A complete architecture includes space-based sensors, ground-based radars, battle management systems, and multiple interceptor tiers. The challenge lies in extremely short engagement timelines and discriminating warheads from decoys. Effective defense requires seamless integration of all network elements.}{catRM}
\glsadd{THAAD}\dictentry{THAAD}{THAAD, or Terminal High Altitude Area Defense, is a US Army missile defense system designed to intercept short-, medium-, and intermediate-range ballistic missiles during their terminal phase. It uses a hit-to-kill interceptor destroying warheads through direct kinetic impact. THAAD fills the gap between lower-altitude Patriot and higher-altitude Ground-Based Interceptors. It has been deployed in several regions to protect against regional ballistic missile threats.}{catRM}
\glsadd{Patriot}\dictentry{Patriot}{Patriot is a US Army surface-to-air missile system originally designed for air defense and later upgraded for limited theater ballistic missile defense. The system uses radar-guided interceptors and sophisticated command and control to engage threats at various altitudes and ranges. Patriot batteries have been deployed worldwide and seen combat in multiple conflicts. The system continues receiving upgrades to improve capability against evolving missile threats.}{catRM}
\glsadd{Aegis}\dictentry{Aegis}{Aegis is a US Navy combat system built around the AN/SPY-1 phased array radar and Standard Missile family, providing area defense against aircraft, cruise missiles, and ballistic missiles. Deployed on guided missile cruisers and destroyers, it provides mobile defensive coverage across the world's oceans. Aegis Ashore extends this capability to land-based sites in Europe and the Pacific. Its simultaneous multi-target tracking makes it a cornerstone of integrated missile defense.}{catRM}
\glsadd{Ground-Based-Interceptor}\dictentry{Ground-Based Interceptor}{The Ground-Based Interceptor is a kinetic kill vehicle launched from fortified silos in the United States as part of the Ground-based Midcourse Defense system protecting the continental United States from limited ICBM attacks. It uses hit-to-kill to collide with incoming warheads during midcourse outside the atmosphere. Interceptors are based at Fort Greely, Alaska and Vandenberg Air Force Base, California. This system is the last line of defense against strategic ballistic missile attacks.}{catRM}
\glsadd{Iron-Dome}\dictentry{Iron Dome}{Iron Dome is a short-range air defense system developed by Rafael Advanced Defense Systems in Israel, designed to intercept rockets, artillery, and mortar shells fired at populated areas. It uses radar to calculate trajectories and launches Tamir interceptors only against projectiles predicted to strike populated areas. Iron Dome has demonstrated high success rates, intercepting thousands of rockets since deployment. The system has been widely studied and adopted by other nations.}{catRM}
\glsadd{David-Sling}\dictentry{David Sling}{David Sling, also known as the Magic Wand, is an Israeli medium-range air defense system developed by Rafael and Raytheon to fill the gap between Iron Dome and the Arrow system. It intercepts tactical ballistic missiles, medium-range rockets, cruise missiles, and aircraft using Stunner interceptors. David Sling uses advanced multi-pulse seeker technology for high kill probability. The system entered operational service in 2017, strengthening Israel's layered defense architecture.}{catRM}
\glsadd{Arrow}\dictentry{Arrow}{The Arrow is an Israeli anti-ballistic missile system developed with the United States, designed to intercept exo-atmospheric and endo-atmospheric ballistic missile threats. The Arrow-2 provides endo-atmospheric interception while the Arrow-3 is designed for exo-atmospheric hit-to-kill above the atmosphere. The system uses the Green Pine radar for detection and tracking. Arrow forms the upper tier of Israel's multi-layered missile defense architecture.}{catRM}
\glsadd{S-400}\dictentry{S-400}{The S-400 Triumf is a Russian long-range surface-to-air missile system designed to engage aircraft, cruise missiles, and ballistic missiles at ranges up to four hundred kilometers. It uses multiple missile types optimized for different target categories and altitudes. The S-400 features rapid deployment and can track hundreds of targets while engaging dozens simultaneously. The system has been exported to several nations and generated significant international attention.}{catRM}
\glsadd{S-500}\dictentry{S-500}{The S-500 Prometey is a Russian next-generation air and missile defense system designed to engage ICBMs, hypersonic glide vehicles, and low Earth orbit satellites. It represents a significant step beyond the S-400 with greater range and higher altitude capability. The S-500 complements the S-400 in layered defense, handling the most demanding targets at greater distances. Its development reflects the growing importance of countering advanced missile and space threats.}{catRM}
\glsadd{Ramjet-Missile}\dictentry{Ramjet Missile}{A Ramjet Missile uses a ramjet engine that compresses incoming air through the missile's forward motion without a mechanical compressor. Ramjets offer higher specific impulse than rockets at supersonic speeds, allowing greater range for a given fuel load. Several advanced anti-ship and surface-to-air missiles use ramjet propulsion for high sustained speeds. The limitation is that ramjets cannot produce thrust at zero velocity, requiring a booster stage for launch.}{catRM}
\glsadd{Scramjet-Missile}\dictentry{Scramjet Missile}{A Scramjet Missile uses a supersonic combustion ramjet engine that operates at hypersonic speeds while maintaining supersonic airflow through combustion. Scramjet technology enables sustained flight at speeds exceeding Mach 5, opening possibilities for long-range strike and rapid global reach. The engineering challenges include managing extreme temperatures and achieving stable combustion at hypersonic velocities. Several nations are actively developing scramjet-powered missiles and vehicles.}{catRM}
\glsadd{Solid-Rocket-Booster}\dictentry{Solid Rocket Booster}{A Solid Rocket Booster is a solid-propellant motor attached to a launch vehicle's sides to provide additional thrust during initial ascent. After propellant exhaustion, boosters are jettisoned to reduce vehicle mass for the remainder of ascent. The Space Shuttle solid rocket boosters were the most prominent examples, but strap-on solid boosters are used on many launch vehicles worldwide. Solid boosters offer high thrust-to-weight ratios and simplicity compared to liquid alternatives.}{catRM}
\glsadd{Dual-Thrust-Motor}\dictentry{Dual Thrust Motor}{A Dual Thrust Motor produces two distinct thrust levels during its burn, typically a high-thrust boost phase followed by lower-thrust sustain. This is achieved through grain geometry design with sections that burn at different rates. Dual thrust motors are common in tactical missiles needing rapid acceleration to flight speed followed by sustained cruise. The two-phase profile optimizes both launch response time and total available impulse.}{catRM}
\glsadd{Boost-Sustain-Motor}\dictentry{Boost-Sustain Motor}{A Boost-Sustain Motor provides an initial high-thrust boost to accelerate the missile to flight speed, followed by lower-thrust sustain maintaining velocity during cruise. The profile is achieved through propellant grain design with varying cross-sectional areas. This motor type is widely used in air-to-air and surface-to-air missiles requiring both rapid acceleration and efficient cruise. The transition between phases must be smooth to avoid flight instabilities.}{catRM}
\glsadd{Thrust-Profile}\dictentry{Thrust Profile}{The Thrust Profile represents how a rocket engine's thrust varies over the duration of its burn. It is determined by propellant grain geometry, combustion chamber design, and nozzle characteristics. Engineers design the thrust profile to match mission requirements, balancing acceleration, structural loads, and fuel efficiency. A well-optimized profile ensures the vehicle achieves required velocity at burnout while minimizing stresses on structure and payload.}{catRM}
\glsadd{Flame-Deflector}\dictentry{Flame Deflector}{A Flame Deflector is a reinforced structure at a launch pad's base designed to redirect intense exhaust gases and heat away from the vehicle and infrastructure. These are typically angled channels of steel and concrete, often with water injection to reduce thermal loads. The deflector must withstand temperatures exceeding thousands of degrees and enormous acoustic energy during launch. Without it, reflected exhaust could damage the vehicle or ignite surrounding structures.}{catRM}
\glsadd{Exhaust-Plume}\dictentry{Exhaust Plume}{The Exhaust Plume is the visible stream of hot gases, particles, and combustion products expelled from a rocket engine nozzle. Its appearance, size, and shape are determined by propellant chemistry, chamber pressure, nozzle design, and ambient conditions. The plume contains information about engine performance and combustion efficiency. In missile defense, the plume's infrared signature is a primary means of detecting and tracking boosting missiles from space.}{catRM}
\glsadd{Mach-Diamond}\dictentry{Mach Diamond}{A Mach Diamond is the luminous diamond-shaped pattern appearing in the exhaust plume of a supersonic rocket or jet engine. These patterns form when exhaust pressure differs from ambient pressure, creating standing shock and expansion waves. The number and spacing of diamonds provide information about the engine's pressure ratio and exhaust velocity. Mach diamonds are commonly visible during ground tests and launches, particularly at night.}{catRM}
\glsadd{Meteorological-Rocket}\dictentry{Meteorological Rocket}{A Meteorological Rocket is a specialized sounding rocket designed to carry instruments measuring atmospheric temperature, pressure, wind speed, and composition in the upper atmosphere. These measurements are critical for weather forecasting, climate research, and understanding atmospheric dynamics above balloon altitude. The instrumented payload descends by parachute, transmitting data during its fall. Meteorological rockets fill an important gap in atmospheric observation capabilities.}{catRM}
\glsadd{Target-Drone}\dictentry{Target Drone}{A Target Drone is an unmanned aircraft or missile used to simulate enemy threats during missile testing and air defense training. Target drones replicate the speed, altitude, radar cross-section, and maneuvering characteristics of the threats they represent. They are expendable by design, destroyed during engagement. Target drones are essential for validating missile system performance and training operators under realistic conditions.}{catRM}
\glsadd{Burn-Rate}\dictentry{Burn Rate}{The Burn Rate is the velocity at which solid rocket propellant is consumed during combustion, typically measured in millimeters per second. It determines the thrust level and burn duration of a solid rocket motor for a given grain geometry. Burn rate is influenced by propellant chemistry, particle size, chamber pressure, and initial temperature. Engineers control it through formulation and grain design to achieve the desired thrust profile.}{catRM}
\glsadd{Strand-Burn-Rate}\dictentry{Strand Burn Rate}{Strand Burn Rate is the burn rate measured on a small standardized propellant sample called a strand, characterizing the propellant's combustion behavior under controlled conditions. The strand test provides baseline data for motor design simulations and quality control. Testing at various pressures and temperatures develops burn rate models predicting motor performance across operating conditions. Strand testing is routine in solid propellant development and production qualification.}{catRM}
\glsadd{Exhaust-Trail}\dictentry{Exhaust Trail}{An Exhaust Trail is the visible streak left in the atmosphere by a missile, composed of condensed water vapor, particulates, and combustion products. It can persist for extended periods depending on atmospheric conditions, sometimes spreading into cirrus-like clouds at high altitudes. Exhaust trails from missile launches are visible over hundreds of kilometers. Their characteristics provide information about propellant type and engine performance.}{catRM}
\glsadd{Liquid-Injection-Thrust-Vectoring}\dictentry{Liquid Injection Thrust Vectoring}{Liquid Injection Thrust Vectoring redirects rocket engine thrust by injecting a secondary fluid into the nozzle divergent section. The injected fluid creates asymmetric pressure that produces a lateral steering force. This avoids the mechanical complexity of gimbaled nozzles while providing adequate control for many applications. Liquid injection thrust vectoring is used on several operational missile systems due to its simplicity and reliability.}{catRM}
\glsadd{Jet-Vanes}\dictentry{Jet Vanes}{Jet Vanes are surfaces placed in a rocket nozzle's exhaust stream to deflect thrust and provide steering. When a vane rotates into the exhaust flow, it creates an asymmetric force producing a turning moment. Jet vanes are a simple, reliable thrust vector control method for solid-propellant rockets where gimbaled nozzles are impractical. The main disadvantage is exposure to extreme thermal and erosive environments from hot exhaust gases.}{catRM}
\glsadd{Jet-Tab}\dictentry{Jet Tab}{A Jet Tab is a small movable surface in a rocket nozzle's exhaust flow that, when extended, deflects exhaust gas to produce a lateral steering force. Tabs are actuated by small motors or pyrotechnic devices for flight steering commands. They offer a compact alternative to larger vane systems for tactical missiles. Jet tabs provide less control authority and can erode during extended operation compared to heavier-duty alternatives.}{catRM}
\glsadd{Flexible-Bearing}\dictentry{Flexible Bearing}{A Flexible Bearing is a compliant joint supporting and pivoting a rocket nozzle for thrust vector control, allowing gimballing while maintaining a pressure seal. It consists of alternating elastomer and metal layers that flex under actuator loads while containing high-pressure gases. This eliminates sliding seals and their friction and leakage problems. Flexible bearings are used on many modern rocket engines as the primary nozzle gimbal mechanism.}{catRM}
\glsadd{Hot-Gas-Valve}\dictentry{Hot Gas Valve}{A Hot Gas Valve controls high-temperature combustion gas flow within rocket engine systems, typically directing gas to actuators or secondary systems. These valves must operate reliably in extreme thermal environments often exceeding a thousand degrees. They are used in gas generator cycles and for directing bleed gas to pneumatic or hydraulic actuators. Materials and sealing mechanisms are carefully engineered to withstand harsh operating conditions.}{catRM}
\glsadd{Secondary-Injection}\dictentry{Secondary Injection}{Secondary Injection is a thrust vector control technique where fluid is injected into the nozzle wall or divergent section to create asymmetric pressure deflecting the exhaust plume. The injected fluid can be liquid, gas, or solid that vaporizes on contact with hot exhaust. This provides steering without mechanical nozzle movement, advantageous for solid rocket motors. The technique offers moderate control authority with relatively simple hardware.}{catRM}
\glsadd{Main-Propellant-Valve}\dictentry{Main Propellant Valve}{The Main Propellant Valve controls fuel and oxidizer flow from tanks to the engine combustion chamber. It must open and close precisely to control engine start, thrust regulation, and shutdown sequences. Designed for high reliability, fast response, and propellant compatibility, its timing and sequencing are critical for safe and efficient engine performance. Main propellant valves are among the most safety-critical components in liquid rocket engines.}{catRM}
\glsadd{Fill-Valve}\dictentry{Fill Valve}{A Fill Valve is used to load propellant into rocket tanks prior to launch. It must provide leak-tight seals during and after loading, compatible with cryogenic or hazardous propellants. Some fill valves also serve as vent paths during loading to allow displaced gas to escape. Their design balances ground crew access with reliable sealing during flight operations.}{catRM}
\glsadd{Drain-Valve}\dictentry{Drain Valve}{A Drain Valve removes propellant from rocket tanks during ground operations such as launch scrubs or maintenance. It must safely handle hazardous propellants and drain tanks completely to prevent residual propellant safety hazards. Typically located at the tank's lowest point for gravity-assisted draining. Quick, reliable drain capability is important for launch campaign timelines and safety.}{catRM}
\glsadd{Vent-Valve}\dictentry{Vent Valve}{A Vent Valve allows gas to escape from propellant tanks to regulate internal pressure as propellant is loaded, consumed, or as thermal conditions change. It prevents dangerous overpressure or vacuum by maintaining pressure within design limits. During loading, vent valves allow displaced gas to escape as liquid fills the tank. In flight, they manage pressure as propellant levels decrease and ambient conditions change.}{catRM}
\glsadd{Relief-Valve}\dictentry{Relief Valve}{A Relief Valve automatically opens to release pressure when it exceeds a predetermined limit, preventing catastrophic failure of pressure vessels. In rocket systems, they protect propellant tanks, gas systems, and other components from overpressure. They are calibrated to open at specific thresholds and close once pressure returns to safe levels. Relief valves are critical safety elements in all pressurized rocket systems.}{catRM}
\glsadd{Check-Valve}\dictentry{Check Valve}{A Check Valve allows fluid or gas flow in one direction only, automatically closing to prevent reverse flow. In rocket systems, check valves maintain proper flow direction and prevent contamination or dangerous fluid mixing throughout propellant feed, pressurization, and pneumatic systems. They are passive devices requiring no external actuation, enhancing reliability. They are particularly important for preventing oxidizer-fuel cross-contamination.}{catRM}
\glsadd{Solenoid-Valve}\dictentry{Solenoid Valve}{A Solenoid Valve is actuated by an electromagnetic solenoid, providing electrical control of fluid or gas flow. When energized, the solenoid moves a plunger that opens or closes the valve passage. Solenoid valves offer fast response times and precise on-off control suitable for rapid actuation applications. They are widely used in rocket engine control systems for propellant isolation, pressurization, and pneumatic management.}{catRM}
\glsadd{Poppet-Valve}\dictentry{Poppet Valve}{A Poppet Valve uses a mushroom-shaped disc raised from a seat to allow flow and pressed against it to stop flow. Poppet valves offer excellent sealing and positive shutoff, commonly used in propellant feed and gas handling systems. They can be mechanically, pneumatically, or electrically actuated. Their simple, robust design makes them reliable for critical rocket system applications requiring dependable shutoff.}{catRM}
\glsadd{Ball-Valve}\dictentry{Ball Valve}{A Ball Valve uses a rotating ball with a through-hole that turns ninety degrees to open or close flow. When aligned with flow, the valve is fully open with minimal restriction; when perpendicular, it blocks flow completely. Ball valves are used in propellant systems for low pressure drop and reliable shutoff. They are favored for main propellant isolation valves where full flow capacity is required.}{catRM}
\glsadd{Butterfly-Valve}\dictentry{Butterfly Valve}{A Butterfly Valve uses a disc on a rotating shaft that turns within the flow path to control rate. When parallel to flow the valve is open; when perpendicular it blocks flow. Butterfly valves are lightweight and compact, suitable for weight- and space-constrained applications. In rocket systems they are used in propellant lines and gas systems where moderate flow control is needed.}{catRM}
\glsadd{Pneumatic-Valve}\dictentry{Pneumatic Valve}{A Pneumatic Valve is actuated by compressed gas rather than electrical or mechanical means. Pneumatic actuation offers fast response, high force, and intrinsic safety in spark-hazardous environments. In rocket systems, pneumatic valves are commonly used for propellant isolation, engine sequencing, and pressurization control. The compressed gas supply is typically stored in high-pressure bottles charged before launch.}{catRM}
\glsadd{Pyro-Valve}\dictentry{Pyro Valve}{A Pyro Valve is opened or closed by a small pyrotechnic charge, providing single-use, extremely reliable actuation. Once fired, pyro valves cannot be reset, making them suitable for events that must occur exactly once. They are used for propellant line isolation, stage separation, and other critical events where failure is unacceptable. The pyrotechnic actuator provides high force in a compact package with no power requirement during storage.}{catRM}
\glsadd{Propellant-Line}\dictentry{Propellant Line}{A Propellant Line is a pipe or hose carrying liquid propellant from tanks to the engine or between system components. It must withstand the propellant's chemical properties, thermal extremes, vibration, and feed system pressure requirements. Materials selection is critical as incompatible materials cause contamination, corrosion, or catastrophic failure. Line routing and support must account for vehicle flexibility and thermal expansion during flight.}{catRM}
\glsadd{Suction-Line}\dictentry{Suction Line}{A Suction Line is the propellant line on a turbopump's inlet side, carrying propellant from tank to pump. Its design is critical because pump inlet conditions directly affect performance and cavitation risk. The line must provide a smooth, low-resistance flow path with minimal pressure drop. Suction line design requires careful attention to flow conditioning, thermal management, and vibration isolation.}{catRM}
\glsadd{Discharge-Line}\dictentry{Discharge Line}{A Discharge Line is the propellant line on a turbopump's outlet side, carrying pressurized propellant to the combustion chamber or injectors. It must withstand full discharge pressure, which can reach hundreds of atmospheres in high-performance engines. The line must deliver propellant at required pressure and flow rate while minimizing losses and maldistribution. Discharge lines are typically constructed from high-strength alloys for extreme internal pressures.}{catRM}
\glsadd{Helium-Pressurization}\dictentry{Helium Pressurization}{Helium Pressurization maintains propellant tank pressure by introducing high-pressure helium into the tank ullage space. Helium is commonly used because it is inert, non-reactive with propellants, and remains gaseous at cryogenic temperatures. As propellant is consumed, helium replaces displaced volume to maintain pressure for turbopump or pressure-fed engine operation. The system includes high-pressure storage bottles, regulators, and control valves.}{catRM}
\glsadd{Nitrogen-Pressurization}\dictentry{Nitrogen Pressurization}{Nitrogen Pressurization maintains tank pressure using high-pressure nitrogen gas in the ullage space. Nitrogen is less expensive and more available than helium, suitable for smaller rockets and tactical missiles. Unlike helium, nitrogen has limited propellant solubility and is unsuitable for cryogenic applications where it would liquefy. Nitrogen systems are simpler but limited to non-cryogenic, non-reactive propellant applications.}{catRM}
\glsadd{Autogenous-Pressurization}\dictentry{Autogenous Pressurization}{Autogenous Pressurization maintains tank pressure by vaporizing a small propellant portion and routing the gas back to the tank ullage. This eliminates separate pressurant gas, reducing system weight and complexity. Common with liquid oxygen and hydrogen engines where heat exchangers vaporize small propellant flows. The technique requires careful thermal management to ensure sufficient gas generation under all operating conditions.}{catRM}
\glsadd{Accumulator}\dictentry{Accumulator}{An Accumulator is a pressure vessel storing fluid or gas under pressure to provide reserve capacity for pneumatic, hydraulic, or propellant systems. Accumulators absorb pressure pulsations, provide emergency power, and compensate for thermal expansion. They are found in hydraulic actuation, pneumatic control, and propellant feed systems. Stored energy can be released rapidly to power critical functions when needed.}{catRM}
\glsadd{Pressure-Regulator}\dictentry{Pressure Regulator}{A Pressure Regulator reduces upstream high pressure to a constant lower downstream pressure regardless of flow or upstream variations. In rocket systems, they maintain precise pressures throughout pressurization, pneumatic, and propellant systems for proper operation. Poor performance can cause engine instability, feed problems, or structural overpressure failures. Designs must account for extreme temperatures, vibration, and specific fluid properties.}{catRM}
\glsadd{Gas-Generator}\dictentry{Gas Generator}{A Gas Generator is a small combustion device producing high-temperature gas to drive turbopumps or provide pressurization. It burns a small propellant amount at fuel-rich or oxidizer-rich ratios to produce required temperature and pressure gas. The gas generator cycle is one of the most common liquid engine cycles, balancing performance and complexity. Turbine exhaust may be discarded or routed to a secondary nozzle for additional thrust.}{catRM}
\glsadd{Heat-Exchanger}\dictentry{Heat Exchanger}{A Heat Exchanger transfers thermal energy between fluid streams without mixing them, used throughout rocket systems for thermal management. In engines, they preheat propellants, cool components, gasify pressurant, and manage thermal environments. Design balances heat transfer efficiency, pressure drop, weight, and reliability. Rocket heat exchangers operate under extreme thermal gradients for short but intense burn durations.}{catRM}
\glsadd{Gas-Cooler}\dictentry{Gas Cooler}{A Gas Cooler reduces hot gas temperature from a gas generator or combustion process before downstream use. Cooling protects turbine blades, seals, and components that cannot withstand full combustion temperatures. Gas coolers may use propellant as coolant, simultaneously conditioning gas and preheating propellant. Cooler effectiveness directly impacts turbopump and downstream hardware life and reliability.}{catRM}
\glsadd{Filter}\dictentry{Filter}{A Filter removes solid particles and contaminants from liquid or gas propellant streams to protect downstream components. Filters are placed throughout propellant systems, with critical locations including turbopump inlets, injector orifices, and valve seats. Filtration must be fine enough to catch harmful particles without excessive pressure drop. Filter integrity is verified during ground operations, as failure can lead to catastrophic engine damage.}{catRM}
\glsadd{Propellant-Management}\dictentry{Propellant Management}{Propellant Management encompasses systems ensuring liquid propellant is properly positioned and delivered to the engine in correct phase, quantity, and orientation. In microgravity, surface tension and wetting effects dominate over gravity. Management devices include baffles, sponges, galleries, and surface tension screens controlling propellant location. Effective management is essential for reliable engine restarts in space.}{catRM}
\glsadd{Propellant-Settling}\dictentry{Propellant Settling}{Propellant Settling uses small auxiliary thrusters to push liquid propellant toward the tank bottom in microgravity before engine restart. Without gravity, propellant floats freely and may not be at the engine inlet. Settling thrusters fire briefly to create acceleration pushing liquid toward the tank outlet. This technique is commonly used on upper stages and spacecraft that must restart engines after orbital coast periods.}{catRM}
\glsadd{Fin-Stabilized-Rocket}\dictentry{Fin-Stabilized Rocket}{A Fin-Stabilized Rocket uses aerodynamic fins on the rear to maintain flight stability through passive aerodynamic forces. Fins shift the center of pressure behind the center of gravity, creating a restoring moment. This is the simplest and most common stabilization form, used in everything from model rockets to tactical missiles. The limitation is that it only works within the atmosphere and adds drag during flight.}{catRM}
\glsadd{Spin-Stabilized-Rocket}\dictentry{Spin-Stabilized Rocket}{A Spin-Stabilized Rocket maintains stability by rotating around its longitudinal axis, using gyroscopic effects to resist perturbations. Spin is imparted by canted nozzles, angled fins, or separate spin motors. Spin stabilization is mechanically simple and doesn't require active guidance for short-range applications, ideal for artillery rockets and some tactical missiles. The tradeoff is that spin can interfere with seekers and guidance, limiting achievable accuracy.}{catRM}
\glsadd{Tail-Fins}\dictentry{Tail Fins}{Tail Fins are aerodynamic surfaces on a rocket's aft end providing passive flight stability through restoring forces when the vehicle deviates from its flight path. They shift the aerodynamic center of pressure behind the center of gravity, ensuring disturbances produce correcting moments. Fins may be fixed or moveable for steering, with size and shape carefully designed for adequate stability without excessive drag. They are essential for unguided rockets and many guided missiles.}{catRM}
\glsadd{Canard-Control}\dictentry{Canard Control}{Canard Control uses small movable surfaces near a missile's nose to generate steering forces. The forward placement provides rapid response to commands because the moment arm produces quick body rotation. Canard control is common in short-range guided and air-to-air missiles requiring fast maneuvering. The disadvantage is potential aerodynamic interference with main wings and reduced effectiveness at high angles of attack.}{catRM}
\glsadd{Guidance-Computer}\dictentry{Guidance Computer}{A Guidance Computer is the onboard flight computer responsible for navigation, attitude control, and trajectory execution in rockets and missiles. It processes sensor data from inertial measurement units and other instruments to compute thrust vectoring, staging events, and orbital insertion burns. The computer runs real-time control loops to keep the vehicle on its intended flight path. Modern versions integrate with ground tracking systems to allow mid-course corrections.}{catRM}
\glsadd{Launch-Escape-System}\dictentry{Launch Escape System}{A Launch Escape System is a safety mechanism designed to pull the crew capsule away from the rocket in the event of a catastrophic failure during ascent. It typically consists of solid-fuel escape motors mounted on a tower or integrated into the capsule's service module. When triggered, the system fires to accelerate the crew module to a safe distance, after which parachutes deploy for a controlled landing.}{catRM}
\glsadd{Abort-System}\dictentry{Abort System}{An Abort System encompasses all hardware and software that enable a rocket mission to be safely terminated at any point during flight. It includes sensors that detect failures, logic controllers that decide when to abort, and propulsion elements that separate the crew or payload from a failing vehicle. Abort modes range from pad escape to high-altitude separation depending on the flight phase.}{catRM}
\glsadd{Ablative-Shield}\dictentry{Ablative Shield}{An Ablative Shield is a heat-resistant covering applied to the nose cone or leading surfaces of a rocket or reentry vehicle. It is composed of materials that absorb thermal energy by charring and eroding in a controlled manner, carrying heat away from the vehicle structure. Ablative shields are simpler and lighter than active cooling systems, making them common on ballistic missile warheads and crew capsules.}{catRM}
\glsadd{Ballistic-Coefficient}\dictentry{Ballistic Coefficient}{The Ballistic Coefficient is a measure of a body's ability to overcome air resistance during flight, defined as mass divided by the product of drag coefficient and cross-sectional area. A high ballistic coefficient means the object retains velocity well in the atmosphere, which is critical for reentry vehicles and ballistic warheads. It directly influences terminal velocity, range, and impact accuracy.}{catRM}
\glsadd{Castor}\dictentry{Castor}{Castor refers to a family of solid rocket boosters originally developed by Thiokol and later produced by ATK and Northrop Grumman. Castor motors have served as strap-on boosters for rockets such as the Delta, Atlas, and Falcon series. Each Castor unit uses a high-performance solid propellant grain to provide additional thrust during launch. The design has been upgraded through multiple variants to meet evolving payload requirements.}{catRM}
\glsadd{Centaur}\dictentry{Centaur}{Centaur is a high-energy upper stage developed in the 1960s that uses liquid hydrogen and liquid oxygen propellants in a pressure-stabilized stainless steel structure. Its remarkable mass fraction made it one of the most efficient upper stages ever built. Centaur has been used atop Atlas and Titan rockets to place satellites into geostationary orbit and send probes on interplanetary missions. The stage pioneered the concept of applying thin-walled, balloon-type tank construction to space vehicles.}{catRM}
\glsadd{Launch-Clamp}\dictentry{Launch Clamp}{A Launch Clamp is a mechanical hold-down device that secures a rocket to the launch pad until full thrust is confirmed. It ensures the vehicle remains stationary during engine ignition and health checks, allowing the crew or computers to verify all systems before liftoff. Once conditions are met, pyrotechnic bolts or hydraulic mechanisms release the clamps and the rocket rises under its own power.}{catRM}
\glsadd{Movable-Nozzle}\dictentry{Movable Nozzle}{A Movable Nozzle is a thrust vector control mechanism in which the bell or nozzle extension of a rocket engine can be pivoted to redirect the exhaust plume. Actuators tilt the nozzle in response to guidance commands, producing lateral force components that steer the vehicle. This approach eliminates the need for separate vernier engines and is widely used on liquid-fueled launch vehicles and strategic missiles.}{catRM}
\glsadd{Retrofire}\dictentry{Retrofire}{Retrofire is the firing of retro-rockets oriented against the direction of flight to decelerate a spacecraft. The maneuver is essential for orbital deorbit burns, atmospheric entry alignment, and landing sequences. Retrofire timing must be precise, as even small errors can shift the landing point by hundreds of kilometers.}{catRM}
\glsadd{Stage}\dictentry{Stage}{A Stage is a self-contained section of a multistage rocket that includes its own engines, propellant tanks, and structural elements. When a stage exhausts its fuel, it is jettisoned to reduce the remaining vehicle mass, improving acceleration and efficiency. Multistaging allows rockets to achieve much higher velocities than single-stage designs by shedding dead weight progressively throughout flight.}{catRM}
\glsadd{Abort-Motor}\dictentry{Abort Motor}{An Abort Motor is a solid rocket motor within a launch escape system that provides the rapid thrust needed to pull a crew capsule away from a failing booster. It must produce extremely high thrust-to-weight ratios to achieve safe separation in fractions of a second. The motor burns for a short duration and is discarded once the capsule reaches a safe distance.}{catRM}
\glsadd{Attitude-Control-Motor}\dictentry{Attitude Control Motor}{An Attitude Control Motor is a small thruster or set of thrusters used to orient a spacecraft or missile in three axes during flight. It may use cold gas, monopropellant, or solid propellant to produce small torques that rotate the vehicle. Attitude control is critical for pointing sensors, aligning engines for burns, and maintaining aerodynamic stability during ascent.}{catRM}
\glsadd{Escape-Tower}\dictentry{Escape Tower}{An Escape Tower is a rigid structure mounted atop a crew capsule that houses the solid-fuel abort motors. In an emergency, the tower's motors ignite and pull the capsule away from the launch vehicle at high speed. After reaching a safe altitude and distance, the tower is jettisoned and the capsule's parachutes deploy for recovery.}{catRM}
\glsadd{Pusher-Rocket}\dictentry{Pusher Rocket}{A Pusher Rocket is a propulsion system mounted on the crew capsule itself rather than on a separate escape tower. When activated, it pushes the capsule away from a failing launch vehicle. This configuration eliminates the need for a jettisonable tower and simplifies the vehicle design, though it requires careful integration of the motors into the capsule's structure.}{catRM}
\glsadd{ullage-Motor}\dictentry{ullage Motor}{An ullage motor is a small auxiliary rocket used to settle propellants at the bottom of a fuel or oxidizer tank before main engine ignition. In microgravity, propellants float freely and can block feed lines, so the ullage motor produces a brief forward thrust to push fluids against the tank outlet. This ensures a gas-free liquid flow into the propulsion system at engine start.}{catRM}
\glsadd{ullage}\dictentry{ullage}{Ullage refers to the empty space above the liquid propellant inside a rocket tank. Managing ullage is critical because bubbles or vapor in the feed lines can cause engine instability or failure. In zero gravity, the propellant may drift away from the outlet, requiring ullage motors or diaphragms to ensure proper propellant delivery.}{catRM}
\glsadd{Propellant-Gauging}\dictentry{Propellant Gauging}{Propellant Gauging is the process of measuring the remaining quantity of fuel and oxidizer in a rocket's tanks during flight. Accurate gauging is essential for mission planning, burn timing, and determining when to shut down engines. Methods include capacitance probes, ultrasonic sensors, and mass flow integration, each chosen based on the propellant type and tank geometry.}{catRM}
\glsadd{Mass-Fraction}\dictentry{Mass Fraction}{Mass Fraction is the ratio of propellant mass to the total initial mass of a rocket stage or vehicle. A higher mass fraction indicates a more structurally efficient design, as more of the vehicle's weight is dedicated to propellant. It is a fundamental parameter in the Tsiolkovsky rocket equation and directly determines the achievable delta-v of a stage.}{catRM}
\glsadd{Nozzle-Efficiency}\dictentry{Nozzle Efficiency}{Nozzle Efficiency, also called characteristic velocity efficiency, measures how effectively a rocket nozzle converts the thermal energy of combustion gases into directed kinetic energy. It depends on the nozzle contour, expansion ratio, and the uniformity of gas flow at the exit plane. Nozzle design optimization is a key factor in maximizing overall engine performance and specific impulse.}{catRM}
\glsadd{Ignition-Overpressure}\dictentry{Ignition Overpressure}{Ignition Overpressure is a sudden pressure spike that occurs when a rocket engine ignites, caused by the rapid expansion of combustion gases inside the combustion chamber. If uncontrolled, it can damage engine components and propagate acoustic vibrations through the vehicle structure. Engineers mitigate it using pre-ignition purges, chamber ventilation, and gradual propellant valve sequencing.}{catRM}
\glsadd{Hard-Start}\dictentry{Hard Start}{A Hard Start is an abnormal engine ignition event characterized by an excessive pressure spike or detonation within the combustion chamber. It is usually caused by propellant pooling before ignition, leading to an uncontrolled combustion event. A hard start can cause severe damage to engine hardware and is one of the most dangerous anomalies in rocket propulsion.}{catRM}
\glsadd{Two-Combustion-Instability}\dictentry{Two-Combustion Instability}{Two-Combustion Instability is a condition in which a rocket engine's combustion process oscillates between two distinct operating modes, typically low-frequency and high-frequency pressure oscillations. The instability can cause thermal stresses, mechanical vibrations, and thrust fluctuations that threaten structural integrity. It is mitigated through injector design changes, baffles, and acoustic cavities.}{catRM}
\glsadd{Acoustic-Mode}\dictentry{Acoustic Mode}{An Acoustic Mode is a resonant vibration pattern within a rocket combustion chamber driven by the interaction between gas pressure waves and the chamber geometry. These modes can couple with structural vibrations to produce destructive oscillations known as combustion instability. Identifying and suppressing acoustic modes is essential for reliable engine operation.}{catRM}
\glsadd{Grain-Burn-Rate}\dictentry{Grain Burn Rate}{The Grain Burn Rate is the speed at which the surface of a solid rocket propellant grain recedes during combustion, typically measured in millimeters per second. It depends on the propellant formulation, oxidizer particle size, and combustion chamber pressure. The burn rate directly determines the thrust profile and must be carefully controlled through grain geometry and propellant chemistry.}{catRM}
\glsadd{Web-Thickness}\dictentry{Web Thickness}{Web Thickness in solid rocketry refers to the minimum distance from the propellant grain surface to the case wall or the next burn surface. It determines the total burn duration of a grain segment, as the flame front must travel through the entire web before that section is consumed. Accurate web design is critical for controlling thrust duration and preventing burn-through to the motor case.}{catRM}
\glsadd{Burn-Rate-Exponent}\dictentry{Burn Rate Exponent}{The Burn Rate Exponent, often denoted as n, describes how the burn rate of a solid propellant grain responds to changes in chamber pressure. A value near zero means the burn rate is largely pressure-insensitive, while higher values indicate stronger pressure dependence. The exponent is a critical parameter in motor stability analysis and grain design.}{catRM}
\glsadd{Neutral-Burn}\dictentry{Neutral Burn}{A Neutral Burn is a solid propellant grain design in which the burning surface area remains approximately constant throughout the motor's firing duration. This produces a relatively steady thrust level from ignition to burnout. Neutral grains are used when a consistent thrust profile is required, such as in booster stages and interceptor missiles.}{catRM}
\glsadd{Progressive-Burn}\dictentry{Progressive Burn}{A Progressive Burn is a solid propellant grain configuration in which the burning surface area increases over time, producing rising thrust. This is achieved by shaping the grain so that the flame front exposes more area as it recedes. Progressive burns are useful in applications where increasing thrust helps overcome growing aerodynamic loads during ascent.}{catRM}
\glsadd{Regressive-Burn}\dictentry{Regressive Burn}{A Regressive Burn is a solid propellant grain geometry where the burning surface area decreases over time, resulting in declining thrust. This occurs when the grain shape causes the flame front to shrink as it progresses. Regressive profiles can help limit structural loads and are sometimes combined with other burn patterns for tailored thrust curves.}{catRM}
\glsadd{Sliver}\dictentry{Sliver}{In solid rocketry, a sliver is a thin remnant of unburned propellant left against the motor case wall after the main grain has been consumed. Slivers represent wasted mass and reduce the motor's overall efficiency. Minimizing sliver mass is a key objective in grain geometry optimization and motor design.}{catRM}
\glsadd{Void-Fraction}\dictentry{Void Fraction}{Void Fraction is the ratio of unoccupied volume to total volume within a solid rocket motor chamber or a packed bed reactor. In rocket propulsion, it affects the internal ballistic behavior and pressure distribution during combustion. A controlled void fraction can help manage gas flow and combustion stability within the motor.}{catRM}
\glsadd{Free-Volume}\dictentry{Free Volume}{Free Volume is the open space within a solid rocket motor chamber that is not occupied by propellant grains, insulation, or other hardware. It allows combustion gases to flow between grain segments and toward the nozzle. The amount of free volume influences chamber pressure dynamics and must be carefully sized during motor design.}{catRM}
\glsadd{Impinging-Jet}\dictentry{Impinging Jet}{An Impinging Jet injector is a rocket engine injector design where separate streams of fuel and oxidizer collide at a designated point to achieve rapid atomization and mixing. The collision breaks the liquid streams into fine droplets that combust efficiently. Impinging jet injectors are valued for their simplicity and effectiveness in smaller rocket engines.}{catRM}
\glsadd{Coaxial-Swirl}\dictentry{Coaxial Swirl}{A Coaxial Swirl injector is a rocket engine injector configuration in which one propellant flows coaxially around the other while being given a tangential velocity component. The swirling motion enhances atomization and mixing of fuel and oxidizer within the combustion chamber. This design is widely used in large liquid-fueled engines for its reliable and uniform combustion characteristics.}{catRM}
\glsadd{Turbopump-Assembly}\dictentry{Turbopump Assembly}{A Turbopump Assembly is a high-speed rotating machine that delivers fuel and oxidizer from low-pressure tanks to the high-pressure combustion chamber of a liquid rocket engine. It consists of a turbine driven by hot gases connected to centrifugal pumps for each propellant. The assembly enables engines to operate at extremely high chamber pressures, which directly improves specific impulse and thrust density.}{catRM}
\glsadd{Oxidizer-Turbine}\dictentry{Oxidizer Turbine}{The Oxidizer Turbine is the turbine stage dedicated to driving the oxidizer pump within a turbopump assembly. It is typically powered by hot gas from a gas generator or preburner and must operate reliably in an oxidizer-rich environment. Its design must prevent ignition or detonation while maintaining high rotational speeds.}{catRM}
\glsadd{Fuel-Turbine}\dictentry{Fuel Turbine}{The Fuel Turbine is the turbine stage that powers the fuel pump in a turbopump assembly. Driven by hot combustion gases, it spins at tens of thousands of revolutions per minute to deliver fuel at high pressure to the combustion chamber. Fuel turbine design must balance power output, weight, and resistance to the thermal environment of the drive gas.}{catRM}
\glsadd{Staged-Combustion}\dictentry{Staged Combustion}{Staged Combustion is a rocket engine cycle in which propellants are partially burned in a preburner to drive the turbopump, and the exhaust products are then routed to the main combustion chamber for full combustion. This approach allows very high chamber pressures and improved specific impulse compared to simpler cycles. The trade-off is significantly increased engine complexity.}{catRM}
\glsadd{Gas-generator-Cycle}\dictentry{Gas-generator Cycle}{The Gas-generator Cycle is a rocket engine cycle in which a small fraction of propellant is burned in a gas generator to drive the turbopump, and the resulting exhaust is dumped overboard rather than re-entering the main combustion chamber. It is mechanically simpler than staged combustion but somewhat less efficient due to the wasted exhaust energy. Many reliable engines, including the F-1 and Merlin, use this cycle.}{catRM}
\glsadd{Full-Flow-Staged-Combustion}\dictentry{Full-Flow Staged Combustion}{Full-Flow Staged Combustion is an advanced rocket engine cycle in which all of the fuel passes through one preburner and all of the oxidizer through another, producing hot gas streams that drive both turbopumps before entering the main chamber. This design achieves the highest possible chamber pressures and specific impulse. It also reduces turbine temperatures since both propellants are present in excess.}{catRM}
\glsadd{Oxidizer-Rich-Staged-Combustion}\dictentry{Oxidizer-Rich Staged Combustion}{Oxidizer-Rich Staged Combustion is a variant of the staged combustion cycle in which the preburner operates with excess oxidizer, producing a hot, oxygen-rich gas to drive the turbines. This approach avoids the need for a separate fuel-rich preburner and simplifies plumbing, but requires materials that resist oxidation at extreme temperatures. The Soviet RD-170 engine family famously employed this technique.}{catRM}
\glsadd{Pressure-Fed-Engine}\dictentry{Pressure-Fed Engine}{A Pressure-Fed Engine is a rocket propulsion system in which high-pressure inert gas forces propellants from the tanks into the combustion chamber without the need for turbopumps. This design is mechanically simple and highly reliable, making it suitable for spacecraft thrusters and upper stages. However, the lower chamber pressures limit achievable specific impulse compared to pump-fed engines.}{catRM}
\glsadd{Monopropellant-Thruster}\dictentry{Monopropellant Thruster}{A Monopropellant Thruster is a spacecraft propulsion device that decomposes a single propellant, such as hydrazine, over a catalyst bed to produce thrust without requiring a separate oxidizer. The simplicity of the design makes it reliable for attitude control and small orbital maneuvers. It offers moderate specific impulse but lacks the performance of bipropellant systems.}{catRM}
\glsadd{Bipropellant-Thruster}\dictentry{Bipropellant Thruster}{A Bipropellant Thruster uses two separate propellants, a fuel and an oxidizer, that ignite either hypergolically or with an igniter to produce thrust. This configuration provides higher specific impulse than monopropellant designs, making it more efficient for orbital maneuvers. Bipropellant thrusters are used on crewed spacecraft, satellites, and interplanetary probes.}{catRM}
\glsadd{Green-Propellant}\dictentry{Green Propellant}{Green Propellant refers to next-generation rocket and spacecraft propellants formulated to be less toxic and less carcinogenic than traditional hydrazine-based fuels. Examples include AF-M315E and LMP-103S, which offer comparable or superior performance to hydrazine with significantly reduced handling hazards. Adoption of green propellants lowers operational costs and environmental impact.}{catRM}
\glsadd{Ion-Thruster}\dictentry{Ion Thruster}{An Ion Thruster is an electric propulsion device that accelerates ionized propellant through an electric field to produce thrust with extremely high specific impulse. While thrust levels are very low compared to chemical rockets, ion engines can operate continuously for months or years. They are ideal for deep-space missions and station-keeping of geostationary satellites.}{catRM}
\glsadd{Hall-Effect-Thruster}\dictentry{Hall Effect Thruster}{A Hall Effect Thruster is a type of ion thruster that uses a magnetic field to trap electrons and create an electric field that accelerates ionized propellant. The trapped electrons ionize the propellant gas, and the resulting ions are expelled at high velocity. Hall thrusters offer a balance of moderate thrust and high efficiency, making them popular for satellite station-keeping and interplanetary missions.}{catRM}
\glsadd{Pulsed-Plasma-Thruster}\dictentry{Pulsed Plasma Thruster}{A Pulsed Plasma Thruster is an electric propulsion device that generates thrust by discharging a capacitor across a solid propellant surface, creating a plasma arc that is accelerated by electromagnetic forces. Each pulse produces a small impulse, and pulses are repeated at high frequency. The design is simple and robust, suitable for small satellite attitude control and orbit maintenance.}{catRM}
\glsadd{Arcjet}\dictentry{Arcjet}{An Arcjet is an electrothermal propulsion device in which an electric arc heats the propellant gas to very high temperatures before it expands through a nozzle. The added thermal energy gives arcjets higher specific impulse than purely chemical thrusters. They have been used extensively for station-keeping and orbit raising of communication satellites.}{catRM}
\glsadd{Resistojet}\dictentry{Resistojet}{A Resistojet is an electrothermal thruster that uses an electric resistance heating element to warm the propellant gas before it exits through a converging-diverging nozzle. It is one of the simplest forms of electric propulsion and provides modest performance improvements over cold gas thrusters. Resistojets are used on small satellites for fine attitude control and drag compensation.}{catRM}
\glsadd{Variable-Specific-Impulse-Magnetoplasma-Rocket}\dictentry{Variable Specific Impulse Magnetoplasma Rocket}{The Variable Specific Impulse Magnetoplasma Rocket, or VASIMR, is an experimental electric propulsion engine that uses radio waves to ionize and heat a propellant, then a magnetic nozzle to direct the plasma exhaust. Its unique feature is the ability to vary the specific impulse and thrust level in flight by adjusting the power input. This flexibility could optimize deep-space transit times for crewed missions.}{catRM}
\glsadd{Laser-Propulsion}\dictentry{Laser Propulsion}{Laser Propulsion is a concept in which external laser beams are used to heat a propellant or ablate a surface to generate thrust, eliminating the need to carry onboard energy sources. Ground-based or orbiting laser arrays could beam energy to the vehicle, enabling very high specific impulse. The concept is still experimental but holds promise for reducing launch costs.}{catRM}
\glsadd{Nuclear-Thermal-Propulsion}\dictentry{Nuclear Thermal Propulsion}{Nuclear Thermal Propulsion uses a fission reactor to heat a propellant such as liquid hydrogen to extremely high temperatures, producing thrust with roughly twice the specific impulse of chemical engines. NTP systems are of great interest for crewed Mars missions because they can significantly reduce transit times. Development programs have demonstrated feasibility, though engineering and safety challenges remain.}{catRM}
\glsadd{Nuclear-Electric-Propulsion}\dictentry{Nuclear Electric Propulsion}{Nuclear Electric Propulsion generates electricity from a nuclear reactor to power electric thrusters such as ion or Hall engines. The combination provides high specific impulse with the ability to operate independent of solar proximity. NEP is particularly attractive for outer-planet missions where solar power is insufficient, though it requires significant mass for the reactor and power conversion systems.}{catRM}
\glsadd{Antimatter-Propulsion}\dictentry{Antimatter Propulsion}{Antimatter Propulsion is a theoretical concept in which matter and antimatter annihilate each other to release energy far exceeding chemical or nuclear reactions. The annihilation products would be directed to produce thrust. While the energy density is unparalleled, producing and storing antimatter in useful quantities remains far beyond current technological capabilities.}{catRM}
\glsadd{Tether-Propulsion}\dictentry{Tether Propulsion}{Tether Propulsion uses long conductive cables deployed from a spacecraft to interact with a planetary magnetic field, generating thrust or electrical power through electromagnetic induction. A conducting tether moving through a magnetic field produces a voltage that can be used to drive current and create force. The concept has been demonstrated in low-Earth orbit but faces practical challenges with deployment and orbital debris.}{catRM}
\glsadd{Mass-Driver}\dictentry{Mass Driver}{A Mass Driver is an electromagnetic launcher that accelerates payloads along a track using sequential magnetic coils, similar in principle to a linear motor. It could be used on the Moon or in space to launch materials into orbit without conventional rockets. The absence of atmospheric drag and the availability of solar power make the Moon a particularly attractive location for such a system.}{catRM}
\glsadd{Launch-Loop}\dictentry{Launch Loop}{A Launch Loop is a proposed electromagnetic launch system consisting of a long, thin structure suspended at high altitude by the momentum of continuously circulated iron pellets. Vehicles would be magnetically accelerated along the track to orbital velocity. The concept avoids the need for a space elevator while providing continuous launch capability, though it remains purely theoretical.}{catRM}
\glsadd{Space-Elevator}\dictentry{Space Elevator}{A space elevator is a theoretical structure extending from a planetary surface to geostationary orbit via a tether, enabling payload transport without rockets. It demands a material with extraordinary tensile strength, like carbon nanotubes. A functioning elevator could dramatically reduce launch costs by replacing chemical rockets with electric climbers. Materials research continues advancing feasibility.}{catRM}
\glsadd{Orbital-Tug}\dictentry{Orbital Tug}{An Orbital Tug is a reusable spacecraft designed to transport payloads between different orbits, such as from a low-Earth parking orbit to geostationary orbit or beyond. It fills the gap between launch vehicle upper stages and final mission orbits. By operating repeatedly in space, an orbital tug can reduce overall mission costs and extend payload flexibility.}{catRM}
\glsadd{Cislunar-Transport}\dictentry{Cislunar Transport}{Cislunar Transport refers to the infrastructure and vehicles needed to move cargo and crew between Earth orbit and the Moon's vicinity. It includes transfer stages, fuel depots, and navigation systems that operate in the Earth-Moon cislunar space. Establishing cislunar transport networks is seen as a critical step toward sustained lunar exploration and eventual Mars missions.}{catRM}
\end{multicols}

\clearpage
\chapter*{Drones}
\addcontentsline{toc}{chapter}{Drones}
\markboth{Drones}{Drones}
\clearpage
\begin{multicols}{2}
\small
Unmanned aerial vehicles: multirotor and fixed-wing drones, flight controllers, autonomy, regulations, and the technology powering the UAV revolution.\par
\vspace{0.5cm}
\glsadd{Drone}\dictentry{Drone}{A Drone is an unmanned aircraft that can fly either autonomously or under remote human control. It uses rotors, wings, or a hybrid system to generate lift, and carries sensors, cameras, or payloads for a wide variety of missions. Drones range in size from palm-sized nano platforms to large fixed-wing systems with wingspans exceeding several meters. They have transformed industries including agriculture, logistics, cinematography, defense, and infrastructure inspection.}{catD}
\glsadd{RPAS}\dictentry{RPAS}{RPAS stands for Remotely Piloted Aircraft System and is the term adopted by the International Civil Aviation Organization for civil unmanned aviation. It emphasizes that a human remote pilot maintains authority over the aircraft at all times. The RPAS framework distinguishes these systems from fully autonomous drones by requiring continuous human oversight. This terminology aligns with international aviation regulations and air traffic management standards.}{catD}
\glsadd{Multirotor}\dictentry{Multirotor}{A Multirotor is a drone configuration using multiple rotors, most commonly four, six, or eight, to achieve vertical flight. Each rotor spins at variable speeds controlled by the flight computer to manage pitch, roll, yaw, and throttle. Multirotors are valued for their mechanical simplicity, as they have no complex swashplate mechanisms like helicopters. They excel at hovering, slow-speed maneuvering, and short-range missions such as photography and inspection.}{catD}
\glsadd{Quadcopter}\dictentry{Quadcopter}{A Quadcopter is a multirotor drone with exactly four rotors arranged in a cross or H configuration. It achieves flight control by varying the speed of each rotor, with diagonal pairs rotating in opposite directions to cancel torque. The quadcopter is the most popular consumer and commercial drone format due to its simple design and reliable performance. Its compact footprint and ease of use have made it the standard platform for photography, mapping, and recreational flying.}{catD}
\glsadd{Hexacopter}\dictentry{Hexacopter}{A Hexacopter is a multirotor drone with six rotors mounted on a central frame. The additional two rotors compared to a quadcopter provide greater lift capacity and motor redundancy. If one motor fails, the hexacopter can often maintain controlled flight and land safely, which is critical for carrying expensive payloads. This configuration is popular for professional aerial photography, industrial inspection, and light cargo delivery.}{catD}
\glsadd{Octocopter}\dictentry{Octocopter}{An Octocopter is a multirotor drone equipped with eight rotors for maximum redundancy and payload capacity. With eight independent propulsion units, it can tolerate the failure of one or even two motors while still maintaining controlled flight. The octocopter can carry heavier camera rigs, sensors, and delivery packages than smaller multirotors. It is commonly used in professional filmmaking, surveying, and commercial applications where reliability is paramount.}{catD}
\glsadd{Fixed-Wing-Drone}\dictentry{Fixed-Wing Drone}{A Fixed-Wing Drone uses rigid wings to generate lift through forward airspeed rather than relying on rotors. This design is far more aerodynamically efficient than multirotors, enabling longer flight times and greater range on a single battery charge. Fixed-wing drones cannot hover in place, but they can cover hundreds of kilometers in mapping, surveying, and surveillance missions. They typically launch by hand toss or catapult and recover via belly landing or parachute.}{catD}
\glsadd{VTOL-Drone}\dictentry{VTOL Drone}{A VTOL Drone combines vertical takeoff and landing capability with the efficient forward flight of a fixed-wing aircraft. It uses rotors to lift off vertically, then transitions to wing-borne flight for long-range cruising. This hybrid approach eliminates the need for a runway while retaining the endurance advantages of fixed-wing designs. VTOL drones are increasingly used for delivery, surveillance, and long-range inspection where both flexibility and efficiency are required.}{catD}
\glsadd{Hybrid-Drone}\dictentry{Hybrid Drone}{A Hybrid Drone integrates both multirotor and fixed-wing configurations into a single airframe. It can hover and maneuver like a multirotor using its vertical rotors, then transition to efficient wing-borne flight for covering long distances. The transition between flight modes is managed automatically by the flight controller using coordinated motor and control surface adjustments. This versatility makes hybrid drones ideal for missions that demand both precise station-keeping and extended range.}{catD}
\glsadd{Nano-Drone}\dictentry{Nano Drone}{A Nano Drone is an extremely small unmanned aircraft, typically weighing under 250 grams. These micro-scale platforms often fit in the palm of a hand and use miniature brushless or brushed motors with small propellers. Despite their size, nano drones can carry basic cameras and are used for indoor exploration, educational projects, and recreational flying. Their low weight often places them in the lightest regulatory category, requiring fewer operational restrictions.}{catD}
\glsadd{Micro-Drone}\dictentry{Micro Drone}{A Micro Drone is a small unmanned aircraft generally weighing under 2 kilograms. It offers more capability than a nano drone, including longer flight times, better cameras, and more stable GPS-assisted flight. Micro drones bridge the gap between hobby-grade toys and professional platforms, making them accessible for education and light commercial work. They are commonly used for indoor inspection, real estate photography, and learning to fly unmanned aircraft.}{catD}
\glsadd{Tactical-Drone}\dictentry{Tactical Drone}{A Tactical Drone is a mid-range unmanned system designed for military and public safety operations at the battalion or regiment level. It provides real-time situational awareness, target identification, and battlefield surveillance to ground commanders. Tactical drones typically operate at ranges of 10 to 50 kilometers with endurance of several hours. They are designed for rapid deployment, often launched by hand or small catapult in forward operating areas.}{catD}
\glsadd{MALE-Drone}\dictentry{MALE Drone}{A MALE Drone operates at medium altitudes between 10,000 and 30,000 feet with endurance measured in many hours. The Medium Altitude Long Endurance designation reflects its ability to remain on station for extended surveillance or strike missions. MALE platforms like the MQ-1 Predator and MQ-9 Reaper have become central to modern military operations. They carry electro-optical and infrared sensors along with precision-guided munitions for persistent intelligence and strike roles.}{catD}
\glsadd{HALE-Drone}\dictentry{HALE Drone}{A HALE Drone flies at high altitudes above 30,000 feet and can remain airborne for more than 24 hours. The High Altitude Long Endurance category includes platforms like the RQ-4 Global Hawk that operate at the edge of the stratosphere. These drones provide wide-area surveillance using advanced radar, signals intelligence, and optical sensors. Their extreme endurance and altitude allow a single platform to monitor vast geographic areas continuously.}{catD}
\glsadd{Combat-Drone}\dictentry{Combat Drone}{A Combat Drone is an armed unmanned aircraft designed to carry and deploy precision-guided munitions against ground targets. It combines long-endurance flight capability with targeting sensors and weapons systems typically including Hellfire missiles or guided bombs. Combat drones allow military forces to conduct strike operations without risking human pilots in the cockpit. They operate under strict rules of engagement with operators controlling weapons release from ground stations.}{catD}
\glsadd{Reconnaissance-Drone}\dictentry{Reconnaissance Drone}{A Reconnaissance Drone is an unmanned aircraft optimized for intelligence gathering and surveillance. It carries electro-optical cameras, infrared sensors, and sometimes synthetic aperture radar to observe ground activity from safe distances. These drones provide critical battlefield intelligence, border monitoring, and disaster assessment without exposing personnel to danger. They range from small hand-launched systems to large strategic platforms capable of monitoring entire regions.}{catD}
\glsadd{Swarm-Drone}\dictentry{Swarm Drone}{A Swarm Drone is designed to operate as part of a coordinated group that communicates and collaborates autonomously. Swarm technology enables dozens or hundreds of small drones to share sensor data and coordinate movements without centralized control. Each drone in the swarm makes local decisions based on rules inspired by biological systems like bird flocks or insect colonies. Swarm drones offer advantages in surveillance, search and rescue, and saturation operations where many platforms are more effective than one.}{catD}
\glsadd{Delivery-Drone}\dictentry{Delivery Drone}{A Delivery Drone is an unmanned aircraft designed to transport parcels, medical supplies, or small cargo from distribution centers to customers. It uses GPS navigation, obstacle avoidance sensors, and parachutes or winches to safely drop off packages at designated locations. Delivery drones are being developed by major logistics companies to reduce last-mile shipping costs and delivery times. Regulatory frameworks are evolving to enable routine beyond-visual-line-of-sight delivery operations over populated areas.}{catD}
\glsadd{Agriculture-Drone}\dictentry{Agriculture Drone}{An Agriculture Drone is an unmanned platform optimized for precision farming tasks including crop monitoring, spraying, and seeding. It carries multispectral cameras assessing plant health by measuring near-infrared reflectance, along with liquid tanks for targeted pesticide or fertilizer application. Agriculture drones can survey hundreds of acres per day, providing detailed maps of crop stress, irrigation needs, and pest infestations. This technology enables variable-rate application that reduces chemical use while improving yields.}{catD}
\glsadd{Surveying-Drone}\dictentry{Surveying Drone}{A Surveying Drone is equipped with high-resolution cameras and sometimes LiDAR sensors to create accurate topographic maps and 3D models. It flies pre-programmed grid patterns over terrain, capturing overlapping images that photogrammetry software stitches into orthomosaics and digital elevation models. Surveying drones have dramatically reduced the time and cost of land surveying compared to traditional ground-based methods. They are used in construction, mining, urban planning, and environmental monitoring.}{catD}
\glsadd{Inspection-Drone}\dictentry{Inspection Drone}{An Inspection Drone carries cameras and sensors to examine infrastructure that is difficult, dangerous, or expensive to access by human crews. It can inspect power lines, wind turbines, bridges, pipelines, and building facades at close range while recording detailed imagery. The drone operator reviews footage in real time or after the flight to identify structural damage, corrosion, or wear. This approach significantly reduces inspection costs and eliminates the need for scaffolding, rope access, or helicopter fly-bys.}{catD}
\glsadd{Search-and-Rescue-Drone}\dictentry{Search and Rescue Drone}{A Search and Rescue Drone carries thermal cameras and other sensors to locate missing persons in wilderness, disaster zones, or maritime environments. It can cover large search areas much faster than ground teams, using infrared imaging to detect body heat through darkness or light vegetation. Some SAR drones can drop emergency supplies like radios, blankets, or flotation devices to people in distress. The technology has proven invaluable in time-critical missions where rapid location of survivors is essential.}{catD}
\glsadd{Racing-Drone}\dictentry{Racing Drone}{A Racing Drone is a high-performance multirotor built for speed, agility, and competitive first-person-view racing events. It uses lightweight carbon-fiber frames, high-KV brushless motors, and aggressive-pitch propellers to achieve speeds exceeding 160 kilometers per hour. Pilots fly these drones through complex three-dimensional courses using goggles that display a live video feed from an onboard camera. Racing drones have spawned an organized competitive sport with professional leagues and international championships.}{catD}
\glsadd{FPV-Drone}\dictentry{FPV Drone}{An FPV Drone is piloted using a First Person View system where the pilot wears video goggles displaying a live feed from the drone's onboard camera. This immersive perspective gives the pilot a cockpit-like experience, enabling precise and dynamic maneuvering through complex environments. FPV drones use low-latency analog or digital video transmitters to minimize the delay between the drone's movement and what the pilot sees. The technology has applications in racing, cinematography, inspection, and recreational flying.}{catD}
\glsadd{Cinematic-Drone}\dictentry{Cinematic Drone}{A Cinematic Drone is an unmanned platform designed specifically for professional filmmaking and broadcast-quality aerial video. It carries large gimbal-stabilized cameras capable of 4K, 6K, or higher resolution with cinema-grade lenses. These drones feature smooth flight characteristics and advanced software modes enabling precise, repeatable camera movements like orbit shots, reveals, and tracking sequences. Cinematic drones have replaced helicopter-mounted systems for most aerial shots in modern film and television.}{catD}
\glsadd{Submersible-Drone}\dictentry{Submersible Drone}{A Submersible Drone is an unmanned underwater vehicle designed for marine inspection, research, and exploration. It uses thrusters and control surfaces to navigate through water, carrying cameras, sonar, and sampling equipment to depths inaccessible to human divers. Submersible drones inspect ship hulls, underwater pipelines, coral reefs, and hydrothermal vents without risking human life. They transmit video and sensor data in real time through a tether or acoustic communication link to operators on the surface.}{catD}
\glsadd{Autonomous-Drone}\dictentry{Autonomous Drone}{An Autonomous Drone is capable of planning and executing missions with minimal or no human intervention. It uses onboard artificial intelligence, computer vision, and advanced navigation algorithms to perceive its environment and make real-time decisions. Autonomous drones can identify obstacles, select landing sites, and adapt flight paths without pilot input. This capability is essential for operations at scale, such as large-area mapping, delivery networks, and persistent surveillance missions.}{catD}
\glsadd{Tethered-Drone}\dictentry{Tethered Drone}{A Tethered Drone is connected to a ground station by a thin cable that provides continuous electrical power and high-bandwidth data transfer. The tether eliminates battery limitations, allowing the drone to remain airborne for hours or even days in a fixed location. It is commonly used for persistent surveillance, temporary communication relay, and event coverage where long-duration loitering is needed. The physical tether also provides a security advantage, as the drone cannot fly beyond a defined radius.}{catD}
\glsadd{Solar-Drone}\dictentry{Solar Drone}{A Solar Drone harvests energy from sunlight using photovoltaic cells mounted on its wings or airframe to extend flight endurance. During daylight, solar panels charge onboard batteries that sustain flight during cloud cover and twilight. Solar-powered drones have achieved flight durations of many hours, approaching perpetual flight in favorable conditions. They are used for atmospheric research, telecommunications relay, and persistent wide-area surveillance missions.}{catD}
\glsadd{Flight-Controller}\dictentry{Flight Controller}{A Flight Controller is the onboard computer that processes sensor data and pilot commands to stabilize and direct a drone's movement. It runs PID control loops that calculate precise motor speeds needed to achieve the desired attitude and position. The flight controller reads data from gyroscopes, accelerometers, magnetometers, barometers, and GPS to maintain awareness of the drone's state. It is the central processing unit of any unmanned aircraft and determines overall flight characteristics and responsiveness.}{catD}
\glsadd{IMU}\dictentry{IMU}{An IMU, or Inertial Measurement Unit, is a sensor module combining accelerometers and gyroscopes to measure a drone's linear acceleration and angular velocity. It provides the foundational inertial data that the flight controller uses to determine orientation and motion in three-dimensional space. Modern drone IMUs often include magnetometers for heading reference and are built as micro-electromechanical systems on a single chip. The quality and sampling rate of the IMU directly influence the stability and precision of drone flight.}{catD}
\glsadd{Gyroscope}\dictentry{Gyroscope}{A gyroscope measures spacecraft rotation rates by exploiting conservation of angular momentum. Mechanical gyroscopes use spinning rotors while fiber-optic and ring-laser gyroscopes use the Sagnac effect with no moving parts. Gyroscope data is integrated over time to track attitude changes, though small biases cause drift. They are combined with star tracker observations in Kalman filters for continuous accurate attitude estimates.}{catD}
\glsadd{Magnetometer}\dictentry{Magnetometer}{A magnetometer measures magnetic field strength and direction in space using fluxgate, Hall-effect, or cesium vapor technologies. It provides attitude information by comparing measured field direction to reference models and conducts scientific studies of planetary magnetic fields.}{catD}
\glsadd{Barometer}\dictentry{Barometer}{A Barometer measures atmospheric pressure, which the flight controller converts to an estimated altitude using a standard pressure-altitude model. Small changes in air pressure correspond to changes in height, allowing the drone to maintain a steady hover altitude without GPS. Barometric sensors in drones typically have resolution sufficient to detect altitude changes of a few centimeters. This sensor is critical for indoor flight and altitude hold modes where GPS may be unavailable or insufficiently accurate.}{catD}
\glsadd{GPS-Module}\dictentry{GPS Module}{A GPS Module receives signals from the United States Global Positioning System satellite constellation to compute the drone's three-dimensional position, velocity, and time. It provides the flight controller with latitude, longitude, and altitude data accurate to within a few meters under normal conditions. GPS data enables position hold, waypoint navigation, return-to-home, and geofencing capabilities. Modern drone GPS modules often support multiple satellite constellations simultaneously for improved accuracy and reliability.}{catD}
\glsadd{GLONASS}\dictentry{GLONASS}{GLONASS is the Global Navigation Satellite System operated by the Russian Federation, providing an alternative to GPS. It consists of 24 satellites in medium Earth orbit transmitting positioning signals that drone receivers can use for navigation. Many modern drone GPS modules support both GPS and GLONASS simultaneously, increasing visible satellites and improving position accuracy. Dual-constellation reception provides better performance in urban canyons and at high latitudes.}{catD}
\glsadd{Compass}\dictentry{Compass}{A Compass in a drone context is a magnetic heading sensor, typically a three-axis magnetometer, that determines the aircraft's orientation relative to the Earth's magnetic field. It provides the yaw reference that the flight controller needs for directional navigation and stable flight. The compass must be mounted away from electromagnetic interference from motors, ESCs, and batteries to maintain accuracy. Regular compass calibration is required before flight to account for local magnetic variations and vehicle-specific interference patterns.}{catD}
\glsadd{ESC}\dictentry{ESC}{An ESC, or Electronic Speed Controller, is an electronic device that regulates the rotational speed of a brushless motor based on commands from the flight controller. It converts DC battery power into three-phase AC signals, switching the motor windings at precise timing to achieve the desired RPM. Each motor on a multirotor has its own ESC, allowing independent speed control for attitude management. ESCs also provide braking, reverse, and telemetry functions, and their response time directly affects flight performance.}{catD}
\glsadd{Motor}\dictentry{Motor}{A Motor on a drone converts electrical energy from the battery into rotational mechanical energy to drive the propeller. Electric brushless motors are the standard choice for modern drones due to their high efficiency, reliability, and favorable power-to-weight ratio. Motor specifications include KV rating, which indicates RPM per volt under no load, and physical size denoted by stator dimensions. The selection of motor size and KV must be matched to the propeller and battery to achieve optimal flight performance.}{catD}
\glsadd{Brushless-Motor}\dictentry{Brushless Motor}{A Brushless Motor uses electronic commutation instead of physical brushes to drive rotation, resulting in higher efficiency and longer lifespan. The motor controller energizes the stator windings in sequence, creating a rotating magnetic field that pulls the permanent-magnet rotor around. Brushless motors generate less heat, produce less electromagnetic interference, and require virtually no maintenance compared to brushed alternatives. They are the dominant motor type in modern drones, offering excellent torque, responsiveness, and durability.}{catD}
\glsadd{Pitch-Propeller}\dictentry{Pitch Propeller}{A Pitch Propeller refers to a propeller blade designed with a specific angle of attack that determines how aggressively it bites into the air. Higher pitch propellers move more air per revolution, generating greater thrust at higher speeds but drawing more power. Lower pitch propellers are more efficient at lower speeds and provide better hovering performance. The pitch rating is measured in inches and represents the theoretical forward distance the propeller would travel in one revolution.}{catD}
\glsadd{Yaw-Propeller}\dictentry{Yaw Propeller}{Yaw Propellers in a multirotor are counter-rotating pairs that cancel the reactive torque that would otherwise cause the drone to spin. On a quadcopter, two motors spin clockwise and two spin counterclockwise to maintain yaw stability. By varying the speed differential between clockwise and counterclockwise pairs, the flight controller can command yaw rotation without affecting altitude. This arrangement is fundamental to multirotor flight dynamics and eliminates the need for a tail rotor like a conventional helicopter.}{catD}
\glsadd{LiPo-Battery}\dictentry{LiPo Battery}{A LiPo Battery, or Lithium Polymer Battery, is the standard energy source for modern drones due to its high energy density and discharge rate capability. It consists of lithium-ion polymer cells enclosed in a flexible pouch that can be formed into various shapes to fit drone frames. LiPo batteries are characterized by their cell count, capacity in milliampere-hours, and C-rating for maximum discharge current. Proper charging, storage, and voltage management are critical because LiPo cells can be damaged by over-discharge or overcharging.}{catD}
\glsadd{Battery-Capacity}\dictentry{Battery Capacity}{Battery Capacity measures the total energy stored in a drone battery, expressed in milliampere-hours. A higher capacity battery stores more energy and extends flight time, but also adds weight that can reduce maneuverability and performance. The relationship between capacity, weight, and flight duration involves trade-offs that pilots must consider for each mission. Typical consumer drone batteries range from 1,000 to 6,000 mAh, while professional platforms may carry much larger packs.}{catD}
\glsadd{S-Rating}\dictentry{S Rating}{The S Rating of a LiPo battery indicates the number of cells connected in series, which determines the battery's nominal voltage. Each lithium polymer cell provides approximately 3.7 volts, so a 3S battery delivers about 11.1 volts and a 6S battery delivers about 22.2 volts. Higher voltage systems can drive larger motors more efficiently and generate more power for a given current draw. The S rating must be compatible with the motors, ESCs, and flight controller to ensure safe and optimal operation.}{catD}
\glsadd{C-Rating}\dictentry{C Rating}{The C Rating of a LiPo battery specifies its maximum safe continuous discharge rate relative to its capacity. A 1,500 mAh battery with a 25C rating can theoretically deliver 37.5 amperes continuously without damage. The C rating determines whether a battery can supply the peak current demanded by the motors during aggressive flight or heavy payload operations. Exceeding the C rating can cause voltage sag, overheating, reduced battery life, and in extreme cases, fire or cell failure.}{catD}
\glsadd{Power-Distribution-Board}\dictentry{Power Distribution Board}{A Power Distribution Board distributes the main battery voltage to each ESC on a multirotor drone. It provides a central connection point that splits the high-current battery bus to individual motor controllers while including filtering capacitors and sometimes voltage regulators. Some PDBs also integrate current and voltage sensors that feed battery monitoring data to the flight controller. The board must handle the total current draw of all motors simultaneously without excessive resistance or heat generation.}{catD}
\glsadd{PDB}\dictentry{PDB}{PDB is the abbreviation for Power Distribution Board, a central hub that routes battery power to the drone's electronic speed controllers. It serves as the electrical backbone of a multirotor frame, providing clean and balanced power distribution to all motors. Modern PDBs often incorporate voltage regulation, battery monitoring, and LED status indicators on a single compact board. Integrated PDB and ESC stack designs have become common in racing and freestyle drone builds.}{catD}
\glsadd{BEC}\dictentry{BEC}{A BEC, or Battery Eliminator Circuit, is a voltage regulator that steps down the main battery voltage to a stable lower voltage for powering the flight controller, receiver, and other avionics. It eliminates the need for a separate receiver battery by deriving regulated power from the main flight pack. BECs are either integrated into ESCs or provided as standalone modules, typically outputting 5 volts or adjustable voltages. A reliable BEC is essential because voltage fluctuations can cause flight controller resets or loss of receiver control.}{catD}
\glsadd{Telemetry}\dictentry{Telemetry}{Telemetry is the automated collection and transmission of measurement data from a spacecraft to ground stations for monitoring. Telemetry includes health parameters like temperatures and voltages plus scientific instrument readings. Data is encoded, modulated onto radio waves, and transmitted in real time or stored for later playback. Telemetry systems are the primary means by which controllers monitor spacecraft status and detect anomalies.}{catD}
\glsadd{Receiver}\dictentry{Receiver}{A Receiver is the radio module on the drone that picks up control commands transmitted by the pilot's radio transmitter. It decodes the incoming signal and converts it into PWM, PPM, or serial protocol outputs that the flight controller interprets as stick inputs. Receivers operate on specific frequency bands and use protocols such as SBUS, CRSF, or IBUS for communication with the flight controller. The receiver's range, latency, and signal reliability are critical factors in maintaining safe pilot control of the aircraft.}{catD}
\glsadd{Transmitter}\dictentry{Transmitter}{A Transmitter is the handheld radio control unit used by a pilot to send commands to the drone. It converts the pilot's stick movements and switch inputs into radio signals transmitted on a specific frequency. Modern transmitters support multiple protocols, model memory, telemetry display, and customizable switch assignments for different aircraft. The transmitter is the primary interface between the pilot and the drone, and its ergonomic design, signal range, and reliability directly affect the flying experience.}{catD}
\glsadd{Radio-Link}\dictentry{Radio Link}{A Radio Link is the wireless communication channel between the ground transmitter and the drone's onboard receiver. It carries control commands from the pilot to the aircraft and often telemetry data back to the ground. Radio links operate on various frequency bands including 2.4 GHz for control and 5.8 GHz for video, each offering different range, penetration, and bandwidth characteristics. Link quality is influenced by antenna choice, transmitter power, environmental interference, and line-of-sight conditions.}{catD}
\glsadd{Control-Link}\dictentry{Control Link}{The Control Link is the radio uplink that transmits pilot commands from the ground transmitter to the drone's receiver and flight controller. It must maintain low latency and high reliability because any interruption can result in loss of aircraft control. Control links use frequency hopping and error correction to maintain connection quality in the presence of interference. The maximum range and robustness of the control link are fundamental to the safe operational envelope of any unmanned aircraft.}{catD}
\glsadd{Video-Link}\dictentry{Video Link}{The Video Link is the downlink that transmits live camera footage from the drone to the pilot's goggles or ground monitor. It enables first-person-view piloting by providing a real-time visual perspective from the aircraft's onboard camera. Video links typically operate on 5.8 GHz using either low-latency analog or higher-quality digital transmission systems. The video link must deliver smooth, low-latency imagery because delays above a few milliseconds can make FPV piloting difficult or unsafe.}{catD}
\glsadd{Data-Link}\dictentry{Data Link}{A Data Link is a bi-directional communication channel that carries both telemetry from the drone to ground and commands from ground to drone. It supports higher-bandwidth data exchange than a simple control link, enabling mission uploads, parameter changes, and detailed system status reporting. Data links may use dedicated radio modems, cellular networks, or satellite connections depending on range and infrastructure requirements. Reliable data links are essential for beyond-visual-line-of-sight operations and autonomous missions.}{catD}
\glsadd{Frequency-Band}\dictentry{Frequency Band}{A Frequency Band is a designated range of the radio spectrum allocated for drone communication between the aircraft and ground control. Different bands offer varying trade-offs between range, data throughput, obstacle penetration, and regulatory restrictions. The most common drone bands include 2.4 GHz for control, 5.8 GHz for video, and lower frequencies like 900 MHz for long-range telemetry. Selecting the appropriate frequency band depends on the mission requirements for range, bandwidth, and interference resilience.}{catD}
\glsadd{24-GHz}\dictentry{2.4 GHz}{The 2.4 GHz band is the most widely used frequency range for drone control links worldwide. It offers a good balance between range, data rate, and antenna size, making it suitable for both short and medium-range operations. The band is shared with Wi-Fi, Bluetooth, and other consumer devices, so drone protocols use frequency hopping to avoid interference. Most modern radio control transmitters and receivers operate in this band using protocols like ExpressLRS, Crossfire, or FrSky.}{catD}
\glsadd{58-GHz}\dictentry{5.8 GHz}{The 5.8 GHz band is commonly used for FPV video transmission on drones due to its higher bandwidth and reduced congestion compared to 2.4 GHz. It supports both analog and digital video transmission with resolutions suitable for real-time piloting and recording. The higher frequency provides shorter wavelength antennas that are compact enough for small drones, but it has reduced range and penetration through obstacles. Many FPV systems use 5.8 GHz for video while simultaneously using 2.4 GHz for control on separate frequencies.}{catD}
\glsadd{900-MHz}\dictentry{900 MHz}{The 900 MHz band is used for long-range telemetry and control links on drones where extended distance is required. Its longer wavelength provides superior obstacle penetration and range compared to higher-frequency bands. The lower data rate limits it primarily to control signals and telemetry rather than video, but for long-range missions this trade-off is favorable. Systems like Crossfire and ExpressLRS operate in this band to provide reliable multi-kilometer control links for autonomous and long-range drone operations.}{catD}
\glsadd{433-MHz}\dictentry{433 MHz}{The 433 MHz band offers ultra-long-range capability for drone telemetry and control due to its excellent propagation characteristics. Its long wavelength diffracts around obstacles and travels great distances with relatively low transmitter power. The band has limited data bandwidth, making it suitable for low-rate telemetry and control signals rather than video. Some long-range FPV and autonomous drone systems use 433 MHz as a backup or primary control link for missions extending tens of kilometers from the operator.}{catD}
\glsadd{FHSS}\dictentry{FHSS}{FHSS, or Frequency Hopping Spread Spectrum, is a modulation technique where the radio signal rapidly switches between many frequency channels within a band. This hopping pattern is synchronized between transmitter and receiver, making the link resistant to interference and difficult to jam or intercept. FHSS provides robust communication in crowded spectrum environments by spreading the signal energy across a wide bandwidth. It is the standard modulation method for modern drone control links including ExpressLRS and FrSky protocols.}{catD}
\glsadd{DSSS}\dictentry{DSSS}{DSSS, or Direct Sequence Spread Spectrum, is a modulation technique that spreads a signal across a wide bandwidth by multiplying the data with a pseudo-random noise code. The receiver uses the same code to despread and recover the original signal, rejecting interference at other frequencies. DSSS provides strong resistance to narrowband interference and multipath fading, making it reliable in complex RF environments. It is used in some drone communication systems and shares fundamental principles with GPS signal reception.}{catD}
\glsadd{LBT}\dictentry{LBT}{LBT, or Listen Before Talk, is a spectrum access protocol where a radio transmitter listens to the channel before transmitting to detect whether it is already in use. If the channel is occupied, the device waits or selects a different frequency, reducing collisions and interference with other users. LBT is required by regulations in certain regions, particularly in Europe, for unlicensed radio operation. It promotes fair and efficient spectrum sharing among multiple drone systems and other wireless devices operating in the same band.}{catD}
\glsadd{RSSI}\dictentry{RSSI}{RSSI, or Received Signal Strength Indicator, is a measurement of the power level being received by the drone's radio antenna from the ground transmitter. It is typically displayed as a percentage or decibel value on the pilot's transmitter or telemetry display. RSSI provides a relative indication of link quality, with lower values indicating a weakening connection that may lead to signal loss. Pilots use RSSI to gauge their operational range and decide when to return before the control link becomes unreliable.}{catD}
\glsadd{Antenna}\dictentry{Antenna}{A spacecraft antenna is a transducer converting electrical signals into radio waves for transmission and vice versa for reception. Antennas include dipoles, horns, arrays, and parabolic reflectors, each optimized for specific frequencies and coverage patterns. The choice depends on data rates, pointing accuracy, beamwidth, and available spacecraft space and power.}{catD}
\glsadd{Omnidirectional-Antenna}\dictentry{Omnidirectional Antenna}{An Omnidirectional Antenna radiates and receives radio signals equally in all directions within a horizontal plane. It provides consistent coverage regardless of the drone's bearing relative to the ground station, making it ideal for general flying where the aircraft changes direction frequently. Omnidirectional antennas typically offer lower gain than directional types, limiting their maximum range. They are the standard antenna choice for drone receivers and transmitters in most consumer and commercial applications.}{catD}
\glsadd{Directional-Antenna}\dictentry{Directional Antenna}{A Directional Antenna focuses its radiation pattern in a specific direction, concentrating radio energy to achieve greater range and signal strength. It provides higher gain than omnidirectional antennas by narrowing the beam, but requires the operator to point it toward the drone during flight. Directional antennas are essential for long-range operations where maximizing link budget is critical. Common types include patch antennas, helical antennas, and Yagi antennas, each offering different beam widths and gain characteristics.}{catD}
\glsadd{Patch-Antenna}\dictentry{Patch Antenna}{A Patch Antenna is a flat, low-profile directional antenna consisting of a metallic patch element on a grounded substrate. It radiates in a broad beam perpendicular to its surface, providing moderate gain in a compact form factor. Patch antennas are commonly used for FPV video reception and GPS signal reception on drones due to their small size and lightweight construction. They are easy to mount on gimbals or tracking platforms that keep the antenna pointed at the aircraft during flight.}{catD}
\glsadd{Yagi-Antenna}\dictentry{Yagi Antenna}{A Yagi Antenna is a high-gain directional antenna consisting of a driven element, a reflector, and one or more directors arranged along a boom. It produces a narrow radiation beam with significantly higher gain than omnidirectional or patch antennas. Yagi antennas are used for long-range drone operations where maximum signal range and penetration are needed. The narrow beam requires aiming toward the drone, which can be done manually or with antenna tracking systems that follow the aircraft's GPS position.}{catD}
\glsadd{3-Axis-Gimbal}\dictentry{3-Axis Gimbal}{A 3-Axis Gimbal stabilizes a camera by actively compensating for movement in pitch, roll, and yaw. Each axis has its own brushless motor and controller that responds to gyroscopic sensor input to counteract the drone's attitude changes. This full three-axis stabilization produces perfectly level and smooth footage even during aggressive flight maneuvers. Three-axis gimbals are the standard for professional aerial cinematography and precision sensor pointing applications.}{catD}
\glsadd{2-Axis-Gimbal}\dictentry{2-Axis Gimbal}{A 2-Axis Gimbal stabilizes a camera by compensating for pitch and roll movements while allowing the yaw axis to rotate freely with the drone. It provides adequate stabilization for many aerial photography applications at a lower weight and cost than three-axis systems. The lack of yaw stabilization means that panning movements must come from the drone's own rotation, which can introduce some horizontal jitter during fast turns. Two-axis gimbals are commonly found on consumer drones and FPV platforms.}{catD}
\glsadd{PID-Tuning}\dictentry{PID Tuning}{PID Tuning adjusts proportional, integral, and derivative gains of the flight controller's control loops to optimize flight performance. The proportional term corrects errors proportionally to magnitude, the integral term eliminates steady-state errors, and the derivative term dampens oscillations. Properly tuned values produce a drone that is responsive yet stable, without excessive overshoot or oscillation. Tuning involves incremental adjustments while observing the drone's response to stick inputs and disturbances.}{catD}
\glsadd{Autotune}\dictentry{Autotune}{Autotune is an automated process where the flight controller systematically tests and adjusts its own PID parameters to achieve optimal flight characteristics. The drone performs small, controlled maneuvers while measuring its response, then calculates and applies tuned gains without pilot intervention. Autotune simplifies setup for new builds and ensures baseline performance even without tuning expertise. Results vary by airframe and conditions, and manual fine-tuning may still be needed for demanding applications.}{catD}
\glsadd{Rate-Mode}\dictentry{Rate Mode}{Rate Mode is a flight mode where the pilot's stick inputs command angular velocity rather than a target angle. The drone rotates at a rate proportional to stick deflection and holds its current angle when the sticks are centered. This mode provides maximum pilot control and is preferred for aerobatic flying, FPV racing, and freestyle maneuvers. Rate mode requires continuous pilot input to maintain orientation, as the drone will hold whatever attitude it reaches rather than leveling itself.}{catD}
\glsadd{Angle-Mode}\dictentry{Angle Mode}{Angle Mode, also called Self-Leveling or Stabilize Mode, automatically returns the drone to a level attitude when the pilot centers the sticks. The pilot's stick inputs command a target tilt angle, with full stick deflection setting a maximum angle limit. This mode is inherently stable and easier to fly, making it ideal for beginners and photography missions that need smooth, predictable behavior. The angle limiting prevents the drone from reaching extreme attitudes that could lead to loss of control.}{catD}
\glsadd{Horizon-Mode}\dictentry{Horizon Mode}{Horizon Mode combines characteristics of both rate mode and angle mode in a single flight mode. At small stick deflections it behaves like rate mode, allowing smooth and precise control for gentle maneuvering. When the sticks are pushed beyond a defined threshold, the drone transitions to angle mode and self-levels upon stick release. This hybrid behavior gives experienced pilots the freedom to perform flips and rolls while retaining an automatic recovery when sticks are centered.}{catD}
\glsadd{Altitude-Hold}\dictentry{Altitude Hold}{Altitude Hold is a flight mode where the drone maintains a constant barometric altitude without pilot throttle input. The flight controller uses data from the barometer and sometimes a downward-facing ultrasonic or LiDAR sensor to regulate motor output and hold vertical position. The pilot can adjust the target altitude with the throttle stick, and the drone will smoothly ascend or descend to the new height. This mode is essential for photography, inspection, and any mission requiring stable altitude without constant pilot attention.}{catD}
\glsadd{Position-Hold}\dictentry{Position Hold}{Position Hold is a flight mode where the drone maintains both its GPS position and altitude, hovering in place without pilot input. The flight controller continuously compares the current GPS coordinates against the locked position and applies corrections to counteract wind and drift. GPS accuracy limitations mean the drone may drift slowly within a small radius, typically one to three meters. This mode is useful for static aerial observation, communication relay, and any scenario requiring the drone to stay over a fixed point.}{catD}
\glsadd{Return-to-Home}\dictentry{Return to Home}{Return to Home is an automated flight mode that navigates the drone back to its takeoff location when activated or triggered by a failsafe condition. The drone typically ascends to a preset safe altitude, then flies a direct or waypoint path back to the home point before descending and landing. RTH can be triggered manually by the pilot, or automatically due to low battery, lost link, or geofence breach. This safety feature provides a critical backup that can recover the aircraft even when the pilot loses situational awareness.}{catD}
\glsadd{RTH}\dictentry{RTH}{RTH is the abbreviation for Return to Home, an automated drone function that returns the aircraft to its takeoff point. It can be activated manually by the pilot or automatically by the flight controller in response to critical conditions like low battery or signal loss. The drone climbs to a safe altitude, navigates to the recorded home coordinates, and descends for landing. RTH is a fundamental safety feature that helps prevent loss of aircraft when normal control becomes difficult or impossible.}{catD}
\glsadd{Loiter}\dictentry{Loiter}{Loiter is a flight mode where the drone holds its current GPS position and altitude, maintaining a steady hover over a fixed point. Unlike Position Hold, which locks to the initial takeoff location, Loiter allows the pilot to move the drone to any point and then command it to stay there. The flight controller uses GPS and barometric data to resist wind drift and maintain the set position. Loiter mode is valuable for surveillance, photography, and observation missions that require the drone to remain stationary over a target area.}{catD}
\glsadd{Waypoint-Navigation}\dictentry{Waypoint Navigation}{Waypoint Navigation is an autonomous flight mode where the drone follows a pre-programmed sequence of GPS coordinates at specified altitudes and speeds. The mission planner defines each waypoint with latitude, longitude, altitude, and optional actions like taking a photo or triggering a sensor. Once the mission is uploaded and launched, the drone flies the entire route without pilot input. This capability enables repeatable survey flights, corridor mapping, and systematic area coverage with consistent precision.}{catD}
\glsadd{Follow-Me}\dictentry{Follow Me}{Follow Me is a flight mode where the drone autonomously tracks and follows a moving target, typically identified by a GPS device carried by the operator or subject. The flight controller continuously adjusts the drone's position to maintain a preset distance, altitude, and angle relative to the tracked person. Follow Me mode uses real-time GPS updates from the tracked device to match speed and direction changes. It is popular for action sports filming, outdoor recreation, and dynamic surveillance applications.}{catD}
\glsadd{Orbit-Mode}\dictentry{Orbit Mode}{Orbit Mode commands the drone to fly in a continuous circle around a specified GPS coordinate at a set radius, altitude, and speed. The drone adjusts its heading to keep the camera or sensor pointed at the center point throughout the orbit. This mode produces smooth, cinematic circling shots that would be extremely difficult for a manual pilot to execute consistently. Orbit Mode is widely used in filmmaking, real estate marketing, and inspection of structures like towers or buildings.}{catD}
\glsadd{Mission-Planning}\dictentry{Mission Planning}{Mission Planning is the process of designing an autonomous flight plan by defining waypoints, altitudes, speeds, and actions that the drone will execute without pilot input. The plan is created using ground station software that provides a map interface for placing waypoints and setting parameters. Missions can include takeoff, multiple waypoints with specific actions, and an automated landing at the end. Effective mission planning ensures complete coverage of survey areas, efficient flight paths, and consistent repeatable operations.}{catD}
\glsadd{Geofence}\dictentry{Geofence}{A Geofence is a virtual three-dimensional boundary programmed into the flight controller that restricts the drone's flight area. If the drone approaches the geofence boundary, the system triggers a warning or automatically executes a protective action such as hovering, returning home, or landing. Geofences can define maximum altitude limits, horizontal distance boundaries, and restricted zones around airports or sensitive areas. They are a critical safety and regulatory compliance tool that prevents unauthorized flight into prohibited airspace.}{catD}
\glsadd{Low-Battery-Warning}\dictentry{Low Battery Warning}{A Low Battery Warning is an alert triggered when the drone's battery voltage drops below a pre-configured threshold during flight. The warning appears on the pilot's transmitter screen, goggles, or mobile application, and may include audible tones or LED indicators on the aircraft. It gives the pilot time to safely return the drone and land before the battery reaches a critically low level that could trigger forced landing. Timely response is essential for preventing crashes and battery damage.}{catD}
\glsadd{Collision-Avoidance}\dictentry{Collision Avoidance}{Collision Avoidance is an onboard system that detects nearby obstacles and automatically adjusts the flight path to prevent impact. It typically uses stereo cameras, ultrasonic sensors, or LiDAR to measure distances to objects in the flight path. When an obstacle is detected within a warning zone, the system can slow the drone, stop it, or reroute around the obstacle depending on flight mode and settings. Collision avoidance is critical for safe operation in complex environments and autonomous flights.}{catD}
\glsadd{Obstacle-Detection}\dictentry{Obstacle Detection}{Obstacle Detection is the sensing capability that identifies physical objects in the drone's flight path using cameras, LiDAR, infrared sensors, or ultrasonic transducers. These sensors create real-time awareness of the surrounding environment that the flight controller can use for collision avoidance and path planning. Detection range, field of view, and processing speed determine effectiveness at various flight speeds. Obstacle detection is the foundation for autonomous flight safety and navigating complex environments.}{catD}
\glsadd{Emergency-Stop}\dictentry{Emergency Stop}{An Emergency Stop is an immediate command that kills all motor rotation, causing the drone to fall from its current position. It is a last resort when the drone poses an immediate danger and continued flight is more hazardous than uncontrolled descent. Emergency stop is typically assigned to a dedicated transmitter switch and must never be activated at high altitude or over people. This function should only be used when damage from a controlled crash is less than the risk of continued flight.}{catD}
\glsadd{Lost-Link}\dictentry{Lost Link}{Lost Link is the condition where the communication between the ground transmitter and the drone's receiver is interrupted or lost entirely. It can result from excessive distance, antenna misalignment, interference, or physical obstructions blocking the radio signal. The flight controller detects the loss of control input and triggers a pre-configured failsafe action, typically Return to Home. Pilots must configure appropriate lost link procedures before every flight to ensure the drone behaves safely when the control connection is broken.}{catD}
\glsadd{Failsafe}\dictentry{Failsafe}{A Failsafe is an automatic protective action executed by the flight controller when a critical condition is detected during flight. Common triggers include low battery voltage, loss of radio link, GPS failure, or geofence breach. The failsafe response can be configured to hover in place, land immediately, or return to home depending on the trigger and pilot preference. Failsafe mechanisms are a fundamental safety layer that protects both the aircraft and people on the ground when normal operation is compromised.}{catD}
\glsadd{Parachute-System}\dictentry{Parachute System}{A Parachute System is a recovery device that deploys a canopy to slow a drone's descent during an emergency or controlled landing. It provides a safe alternative to an uncontrolled freefall when the aircraft experiences motor failure, structural damage, or a commanded flight termination. The parachute deploys using a spring mechanism, compressed gas, or pyrotechnic charge and can reduce descent rates to survivable levels. Parachute systems are increasingly mandated for operations over people and in urban environments.}{catD}
\glsadd{VLOS}\dictentry{VLOS}{Visual Line of Sight (VLOS) requires the remote pilot to maintain continuous, unaided visual contact with the drone throughout flight. The pilot must discern the drone's orientation and altitude without optical aids. This is the baseline rule under FAA Part 107 and EASA regulations, limiting effective range to roughly 500 meters.}{catD}
\glsadd{BVLOS}\dictentry{BVLOS}{Beyond Visual Line of Sight (BVLOS) allows drone operations at distances where the pilot cannot see the aircraft unaided. BVLOS requires detect-and-avoid systems, reliable command links, and regulatory authorization. It is essential for long-range pipeline inspection, surveying, and cargo delivery missions.}{catD}
\glsadd{EVLOS}\dictentry{EVLOS}{Extended Visual Line of Sight (EVLOS) operates a drone beyond the pilot's range but within sight of ground observers along the route. Observers relay the drone's position to the pilot via radio. EVLOS extends range while keeping a human safety net, used in corridor surveying and infrastructure inspection.}{catD}
\glsadd{Remote-ID}\dictentry{Remote ID}{Remote ID requires drones to broadcast identity, location, and operator information during flight via Bluetooth or Wi-Fi signals. It functions as a digital license plate for accountability. The FAA and EASA mandate Remote ID as a prerequisite for routine operations in shared airspace with manned aviation.}{catD}
\glsadd{No-Fly-Zone}\dictentry{No-Fly Zone}{A No-Fly Zone (NFZ) is restricted airspace where drone flight is prohibited without special authorization. NFZs typically surround airports, military bases, and government buildings. Consumer drones include geofencing software that prevents flight in these areas. Violations can result in fines, criminal charges, and equipment confiscation.}{catD}
\glsadd{Airspace}\dictentry{Airspace}{Airspace is the atmospheric volume where drones operate, classified into controlled and uncontrolled categories by altitude and airport proximity. Controlled airspace requires authorization before entry, while uncontrolled airspace permits standard flight. ICAO divides airspace into classes A through G with distinct access requirements. Understanding classification is essential for legal operations.}{catD}
\glsadd{UTM}\dictentry{UTM}{Unmanned Aircraft System Traffic Management (UTM) is a digital framework managing drone air traffic in low-altitude airspace. It enables automated authorizations, flight planning, and real-time conflict detection without human controllers. NASA and the FAA have developed UTM through field trials, critical for scaling urban delivery and large-area surveying operations.}{catD}
\glsadd{NOTAM}\dictentry{NOTAM}{A Notice to Air Missions (NOTAM) is an official advisory alerting pilots to hazards, restrictions, or airspace changes. NOTAMs cover temporary flight restrictions, runway closures, and military areas. Drone operators must check current NOTAMs before every flight. They are distributed through official systems and updated in real time.}{catD}
\glsadd{Flight-Authorization}\dictentry{Flight Authorization}{Flight Authorization is formal permission required before operating a drone in controlled airspace or under non-standard conditions. It may be granted through automated systems like LAANC or via direct application to the aviation authority. Authorization ensures operations do not conflict with manned aircraft.}{catD}
\glsadd{LAANC}\dictentry{LAANC}{Low Altitude Authorization and Notification Capability (LAANC) provides near-instant drone flight authorizations in controlled airspace near airports. Developed by the FAA with private providers, it replaces weeks of manual processing with real-time digital approval. Requests are submitted through approved apps and authorized within seconds.}{catD}
\glsadd{Part-107}\dictentry{Part 107}{Part 107 is the FAA regulation governing small unmanned aircraft for commercial use in the United States. It covers pilot certification, weight limits under 55 pounds, a 400-foot altitude cap, and visual line of sight requirements. The Remote Pilot Certificate is obtained through an FAA knowledge exam. Amendments now allow night flights and operations over people.}{catD}
\glsadd{Drone-Registration}\dictentry{Drone Registration}{Drone Registration records a drone with the national aviation authority, which issues a unique identification number. The FAA requires registration for drones between 0.55 and 55 pounds. The number must be visibly marked on the aircraft and renewed every three years. This system enables authorities to trace drones to their owners.}{catD}
\glsadd{Operator-License}\dictentry{Operator License}{An Operator License certifies an individual or organization to conduct commercial drone operations. Requirements include passing a knowledge exam, demonstrating piloting competency, and maintaining safety protocols. The license may impose altitude limits and geographic restrictions. It is a prerequisite for most commercial drone services worldwide.}{catD}
\glsadd{Remote-Pilot-Certificate}\dictentry{Remote Pilot Certificate}{A Remote Pilot Certificate is an FAA credential authorizing commercial small unmanned aircraft operation under Part 107. Candidates pass a 60-question aeronautical knowledge test covering airspace, weather, and operational safety. Renewal requires a recurrent exam every 24 months. It is the legal equivalent of a pilot's license for U.S. drone operators.}{catD}
\glsadd{Aerial-Photography}\dictentry{Aerial Photography}{Aerial Photography captures images or video from a camera mounted on a drone during flight. Stabilized gimbals and high-resolution sensors produce professional-quality imagery from low altitudes. Applications span real estate, filmmaking, journalism, and construction documentation. Drone photography has revolutionized visual content creation across many industries.}{catD}
\glsadd{Precision-Agriculture}\dictentry{Precision Agriculture}{Precision Agriculture uses drones with multispectral and thermal sensors to monitor crop health at a granular level. Imagery reveals variations in plant vigor and nutrient deficiency invisible to the naked eye. This enables targeted application of water, fertilizer, and pesticides, reducing waste. Workflows include automated flight planning and vegetation index computation.}{catD}
\glsadd{Crop-Monitoring}\dictentry{Crop Monitoring}{Crop Monitoring is systematic aerial surveillance of fields using drones to track plant health and growth over time. Time-series data reveals trends in development. Vegetation indices like NDVI from multispectral sensors quantify plant vigor. Farmers use this data for targeted irrigation and pest treatment, improving efficiency and reducing losses.}{catD}
\glsadd{Mapping}\dictentry{Mapping}{Mapping creates orthophoto maps and georeferenced mosaics from overlapping drone images captured in systematic flight patterns. Software stitches hundreds of photographs into seamless, distortion-free images with accurate ground coordinates. Outputs achieve centimeter-level resolution for surveying, urban planning, and construction tracking.}{catD}
\glsadd{Surveying}\dictentry{Surveying}{Surveying measures and maps land features using drone photogrammetry or LiDAR to produce precise topographic data. Drones collect dense point clouds and orthophotos rapidly, reducing field time compared to ground crews. Digital elevation models and volume calculations meet survey-grade accuracy with ground control points. It is standard in civil engineering and mining.}{catD}
\glsadd{Photogrammetry}\dictentry{Photogrammetry}{Photogrammetry extracts three-dimensional measurements from overlapping two-dimensional photographs. Drone photogrammetry computes 3D point clouds, mesh models, and textured surfaces from images captured at multiple angles. Accuracy reaches centimeters with proper ground control points. It is widely used in construction, archaeology, and forensic documentation.}{catD}
\glsadd{Infrastructure-Inspection}\dictentry{Infrastructure Inspection}{Infrastructure Inspection uses drones to examine bridges, power lines, pipelines, and buildings without human access to dangerous locations. High-resolution and thermal sensors capture detailed structural imagery. Automated flight paths ensure consistent, repeatable data collection. Analysis detects corrosion, cracks, and defects early, reducing costs 40-60\% versus traditional methods.}{catD}
\glsadd{Search-and-Rescue}\dictentry{Search and Rescue}{Search and Rescue (SAR) uses drones to locate missing persons in wilderness, urban, or disaster environments. Thermal cameras detect heat signatures through foliage, at night, or in low visibility. Drones rapidly cover large areas and relay real-time video to ground teams. Autonomous grid-search missions reduce pilot fatigue. Drone SAR has saved lives by locating victims hours faster than ground teams.}{catD}
\glsadd{Disaster-Response}\dictentry{Disaster Response}{Disaster Response deploys drones after catastrophes to assess damage, locate survivors, and coordinate relief. Drones generate orthomosaics and 3D models providing situational awareness to emergency managers. Thermal sensors identify trapped individuals and assess structural integrity. Relay drones restore communication where infrastructure is destroyed.}{catD}
\glsadd{Wildlife-Monitoring}\dictentry{Wildlife Monitoring}{Wildlife Monitoring uses drones to track and count animal populations with minimal disturbance. Thermal and high-resolution cameras survey large, inaccessible areas. Machine learning algorithms enable species identification from aerial imagery. Non-invasive surveys reduce animal stress compared to ground counting. This technology supports conservation and anti-poaching efforts worldwide.}{catD}
\glsadd{Parcel-Delivery}\dictentry{Parcel Delivery}{Parcel Delivery by drone autonomously transports small packages from distribution centers to customers. Delivery drones follow GPS routes with detect-and-avoid systems. Winches or cargo pods enable contactless drop-off at landing zones. Companies like Wing and Zipline have pioneered commercial delivery. The technology promises reduced last-mile costs and lower emissions.}{catD}
\glsadd{Environmental-Monitoring}\dictentry{Environmental Monitoring}{Environmental Monitoring uses drone-mounted sensors to track air quality, water conditions, and vegetation health over time. Drones collect data at resolutions impossible with satellites. Multispectral, thermal, and gas sensors measure pollution, soil moisture, and deforestation. Automated flights ensure consistent longitudinal datasets for trend analysis and regulatory compliance.}{catD}
\glsadd{Mission-Planner}\dictentry{Mission Planner}{Mission Planner is ground control station software designing, simulating, and executing autonomous drone flights. Users define waypoints, altitudes, and survey patterns through a graphical interface. The planner computes optimal paths, estimates battery needs, and uploads missions to the flight controller. It is essential for systematic mapping, surveying, and agricultural operations.}{catD}
\glsadd{QGroundControl}\dictentry{QGroundControl}{QGroundControl is a free, open-source ground control station for flight planning, vehicle configuration, and real-time telemetry. It supports PX4 and ArduPilot on Windows, macOS, Linux, Android, and iOS. The software provides intuitive interfaces for mission setup and sensor calibration. It is the reference ground station for the Dronecode ecosystem.}{catD}
\glsadd{PX4}\dictentry{PX4}{PX4 is a free, open-source flight control software providing a complete autopilot with attitude estimation, position control, sensor fusion, and mission execution. It runs on various hardware using a modular uORB messaging architecture. PX4 integrates with MAVLink and QGroundControl, maintained by the Dronecode Foundation and a global contributor community.}{catD}
\glsadd{ArduPilot}\dictentry{ArduPilot}{ArduPilot is a free, open-source autopilot supporting multicopters, fixed-wing aircraft, helicopters, rovers, and submarines. It provides waypoint navigation, terrain following, and geofencing capabilities. ArduPilot runs on various hardware with extensive parameter customization. It has one of the largest open-source robotics communities worldwide.}{catD}
\glsadd{MAVLink}\dictentry{MAVLink}{MAVLink (Micro Air Vehicle Link) is a lightweight, open-source protocol for drone-to-ground station communication. It uses compact binary messages for telemetry, commands, and parameter updates. MAVLink supports serial, UDP, and TCP transports. It is the standard messaging format for PX4, ArduPilot, and most commercial drone systems.}{catD}
\glsadd{MAVSDK}\dictentry{MAVSDK}{MAVSDK is a C++ and Python library providing a clean API for drone communication via MAVLink. It abstracts protocol complexity into functions for takeoff, landing, navigation, and camera control. MAVSDK supports UDP, TCP, and serial backends. It is designed for building drone applications, maintained by the Dronecode Foundation.}{catD}
\glsadd{Gazebo}\dictentry{Gazebo}{Gazebo is an open-source 3D robotics simulator providing realistic physics, sensor simulation, and environmental rendering for drone development. Developers test flight controllers and autonomous algorithms in virtual environments before hardware deployment. It supports multi-vehicle simulation and integrates with PX4 and ROS for hardware-in-the-loop testing.}{catD}
\glsadd{AirSim}\dictentry{AirSim}{AirSim is a Microsoft-developed open-source simulator for autonomous vehicles including drones. Built on Unreal Engine, it provides photo-realistic environments and accurate aerodynamic models. AirSim simulates cameras, LiDAR, and IMU sensors, integrating with PX4 and ArduPilot. It is widely used in AI research for training perception and navigation algorithms.}{catD}
\glsadd{SITL}\dictentry{SITL}{Software-in-the-Loop (SITL) runs actual flight controller software on a general-purpose computer without physical hardware. The autopilot executes natively while simulated sensors and motor outputs are computed by a companion simulator. SITL enables rapid iteration without real flight risk. It is the primary testing method for PX4 and ArduPilot development.}{catD}
\glsadd{HITL}\dictentry{HITL}{Hardware-in-the-Loop (HITL) connects actual flight controller hardware to a simulator providing synthetic sensor data. The autopilot runs on its real processor while physics execute on a connected computer. HITL validates hardware-specific factors like processing latency and sensor interfaces. It bridges software simulation and real-world flight testing.}{catD}
\glsadd{ROS}\dictentry{ROS}{Robot Operating System (ROS) is a middleware framework for inter-process communication, sensor integration, and algorithm deployment in drone development. ROS 2 adds real-time capabilities and production reliability. In drones, it handles perception pipelines, path planning, and multi-vehicle coordination. ROS is the de facto standard for robotics research and commercial drone development.}{catD}
\glsadd{DroneKit}\dictentry{DroneKit}{DroneKit is a Python library for drone communication and control through a high-level scriptable API. It provides functions for arming, takeoff, waypoint navigation, and telemetry access. DroneKit abstracts MAVLink complexity into simple Python commands. While superseded by MAVSDK for new projects, it remains used in education and prototyping.}{catD}
\glsadd{Camera-Payload}\dictentry{Camera Payload}{A Camera Payload is a stabilized camera system for high-resolution photography or video from a drone. Modern cameras offer 12-megapixel to 8K resolution with optical zoom and RAW capture. Motorized gimbals provide multi-axis vibration isolation. Payload selection balances resolution, weight, and sensor characteristics for photography, mapping, and inspection applications.}{catD}
\glsadd{LiDAR-Payload}\dictentry{LiDAR Payload}{A LiDAR Payload generates dense 3D point clouds by measuring distances through laser time-of-flight reflection. It pulses thousands of beams per second to capture terrain, vegetation, and structures at centimeter accuracy. Drone LiDAR is essential for surveying vegetated areas where photogrammetry cannot reach the ground. It is used in forestry, flood mapping, and infrastructure corridor surveys.}{catD}
\glsadd{Sprayer}\dictentry{Sprayer}{A Sprayer is an agricultural payload dispensing pesticides, herbicides, or fertilizers in controlled quantities. It consists of a tank, pump, nozzles, and flow control electronics regulated by flight speed and altitude. GPS-guided auto-spraying ensures even coverage without overlaps. Drone sprayers reduce chemical usage and keep operators away from hazardous substances.}{catD}
\glsadd{Cargo-Pod}\dictentry{Cargo Pod}{A Cargo Pod is an enclosed container on a delivery drone for transporting parcels or medical supplies. Designs range from open frames to climate-controlled enclosures. Pods minimize aerodynamic drag while protecting contents during flight and release. Quick-release mechanisms allow rapid loading. Pod design directly affects range, speed, and payload capacity.}{catD}
\glsadd{Loudspeaker}\dictentry{Loudspeaker}{A Loudspeaker Payload is an audio broadcasting system on a drone for public address and emergency alerts. It amplifies pre-recorded or live messages from an elevated position. Applications include disaster evacuation direction and search and rescue communication. The aerial position provides superior sound propagation over rough terrain. These payloads are lightweight and battery-powered.}{catD}
\glsadd{Spotlight}\dictentry{Spotlight}{A Spotlight Payload is a high-intensity LED light illuminating areas during nighttime drone operations. Spotlights provide focused beams for night search and rescue, surveillance, and inspection. The light can be gimballed independently of the drone's heading. Battery life and heat dissipation are key design considerations. Spotlights extend operational capability into nighttime hours.}{catD}
\glsadd{Winch}\dictentry{Winch}{A Winch is a mechanical hoisting system lowering or raising payloads without the drone landing. It includes a motorized drum, cable, and release hook controlled by pilot or automation. Winches deliver cargo to unsafe locations like rooftops or dense vegetation. They support SAR extraction from hazardous terrain, eliminating landing risks.}{catD}
\glsadd{Delivery-Mechanism}\dictentry{Delivery Mechanism}{A Delivery Mechanism is any system on a drone designed to release cargo at a designated location. Options include winches, trap doors, robotic arms, and magnetic hooks. The mechanism must secure payloads during flight and release precisely without disturbing stability. Reliability is critical for autonomous logistics where human intervention at the drop point is impossible.}{catD}
\glsadd{Scientific-Instrument}\dictentry{Scientific Instrument}{A Scientific Instrument is a specialized sensor mounted on a drone for collecting research-grade data. Examples include gas analyzers, radiometers, magnetometers, and water samplers. These instruments are lightweight and calibrated for airborne accuracy. Drone-mounted instruments enable measurements in hazardous environments without risking personnel.}{catD}
\glsadd{Hyperspectral-Sensor}\dictentry{Hyperspectral Sensor}{A Hyperspectral Sensor captures data across hundreds of narrow, contiguous spectral bands from visible to shortwave infrared. Each pixel reveals chemical and physical surface properties through detailed spectral signatures. Mounted on drones, these sensors identify crop species, minerals, and pollutants at high resolution. Hyperspectral imaging exceeds the discriminatory power of standard cameras for agriculture and environmental assessment.}{catD}
\glsadd{Thermal-Camera}\dictentry{Thermal Camera}{A Thermal Camera detects infrared radiation, producing images based on temperature differences rather than visible light. Drone-mounted thermal cameras identify heat leaks, locate SAR survivors, and monitor equipment overheating. Uncooled microbolometer sensors remain lightweight for small payloads. Thermal imagery works in darkness and through smoke, making cameras indispensable for inspection and emergency response.}{catD}
\glsadd{Multispectral-Sensor}\dictentry{Multispectral Sensor}{A Multispectral Sensor captures images in several discrete wavelength bands including red, green, blue, and near-infrared. Near-infrared is especially valuable for vegetation health assessment. Drone-mounted sensors enable vegetation index computation like NDVI to quantify plant vigor. They are lighter and less expensive than hyperspectral systems, standard for precision agriculture.}{catD}
\glsadd{Flight-Time}\dictentry{Flight Time}{Flight Time is the maximum duration a drone can remain airborne on a single battery charge. It depends on battery capacity, weight, payload, and flying style. Consumer multicopters achieve 20-40 minutes while fixed-wing drones exceed several hours. Flight time determines area covered per mission. Pilots must plan with sufficient return reserve.}{catD}
\glsadd{Hover-Time}\dictentry{Hover Time}{Hover Time is the maximum duration a drone can maintain a stationary hover on a single battery charge. Hovering is the most power-intensive state for multicopters since rotors must continuously generate lift equal to the aircraft's weight. Hover time is shorter than forward flight time. It matters for inspection, surveillance, and communication relay requiring extended stationary positioning.}{catD}
\glsadd{Endurance}\dictentry{Endurance}{Endurance is the maximum total time a drone can stay airborne from takeoff to landing. It depends on energy density, aerodynamic efficiency, and system weight. Fixed-wing drones can exceed 24 hours using gasoline or solar power. Multicopter endurance typically ranges from 20 to 45 minutes. Endurance is a primary design driver for surveillance and long-range delivery.}{catD}
\glsadd{Range}\dictentry{Range}{Range is the maximum round-trip distance a drone can travel from its takeoff point and return safely. It depends on battery capacity, speed, wind, and payload. Consumer drones range 2-10 kilometers while long-range platforms exceed 100 kilometers. Regulations often impose additional restrictions based on visual line of sight requirements.}{catD}
\glsadd{Payload-Capacity}\dictentry{Payload Capacity}{Payload Capacity is the maximum weight a drone can carry beyond its own structure, battery, and standard equipment. It is determined by total propulsion thrust minus empty weight. Exceeding capacity degrades performance and compromises safety. Capacity ranges from grams on consumer drones to hundreds of kilograms on heavy-lift platforms. It is the primary constraint for sensor and cargo selection.}{catD}
\glsadd{MTOM}\dictentry{MTOM}{Maximum Takeoff Mass (MTOM) is the total mass of a drone at takeoff including airframe, battery, payload, and all systems. MTOM determines which airspace rules apply. In the EU, drones under 250 grams face minimal regulation while those above 25 kilograms need type certification. MTOM must be verified before each mission for compliance.}{catD}
\glsadd{MTOW}\dictentry{MTOW}{Maximum Takeoff Weight (MTOW) is the certified maximum weight at which a drone can lift off, measured in force units. Exceeding MTOW increases wing loading, reduces climb rate, and may destabilize the flight controller. Manufacturers specify MTOW through structural and propulsion testing. Pilots must calculate total weight before every flight to confirm compliance.}{catD}
\glsadd{Service-Ceiling}\dictentry{Service Ceiling}{Service Ceiling is the maximum altitude at which a drone can maintain a specified climb rate, typically 100 feet per minute. Beyond this, air is too thin for sufficient rotor thrust. Most consumer multicopters have a ceiling of 4,000-6,000 meters. Regulatory altitude limits usually prevent operation near this ceiling. It matters for operations in mountainous or high-elevation areas.}{catD}
\glsadd{Maximum-Speed}\dictentry{Maximum Speed}{Maximum Speed is the highest forward velocity a drone achieves in level flight under no-wind conditions. It is limited by motor power, drag, and structural integrity. Consumer drones reach 40-80 km/h while racing drones exceed 160 km/h. Maximum speed is critical for time-critical delivery and emergency response missions.}{catD}
\glsadd{Cruise-Speed}\dictentry{Cruise Speed}{Cruise Speed is the forward velocity at which a drone achieves its greatest range or endurance, representing peak efficiency. For fixed-wing drones, cruise speed falls between 60-120 km/h. Mission planning uses it to estimate flight times and battery requirements. Operating at cruise speed maximizes area covered per energy unit.}{catD}
\glsadd{Wind-Resistance}\dictentry{Wind Resistance}{Wind Resistance is the maximum sustained wind speed in which a drone operates safely and maintains stability. It depends on weight, rotor size, and flight controller capabilities. Consumer drones handle 10-15 m/s while industrial platforms tolerate 20 m/s. Exceeding limits causes instability and potential loss of control. Pilots must check wind forecasts against the drone's rating before every flight.}{catD}
\glsadd{IP-Rating}\dictentry{IP Rating}{An IP Rating (Ingress Protection) classifies a drone's enclosure protection against solids and liquids. Two digits indicate particle and moisture resistance respectively. IP54, for example, means limited dust protection and splash resistance from any direction. Agricultural and marine drones require higher IP ratings. IP rating is key for selecting drones for outdoor industrial applications.}{catD}
\glsadd{Drone-Swarm}\dictentry{Drone Swarm}{A Drone Swarm is a group of drones operating as a coordinated entity through autonomous communication and decentralized decisions. Each drone makes local decisions from neighboring aircraft data, producing emergent collective behavior. Swarms perform tasks impossible for single drones such as simultaneous large-area coverage. Applications span military strike, agricultural fleets, and light shows.}{catD}
\glsadd{Swarm-Intelligence}\dictentry{Swarm Intelligence}{Swarm Intelligence is collective behavior emerging from local interactions of simple agents following decentralized rules. Inspired by bird flocks and fish schools, each drone follows basic separation, alignment, and cohesion rules. These produce robust, adaptive group navigation without central control. Swarm algorithms are applied to drone coordination and logistics optimization.}{catD}
\glsadd{Formation-Flight}\dictentry{Formation Flight}{Formation Flight coordinates multiple drones maintaining specific spatial relationships along predetermined paths. Formations include lines, V-shapes, or grids via relative positioning algorithms. Military formations optimize sensor coverage while commercial ones enable simultaneous area coverage. Maintaining formation requires precise position sensing and rapid inter-drone communication.}{catD}
\glsadd{Cooperative-Navigation}\dictentry{Cooperative Navigation}{Cooperative Navigation has multiple drones sharing sensor data to improve the entire group's navigation accuracy. Each aircraft contributes observations fused with neighbors' data to build a richer environmental model. This reduces uncertainty in GPS-denied environments and improves localization. It is essential for indoor and urban canyon operations where satellite signals are lost.}{catD}
\glsadd{Distributed-Control}\dictentry{Distributed Control}{Distributed Control makes flight decisions locally by individual drones rather than a central station. Each aircraft runs its own autopilot and communicates with neighbors to coordinate. This eliminates single points of failure and adapts to member losses. It scales naturally to large agent counts. Distributed control is the standard paradigm for resilient drone swarm operations.}{catD}
\glsadd{Mesh-Network}\dictentry{Mesh Network}{A Mesh Network is a decentralized topology where every drone relays messages to every other, creating redundant data paths. Unlike hub-and-spoke, mesh networks reroute around lost nodes automatically. This makes them ideal for swarms in dynamic environments. Mesh networking enables low-latency communication for formation keeping and cooperative sensing in aerial networks.}{catD}
\glsadd{Multi-Agent-System}\dictentry{Multi-Agent System}{A Multi-Agent System (MAS) is a framework where multiple autonomous drones interact in a shared environment to achieve goals. Each agent perceives, decides locally, and communicates with others. MAS enables task allocation, negotiation, and coordination without centralized planning. It is used in drone fleets for logistics, monitoring, and security operations.}{catD}
\glsadd{Consensus-Algorithm}\dictentry{Consensus Algorithm}{A Consensus Algorithm enables multiple drones to reach agreement on shared values through message passing. Common algorithms include Raft, Paxos, and gossip protocols adapted for aerial networks. In swarms, consensus coordinates task selection, target assignment, and time synchronization. The algorithms must tolerate delays, packet loss, and disconnections in mobile aerial networks.}{catD}
\glsadd{Swarm-Coordination}\dictentry{Swarm Coordination}{Swarm Coordination manages multiple drones to achieve common objectives while avoiding conflicts and redundancies. It involves task allocation, path deconfliction, and temporal scheduling. Coordination can be centralized with a leader or decentralized through peer-to-peer negotiation. Effective coordination ensures the swarm's output exceeds individual drone capabilities.}{catD}
\glsadd{Collision-Avoidance-Swarm}\dictentry{Collision Avoidance Swarm}{Collision Avoidance Swarm prevents drones within a swarm from colliding while maintaining formation and performing tasks. Each drone uses relative position data from sensors or inter-drone communication to detect conflicts. Avoidance operates at multiple timescales from reactive maneuvers to path planning. The system must handle hundreds of agents without excessive computational overhead.}{catD}
\glsadd{ASTM-International}\dictentry{ASTM International}{ASTM International publishes technical standards for drone design, testing, and operations. Standards cover airworthiness, detect-and-avoid systems, and delivery performance. F3322 defines flight-over-people standards while F3003 establishes pilot competency. Compliance is often required for type certification. ASTM's consensus process ensures industry best practices and engineering rigor.}{catD}
\glsadd{RTCA}\dictentry{RTCA}{RTCA is a nonprofit developing consensus aviation recommendations including unmanned aircraft standards. Committee 228 produced the DO-365 detect-and-avoid standard and DO-362 UAS performance standards. These define requirements for drones to share airspace safely with manned aircraft. RTCA standards are referenced by the FAA for type certification of advanced platforms.}{catD}
\glsadd{Dronecode-Foundation}\dictentry{Dronecode Foundation}{The Dronecode Foundation is a Linux Foundation project supporting open-source drone software and hardware. It coordinates PX4, QGroundControl, MAVLink, and MAVSDK under unified governance. The foundation fosters an ecosystem building on open standards and interoperable protocols. It is the central hub for the open-source drone community.}{catD}
\glsadd{PX4-Community}\dictentry{PX4 Community}{The PX4 Community is the global network developing and supporting the PX4 autopilot software. Contributors develop features, fix bugs, write documentation, and provide forum support. Contributions span core flight algorithms to peripheral drivers and simulation tools. The community follows transparent development with public reviews. It is one of the most active open-source robotics communities.}{catD}
\glsadd{ArduPilot-Project}\dictentry{ArduPilot Project}{The ArduPilot Project is the community effort behind ArduPilot autopilot supporting multicopters, fixed-wings, helicopters, and rovers. It maintains extensive hardware support, parameter references, and an active support forum. Hundreds of contributors develop it under GPLv3. ArduPilot is one of the longest-running and most widely deployed open-source autopilot projects.}{catD}
\glsadd{Flying-Weight}\dictentry{Flying Weight}{Flying Weight is the total mass of a drone at takeoff including airframe, batteries, payloads, and accessories. It directly affects flight time, maneuverability, and regulatory compliance. Pilots must calculate flying weight before every mission to confirm certified limits. Changes in payload configuration require re-evaluation of performance parameters.}{catD}
\glsadd{Frame}\dictentry{Frame}{The Frame is the structural chassis providing the rigid platform to which all drone components attach. Frames use carbon fiber, fiberglass, aluminum, or plastic depending on weight and strength needs. Frame geometry defines drone type from quadcopter to octocopter layouts. It must withstand flight loads, vibrations, and landing impacts while maintaining motor alignment.}{catD}
\glsadd{Arm}\dictentry{Arm}{An Arm is a structural member connecting the drone's central frame to a motor and propeller assembly. A quadcopter typically has four arms arranged in X or + configuration. Each arm must be rigid enough to transmit thrust without flexing while remaining lightweight to minimize total mass. Some designs incorporate folding arms for compact transport. Arm length directly affects moment of inertia and maximum propeller size that can be fitted.}{catD}
\glsadd{Landing-Skid}\dictentry{Landing Skid}{Landing Skids are the support structures at a drone's base absorbing touchdown impact and stabilizing the aircraft on the ground. They are typically lightweight plastic, composite, or foam designed to flex on impact. Skids elevate the body above the ground, providing clearance for belly-mounted sensors. Some integrate vibration dampening. They protect the airframe during every landing cycle.}{catD}
\glsadd{Motor-Mount}\dictentry{Motor Mount}{A Motor Mount is the bracket attaching an electric motor to a drone arm. It holds the motor securely while allowing precise shaft alignment with the propeller. Typically machined from aluminum or molded from reinforced polymer, it transfers thrust and vibration to the arm. The mount must resist fatigue from millions of vibration cycles to prevent premature motor failure.}{catD}
\glsadd{Propeller-Guard}\dictentry{Propeller Guard}{A Propeller Guard is a protective ring surrounding a drone's propeller preventing blade contact with people or objects. Made from lightweight plastic or carbon fiber, guards add minimal weight. They are required for flights over people under Part 107 and standard on many consumer drones. Guards reduce injury risk during indoor and close-proximity operations with a slight aerodynamic cost.}{catD}
\glsadd{Camera-Mount}\dictentry{Camera Mount}{A Camera Mount is the mechanical interface securing a camera to a drone while providing vibration isolation and positioning. Simple mounts are rigid brackets; advanced versions use motorized gimbals with multi-axis stabilization. The mount must maintain camera orientation despite maneuvers and vibrations. Quick-release designs allow rapid camera swaps. Mount quality directly affects image sharpness.}{catD}
\glsadd{Vibration-Damping}\dictentry{Vibration Damping}{Vibration Damping isolates sensitive drone components from mechanical vibrations generated by motors and propellers. Uncontrolled vibrations degrade images, corrupt sensors, and cause fatigue. Solutions include elastomeric mounts, foam pads, and tuned mass dampers. The goal is attenuating vibrations at motor frequencies while maintaining rigid structural connections. Effective damping is essential for sensor fusion and imaging quality.}{catD}
\glsadd{Anti-Vibration-Mount}\dictentry{Anti-Vibration Mount}{An Anti-Vibration Mount reduces vibration transmission between sources and sensitive components. Designs use silicone dampers, rubber grommets, or spring-loaded platforms tuned to motor frequencies. These protect flight controller IMUs, cameras, and GPS modules from high-frequency oscillations. The mount must balance isolation with positional stability to prevent component shifting during flight.}{catD}
\glsadd{Foam-Damping}\dictentry{Foam Damping}{Foam Damping uses foam materials to absorb mechanical vibrations in drone systems. Closed-cell and open-cell foams of various densities are placed between vibrating components and electronics. Foam is lightweight, inexpensive, and effective across broad frequencies. It commonly mounts flight controllers and cameras. Foam damping is one of the most widely used vibration techniques in small drones, requiring no tuning or maintenance.}{catD}
\glsadd{GPS-Mast}\dictentry{GPS Mast}{A GPS Mast is a raised post elevating the GPS antenna above the drone's frame, motors, and electronics. Elevation reduces electromagnetic interference from motor controllers and power wiring. The mast positions the antenna in clear sky view. Typically made from carbon fiber or plastic, proper height and placement are critical for reliable satellite lock and accurate position hold.}{catD}
\glsadd{Compass-Mount}\dictentry{Compass Mount}{A Compass Mount positions the magnetometer away from magnetic interference from power wires, motors, and batteries. Compass interference causes heading errors degrading navigation in GPS-assisted modes. The mount is typically integrated with the GPS mast above the airframe. Calibration verifies proper mounting before each flight. Accurate heading data is essential for waypoint navigation and return-to-home.}{catD}
\glsadd{LED-Indicator}\dictentry{LED Indicator}{An LED Indicator is a visual status light on a drone communicating the aircraft's state to the pilot. Patterns indicate GPS lock, arming status, flight mode, battery level, and errors. Color codes vary by manufacturer. LEDs are essential for pre-flight verification and monitoring at distances where other feedback is invisible. Some drones use programmable RGB LEDs for customizable display.}{catD}
\glsadd{Buzzer}\dictentry{Buzzer}{A Buzzer is an audible alert producing tones for status changes, warnings, and errors. Functions include arming confirmation, low battery warnings, GPS notifications, and lost-model alarms. Buzzers sound at volumes audible from tens of meters, helping locate downed drones. Some connect to a transmitter button for activation. This simple component provides critical feedback when visual indicators are not visible.}{catD}
\glsadd{Current-Sensor}\dictentry{Current Sensor}{A Current Sensor measures real-time electrical current from the battery to propulsion and avionics systems. It uses hall-effect or shunt resistor mechanisms to detect amperage. The flight controller calculates consumed milliamp-hours and estimates remaining flight time. Accurate sensing enables low-battery warnings and return-to-home triggers. Calibration is important for reliable energy monitoring.}{catD}
\glsadd{Voltage-Sensor}\dictentry{Voltage Sensor}{A Voltage Sensor measures real-time battery voltage, providing data to estimate charge and detect low-voltage conditions. Voltage drops under load and recovers when load reduces, requiring rapid sampling. Low voltage alerts trigger return-to-home before crash damage. Voltage sensing enables cell-level monitoring in smart batteries. Accurate data is fundamental to safe drone operations.}{catD}
\glsadd{Power-Module}\dictentry{Power Module}{A Power Module provides regulated voltage and current sensing from the battery to the flight controller. It combines a voltage divider, current sensor, and regulator in a compact board. The module powers the flight controller while reporting battery telemetry. Standard in PX4 and ArduPilot systems, it connects directly between battery and flight controller, simplifying wiring and providing calibrated energy monitoring.}{catD}
\glsadd{Pixhawk}\dictentry{Pixhawk}{Pixhawk is a family of open-source flight controller hardware for autonomous drone flight. It features redundant IMUs, a powerful processor, and extensive I/O interfaces. Pixhawk runs PX4 and ArduPilot firmware supported by QGroundControl. Multiple generations from Pixhawk 1 through 6X offer improved sensors and processing. It is the most widely deployed open-source flight controller.}{catD}
\glsadd{Cube}\dictentry{Cube}{Cube (formerly CubePilot) is a modular flight controller using a stackable enclosure architecture. The Cube module containing processor and sensors plugs into carrier boards providing power and I/O. This design allows swapping between carrier boards without rewiring. Cube supports PX4 and ArduPilot firmware for commercial, military, and research applications. Its ruggedized design and hot-swap capability suit demanding professional platforms.}{catD}
\glsadd{Kakute}\dictentry{Kakute}{The Kakute is a line of flight controllers designed for racing and freestyle drones. Manufactured by Holybro, Kakute boards feature integrated OSD, current sensing, and support for high PID loop rates. They are known for reliability and clean layout, making them popular among competitive pilots. The series supports Betaflight and EmuFlight, offering flexibility across disciplines.}{catD}
\glsadd{F4-Processor}\dictentry{F4 Processor}{The F4 Processor is an STM32F4-series microcontroller widely used in drone flight controllers. It runs at up to 168 MHz with a floating-point unit, enabling fast PID loop calculations and sensor fusion. The F4 was a major leap from the older F1 and F3 chips, supporting higher loop rates, more features, and dual gyro sampling. It remains a cost-effective choice for pilots who do not need the absolute highest performance.}{catD}
\glsadd{F7-Processor}\dictentry{F7 Processor}{The F7 Processor is an STM32F7-series microcontroller found in high-end flight controllers. Running at up to 216 MHz with a Cortex-M7 core and larger RAM, it handles complex filtering and high PID loop rates with ease. The F7 enables advanced features like RPM filtering and bidirectional DShot without performance compromise. It offers a noticeable upgrade in processing headroom over the F4 for demanding builds.}{catD}
\glsadd{H7-Processor}\dictentry{H7 Processor}{The H7 Processor is an STM32H7-series microcontroller representing the top tier of flight controller hardware. With a dual-core Cortex-M7 running at up to 480 MHz, it provides massive headroom for filtering, logging, and telemetry simultaneously. The H7 supports the highest PID loop rates and complex filtering without performance compromise. It is the preferred choice for cutting-edge racing and research platforms.}{catD}
\glsadd{Betaflight}\dictentry{Betaflight}{Betaflight is the most widely used open-source firmware for racing and freestyle drones. It provides a highly tunable PID controller, advanced filtering pipeline, and extensive configurator software for setup. Betaflight supports a vast range of flight controllers, ESCs, and peripherals, making it the default choice for the FPV community. Regular updates continually improve flight performance and introduce new features.}{catD}
\glsadd{INAV}\dictentry{INAV}{INAV is an open-source firmware focused on GPS-assisted navigation and autonomous flight for multirotors and fixed-wing aircraft. It provides position hold, return to launch, waypoint missions, and altitude hold with barometric and GPS fusion. INAV diverges from Cleanflight to prioritize navigation over pure acrobatic performance. It is a popular choice for long-range builds and autonomous platforms.}{catD}
\glsadd{Cleanflight}\dictentry{Cleanflight}{Cleanflight is an early open-source flight controller firmware that laid the groundwork for modern drone software. It forked from MultiWii and introduced a modular architecture supporting F1 and F3 processors. Cleanflight pioneered configurable PID tuning, serial passthrough, and LED strip control that later became standard. Though superseded by Betaflight, its codebase was the direct ancestor of several major forks.}{catD}
\glsadd{Kiss}\dictentry{Kiss}{Kiss is a commercial flight controller firmware developed by Flyduino for their proprietary KISS FC hardware. It is known for exceptionally smooth flight characteristics and a streamlined setup requiring minimal tuning. The KISS ecosystem includes integrated ESC firmware, creating a tightly coupled platform. It has a dedicated following among pilots who value simplicity and polish over open-source customization.}{catD}
\glsadd{ButterFlight}\dictentry{ButterFlight}{ButterFlight is a fork of Betaflight optimized specifically for smooth and cinematic flight feel. It modifies filtering, PID scaling, and feedforward algorithms to reduce jello and vibration artifacts in video. ButterFlight appeals to pilots who prioritize smoothness over aggressive responsiveness. The project demonstrates how firmware forks can target niche performance characteristics within the broader community.}{catD}
\glsadd{EmuFlight}\dictentry{EmuFlight}{EmuFlight is a fork of Betaflight focused on improving flight controller performance and reducing latency. It incorporates modified filtering, updated PID algorithms, and optimizations for newer processor architectures. EmuFlight aims to squeeze out incremental gains that matter in competitive racing environments. The project maintains compatibility with the Betaflight configurator while offering a distinct tuning philosophy.}{catD}
\glsadd{Raceflight}\dictentry{Raceflight}{Raceflight is a firmware designed specifically for racing drones with an emphasis on minimal latency and maximum responsiveness. It was one of the first firmwares to target sub-millisecond loop times on dedicated racing hardware. Raceflight features a simplified tuning approach and aggressive default profiles suited to competitive racing. While its active development has slowed, it influenced many design decisions in later firmware projects.}{catD}
\glsadd{ExpressLRS}\dictentry{ExpressLRS}{ExpressLRS is an open-source radio control link achieving ultra-low latency and long range using LoRa modulation. It operates on 2.4 GHz or 900 MHz bands and supports packet rates up to 500 Hz or higher. The open-source nature allows hobbyists to build and flash their own transmitters and receivers at minimal cost. It has rapidly become the dominant RC link in the FPV community.}{catD}
\glsadd{Crossfire}\dictentry{Crossfire}{Crossfire is a long-range radio control link developed by Team BlackSheep operating on the 868 or 915 MHz band. It provides reliable control at ranges exceeding 40 kilometers with low latency telemetry. Crossfire uses LoRa modulation for robust signal integrity even in challenging RF environments. It was one of the first commercial long-range links to gain widespread adoption in the FPV community.}{catD}
\glsadd{Tracer}\dictentry{Tracer}{Tracer is a low-latency radio control link from Team BlackSheep optimized for racing and proximity flying. It operates on the 2.4 GHz band and achieves end-to-end latency as low as 2.5 milliseconds. Tracer combines the reliability of Spread Spectrum technology with the speed demands of competitive racing. It targets pilots who want Crossfire-level reliability with racing-oriented response times.}{catD}
\glsadd{Ghost}\dictentry{Ghost}{Ghost is an ultra-low latency radio control link from TBS that operates on the 2.4 GHz band. It uses proprietary modulation to achieve sub-two-millisecond latency while maintaining long-range capability. Ghost is designed for pilots who demand the absolute fastest control response without sacrificing link reliability. The protocol supports bidirectional telemetry for real-time link quality monitoring.}{catD}
\glsadd{DSMX}\dictentry{DSMX}{DSMX is a frequency-hopping spread spectrum protocol developed by Spektrum for RC control links. It operates in the 2.4 GHz band and provides interference-resistant communication between transmitter and receiver. DSMX supports multiple aircraft binding and bidirectional telemetry for link quality and battery data. It is the standard protocol across Spektrum's ecosystem of transmitters and receivers.}{catD}
\glsadd{SBUS}\dictentry{SBUS}{SBUS is a serial bus protocol originally developed by Futaba for connecting RC receivers to flight controllers. It carries up to 16 or 18 channels over a single digital signal wire using inverted UART communication. SBUS is supported by virtually all modern flight controllers and provides low-latency channel data. Its simplicity and widespread compatibility have made it one of the most commonly used RC protocols.}{catD}
\glsadd{CRSF}\dictentry{CRSF}{CRSF, or Crossfire Serial Protocol, is the communication standard used between TBS receivers and flight controllers. It carries RC channel data, telemetry, and configuration commands over a single UART connection. CRSF supports high update rates and bidirectional communication, enabling real-time telemetry back to the transmitter. It has become the dominant protocol for high-performance FPV builds due to its speed and feature set.}{catD}
\glsadd{PWM}\dictentry{PWM}{PWM, or Pulse Width Modulation, is a signaling method that encodes a value in the width of a repeated pulse. In drones, PWM is used to send throttle and control commands from the flight controller to ESCs and servos. Each channel requires its own signal wire, making it less efficient than serial protocols. Despite its simplicity, PWM remains relevant for servo control and legacy ESC connections.}{catD}
\glsadd{DShot}\dictentry{DShot}{DShot is a digital ESC communication protocol that sends throttle commands as encoded bit patterns. Unlike analog protocols, DShot is immune to electrical noise and requires no calibration. It supports bidirectional variants that return telemetry data such as RPM, temperature, and current. DShot has become the standard protocol for modern ESCs due to its precision and reliability.}{catD}
\glsadd{Multishot}\dictentry{Multishot}{Multishot is an analog ESC protocol that encodes throttle values using pulse timing at higher resolution than PWM. It operates at much shorter cycle times, enabling faster communication between flight controller and ESC. Multishot was a transitional protocol between PWM and digital standards like DShot. It offered improved resolution and speed but has largely been replaced by DShot in modern builds.}{catD}
\glsadd{Oneshot}\dictentry{Oneshot}{Oneshot is an analog ESC protocol that sends a single pulse per command cycle, reducing the time between updates. It provides faster and more consistent throttle response compared to traditional PWM. Oneshot was one of the first high-frequency ESC protocols and popularized the move toward faster digital communication. It served as an important stepping stone in the evolution of ESC protocols.}{catD}
\glsadd{Bidirectional-DShot}\dictentry{Bidirectional DShot}{Bidirectional DShot extends the DShot protocol to allow ESCs to send telemetry data back to the flight controller. It enables real-time motor RPM reporting, which powers RPM-based dynamic notch filtering for vibration suppression. The bidirectional communication uses the same signal wire, requiring compatible ESC firmware. This feature has become essential for achieving the cleanest possible flight performance.}{catD}
\glsadd{RPM-Filtering}\dictentry{RPM Filtering}{RPM Filtering uses motor RPM data from bidirectional DShot to dynamically track and notch out vibration frequencies generated by the motors. As motor speeds change during flight, the filter adapts in real time to suppress the dominant noise peaks. This approach is far more effective than static notch filters, which cannot track shifting frequencies. RPM filtering is one of the most significant advances in drone flight smoothness in recent years.}{catD}
\glsadd{Dynamic-Filter}\dictentry{Dynamic Filter}{The Dynamic Filter is an adaptive noise filter that automatically detects and removes gyro noise peaks in real time. It uses a self-tuning notch filter that sweeps the frequency spectrum to identify dominant vibrations. The filter adjusts continuously as motor speeds and aerodynamic loads change during flight. It dramatically reduces tuning effort and improves flight quality across different configurations.}{catD}
\glsadd{Notch-Filter}\dictentry{Notch Filter}{A Notch Filter is a signal processing filter that attenuates a narrow band of frequencies while passing all others. In drones, notch filters target specific vibration frequencies generated by motors and propellers. Static notch filters are configured manually at known frequencies, while dynamic notch filters track shifting noise peaks automatically. Proper notch filtering is critical for stable flight and clean video output.}{catD}
\glsadd{Low-Pass-Filter}\dictentry{Low Pass Filter}{A Low Pass Filter allows low-frequency signals to pass through while attenuating higher frequencies. In drone flight controllers, low pass filters are applied to gyro and accelerometer data to remove high-frequency vibration noise. The cutoff frequency determines how much noise is removed versus how much signal delay is introduced. Balancing filter strength against latency is one of the key tuning challenges in drone firmware.}{catD}
\glsadd{Gyro-Sampling}\dictentry{Gyro Sampling}{Gyro Sampling refers to the rate at which the gyroscope sensor is read by the flight controller's processor. Higher sampling rates provide more data points for the PID loop, enabling smoother and more responsive control. Modern flight controllers sample the gyro at 8 kHz or higher, with some supporting dual-gyro configurations. The sampling rate directly impacts the maximum achievable PID loop frequency and filtering resolution.}{catD}
\glsadd{PID-Loop-Frequency}\dictentry{PID Loop Frequency}{The PID Loop Frequency is the rate at which the flight controller calculates and outputs control corrections, measured in Hz. Higher frequencies allow the controller to react faster to disturbances and maintain tighter flight characteristics. Common rates range from 4 kHz to 8 kHz, with some systems supporting even higher. The loop frequency is limited by the processor speed, sensor sampling rate, and filtering overhead.}{catD}
\glsadd{Desync}\dictentry{Desync}{Desync occurs when an ESC loses synchronization with the motor's rotor position, causing erratic motor behavior or a complete stall. It typically happens during rapid throttle changes or at extreme motor timing settings. A desync often manifests as a sudden loss of thrust on one motor, resulting in a spin or tumble. Proper ESC calibration, motor timing, and throttle ramp settings help prevent desync events.}{catD}
\glsadd{Motor-Timing}\dictentry{Motor Timing}{Motor Timing refers to the electrical advance angle at which the ESC energizes the next motor phase during commutation. Higher timing advances improve motor responsiveness and torque but reduce efficiency and increase heat. Lower timing improves efficiency at the cost of reduced high-RPM performance. Finding the optimal timing setting depends on the motor, propeller, and flight application.}{catD}
\glsadd{Demag-Compensation}\dictentry{Demag Compensation}{Demag Compensation is an ESC feature that accounts for the voltage spike caused by motor back-EMF during commutation. Without compensation, rapid throttle changes can cause demagnetization of the motor magnets or timing errors. The ESC adjusts its switching behavior to safely handle these voltage transients. Proper demag compensation is essential for preventing desyncs during aggressive flying.}{catD}
\glsadd{Throttle-Scale}\dictentry{Throttle Scale}{Throttle Scale is a firmware parameter that maps the throttle input range to a different output range for more linear motor response. It compensates for the non-linear thrust curve inherent in most propeller and motor combinations. By adjusting the throttle scale, pilots can achieve more predictable power delivery across the throttle range. This parameter is particularly useful for matching different battery voltages or propeller sizes.}{catD}
\glsadd{Throttle-Limit}\dictentry{Throttle Limit}{Throttle Limit caps the maximum throttle output sent to the motors regardless of transmitter input. It serves as a safety feature to prevent over-powering the aircraft or exceeding structural limits. Pilots may set a lower throttle limit when learning to fly or when operating in confined spaces. The limit can be configured as a percentage of total available thrust.}{catD}
\glsadd{Motor-Output}\dictentry{Motor Output}{Motor Output is the signal sent by the flight controller to the ESC commands for each motor. It represents the PID controller's calculated thrust requirement for each motor after mixing. The output is a normalized value between minimum and maximum that the ESC translates into actual motor power. Consistent and accurate motor output is fundamental to stable and responsive flight.}{catD}
\glsadd{Direction-Lock}\dictentry{Direction Lock}{Direction Lock is a firmware feature that fixes the rotation direction of each motor in software rather than relying on physical wiring. By configuring direction lock, pilots can reverse motor rotation without resoldering ESC signal wires. This feature is especially useful during field repairs or when swapping motors between positions. It adds convenience and reduces the risk of wiring errors.}{catD}
\glsadd{Resource-Remapping}\dictentry{Resource Remapping}{Resource Remapping allows pilots to reassign hardware pin functions on the flight controller through software configuration. If a pad is damaged or a different pin layout is needed, resource remapping can redirect signals to alternative pins. This feature provides flexibility to salvage damaged boards or adapt to custom wiring configurations. It is a powerful tool for advanced builders working with non-standard hardware.}{catD}
\glsadd{Serial-Port}\dictentry{Serial Port}{A Serial Port is a communication interface that transmits data sequentially, one bit at a time, between the flight controller and external devices. Drones use serial ports to connect GPS modules, telemetry radios, ESCs, and other peripherals. Each serial port has a configurable baud rate that determines its communication speed. Modern flight controllers typically offer multiple UART serial ports to support several peripherals simultaneously.}{catD}
\glsadd{I2C}\dictentry{I2C}{I2C, or Inter-Integrated Circuit, is a two-wire communication bus used to connect low-speed peripherals to the flight controller. It uses a data line and a clock line to transfer addresses and data between a master and multiple slave devices. I2C is commonly used for barometers, magnetometers, and OLED displays on drones. While slower than SPI, I2C requires fewer pins and supports multiple devices on the same bus.}{catD}
\glsadd{SPI}\dictentry{SPI}{SPI, or Serial Peripheral Interface, is a high-speed synchronous communication bus between the flight controller and its peripherals. It uses separate clock, data, and chip-select lines to achieve fast data transfer rates. SPI is used for gyroscopes, accelerometers, and other time-critical sensors that demand low latency. Its speed advantage makes it the preferred bus for primary flight sensors.}{catD}
\glsadd{UART}\dictentry{UART}{UART, or Universal Asynchronous Receiver-Transmitter, is a serial communication interface that transmits data without a shared clock signal. In drones, UARTs connect GPS modules, telemetry radios, RC receivers, and other serial devices. Each UART requires two wires for bidirectional communication at a configured baud rate. Flight controllers typically provide multiple UARTs to support various simultaneous connections.}{catD}
\glsadd{CAN-Bus}\dictentry{CAN Bus}{CAN Bus, or Controller Area Network, is a robust serial communication standard originally developed for automotive applications. In drones, it connects ESCs, GPS modules, and other peripherals on a shared two-wire bus with built-in error detection. CAN provides reliable communication in electrically noisy environments where UART or I2C may fail. It is increasingly adopted in professional and industrial drone platforms for its reliability.}{catD}
\glsadd{SD-Card-Logging}\dictentry{SD Card Logging}{SD Card Logging is the process of recording flight data to an onboard SD card for later analysis. It captures high-frequency sensor data, PID outputs, motor commands, and other diagnostic information. SD card logging provides far more data bandwidth than Blackbox flash memory, enabling longer and higher-resolution recordings. Pilots and engineers use these logs to diagnose vibration issues, tune PID controllers, and verify system performance.}{catD}
\glsadd{Blackbox}\dictentry{Blackbox}{Blackbox is the flight data logging system built into Betaflight and related firmware. It records gyroscope, PID, motor, and RC input data at the PID loop rate for post-flight analysis. Blackbox data can be visualized using the Blackbox Explorer tool to identify vibration issues and fine-tune filter and PID settings. The logs are stored in onboard flash memory or on an SD card depending on the flight controller hardware.}{catD}
\glsadd{OSD}\dictentry{OSD}{OSD, or On-Screen Display, overlays flight data such as battery voltage, altitude, GPS coordinates, and flight time onto the FPV video feed. It provides the pilot with real-time situational awareness without looking away from the video. OSD is generated by the flight controller and injected into the analog or digital video signal. It is an essential tool for both casual flying and long-range missions.}{catD}
\glsadd{Betaflight-OSD}\dictentry{Betaflight OSD}{Betaflight OSD is the on-screen display system integrated into Betaflight firmware. It renders configurable text overlays including voltage, RSSI, timer, and artificial horizon on the pilot's video feed. The OSD is customized through the Betaflight Configurator, where pilots can position and select which elements to display. It uses a dedicated MAX7456 chip on the flight controller to overlay text onto analog video signals.}{catD}
\glsadd{DJI-OSD}\dictentry{DJI OSD}{DJI OSD is the on-screen display provided by DJI's digital FPV system including the Air Unit and Vista. It renders flight data overlays through the digital video pipeline rather than through a dedicated chip. DJI OSD elements are configured through the DJI goggles interface and are limited to DJI-supported parameters. It provides a clean digital overlay optimized for DJI's ecosystem of cameras and goggles.}{catD}
\glsadd{FPV-Camera}\dictentry{FPV Camera}{An FPV Camera is a small, lightweight camera that transmits live video to the pilot goggles or monitor. It is optimized for low latency rather than image quality, ensuring the pilot sees events in near real-time. FPV cameras come in various sensor sizes and lens options for different flying styles. The camera is the pilot eyes in the air, making it one of the most critical FPV components.}{catD}
\glsadd{VTX}\dictentry{VTX}{VTX, or Video Transmitter, is the device on the drone that broadcasts the FPV video signal to a ground receiver. It converts the camera's video output into a radio frequency signal on the 5.8 GHz band or other supported frequencies. VTX power output typically ranges from 25 mW to 800 mW, with higher power extending range at the cost of heat and battery drain. Proper antenna selection and power management are essential for reliable video links.}{catD}
\glsadd{Video-Receiver}\dictentry{Video Receiver}{A Video Receiver is a ground-side device that picks up the FPV video signal broadcast by the drone's VTX. It demodulates the RF signal into a video feed displayed on goggles or a monitor. Receivers come in single-channel and scanning variants, with diversity receivers offering the best performance through multiple antennas. Video receiver quality directly impacts the clarity and range of the pilot's video link.}{catD}
\glsadd{FPV-Goggles}\dictentry{FPV Goggles}{FPV Goggles are head-mounted displays that show the live video feed from the drone's FPV camera, immersing the pilot in the flight experience. They contain built-in receivers, displays, and often DVR recording capability for capturing flight footage. Goggles range from basic analog models to high-end digital systems with HD resolution and head tracking. They are the primary interface through which FPV pilots experience flight.}{catD}
\glsadd{FPV-Monitor}\dictentry{FPV Monitor}{An FPV Monitor is a standalone screen used to view the drone's live video feed as an alternative to FPV goggles. It typically includes a built-in receiver and battery for portable, hands-free viewing. Monitors are useful for spectators, spotter assistance, and pilots who prefer not to wear goggles. They provide a shared viewing experience that goggles cannot offer.}{catD}
\glsadd{Diversity-Receiver}\dictentry{Diversity Receiver}{A Diversity Receiver uses two or more antennas on different polarization or spatial orientations to improve FPV video reception. It continuously monitors signal strength on each antenna and switches to the stronger one in real time. This approach significantly reduces video dropouts caused by multipath interference and antenna orientation mismatches. Diversity reception is standard equipment for reliable FPV video links.}{catD}
\glsadd{Antenna-Polarization}\dictentry{Antenna Polarization}{Antenna Polarization describes the orientation of the electric field in the electromagnetic wave radiated by the antenna. Matching polarization between transmitter and receiver antennas is critical for maximum signal transfer. Mismatched polarization can result in up to 30 dB of signal loss, drastically reducing video range. Common FPV antenna polarizations include linear, circular left-hand, and circular right-hand.}{catD}
\glsadd{Linear-Polarization}\dictentry{Linear Polarization}{Linear Polarization means the antenna radiates its electric field in a single fixed plane, either horizontal or vertical. Linearly polarized antennas are simple and inexpensive to manufacture. However, they are sensitive to antenna orientation and suffer significant signal loss when the drone banks or tilts relative to the ground antenna. For this reason, linear polarization is less commonly used for FPV video than circular polarization.}{catD}
\glsadd{Circular-Polarization}\dictentry{Circular Polarization}{Circular Polarization means the antenna radiates an electric field that rotates as the wave propagates through space. This rotation makes circular antennas more tolerant of orientation changes and multipath interference compared to linear antennas. Circular polarization is the standard for FPV video links because it provides consistent signal quality during dynamic flight. Antennas come in left-hand and right-hand variants that must match.}{catD}
\glsadd{RHCP}\dictentry{RHCP}{RHCP, or Right-Hand Circular Polarization, is a circular polarization where the electric field rotates clockwise as viewed from the receiver. All antennas in an FPV link must share the same polarization for optimal performance. Mixing RHCP and LHCP antennas on the same link causes severe signal loss. RHCP is one of the two standard polarizations used in the FPV community.}{catD}
\glsadd{LHCP}\dictentry{LHCP}{LHCP, or Left-Hand Circular Polarization, is a circular polarization where the electric field rotates counterclockwise as viewed from the receiver. It is functionally identical to RHCP but with opposite rotation direction. LHCP and RHCP are not compatible on the same link and must be matched between VTX and receiver antenna. Both polarizations perform equally well when properly matched.}{catD}
\glsadd{SMA-Connector}\dictentry{SMA Connector}{The SMA Connector is a threaded coaxial RF connector used to attach antennas to VTX modules and video receivers. It features a center pin surrounded by a threaded barrel that provides a secure, vibration-resistant connection. SMA connectors are rated for frequencies up to 18 GHz and are widely used in FPV equipment. Their threaded design prevents accidental disconnection during flight vibrations.}{catD}
\glsadd{RP-SMA}\dictentry{RP-SMA}{RP-SMA, or Reverse Polarity SMA, is a variant of the SMA connector with the center pin and socket reversed. The threaded barrel remains the same, but the gender of the center contact is swapped. RP-SMA connectors are common on WiFi equipment and some FPV gear. It is essential to verify connector type before purchasing antennas, as SMA and RP-SMA are physically incompatible despite looking similar.}{catD}
\glsadd{MMCX-Connector}\dictentry{MMCX Connector}{The MMCX Connector is a small push-on coaxial RF connector designed for quick antenna attachment and removal. It provides a secure snap-fit connection without the need for threading, making field swaps fast and tool-free. MMCX is increasingly popular on modern VTX modules and receivers for its compact size and ease of use. It supports frequencies up to 6 GHz with low signal loss.}{catD}
\glsadd{UFL-Connector}\dictentry{UFL Connector}{The UFL Connector, also known as IPEX or MHF, is an extremely small surface-mount RF connector used in compact electronics. In FPV drones, it connects internal antennas to VTX modules and video receivers on the circuit board. UFL connectors are delicate and not designed for repeated mating cycles, making them suitable for permanent installations. Their tiny footprint enables the miniaturization of FPV equipment.}{catD}
\glsadd{IPX-Connector}\dictentry{IPX Connector}{The IPX Connector is a miniature push-on RF connector similar in function to the UFL but with a slightly different mechanical design. It provides a compact, low-profile connection for antennas on small circuit boards. IPX connectors are commonly found on VTX modules, receivers, and other compact FPV electronics. They offer reliable RF connections in applications where space is at a premium.}{catD}
\glsadd{LQ}\dictentry{LQ}{LQ, or Link Quality, is a metric that indicates the reliability of the radio control link between transmitter and receiver. It is typically expressed as a percentage or a count of successful versus total packets received. A high LQ value indicates a strong, stable connection, while a low value signals potential range or interference issues. Pilots monitor LQ to gauge how safely they can fly at distance or in challenging RF environments.}{catD}
\glsadd{SNR}\dictentry{SNR}{SNR, or Signal-to-Noise Ratio, measures the strength of the desired signal relative to background noise. A higher SNR indicates a cleaner, more reliable communication link. In FPV systems, SNR is relevant for both the video and radio control links, directly impacting range and quality. Understanding SNR helps pilots diagnose interference issues and optimize antenna placement.}{catD}
\glsadd{Latency}\dictentry{Latency}{Latency is the time delay between a physical event occurring and its corresponding output being perceived. In FPV drones, it encompasses the chain from camera capture to display and from pilot input to motor response. Lower latency provides a more immediate and intuitive flying experience, critical for racing and proximity flying. Total system latency is the sum of delays from every component in the signal chain.}{catD}
\glsadd{End-to-End-Latency}\dictentry{End to End Latency}{End to End Latency is the total delay from the moment the FPV camera captures a frame to the moment it appears on the pilot's goggles or monitor. It includes camera processing, encoding, transmission, reception, decoding, and display times. Minimizing end-to-end latency is a primary engineering goal in FPV video systems. Digital FPV systems have made significant progress, achieving latencies comparable to or better than analog systems.}{catD}
\glsadd{Frame-Rate}\dictentry{Frame Rate}{Frame Rate is the number of individual video frames displayed per second, measured in frames per second (fps). Higher frame rates produce smoother video motion and reduce perceived latency. FPV cameras typically operate at 30 fps for analog systems and up to 60 fps or higher for digital systems. The frame rate interacts with display refresh rate to determine the overall smoothness of the video experience.}{catD}
\glsadd{Resolution}\dictentry{Resolution}{Resolution defines the pixel dimensions of a video image, commonly expressed as 720p, 1080p, or 4K. Higher resolution provides sharper, more detailed images that allow pilots to see finer obstacles at distance. In FPV, resolution must be balanced against latency and bandwidth constraints of the transmission link. Digital FPV systems now offer HD resolutions that dramatically improve visual clarity over traditional analog.}{catD}
\glsadd{Bitrate}\dictentry{Bitrate}{Bitrate is the amount of data transmitted per second in a video stream, typically measured in megabits per second (Mbps). Higher bitrates allow more detailed video with fewer compression artifacts. FPV video bitrates are constrained by available RF bandwidth and transmission power. Digital FPV systems use adaptive bitrate algorithms to maintain video quality as link conditions change.}{catD}
\glsadd{Codec}\dictentry{Codec}{A Codec, short for coder-decoder, compresses video for transmission and decompresses it for display. Video codecs reduce massive raw data rates into manageable bitrates suitable for wireless transmission. Common FPV codecs include H.264 and H.265, each offering different tradeoffs between efficiency and processing requirements. The choice of codec directly affects video quality, latency, and hardware requirements.}{catD}
\glsadd{H264}\dictentry{H.264}{H.264, also known as AVC, is a widely adopted video compression standard used in many FPV digital video systems. It provides good compression efficiency with relatively low encoding and decoding complexity. H.264 enables HD video transmission at reasonable bitrates while maintaining acceptable latency. It remains the most commonly supported codec across FPV equipment manufacturers.}{catD}
\glsadd{H265}\dictentry{H.265}{H.265, also known as HEVC, is an advanced video compression standard that achieves roughly double the compression efficiency of H.264. It produces higher quality video at the same bitrate or the same quality at half the bitrate. H.265 requires more processing power for encoding and decoding, which can increase latency. Newer FPV systems are adopting H.265 to deliver better video quality within bandwidth constraints.}{catD}
\glsadd{Mavlink-Protocol}\dictentry{Mavlink Protocol}{MAVLink is a lightweight messaging protocol designed for communication between unmanned vehicles and ground control stations. It defines a standardized set of messages for telemetry, commands, and parameter exchange across different autopilot systems. MAVLink supports multiple transport layers including serial, UDP, and TCP connections. It is the backbone of communication for ArduPilot, PX4, and many other autopilot platforms.}{catD}
\glsadd{Mavlink-2}\dictentry{Mavlink 2}{MAVLink 2 is the second major version of the MAVLink protocol, introducing improved security and flexibility over the original version. It adds message signing for authentication, variable-length payloads, and support for larger system IDs. MAVLink 2 is backward compatible with MAVLink 1, allowing gradual migration. The security enhancements make it suitable for applications requiring verified command integrity.}{catD}
\glsadd{Mavlink-Inspector}\dictentry{Mavlink Inspector}{A Mavlink Inspector is a software tool that monitors and displays raw MAVLink messages flowing between a drone and ground station. It decodes message fields, timestamps, and source addresses for debugging and analysis. Pilots and developers use Mavlink Inspectors to verify telemetry data, diagnose communication issues, and validate custom messages. It is an essential diagnostic tool for anyone developing or troubleshooting MAVLink-based systems.}{catD}
\glsadd{Parameter}\dictentry{Parameter}{A Parameter is a configurable value stored on the autopilot that controls some aspect of its behavior. Parameters include settings for PID gains, flight mode options, failsafe thresholds, and hardware configurations. They can be read, written, and saved using ground control station software. Proper parameter management is fundamental to drone setup, tuning, and safe operation.}{catD}
\glsadd{Mission}\dictentry{Mission}{A Mission is a predefined sequence of waypoints and actions that the drone executes autonomously. It specifies locations, altitudes, speeds, and specific actions such as taking photos or dropping payloads at each waypoint. Missions are planned in ground control station software and uploaded to the autopilot before flight. They enable repeatable autonomous operations for surveying, mapping, and delivery applications.}{catD}
\glsadd{Flight-Mode}\dictentry{Flight Mode}{A Flight Mode defines the autopilot's control strategy and level of autonomy during flight. Common modes include manual, stabilize, altitude hold, position hold, and auto. Each mode determines which sensors are used and how much authority the autopilot has over pilot inputs. Selecting the appropriate flight mode is essential for matching the drone's behavior to the current mission requirements.}{catD}
\glsadd{Arming}\dictentry{Arming}{Arming is the process of enabling the drone's motors so they can respond to throttle input. It is a critical safety feature that prevents accidental motor spins while handling the aircraft. Arming typically requires a specific stick combination, switch position, or button press. Once armed, the drone is ready to fly and all control surfaces and motors become active.}{catD}
\glsadd{Disarming}\dictentry{Disarming}{Disabling the motors after flight by deactivating the armed state is called disarming. It ensures the motors cannot spin accidentally while the pilot handles or transports the aircraft. Disarming is performed with a specific stick command, switch position, or after landing detection. Proper disarm procedure is an essential part of post-flight safety protocols.}{catD}
\glsadd{Pre-Arm-Check}\dictentry{Pre-Arm Check}{Pre-Arm Checks are automated safety verifications that run before the autopilot allows arming. They verify that sensors are calibrated, GPS has a fix, battery voltage is sufficient, and no critical faults exist. Failed checks prevent arming and display error messages to the pilot through the ground station or OSD. These checks are a vital safeguard against flying in an unsafe or misconfigured state.}{catD}
\glsadd{Calibration}\dictentry{Calibration}{Calibration is the process of determining sensor offsets and scale factors so the autopilot reads accurate data. It typically involves placing the aircraft in known orientations and recording sensor values to compute corrections. Calibration must be performed after hardware changes or when sensor readings drift. Accurate calibration is foundational to reliable autonomous flight and precise control.}{catD}
\glsadd{Level-Calibration}\dictentry{Level Calibration}{Level Calibration sets the accelerometer's zero point by recording its output when the aircraft is perfectly level. The pilot places the drone on a known flat surface and initiates the calibration procedure. This ensures the autopilot correctly understands what level means for altitude hold and stabilized flight modes. An incorrect level calibration results in the drone drifting or holding incorrect angles.}{catD}
\glsadd{Compass-Calibration}\dictentry{Compass Calibration}{Compass Calibration determines the offsets and scale factors for the magnetometer to provide accurate heading information. The procedure requires rotating the aircraft through multiple orientations to map the local magnetic field. Hard and soft iron distortions from the aircraft's structure and electronics are compensated through this process. Accurate compass calibration is essential for GPS-assisted flight modes and reliable heading hold.}{catD}
\glsadd{ESC-Calibration}\dictentry{ESC Calibration}{ESC Calibration synchronizes the throttle range between the flight controller and all electronic speed controllers. It teaches each ESC the minimum and maximum pulse widths that correspond to zero and full throttle. Proper calibration ensures all motors respond uniformly to the same throttle command. Uneven calibration can cause asymmetric thrust and unstable flight characteristics.}{catD}
\glsadd{Radio-Calibration}\dictentry{Radio Calibration}{Radio Calibration records the minimum, maximum, and center pulse widths for each transmitter channel. The pilot moves each stick and switch through its full range while the autopilot captures the endpoints. This ensures the autopilot accurately interpreates pilot inputs for all control axes. Proper radio calibration is a prerequisite for responsive and predictable flight control.}{catD}
\glsadd{Throttle-Calibration}\dictentry{Throttle Calibration}{Throttle Calibration specifically maps the throttle channel's range from minimum to maximum output. It ensures the autopilot correctly interpreates the full range of throttle commands from the transmitter. This is important for consistent hover throttle point and maximum thrust delivery. Accurate throttle calibration contributes to predictable power response across all flight conditions.}{catD}
\glsadd{Safety-Switch}\dictentry{Safety Switch}{A Safety Switch is a physical button or switch that must be pressed or toggled before the drone can be armed. It adds an extra layer of protection against accidental motor startup during handling or transport. The safety switch is commonly found on boards with integrated power distribution. It provides a simple but effective hardware-level safety interlock.}{catD}
\glsadd{Beacon}\dictentry{Beacon}{A Beacon is an audible or visual signal activated to help locate a downed or lost drone. It typically uses the VTX's LED, motor beeper, or a dedicated buzzer to produce a loud sound or flashing light. Beacons are often automatically triggered by failsafe conditions or activated remotely through the transmitter. They are invaluable for finding crashed aircraft in tall grass, dense vegetation, or remote areas.}{catD}
\glsadd{Return-to-Launch}\dictentry{Return to Launch}{Return to Launch, or RTL, is a failsafe or pilot-initiated mode that commands the drone to fly back to its takeoff position autonomously. It uses GPS coordinates to navigate a programmed return path, ascending to a safe altitude if needed. RTL is a critical safety feature that enables recovery of the aircraft when the pilot loses situational awareness or control. The drone typically lands automatically upon reaching the launch point.}{catD}
\glsadd{Land-Immediately}\dictentry{Land Immediately}{Land Immediately is a failsafe action that commands the drone to descend and land at its current position as quickly as possible. It is triggered in emergency situations where continued flight is unsafe, such as critical battery or system failure. The autopilot reduces throttle and adjusts attitude to achieve a controlled descent. This action prioritizes getting the aircraft on the ground safely over returning to home.}{catD}
\glsadd{Geofence-Action}\dictentry{Geofence Action}{Geofence Action is the autopilot's response when the drone breaches a predefined virtual boundary around the operating area. Common actions include hovering in place, returning to launch, or landing immediately. Geofences prevent the aircraft from flying into restricted airspace or beyond communication range. They are a standard safety feature in both consumer and commercial drone operations.}{catD}
\glsadd{Battery-Failsafe}\dictentry{Battery Failsafe}{Battery Failsafe is an automatic response triggered when the battery voltage or remaining capacity drops below a configured threshold. It typically initiates a return to launch or controlled landing to prevent a power-related crash. Multiple voltage thresholds can be set for warning, low, and critical levels. Battery failsafe is essential for preventing hard landings and potential damage from battery exhaustion.}{catD}
\glsadd{RC-Failsafe}\dictentry{RC Failsafe}{RC Failsafe is the autopilot's response when the radio control link between transmitter and receiver is lost. It can be configured to hold the last known position, return to launch, land, or descend depending on the situation. The failsafe activates after a configurable timeout period without valid RC signals. Proper RC failsafe configuration is crucial for preventing flyaways when the pilot loses control link.}{catD}
\glsadd{GPS-Failsafe}\dictentry{GPS Failsafe}{GPS Failsafe is the autopilot's response when GPS satellite lock is lost during flight. It typically triggers a transition from GPS-dependent modes to a non-GPS mode such as altitude hold or manual control. The failsafe prevents the autopilot from using inaccurate position data for navigation. Pilots should understand the GPS failsafe behavior to maintain safe flight when satellite coverage is poor.}{catD}
\glsadd{Compass-Error}\dictentry{Compass Error}{Compass Error is a fault condition detected when the magnetometer provides inconsistent or unreliable heading data. It can result from electromagnetic interference, uncalibrated sensors, or hardware failure. The autopilot may refuse to arm or switch to a non-compass flight mode to maintain safety. Addressing compass errors promptly is important for maintaining reliable autonomous navigation.}{catD}
\glsadd{Gyro-Calibration}\dictentry{Gyro Calibration}{Gyro Calibration determines the zero-rate offset of the gyroscope when the aircraft is stationary. The autopilot averages multiple readings while the drone is still to compute the bias for each axis. This ensures the gyroscope reads zero rotation when the aircraft is not turning. Accurate gyro calibration is fundamental to stable flight and precise rate control.}{catD}
\glsadd{Accel-Calibration}\dictentry{Accel Calibration}{Accel Calibration determines the offset and scale factors for the accelerometer across all three axes. The procedure requires placing the aircraft in multiple known orientations to measure gravity's component on each axis. This enables the autopilot to correctly determine the aircraft's attitude relative to the Earth. Proper accelerometer calibration is necessary for accurate tilt estimation in stabilized and autonomous flight modes.}{catD}
\glsadd{ADS-B}\dictentry{ADS-B}{Automatic Dependent Surveillance-Broadcast is a surveillance technology in which an aircraft determines its position via GPS and periodically broadcasts it to ground stations and other aircraft. The transmission includes identity, position, altitude, velocity, and intent data. ADS-B enables real-time air traffic awareness without radar and is mandated in many controlled airspaces. For drones, ADS-B receivers allow unmanned systems to detect nearby manned aircraft and maintain safe separation.}{catD}
\glsadd{Beyond-Visual-Line-of-Sight}\dictentry{Beyond Visual Line of Sight}{Beyond Visual Line of Sight, commonly abbreviated BVLOS, refers to drone operations conducted beyond the natural visual range of the remote pilot or visual observer. BVLOS flights require robust detect-and-avoid capabilities and reliable command-and-control datalinks to ensure airspace safety. Regulatory agencies such as the FAA and EASA have developed specific frameworks to authorize BVLOS operations. Enabling BVLOS is considered a key milestone for expanding commercial drone applications such as delivery and infrastructure inspection.}{catD}
\glsadd{Detect-and-Avoid}\dictentry{Detect and Avoid}{Detect and Avoid is a technology and operational concept that enables drones to autonomously sense other aircraft or obstacles and take appropriate action to prevent collisions. DAA systems typically combine radar, cameras, acoustic sensors, and ADS-B receivers to build a picture of surrounding air traffic. The system must detect threats, track their trajectories, and execute avoidance maneuvers within required timeframes. DAA is a critical enabler for BVLOS operations and integration into shared airspace with manned aviation.}{catD}
\glsadd{Electronic-Speed-Controller}\dictentry{Electronic Speed Controller}{An Electronic Speed Controller is a device that regulates the rotational speed of an electric brushless motor by switching power from the battery at a controlled frequency. The ESC receives throttle commands from the flight controller and adjusts the three-phase alternating current delivered to the motor accordingly. Key specifications include continuous current rating, voltage range, and firmware protocol support. Modern ESCs also provide features such as regenerative braking, telemetry feedback, and programmable timing advance.}{catD}
\glsadd{Ground-Control-Station}\dictentry{Ground Control Station}{A Ground Control Station is a hardware and software system used by operators to command, control, and monitor a drone during flight. The GCS provides a user interface displaying real-time telemetry such as altitude, speed, battery status, and GPS position alongside mission planning tools. Communication between the GCS and the aircraft is maintained through a radio datalink operating on designated frequency bands. Advanced GCS platforms support waypoint mission editing, geofence configuration, and autonomous flight execution.}{catD}
\glsadd{On-Screen-Display}\dictentry{On-Screen Display}{An On-Screen Display is an overlay rendered on a video feed that presents real-time flight data to the pilot or operator during drone operations. Typical OSD information includes battery voltage, current draw, GPS coordinates, altitude, airspeed, and flight mode. The OSD is embedded into the analog or digital video signal before transmission to the ground. It provides critical situational awareness that allows the pilot to monitor aircraft status without looking away from the video feed.}{catD}
\glsadd{Received-Signal-Strength-Indicator}\dictentry{Received Signal Strength Indicator}{Received Signal Strength Indicator is a measurement of the power level of a radio signal received by a drone's communication system. RSSI values are typically reported in decibels relative to one milliwatt and indicate the quality of the wireless link between the aircraft and the ground station. Pilots and autonomous systems use RSSI readings to assess link health and trigger failsafe actions such as return-to-home when signal quality degrades. RSSI is influenced by antenna orientation, obstacles, interference, and the distance between transmitter and receiver.}{catD}
\glsadd{Collision-Avoidance-System}\dictentry{Collision Avoidance System}{A Collision Avoidance System is an onboard technology that detects nearby aircraft and obstacles, assesses collision risk, and executes evasive maneuvers to maintain safe separation. These systems use sensors such as radar, ultrasonic detectors, stereo cameras, or LiDAR to perceive the environment. When a potential conflict is identified, the system calculates an avoidance trajectory and commands the flight controller accordingly. Collision avoidance is essential for autonomous drone operations, particularly in populated areas and shared airspace.}{catD}
\glsadd{Datalink}\dictentry{Datalink}{A Datalink is a communication channel that transmits data between a drone and its ground control station or other network nodes. Datalinks carry command-and-control instructions, telemetry, payload data, and video feeds using radio frequencies such as 2.4 GHz or 900 MHz. Link quality depends on transmitter power, antenna gain, frequency selection, and environmental conditions. Redundant and encrypted datalinks are critical for reliable operations, especially for BVLOS missions where maintaining continuous communication is a regulatory requirement.}{catD}
\glsadd{Flight-Envelope-Protection}\dictentry{Flight Envelope Protection}{Flight Envelope Protection is a set of software safeguards built into a drone's flight controller that prevents the aircraft from exceeding safe operating limits. These limits include maximum airspeed, maximum altitude, maximum tilt angle, and maximum climb or descent rate. When the pilot or autonomous system issues a command that would breach these boundaries, the controller constrains the response to remain within the safe envelope. This feature reduces the risk of structural failure, loss of control, and aerodynamic stalls during aggressive maneuvers.}{catD}
\glsadd{Geofencing}\dictentry{Geofencing}{Geofencing is a technology that uses GPS coordinates to define virtual geographic boundaries around a drone's operating area. When the aircraft approaches or crosses a geofence perimeter, the flight controller can trigger automated responses such as altitude limits, hover-in-place, or return-to-home. Geofences are used to enforce regulatory no-fly zones around airports, military installations, and other restricted areas. Manufacturers maintain cloud-based geofence databases that are updated regularly to reflect current airspace restrictions.}{catD}
\glsadd{GPS-Denial}\dictentry{GPS Denial}{GPS Denial refers to conditions under which a drone cannot reliably receive or process signals from Global Positioning System satellites. Denial may result from intentional jamming or spoofing, natural ionospheric disturbances, urban canyon signal blockage, or equipment malfunction. When GPS is unavailable, the drone must rely on alternative navigation methods such as inertial measurement units, visual odometry, or optical flow sensors. GPS-denied environments pose significant challenges for autonomous flight, as position hold, waypoint navigation, and geofencing capabilities are degraded.}{catD}
\glsadd{Gyro-Stabilization}\dictentry{Gyro Stabilization}{Gyro Stabilization is the use of gyroscope sensors within an inertial measurement unit to detect and correct unwanted rotational movements of a drone in flight. The gyroscope measures angular velocity along the pitch, roll, and yaw axes, and the flight controller generates corrective motor commands to maintain attitude stability. Modern multi-axis MEMS gyroscopes provide high sampling rates and low noise, enabling precise stabilization even in turbulent wind conditions. Gyro stabilization is fundamental to stable flight and underpins all autonomous and assisted flight modes.}{catD}
\glsadd{LiDAR}\dictentry{LiDAR}{LiDAR, which stands for Light Detection and Ranging, is a remote sensing technology that uses pulsed laser light to measure distances to objects and create detailed three-dimensional maps of the environment. On drones, LiDAR systems emit rapid laser pulses and measure the time of flight for each pulse to return after reflecting off surfaces. The resulting point cloud data can be processed to generate elevation models, detect obstacles, and enable terrain-following flight. LiDAR-equipped drones are widely used in surveying, forestry, archaeology, and autonomous navigation.}{catD}
\glsadd{Obstacle-Avoidance}\dictentry{Obstacle Avoidance}{Obstacle Avoidance is the capability of a drone to detect physical objects in its flight path and maneuver around them without human intervention. Sensor modalities used for obstacle avoidance include ultrasonic rangefinders, infrared sensors, stereo vision cameras, time-of-flight sensors, and LiDAR. The flight controller processes sensor data to compute safe trajectories and adjusts motor outputs to steer the aircraft clear of obstacles. This capability is essential for low-altitude operations in cluttered environments such as warehouses, forests, and urban areas.}{catD}
\glsadd{Payload-Release}\dictentry{Payload Release}{Payload Release is the mechanism or procedure by which a drone delivers or jettisons a carried payload at a designated location during flight. Release systems may use servo-actuated clamps, electromagnetic latches, or pyrotechnic cutters depending on the weight and type of payload. Precise GPS waypoint targeting and descent management ensure the payload is delivered to the intended area. Common applications include package delivery, agricultural seed dispersal, life raft deployment, and scientific instrument deployment.}{catD}
\glsadd{Sense-and-Avoid}\dictentry{Sense and Avoid}{Sense and Avoid is an integrated system that combines environmental sensing with automated decision-making to prevent mid-air collisions and ground impacts during drone operations. The sensing component uses cameras, radar, ADS-B receivers, or acoustic arrays to detect nearby aircraft and obstacles. The avoid component evaluates threat trajectories and executes appropriate evasion maneuvers or speed adjustments. Sense and Avoid is a regulatory prerequisite for integrating unmanned aircraft into non-segregated airspace.}{catD}
\glsadd{Unmanned-Traffic-Management}\dictentry{Unmanned Traffic Management}{Unmanned Traffic Management is a system of services and infrastructure designed to safely manage drone traffic in low-altitude airspace, analogous to air traffic control for manned aircraft. UTM provides capabilities for flight authorization, real-time traffic monitoring, geofence enforcement, and conflict detection between drones. Key stakeholders include airspace regulators, drone operators, technology providers, and communication service providers. UTM frameworks are being developed worldwide to enable scalable and safe drone operations in urban and rural environments.}{catD}
\glsadd{Visual-Line-of-Sight}\dictentry{Visual Line of Sight}{Visual Line of Sight refers to the requirement that a drone operator or designated visual observer maintain unaided visual contact with the aircraft throughout its flight. VLOS operations ensure the pilot can perceive the drone's attitude, altitude, and surrounding airspace to avoid hazards. Most current regulations restrict drone operations to VLOS, typically limiting flight distances to a few hundred meters. VLOS is the baseline operational category for recreational and commercial drone pilots and serves as the foundation for higher-risk operational approvals.}{catD}
\glsadd{Circle-Mode}\dictentry{Circle Mode}{Circle Mode is an autonomous or semi-autonomous flight mode in which a drone orbits a specified point of interest at a set radius, altitude, and speed. The flight controller continuously adjusts heading and bank angle to maintain a consistent circular trajectory around the target. This mode is commonly used for aerial photography and surveillance, allowing the camera to remain pointed at the center of the orbit. Advanced implementations allow the pilot to dynamically adjust radius, speed, and camera gimbal angle during the orbit.}{catD}
\glsadd{Follow-Me-Mode}\dictentry{Follow Me Mode}{Follow Me Mode is a flight mode in which a drone autonomously tracks and follows a moving target, typically the operator or a GPS-equipped device. The system uses GPS coordinates from the target device, sometimes supplemented by visual tracking algorithms, to maintain a preset distance and relative position. The drone adjusts its heading, altitude, and speed to match the target's movement in real time. Follow Me Mode is popular for action sports filming, outdoor recreation, and solo field operations where a second operator is unavailable.}{catD}
\glsadd{Land-Mode}\dictentry{Land Mode}{Land Mode is an automated flight mode that initiates a controlled vertical descent and touchdown at the drone's current horizontal position. The flight controller reduces thrust progressively while maintaining attitude stability until the aircraft detects ground contact or reaches its landing gear. Some implementations use barometric altitude, ultrasonic rangefinders, or LiDAR to gauge descent rate and trigger motor shutoff at touchdown. Land Mode provides a safe and predictable landing procedure and can be triggered manually or as part of an automated mission sequence.}{catD}
\glsadd{Takeoff-Mode}\dictentry{Takeoff Mode}{Takeoff Mode is an automated flight sequence that raises a drone from the ground to a predetermined hover altitude with minimal pilot input. The flight controller increases motor thrust to overcome gravity, stabilizes the aircraft using gyro and accelerometer data, and levels off at the target altitude. Once at hover altitude, the drone typically holds position using GPS and barometric sensors. Automated takeoff reduces pilot workload and ensures a consistent, safe departure procedure for both manual and autonomous missions.}{catD}
\glsadd{Data-Logger}\dictentry{Data Logger}{A Data Logger is an onboard device or software module that records flight data from a drone's sensors and systems for later analysis. Logged parameters typically include GPS coordinates, altitude, attitude, motor RPM, battery voltage, current draw, and control inputs. Data loggers may store information on onboard flash memory, SD cards, or transmit it to ground-based storage via telemetry. Reviewing logged data is essential for post-flight diagnostics, performance optimization, incident investigation, and regulatory compliance.}{catD}
\glsadd{Black-Box}\dictentry{Black Box}{A Black Box on a drone is a ruggedized onboard recording device that captures critical flight data and, in some cases, video footage for post-incident analysis. Unlike a standard data logger, a black box is typically designed to survive crashes, fires, and water immersion through protective enclosures. Recorded parameters include flight controls, navigation data, motor status, and system health indicators. Black boxes are increasingly required for commercial and regulatory operations to support accident reconstruction and accountability.}{catD}
\glsadd{Flight-Log}\dictentry{Flight Log}{A Flight Log is a comprehensive record of all events, parameters, and actions that occurred during a single drone flight session. It includes takeoff time, landing time, total flight duration, GPS track history, battery usage, any triggered failsafes, and pilot commands. Flight logs are generated automatically by the flight controller and stored onboard or transmitted to a ground station. They serve as essential documentation for maintenance scheduling, mission verification, accident investigation, and regulatory record-keeping requirements.}{catD}
\glsadd{Telemetry-Link}\dictentry{Telemetry Link}{A Telemetry Link is a one-way or bidirectional radio connection that transmits real-time flight data from a drone to a ground control station or other monitoring system. Telemetry data includes aircraft position, attitude, battery health, motor status, and sensor readings. The link operates on designated radio frequencies and uses protocols optimized for low latency and reliable delivery. High-quality telemetry links are critical for pilot situational awareness, autonomous operations, and remote fleet management in commercial drone programs.}{catD}
\glsadd{Channel-Mapping}\dictentry{Channel Mapping}{Channel Mapping is the process of assigning specific control functions to individual radio transmitter channels so the drone's flight controller interprets commands correctly. Standard channel assignments include throttle, aileron, elevator, and rudder, with additional channels for flight mode selection, camera gimbal control, and auxiliary functions. The mapping is configured in the transmitter and matched in the flight controller software. Incorrect channel mapping can result in reversed or unresponsive controls, making proper configuration essential for safe flight operations.}{catD}
\glsadd{Receiver-Binding}\dictentry{Receiver Binding}{Receiver Binding is the process of pairing a drone's radio receiver with a specific transmitter so they communicate exclusively with each other. Binding establishes a unique identification link using protocols specific to the radio manufacturer, such as FrSky, ELRS, or IBUS. The process typically involves placing both the receiver and transmitter into bind mode within a defined time window. A successful bind ensures reliable command reception and prevents interference from other transmitters operating on the same frequency.}{catD}
\glsadd{Lost-Link}\dictentry{Lost Link}{Lost Link refers to the condition where a drone loses its command-and-control radio connection with the ground control station during flight. Causes include excessive distance, antenna misalignment, signal interference, or physical obstruction between the transmitter and receiver. When a lost link event is detected, the flight controller activates a preconfigured failsafe procedure such as return-to-home, hover in place, or land immediately. Reliable lost link recovery procedures are critical for preventing flyaways and ensuring safe drone operations.}{catD}
\glsadd{Low-Voltage-Cutoff}\dictentry{Low Voltage Cutoff}{Low Voltage Cutoff is a protective feature that automatically reduces or shuts down motor output when the battery voltage drops below a safe threshold. The cutoff prevents deep discharge of lithium polymer batteries, which can cause permanent cell damage, swelling, or fire hazards. The threshold is configurable in the flight controller and is typically set to allow sufficient reserve for a safe return-to-home maneuver. Proper LVC configuration balances maximum flight time with battery longevity and flight safety.}{catD}
\glsadd{Battery-Monitoring}\dictentry{Battery Monitoring}{Battery Monitoring is the continuous measurement and reporting of a drone's battery health parameters during flight, including voltage, current draw, remaining capacity, and cell balance. Monitoring systems use current sensors and voltage dividers to calculate real-time energy consumption and estimated remaining flight time. Battery data is displayed on the ground station or OSD and can trigger failsafe actions when critical thresholds are reached. Effective battery monitoring prevents in-flight power loss, extends battery lifespan, and supports safe mission planning.}{catD}
\glsadd{Arm-Switch}\dictentry{Arm Switch}{An Arm Switch is a physical or software-based control that enables or disables motor output on a drone. When disarmed, the flight controller prevents the motors from spinning regardless of throttle input, providing a critical safety mechanism during handling and setup. Arming typically requires a specific stick combination or switch position, and some systems enforce additional pre-arm checks such as GPS lock and sensor calibration. The arm switch concept ensures that a drone cannot accidentally start its motors, reducing the risk of injury during preflight preparation.}{catD}
\glsadd{Motor-Protocol}\dictentry{Motor Protocol}{A Motor Protocol is the communication standard used between a flight controller and electronic speed controllers to command motor speeds. Common protocols include PWM, OneShot, DShot, and Bidirectional DShot, each differing in signal resolution, speed, and digital vs analog transmission. Higher-performance protocols offer faster update rates and higher throttle resolution, resulting in smoother and more responsive flight characteristics. Protocol selection affects flight controller compatibility, wiring requirements, and overall system performance.}{catD}
\glsadd{IBUS}\dictentry{IBUS}{IBUS is a digital serial communication protocol developed by FlySky for transmitting control signals and telemetry data between a radio receiver and flight controller. It operates at a fixed baud rate and supports multiple channels of control input plus bidirectional telemetry for battery voltage, current, and signal quality. IBUS offers lower latency and simpler wiring compared to traditional PWM connections, as a single cable carries all channel data. The protocol is widely used in budget and intermediate drone builds due to its simplicity and compatibility with FlySky transmitter systems.}{catD}
\glsadd{ELRS}\dictentry{ELRS}{ELRS, or ExpressLRS, is an open-source radio control link designed for high-performance drone operations with low latency and long range. It operates on commonly available ISM band frequencies such as 900 MHz and 2.4 GHz, using advanced modulation techniques like LoRa and FLRC. ELRS supports configurable packet rates up to 500 Hz, providing extremely responsive control input for racing and freestyle applications. The open-source nature of ELRS allows community-driven firmware development, wide hardware compatibility, and continuous performance improvements.}{catD}
\glsadd{FrSky}\dictentry{FrSky}{FrSky is a manufacturer of radio control systems and protocols widely used in the drone and RC aviation community. Their systems include transmitters, receivers, and telemetry modules that communicate using protocols such as ACCST and ACCESS. FrSky receivers offer features including channel mapping, telemetry feedback, and failsafe configuration at competitive price points. The ecosystem supports long-range links, multi-protocol compatibility, and integration with popular flight controller firmware such as Betaflight and INAV.}{catD}
\glsadd{Mission-Planner-GCS}\dictentry{Mission Planner GCS}{Mission Planner is an open-source ground control station software designed for autonomous drone operations, primarily used with the ArduPilot firmware ecosystem. It supports full mission planning with waypoint editing, survey grid generation, and terrain-following flight path design. The software provides real-time telemetry display, flight data logging, and post-flight analysis tools including 3D track visualization. Mission Planner is a widely adopted tool in the commercial and research drone communities due to its extensive feature set and active development community.}{catD}
\glsadd{Video-Transmitter}\dictentry{Video Transmitter}{A Video Transmitter, commonly abbreviated VTX, is an onboard device that wirelessly transmits live video from a drone's camera to a ground-based receiver or goggles. VTX units operate on various frequency bands including 5.8 GHz, 2.4 GHz, and 1.3 GHz, with adjustable power levels to balance range and regulatory compliance. Video quality is affected by antenna selection, transmitter power, channel interference, and environmental conditions such as obstacles and weather. The video transmitter is essential for first-person view piloting, providing the pilot with real-time visual feedback from the aircraft's perspective.}{catD}
\end{multicols}

\clearpage
\chapter*{Space Science}
\addcontentsline{toc}{chapter}{Space Science}
\markboth{Space Science}{Space Science}
\clearpage
\begin{multicols}{2}
\small
Orbital mechanics, spacecraft engineering, and the space environment. From LEO to Lagrange points, re-entry to radiation belts.\par
\vspace{0.5cm}
\glsadd{Space-Science}\dictentry{Space Science}{Space Science is the study of celestial bodies, the space environment, and the universe. It draws on physics, chemistry, biology, and geology to understand phenomena from planetary surfaces to distant galaxies. Research uses ground-based observatories, space telescopes, and robotic missions. Advances in this field deepen our understanding of how the cosmos formed and operates.}{catSS}
\glsadd{Astronomy}\dictentry{Astronomy}{Astronomy is the scientific study of celestial objects such as stars, planets, and galaxies, plus phenomena outside Earth's atmosphere. It relies on observing electromagnetic radiation collected by telescopes and space probes. Modern astronomy includes observational, theoretical, and radio sub-disciplines. It has revealed the scale of the universe, stellar life cycles, and thousands of exoplanets.}{catSS}
\glsadd{Astrophysics}\dictentry{Astrophysics}{Astrophysics applies physics principles to understand the nature and behavior of celestial objects. It studies physical properties of stars including luminosity, temperature, and energy generation through nuclear fusion. Astrophysicists model black holes, neutron stars, and the large-scale structure of the universe. It combines theoretical modeling with observational data from telescopes across the electromagnetic spectrum.}{catSS}
\glsadd{Cosmology}\dictentry{Cosmology}{Cosmology studies the origin, evolution, and fate of the universe as a whole. It investigates the Big Bang, cosmic expansion, dark matter, and dark energy through observation and theory. Cosmologists use data from the cosmic microwave background and galaxy surveys to test models. This field addresses questions about the age, geometry, and composition of the universe.}{catSS}
\glsadd{Planetary-Science}\dictentry{Planetary Science}{Planetary Science is the interdisciplinary study of planets, moons, and planetary systems including their formation and evolution. It encompasses geology, atmospheric science, and astrobiology applied to worlds beyond Earth. Researchers use spacecraft missions, telescopes, and laboratory analysis of meteorites. It informs our understanding of how habitable environments arise elsewhere in the universe.}{catSS}
\glsadd{Exoplanet}\dictentry{Exoplanet}{An Exoplanet is a planet orbiting a star outside our solar system, also known as an extrasolar planet. Thousands have been discovered using transit and radial velocity methods since the first confirmed detection in 1992. Exoplanets range from rocky Earth-sized worlds to massive gas giants. Studying their atmospheres helps scientists assess potential habitability and planetary diversity.}{catSS}
\glsadd{Solar-System}\dictentry{Solar System}{The Solar System is the gravitationally bound system of the Sun and all objects orbiting it, including eight planets, moons, asteroids, and comets. It formed roughly 4.6 billion years ago from the collapse of a molecular cloud. Planets divide into inner rocky terrestrials and outer gas and ice giants. Studying it provides a baseline for understanding other planetary systems.}{catSS}
\glsadd{Star}\dictentry{Star}{A Star is a massive luminous sphere of plasma held together by gravity, powered by nuclear fusion of hydrogen into helium. Stars vary enormously in size, temperature, and lifespan, from small red dwarfs to supergiants. A star's mass determines its lifecycle from birth in a nebula to death as a white dwarf, neutron star, or black hole. Stars are the fundamental building blocks of galaxies.}{catSS}
\glsadd{Planet}\dictentry{Planet}{A Planet is a large celestial body orbiting a star, massive enough for gravity to shape it into a rough sphere, and has cleared its orbital neighborhood. The eight Solar System planets range from small rocky Mercury to massive gas giant Jupiter. Planets may have atmospheres, moons, rings, and geological activity. Studying planets across star systems reveals wide diversity in planetary formation.}{catSS}
\glsadd{Moon}\dictentry{Moon}{A Moon is a natural satellite orbiting a planet or dwarf planet, held in orbit by gravitational attraction. The Solar System contains hundreds of known moons, from tiny irregular bodies to Titan, which is larger than Mercury. Moons can influence their parents through tidal forces that drive geological activity. Some moons, such as Europa and Enceladus, are candidates for extraterrestrial life searches.}{catSS}
\glsadd{Asteroid}\dictentry{Asteroid}{An Asteroid is a small rocky body orbiting the Sun, primarily found between Mars and Jupiter. They are remnants from the Solar System's formation about 4.6 billion years ago. Sizes range from pebbles to Vesta, over 500 kilometers across. Some cross Earth's orbit as near-Earth objects, making them subjects of planetary defense research.}{catSS}
\glsadd{Comet}\dictentry{Comet}{A Comet is an icy body that releases gas and dust near the Sun, forming a visible coma and tails. Comets originate from the Kuiper Belt and Oort Cloud where they remain frozen until perturbed inward. Solar heating sublimes their volatile ices, releasing primitive material from the early Solar System. They help scientists understand the composition of the early solar nebula.}{catSS}
\glsadd{Meteor}\dictentry{Meteor}{A Meteor is the streak of light produced when a meteoroid enters Earth's atmosphere and heats surrounding air through friction. Meteors are commonly called shooting stars and occur at altitudes of 80 to 120 kilometers. Meteor showers happen when Earth passes through comet debris trails. Observing meteors helps scientists map the distribution of small particles in the inner Solar System.}{catSS}
\glsadd{Meteorite}\dictentry{Meteorite}{A Meteorite is a fragment of an asteroid or other body that survives atmospheric entry and reaches Earth's surface. They classify into stony, iron, and stony-iron types reflecting different origins. Meteorites provide invaluable data about the age and composition of the Solar System. Some contain presolar grains predating the Sun itself. They are studied to understand planetary building blocks.}{catSS}
\glsadd{Asteroid-Belt}\dictentry{Asteroid Belt}{The Asteroid Belt is a torus-shaped region between Mars and Jupiter containing millions of rocky bodies. Their combined mass is less than that of Earth's Moon. Jupiter's strong gravity prevented this material from coalescing into a planet. Key objects include the dwarf planet Ceres and the large asteroid Vesta. The belt is a source of near-Earth asteroids and a target for exploration.}{catSS}
\glsadd{Kuiper-Belt}\dictentry{Kuiper Belt}{The Kuiper Belt is a vast region of icy bodies extending beyond Neptune to about 50 astronomical units from the Sun. It contains dwarf planets like Pluto, Eris, Haumea, and Makemake, plus numerous small icy objects. The belt is a source of short-period comets deflected inward toward the Sun. Studying its objects reveals early Solar System conditions and planetary migration processes.}{catSS}
\glsadd{Oort-Cloud}\dictentry{Oort Cloud}{The Oort Cloud is a hypothetical spherical shell of icy bodies surrounding the Solar System at roughly 2,000 to 100,000 astronomical units. It is believed to be the source of long-period comets on highly elliptical orbits. The cloud has never been directly observed and is inferred from incoming comets' orbital characteristics. Its total mass is estimated at several Earth masses of icy material.}{catSS}
\glsadd{Dwarf-Planet}\dictentry{Dwarf Planet}{A Dwarf Planet orbits the Sun with enough mass for a nearly round shape, but has not cleared its orbital neighborhood. This distinguishes it from full planets under the 2006 IAU definition. Known examples include Pluto, Eris, Ceres, Haumea, and Makemake. Many have moons and complex surface features. They represent an important class of bodies for understanding outer Solar System diversity.}{catSS}
\glsadd{Pluto}\dictentry{Pluto}{Pluto is a dwarf planet in the Kuiper Belt, discovered in 1930 by Clyde Tombaugh and reclassified in 2006. It has a thin nitrogen-methane atmosphere and a surface of nitrogen, methane, and water ice. Pluto has five known moons, with Charon being nearly half its diameter. The New Horizons flyby in 2015 revealed nitrogen glaciers and possible cryovolcanism. Its orbit is highly eccentric and inclined.}{catSS}
\glsadd{Ceres}\dictentry{Ceres}{Ceres is the largest object in the asteroid belt and the only dwarf planet in the inner Solar System, with a diameter of roughly 940 kilometers. NASA's Dawn spacecraft orbited it from 2015 to 2018, revealing bright salt deposits, cryovolcanism evidence, and a subsurface ocean. Its low density suggests a rock-and-ice composition. Ceres provides insights into planetary differentiation and early Solar System conditions.}{catSS}
\glsadd{Sun}\dictentry{Sun}{The Sun is a G-type main-sequence star at the Solar System's center, containing 99.86 percent of the system's mass. It generates energy through hydrogen-to-helium fusion in its core, producing temperatures around 15 million degrees Celsius. Its magnetic activity drives sunspots, solar flares, and coronal mass ejections. The Sun is about 4.6 billion years old and will remain stable for another 5 billion years.}{catSS}
\glsadd{Solar-Flare}\dictentry{Solar Flare}{A Solar Flare is a sudden intense burst of radiation and energy from the Sun's surface, caused by magnetic energy release. Flares emit radiation across the electromagnetic spectrum from radio waves to gamma rays and can last minutes to hours. They can disrupt radio communications, GPS, and satellite operations on Earth. Solar flares are often linked to coronal mass ejections and are key space weather phenomena.}{catSS}
\glsadd{Coronal-Mass-Ejection}\dictentry{Coronal Mass Ejection}{A Coronal Mass Ejection is a massive expulsion of magnetized plasma from the Sun's corona into interplanetary space. CMEs carry billions of tons of solar material at speeds up to 3,000 kilometers per second. Directed at Earth, they can trigger geomagnetic storms affecting power grids and satellites. CMEs are often associated with solar flares but can occur independently. Predicting them is essential for protecting infrastructure.}{catSS}
\glsadd{Heliosphere}\dictentry{Heliosphere}{The Heliosphere is the vast bubble-like region dominated by the solar wind and the Sun's magnetic field. It extends beyond the outer planets with the heliopause boundary at roughly 120 astronomical units. The heliosphere shields the inner Solar System from galactic cosmic radiation. Voyager 1 and 2 have crossed the heliopause, providing direct measurements. Its size and shape depend on solar wind pressure and the interstellar medium.}{catSS}
\glsadd{Magnetosphere}\dictentry{Magnetosphere}{A Magnetosphere is the region around a planet where its magnetic field dominates over the solar wind. It deflects most incoming particles and traps charged particles in radiation belts. Earth's extends tens of thousands of kilometers, compressed sunward and elongated into a magnetotail. It protects atmospheres from being stripped away. Understanding magnetospheric dynamics is critical for satellite and astronaut safety.}{catSS}
\glsadd{Van-Allen-Belt}\dictentry{Van Allen Belt}{A Van Allen Belt is a zone of energetic charged particles trapped within a planet's magnetic field. Earth has two primary belts: an inner belt at 1,000 to 6,000 kilometers and an outer belt at 13,000 to 60,000 kilometers altitude. Discovered by James Van Allen in 1958, they pose radiation hazards to spacecraft and astronauts. Their intensity varies with solar activity and geomagnetic storms.}{catSS}
\glsadd{Aurora}\dictentry{Aurora}{An Aurora is a natural light display caused by solar wind particles interacting with atoms in a planet's upper atmosphere. Guided by magnetic field lines toward the poles, these particles excite atmospheric gases, causing them to emit light. Oxygen produces green and red hues, while nitrogen creates blue and purple. Auroras occur on other magnetized planets including Jupiter and Saturn. They are visible indicators of space weather activity.}{catSS}
\glsadd{Northern-Lights}\dictentry{Northern Lights}{The Northern Lights, also known as Aurora Borealis, are auroral displays visible in the Northern Hemisphere's high-latitude regions. They result from solar wind particles interacting with Earth's magnetic field and upper atmosphere. The lights appear as shimmering curtains, arcs, or rays of green, red, and purple. They are most active during high solar activity and best observed in Norway, Iceland, and northern Canada.}{catSS}
\glsadd{Southern-Lights}\dictentry{Southern Lights}{The Southern Lights, also known as Aurora Australis, are auroral displays in the Southern Hemisphere's high-latitude regions. They form through the same mechanism as the Northern Lights, with charged particles interacting with the atmosphere along magnetic field lines. Displays are visible from Antarctica, southern New Zealand, and Tasmania. They mirror the Northern Hemisphere's auroral oval and occur simultaneously.}{catSS}
\glsadd{Ionosphere}\dictentry{Ionosphere}{The Ionosphere is a region of Earth's upper atmosphere from about 60 to 1,000 kilometers where solar radiation ionizes atmospheric gases. The resulting free electrons and ions affect radio wave propagation, enabling long-distance communication. It varies with time of day, solar activity, and geomagnetic conditions. The ionosphere is critical for satellite navigation accuracy and high-frequency radio communication.}{catSS}
\glsadd{Atmosphere}\dictentry{Atmosphere}{The Atmosphere is the layer of gases gravitationally bound to a planet, providing radiation protection and regulating temperature. Earth's atmosphere is mainly nitrogen and oxygen with trace gases. It divides into layers based on temperature gradients: troposphere, stratosphere, mesosphere, thermosphere, and exosphere. The atmosphere sustains life by maintaining pressure, distributing heat, and filtering harmful solar radiation.}{catSS}
\glsadd{Exosphere}\dictentry{Exosphere}{The Exosphere is the outermost atmospheric layer extending from the thermosphere's top to where particles escape into space. On Earth it begins at roughly 500 kilometers and merges with the interplanetary medium. It is composed mainly of hydrogen and helium atoms at very low densities. Many high-orbit satellites operate within the exosphere. This region is important for understanding atmospheric escape and long-term planetary evolution.}{catSS}
\glsadd{Thermosphere}\dictentry{Thermosphere}{The Thermosphere extends from about 80 to 600 kilometers above Earth, above the mesosphere. Temperatures can exceed 2,000 degrees Celsius due to extreme ultraviolet absorption, though the thin air would feel cold. The International Space Station and many low-Earth-orbit satellites operate here. Auroral displays also occur in this region as charged particles collide with atmospheric gases.}{catSS}
\glsadd{Mesosphere}\dictentry{Mesosphere}{The Mesosphere lies between the stratosphere and thermosphere, extending from about 50 to 85 kilometers. It is the coldest atmospheric layer, with temperatures dropping to minus 90 degrees Celsius at the mesopause. Most meteors burn up here, producing visible shooting stars. Noctilucent clouds, the highest clouds in the atmosphere, form near the mesopause. This layer is difficult to study directly with balloons or satellites.}{catSS}
\glsadd{Stratosphere}\dictentry{Stratosphere}{The Stratosphere extends from about 12 to 50 kilometers above Earth's surface, above the troposphere. It contains the ozone layer, which absorbs ultraviolet radiation and protects life on Earth. Temperature increases with altitude here due to ozone absorbing solar energy. Commercial aircraft sometimes fly in the lower stratosphere to avoid turbulent weather. The stratosphere plays a key role in Earth's energy balance and climate.}{catSS}
\glsadd{Troposphere}\dictentry{Troposphere}{The Troposphere is the lowest atmospheric layer, extending from the surface to about 12 kilometers. It contains approximately 75 percent of the atmosphere's mass and nearly all its water vapor. Nearly all weather phenomena occur within this layer. Temperature decreases with altitude, driving convection and atmospheric circulation. This is where human activity takes place and where greenhouse gases most directly affect climate.}{catSS}
\glsadd{Ozone-Layer}\dictentry{Ozone Layer}{The Ozone Layer is a stratospheric region containing concentrated ozone molecules that absorb harmful ultraviolet-B and ultraviolet-C radiation. This absorption protects surface life from DNA damage and skin cancer. The layer is most concentrated at 15 to 35 kilometers altitude and thinnest over the poles. Human-produced CFCs caused significant depletion, leading to the Montreal Protocol. Recovery is expected to continue through coming decades.}{catSS}
\glsadd{Greenhouse-Effect}\dictentry{Greenhouse Effect}{The Greenhouse Effect is the process by which atmospheric gases trap outgoing infrared radiation, warming the surface and lower atmosphere. On Earth, key greenhouse gases include water vapor, carbon dioxide, methane, and nitrous oxide. This natural process maintains Earth's average temperature at about 15 degrees Celsius. Human activities have increased these concentrations, enhancing the effect and driving global warming.}{catSS}
\glsadd{Carbon-Dioxide}\dictentry{Carbon Dioxide}{Carbon Dioxide is a colorless, odorless greenhouse gas significant in planetary atmospheres. On Earth, it is produced by respiration, volcanic eruptions, fossil fuel combustion, and deforestation. Plants absorb it during photosynthesis, forming a key part of the carbon cycle. Rising levels are the primary driver of contemporary climate change. On Venus, its dense atmosphere creates a runaway greenhouse effect with extreme surface temperatures.}{catSS}
\glsadd{Orbital-Mechanics}\dictentry{Orbital Mechanics}{Orbital Mechanics is the study of object motion in space under gravity and other forces. It applies Newton's laws to predict and plan trajectories for spacecraft, satellites, and celestial bodies. Key concepts include orbital elements, conic section trajectories, and delta-v maneuver planning. It is essential for mission design, station-keeping, and collision avoidance. It underpins all spaceflight operations.}{catSS}
\glsadd{Celestial-Mechanics}\dictentry{Celestial Mechanics}{Celestial Mechanics is the branch of astronomy dealing with motions and gravitational interactions of celestial bodies. It extends orbital mechanics to include perturbations from multiple bodies, tidal forces, and long-term orbital evolution. Its mathematical foundations were developed by Lagrange, Laplace, and others. It explains phenomena like orbital resonances, precession, and planetary system stability.}{catSS}
\glsadd{Keplers-Laws}\dictentry{Kepler's Laws}{Kepler's Laws are three empirical rules describing planetary motion around the Sun. Formulated by Johannes Kepler in the early 17th century from Tycho Brahe's data, they describe elliptical orbits, equal area sweeping, and the period-distance relationship. These laws were later explained by Newton's gravitation. They form the foundation of orbital mechanics for any two-body gravitational system.}{catSS}
\glsadd{Keplers-First-Law}\dictentry{Kepler's First Law}{Kepler's First Law states that a planet's orbit around the Sun is an ellipse with the Sun at one focus. This replaced ancient circular orbit models and explained observed speed variations. The ellipse's shape is described by eccentricity, ranging from zero for circles to near one for elongated orbits. This law applies to any two-body gravitational system and fundamentally changed planetary motion understanding.}{catSS}
\glsadd{Keplers-Second-Law}\dictentry{Kepler's Second Law}{Kepler's Second Law states that a line from a planet to the Sun sweeps equal areas in equal time intervals. This means planets move faster near the Sun at perihelion and slower at aphelion. The law follows from angular momentum conservation in a gravitational field. It explains observed orbital speed variations without external forces. It is essential for calculating spacecraft velocities at different orbital points.}{catSS}
\glsadd{Keplers-Third-Law}\dictentry{Kepler's Third Law}{Kepler's Third Law states that the square of a planet's orbital period is proportional to the cube of its semi-major axis. More distant planets take longer to orbit the Sun. Newton generalized this law to include the orbiting bodies' masses. It allows astronomers to determine a planet's distance from its orbital period alone. This law is widely used to find exoplanet and binary star system masses.}{catSS}
\glsadd{Newtons-Law-of-Gravitation}\dictentry{Newton's Law of Gravitation}{Newton's Law of Universal Gravitation states that every mass attracts every other mass with force proportional to their product and inversely proportional to the distance squared. Published in 1687, it unified terrestrial and celestial mechanics. The law explains Kepler's planetary motion laws and predicts tidal forces. It remained accepted until Einstein's general relativity and is still used for most spaceflight calculations.}{catSS}
\glsadd{Gravitational-Constant}\dictentry{Gravitational Constant}{The Gravitational Constant, denoted G, is the proportionality constant in Newton's law of universal gravitation. Its value is approximately 6.674 times 10 to the minus 11 newton square meters per square kilogram. It determines the strength of gravitational attraction between any two masses. First measured accurately by Henry Cavendish in 1798, its precise value remains one of physics' less well-determined fundamental constants.}{catSS}
\glsadd{Gravity}\dictentry{Gravity}{Gravity is the fundamental attractive force between any objects with mass or energy. It is the weakest of the four fundamental forces but acts over infinite distances and is always attractive. Einstein's general relativity describes gravity as spacetime curvature caused by mass and energy. Gravity governs planetary, stellar, and galactic motions and gives objects weight on Earth. It also drives tides and atmospheric convection.}{catSS}
\glsadd{Weightlessness}\dictentry{Weightlessness}{Weightlessness is the condition of having no apparent weight, occurring when no contact forces act on a body. In orbit, astronauts experience it because both they and their spacecraft are in free fall together. True weightlessness occurs only in free fall or at gravitational null points. It affects fluid distribution, bone density, and muscle mass. Understanding it is essential for designing spacecraft and medical countermeasures.}{catSS}
\glsadd{Free-Fall}\dictentry{Free Fall}{Free Fall is motion under the sole influence of gravity with no other forces acting. Near Earth's surface, an object accelerates at 9.81 meters per second squared. In space, free fall produces the weightlessness astronauts experience in orbit. Satellites are perpetually falling toward Earth while moving forward fast enough to miss it. Free fall is fundamental to orbital mechanics, spacecraft design, and gravitational physics studies.}{catSS}
\glsadd{Circular-Velocity}\dictentry{Circular Velocity}{Circular Velocity is the specific speed for maintaining a perfectly circular orbit at a given distance from a central body. It equals the orbital velocity when eccentricity is zero. On Earth, circular velocity at the surface is roughly 7.9 kilometers per second ignoring drag. Higher altitudes require less circular velocity as gravitational force weakens. Circular orbits are common for communication and Earth-observation satellites.}{catSS}
\glsadd{Hohmann-Transfer-Orbit}\dictentry{Hohmann Transfer Orbit}{A Hohmann Transfer Orbit is the most fuel-efficient two-burn trajectory between two circular, coplanar orbits. It consists of an elliptical orbit tangent to both orbits, with burns at each end. The first burn enters the transfer ellipse, and the second circularizes at the destination. First proposed by Walter Hohmann in 1925, it is widely used for interplanetary missions. It provides the minimum energy path between two orbits.}{catSS}
\glsadd{Bi-Elliptic-Transfer}\dictentry{Bi-Elliptic Transfer}{A Bi-Elliptic Transfer uses two intermediate elliptical arcs with three burns to transfer between circular orbits. It becomes more fuel-efficient than a Hohmann transfer when the final-to-initial orbit ratio exceeds about 11.94. The spacecraft raises its apoapsis beyond the target, adjusts the periapsis, then circularizes. Though more propellant-efficient for large orbit changes, it requires significantly longer transfer times.}{catSS}
\glsadd{Gravity-Assist}\dictentry{Gravity Assist}{A Gravity Assist uses a planet's gravitational field and orbital motion to change a spacecraft's speed and trajectory. The spacecraft exchanges momentum with the planet, gaining or losing kinetic energy relative to the Sun. This technique reaches distant targets with much less fuel than direct trajectories. It has been used in Voyager, Cassini, and Juno missions. The maneuver requires precise timing and navigation.}{catSS}
\glsadd{Swing-By}\dictentry{Swing-By}{A Swing-By is an orbital maneuver where a spacecraft passes close to a body to alter its trajectory through gravitational interaction. It is synonymous with gravity assist and relies on momentum exchange between spacecraft and planet. The spacecraft follows a hyperbolic trajectory, gaining speed relative to the Sun. Swing-bys enabled missions to explore the outer Solar System in single missions. Careful approach geometry planning is essential.}{catSS}
\glsadd{Slingshot}\dictentry{Slingshot}{A Slingshot is an informal term for a gravity assist, describing how a spacecraft is accelerated or redirected by a planet's gravitational field. In reality, the spacecraft follows a hyperbolic trajectory relative to the planet, gaining speed relative to the Sun. The slingshot effect enabled outer Solar System exploration impossible with available propulsion alone. It is a foundational technique in interplanetary mission design.}{catSS}
\glsadd{Lagrange-Point}\dictentry{Lagrange Point}{A Lagrange Point is a position where gravitational forces of two large bodies balance with orbital centrifugal force. Five such points exist in any two-body system, labeled L1 through L5. Spacecraft at Lagrange points require minimal station-keeping fuel, making them ideal for observatories. James Webb Space Telescope operates at Sun-Earth L2, while SOHO is at L1. They were derived mathematically by Joseph-Louis Lagrange in 1772.}{catSS}
\glsadd{L1}\dictentry{L1}{L1 is the first Lagrange point, about 1.5 million kilometers from Earth toward the Sun, on the Sun-Earth line. At this point, combined gravitational pull of Sun and Earth with centrifugal force maintains a fixed spacecraft position. L1 is ideal for solar observatories providing uninterrupted Sun views. The Advanced Composition Explorer and Deep Space Climate Observatory operate there for space weather monitoring.}{catSS}
\glsadd{L2}\dictentry{L2}{L2 is the second Lagrange point, about 1.5 million kilometers from Earth on the Sun's opposite side. A spacecraft here orbits the Sun with the same period as Earth. L2 is ideal for telescopes requiring stable thermal conditions and unobstructed deep-space views. James Webb Space Telescope operates there, observing in infrared without Earth's heat interference. L2 has become one of modern astronomy's most important locations.}{catSS}
\glsadd{Lissajous-Orbit}\dictentry{Lissajous Orbit}{A Lissajous Orbit is a quasi-periodic trajectory around a Lagrange point that follows a bounded but non-repeating path. Named after Lissajous curves in mathematics, it describes complex harmonic motions. Spacecraft in these orbits need periodic station-keeping to maintain position. James Webb Space Telescope initially entered a Lissajous orbit around Sun-Earth L2. These orbits are favored when strict periodicity is not required.}{catSS}
\glsadd{Halo-Orbit}\dictentry{Halo Orbit}{A Halo Orbit is a periodic three-dimensional orbit around a Lagrange point that appears as a closed loop from one of the main bodies. Unlike Lissajous orbits, it repeats regularly and requires active station-keeping. The name comes from its halo-like shape when projected onto the sky. Halo orbits were first proposed for solar observatories at Sun-Earth L1. They enable continuous Earth communication while maintaining advantageous positioning.}{catSS}
\glsadd{Patched-Conics}\dictentry{Patched Conics}{Patched Conics is an approximation method for interplanetary trajectory calculation through multiple gravitational fields. The trajectory is divided into segments each dominated by one body's gravity, treated as conic sections. At boundaries, velocity is adjusted for the new gravitational regime. This provides a useful first-order estimate before detailed numerical integration. It is standard in introductory orbital mechanics courses.}{catSS}
\glsadd{Sphere-of-Influence}\dictentry{Sphere of Influence}{The Sphere of Influence is the region around a body where its gravity dominates over a more distant massive body. Earth's extends to roughly 925,000 kilometers. Within this region, the planet's gravity is the primary force on a spacecraft. The concept defines gravitational regime boundaries in the patched-conics method. Understanding it is essential for planning orbital insertions, departures, and gravity-assist flybys.}{catSS}
\glsadd{Launch-Window}\dictentry{Launch Window}{A Launch Window is the specific time period for launching a spacecraft to achieve a desired trajectory to its target. It depends on relative positions and motions of the launch site, Earth, and target body. Interplanetary windows occur at intervals like the 26-month Earth-to-Mars opportunity. Missing a window may delay a mission by months or years. Planning requires precise celestial mechanics and vehicle performance calculations.}{catSS}
\glsadd{Interplanetary-Transfer}\dictentry{Interplanetary Transfer}{An Interplanetary Transfer is a trajectory carrying a spacecraft between planets within the Solar System. The most common method is the Hohmann transfer using a minimum-energy elliptical orbit. Complex trajectories may employ gravity assists from intermediate planets. These transfers require precise timing, navigation, and adequate propulsion. Missions to Mars, Jupiter, and beyond have demonstrated practical application over decades.}{catSS}
\glsadd{Earth-Orbit}\dictentry{Earth Orbit}{Earth Orbit is the trajectory of objects revolving around Earth under its gravitational field. Objects include the Moon and artificial satellites, space stations, and debris. Orbit type depends on altitude, eccentricity, inclination, and period. Earth orbit is the starting point for most missions, from communications satellites to crewed ISS operations. Understanding orbital dynamics is critical for operations and collision avoidance.}{catSS}
\glsadd{Geosynchronous-Orbit}\dictentry{Geosynchronous Orbit}{A Geosynchronous Orbit has a period equal to Earth's rotation of about 23 hours 56 minutes, placing a satellite over the same sky position daily. At about 35,786 kilometers, it can have any inclination. Inclined geosynchronous satellites trace figure-eight ground patterns. This orbit suits communication, weather, and surveillance satellites. The region is a limited resource requiring careful satellite placement coordination.}{catSS}
\glsadd{Geostationary-Transfer-Orbit}\dictentry{Geostationary Transfer Orbit}{A Geostationary Transfer Orbit, or GTO, is an elliptical orbit transferring a spacecraft from a low Earth parking orbit to geosynchronous orbit. Perigee is at LEO altitude while apogee reaches 35,786 kilometers. A circularization burn at apogee and inclination change achieve the final GEO. GTO is the standard deployment orbit for commercial communication satellites. Required delta-v depends on initial orbit inclination and altitude.}{catSS}
\glsadd{Highly-Elliptical-Orbit}\dictentry{Highly Elliptical Orbit}{A Highly Elliptical Orbit, or HEO, has high eccentricity creating a long narrow path with low perigee and very high apoapsis. Satellites spend most of their period near apoapsis moving slowly, providing extended high-latitude coverage. This makes HEO useful for communications where geostationary satellites have poor polar visibility. The Molniya orbit is a well-known HEO example. These orbits exploit the period-altitude relationship.}{catSS}
\glsadd{Molniya-Orbit}\dictentry{Molniya Orbit}{A Molniya Orbit is a highly elliptical orbit with 63.4-degree inclination and approximately 12-hour period. Apogee is positioned over high northern latitudes for extended coverage. Designed by Soviet engineers in the 1960s, it provides continuous communications for Russia's northern territories. The specific inclination prevents apsidal precession, keeping apogee fixed. Molniya orbits remain in use for high-latitude communications.}{catSS}
\glsadd{Tundra-Orbit}\dictentry{Tundra Orbit}{A Tundra Orbit is a highly elliptical geosynchronous orbit with about 63.4-degree inclination and 24-hour period. It provides continuous coverage of high-latitude regions, particularly the northern hemisphere. The satellite traces a figure-eight ground track, spending most of its period near apogee. Used by communication satellites serving northern countries, it prevents apsidal rotation to maintain stable apogee position.}{catSS}
\glsadd{Dawn-Dusk-Orbit}\dictentry{Dawn-Dusk Orbit}{A Dawn-Dusk Orbit is a Sun-synchronous orbit crossing the equator at approximately 6:00 AM and 6:00 PM local solar time. This places the orbit along Earth's terminator, the day-night dividing line. Satellites receive nearly continuous sunlight, advantageous for solar-powered spacecraft. The orbit also provides consistent lighting for day-night boundary imaging. It is commonly used for Earth-observation and some radar satellite constellations.}{catSS}
\glsadd{Frozen-Orbit}\dictentry{Frozen Orbit}{A Frozen Orbit maintains nearly constant orbital elements despite gravitational perturbations. This is achieved by selecting specific eccentricity and argument of perigee where perturbative effects cancel. Frozen orbits are valuable for Earth-observation requiring consistent altitude over specific latitudes. They reduce station-keeping fuel needs and extend satellite life. The concept is important for altimetry and gravity-mapping missions.}{catSS}
\glsadd{Graveyard-Orbit}\dictentry{Graveyard Orbit}{A Graveyard Orbit, also called a disposal orbit, is a supersynchronous orbit for safely decommissioning satellites at end of life. GEO satellites are moved approximately 200 to 300 kilometers above their operational orbit to prevent collisions. The move requires a final propulsive altitude-raising maneuver. International guidelines recommend this practice to preserve limited GEO slots. It is an important space debris mitigation strategy.}{catSS}
\glsadd{Parking-Orbit}\dictentry{Parking Orbit}{A Parking Orbit is a temporary orbit occupied after launch before a subsequent maneuver to the final destination. It provides a holding position for system verification and next-phase planning. Parking orbits are typically low Earth orbits but can vary by mission. Spacecraft may remain for minutes, hours, or days before the next burn. The parking orbit allows flexibility in launch timing and trajectory planning.}{catSS}
\glsadd{Phasing-Orbit}\dictentry{Phasing Orbit}{A Phasing Orbit adjusts a spacecraft's position relative to a target by spending one or more revolutions at a different orbital period. The period difference creates predictable angular displacement per revolution. This technique is used for rendezvous, docking, and satellite constellation positioning. It is a fundamental tool in orbital rendezvous and mission planning. The phasing orbit ensures correct positioning for the next maneuver.}{catSS}
\glsadd{Co-orbital}\dictentry{Co-orbital}{Co-orbital describes two or more objects sharing the same orbit around a central body, maintaining similar altitude, inclination, and period. Objects may be separated by position along the orbit, their mean anomaly. Natural Trojan asteroids share a planet's orbit at stable Lagrange points. In satellite operations, co-orbital configurations serve constellation management and proximity operations. The concept is also studied for debris avoidance.}{catSS}
\glsadd{Retrograde-Orbit}\dictentry{Retrograde Orbit}{A Retrograde Orbit moves opposite to its central body's rotation direction. For Earth, retrograde orbits have inclinations greater than 90 degrees. Launching into retrograde is more energy-intensive as it cannot use Earth's rotational velocity. However, retrograde orbits offer unique advantages for specific missions including certain Sun-synchronous configurations. The concept applies to any orbiting body and is important in mission design.}{catSS}
\glsadd{Prograde-Orbit}\dictentry{Prograde Orbit}{A Prograde Orbit moves in the same direction as its central body's rotation. For Earth, prograde orbits have inclinations less than 90 degrees. Launching prograde benefits from Earth's rotational velocity, reducing energy required. Most satellites, including geosynchronous and low Earth orbit types, use prograde orbits. The prograde direction is the most common and fuel-efficient orientation for most space missions.}{catSS}
\glsadd{Ascending-Node}\dictentry{Ascending Node}{The Ascending Node is where an orbit crosses the reference plane moving from south to north. It is one of two nodes defining the orbital plane's orientation in space. Its position is specified by the longitude of the ascending node, one of six classical orbital elements. Nodal precession from Earth's oblateness causes it to drift slowly over time. The ascending node is a key reference point for orbital element definition.}{catSS}
\glsadd{Descending-Node}\dictentry{Descending Node}{The Descending Node is where an orbit crosses the reference plane moving from north to south, opposite the ascending node. Together they define the line of nodes. The time between ascending and descending node passage is exactly half an orbital period. Like the ascending node, it precesses over time due to gravitational perturbations. The descending node is used with other parameters to fully specify orbit orientation and position.}{catSS}
\glsadd{Line-of-Nodes}\dictentry{Line of Nodes}{The Line of Nodes connects the ascending and descending nodes of an orbit, lying in the reference plane. It defines where the orbital plane intersects the reference plane. Its orientation is specified by the longitude of the ascending node. Nodal precession from Earth's oblateness is key to Sun-synchronous orbit design. The line of nodes is important for planning eclipse periods and solar-powered spacecraft shadow times.}{catSS}
\glsadd{Argument-of-Perigee}\dictentry{Argument of Perigee}{The Argument of Perigee is the angle measured in the orbital plane from the ascending node to perigee, the closest approach point. It defines the orbit ellipse's orientation within the orbital plane. Together with other classical elements, it fully specifies orbit geometry. Changes from gravitational perturbations are significant for long-duration missions. It is critical for frozen orbit design where it is held nearly constant.}{catSS}
\glsadd{Right-Ascension-of-Ascending-Node}\dictentry{Right Ascension of Ascending Node}{The Right Ascension of Ascending Node, or RAAN, is the angle in the reference plane from a reference direction to the ascending node. It defines the orbital plane's orientation in inertial space. RAAN is one of the six classical Keplerian orbital elements. It changes over time due to gravitational perturbations, especially Earth's oblateness. RAAN is essential for constellation spacing coordination and inter-satellite link planning.}{catSS}
\glsadd{True-Anomaly}\dictentry{True Anomaly}{True Anomaly is the angular position of a satellite along its orbit measured from periapsis to the satellite's current position. It ranges from 0 to 360 degrees and increases as the satellite moves. Unlike mean anomaly, true anomaly directly describes the satellite's actual position. The relationship between true anomaly and time follows Kepler's equation. It is one of the six classical orbital elements.}{catSS}
\glsadd{Mean-Anomaly}\dictentry{Mean Anomaly}{Mean Anomaly is an angular parameter describing a satellite's orbital position assuming constant angular speed. It increases linearly with time, ranging from 0 to 360 degrees per orbit. It provides simplified position tracking without accounting for eccentric orbit speed variations. Mean anomaly relates to true anomaly through Kepler's equation and the eccentric anomaly. It is widely used in orbit propagation for its simple time dependence.}{catSS}
\glsadd{Eccentric-Anomaly}\dictentry{Eccentric Anomaly}{Eccentric Anomaly is an angular parameter relating mean and true anomalies in Kepler's equation. Measured from the auxiliary circle's center to the satellite's projection, it serves as an intermediate variable in position calculations. It ranges from 0 to 360 degrees per orbit and varies nonlinearly with time. The parameter is fundamental to analytical orbit propagation used in spacecraft navigation.}{catSS}
\glsadd{Semi-Major-Axis}\dictentry{Semi-Major Axis}{The Semi-Major Axis is half the longest diameter of an elliptical orbit from center to farthest point. It represents the average distance from the central body and relates directly to orbital period through Kepler's Third Law. One of six classical orbital elements, it determines orbit size. Larger semi-major axes correspond to longer periods and higher total energy. It is fundamental to mission planning and orbit determination.}{catSS}
\glsadd{Semi-Minor-Axis}\dictentry{Semi-Minor Axis}{The Semi-Minor Axis is half the shortest diameter of an elliptical orbit, perpendicular to the semi-major axis. Together with the semi-major axis, it defines the ellipse's shape and size. It relates to the semi-major axis and eccentricity through the square root of one minus eccentricity squared. A circular orbit has equal semi-minor and semi-major axes. The parameter is important for understanding non-circular orbit geometry.}{catSS}
\glsadd{Inclination}\dictentry{Inclination}{Inclination is the tilt of an orbital plane relative to a reference plane, measured in degrees. For Earth orbits the reference is the equatorial plane; for solar orbits it is the ecliptic. Values range from zero for equatorial orbits to ninety degrees for polar orbits. Inclination directly determines the latitudes a satellite can observe and governs mission design and ground coverage patterns.}{catSS}
\glsadd{Periapsis}\dictentry{Periapsis}{Periapsis is the point in an elliptical orbit where a spacecraft is closest to the body it orbits. At periapsis, orbital velocity reaches its maximum and gravitational potential energy is at its minimum. The altitude of periapsis determines closest approach distance and affects atmospheric drag in low orbits. Mission planners carefully select periapsis altitude to balance scientific objectives with orbital lifetime.}{catSS}
\glsadd{Apoapsis}\dictentry{Apoapsis}{Apoapsis is the point in an elliptical orbit where a spacecraft is farthest from the body it orbits. Velocity is at its minimum and gravitational potential energy at its maximum at this point. The apoapsis altitude defines the outer extent of the orbit and influences the overall orbital period. Understanding apoapsis is essential for planning orbit insertion and circularization maneuvers.}{catSS}
\glsadd{Perigee}\dictentry{Perigee}{Perigee is the closest point in an Earth-centered orbit to Earth's center. It is the Earth-specific term for periapsis, where a satellite travels fastest. Perigee altitude is critical for low Earth orbit missions, governing atmospheric drag exposure and radiation environment intensity. Satellites in elliptical orbits experience their strongest gravitational effects at perigee.}{catSS}
\glsadd{Apogee}\dictentry{Apogee}{Apogee is the farthest point in an Earth-centered orbit from Earth's center. It is the Earth-specific term for apoapsis, where a satellite moves slowest. Apogee altitude is a key parameter for geostationary transfer orbits, where a burn at perigee raises the apogee to geostationary height. The apogee distance also influences the radiation environment a spacecraft encounters.}{catSS}
\glsadd{Perihelion}\dictentry{Perihelion}{Perihelion is the point in a solar orbit closest to the Sun. Earth reaches perihelion in early January at roughly 147.1 million kilometers. At perihelion an orbiting body attains its highest velocity due to stronger gravitational pull. This parameter is important for understanding seasonal variations and planning interplanetary trajectories that exploit gravitational assists.}{catSS}
\glsadd{Aphelion}\dictentry{Aphelion}{Aphelion is the point in a solar orbit farthest from the Sun. Earth reaches aphelion in early July at approximately 152.1 million kilometers. Orbital velocity is at its minimum for a given elliptical orbit at this point. Aphelion distance affects solar flux received by a spacecraft and is a key consideration for thermal planning in interplanetary missions.}{catSS}
\glsadd{Synodic-Period}\dictentry{Synodic Period}{Synodic period is the time between successive identical alignments of a body relative to the Sun as seen from Earth. It differs from the sidereal period because Earth simultaneously moves in its own orbit. Mars has a synodic period of about 780 days, which defines the launch window cycle for Mars missions. The synodic period governs planetary oppositions and conjunctions and is essential for interplanetary planning.}{catSS}
\glsadd{Sidereal-Period}\dictentry{Sidereal Period}{Sidereal period is the time a body takes to complete one orbit relative to the fixed background stars. Unlike the synodic period it represents the true orbital period independent of Earth's motion. Earth's sidereal year is approximately 365.256 days, slightly longer than the solar year. This measurement is fundamental to celestial mechanics and serves as the baseline for comparing orbital characteristics of all solar system bodies.}{catSS}
\glsadd{Nodal-Precession}\dictentry{Nodal Precession}{Nodal precession is the gradual rotation of the line of nodes where an orbit crosses a reference plane. It arises from gravitational perturbations, most commonly the oblateness of the central body which torques the orbital plane. For Earth orbits the J2 effect causes the ascending node to drift several degrees per day at low altitudes.}{catSS}
\glsadd{Apsidal-Precession}\dictentry{Apsidal Precession}{Apsidal precession is the gradual rotation of the line of apsides connecting periapsis and apoapsis within the orbital plane. It results from perturbations including the central body's oblateness, third-body gravity, and general relativistic effects. Mercury's apsidal precession famously confirmed Einstein's general theory of relativity. Understanding this effect is essential for precise orbit determination and long-term trajectory prediction.}{catSS}
\glsadd{Perturbation}\dictentry{Perturbation}{Perturbation is any deviation from ideal two-body Keplerian motion caused by additional forces on a spacecraft. These include atmospheric drag, solar radiation pressure, third-body gravity, and the nonspherical shape of the central body. Perturbations cause gradual changes in orbital elements such as eccentricity, inclination, and argument of perigee.}{catSS}
\glsadd{J2-Perturbation}\dictentry{J2 Perturbation}{J2 perturbation is the dominant orbital disturbance for Earth-orbiting satellites, caused by Earth's equatorial bulge. The J2 zonal harmonic describes the deviation from a perfect sphere, producing torques that drive nodal and apsidal precession. Effects are strongest at low altitudes where the bulge influence is greatest.}{catSS}
\glsadd{Third-Body-Perturbation}\dictentry{Third Body Perturbation}{Third body perturbation is the gravitational disturbance from a body other than the primary central body. The Sun, Moon, and planets each exert tugs that alter a spacecraft's orbit over time. For high-altitude Earth orbits, lunar and solar perturbations can significantly change eccentricity and inclination over months to years. These perturbations are critical for geostationary station-keeping calculations.}{catSS}
\glsadd{Drag-Perturbation}\dictentry{Drag Perturbation}{Drag perturbation is the orbital decay caused by atmospheric drag on a spacecraft in low Earth orbit. Trace atmospheric particles create a persistent retarding force that reduces orbital energy. This causes the orbit to circularize and the satellite to gradually descend toward reentry. Drag perturbation is the dominant factor determining orbital lifetime for low-altitude satellites and space debris.}{catSS}
\glsadd{Radiation-Pressure}\dictentry{Radiation Pressure}{Solar radiation pressure is the force exerted on a spacecraft by photons from the Sun. Though tiny compared to gravity, it becomes significant for large, low-mass spacecraft or over extended durations. The force depends on solar flux, cross-sectional area, and surface reflectivity. Radiation pressure must be accounted for in precise orbit determination and is exploited as a propellantless propulsion mechanism in solar sail missions.}{catSS}
\glsadd{Poynting-Robertson-Effect}\dictentry{Poynting-Robertson Effect}{The Poynting-Robertson effect is a drag force on small orbiting dust particles caused by absorption and re-emission of solar radiation. As particles absorb sunlight and re-emit it isotropically, the forward component of motion creates a slight braking force. Dust particles slowly spiral inward toward the Sun over thousands to millions of years. The effect shapes circumstellar disks and governs the clearing of interplanetary dust from the solar system.}{catSS}
\glsadd{Yarkovsky-Effect}\dictentry{Yarkovsky Effect}{The Yarkovsky effect is a small force on rotating asteroids caused by anisotropic thermal emission of absorbed solar radiation. As an asteroid rotates, heated surfaces radiate infrared photons producing a tiny thrust that slowly alters the orbit over millions of years. The magnitude depends on size, spin, thermal properties, and orbital characteristics.}{catSS}
\glsadd{Roche-Limit}\dictentry{Roche Limit}{The Roche limit is the minimum distance at which a self-gravitating body can orbit a more massive body without being torn apart by tidal forces. Inside this limit tidal forces exceed self-gravity, causing disintegration. The limit depends on the densities of both bodies and is typically about 2.44 times the larger body's radius for fluid satellites. This concept explains planetary ring formation and constrains natural satellite origins.}{catSS}
\glsadd{Hill-Sphere}\dictentry{Hill Sphere}{The Hill sphere defines the region around a celestial body where it dominates gravitationally despite a more massive parent body nearby. Its radius depends on the mass ratio and orbital separation of the two bodies. Earth's Hill sphere extends roughly 1.5 million kilometers, well beyond the Moon's orbit. Natural satellites must orbit within the Hill sphere for long-term stability.}{catSS}
\glsadd{Tidal-Force}\dictentry{Tidal Force}{Tidal force is the differential gravitational pull across a body due to the variation in gravitational attraction from a nearby massive object. The near side experiences a stronger pull than the far side, creating a stretching effect. Tidal forces cause ocean tides on Earth, form planetary rings, and gradually circularize satellite orbits. Over geological time they can significantly alter the shapes and rotations of moons and planets.}{catSS}
\glsadd{Tidal-Locking}\dictentry{Tidal Locking}{Tidal locking occurs when a body's rotation period equals its orbital period, causing one hemisphere to permanently face its parent body. Gravitational torques gradually slow the rotation until it synchronizes with orbital motion. The Moon is tidally locked to Earth, which is why we always see the same face. Most large moons in the solar system are tidally locked to their host planets.}{catSS}
\glsadd{Synchronous-Rotation}\dictentry{Synchronous Rotation}{Synchronous rotation describes a body that rotates on its axis in exactly the time it takes to complete one orbit. This is the end state of tidal locking, where gravitational torques have synchronized spin and orbital periods. One hemisphere permanently faces the parent body while the opposite side always faces away. This condition is common among major moons and is a design target for certain communication satellite orbits.}{catSS}
\glsadd{Space-Station}\dictentry{Space Station}{A space station is a large, habitable artificial satellite designed for long-term human presence in orbit. It provides a microgravity laboratory where crews conduct experiments in physics, biology, and materials science. Stations feature life support systems, docking ports for visiting vehicles, and robotic arms for assembly and maintenance. The International Space Station has been continuously occupied since November 2000.}{catSS}
\glsadd{ISS}\dictentry{ISS}{The International Space Station is a modular station in low Earth orbit at roughly 408 kilometers, assembled by fifteen nations. It has been continuously inhabited since November 2000 as a microgravity research laboratory. The station orbits every 90 minutes at about 28,000 kilometers per hour and serves as a testbed for deep space technologies. With mass exceeding 400,000 kilograms it is the largest artificial object in Earth orbit.}{catSS}
\glsadd{Capsule}\dictentry{Capsule}{A space capsule is a crewed or cargo spacecraft with a blunt-body design that reenters the atmosphere using a heat shield. Capsules are launched atop rockets and return by splashing down in the ocean or landing on solid ground. Their simple, robust design makes them well-suited for transporting crew and supplies to orbiting spacecraft. Modern capsules like Crew Dragon and Starliner incorporate advanced avionics and abort systems for crew safety.}{catSS}
\glsadd{Space-Probe}\dictentry{Space Probe}{A space probe is an unmanned spacecraft designed to explore destinations beyond Earth orbit including planets, moons, asteroids, and comets. Probes carry instruments to study composition, atmospheres, magnetic fields, and geology of their targets. They communicate findings through deep space network antennas using radio signals. Notable probes include Voyager, which has traveled beyond the solar system, and New Horizons which flew by Pluto.}{catSS}
\glsadd{Lander}\dictentry{Lander}{A lander is a spacecraft designed to descend from orbit or transfer trajectory and make a controlled landing on a celestial body. Landers manage descent propulsion, navigation, and hazard avoidance to achieve a safe touchdown. They often carry instruments for surface analysis and may deploy rovers or other payloads. Successful landers include Mars Pathfinder, Philae on comet 67P, and the Chang'e lunar landers.}{catSS}
\glsadd{Rover}\dictentry{Rover}{A rover is a mobile robotic vehicle designed to traverse and explore a celestial body's surface. Rovers carry cameras, spectrometers, drills, and other instruments to study geology and search for signs of water. They navigate autonomously or under remote control over terrain that fixed landers cannot access. NASA's Mars rovers including Spirit, Opportunity, Curiosity, and Perseverance have collectively traveled tens of kilometers across Mars.}{catSS}
\glsadd{Orbiter}\dictentry{Orbiter}{An orbiter is a spacecraft designed to enter and remain in orbit around a celestial body for extended scientific observation. Orbiters carry remote sensing instruments including cameras, spectrometers, and radar to map surfaces and study atmospheres from above. They serve as communication relays between Earth and surface assets like landers and rovers. Orbiters have mapped every planet in the solar system and several moons.}{catSS}
\glsadd{Flyby-Spacecraft}\dictentry{Flyby Spacecraft}{A flyby spacecraft passes close to a target body without entering orbit, collecting data during a brief encounter. Flyby missions are fuel-efficient since they require no orbital insertion burn, allowing visits to multiple targets via gravity assists. Data collection occurs over hours or days as the spacecraft approaches, passes, and recedes. Voyager executed flybys of all four giant planets, while New Horizons flew by Pluto and Arrokoth.}{catSS}
\glsadd{Sample-Return-Mission}\dictentry{Sample Return Mission}{A sample return mission collects material from a celestial body and returns it to Earth for laboratory analysis. These missions are among the most complex in exploration, requiring precise landing, sample collection, ascent, rendezvous, and reentry. Returned samples can be studied with instruments far more capable than any robotic payload. Examples include the Hayabusa asteroid missions and the Apollo lunar sample returns.}{catSS}
\glsadd{Crewed-Spacecraft}\dictentry{Crewed Spacecraft}{A crewed spacecraft is a vehicle designed to transport humans safely through space and return them to Earth. These vehicles require life support systems, radiation shielding, thermal control, and robust abort capabilities. They must withstand extreme launch and reentry accelerations while maintaining a habitable interior. Modern crewed spacecraft include SpaceX Crew Dragon, Boeing Starliner, and China's Shenzhou.}{catSS}
\glsadd{Space-Suit}\dictentry{Space Suit}{A space suit is a pressurized garment that protects astronauts from the vacuum, temperature extremes, and radiation of space. It provides oxygen, removes carbon dioxide, maintains pressure, and regulates temperature. A full suit includes a torso assembly, helmet, gloves, boots, and a portable life support system. Modern suits allow several hours of work outside a spacecraft with advanced joints and communication systems.}{catSS}
\glsadd{EVA-Suit}\dictentry{EVA Suit}{An EVA suit is a specialized space suit designed for extravehicular activity in the vacuum of space. It provides complete life support including oxygen supply, carbon dioxide removal, and temperature regulation for typically six to eight hours. Bearing joints at major articulation points allow pressurized mobility. Current NASA designs include the EMU for orbital operations and the xEMU planned for lunar surface EVAs.}{catSS}
\glsadd{Life-Support-System}\dictentry{Life Support System}{A life support system provides essential environmental conditions to sustain human life aboard a spacecraft. It manages oxygen generation, carbon dioxide removal, water recycling, temperature regulation, and waste management. These systems must operate reliably for months or years with minimal maintenance. Advances in life support are critical for enabling long-duration lunar and Mars missions where resupply from Earth is impractical.}{catSS}
\glsadd{ECLSS}\dictentry{ECLSS}{The Environmental Control and Life Support System manages the habitable environment aboard the International Space Station. It generates oxygen, removes carbon dioxide, recycles water from humidity and urine, and controls temperature and humidity. ECLSS represents decades of engineering refinement in closed-loop environmental management. Its technologies provide the foundation for life support on future lunar and Martian habitats.}{catSS}
\glsadd{Habitat-Module}\dictentry{Habitat Module}{A habitat module is a pressurized spacecraft component providing living quarters for crew during extended missions. It includes sleeping areas, workstations, galley facilities, hygiene stations, and exercise equipment. Habitat modules are launched separately and assembled in orbit using docking or berthing mechanisms. ISS modules like Unity, Destiny, and Columbus sustain continuous human presence in space.}{catSS}
\glsadd{Docking-Mechanism}\dictentry{Docking Mechanism}{A docking mechanism enables two spacecraft to join in orbit, creating a rigid and airtight connection. It must align approaching vehicles, capture them, and retract to form a hard structural mate. Modern docking systems include sensors, latches, and seals for alignment, capture, and structural connection. Systems must meet international standards to enable cooperation between different agencies and commercial spacecraft.}{catSS}
\glsadd{Androgynous-Docking}\dictentry{Androgynous Docking}{Androgynous docking uses identical interfaces on both mating spacecraft, allowing any two equipped vehicles to dock. This eliminates separate active and passive components, simplifying logistics and increasing mission flexibility. The International Docking Adapter on the ISS is androgynous, compatible with NASA, ESA, JAXA, and commercial vehicles.}{catSS}
\glsadd{Probe-and-Drogue}\dictentry{Probe and Drogue}{Probe and drogue is a docking system where one spacecraft extends a probe into a receiving cone on the other vehicle. The probe captures inside the drogue and retraction pulls the spacecraft together for a hard mate. This asymmetric system was used extensively during Apollo-Soyuz and on early space stations. While reliable, it requires designated active and passive roles for each vehicle.}{catSS}
\glsadd{Capture}\dictentry{Capture}{Capture is the moment of initial mechanical contact and engagement between two docking spacecraft. During capture, latches or hooks from one vehicle engage corresponding structures on the other, establishing a temporary connection. The capture phase must absorb relative motion without causing structural damage. Successful capture is the critical first step leading to pressurized seal and structural connection between spacecraft.}{catSS}
\glsadd{Soft-Capture}\dictentry{Soft Capture}{Soft capture is the initial docking phase where two spacecraft make gentle contact with minimal relative force. Springs and dampers in the docking mechanism absorb approach velocity and misalignment, preventing structural damage. Capture latches temporarily hold the vehicles together while motion is damped. This phase precedes retraction and establishment of the rigid, pressurized connection known as hard capture.}{catSS}
\glsadd{Hard-Capture}\dictentry{Hard Capture}{Hard capture is the final docking phase where two spacecraft are pulled into a rigid, structurally connected, pressurized seal. After soft capture and damping, the mechanism retracts and structural latches engage. Hard capture creates an airtight seal allowing crew and cargo transfer. This connection must withstand maneuver loads and thermal cycling while maintaining pressure integrity throughout the docked period.}{catSS}
\glsadd{Berthing}\dictentry{Berthing}{Berthing attaches a spacecraft module to a station using a robotic arm for precise positioning into a docking port. Unlike autonomous docking, the arm captures the arriving vehicle and guides it into alignment with the berthing port. The Common Berthing Mechanism provides a larger hatch than docking adapters, allowing oversized cargo and laboratory rack transfer. Berthing is used for large ISS laboratory and logistics vehicles.}{catSS}
\glsadd{Robotic-Arm}\dictentry{Robotic Arm}{A robotic arm is a multi-jointed manipulator on a spacecraft or station for handling payloads and assisting assembly. It provides precise force-controlled movement in microgravity for satellite deployment, module installation, and external instrument maintenance. Robotic arms are operated by crew inside the station or by ground controllers on Earth. They are essential for constructing large space structures that cannot be assembled by EVA alone.}{catSS}
\glsadd{Canadarm}\dictentry{Canadarm}{Canadarm is the name of the robotic arms on the Space Shuttle and International Space Station, built by the Canadian Space Agency. The original Shuttle arm was 15.2 meters long for satellite deployment and retrieval. Canadarm2 on the ISS is 17.6 meters long and walks end-over-end between grapple fixtures. It has been instrumental in ISS assembly, berthing visiting vehicles, and supporting spacewalks.}{catSS}
\glsadd{Remote-Manipulator-System}\dictentry{Remote Manipulator System}{A remote manipulator system is a robotic arm operated from inside a spacecraft that grapples, manipulates, and deploys objects outside the vehicle. It translates crew hand controller inputs into precise robotic movements in real time. Applications include satellite deployment, free-flyer capture, and external hardware installation. The Shuttle and Station Remote Manipulator Systems have demonstrated robotic operations versatility in orbit.}{catSS}
\glsadd{Space-Tether}\dictentry{Space Tether}{A space tether is a long cable connecting two objects in orbit, typically extending several kilometers. Tethers generate electricity via electromagnetic interaction with a planetary field, provide artificial gravity through rotation, or serve as structural connections. Tethered systems involve complex orbital mechanics since endpoints trace non-Keplerian paths.}{catSS}
\glsadd{Electrodynamic-Tether}\dictentry{Electrodynamic Tether}{An electrodynamic tether generates current by cutting through a planetary magnetic field as it moves in orbit. The interaction between induced voltage and surrounding ionosphere creates current that powers spacecraft or produces thrust without propellant. Driving current in reverse acts as an electromagnetic thruster for orbit changes. This technology promises propellantless propulsion and power for future space missions.}{catSS}
\glsadd{Solar-Sail}\dictentry{Solar Sail}{A solar sail uses pressure of sunlight on a large reflective surface to generate continuous thrust without propellant. Individual photon pressure is tiny, but cumulative force on a large lightweight sail produces measurable acceleration over time. Solar sails can achieve high velocities for inner solar system missions or non-Keplerian orbits. IKAROS and LightSail 2 have demonstrated successful solar sail propulsion in space.}{catSS}
\glsadd{Photovoltaic-Cell}\dictentry{Photovoltaic Cell}{A photovoltaic cell is a semiconductor device that converts light energy directly into electrical voltage through the photovoltaic effect. Photons excite electrons across the semiconductor bandgap, creating a current flow. Space-grade cells use gallium arsenide and multi-junction designs for maximum efficiency under the solar spectrum.}{catSS}
\glsadd{Solar-Array}\dictentry{Solar Array}{A solar array is a complete assembly of panels, support structures, wiring, and power management hardware for spacecraft electrical power. Arrays deploy after launch, unfold to full area, and track the Sun for maximum generation. They range from small CubeSat panels to ISS wings generating over 120 kilowatts. Design must balance power output, mass, deployable area, and resilience to the space environment.}{catSS}
\glsadd{RTG}\dictentry{RTG}{An RTG is a Radioisotope Thermoelectric Generator that converts heat from radioactive decay into electricity using thermocouples. It uses plutonium-238 as its heat source with an 87.7-year half-life. RTGs provide reliable long-lasting power independent of sunlight for deep space missions. They have powered Voyager, Cassini, and the Curiosity and Perseverance Mars rovers.}{catSS}
\glsadd{Radioisotope-Thermoelectric-Generator}\dictentry{Radioisotope Thermoelectric Generator}{A Radioisotope Thermoelectric Generator converts heat from radioactive isotope decay into electrical power using the Seebeck effect in thermocouples. Plutonium-238 produces heat as it decays, and thermocouple junctions at different temperatures generate voltage. RTGs have no moving parts, require no maintenance, and operate reliably for decades. They are the power choice for missions beyond Jupiter and surface missions on Mars.}{catSS}
\glsadd{Nuclear-Power-in-Space}\dictentry{Nuclear Power in Space}{Nuclear power in space encompasses radioisotope systems like RTGs and fission reactors for spacecraft and surface installations. These systems offer high energy density, independence from sunlight, and operation in extreme environments. They are essential for deep space exploration, lunar operations during the long night, and high-power electric propulsion.}{catSS}
\glsadd{Fuel-Cell}\dictentry{Fuel Cell}{A fuel cell generates electricity by combining hydrogen and oxygen through an electrochemical reaction, producing water as a byproduct. Unlike batteries, fuel cells produce power continuously as long as fuel is supplied. The Space Shuttle used fuel cells that also provided drinking water for the crew. Fuel cells offer higher energy density than batteries for extended missions and are being developed for lunar surface power systems.}{catSS}
\glsadd{Attitude-Control-System}\dictentry{Attitude Control System}{An attitude control system determines and controls a spacecraft's orientation in three-dimensional space. It uses sensors like star trackers and gyroscopes to measure orientation and actuators like reaction wheels and thrusters to apply torques. Precise control is essential for pointing antennas, solar panels, and instruments at their targets. The system must maintain accuracy while managing momentum and minimizing propellant consumption.}{catSS}
\glsadd{Control-Moment-Gyro}\dictentry{Control Moment Gyro}{A control moment gyro is a gyroscope whose spin axis is gimbaled, producing large torques by redirecting angular momentum. Unlike reaction wheels that change speed, CMGs redirect a constant-speed rotor's momentum to generate much higher torque for the same mass. This makes them ideal for large stations and agile spacecraft requiring rapid slewing. The ISS uses CMGs as its primary attitude control actuators.}{catSS}
\glsadd{Thruster}\dictentry{Thruster}{A spacecraft thruster is a small rocket engine for attitude control, orbital maneuvers, and velocity corrections. Thrusters use monopropellant or bipropellant chemical systems producing controlled thrust bursts. They maintain orbit, perform station-keeping, execute proximity operations, and desaturate reaction wheels. Modern spacecraft may use chemical, ion, and Hall-effect engines, each optimized for different thrust and efficiency requirements.}{catSS}
\glsadd{Horizon-Sensor}\dictentry{Horizon Sensor}{A horizon sensor detects a planet's curved atmospheric edge to determine the spacecraft's local vertical direction. By measuring infrared radiation contrast between the warm planet and cold space, it identifies the horizon and calculates nadir direction. Horizon sensors provide attitude reference for Earth-pointing spacecraft, enabling accurate antenna and instrument pointing.}{catSS}
\glsadd{Earth-Sensor}\dictentry{Earth Sensor}{An Earth sensor detects Earth's horizon to establish a local vertical reference for the spacecraft. It uses infrared detectors sensing the thermal contrast between Earth's atmosphere and cold space. Earth sensors are essential for maintaining precise Earth-pointing orientation required by communication, navigation, and observation satellites.}{catSS}
\glsadd{Inertial-Measurement-Unit}\dictentry{Inertial Measurement Unit}{An inertial measurement unit combines three accelerometers and three gyroscopes to measure acceleration and rotation in all axes. Integrating these measurements over time tracks velocity and attitude changes without external references. IMUs are essential for launch vehicle guidance, spacecraft navigation during burns, and attitude determination during sensor outages.}{catSS}
\glsadd{Attitude-Propagation}\dictentry{Attitude Propagation}{Attitude propagation predicts a spacecraft's future orientation from a known state and measured angular velocities. Gyroscope data provides rotation rates integrated forward in time between sensor updates. Propagation accounts for internal dynamics like wheel speed changes and external torques from gravity gradients and magnetic interactions.}{catSS}
\glsadd{Despin}\dictentry{Despin}{Despin is the process of reducing the spin rate of a spacecraft that was spin-stabilized during launch. It is accomplished using yo-yo weights, despun platforms, or thruster firings to counteract angular momentum. Transition from spinning to three-axis stabilized mode is a critical commissioning phase. Successful despin allows fixed-pointing of antennas, solar panels, and instruments for the primary mission.}{catSS}
\glsadd{Spin-Stabilization}\dictentry{Spin Stabilization}{Spin stabilization maintains spacecraft orientation by spinning about the major axis, exploiting gyroscopic rigidity of angular momentum. A spinning spacecraft resists external torques, providing simple and reliable attitude control. This technique was widely used in early communications satellites and interplanetary probes.}{catSS}
\glsadd{Three-Axis-Stabilization}\dictentry{Three-Axis Stabilization}{Three-axis stabilization actively maintains orientation about all three axes simultaneously without spinning. Sensors measure attitude while reaction wheels and thrusters correct deviations in real time. This provides continuous pointing in any direction for Earth observation, astronomy, and communications.}{catSS}
\glsadd{Gravity-Gradient-Stabilization}\dictentry{Gravity Gradient Stabilization}{Gravity gradient stabilization is a passive method using the difference in gravitational force across an extended spacecraft to maintain orientation. A long boom with an end mass creates a restoring torque aligning the long axis toward the planet. This technique requires no propellant or power, making it attractive for simple long-life missions.}{catSS}
\glsadd{Magnetic-Torquer}\dictentry{Magnetic Torquer}{A magnetic torquer is an electromagnet that interacts with a planet's magnetic field to produce controlled torque on a spacecraft. Passing current through wire coils creates a magnetic dipole that aligns or precesses the spacecraft relative to the local field. Magnetic torquers commonly desaturate reaction wheels by bleeding accumulated angular momentum.}{catSS}
\glsadd{Thermal-Control-System}\dictentry{Thermal Control System}{A thermal control system maintains spacecraft components within required temperature ranges despite the extreme space thermal environment. It manages heat from solar radiation, electronics, and albedo while rejecting excess heat to space. The system combines passive techniques like insulation and coatings with active methods including heaters, heat pipes, and fluid loops.}{catSS}
\glsadd{Passive-Thermal-Control}\dictentry{Passive Thermal Control}{Passive thermal control uses materials and design features to regulate spacecraft temperature without power or moving parts. Techniques include multi-layer insulation, surface coatings with specific optical properties, heat pipes, and thermal straps. These methods are highly reliable with no mechanical or electrical failure modes.}{catSS}
\glsadd{Active-Thermal-Control}\dictentry{Active Thermal Control}{Active thermal control uses powered devices for precise temperature regulation where passive methods fall short. Components include electric heaters, pumped fluid loops, cryocoolers, and variable-conductance heat pipes. Active systems respond dynamically to changing thermal loads, maintaining tight tolerances for sensitive instruments.}{catSS}
\glsadd{Multi-Layer-Insulation}\dictentry{Multi-Layer Insulation}{Multi-layer insulation consists of many thin reflective film layers separated by low-conductivity spacers wrapped around spacecraft to minimize heat transfer. Reflective layers of aluminized Mylar or Kapton bounce thermal radiation between layers, reducing radiative exchange. MLI blankets are custom-fitted to each component and are among the most effective passive thermal controls.}{catSS}
\glsadd{Radiator}\dictentry{Radiator}{A radiator is a dedicated spacecraft surface designed to reject excess heat to space through thermal radiation. Radiators use high-emissivity coatings to efficiently radiate infrared energy away from the spacecraft. They connect to internal heat sources through heat pipes or fluid loops transporting waste heat from electronics and power systems. Proper radiator sizing determines the maximum continuous power dissipation a spacecraft can sustain.}{catSS}
\glsadd{Heat-Pipe}\dictentry{Heat Pipe}{A heat pipe is a passive device using evaporation and condensation of a working fluid to transport thermal energy efficiently. The evaporator absorbs heat, vaporizing the fluid, which travels to the condenser to release heat and return to liquid. Capillary wicks or gravity return the liquid to the evaporator. Heat pipes provide high thermal conductivity with no moving parts, making them widely used in spacecraft thermal control.}{catSS}
\glsadd{Louver}\dictentry{Louver}{A louver is a variable-geometry thermal control device adjusting radiator surface emissivity to regulate heat rejection. Movable blades open to expose high-emissivity surfaces or close to reduce heat loss. They provide active thermal control without propellant, using small motors for blade position adjustment based on temperature feedback.}{catSS}
\glsadd{Phase-Change-Material}\dictentry{Phase Change Material}{A phase change material absorbs or releases large amounts of thermal energy during melting or solidification at a relatively constant temperature. Spacecraft use PCMs as thermal buffers smoothing temperature fluctuations from eclipse cycles and varying power loads. As the spacecraft heats up the PCM melts absorbing heat; as it cools the PCM solidifies releasing stored heat. Common space PCMs include paraffin waxes and metallic alloys.}{catSS}
\glsadd{Thermal-Blanket}\dictentry{Thermal Blanket}{A thermal blanket is a removable or permanent insulating cover applied to spacecraft components for temperature control. It may consist of multi-layer insulation, foam, or specialized fabrics for specific thermal properties. Thermal blankets protect sensitive components from extreme temperatures during launch, transit, and orbital operations.}{catSS}
\glsadd{Cold-Plate}\dictentry{Cold Plate}{A cold plate is a metal plate with internal fluid channels that absorbs heat from mounted electronic components. Components are thermally bonded using interface materials that minimize contact resistance. Circulating coolant carries absorbed heat to radiators for space rejection. Cold plates are the primary cooling method for high-power electronics, amplifiers, and computing systems on spacecraft and space stations.}{catSS}
\glsadd{Telecommand}\dictentry{Telecommand}{Telecommand is the process of sending encoded instructions from a ground station to a spacecraft. Commands range from instrument settings and orbital maneuvers to software updates and contingency responses. Commands are encrypted, error-protected, and sequenced for reliable onboard execution. Telecommand links must maintain connectivity at vast interplanetary distances, requiring powerful ground transmitters and sensitive spacecraft receivers.}{catSS}
\glsadd{Ground-Station}\dictentry{Ground Station}{A ground station is a terrestrial facility with antennas and communication systems that transmit and receive data to and from spacecraft. Stations are strategically located worldwide to provide continuous contact with satellites and interplanetary missions. Large stations form deep space networks communicating with spacecraft billions of kilometers away.}{catSS}
\glsadd{High-Gain-Antenna}\dictentry{High Gain Antenna}{A high-gain antenna is a directional antenna concentrating radio energy into a narrow beam for high data rate communication. It typically uses a parabolic reflector or phased array, requiring precise pointing toward Earth or a relay. High-gain antennas enable transmission of large data volumes including imagery and video from spacecraft. The trade-off is the need for accurate attitude control and the limited coverage of the narrow beam.}{catSS}
\glsadd{Low-Gain-Antenna}\dictentry{Low Gain Antenna}{A low-gain antenna is an omnidirectional antenna providing broad communication coverage without precise pointing. It has lower data rates than high-gain antennas but is invaluable during emergencies, launch, and orbit acquisition when orientation may be uncertain. It ensures basic communication regardless of spacecraft attitude. Most spacecraft carry both high-gain and low-gain antennas to balance data throughput with link reliability.}{catSS}
\glsadd{Parabolic-Antenna}\dictentry{Parabolic Antenna}{A parabolic antenna uses a dish-shaped reflector to focus radio waves onto a feed horn, creating a highly directional high-gain beam. Parabolic geometry ensures reflected waves arrive at the focal point in phase, producing constructive interference and a narrow beam. These antennas are standard for deep space communication, radar, and high-capacity satellite links. Gain increases with the dish-diameter-to-wavelength ratio.}{catSS}
\glsadd{Phased-Array}\dictentry{Phased Array}{A phased array is an antenna system that uses multiple radiating elements arranged in a pattern. By controlling the phase of the signal fed to each element, the beam can be steered electronically without physically moving the antenna. This allows rapid scanning and the ability to track multiple targets simultaneously. Phased arrays are widely used in space communications and radar for their speed and lack of mechanical wear.}{catSS}
\glsadd{RF-Communication}\dictentry{RF Communication}{RF Communication refers to the transmission and reception of information using radio frequency electromagnetic waves. Spacecraft use RF links to send telemetry data back to Earth and receive commands from ground controllers. Different frequency bands are chosen based on data rate requirements, antenna size, and atmospheric attenuation. RF communication remains the primary method for maintaining contact with spacecraft across all mission phases.}{catSS}
\glsadd{S-Band}\dictentry{S Band}{The S Band is a portion of the electromagnetic spectrum spanning approximately 2 to 4 GHz. It is commonly used for spacecraft telemetry, tracking, and command links because it offers a good balance between data rate and atmospheric penetration. S Band signals pass through rain and clouds with relatively low attenuation. This band is widely employed for near-Earth and deep-space communication by agencies such as NASA and ESA.}{catSS}
\glsadd{X-Band}\dictentry{X Band}{The X Band covers the frequency range from approximately 8 to 12 GHz. It supports higher data rates than lower frequency bands while still maintaining reasonable atmospheric penetration. X Band is frequently used for deep-space communication, synthetic aperture radar, and military satellite links. Its narrower beamwidth also provides better interference rejection and improved signal quality.}{catSS}
\glsadd{Ka-Band}\dictentry{Ka Band}{The Ka Band spans approximately 26 to 40 GHz and enables very high data rate transmission. It is used for high-bandwidth applications such as video relay, broadband satellite internet, and high-resolution data downlinks. Ka Band signals are more susceptible to rain fade than lower bands but offer significantly greater bandwidth capacity. Advances in technology have made Ka Band increasingly important for modern satellite communication systems.}{catSS}
\glsadd{Deep-Space-Network}\dictentry{Deep Space Network}{The Deep Space Network is NASA's international array of large antenna complexes used for communicating with interplanetary spacecraft. It consists of three facilities roughly 120 degrees apart in longitude, ensuring continuous coverage as Earth rotates. The DSN provides tracking, telemetry, command, and radio science services throughout the solar system. It is one of the most critical ground-based assets for deep-space exploration.}{catSS}
\glsadd{DSN}\dictentry{DSN}{DSN stands for Deep Space Network, NASA's global system of antenna complexes for communicating with spacecraft beyond Earth orbit. The network's three facilities in California, Spain, and Australia provide round-the-clock coverage. The DSN supports missions ranging from lunar orbiters to interstellar probes, enabling data collection and mission control at extreme distances.}{catSS}
\glsadd{Data-Rate}\dictentry{Data Rate}{Data rate is the speed at which information is transmitted from a spacecraft to a ground station, typically measured in bits per second. It depends on factors such as transmitter power, antenna gain, distance, and the frequency band used. Higher data rates allow more scientific data and higher-resolution imagery to be returned. Data rate is a key design parameter that influences spacecraft power systems and communication architecture.}{catSS}
\glsadd{Compression}\dictentry{Compression}{Compression is the process of reducing data size before transmission to conserve bandwidth and storage. Spacecraft use both lossless and lossy algorithms depending on the data type. Lossless methods preserve all original information, while lossy methods achieve higher ratios by discarding less critical data. Effective compression maximizes the scientific return of data-limited space missions.}{catSS}
\glsadd{Error-Correction}\dictentry{Error Correction}{Error correction is a set of techniques used to detect and recover from data corruption during transmission. Space communication channels are subject to noise, interference, and signal degradation over vast distances. Error correction codes add redundant bits to the transmitted data, allowing the receiver to reconstruct the original message. These codes are vital for ensuring reliable communication with spacecraft across the solar system.}{catSS}
\glsadd{Reed-Solomon-Code}\dictentry{Reed-Solomon Code}{A Reed-Solomon code is a powerful error correction scheme that operates on blocks of data symbols. It can correct multiple symbol errors within a block, making it particularly effective at handling burst errors. Reed-Solomon codes are widely used in deep-space communication, satellite broadcasting, and data storage systems. They form part of the standard coding architecture for many NASA and ESA missions.}{catSS}
\glsadd{Convolutional-Code}\dictentry{Convolutional Code}{A convolutional code is an error correction code that processes data as a continuous stream rather than in fixed blocks. It uses shift registers and modulo arithmetic to generate redundant bits from the input. Convolutional codes are often decoded using the Viterbi algorithm to find the most likely transmitted sequence. They are commonly paired with Reed-Solomon codes in concatenated schemes for space missions.}{catSS}
\glsadd{Turbo-Code}\dictentry{Turbo Code}{A turbo code is a high-performance error correction code approaching the theoretical Shannon limit for channel capacity. It uses two or more component codes connected through an interleaver with iterative decoding. Turbo codes achieve excellent performance at low signal-to-noise ratios. They are used in deep-space missions and modern satellite communication standards.}{catSS}
\glsadd{LDPC}\dictentry{LDPC}{LDPC stands for Low-Density Parity-Check code, a class of linear error correction codes with sparse parity-check matrices. LDPC codes can achieve performance very close to the Shannon limit through iterative belief-propagation decoding. They are used in modern space communication standards, including CCSDS recommendations for deep-space and near-Earth links. LDPC codes offer a flexible trade-off between complexity and error correction capability.}{catSS}
\glsadd{Structure}\dictentry{Structure}{The structure is the mechanical framework that forms the skeleton of a spacecraft. It must withstand launch loads, thermal stresses, and operational forces while maintaining precise alignment of all components. Spacecraft structures are designed to be as lightweight as possible while meeting strict strength and stiffness requirements. Materials such as aluminum alloys, titanium, and carbon fiber composites are commonly used in their construction.}{catSS}
\glsadd{Aluminum-Honeycomb}\dictentry{Aluminum Honeycomb}{Aluminum honeycomb is a core material made from thin aluminum foil formed into a hexagonal cell pattern. When sandwiched between face sheets, it creates panels with an outstanding strength-to-weight ratio. Aluminum honeycomb is widely used in spacecraft structures, launch vehicle adapters, and satellite panels. It offers good thermal stability, crash energy absorption, and is compatible with standard bonding techniques.}{catSS}
\glsadd{Carbon-Fiber-Structure}\dictentry{Carbon Fiber Structure}{A carbon fiber structure is a spacecraft component made from carbon fiber reinforced polymer composite material. Carbon fiber composites offer very high specific strength and stiffness, making them ideal for weight-critical space applications. They are used in antenna reflectors, solar array structures, and primary bus structures. The material also provides excellent thermal stability, which is important for precision-pointing instruments.}{catSS}
\glsadd{Deployable-Structure}\dictentry{Deployable Structure}{A deployable structure is a mechanical system that unfolds from a compact stowed configuration to its full operational shape in space. Deployment is necessary because many components are too large for launch vehicle fairings in their deployed state. Reliable mechanisms including hinges, latches, and actuators are critical to mission success. Examples include solar arrays, antennas, booms, and sunshields.}{catSS}
\glsadd{Solar-Array-Deployment}\dictentry{Solar Array Deployment}{Solar array deployment is the process of unfolding solar panel assemblies from their stowed configuration to their full extended position in space. This typically occurs shortly after spacecraft separation from the launch vehicle. The deployment mechanism must function reliably in the vacuum and thermal environment of space. Successful deployment is essential for providing the electrical power needed for all spacecraft operations.}{catSS}
\glsadd{Antenna-Deployment}\dictentry{Antenna Deployment}{Antenna deployment is the process of unfurling a communication or sensing antenna to its operational geometry in space. Many high-gain antennas must stow compactly and deploy reliably since they are too large for the launch fairing when fully extended. The process may involve mechanical hinges, inflatable structures, or shape-memory materials. Accurate deployment is critical for establishing communication links.}{catSS}
\glsadd{Boom}\dictentry{Boom}{A boom is a long, slender deployable arm or mast extending from the spacecraft body. Booms are used to separate sensitive instruments from the electromagnetic interference of the main bus, to position sensors, or to provide structural support for appendages. They may be rigid, articulated, or inflatable in design. Booms play an important role in missions requiring precise instrument placement or gravity-gradient stabilization.}{catSS}
\glsadd{Mast}\dictentry{Mast}{A mast is a deployable vertical structure extending from a spacecraft to elevate instruments, antennas, or other equipment above the main body. Masts are commonly used for magnetometers, antennas, and solar wind sensors to reduce interference from the spacecraft. They can be rigid telescoping structures or lattice-type booms. Mast deployment is a critical event that must be carefully timed and verified during mission operations.}{catSS}
\glsadd{Mechanism}\dictentry{Mechanism}{A mechanism is any moving or articulating component on a spacecraft designed to perform a specific mechanical function. Mechanisms include hinges, latches, actuators, gears, and deployment devices. They must operate reliably in the harsh space environment of vacuum, extreme temperatures, and radiation with no opportunity for repair. Mechanism design and testing are among the most demanding aspects of spacecraft engineering.}{catSS}
\glsadd{Pyrotechnic-Device}\dictentry{Pyrotechnic Device}{A pyrotechnic device is a small explosive component used on spacecraft to produce controlled mechanical action. These devices are used for separation, deployment, and release operations where a single, reliable, high-force event is required. Examples include bolt cutters, nut crackers, and linear shaped charges. Pyrotechnic devices are valued for their high reliability, fast response, and simplicity, though they are one-shot devices.}{catSS}
\glsadd{Frangible-Nut}\dictentry{Frangible Nut}{A frangible nut is a pyrotechnic fastener designed to fracture and release when an explosive charge is initiated. It is commonly used to release hold-down mechanisms, separate structural joints, or initiate deployment sequences. The nut splits apart along predetermined fracture lines, allowing the bolt or stud to pass through freely. Frangible nuts are widely used in spacecraft and launch vehicle separation systems.}{catSS}
\glsadd{Cable-Cutter}\dictentry{Cable Cutter}{A cable cutter is a pyrotechnic device that uses a shaped charge or blade to sever a cable or line on command. It is used to release tensioned cables, release mechanisms, or cut tie-down straps during deployment operations. The cutting action is nearly instantaneous and highly reliable. Cable cutters are commonly employed in spacecraft separation and deployment systems where a clean, definitive release is required.}{catSS}
\glsadd{Pin-Puller}\dictentry{Pin Puller}{A pin puller is a pyrotechnic actuator that retracts a locking pin from a mechanism on command. Once the pin is withdrawn, the mechanism it was restraining is free to move, such as a deployment hinge or separation joint. Pin pullers are compact, reliable, and provide a positive release action. They are used extensively in spacecraft deployment and release mechanisms where precise timing and high reliability are essential.}{catSS}
\glsadd{Separation-Nut}\dictentry{Separation Nut}{A separation nut is a pyrotechnic device that releases a bolted connection when initiated. It typically consists of a nut with a weakened or split design that fractures under the force of an explosive charge. Separation nuts are used for stage separation, payload release, and structural disconnection events. They provide a clean, high-force separation with no residual connection between the separated bodies.}{catSS}
\glsadd{Launch-Adapter}\dictentry{Launch Adapter}{A launch adapter is the mechanical interface structure that connects a spacecraft to the upper stage of its launch vehicle. It must transmit all launch loads from the rocket to the spacecraft while providing a clean separation interface. The adapter typically includes separation mechanisms, electrical connectors, and hold-down points. Its design is critical to ensuring safe and reliable spacecraft deployment into orbit.}{catSS}
\glsadd{Payload-Adapter}\dictentry{Payload Adapter}{A payload adapter is a structure that attaches a satellite to the launch vehicle's payload attach fitting. It serves as the mechanical, electrical, and thermal interface between the payload and rocket. The adapter supports the payload during launch and facilitates clean separation once orbit is reached. Standardized interfaces allow different payloads to be accommodated on various launch vehicles.}{catSS}
\glsadd{Clamp-Band}\dictentry{Clamp Band}{A clamp band is a tensioned metal band that holds a spacecraft or upper stage securely to its separation interface. When released, the band loosens and allows the two bodies to separate cleanly. Clamp bands are simple, reliable, and have been used in hundreds of space missions. They provide a low-shock separation compared to some alternative systems, making them suitable for sensitive payloads.}{catSS}
\glsadd{Marman-Clamp}\dictentry{Marman Clamp}{A Marman clamp is a specific type of clamp band separation system widely used in the space industry. It consists of a V-shaped band that wraps around a flange on the separation interface, held in tension by bolts. When the bolts are released pyrotechnically, the band opens and the spacecraft separates from the launch vehicle. Marman clamps have an extensive flight heritage and are known for their reliability and relatively low separation shock.}{catSS}
\glsadd{Lightband}\dictentry{Lightband}{A Lightband is a low-shock spacecraft separation system developed by Planetary Systems Corporation. It uses a motorized rotary actuator to release a clamp band, resulting in very low separation shock levels. The low-shock characteristic makes Lightband ideal for sensitive payloads such as large antennas and optical instruments. It also allows the separation system to be reused, which is beneficial for servicing missions.}{catSS}
\glsadd{Separation-System}\dictentry{Separation System}{A separation system is the collection of mechanisms that physically release a spacecraft from its launch vehicle or another spacecraft. It must provide clean, controlled separation with minimal disturbance to attitude and trajectory. Systems may use pyrotechnic, mechanical, or electromechanical actuators. The separation system design is critical to mission success, as failure to separate would end the mission.}{catSS}
\glsadd{Yo-Yo-Deploy}\dictentry{Yo-Yo Deploy}{A yo-yo deploy is a technique used to despin a spinning spacecraft before or after separation. It consists of small masses on wires released from the spinning body. As the masses extend outward, angular momentum is transferred to the wires, reducing the spacecraft's spin rate. Once the desired spin rate is achieved, the masses and wires are released. This simple mechanical method is an effective and reliable way to control spacecraft spin.}{catSS}
\glsadd{Spin-Table}\dictentry{Spin Table}{A spin table is a rotating platform used to spin up spacecraft or payloads prior to launch or deployment. The payload is mounted on the table and spun to the required rate to impart gyroscopic stability. Spin tables are used during ground testing and sometimes during the launch sequence itself. They ensure uniform and precise spin rates for missions that require spin-stabilized configurations.}{catSS}
\glsadd{Center-of-Mass}\dictentry{Center of Mass}{The center of mass is the point where the entire mass of a spacecraft can be considered concentrated for motion analysis. It is the balance point where the spacecraft would remain stationary if supported. The center of mass must be carefully controlled throughout the mission, as shifts from propellant consumption or deployment affect stability. Accurate knowledge of its location is essential for attitude control and trajectory.}{catSS}
\glsadd{Margin-of-Safety}\dictentry{Margin of Safety}{Margin of safety is a factor applied in structural design to ensure a component can withstand loads beyond maximum expected levels. It is expressed as the ratio of failure load to maximum expected load, minus one. A positive margin indicates the structure can handle loads beyond those anticipated. Margins are calculated for all critical structural elements and verified during qualification testing.}{catSS}
\glsadd{Qualification-Model}\dictentry{Qualification Model}{A qualification model is a spacecraft built to the same design as the flight unit but intended solely for ground testing. It is subjected to loads, thermal, vibration, and environmental tests at levels exceeding expected flight conditions. The purpose is to prove the design is robust and capable of surviving the mission environment. Results give confidence that the flight hardware will perform as intended.}{catSS}
\glsadd{Flight-Model}\dictentry{Flight Model}{The flight model is the actual spacecraft or hardware unit that will be launched and operated in space. It is built to the final, qualified design and undergoes acceptance testing to verify workmanship and functionality. The flight model is the culmination of all design, development, and testing efforts. Once it reaches orbit, its performance determines the success of the mission.}{catSS}
\glsadd{Engineering-Model}\dictentry{Engineering Model}{An engineering model is a spacecraft built to the same functional design as the flight unit but used for ground-based development and testing. It allows engineers to verify electrical interfaces, software, and thermal designs before committing to flight hardware. Engineering models are less expensive and can be modified more readily. They play a critical role in identifying and resolving design issues early.}{catSS}
\glsadd{Structural-Model}\dictentry{Structural Model}{A structural model is a spacecraft component built to replicate the mechanical properties of the flight design. It is subjected to static load, vibration, acoustic, and shock testing to verify structural integrity. The model may not include functional electronics, focusing instead on material properties and joint strengths. Test results validate analytical models and confirm the structure meets design requirements.}{catSS}
\glsadd{Thermal-Model}\dictentry{Thermal Model}{A thermal model is a representation of a spacecraft used to verify its thermal control design under simulated space conditions. It may include flight-representative hardware or simplified thermal simulators. The model is tested in thermal vacuum chambers where temperature extremes and vacuum conditions are recreated. Results confirm that all components will remain within their operational temperature ranges throughout the mission.}{catSS}
\glsadd{Protoflight-Model}\dictentry{Protoflight Model}{A protoflight model is a spacecraft unit that combines the roles of both qualification and flight hardware. It undergoes a reduced-level qualification test campaign followed by acceptance testing, and is then flown on the actual mission. This approach saves cost and schedule by building only one unit instead of two. Protoflight models are common in smaller missions and programs with limited budgets.}{catSS}
\glsadd{Clean-Room}\dictentry{Clean Room}{A clean room is a controlled environment where airborne particulate contamination is strictly regulated. Spacecraft assembly and integration are performed in clean rooms to prevent dust and debris from contaminating sensitive hardware. Classifications specify the maximum allowable number and size of particles per unit volume. Maintaining cleanliness is essential for protecting optical instruments and electronics.}{catSS}
\glsadd{Cleanliness}\dictentry{Cleanliness}{Cleanliness refers to the control of particulate and molecular contamination on spacecraft hardware. Contamination can degrade optical surfaces, contaminate scientific instruments, and cause electrical failures. Cleanliness levels are specified for different areas of the spacecraft based on sensitivity. Procedures include cleaning, protective covers, and ongoing monitoring throughout assembly, test, and launch operations.}{catSS}
\glsadd{Outgassing}\dictentry{Outgassing}{Outgassing is the release of volatile gases from materials when exposed to the vacuum of space. These gases can condense on cold surfaces such as optics and thermal control coatings, degrading their performance. Materials used in space are screened for outgassing properties using standards such as NASA ASTM E595. Low-outgassing materials are required for most spacecraft applications to ensure long-term reliability.}{catSS}
\glsadd{Vacuum}\dictentry{Vacuum}{Vacuum refers to a space with significantly reduced atmospheric pressure, approaching the near-total absence of gas molecules. In low Earth orbit, the vacuum is extremely thin, and in deep space it is even more so. The space vacuum presents challenges such as outgassing, arcing, and the inability to conduct heat by convection. Understanding and accounting for vacuum conditions is fundamental to spacecraft design and testing.}{catSS}
\glsadd{Acoustic-Test}\dictentry{Acoustic Test}{An acoustic test simulates the intense sound pressure levels generated by rocket engines during launch. Large spacecraft and lightweight structures are particularly susceptible to acoustic loads, which can cause vibration and fatigue. The test article is placed in a reverberant chamber and exposed to high-intensity broadband noise. Acoustic testing ensures that sensitive components and structures can withstand the acoustic environment of launch.}{catSS}
\glsadd{Shock-Test}\dictentry{Shock Test}{A shock test simulates the transient mechanical shock produced by pyrotechnic events such as stage separation and deployment. These shock events have high peak accelerations and short durations, which can damage sensitive components. The test verifies that electronics, mechanisms, and structures can survive these sudden impulses. Shock testing is a standard part of the environmental test campaign for all spacecraft hardware.}{catSS}
\glsadd{EMIEMC-Test}\dictentry{EMI/EMC Test}{EMI/EMC testing verifies that a spacecraft's electronics can operate correctly in electromagnetic interference and do not emit excessive radiation. Electromagnetic compatibility ensures different subsystems do not interfere with each other. Tests include radiated and conducted emissions measurements, plus susceptibility testing. EMC compliance is mandatory for all spacecraft to ensure reliable mission operations.}{catSS}
\glsadd{Mass-Properties-Measurement}\dictentry{Mass Properties Measurement}{Mass properties measurement is the process of determining a spacecraft's mass, center of mass location, and moments of inertia. These parameters are essential for attitude control design, trajectory analysis, and structural load calculations. Measurements are taken using precision scales, pendulum methods, and spin tables. Accurate mass property data is required before any flight operation, including launch vehicle integration.}{catSS}
\glsadd{Spin-Balance}\dictentry{Spin Balance}{Spin balance is the process of adjusting the mass distribution of a spinning spacecraft so that its spin axis coincides with its principal axis of inertia. An unbalanced spinning body will exhibit nutation, which can degrade attitude control performance. Balance weights are added or relocated to minimize the transverse moment of inertia. Spin balance is a critical step in preparing a spin-stabilized spacecraft for flight.}{catSS}
\glsadd{Alignment}\dictentry{Alignment}{Alignment is the process of precisely orienting spacecraft instruments and structural axes relative to a reference coordinate system. Optical tools such as theodolites and autocollimators are used to measure and adjust pointing directions. Accurate alignment ensures instruments point at intended targets and sensors are correctly referenced. Misalignment can lead to pointing errors that degrade mission performance.}{catSS}
\glsadd{Integration}\dictentry{Integration}{Integration is the process of assembling and interconnecting all spacecraft subsystems into a complete vehicle. It involves mechanical assembly, electrical wiring, software loading, and functional verification of interfaces. Integration is performed in a clean room following detailed procedures and workmanship standards. This phase transforms individually tested components into a fully operational spacecraft.}{catSS}
\glsadd{Test}\dictentry{Test}{Test is the comprehensive verification process ensuring a spacecraft and its subsystems meet mission requirements. It includes functional checks, environmental tests such as vibration and thermal vacuum, and end-to-end system tests. Each phase has specific acceptance criteria that must be met before proceeding. Thorough testing is the primary means of reducing risk and ensuring mission success.}{catSS}
\glsadd{Launch-Campaign}\dictentry{Launch Campaign}{A launch campaign encompasses all activities from the spacecraft's arrival at the launch site through liftoff. This includes final assembly, integration with the launch vehicle, propellant loading, and the countdown sequence. The campaign involves coordination between the spacecraft team, launch vehicle provider, and range operations. Meticulous planning and execution of hundreds of procedures is required.}{catSS}
\glsadd{Launch-Site}\dictentry{Launch Site}{A launch site is a ground facility specifically designed and equipped for the preparation and launch of spacecraft and rockets. It includes launch pads, integration buildings, propellant storage, tracking systems, and range safety infrastructure. Launch sites are typically located near coastlines to allow trajectories over open ocean for safety. The choice of launch site depends on orbital requirements, logistics, and political considerations.}{catSS}
\glsadd{Cape-Canaveral}\dictentry{Cape Canaveral}{Cape Canaveral is a major United States launch site located on the Atlantic coast of Florida. It is home to multiple launch complexes used by both government and commercial launch providers. Cape Canaveral Space Force Station operates alongside the adjacent Kennedy Space Center. Its equatorial-advantage latitude and easterly launch azimuth over the Atlantic make it ideal for a wide range of orbital missions.}{catSS}
\glsadd{Kennedy-Space-Center}\dictentry{Kennedy Space Center}{Kennedy Space Center is NASA's primary launch facility, located on Merritt Island, Florida. It has been the departure point for all human spaceflight missions from the United States, including the Apollo Moon landings and Space Shuttle launches. The center includes the Vehicle Assembly Building, launch pads 39A and 39B, and extensive processing facilities. It remains a central hub for NASA and commercial crew and cargo missions.}{catSS}
\glsadd{Vandenberg-SFB}\dictentry{Vandenberg SFB}{Vandenberg Space Force Base is a United States launch facility located on the central coast of California. It is the primary site for launches into polar and sun-synchronous orbits, which are difficult to achieve from east-coast sites due to population safety concerns. The base operates multiple launch complexes for both government and commercial missions. Its southerly trajectory over the Pacific Ocean makes it ideal for high-inclination orbits.}{catSS}
\glsadd{Wallops-Flight-Facility}\dictentry{Wallops Flight Facility}{Wallops Flight Facility is a NASA launch site located on the Eastern Shore of Virginia. It supports suborbital rocket launches, small orbital missions, and scientific balloon flights. The facility also serves as a range for aerospace research and technology development. Its location and smaller infrastructure make it well-suited for research missions and smaller launch vehicles.}{catSS}
\glsadd{Baikonur}\dictentry{Baikonur}{Baikonur is a Russian space launch facility located in the desert steppe of Kazakhstan. It is the world's first and one of the most active space launch sites, having been used since the beginning of the space age. Baikonur was the launch site for Sputnik, Yuri Gagarin's historic flight, and continues to support Soyuz and Proton launches. The facility is leased by Russia and remains a critical element of international space access.}{catSS}
\glsadd{Kourou}\dictentry{Kourou}{Kourou is a European Space Agency launch site located in French Guiana, South America. Its proximity to the equator provides a significant performance advantage for geostationary orbit launches. The Guiana Space Centre operates Ariane, Vega, and Soyuz launch vehicles from this facility. Kourou's location over the Atlantic Ocean ensures safe launch trajectories for a wide variety of orbital missions.}{catSS}
\glsadd{Spaceport}\dictentry{Spaceport}{A spaceport is a facility designed for the launch and recovery of spacecraft and space vehicles. It includes launch pads, processing buildings, control centers, and support infrastructure. Modern spaceports serve both government and commercial customers, supporting a growing range of launch vehicle types. The term has expanded to include facilities for suborbital tourism and reusable vehicle operations.}{catSS}
\glsadd{Launch-Complex}\dictentry{Launch Complex}{A launch complex is a complete launch facility that includes one or more launch pads along with associated support buildings, fuel storage, and control centers. Each complex is typically designed to support specific launch vehicles or families of vehicles. The complex provides all the ground infrastructure needed to prepare, fuel, and launch a rocket. Major spaceports contain multiple launch complexes to support concurrent operations.}{catSS}
\glsadd{Integration-Facility}\dictentry{Integration Facility}{An integration facility is a specialized building where spacecraft and launch vehicles are assembled and connected. It provides a clean room environment for mating the payload with the rocket and performing final checkouts. Facilities include overhead cranes, environmental control systems, and access platforms. The facility ensures all interfaces are properly connected before transport to the launch pad.}{catSS}
\glsadd{Payload-Processing}\dictentry{Payload Processing}{Payload processing is the series of activities performed to prepare a spacecraft for integration with its launch vehicle. This includes final inspections, fuel loading for bi-propellant systems, battery charging, and protective cover installation. Processing is performed in clean room facilities with strict contamination control. The goal is to deliver a flight-ready payload that can be safely mated with the rocket and launched.}{catSS}
\glsadd{Propellant-Loading}\dictentry{Propellant Loading}{Propellant loading is the process of filling a spacecraft or launch vehicle stage with its fuel and oxidizer. This is one of the most hazardous operations in a launch campaign due to the toxic, corrosive, or cryogenic nature of many propellants. Loading procedures follow strict safety protocols with multiple verification steps. Propellant loading typically occurs during the final hours before launch as part of the countdown sequence.}{catSS}
\glsadd{Big-Bang}\dictentry{Big Bang}{The Big Bang is the cosmological event marking the beginning of the universe approximately 13.8 billion years ago. It represents the rapid expansion of space-time from an extremely hot, dense state. As the universe expanded and cooled, matter condensed into the fundamental particles and forces we observe today. Evidence includes the observed expansion of the universe, cosmic microwave background radiation, and light element abundance.}{catSS}
\glsadd{Dark-Matter}\dictentry{Dark Matter}{Dark matter is a form of matter that does not emit, absorb, or reflect electromagnetic radiation, making it invisible to direct observation. Its existence is inferred from gravitational effects on visible matter, such as galaxy rotation curves and gravitational lensing. Dark matter constitutes about 27 percent of the universe's total mass-energy content. Understanding its nature is one of the most important open questions in physics.}{catSS}
\glsadd{Dark-Energy}\dictentry{Dark Energy}{Dark energy is a hypothetical form of energy that permeates all of space and is responsible for the observed accelerated expansion of the universe. It is thought to constitute roughly 68 percent of the total mass-energy content of the universe. Unlike dark matter, dark energy appears to be uniformly distributed and does not clump under gravity. Its fundamental nature remains one of the greatest mysteries in cosmology and astrophysics.}{catSS}
\glsadd{Black-Hole}\dictentry{Black Hole}{A black hole is a region of spacetime where gravity is so intense that nothing, not even light, can escape from within its event horizon. Black holes form when massive stars collapse at the end of their lives or through other extreme gravitational processes. They are characterized by their mass, spin, and electric charge. Supermassive black holes reside at the centers of most galaxies and play a key role in galactic evolution.}{catSS}
\glsadd{Neutron-Star}\dictentry{Neutron Star}{A neutron star is an extremely dense stellar remnant formed from the core collapse of a massive star during a supernova. It packs more mass than the Sun into a sphere roughly 20 kilometers in diameter. Its matter is primarily composed of neutrons compressed to nuclear density. Neutron stars exhibit extraordinary properties including rapid rotation, intense magnetic fields, and radiation emission as pulsars.}{catSS}
\glsadd{White-Dwarf}\dictentry{White Dwarf}{A white dwarf is the dense, compact remnant left behind after a low- to medium-mass star has exhausted its nuclear fuel and shed its outer layers. A white dwarf contains roughly the mass of the Sun compressed into a volume the size of Earth. It is supported against further gravitational collapse by electron degeneracy pressure. White dwarfs slowly cool and fade over billions of years, eventually becoming black dwarfs.}{catSS}
\glsadd{Supernova}\dictentry{Supernova}{A supernova is a catastrophic stellar explosion occurring at the end of a massive star's life or in a white dwarf binary system. The star can briefly outshine its entire host galaxy and release enormous energy. The explosion disperses heavy elements into the interstellar medium, enriching material from which new stars and planets form. Supernovae are critical to the chemical evolution of the universe.}{catSS}
\glsadd{Nebula}\dictentry{Nebula}{A nebula is a vast cloud of interstellar gas and dust found throughout galaxies. Nebulae can serve as star-forming regions where gravity compresses gas to ignite new stars, or as remnants of dead stars such as planetary nebulae and supernova remnants. They are composed primarily of hydrogen and helium with trace heavier elements. Nebulae are among the most visually spectacular objects in astronomy.}{catSS}
\glsadd{Galaxy}\dictentry{Galaxy}{A galaxy is a massive, gravitationally bound system consisting of stars, stellar remnants, gas, dust, and dark matter. Galaxies range from small dwarfs with a few billion stars to giants containing trillions of stars. They are classified by shape into types such as spiral, elliptical, and irregular. Our Sun is one of hundreds of billions of stars in the Milky Way galaxy.}{catSS}
\glsadd{Milky-Way}\dictentry{Milky Way}{The Milky Way is the barred spiral galaxy that contains our Solar System. It spans roughly 100,000 light-years in diameter and contains an estimated 100 to 400 billion stars. The Solar System is located in one of the Milky Way's spiral arms, approximately 26,000 light-years from the galactic center. The galaxy is named for the hazy band of light visible in the night sky, formed by the combined glow of its distant stars.}{catSS}
\glsadd{Light-Year}\dictentry{Light-Year}{A light-year is the distance that light travels in a vacuum in one Julian year, approximately 9.46 trillion kilometers or 5.88 trillion miles. It is a standard unit of measurement used to express the vast distances between objects in space. For example, the nearest star to the Sun, Proxima Centauri, is about 4.24 light-years away. The light-year provides a convenient scale for discussing the immense reaches of the cosmos.}{catSS}
\glsadd{Astronomical-Unit}\dictentry{Astronomical Unit}{An astronomical unit is the mean distance from the center of the Earth to the center of the Sun, defined as exactly 149,597,870.7 kilometers. It is commonly used as a convenient unit for measuring distances within the Solar System. For example, Jupiter orbits the Sun at about 5.2 astronomical units. The astronomical unit provides a practical scale for expressing the relative distances of planets and other Solar System bodies.}{catSS}
\glsadd{Parsec}\dictentry{Parsec}{A parsec is a unit of distance equal to approximately 3.26 light-years or 30.9 trillion kilometers. It is defined as the distance at which one astronomical unit subtends an angle of one arcsecond. The parsec is widely used in professional astronomy because it relates directly to stellar parallax measurements. It provides a practical scale for expressing interstellar and intergalactic distances.}{catSS}
\glsadd{Redshift}\dictentry{Redshift}{Redshift is the phenomenon in which the wavelength of light from an object is increased, shifting toward the red end of the spectrum. It occurs when a source is moving away from the observer or when the universe is expanding. Cosmological redshift is key evidence for the expansion of the universe. Measuring redshift of distant galaxies allows estimation of their distance and recession velocity.}{catSS}
\glsadd{Cosmic-Microwave-Background}\dictentry{Cosmic Microwave Background}{The cosmic microwave background is residual thermal radiation from the Big Bang permeating all of space. It was released about 380,000 years after the Big Bang when the universe cooled enough for atoms to form. The CMB is observed as a nearly uniform microwave signal at about 2.7 Kelvin. Small fluctuations provide crucial information about early universe density variations and the seeds of galaxy formation.}{catSS}
\glsadd{Terrestrial-Planet}\dictentry{Terrestrial Planet}{A terrestrial planet is a rocky planet composed primarily of silicate rocks and metals, with a solid surface. The four inner planets of our Solar System, Mercury, Venus, Earth, and Mars, are terrestrial planets. They are generally smaller and denser than gas giants, with iron cores and rocky mantles. Terrestrial planets are the primary targets in the search for potentially habitable worlds beyond our Solar System.}{catSS}
\glsadd{Gas-Giant}\dictentry{Gas Giant}{A gas giant is a large planet composed mostly of hydrogen and helium, with no well-defined solid surface. In our Solar System, Jupiter and Saturn are the gas giants, while Uranus and Neptune are sometimes classified separately as ice giants. Gas giants have deep atmospheres, strong magnetic fields, and extensive moon systems. They play an important role in shaping the architecture of planetary systems through their gravitational influence.}{catSS}
\glsadd{Impact-Crater}\dictentry{Impact Crater}{An impact crater is a bowl-shaped depression on a celestial body formed by the collision of a meteoroid, asteroid, or comet. Crater size depends on the impactor's mass, velocity, angle, and the target surface properties. Impact craters are the most common geological feature on many Solar System bodies. Studying them provides insights into the impact history and surface age of planetary bodies.}{catSS}
\glsadd{Regolith}\dictentry{Regolith}{Regolith is the layer of loose, unconsolidated material covering the solid bedrock of a celestial body. Lunar regolith consists of fine dust, rock fragments, and glassy particles produced by billions of years of micrometeorite impacts. Mars, asteroids, and other bodies also have regolith of varying thickness and composition. Understanding regolith properties is essential for surface operations and resource utilization.}{catSS}
\glsadd{Cubesat}\dictentry{Cubesat}{A CubeSat is a standardized small satellite format based on a 10-centimeter unit side with a mass of approximately 1.3 kilograms per unit. They use commercial off-the-shelf components, reducing development time and cost. CubeSats are deployed from a standard dispenser attached to a launch vehicle. They have democratized access to space, enabling universities, small companies, and developing nations to conduct orbital missions.}{catSS}
\glsadd{SmallSat}\dictentry{SmallSat}{A SmallSat is a category of satellites with reduced mass and size compared to traditional spacecraft, typically ranging from 1 to 500 kilograms. SmallSats use miniaturized components and modern manufacturing to deliver capable missions at lower cost. They are used for Earth observation, communications, scientific research, and technology demonstration. The growth of SmallSat missions has transformed the space industry.}{catSS}
\glsadd{Space-Tourism}\dictentry{Space Tourism}{Space tourism refers to commercial spaceflight for recreational, leisure, or business purposes rather than scientific research or government missions. Companies offer suborbital flights, orbital stays, and missions to the International Space Station. The industry is driven by advances in reusable launch vehicle technology and growing private investment. Space tourism represents a significant expansion of the space economy.}{catSS}
\glsadd{Planetary-Protection}\dictentry{Planetary Protection}{Planetary protection is the practice of preventing biological contamination of celestial bodies during space missions. It involves strict protocols for cleaning spacecraft and managing operations to avoid contaminating target worlds with Earth organisms. Forward contamination protects the integrity of extraterrestrial life searches, while back contamination protects Earth. Policies are established internationally through COSPAR.}{catSS}
\glsadd{Eclipse}\dictentry{Eclipse}{An eclipse occurs when one celestial body passes into the shadow of another, blocking its light. A solar eclipse happens when the Moon intercepts sunlight from reaching Earth, while a lunar eclipse occurs when Earth's shadow falls on the Moon. Total eclipses reveal the Sun's corona, and partial eclipses produce varying dimming. Eclipses follow predictable saros cycles and historically validated celestial mechanics models.}{catSS}
\glsadd{Ephemeris}\dictentry{Ephemeris}{An ephemeris is a table of calculated celestial positions at regular time intervals. It uses gravitational models and orbital mechanics to predict where planets, moons, or spacecraft will be at future dates. Astronomers rely on ephemerides for telescope pointing and observation scheduling, while mission planners use them for navigation. Modern ephemerides are continuously refined with new tracking data.}{catSS}
\glsadd{Oberth-Effect}\dictentry{Oberth Effect}{The Oberth Effect describes the efficiency gained by firing a rocket engine at high orbital velocity, typically near periapsis. Because kinetic energy scales with velocity squared, a burn at speed converts propulsive energy more effectively than at low velocity. This principle is central to gravity-assist trajectory design and orbital transfers. Mission planners time burns deep within gravity wells to maximize velocity change.}{catSS}
\glsadd{Orbital-Debris}\dictentry{Orbital Debris}{Orbital debris consists of defunct satellites, spent rocket stages, and collision fragments orbiting Earth at over seven kilometers per second. Even small pieces pose serious collision risks to operational spacecraft. The Kessler Syndrome describes a cascade where collisions generate ever more debris, potentially making certain orbits unusable. Agencies track debris with radar networks and perform avoidance maneuvers.}{catSS}
\glsadd{Rendezvous}\dictentry{Rendezvous}{Rendezvous brings two spacecraft to the same orbital position at the same time for docking or capture. It requires precise phasing burns and relative navigation using radar or optical sensors. The technique enables crewed station visits, satellite servicing, and assembly of large orbital structures. The first successful space rendezvous occurred during NASA's Gemini program in 1965.}{catSS}
\glsadd{Tsiolkovsky-Rocket-Equation}\dictentry{Tsiolkovsky Rocket Equation}{The Tsiolkovsky Rocket Equation relates velocity change to exhaust velocity and the ratio of initial to final mass. It shows delta-v depends logarithmically on mass ratio, meaning exponentially more propellant is needed for each speed increment. This relationship governs all propulsion design and explains the necessity of staging. Konstantin Tsiolkovsky derived it in 1897.}{catSS}
\glsadd{Aberration}\dictentry{Aberration}{Aberration is the apparent displacement of a celestial object's position caused by the observer's motion relative to light speed. Annual aberration from Earth's orbital velocity shifts star positions by up to 20.5 arcseconds. The effect must be corrected in precision astrometry to obtain true coordinates. James Bradley first measured it in 1729, providing early evidence of light's finite speed.}{catSS}
\glsadd{Astrobiology}\dictentry{Astrobiology}{Astrobiology studies the origin, evolution, and distribution of life in the universe. It investigates how life arises from non-living matter and what conditions support biological processes beyond Earth. Researchers examine extremophiles in harsh terrestrial environments and analyze planetary chemistry. The field combines biology, chemistry, geology, and planetary science to address fundamental questions about life in the cosmos.}{catSS}
\glsadd{Asteroid-Mining}\dictentry{Asteroid Mining}{Asteroid mining extracts minerals and resources from asteroids for space manufacturing or use on Earth. Many asteroids contain valuable metals like platinum and nickel, plus water ice for propellant and life support. Economic viability depends on reducing launch costs and developing autonomous microgravity extraction. Several companies and agencies target near-Earth asteroids as accessible early resources.}{catSS}
\glsadd{Atmospheric-Entry}\dictentry{Atmospheric Entry}{Atmospheric entry is a spacecraft's transition from space into a planetary atmosphere, generating extreme aerodynamic heating from hypersonic compression. Ablative or ceramic heat shields protect vehicles by absorbing or radiating thermal energy. Precise trajectory control manages deceleration forces and heating rates. The entry, descent, and landing sequence is critical for Earth return and Mars or Venus missions.}{catSS}
\glsadd{Colony}\dictentry{Colony}{A space colony is a permanent human settlement beyond Earth supporting self-sustaining habitation. Concepts range from rotating orbital habitats to Moon or Mars surface bases, all requiring closed-loop life support, radiation shielding, and in-situ resource use. Psychological and biological isolation challenges make human factors as critical as engineering. Establishing colonies is considered essential for humanity's long-term survival.}{catSS}
\glsadd{Conjunction}\dictentry{Conjunction}{In orbital mechanics, conjunction is the closest approach of two orbiting bodies. For spacecraft, conjunction events trigger collision avoidance maneuvers in crowded regimes. A Mars conjunction occurs when the Sun lies between Earth and Mars, temporarily blocking communications. Predicting conjunctions requires accurate orbital determination and propagation for all objects involved.}{catSS}
\glsadd{Delta-V-Budget}\dictentry{Delta-V Budget}{A delta-v budget is the total velocity change available to a spacecraft, determined by propellant mass and engine efficiency. Every maneuver from launch to landing consumes part of this finite resource. Mission designers use it to evaluate feasibility, compare trajectories, and size propulsion systems. Once exhausted, a spacecraft cannot perform additional maneuvers without external assistance.}{catSS}
\glsadd{Deorbit}\dictentry{Deorbit}{Deorbiting reduces a spacecraft's orbital velocity to lower its orbit for atmospheric reentry or controlled disposal. It involves firing thrusters against the direction of travel at a precise point. Controlled deorbiting prevents spacecraft from becoming debris and ensures landing in designated areas. The International Space Station has a planned deorbit corridor over the Pacific Ocean.}{catSS}
\glsadd{Equatorial-Orbit}\dictentry{Equatorial Orbit}{An equatorial orbit has zero inclination, passing directly above Earth's equator. Geostationary satellites use equatorial orbits at about 35,786 kilometers to remain stationary relative to the surface. This orientation eliminates north-south station-keeping fuel needs, extending satellite life. Launching into equatorial orbit from high latitudes requires significant plane-change maneuvers.}{catSS}
\glsadd{Habitable-Zone}\dictentry{Habitable Zone}{The habitable zone is the orbital region around a star where liquid water can exist on a planetary surface. Its boundaries depend on stellar luminosity, with more massive stars having wider zones at greater distances. Atmospheric composition and greenhouse effects further modify habitability. Identifying planets in this zone is a primary goal of exoplanet research and life-detection missions.}{catSS}
\glsadd{Keplerian-Elements}\dictentry{Keplerian Elements}{Keplerian elements are six orbital parameters describing an orbit's size, shape, orientation, and the body's position. They include semi-major axis, eccentricity, inclination, right ascension of the ascending node, argument of periapsis, and true anomaly. Named after Johannes Kepler, these elements are standard in satellite tracking, astrodynamics, and mission planning computations.}{catSS}
\glsadd{Orbital-Plane}\dictentry{Orbital Plane}{An orbital plane is the two-dimensional surface defined by an orbiting body's path around its primary. The ecliptic is Earth's orbital plane around the Sun, and satellite inclinations are measured relative to it. Changing orbital planes requires altering the angular momentum vector, demanding significant propulsive effort. Understanding orbital planes is essential for constellation design and intercept trajectories.}{catSS}
\glsadd{Probe}\dictentry{Probe}{A space probe is an unmanned spacecraft exploring regions beyond Earth's orbit, transmitting data to ground stations. Probes may flyby, orbit, or land on targets, with designs tailored to specific scientific objectives. Notable examples include Voyager and New Horizons. Due to vast communication delays, probes must operate with considerable autonomy.}{catSS}
\glsadd{Radiation-Belt}\dictentry{Radiation Belt}{Radiation belts are zones of energetic charged particles trapped by a planet's magnetic field, known as Van Allen belts around Earth. The inner belt contains protons and the outer belt dominates with electrons. These belts pose radiation hazards to electronics and astronaut health, requiring shielding and trajectory planning. Belt intensity varies with solar activity.}{catSS}
\glsadd{Station-Keeping}\dictentry{Station-Keeping}{Station-Keeping performs small orbital corrections to maintain a satellite in its assigned position. Perturbations from drag, gravitational anomalies, solar radiation, and third-body effects cause orbital drift. Geostationary satellites must maintain both east-west position and inclination within narrow limits. Propellant consumed for station-keeping is a major factor in satellite operational lifetime.}{catSS}
\glsadd{Venus-Flyby}\dictentry{Venus Flyby}{A Venus flyby passes close to Venus to alter spacecraft velocity through gravitational assistance. Venus is a common gravity-assist target for outer solar system missions due to its significant mass and orbital velocity. The technique transfers orbital energy from Venus to the spacecraft, reducing propellant needs. Multiple flybys can be chained to reach distant destinations.}{catSS}
\glsadd{Aerocapture}\dictentry{Aerocapture}{Aerocapture uses atmospheric drag to decelerate a spacecraft from hyperbolic approach into a captured orbit. The vehicle enters the atmosphere at high speed, loses energy aerodynamically, and exits into the desired orbit. This saves significant delta-v compared to propulsive insertion, enabling lighter spacecraft. NASA has studied aerocapture for Mars and Titan missions.}{catSS}
\glsadd{Apsis}\dictentry{Apsis}{An apsis is the point in an orbit where a body is closest to or farthest from its primary. Periapsis marks nearest approach and apoapsis the most distant point. The line connecting both is the line of apsides, representing the elliptical orbit's major axis. Gravitational perturbations cause apsidal precession over time.}{catSS}
\glsadd{Argument-of-Periapsis}\dictentry{Argument of Periapsis}{The argument of periapsis is the angle from the ascending node to periapsis measured within the orbital plane in the direction of motion. It defines the orbit's elliptical orientation within its plane and is one of six classical orbital elements. This parameter determines where a satellite reaches closest approach. Earth's J2 oblateness causes it to precess gradually.}{catSS}
\glsadd{Longitude-of-Ascending-Node}\dictentry{Longitude of Ascending Node}{The longitude of the ascending node is the angle from a reference direction to where the orbit crosses the reference plane moving northward. It defines the orbital plane's orientation relative to an inertial frame and is one of six Keplerian elements. Solar radiation pressure and third-body effects cause this angle to shift over time.}{catSS}
\glsadd{Gravitational-Parameter}\dictentry{Gravitational Parameter}{The gravitational parameter is the product of the universal gravitational constant and a central body's mass, denoted mu. It governs gravitational influence strength on orbiting objects throughout orbital mechanics. Using this parameter simplifies calculations because the central body's mass is known more precisely than the gravitational constant. Earth's mu is about 3.986e14 m3/s2.}{catSS}
\glsadd{Two-Body-Problem}\dictentry{Two-Body Problem}{The two-body problem describes the motion of two objects interacting solely through mutual gravitational attraction. Its analytical solution yields conic section orbits following Kepler's laws, including ellipses, parabolas, and hyperbolas. This foundational problem provides the baseline from which perturbations are developed. Real missions must incorporate additional forces like drag and third-body gravity.}{catSS}
\glsadd{Restricted-Three-Body-Problem}\dictentry{Restricted Three-Body Problem}{The restricted three-body problem examines a small mass's motion under two large bodies' gravitational influence, where the small mass doesn't affect the primaries. It reveals five equilibrium Lagrange points where small objects can maintain stable positions. This model guides missions to Lagrange points, Trojan asteroid studies, and stability analysis. No general closed-form solution exists.}{catSS}
\glsadd{Resonance}\dictentry{Resonance}{Orbital resonance occurs when orbiting bodies exert regular gravitational influence through orbital period ratios of small integers. Mean-motion resonances can stabilize or destabilize orbits depending on configuration. Resonances shape planetary systems, clear asteroid belt gaps, and trap Trojan asteroids. Spacecraft can exploit resonance to maintain orbits with minimal fuel.}{catSS}
\glsadd{Secular-Perturbation}\dictentry{Secular Perturbation}{Secular perturbations are long-term, non-periodic changes in orbital elements from gravitational interactions. Unlike periodic variations, they accumulate monotonically, shifting eccentricity, inclination, and orientation gradually. These changes predict long-term orbit evolution for satellites and planetary systems. Lagrange's planetary equations compute secular rates of change analytically.}{catSS}
\glsadd{Periodic-Perturbation}\dictentry{Periodic Perturbation}{Periodic perturbations are oscillatory variations in orbital elements repeating over defined intervals from gravitational and non-gravitational forces. They arise from planetary gravity, solar radiation pressure, and Earth's oblateness. Unlike secular effects, periodic perturbations average to zero over complete cycles. Accurate prediction requires both short-period and long-period periodic terms.}{catSS}
\glsadd{J2-Effect}\dictentry{J2 Effect}{The J2 effect is the dominant perturbation on Earth-orbiting satellites from the planet's equatorial bulge creating oblate gravity. It produces secular changes in the ascending node and argument of periapsis, causing orbital plane precession and ellipse rotation. Magnitude depends on altitude and inclination. J2 is corrected through station-keeping or incorporated into specialized orbit designs.}{catSS}
\glsadd{Precession}\dictentry{Precession}{Precession is the gradual change in orientation of a rotational or orbital axis due to gravitational torques. Earth's axial precession traces a 26,000-year circle, shifting star positions relative to seasons. Orbital precession from perturbations like J2 rotates the major axis. Accurate precession models are essential for long-term astronomical predictions and satellite operations.}{catSS}
\glsadd{Nutation}\dictentry{Nutation}{Nutation is a small periodic oscillation superimposed on the larger precession of a rotating body's axis. Earth's dominant 18.6-year nutation period results from the Moon's gravity on the equatorial bulge. This wobble affects precise astronomical observations and requires correction in telescope pointing. Nutation and precession together define Earth's complete axial motion in space.}{catSS}
\glsadd{Libration}\dictentry{Libration}{Libration is the apparent oscillation of a tidally locked body, allowing observers to see slightly more than half its surface. Lunar librations from the Moon's elliptical orbit and axial tilt expose about 59 percent of its surface. This differs from true rotation and arises because tidal locking imperfectly synchronizes rotation and revolution. Libration measurements reveal internal structure and tidal properties.}{catSS}
\glsadd{Albedo}\dictentry{Albedo}{Albedo is the fraction of incident light reflected by a celestial body, ranging from zero for total absorption to one for perfect reflection. It is critical in climate models, determining how much solar energy a planet absorbs versus reflects. High-albedo bodies like Enceladus reflect most sunlight, while low-albedo asteroids absorb most. Measuring albedo infers surface composition and texture.}{catSS}
\glsadd{Bond-Albedo}\dictentry{Bond Albedo}{Bond albedo is the fraction of total electromagnetic energy reflected by a body in all directions across all wavelengths and phase angles. It differs from geometric albedo by accounting for light scattered across the entire sphere. Bond albedo calculates a planet's energy balance and equilibrium temperature. Earth's Bond albedo is approximately 0.306.}{catSS}
\glsadd{Geometric-Albedo}\dictentry{Geometric Albedo}{Geometric albedo measures a body's brightness at zero phase angle relative to a perfectly diffusing disk of equal cross-section. It quantifies reflectivity when illuminated and observed from the same direction. Unlike Bond albedo, it can exceed one for efficient backscatterers like Enceladus. This quantity is fundamental for characterizing asteroids, moons, and planets photometrically.}{catSS}
\glsadd{Apparent-Magnitude}\dictentry{Apparent Magnitude}{Apparent magnitude quantifies how bright a celestial object appears from Earth, with lower numbers being brighter. The logarithmic scale means a five-magnitude difference equals a hundredfold brightness change. Atmospheric conditions, distance, and intrinsic luminosity all influence apparent magnitude. This ancient Greek classification system remains astronomy's standard brightness measure.}{catSS}
\glsadd{Absolute-Magnitude}\dictentry{Absolute Magnitude}{Absolute magnitude is the apparent magnitude an object would have at a standard distance of ten parsecs. For solar system bodies, it is defined at one astronomical unit from both Sun and Earth at zero phase angle. This standardization enables direct intrinsic brightness comparison between distant objects. It is fundamental for estimating asteroid sizes and stellar luminosities.}{catSS}
\glsadd{Stellar-Parallax}\dictentry{Stellar Parallax}{Stellar parallax is the apparent shift in a star's position against background stars from Earth's orbital motion. Measuring this tiny angular displacement over six months enables direct distance calculation via trigonometry. Parallax forms the cosmic distance ladder's foundation for nearby stars. Hipparcos and Gaia missions have measured parallaxes for millions of stars.}{catSS}
\glsadd{Light-Curve}\dictentry{Light Curve}{A light curve graphs how a celestial object's brightness varies over time. It studies variable stars, exoplanet transits, asteroid rotations, and supernovae. Shape, period, and amplitude reveal physical properties like size, rotation rate, and orbital configuration. Automated surveys generate millions of light curves for large-scale stellar and planetary studies.}{catSS}
\glsadd{Spectroscopy}\dictentry{Spectroscopy}{Spectroscopy analyzes light spectra emitted, absorbed, or scattered by matter to determine composition, temperature, density, and motion. Spectral lines serve as chemical fingerprints for elements in stars and atmospheres. Doppler shifts reveal radial velocities and Zeeman splitting indicates magnetic fields. Spectroscopy is among astrophysics' most powerful analytical tools.}{catSS}
\glsadd{Blueshift}\dictentry{Blueshift}{Blueshift is the displacement of spectral features toward shorter wavelengths, indicating a source is moving toward the observer. It is redshift's opposite and occurs when objects approach at significant light-speed fractions. Cosmological blueshift is rare due to universal expansion, but nearby galaxies like Andromeda exhibit it. It measures approach velocities in local galactic dynamics.}{catSS}
\glsadd{Doppler-Effect}\dictentry{Doppler Effect}{The Doppler Effect is the change in wave frequency from relative motion between source and observer. For light, approaching sources shift spectra blue while receding ones shift red. It measures stellar radial velocities, detects exoplanets, and reveals the universe's expansion rate. Sound waves show the same effect, explaining pitch changes of passing vehicles.}{catSS}
\glsadd{Cosmic-Rays}\dictentry{Cosmic Rays}{Cosmic rays are high-energy charged particles, mostly protons and nuclei, traveling at near light speed. They originate from solar flares, supernovae, and galactic astrophysical processes. Striking Earth's atmosphere, they produce secondary particle showers affecting electronics and posing radiation hazards. Studying them reveals extreme conditions in distant astrophysical environments.}{catSS}
\glsadd{Magnetic-Reconnection}\dictentry{Magnetic Reconnection}{Magnetic reconnection breaks and reconnects oppositely directed magnetic field lines in plasma, releasing stored energy. It drives solar flares, coronal mass ejections, and geomagnetic storms when solar fields interact with Earth's magnetosphere. The process converts magnetic energy into kinetic energy, heat, and particle acceleration. Understanding it is essential for space weather prediction.}{catSS}
\glsadd{Krmn-Line}\dictentry{Kármán Line}{The Kármán Line at one hundred kilometers altitude is the internationally recognized boundary between Earth's atmosphere and outer space. At this height, aerodynamic flight becomes impractical due to insufficient atmospheric density for lift. The Fédération Aéronautique Internationale uses it for astronaut certification. Some agencies, including NASA, recognize eighty kilometers as space's edge.}{catSS}
\glsadd{Coronal-Hole}\dictentry{Coronal Hole}{A coronal hole is a Sun region with open magnetic field lines allowing high-speed solar wind to escape into space. These areas appear darker in X-ray and ultraviolet observations from lower temperature and density. Their solar wind streams can trigger geomagnetic storms at Earth. Coronal holes are most prominent during solar minimum but persist throughout the cycle.}{catSS}
\glsadd{Solar-Maximum}\dictentry{Solar Maximum}{Solar maximum is the peak of the approximately eleven-year solar cycle, with highest sunspot counts, flares, and coronal mass ejections. The Sun's magnetic field is most complex during this phase, driving intense space weather. Increased activity disrupts satellite communications, navigation, and power grids. Maximum typically lasts one to two years before declining.}{catSS}
\glsadd{Solar-Minimum}\dictentry{Solar Minimum}{Solar minimum is the quietest solar cycle period with fewest sunspots and flares. The simpler, more stable magnetic field produces fewer geomagnetic disturbances. However, weakened heliospheric shielding allows more galactic cosmic rays to reach Earth. Minimum typically lasts several years before activity rises toward the next maximum.}{catSS}
\glsadd{Sunspot}\dictentry{Sunspot}{A sunspot is a temporary dark photospheric region where concentrated magnetic fields inhibit convection. At roughly 3,500 Kelvin versus the surrounding 5,800 Kelvin, they appear dark by contrast. Sunspots often form magnetically paired groups and are sites of flares and coronal mass ejections. Their count and distribution indicate solar cycle phase.}{catSS}
\glsadd{Solar-Cycle}\dictentry{Solar Cycle}{The solar cycle is the approximately eleven-year periodic variation in solar magnetic activity, visible in sunspot numbers and flare frequency. The solar dynamo drives it through differential rotation and convection. The cycle transitions between maximum activity and minimum. Understanding it is crucial for predicting space weather impacts on technological systems.}{catSS}
\glsadd{Stellar-Remnant}\dictentry{Stellar Remnant}{A stellar remnant is the dense object left after a star exhausts nuclear fuel and sheds outer layers. Depending on progenitor mass, remnants include white dwarfs, neutron stars, and black holes. They represent stellar evolution's final stage and provide extreme physics laboratories. Studying remnants illuminates end states of stellar evolution and galactic chemical enrichment.}{catSS}
\glsadd{Pulsar}\dictentry{Pulsar}{A pulsar is a rapidly rotating, highly magnetized neutron star emitting electromagnetic beams from its magnetic poles. Rotation sweeps these beams through space, producing periodic pulses toward Earth. Some pulsars serve as extraordinarily precise natural clocks. They study neutron star physics, test general relativity, and detect gravitational waves through timing arrays.}{catSS}
\glsadd{Quasar}\dictentry{Quasar}{A quasar is an extremely luminous active galactic nucleus powered by a supermassive black hole accreting matter rapidly. Its accretion disk radiation can outshine all stars in its host galaxy. Among the most distant observable objects, quasars probe early universe conditions. Their study advanced understanding of galaxy evolution and black hole growth.}{catSS}
\glsadd{Event-Horizon}\dictentry{Event Horizon}{The event horizon is a black hole's boundary beyond which nothing, including light, can escape gravity. It is the point where escape velocity exceeds light speed. Its size is proportional to the black hole's mass, described by the Schwarzschild radius for non-rotating holes. While unobservable directly, its effects on surrounding matter and light are detectable.}{catSS}
\glsadd{Singularity}\dictentry{Singularity}{A singularity is a theoretical point of infinite density and zero volume at a black hole's center where general relativity breaks down. Matter is compressed beyond the theory's limits, requiring quantum gravity for accurate modeling. Singularities also appear in Big Bang cosmology as the universe's initial state. Whether they physically exist or are resolved by quantum effects remains open.}{catSS}
\glsadd{Gamma-Ray-Burst}\dictentry{Gamma-Ray Burst}{A gamma-ray burst is the most luminous electromagnetic event in the universe, observed as extremely energetic cosmic explosions. Long-duration bursts result from massive star collapse, short-duration from neutron star mergers. Most energy releases in gamma rays within seconds, followed by fading multi-wavelength afterglows. These events probe the early universe and interstellar medium.}{catSS}
\glsadd{Gravitational-Wave}\dictentry{Gravitational Wave}{Gravitational waves are spacetime ripples from accelerating massive objects, predicted by general relativity. They propagate at light speed, stretching and compressing space perpendicularly. LIGO's 2015 detection from merging black holes confirmed their existence and opened a new observational window. Gravitational wave astronomy reveals events invisible to electromagnetic telescopes.}{catSS}
\glsadd{Gravitational-Lens}\dictentry{Gravitational Lens}{A gravitational lens is a massive object whose gravity bends light from distant sources, distorting and magnifying images. Galaxy clusters produce multiple images, arcs, or Einstein rings of background objects. This general relativity phenomenon maps dark matter and reveals otherwise invisible distant galaxies. Gravitational lensing is indispensable in observational cosmology.}{catSS}
\glsadd{Exoplanet-Detection}\dictentry{Exoplanet Detection}{Exoplanet detection uses various methods to discover planets orbiting other stars. Techniques include radial velocity, transit photometry, direct imaging, and microlensing, each sensitive to different planet masses and distances. The field has advanced rapidly since the 1992 first discovery, revealing thousands of systems. Next-generation telescopes aim to characterize habitable world atmospheres.}{catSS}
\glsadd{Radial-Velocity-Method}\dictentry{Radial Velocity Method}{The radial velocity method detects exoplanets through periodic Doppler shifts in a star's spectrum from an orbiting planet's gravitational tug. Spectral lines shift blue and red as the star oscillates. Most sensitive to massive close-in planets, it yields minimum mass and orbital period. Combined with transit data, it provides both mass and radius for density calculations.}{catSS}
\glsadd{Transit-Method}\dictentry{Transit Method}{The transit method detects exoplanets by observing periodic star dimming when a planet crosses in front. Dip depth reveals the planet's relative size, while duration and frequency provide orbital parameters. Combined with radial velocity, it yields radius and mass for density determination. Kepler and TESS missions discovered thousands of exoplanets using this approach.}{catSS}
\glsadd{Direct-Imaging}\dictentry{Direct Imaging}{Direct imaging captures photons from an exoplanet directly, separating them from the host star's overwhelming light. It requires coronagraphs, adaptive optics, and space telescopes for necessary contrast ratios. Most effective for young, massive planets in wide infrared-emitting orbits, it provides atmospheric spectroscopic data. The technique is challenging but scientifically valuable.}{catSS}
\glsadd{Habitable-Exoplanet}\dictentry{Habitable Exoplanet}{A habitable exoplanet is a world potentially supporting liquid water on its surface, a key life requirement. Orbital distance, atmosphere, and stellar type determine habitability. Transit and radial velocity surveys identify candidates, but atmospheric characterization confirms true potential. Discovering and studying such planets drives next-generation space telescope design.}{catSS}
\glsadd{Goldilocks-Zone}\dictentry{Goldilocks Zone}{The Goldilocks zone is the orbital region around a star where temperatures allow liquid water on a planetary surface. Named for the fairy tale, boundaries depend on stellar luminosity, atmosphere, and greenhouse effects. Finding planets in this zone around Sun-like stars is a primary exoplanet survey objective. It is synonymous with the habitable zone.}{catSS}
\glsadd{Biosignature}\dictentry{Biosignature}{A biosignature is any measurable substance or pattern providing evidence of life. In astrobiology, biosignatures include atmospheric oxygen-methane disequilibrium, organic molecules, and biologically altered isotopic ratios. Remote spectroscopic detection in exoplanet atmospheres is a future mission goal. Distinguishing true biological signatures from abiotic false positives remains challenging.}{catSS}
\glsadd{Panspermia}\dictentry{Panspermia}{Panspermia hypothesizes life distributes throughout the universe via meteoroids, asteroids, and comets carrying biological material. Earth's life may have originated from space-delivered organic molecules rather than forming locally. Lithopanspermia involves rock fragments, radiopanspermia radiation pressure. Supported by evidence some microorganisms survive space conditions during meteorite transit simulations.}{catSS}
\glsadd{Fermi-Paradox}\dictentry{Fermi Paradox}{The Fermi Paradox highlights the contradiction between high extraterrestrial civilization probability and absent observed evidence. Enrico Fermi asked where all alien civilizations are given the galaxy's vast star and planet numbers. Solutions range from rare Earth to self-destruction or undetectability. The paradox drives ongoing astrobiology and SETI debate.}{catSS}
\glsadd{Drake-Equation}\dictentry{Drake Equation}{The Drake Equation estimates communicative extraterrestrial civilizations in the Milky Way by multiplying star formation rates, planet fractions, habitability, and intelligence probabilities. While poorly constrained, it organizes thinking about extraterrestrial intelligence. Frank Drake formulated it in 1961 for the first SETI conference. It remains a foundational framework for the field.}{catSS}
\glsadd{SETI}\dictentry{SETI}{SETI, the Search for Extraterrestrial Intelligence, detects signals or evidence of technological civilizations beyond Earth. It monitors radio and optical wavelengths for artificial transmissions. Projects range from targeted star observations to distributed computing surveys. Despite decades without confirmed detections, SETI continues with increasingly sensitive instruments and expanded search parameters.}{catSS}
\glsadd{Von-Neumann-Probe}\dictentry{Von Neumann Probe}{A Von Neumann probe is a hypothetical self-replicating spacecraft using interstellar raw materials to manufacture copies. John Von Neumann proposed it could rapidly colonize a galaxy through exponential expansion. Its absence is one Fermi Paradox resolution. Theoretical studies explore engineering challenges, replication limits, and detection signatures of self-replicating systems.}{catSS}
\glsadd{Dyson-Sphere}\dictentry{Dyson Sphere}{A Dyson Sphere is a hypothetical megastructure completely enclosing a star to capture its total energy output. Variants include solid shells, orbiting swarms, and ring structures. Freeman Dyson proposed it as a possible explanation for unusual infrared stellar signatures. Detecting such structures is a goal of technosignature surveys in astronomical data.}{catSS}
\glsadd{Kardashev-Scale}\dictentry{Kardashev Scale}{The Kardashev Scale classifies civilizations by energy-harvesting capability. Type I uses all planetary energy, Type II harnesses their star's total output, Type III an entire galaxy. Nikolai Kardashev proposed it in 1964. The scale provides a framework for discussing advanced extraterrestrial civilizations' capabilities and detectability.}{catSS}
\glsadd{Light-Sail}\dictentry{Light Sail}{A light sail uses radiation pressure from sunlight or lasers to accelerate a large, thin reflective surface for propulsion. It requires extremely low-mass, high-reflectivity sails for meaningful acceleration. Unlike rockets, light sails carry no propellant, enabling potentially extreme speeds over interstellar distances. Japan's IKAROS mission validated the basic principle in space.}{catSS}
\glsadd{Breakthrough-Starshot}\dictentry{Breakthrough Starshot}{Breakthrough Starshot develops laser-accelerated light sail technology for interstellar travel. The concept envisions gram-scale probes reaching Alpha Centauri at light-speed fractions within decades. Key challenges include heat-resistant materials, precision navigation, and interstellar communication. It represents one of humanity's most ambitious interstellar exploration efforts.}{catSS}
\end{multicols}

\clearpage
\printglossary[style=altlist,title=Index]
\end{document}